\documentclass[11pt]{article}

\usepackage[T1]{fontenc}
\usepackage[utf8]{inputenc}
\usepackage{lmodern}
\usepackage{microtype}
\usepackage[margin=1.05in]{geometry}
\usepackage{amsmath,amssymb,amsthm,mathtools}
\usepackage{bm,mathrsfs}
\usepackage{booktabs,longtable,array}
\usepackage{enumitem}
\usepackage{aliascnt}
\usepackage{placeins}
\usepackage{tocloft}
\usepackage{xcolor}
\usepackage{tikz}
\usetikzlibrary{arrows.meta,positioning,fit,calc,shapes.geometric}
\usepackage[numbers,sort&compress]{natbib}
\usepackage[hidelinks,unicode]{hyperref}

\allowdisplaybreaks
\emergencystretch=2em
\setcounter{secnumdepth}{3}
\setcounter{tocdepth}{2}
\setlist{nosep,leftmargin=*}
\setlength{\cftsecnumwidth}{4.5em}
\setlength{\cftsubsecnumwidth}{5.5em}
\numberwithin{section}{part}

\theoremstyle{plain}
\newtheorem{theorem}{Theorem}[section]
\newaliascnt{proposition}{theorem}
\newtheorem{proposition}[proposition]{Proposition}
\aliascntresetthe{proposition}
\newaliascnt{lemma}{theorem}
\newtheorem{lemma}[lemma]{Lemma}
\aliascntresetthe{lemma}
\newaliascnt{corollary}{theorem}
\newtheorem{corollary}[corollary]{Corollary}
\aliascntresetthe{corollary}

\theoremstyle{definition}
\newaliascnt{definition}{theorem}
\newtheorem{definition}[definition]{Definition}
\aliascntresetthe{definition}
\newaliascnt{example}{theorem}
\newtheorem{example}[example]{Example}
\aliascntresetthe{example}

\theoremstyle{remark}
\newaliascnt{remark}{theorem}
\newtheorem{remark}[remark]{Remark}
\aliascntresetthe{remark}

\usepackage[nameinlink,capitalise]{cleveref}
\crefname{part}{Part}{Parts}
\crefname{theorem}{Theorem}{Theorems}
\crefname{proposition}{Proposition}{Propositions}
\crefname{lemma}{Lemma}{Lemmas}
\crefname{corollary}{Corollary}{Corollaries}
\crefname{definition}{Definition}{Definitions}
\crefname{example}{Example}{Examples}
\crefname{remark}{Remark}{Remarks}

\newcommand{\dd}{\,\mathrm d}
\newcommand{\given}{\,\middle|\,}
\newcommand{\Geom}{\mathsf{Geom}}
\newcommand{\Caus}{\mathsf{Caus}}
\newcommand{\Abs}{\mathsf{Abs}}
\newcommand{\Lift}{\mathsf{Lift}}
\newcommand{\Bdry}{\mathsf{Bdry}}
\newcommand{\Res}{J_{\mathrm{res}}}
\newcommand{\Channels}{\mathfrak C}
\newcommand{\Modes}{\mathfrak M}
\newcommand{\Anc}{\operatorname{Anc}}
\newcommand{\Src}{\operatorname{Src}}
\newcommand{\Reach}{\operatorname{Reach}}
\newcommand{\Price}{\operatorname{Price}}
\newcommand{\Align}{\operatorname{Align}}
\newcommand{\Null}{\mathscr N}
\newcommand{\Orbit}{\mathscr O}
\newcommand{\Records}{\mathscr R}
\newcommand{\Spine}{\mathscr T}
\newcommand{\Phys}{\mathrm{phys}}
\newcommand{\Sat}{\mathrm{sat}}
\newcommand{\Budget}{B_{\Phys}}
\newcommand{\CausD}{\mathfrak D_t^{\Caus}}
\newcommand{\Energy}{\mathcal E}
\newcommand{\Action}{\mathcal A}
\newcommand{\Production}{\mathfrak D}
\newcommand{\Entropy}{\mathcal M}
\newcommand{\Hspace}{\mathcal H}
\newcommand{\Vspace}{\mathcal V}
\newcommand{\Capacity}{\operatorname{Cap}}
\newcommand{\dist}{\operatorname{dist}}
\newcommand{\Ran}{\operatorname{Ran}}
\DeclareMathOperator{\supp}{supp}

\definecolor{fwInk}{HTML}{2F3E46}
\definecolor{fwBlue}{HTML}{4C78A8}
\definecolor{fwTeal}{HTML}{4C9F92}
\definecolor{fwAmber}{HTML}{C58B2A}
\definecolor{fwCoral}{HTML}{B9655A}
\definecolor{fwViolet}{HTML}{8064A2}
\definecolor{fwGray}{HTML}{E9ECEF}
\tikzset{
  fw diagram/.style={>=Stealth,font=\small,line width=0.75pt},
  fw box/.style={draw=fwInk,rounded corners=2pt,align=center,
    minimum height=7mm,inner xsep=5pt,inner ysep=3pt,fill=white},
  fw source/.style={fw box,draw=fwBlue,fill=fwBlue!13},
  fw geom/.style={fw box,draw=fwTeal,fill=fwTeal!14},
  fw use/.style={fw box,draw=fwAmber,fill=fwAmber!14},
  fw provenance/.style={fw box,draw=fwViolet,fill=fwViolet!12},
  fw terminal/.style={fw box,draw=fwCoral,fill=fwCoral!12},
  fw residual/.style={fw box,draw=fwInk!65,fill=fwGray},
  fw arrow/.style={-{Stealth[length=2.2mm,width=1.5mm]},draw=fwInk},
  fw dashed arrow/.style={fw arrow,dashed},
  box/.style={rectangle,rounded corners,draw=fwInk,fill=white,
    align=center,inner sep=3pt,text width=2.65cm,font=\scriptsize},
  dec/.style={diamond,draw=fwAmber,fill=fwAmber!14,aspect=2.35,
    align=center,inner sep=1.5pt,text width=2.25cm,font=\scriptsize},
  term/.style={ellipse,draw=fwCoral,fill=fwCoral!12,
    align=center,inner sep=3pt,text width=2.7cm,font=\scriptsize},
  route/.style={ellipse,draw=fwInk!65,dashed,fill=fwGray,
    align=center,inner sep=2pt,text width=2.7cm,font=\scriptsize},
  arrow/.style={-{Latex[length=2mm]},thick,draw=fwInk}
}

\title{Continuum Structural Complexity:\\
Target-Aligned Source Algebra, Physical Expenditure, and Singular Reachability}
\author{Guillem Duran Ballester\\
\href{mailto:guillem@fragile.tech}{\texttt{guillem@fragile.tech}}}
\date{}

\hypersetup{
  pdftitle={Continuum Structural Complexity: Target-Aligned Source Algebra, Physical Expenditure, and Singular Reachability},
  pdfauthor={Guillem Duran Ballester}
}

\begin{document}
\maketitle

\begin{abstract}
A public continuum presentation consists of a weak PDE or flow, its datum and
boundary conditions, a rational approximation grammar, its continuation and
native-operation syntax, a canonical weak test core, and the formal physical
carrier identity.  No target-relative estimate, compactness statement,
rigidity theorem, or continuation conclusion is input.  The presentation
generates an Archimedean ordered cylinder algebra.  Its state space contains
ordinary solutions, every approximation limit, and the first graph value at
which an equation, carrier, chart, trace, or covariance relation fails.

The continuum extension preserves the five source families of Part~I.  They
are \(\Geom,\Caus,\Abs,\Lift\), and \(\Bdry\).  Limits, concentrations, oscillations,
defect measures, tails, traces, and blow-up germs retain the ancestry of the
records that generate them.  A primitive-role theorem partitions every
current path of the public weak expression graph into canonical bulk
form/directed, quotient, lift, and boundary roles.  The free totalized
constructor algebra has a unique provenance morphism into the five-letter
support semilattice.  Thus every target-active atomic current occurrence
belongs to one of the \(2^5-1=31\) possible nonempty provenance cells, while
a total current retains their joint vector, including every nonordinary graph
value and sourcewise limit.  The residual is taken only after this sourcewise
completion and is exactly the common kernel of the prospective target
functionals.  A target-visible residual source is therefore impossible.
Finite-cylinder target shells carry an exact derived Stieltjes-current
identity.  Contact forces unit shell progress, hence a positive contribution
from at least one inherited source on every sufficiently small pre-contact
shell.  A canonical disjoint cascade then converts finite total physical
expenditure into either divergent source price, which excludes contact, or a
zero-price target-aligned successor path.

A distinguished nonreversible subfamily applies a Wick transform to the
self-adjoint \(\Abs\) quotient field of an existing Part~I compensated
record and leaves its real \(\Geom\) and authorized \(\Caus\) fields
unchanged.  The transformed causal
bridge produces an exact Feshbach self-energy, a Mori--Zwanzig memory kernel,
and a completed-square structural action.  Its positive response kernel
transfers the effective-dimension, Rademacher, PAC--Bayes, sequential, control,
and continuum-learning bounds of Part~I.  Empirical trace terms use the
finite Part~I evaluation carriers; the population effective dimension is the
extended supremum of those finite compressions.  Structure-preserving likelihood,
inverse-response kernel ridge, and action-safe fitted Bellman/Gibbs losses
train the dynamics, unrestricted bridge, and RL branches, respectively.  The
inverse-response norm is the exact minimum energy of a latent response lift.
The Part~I exact-cell Rademacher union and learned-operator cover, or the
cell-mixture PAC--Bayes divergence and coefficient-selection charge, select
the model family without an additional simplicity criterion.  The learner
restricts the Part~I saturated usefulness profile to records at a requested
risk tolerance, then returns the Pareto-minimal points of its strict inverse
at the inherited visibility threshold.  Minimum-capacity covering readout
remains an inner coordinate within each fixed geometry; it does not replace
saturated-profile selection.

The remaining path problem is evaluated by an explicit ordered calculus.
Every identity is first restricted to the same-origin target/source stratum
that invoked it and decomposed into signed five-channel remainders, carrier
coboundary, admissible physical production, exterior input, target-null
current, and total-graph defects.  No equation-wide identity has an
independent reachability consequence.
For each source support and rational target floor, a quadratic module contains
the weak equation, signed carrier balance, total graph relations, same-origin
successor relations, nonnegative physical production, and shift
coboundaries.  Exactly one algebraic alternative holds: either a finite
rational identity represents \(-1\) in the ordered cone and excludes that
target-aligned cell, or an Archimedean separation theorem produces a positive
shift-invariant functional satisfying all retained relations.  A nested
cylinder-moment hierarchy enumerates the finite identities and reconstructs
the positive functional from compatible finite moments.  Hamiltonian
resolvents, signed carrier cohomology, entropy cocycles, commutator squares,
and compactness--rigidity formulae are consequence generators inside this
calculus, not additional hypotheses.  A complementary positive functional is
an equality successor: it re-enters the target-cyclic compiler generated by
the selected shell covector.  The compiler restricts every operation to the
same-origin target cell.  Projection-first saturation proves that quotienting
by the common target kernel commutes with every constructor; a failed descent
becomes a projected graph defect.  The resulting target-cyclic spine
assembles one feeding--storage current before a
canonical zero/signed-dyadic/graph-boundary polarity split.  A new coordinate
increases resolved-prefix rank; an equivalent positive cell is followed by
its induced first-return law.  On each recurrent cell, a primal--dual
circulation hierarchy computes the sharp physical price of unit target
return.  Positive price produces a carrier-potential identity; zero price
reconstructs the invariant circulation that generates the next target-cyclic
refinement.  Adjoint transport, constraint descent, defect
polarization, commutator production, and covariant scale valuation are
internal operations of this one queue.  Its contact-seeded outputs are
divergent physical charge, target-nullity, or a joint-saturated five-source
witness for a nonzero saturated prospective source cell.  Every
cofinal refinement is cumulatively joined, closed under mixed operations,
and reprojected before this last output.  Origin polarity prevents that
hypothetical witness from being recycled as singular existence.  An
independently generated ordinary contacting orbit is a separate constructive
output.  A formal model,
equality kernel, failed analytic property, or equation-wide estimate is an
internal coordinate rather than a terminal.
\end{abstract}

\tableofcontents

\section*{Introduction}
\addcontentsline{toc}{section}{Introduction}

Let \(u\) be an ordinary maximal branch of a fixed continuum evolution and let
\(\Sigma\) be a target in its state or germ space.  The central question is
whether the equation can reach \(\Sigma\) while spending only a finite amount
of the physical quantity controlled by its evolution law.  Each orbit carries
one nonnegative countably additive expenditure measure \(\mu_u\) on a
generated event space \(\mathfrak Q_u\).  The finite-budget stratum is
\[
  \mu_u(\mathfrak Q_u)\le \Budget<\infty.
\]
The target-blind public multiplier grammar generates candidate energy,
action, entropy, dissipation, forcing-work, and boundary-flux balances.  Its
finite checker admits a carrier stratum only after verifying the exact signed
balance, positivity, lower bound, localizability, and exterior-work rows.  On
one admitted stratum the mutually compatible positive variations constitute
one measure.  Signed internal
currents are glued first, and only carrier-dominated target-local production
may spend it.  The admitted carrier identity then bounds its total mass before
the ceiling \(\Budget\) is set.  No carrier or favorable estimate is supplied
by the datum or chosen after the target is inspected.

The source alphabet, ancestry discipline, saturation rule, and
target-relative residual extend the structural algebra of Part~I
\citep{DuranBallesterStructuralComplexityI}.  The continuum construction
below derives its physical budget, source-preserving limit closure,
prospective contact identity, signed carrier calculus, successor recurrence,
Hamiltonian formulation, and saturation equations within this paper.  The
remaining citations identify standard analytic results used at explicitly
named steps.  These references and Part~I constitute the complete logical
dependency set.

The structural source alphabet records how target-directed behavior is
generated.  The native dissipative or reversible mechanism is \(\Geom\); the
fixed directed or nonreversible transport is \(\Caus\); quotient and symmetry
reduction are \(\Abs\); charts, renormalizations, and augmented state
variables are \(\Lift\); and boundary, terminal, killing, and exterior-flux
records are \(\Bdry\).  A composite continuum object may carry several
labels, but every elementary source belongs to this alphabet.

The completion is sourcewise.  A limit object is placed in the closure of the
cell that generated it.  Thus a concentration measure retains the ancestry of
the concentrating quantity, a blow-up profile retains the ancestry of the
physical orbit together with its \(\Abs/\Lift\) chart, and a terminal tail
trace acquires \(\Bdry\) ancestry.  Every source cell has an exact structural
row value; its complementary value remains in the same cell and does not
create an additional continuum source.

The target-relative residual is correspondingly strict.  Let
\(\Null_\Sigma\) be the common kernel of the target shell derivatives, target
fluxes, and other declared prospective contact functionals.  The residual
coordinate is contained in \(\Null_\Sigma\).  Hence
\[
 \boxed{
 \text{target-visible}\Longrightarrow
 \text{ancestry in }\{\Geom,\Caus,\Abs,\Lift,\Bdry\};
 \qquad
 \Res\Longrightarrow\text{target-null}.}
\]
Every target-visible limit therefore retains ancestry in one of the five
source families.

The central target-relative question is therefore not whether the entire
solution is large, but which source current contributes positively to the
formation of the singularity.  For a normalized shell crossing \(I_r\), set
\[
 Q_{c,r}:=\int_{I_r}\langle D\Phi_r(u_t),J_t^c\rangle\,dt.
\]
The compiled current identity gives the exact pre-contact identity on the
finite-value ordinary shell branch
\[
 1=\sum_{c\in\Channels}Q_{c,r},
 \qquad
 \sum_{c\in\Channels}(Q_{c,r})_+\ge1,
 \qquad
 \max_{c\in\Channels}(Q_{c,r})_+\ge\frac15.
\]
On an infinite-variation value, the joint shell graph instead retains a
source owner \(o\) with \(\overline Q^o(1)\ge1/5\), without forming an
extended signed sum.  The residual source term vanishes.  Thus contact itself supplies the unit
progress and proves the existence of target-aligned structure.  If all source
currents vanish in the target-active quotient, their magnitude is irrelevant:
every shell observable is constant along the orbit, so contact is excluded
without compactness, a physical-price gap, or a source-transverse estimate.

The same accounting is stable under every generated continuum composition,
not by separate estimates for time, scale, nonlinearities, or limits, but by
the expanded-derivation rule inherited from Part~I.  Before a residual is
formed, the complete derivation DAG of the entire same-origin prefix is
joined and closed under all mixed public operations.  Every physical
quantity remains attached to its owner and domination graph.  Normalized
values retain their numerator, denominator, boundary cell, and unnormalized
target cocycle.  The target projection on these increasing
joint algebras has orthogonal increments \(d_n\), so
\[
 \|P_\infty y_\Phi\|_2^2=\sum_n\|d_n\|_2^2.
\]
For actual crossings, the canonical source-current frontier gives a
two-record contrast of squared norm \(1/4\).  The finite occurrence direct
sum therefore satisfies the exact identities
\[
 \frac N4=\sum_{j<N}\sum_\ell\Delta_{j,\ell},
 \qquad
 N=\sum_{j<N}\sum_\ell\widehat\Delta_{j,\ell}.
\]
No increment is discarded for being small.  The same \(N\) crossings retain
the unnormalized cocycle
\[
 N=\sum_{c\in\Channels}Q^c_{\Phi,N},
 \qquad
 \max_c(Q^c_{\Phi,N})_+\ge N/5.
\]
Thus a countable accumulation of locally target-null-looking operations is
either genuinely null after joint saturation or reappears as cumulative
target-aligned structure with one of the same five ancestries.  It cannot
become a target-bearing residual at the limit.

This is the root no-emergence theorem.  Scale collapse, long-time
accumulation, weak defects, oscillation, concentration, recurrent averaging,
tail escape, boundary passage, and nonlinear cross terms are constructor
instances, not separate loopholes.  Individually regular inputs may have a
target-aligned joint effect, just as individually uninformative observables
may have informative joint fibers in Part~I; the joint effect is retained as
an existing source cell and cannot be declared free merely because it was
absent from each input separately.

The aligned branch is simplified before any size estimate is attempted.  A
target-null or gauge direction is passed to the generated quotient and never
reconsidered.  Projection-first saturation generates every descendant inside
the target-cyclic spine and retains each failed descent as a graph coordinate.
For every surviving total-graph value the algebra forms its
positive target-current measure, compact activity lift, and witness-indexed
successors.  Every subsequent analytic identity is restricted to that cell
and compiled into \eqref{eq:target-local-fact-normal-form} before it can close
or charge the branch.  The topological sets
\(Z_{n+1}=\mathcal V_\eta(Z_n)\) describe whether an occurrence can sustain
one more target-retaining zero-price step.  Proof is supplied separately by
the ordered cone: a checked finite identity
\(-1\in\mathcal K_{\mathfrak b}\) excludes the active cell, while the complementary
order-unit functional is a stationary equality successor satisfying every
generated relation at that depth.  Its first unprocessed local coordinate is
split into zero, signed dyadic, and graph-boundary cells; finite local
identities prune target-null cells, while positive-moment cells re-enter
source selection.  New coordinates increase resolved-prefix rank; recurrent
equivalent cells use their induced return maps.  The target-aligned recurrent
closure compiler performs all of these transitions internally and returns
on a contact-seeded execution only divergent physical charge, a target-null
path, or a saturated prospective source-profile witness.  An independently
generated ordinary orbit may separately realize the same structural atom and
thereby establish a singular mechanism.  An infinite rank chain is first joined
with its complete transcript; every unprocessed joint increment re-enters
the five-source compiler.  Semantic emptiness is never used as a proof object.

The main expenditure argument then uses the same measure \(\mu_u\).  A
singular contact compiler produces infinitely many pairwise disjoint
pre-contact events and selects a persistent aligned source cell.  A
divergent-price value contradicts countable additivity and
\(\mu_u(\mathfrak Q_u)\le\Budget\).  A finite-price value produces
vanishing-price windows and activates the ordered successor calculus.  It is
therefore a reduction step, not a terminal singularity class.
\[
 \boxed{
 \begin{gathered}
 \text{contact}
 \xLongrightarrow{\text{normalized shell}}
 \text{target-aligned source},\\
 \text{target-aligned source}
 \xLongrightarrow{\text{one physical measure}}
 \begin{cases}
  \text{divergent expenditure},\\
  \text{zero-price successor path},
 \end{cases}\\
 \text{first branch contradicts }\Budget;\qquad
 Z_{n+1}=\mathcal V_\eta(Z_n)\text{ evaluates the second};\\
 \text{a nonpruned fixed atom remains prospective.}
 \end{gathered}}
\]

Hamiltonian, resolvent, commutator, constraint-tail, scale, and entropy
identities are functorial refinements generated when their typed expressions
occur in the public equation.  The symmetric form yields a spectral
Hamiltonian and resistance; the full fixed-Caus generator retains its
nonnormal current; signed carrier gluing cancels internal transport and
exposes scale-feed cohomology; and a covariant entropy cocycle telescopes
production across changing charts.  A failed identity is lifted as a
successor with the same ancestry.  These modules operate on the five source
families; they neither enlarge the public presentation nor alter the source algebra.

The Wick-rotated compensated-quotient module treats a self-adjoint
\(\Abs\) generator in spectral time and a real \(\Geom\) generator in
dissipative time.  The actual selected \(\Caus\) bridge determines both the
Feshbach response and the structural action.  The resulting positive response
kernel imports the learning theory of Part~I without diagonalizing a
nonnormal generator, while free, hybrid, and imposed authority remain
distinct optimization domains.  The unrestricted branch optimizes spectral
kernel ridge risk on the zero-waste bridge fiber; the RL branch uses a frozen
Bellman target, safe response compression, and a Gibbs actor calibrated to
the same action-valued critic.

\paragraph{Scope discipline.}
The paper uses exhaustive quantification only to construct the closed
five-source algebra: enumeration of tests, total graphs, sourcewise closures,
and the common target kernel.  Such a construction statement has no
reachability conclusion.  An analytic identity becomes a reachability fact
only after contact has selected a same-origin active branch
\(\mathfrak b\), the identity has been restricted to
\(X_{\mathfrak b}\), and its terms have been placed in the target-local
normal form \eqref{eq:target-local-fact-normal-form}.  A physical consequence
additionally requires the charging rule
\eqref{eq:target-aligned-charge-rule}.  Thus full-algebra quantifiers provide
source exhaustion, whereas every exclusion, price, successor, and
continuation inference is target-local and source-resolved.  This distinction
is formalized in
\cref{thm:target-local-fact-factorization,cor:no-unscoped-reachability-inference}.


\part{Closed Continuum Structural Algebra}
\label{part:closed-algebra}

\section{Physical continuum presentations}
\label{sec:physical-presentations}

\begin{definition}[Canonical physical continuum presentation]
\label{def:physical-presentation}
A canonical physical continuum presentation is the public tuple
\[
 \mathfrak P_{\Phys}
 =\bigl(\mathsf{Eq},u_0,\mathcal D,\mathsf{Appr},\mathsf{Cont},
        \mathsf{Op},\mathcal K\bigr)
\]
with the following data, all transcribed directly from a public textbook or
standard reference statement of the initial--boundary value problem.  A
nonliteral entry below is admissible only when it is target-blind and already
part of the standard local theory of the displayed equation; establishing a
new estimate or theorem is not part of the specialization datum.
\begin{enumerate}[label=\textup{(P\arabic*)}]
\item \(\mathsf{Eq}\) is the displayed differential equation with its domain,
coefficients, boundary conditions, and declared regularity class;
\(u_0\) is its datum; and
\(\mathcal D=\{\varphi_n\}_{n\ge1}\) is the canonical rational smooth test
core.  A standard local solution theorem may identify an ordinary component
of the generated solution graph, but it is not needed to define that graph.
Any local theorem used for this identification is target-blind in the sense
of \cref{def:target-blind-public-input}.
\item The \emph{total local-solution graph} contains every maximal path
\(u:[0,T_*(u))\to X_{\rm reg}\) satisfying \(\mathsf{Eq}\) against every
\(\varphi_n\), with the prescribed datum and physical boundary data.  If the
fiber over \(u_0\) is empty, its value is the explicit record
\(\bot_{\rm loc}\).  Nonuniqueness retains all maximal paths.  Consequently an
empty solution fiber is a form defect and can never be read as target
avoidance.  The approximation graph below independently tests whether the
ordinary fiber is nonempty.
\item \(\mathsf{Appr}\) is the least countable grammar generated by rational
Galerkin projections, rational mollifiers, rational implicit time steps,
rational domain exhaustions, and the typed gauge or boundary operations in
\(\mathsf{Op}\).  A grammar word is retained with its finite-dimensional or
finite-step residual relation against the first finitely many tests and with
the signed carrier identity obtained by applying the same word to
\(\mathcal K\).  Every compatible grammar word is enumerated; no convergence
property or preferred scheme belongs to the presentation.
\item \(\mathsf{Cont}\) is the syntactic extension relation in the declared
regularity class: two local paths are related when their restrictions agree.
No continuation theorem is supplied.  Its ordinary values are actual regular
extensions and its graph-boundary values are terminal germs.  A terminal germ
is stored together with its orbit-origin tag and the cofinal family of
restrictions of that same orbit which generates it.  In the completed
restriction-record space, each member of this cofinal family carries the
same origin and terminal-status tags; by definition the family converges to
its terminal germ.  These tags belong to the retrospective proof record and
do not alter the PDE current on the preterminal interval.  Thus the continuation
graph is closed fiberwise over the orbit origin: two distinct orbits are not
identified merely because some evaluation coordinates agree.  The singular
target is generated from this same-origin graph by
\cref{def:singular-target}; it is not an independently chosen favorable
subset.
\item \(\mathsf{Op}=\{O_n\}_{n\ge1}\) is a countable typed list of the
restriction, rational truncation, quotient, chart, augmentation, trace, and
boundary operations syntactically available for \(\mathsf{Eq}\).  An operation
acts as a symmetry or gauge only on the ordinary value of its total covariance
graph, where it preserves the weak equation, the carrier, and the continuation
target.  Thus the generated gauge family is the equalizer of these three
actions; it is not declared independently.
\item
\[
 \Channels=\{\Geom,\Caus,\Abs,\Lift,\Bdry\}
\]
is the inherited source alphabet of \cref{def:source-alphabet}.  Native terms
are routed by the sector/gradient compiler and subsequent orbit-record
operations acquire their source words by the constructor grammar; this is
proved in \cref{thm:continuum-normal-form-compilation}.
\item \(\mathcal K\) is the target-blind public multiplier grammar generated
from the displayed equation.  It enumerates the equation's native energy,
entropy, action, Hamiltonian, and other Lyapunov balance candidates before the
continuation target is generated.  Testing \(\mathsf{Eq}\), integration by
parts, summing with the native signs, and cancelling common internal faces
generates each total signed carrier relation.  A candidate becomes a budget
carrier only when the finite public checker verifies its exact balance,
nonnegative production, lower bound, localizability, and legitimate exterior
work decomposition.  These are graph values, not presentation data.  No
datum or user field names the successful carrier or supplies any favorable
estimate for it.
\item Let \(\mathcal A_{\mathbb Q}(\mathfrak P)\) be the unital commutative
rational cylinder algebra generated by bounded truncations of all weak
pairings, continuation coordinates, approximation residuals, carrier
coordinates, and input--output coordinates of the typed operations in
\textup{(P1)}--\textup{(P7)}.  For each typed operation it also contains
graph-defect, oscillation, concentration, and exterior-tail coordinates.
Let \(M_{\mathfrak P}\) be the quadratic module generated by squares, the
bounds \(1-e_n^2\), nonnegative defect masses, carrier-production atoms, and
the positive and negative forms of every exact rational cylinder identity.
The graph-complete state space is
\begin{equation}
 \label{eq:graph-complete-state-space}
 X_{\rm gc}:=
 \left\{L:\mathcal A_{\mathbb Q}(\mathfrak P)\to\mathbb R:
 L(1)=1,\quad L(M_{\mathfrak P})\subseteq[0,\infty)\right\}.
\end{equation}
It carries pointwise convergence on the countable cylinder algebra.
Ordinary states and germs embed by evaluation.  The notation \(X_{\rm ev}\)
below means \(X_{\rm gc}\), and
\begin{equation}
 \label{eq:canonical-evaluation-metric}
 d_{\rm ev}(L_1,L_2)
 :=\sum_{n\ge1}2^{-n}|L_1(e_n)-L_2(e_n)|
\end{equation}
for a fixed enumeration of normalized generators and rational products.
\end{enumerate}
The presentation is therefore fixed by public equation data and a formal
multiplier calculation.  The specialization step consists only of recording
these statements and their public references; it contains no new mathematical
argument.  It may not contain a compactness theorem, coercive
estimate, price lower bound, Liouville statement, target-exclusion lemma,
special blow-up theorem, or exact invariant value.  Such
an insertion changes the equation signature and is rejected by the compiler.
\end{definition}

\begin{definition}[Effective canonical continuum presentation]
\label{def:effective-canonical-continuum-presentation}
An \emph{effective canonical continuum presentation} is a canonical physical
presentation \(\mathfrak P_{\Phys}\) together with a finite code
\(\mathsf{Code}(\mathfrak P_{\Phys})\) in the following typed instruction
language.
\begin{enumerate}[label=\textup{(E\arabic*)}]
\item The rational weak test core and the approximation words have total
enumerators.  At a finite approximation word the equation, datum, boundary,
and carrier residuals against the first \(m\) tests are rational intervals
computed to any prescribed rational accuracy.
\item Each constructor in \(\mathsf{Op}\), and each total graph used by that
constructor, has an enumerator of nested rational finite-cylinder graph
enclosures.  The enclosure diameter is itself a public nonnegative graph
coordinate.  If the diameters have a computable decay modulus, the graph has
its ordinary effective value.  Otherwise the first persistent diameter,
oscillation, concentration, or boundary value is retained by the total graph
and becomes target-null or the next target-aligned successor after cumulative
joining and projection.
\item The displayed carrier derivation is a finite typed transcript.  Its
instructions are substitution from \(\mathsf{Eq}\), rational linear
combination, the declared product or chain rule, signed integration by parts,
and cancellation of a pair of equal oppositely oriented internal faces.
Every positivity assertion is a separate rational quadratic-module row.
\item The continuation coordinate has a decidable finite-prefix type in
\[
 \{\mathsf{ordinary},\ \mathsf{noncontinuation},\
   \ \mathsf{graph\text{-}boundary}\}.
\]
Origin and restriction maps are computable on codes and preserve their tags.
The normalized target-shell coordinates have nested rational cylinder
enclosures.  Their ordinary shrinking value carries the modulus generated by
the enclosure stream; every nonshrinking value is retained as a shell-graph
boundary coordinate.
\item Every normalization denominator is represented by its typed total
graph.  Rational generator bounds, finite graph residuals, and every line of
the carrier transcript are decidable from their finite codes.
\end{enumerate}
The instruction language contains only total primitive constructors.  An
enumerator is represented by a program in this language, rather than by an
arbitrary partial program.  A finite package which fails a type check, omits
a required enclosure enumerator, or has an invalid transcript is denoted
\(\mathsf{Reject}(\mathsf{Code})\).  Such a package lies outside the typed
presentation syntax; rejection is not a regularity conclusion.  Failure of
an enclosure stream to furnish a favorable modulus is never a rejection.
\end{definition}

\begin{theorem}[Effective canonical-presentation theorem]
\label{thm:effective-canonical-presentation}
There is a terminating verifier \(\mathsf{Check}_{\rm pres}\) for the finite
codes of \cref{def:effective-canonical-continuum-presentation}.  It returns
exactly one of
\[
 \mathsf{Accept}(\mathfrak P_{\rm eff})
 \qquad\text{or}\qquad
 \mathsf{Reject}(\mathsf{Code}).
\]
For an accepted code, the weak test core, approximation grammar, constructor
graphs, continuation-prefix rows, and target-shell cylinders are uniformly
computably enumerable.  Moreover, every element of the cylinder algebra,
every defining row of \(M_{\mathfrak P}\), and every later coordinate obtained
by a finite composition of public constructors has nested rational
finite-cylinder enclosures.  Exactly one of two generated values occurs: the
enclosures shrink with a computable composed modulus, or their least
nonordinary diameter/oscillation/boundary coordinate is retained by the total
graph.  After restriction to an affordable target-aligned cell, the latter is
target-null or a typed five-source successor.  The accepted carrier identity
has a finite checkable derivation from the displayed weak equation and its
declared boundary data.
\end{theorem}

\begin{proof}
Parse the finite code and type-check its instruction tree.  Totality of an
enumerator is syntactic because the instruction language permits composition
and primitive recursion only over finite words and natural-number indices.
Evaluate each line of the carrier transcript in order.  Substitution and
rational linear combination are checked symbolically; the product, chain,
and integration-by-parts rules are checked against the declared types; and an
internal-face cancellation is accepted only when its two boundary words and
coefficients agree with opposite orientation.  Check separately that every
claimed nonnegative row is an explicit sum-of-squares or a declared positive
carrier atom.  These are finite computations, so the verifier terminates and
rejects at the first failed check.

For an accepted code, the enumerability assertions in
\textup{(E1)} and \textup{(E4)} follow by running the corresponding total
enumerators.  For a typed constructor, compose its nested graph enclosures
with those of its inputs.  Addition adds enclosure radii, multiplication uses
the recorded rational cylinder bounds, and composition propagates the
declared graph enclosures.  When all radii shrink, these formulas compute the
composed modulus.  Otherwise the fixed enumeration selects the least
persistent diameter, oscillation, concentration, or boundary coordinate.
Quotients and other potentially partial
operations are treated through their total graphs, so no reciprocal is
formed outside an ordinary denominator row.  Structural induction on a
finite constructor word therefore produces both a rational cylinder
enclosure and either its error modulus or its least total-graph failure
coordinate.  The same induction enumerates the
generators and relations of \(\mathcal A_{\mathbb Q}(\mathfrak P)\) and
\(M_{\mathfrak P}\).  Expanded-derivation no-creation and universal
provenance route every nonzero target projection of a failure coordinate.
The transcript check proves the final assertion.
\end{proof}

\begin{remark}[Scope of presentation syntax]
\label{rem:effective-presentation-scope}
The theorem verifies a finite transcription in the total instruction
language.  It does not quantify over arbitrary programs or arbitrary
analytic objects.  A public enclosure enumerator is syntactic data; favorable
convergence is not.  Every failure of effective graph or shell convergence
is retained as a generated nonordinary coordinate and is evaluated only
after restriction to the public target spine.
\end{remark}

\begin{theorem}[Graph-complete compactness and separation]
\label{thm:graph-complete-compactness}
The module \(M_{\mathfrak P}\) is Archimedean.  Its state space
\(X_{\rm gc}\) is compact and metrizable.  The cylinder coordinates separate
its points, and every ordinary state, germ, orbit segment, or approximation
record defines an evaluation state in \(X_{\rm gc}\).  If
\(X_{\rm gc}=\varnothing\), then \(-1\in M_{\mathfrak P}\); this finite
algebraic inconsistency is the total form-defect value.
\end{theorem}

\begin{proof}
Every cylinder polynomial contains only finitely many normalized generators.
The relations \(1-e_n^2\in M_{\mathfrak P}\), together with closure of a
quadratic module under addition and multiplication by squares, give
\(N\pm a\in M_{\mathfrak P}\) for some rational \(N\) and every
\(a\in\mathcal A_{\mathbb Q}(\mathfrak P)\).  Thus the module is
Archimedean.  Positivity gives
\(|L(a)|\le N\) whenever \(N\pm a\in M_{\mathfrak P}\).  The state space is
therefore a closed subset of a countable product of compact intervals; it is
compact and metrizable.  Pointwise equality on the cylinder algebra is
equality of states.  Ordinary records satisfy every defining graph and
positivity relation, so evaluation is positive and normalized.

For the final assertion, the Archimedean separation theorem for quadratic
modules says that a proper Archimedean module admits a normalized positive
linear functional \citep{Schmudgen2017}.  Hence an empty state
space forces \(M_{\mathfrak P}\) to be improper, equivalently
\(-1\in M_{\mathfrak P}\).  Membership uses one finite algebraic expression
and hence only finitely many presentation relations.
\end{proof}

\begin{definition}[Exhaustive approximation tower]
\label{def:exhaustive-approximation-tower}
An approximation record is a tuple
\[
 a=(g,N,h,\mathsf r_a,\mathsf k_a),
\]
where \(g\) is a word of \(\mathsf{Appr}\), \(N\) is the number of tested
weak coordinates, \(h\in\mathbb Q_+\) is its resolution,
\(\mathsf r_a\) is the vector of equation and operation-graph residuals, and
\(\mathsf k_a\) is its exact signed carrier record.  Write \(a'\preceq a\)
when \(a'\) is obtained from \(a\) by forgetting tests, coarsening resolution,
or restricting its time or spatial domain.  The exhaustive approximation
tower is the countable inverse graph
\[
 \mathfrak A(\mathfrak P)
 :=\{a:\text{\(a\) is generated by }\mathsf{Appr}\}
\]
with all maps \(a\mapsto a'\) for \(a'\preceq a\).  A compatible tower path
is \emph{zero defect} when every fixed residual coordinate tends to zero and
the signed carrier coordinates converge in the graph-complete topology.
For each enumerated normalized residual \(r_k\in[-1,1]\), the grammar also contains the monotone
tail coordinate
\[
 d_{k,N}:=\sup\{|r_k(a)|:a\text{ occurs after refinement level }N\},
 \qquad d_k:=\inf_N d_{k,N}.
\]
Suprema are totalized through bounded rational upper-envelope coordinates.
Thus \(d_k=\limsup|r_k|\), and failure of convergence to zero is a nonzero
closed graph coordinate rather than an unrecorded oscillation.
\end{definition}

\begin{theorem}[Approximation relaxation and first failure]
\label{thm:approximation-realization-first-failure}
Every zero-defect compatible tower path has a graph-complete positive-functional
limit satisfying the public weak equation and native carrier relation in every
rational test coordinate.  This is a relaxed equation state; it is an
ordinary solution only after the Dirac-realization gate below.  Conversely,
every ordinary weak solution embeds in
\(X_{\rm gc}\) by evaluation and is a zero-defect value of the total weak
solution graph.  For any compatible tower path which is not zero defect,
there is a least operation index whose residual, oscillation,
concentration, tail, or carrier-defect coordinate is nonzero.  That coordinate
is the first-failure record of the path.
\end{theorem}

\begin{proof}
Compactness gives a convergent subnet, and metrizability gives a convergent
subsequence.  Fix a rational test coordinate.  Compatibility makes that
coordinate present on every sufficiently refined record, while zero defect
makes its residual tend to zero.  Pointwise convergence of states therefore
passes the tested weak equation and the signed carrier relation to the limit.
Here satisfaction of the public weak equation means satisfaction on the
declared separating rational test core.  When the conventional formulation
uses all smooth tests, the approximation grammar also records the standard
distribution-order seminorm bound; on its ordinary graph that bound extends
the identity by density, while failure of the bound is its typed form
coordinate.
An ordinary weak solution satisfies these relations directly, so its
evaluation state lies on the zero-defect total graph.

The tail-limsup failure coordinates are countable and nonnegative.  If the
path is not zero defect, at least one \(d_k\) is nonzero.  The fixed enumeration
therefore has a least nonzero coordinate.  Its typed graph already contains
its input records and constructor ancestry, which makes it a first-failure
record rather than an untyped remainder.
\end{proof}

\begin{lemma}[Generator variances determine a Dirac cylinder state]
\label{lem:generator-variance-dirac-state}
Let \(L\) be a normalized positive functional on the Archimedean real
cylinder algebra \(\mathcal A_{\mathbb Q}(\mathfrak P)\), and let
\((e_n)_{n\ge1}\) be its bounded algebra generators.  If
\[
 L(e_n^2)=L(e_n)^2\qquad(n\ge1),
\]
then, for every cylinder polynomial \(p\),
\begin{equation}
 \label{eq:generator-variance-dirac-evaluation}
 L\bigl(p(e_1,\ldots,e_m)\bigr)
 =p\bigl(L(e_1),\ldots,L(e_m)\bigr).
\end{equation}
Consequently \(L(fg)=L(f)L(g)\) for all cylinder polynomials \(f,g\).
\end{lemma}

\begin{proof}
Use the GNS seminorm \(\|a\|_L^2=L(a^2)\).  The variance hypothesis gives
\(\|e_n-L(e_n)1\|_L=0\).  The null space of this seminorm is a left ideal:
for every cylinder polynomial \(b\), Archimedean boundedness supplies
\(C_b>0\) with \(C_b^2-b^2\in M_{\mathfrak P}\).  Indeed, from
\(C_b\pm b\in M_{\mathfrak P}\), the quadratic-module identity
\[
 2C_b(C_b^2-b^2)
 =(C_b-b)(C_b+b)^2+(C_b+b)(C_b-b)^2
\]
gives the assertion after increasing \(C_b\) to a positive rational value.
Consequently
\[
 \|ba\|_L^2=L(b^2a^2)\le C_b^2L(a^2)=0
\]
whenever \(\|a\|_L=0\).  Therefore, in the GNS quotient,
\(e_n=L(e_n)1\) for every generator.  Substitution through a finite product
gives
\[
 p(e_1,\ldots,e_m)=p(L(e_1),\ldots,L(e_m))1
\]
in that quotient.  Pairing with the cyclic vector proves
\eqref{eq:generator-variance-dirac-evaluation}; applying it to \(fg\) gives
multiplicativity.
\end{proof}

\begin{theorem}[Dirac realization gate]
\label{thm:dirac-realization-gate}
Let \(L\in X_{\rm gc}\) be a zero-defect limit of one compatible
approximation-tower path.  Enumerate the real state and path generators by
\((e_n)_{n\ge1}\), and put
\[
 \operatorname{Var}_L(e_n):=L(e_n^2)-L(e_n)^2\ge0.
\]
Exactly one of the following occurs.
\begin{enumerate}[label=\textup{(R\arabic*)}]
\item Some least variance, operation-graph, distribution-order, origin, or
  continuation coordinate is nonordinary or nonzero.  That coordinate, its
  detecting cylinder word, and its inherited five-source ancestry form the
  next realization successor.
\item Every variance vanishes, every public operation graph has its ordinary
  value, the distribution-order bounds are finite, and the approximation
  indices have independent origin.  Then \(L\) is evaluation on one ordinary
  same-origin weak path satisfying the public carrier law.  It is a singular
  terminal exactly when its ordinary continuation rows cross a cofinal target
  shell family.
\end{enumerate}
In particular, a positive functional, invariant law, compatible moment family,
or zero-defect relaxed state is not an ordinary or singular realization before
this gate.
\end{theorem}

\begin{proof}
Apply \cref{lem:generator-variance-dirac-state}; the generator-variance rows
make \(L\) multiplicative on the entire cylinder algebra.  Its linear test coordinates, together with the finite
distribution-order bounds, define one distributional path; ordinary values of
all typed operation graphs identify every nonlinear and boundary coordinate
with the corresponding public operation on that path.  Zero approximation
residuals give the weak equation and carrier law.  The retained origin indices
give independence, and the ordinary continuation-shell rows give contact.
If any clause fails, countability of the fixed enumeration selects its least
coordinate.  Total-graph and source-preservation rules make it the successor
in \textup{(R1)}.
\end{proof}

\begin{theorem}[Approximation--orbit realization criterion]
\label{thm:approximation-orbit-realization-criterion}
Let \(\mathfrak a\) be any graph-complete limit record, positive-functional
state, carrier-polar state, or joint-saturated prospective source atom.  If
there is one compatible path \((a_N)_{N\ge1}\) in the public approximation
tower satisfying
\begin{enumerate}[label=\textup{(D\arabic*)}]
\item every fixed equation, carrier, initial-data, operation-graph, trace,
  chart, and continuation residual tends to zero along this same path;
\item the path evaluations converge to \(\mathfrak a\) on every rational
  cylinder coordinate;
\item the limiting positive functional has zero variance on every real state
  and path generator; and
\item every distribution-order and domain coordinate has its ordinary finite
  value.
\end{enumerate}
then \(\mathfrak a\) represents one ordinary public PDE orbit.  Conversely,
an ordinary orbit claimed as a limit of the public approximation grammar has
such a path if and only if that approximation claim satisfies
\textup{(D1)}--\textup{(D4)}.  An ordinary orbit already present in the total
solution graph is realized directly by its evaluation record and need not be
assumed approximable by this particular grammar.

The tower-reconstructed orbit has independent origin precisely when the tower path
is generated without a target-contact or noncontinuation premise.  It reaches
the continuation target precisely when its ordinary continuation coordinates
cross every member of the generated cofinal shell family.  Thus neither a
compatible moment table, a measure on the graph compactification, a polar
separator, a nonordinary total-graph value, nor a contact-derived tower is a
singular-existence terminal.
\end{theorem}

\begin{proof}
Conditions \textup{(D1)}--\textup{(D2)} give the relaxed weak equation and
carrier law by
\cref{thm:approximation-realization-first-failure}.  Because the convergence
occurs along one compatible path, all coordinates have the same origin; no
coordinatewise change of subsequence is permitted.  By \textup{(D3)},
Cauchy--Schwarz makes the limit functional multiplicative on the cylinder
algebra.  Condition \textup{(D4)} and the ordinary values in \textup{(D1)}
then invoke the Dirac realization gate
\cref{thm:dirac-realization-gate}, reconstructing one distributional path and
identifying every nonlinear, boundary, chart, and continuation coordinate
with the corresponding public operation.  The zero residuals give the public
weak equation and initial datum.  Origin polarity is inherited from the one
tower path, and shell contact is exactly the stated ordinary continuation
condition.  This proves sufficiency.

If an ordinary orbit is asserted to arise through the approximation grammar,
the restrictions of that one asserted approximating path necessarily satisfy
\textup{(D1)}--\textup{(D2)}; multiplicativity of its ordinary evaluation
gives \textup{(D3)}, and ordinary domain values give \textup{(D4)}.  The
reverse implication is the sufficiency just proved.  A direct evaluation
record from the total solution graph already is the orbit it represents, so
no density or convergence property of \(\mathsf{Appr}\) is inferred from its
existence.
\end{proof}

\begin{corollary}[No compactification-to-existence promotion]
\label{cor:no-compactification-existence-promotion}
Every point of the graph-complete compactification that is neither a direct
ordinary evaluation from the total solution graph nor reconstructed by a
tower satisfying \textup{(D1)}--\textup{(D4)} remains an internal typed
successor.  Compactness
may be used to exclude a target cell in the larger relaxed space, but a point
of that space proves physical attainability only through
\cref{thm:approximation-orbit-realization-criterion}.  In particular, graph
completeness proves that no coordinate was forgotten; it does not prove that
an arbitrary compatible coordinate family is dynamically attainable.
\end{corollary}

\begin{definition}[Actual finite-prefix realization tree]
\label{def:actual-finite-prefix-realization-tree}
Fix a target-aligned structural branch signature \(\mathfrak b\) generated by
\cref{def:target-aligned-structural-branch-bound}.  Thus \(\mathfrak b\)
already contains the public presentation \(\mathfrak P_{\Phys}\), one
same-origin target-local stratum \(X_{\mathfrak b}\), its canonical shell
covector \(\lambda=D\Phi\), origin polarity, unnormalised shell cocycle,
exact provenance word \(A\), unique physical-event indices, and the processed
prefix of its cumulative joint record.  At depth \(n\), let
\(W_n(\mathfrak b)\) be the finite set of codes of resource at most \(n\),
over the first \(n\) rational symbols of the public approximation grammar,
whose stopped records lie in
\(\mathfrak D_{A,\Sigma,n}^{\rm pre}(\Budget)\) of
\cref{def:presentation-generated-prospective-source-domain}.  The
resource convention is prefix closed, every finite public code compatible
with \(\mathfrak b\) belongs to some \(W_n\), and the restriction of a code in \(W_{n+1}\) belongs to
\(W_n\).  The approximation enumerator producing \(W_n\) takes only
\(\mathsf{Pre}(\mathsf{Eq})\) as input.  In the independent-origin fiber, the
target, shell, support-floor, and continuation rows are evaluated only after a
code has been generated; they filter the exhaustive target-blind list and are
never arguments to the code generator or its refinement map.  In the
hypothetical fiber, the same rows retain the contact-selected prefix and its
\(\mathsf{Hyp}\) tag.  The two lists share a row checker but never exchange
records or parent maps.

The rooted tree \(\mathcal T^{\rm act}_{\mathfrak b}\) has at depth
\(n\) precisely those codes in \(W_n(\mathfrak b)\) whose finite transcripts
verify, to the certified tolerance \(2^{-n}\), the first \(n\) public weak
and total-graph rows, target-shell and continuation rows, target-local
factorization and carrier rows, branch-polarity rows, exact-support and
target-floor rows, domain rows, and generator-variance rows.  Every transcript
is formed from the corresponding finite prefix of
\(\mathsf{Proof}_{\Sigma}(\mathfrak P)\): it first joins all public inputs,
then applies \(q_\Sigma\), and only then restricts to the exact word \(A\).
A code is an actual finite record only when this transcript verifies.  The
parent map is public restriction or coarsening of the same joint record and
preserves \(\lambda\), \(A\), origin, event indices, and physical-carrier
ownership; for each of the first \(n\) cylinder coordinates its
transcript also verifies the refinement bound \(2^{-n}\) between parent and
child.  An ordinary signature includes fixed rational bounds away from every
declared domain boundary and a fixed finite distribution-order bound; failure
of either bound changes the signature to its nonordinary successor.  A node
retains its event order, finite residual vector, graph witnesses, and source
ancestry.  In particular, a positive moment
functional is not a node unless it is the evaluation of such an actual finite
record.

The effective public presentation generates the following finite realization
machinery from its fixed target-blind approximation grammar.  This machinery
is part of \(\mathsf{Proof}_{\Sigma}(\mathfrak P)\); it contains no separate
realization-completeness field or equation-class admission certificate.
\begin{enumerate}[label=\textup{(RC\arabic*)}]
\item The sets \(W_n\) are given by terminating finite enumerators.  The row
  checker returns, for every code, either a verified actual-node transcript or
  the least failed row with a checkable rejection transcript.  Parent maps
  preserve datum, origin, event order, all earlier verified rows, the
  summable refinement bounds, the ordinary-domain margins when present, the
  target covector, exact support, carrier owner, and cumulative joint-prefix
  restriction.
\item For every ordinary public orbit, finite cylinder, and generated branch
  signature, the grammar enumerates all rational approximations to the
  orbit's finite public coordinate record.  After a code has been generated,
  the totalized weak, nonlinear-graph, boundary, domain, carrier, and
  continuation relations return exactly one of: an accepted actual-prefix
  row; a target-null discrepancy; or the least nonzero target-aligned
  discrepancy with its graph witness and five-source ancestry.  The target,
  shell, support, and continuation coordinates are outputs of this check and
  never inputs to the enumeration or refinement map.
\item For every node \(v\), all codes in the finite extension list
  \(\operatorname{Ext}_n(v)\subseteq W_{n+1}\) are checked.  The list is
  exhaustive for the next resource level, and each rejected code retains its
  least failed row and graph-error enclosure.  The indicator partition of
  this list is part of the finite ordered proof language.
\item The level algebra is the full function algebra
  \(\mathbb R^{\mathcal T_n^{\rm act}}\), obtained as the restriction of the
  finite branch of \(\mathsf{Proof}_{\Sigma}(\mathfrak P)\).  The public
  approximation-code coordinate is rational and injective on the finite code
  universe; Lagrange interpolation in that coordinate computes the point
  indicators.  The checker verifies their multiplication table,
  sum-to-one identity, and restriction maps.  No target-relative function is
  appended after a row is evaluated, and no representation theorem for
  positive functionals is included as presentation data.
\end{enumerate}
The machinery is \emph{effective} because the records, parent maps, graph
enclosures, discrepancy rows, and finite refinement bounds in
\textup{(RC1)}--\textup{(RC4)} are recursively checkable.  Its finite proof
language is \emph{complete} when the exhaustive
extension transcript and indicator partition mechanically produce the
\(-1\)-identity for every empty actual-child list, and every displayed
independent realization has an enumerated verified orbit-and-contact
certificate.  Thus neither extension-cone exactness nor positive-functional
representation is an admission hypothesis.
\end{definition}

\begin{remark}[Scope of public-grammar realization]
\label{rem:realization-completeness-scope}
The construction uses density of rational finite-cylinder records in the
public coordinate topology, together with the total graphs and moduli already
present in \(\mathfrak P_{\Phys}\).  It assumes no convergence theorem for a
preferred Galerkin or time-discretization scheme.  Failure of approximation
in an ordinary graph coordinate is itself a public graph value: after
cumulative joining and target projection it is target-null or becomes the
next source-resolved successor.  Thus the realization stack is complete for
target-aligned structure without asserting density in a stronger ambient
solution topology.
\end{remark}

\begin{theorem}[Actual finite-prefix tree]
\label{thm:actual-finite-prefix-tree}
Every node of \(\mathcal T^{\rm act}_{\mathfrak b}\) is an actual
same-origin finite public record satisfying exactly the finite branch rows
recorded at its depth; its parent is an actual record with the same initial
datum, origin, and event order.  Conversely every actual finite record of
resource at most \(n\) satisfying those rows occurs as a node at depth \(n\).
\end{theorem}

\begin{proof}
Enumerate the finite set \(W_n\) and run its terminating row checker.  By
\textup{(RC1)}, acceptance is equivalent to possessing all required verified
rows, and restriction preserves every earlier row.  Exhaustiveness of the
resource convention places every finite public record at some level; at a
fixed level the accepted codes are exactly the stated nodes.  No functional
or moment table enters this filtering operation.  Membership in
\(\mathfrak D_{A,\Sigma,n}^{\rm pre}(\Budget)\) supplies finite public
ancestry; the branch signature supplies \(X_{\mathfrak b}\), \(\lambda\),
origin, and carrier ownership; and the checked joint-prefix restriction
supplies the target-aligned evaluation.  Thus acceptance introduces no node
outside \(\mathsf{Proof}_{\Sigma}(\mathfrak P)\).
\end{proof}

\begin{theorem}[Approximation-grammar completeness]
\label{thm:approximation-grammar-completeness}
Let \(u\) be an ordinary public orbit and \(\mathfrak b\) a generated
target-aligned branch signature.  Under the priority order displayed below,
exactly one of the following occurs:
\begin{enumerate}[label=\textup{(A\arabic*)}]
\item \(u\) has an enumerated direct-evaluation code in the ordinary
  total-solution graph;
\item in the absence of \textup{(A1)}, the fixed public grammar supplies one compatible path in
  \(\mathcal T^{\rm act}_{\mathfrak b}\) satisfying the residual, graph,
  domain, carrier, and continuation rows;
\item in the absence of \textup{(A1)}--\textup{(A2)}, the tree has a least
  empty finite level.  The exhaustive rejection transcript at that level
  gives a cumulative approximation-discrepancy record which is target-null or
  has a least nonzero target-aligned owner, with its total-graph witness,
  shell covector, exact provenance, origin, event index, and joint prefix.
  The nonnull value is the next typed structural successor.
\end{enumerate}
Consequently failure of full-orbit approximation is never a realization
admission failure.  It either has no effect on the selected target or is one
of the exhaustively generated mechanisms that can produce target alignment.
\end{theorem}

\begin{proof}
An enumerated direct evaluation gives \textup{(A1)}.  Otherwise fix a finite public
cylinder and rational tolerance.  Density of rational Galerkin coefficients,
mollifier scales, time meshes, and domain exhaustions produces a grammar code
approximating the orbit's primitive cylinder coordinates.  The computable
moduli generated on the shrinking-enclosure cells propagate this
approximation through every ordinary constructor.  If an enclosure does not
shrink, its least diameter, oscillation, concentration, or graph-boundary
coordinate is included in the cumulative discrepancy record of
\textup{(A3)}.  The code is generated from
\(\mathsf{Pre}(\mathsf{Eq})\) before any target row is evaluated.

Run the finite row checker on the complete resource-bounded code universe.  If
the accepted tree has nodes at every depth, finite branching and the parent
refinement bounds give one compatible same-origin path, proving
\textup{(A2)}.  Otherwise choose the least empty depth and combine the least
failed rows of every rejected code with the public indicator partition.  This
is the finite exhaustive approximation-discrepancy transcript of
\textup{(A3)}.  Retain its complete expanded derivation and join it to the
preceding public record before applying \(q_\Sigma\).  A zero projection
gives the null cell.  Otherwise universal five-source provenance and
recursive carrier--polar exhaustion assign the least failed graph owner and
its typed successor.  Continuum composition no-creation exposes any nonzero
limiting discrepancy on a finite joint prefix, so there is no fourth cell.
All three outcomes are derived from the fixed public presentation.
\end{proof}

\begin{theorem}[Finite extension obstruction]
\label{thm:finite-extension-obstruction}
Let \(v\) be a node of \(\mathcal T^{\rm act}_{\mathfrak b}\) for an engine
instance generated by an effective public presentation.  Exactly one holds:
\(v\) has an actual
child satisfying the next generated row, or a finite ordered certificate in
the actual-child cone proves that no such child exists.  Thus a parent cannot
survive merely because a relaxed positive functional satisfies the next row.
\end{theorem}

\begin{proof}
Run the checker on the finite exhaustive list
\(\operatorname{Ext}_n(v)\).  If one transcript is accepted, its code is an
actual child.  Otherwise every candidate has a least failed-row transcript.
Let \(\mathbf1_w\) be the verified point indicators from \textup{(RC4)}.
The indicator identity \(1=\sum_w\mathbf1_w\), together with the rejection
row forcing \(\mathbf1_w=0\) on the feasible locus for every candidate,
gives the mechanically checkable ordered-cone identity \(-1\in M_v\).
All sums are finite, and each nonpolynomial row uses its retained graph-error
enclosure.  Acceptance and exhaustive rejection are disjoint and exhaustive;
no exactness of an externally declared extension cone is invoked.  Every
summand is the restriction of a checked public row or the rational code
indicator generated by \(\mathsf{Appr}\), so \(M_v\) is a finite branch of
\(\mathsf{Proof}_{\Sigma}(\mathfrak P)\), not a new proof cone.
\end{proof}

\begin{theorem}[No relaxation gap on target-local branches]
\label{thm:no-relaxation-gap-target-local}
At every finite level of a target-local branch generated by an effective
public presentation,
\[
 \{\text{positive functionals satisfying the retained rows}\}
 =
 \overline{\operatorname{conv}}
 \{\text{evaluations of actual same-origin nodes}\}.
\]
Consequently every positive target detector is positive on an actual node.
\end{theorem}

\begin{proof}
Write \(K=\mathcal T_n^{\rm act}\).  If \(K\) is empty,
\cref{thm:finite-extension-obstruction} gives the obstruction certificate.
Otherwise the public rational code coordinate separates \(K\), and its
verified Lagrange indicators \((\mathbf1_x)_{x\in K}\) are a basis of the
restricted public branch algebra \(\mathbb R^K\).  For a normalized positive functional \(L\), set
\(\lambda_x=L(\mathbf1_x)\).  Positivity gives \(\lambda_x\ge0\), and the
sum-to-one identity gives \(\sum_x\lambda_x=1\).  Hence, for every
\(f\in\mathbb R^K\),
\[
 L(f)=\sum_{x\in K}\lambda_x f(x)
     =\sum_{x\in K}\lambda_x\operatorname{ev}_x(f).
\]
This proves both inclusions in the displayed equality without a relaxation
or closure gap.  If \(L(d)>0\) for a nonnegative detector \(d\), at least one
\(x\) with \(\lambda_x>0\) satisfies \(d(x)>0\).
\end{proof}

\begin{theorem}[Coherent actual-branch compactness]
\label{thm:coherent-actual-branch-compactness}
If every depth of \(\mathcal T^{\rm act}_{\mathfrak b}\) is nonempty, then
there is one infinite same-origin compatible path
\(a_1\preceq a_2\preceq\cdots\) through the tree.
\end{theorem}

\begin{proof}
Each level is finite, hence the tree is finitely branching.  Choose the root
component containing nodes at arbitrarily large depth.  At each step only
finitely many children are available, so one child again has descendants at
arbitrarily large depth.  K\H{o}nig's lemma, in this explicit form, gives one
compatible infinite path.  Dead side nodes need not be retained, and there is
no coordinatewise subsequence change.
\end{proof}

\begin{theorem}[Automatic Dirac realization along an actual branch]
\label{thm:automatic-dirac-realization}
Every infinite path supplied by
\cref{thm:coherent-actual-branch-compactness} satisfies
\textup{(D1)}--\textup{(D4)} of
\cref{thm:approximation-orbit-realization-criterion}.  Hence it reconstructs
one ordinary public PDE orbit when its realization rows are ordinary.
\end{theorem}

\begin{proof}
At depth \(n\), the path satisfies the first \(n\) residual, graph, domain,
and variance rows to tolerance \(2^{-n}\).  Each fixed residual therefore
tends to zero.  For a fixed cylinder coordinate, the parent--child
refinement estimates have summable tail \(\sum_{k\ge n}2^{-k}\), so its
values are Cauchy along this one path and converge without selecting a new
subsequence.  Because every node is an actual evaluation, the identities
\(L_n(fg)=L_n(f)L_n(g)\) hold whenever the three functions have entered the
finite algebra; passage to their common limit gives multiplicativity and
zero generator variance.  Finally, the fixed rational domain margins and
distribution-order bound are part of the ordinary signature and persist
under every parent map.  A failure changes the node to the declared
nonordinary successor rather than remaining on this path.  These are,
respectively, \textup{(D1)}--\textup{(D4)}.
The limits retain the common origin, shell covector, exact source word, and
physical carrier because these are fixed parent-map coordinates.  The weak
and total-graph limits are therefore evaluated on the same cumulative joint
record before \(q_\Sigma\), as required by
\cref{thm:closed-target-aligned-proof-discipline}.
\end{proof}

\begin{theorem}[Effective moduli from a public actual-path program]
\label{thm:public-actual-path-effective-moduli}
Let \(\gamma:\mathbb N\to\bigcup_nW_n(\mathfrak b)\) be a total program in
the public instruction language such that \(\gamma(n)\) is an accepted node
at depth \(n\) and is the parent restriction of \(\gamma(n+1)\).  Then the
program and the public row checker generate the complete effective
realization certificate.  At depth \(n\), every entered residual and graph
error is at most \(2^{-n}\), every fixed cylinder coordinate has remaining
Cauchy tail at most
\[
 \sum_{k\ge n}2^{-k}=2^{1-n},
\]
the multiplicativity and generator-variance defects vanish on the entered
finite algebra, and the ordinary domain margin, distribution-order bound,
origin, shell covector, provenance, event owner, and cumulative joint prefix
are inherited from the parent maps.  If \(\gamma(n)\) verifies the first
\(n\) continuation shells, the same program supplies the cofinal-contact
modulus.
\end{theorem}

\begin{proof}
The residual and graph bounds are acceptance rows of the depth-\(n\) node.
Summing the certified parent--child refinement bounds gives the displayed
Cauchy modulus.  Every node is an actual evaluation on its entered finite
algebra, so multiplicativity and zero variance are exact there.  The ordinary
domain and dependency fields are preserved by \textup{(RC1)}.  The shell
claim follows because the depth counter is a modulus for the checked cofinal
family.  All operations used to compute these bounds are primitive
constructors of the public instruction language.
\end{proof}

\begin{theorem}[Target-aligned accountability of public algorithms]
\label{thm:target-aligned-public-algorithm-saturation}
Fix a presentation-generated affordable target-aligned cell and enumerate all
finite public instruction words acting on its stopped records.  Close each
word under expanded derivation, cumulative same-origin joining, all declared
total graphs, carrier--polar descendants, and actual-prefix restriction
before applying \(q_\Sigma\).  Then every algorithmic contribution to contact
has exactly one disposition:
\begin{enumerate}[label=\textup{(PA\arabic*)}]
\item its complete target projection is zero;
\item a finite charged prefix exceeds the generated admissible-carrier ceiling of the
  cell;
\item a least finite nonzero target projection has a complete five-source
  owner and is the next public structural successor; or
\item every finite contribution has been processed, and the cumulative
  record is a joint-saturated atom carrying its complete actual-prefix tree.
\end{enumerate}
These alternatives classify every public algorithmic contribution to contact.
They do not decide whether an infinite finitely branching actual-prefix tree
has a computable path.
\end{theorem}

\begin{proof}
At resource \(n\), there are finitely many instruction words and finitely
many codes in \(W_n\).  Form their complete joint records before target
projection.  Continuum composition no-creation places every nonzero mixed
value, including a graph, approximation, extension, or selector failure, on
a finite joint prefix with its inherited source word.  Recursive
carrier--polar exhaustion then gives \textup{(PA1)}, \textup{(PA2)}, or the
least successor in \textup{(PA3)}.  On an infinite route, the
no-emergence trichotomy gives another finite unprocessed increment unless the
cumulative record is fixed, which is \textup{(PA4)}.  Thus saturation covers
instruction words as expanded public constructors; it does not quantify over
an ambient class of unrelated algorithms.  In \textup{(PA4)},
\cref{thm:coherent-actual-branch-compactness,thm:automatic-dirac-realization}
give a semantic ordinary orbit when all levels are nonempty and ordinary.
Computable path extraction requires the separate finite-progress condition
below.
\end{proof}

\begin{definition}[Generated Noetherian target cell]
\label{def:generated-noetherian-target-cell}
An affordable target-aligned cell is \emph{generated Noetherian} when the
fixed public rank compiler constructs from its successor graph a computable
rank
\[
 \rho:\mathsf{State}_\Sigma\longrightarrow W
\]
into a represented well-founded order \(W\), and the finite checker verifies
the following facts from the public transition rows:
\begin{enumerate}[label=\textup{(N\arabic*)}]
\item every state is a finite regularity terminal, a finite independently
  realized singular terminal, or has a finitely enumerable nonempty
  successor list;
\item every successor \(s'\) of a nonterminal state \(s\) satisfies
  \(\rho(s')\prec_W\rho(s)\); and
\item the terminal tags and every rank comparison are decidable from the
  finite public records.
\end{enumerate}
The rank compiler may use only a finite target quotient, a discrete carrier
budget, a finite support/refinement lattice, and lexicographic or multiset
combinations of these already generated public coordinates.  The
presentation, datum, and output-package schemas contain no field for a rank,
termination argument, or well-foundedness certificate.  Classification of
target-aligned structure alone does not supply this rank.
\end{definition}

\begin{theorem}[Finite target-aligned terminalization]
\label{thm:generated-noetherian-target-terminalization}
Every execution on a generated Noetherian target cell reaches after finitely
many transitions exactly one checked terminal: a finite regularity
certificate or a finite independently realized singular certificate.
\end{theorem}

\begin{proof}
If a checked state is nonterminal, \textup{(N1)} gives its finite nonempty
successor list and \textup{(N2)} strictly decreases \(\rho\) along every
choice.  An infinite execution would therefore produce an infinite descending
chain in \(W\), contradicting well-foundedness.  Hence every execution reaches
a terminal after finitely many steps.  The two terminal tags are disjoint by
the origin and contact checks, and their soundness follows from the
proof-carrying output verifier.
\end{proof}

\begin{corollary}[Target accountability does not imply a universal finite rank]
\label{cor:no-universal-target-halting-rank}
There is no construction which assigns the generated Noetherian rank of
\cref{def:generated-noetherian-target-cell} to every effective affordable
target-aligned public process solely from its five-source accountability
record.
\end{corollary}

\begin{proof}
Part~I embeds every deterministic Turing machine and finite input as a
countably represented process with exact finite stopped slices, including
nonhalting computations.  Declare the halting configuration to be the target.
Its indicator is an affordable public boundary readout, and every finite
transition has complete target-aligned provenance.  A universal generated
Noetherian rank would make
\cref{thm:generated-noetherian-target-terminalization} return in finite time
either a target-reaching certificate or a certificate that the target is
never reached.  This would decide the halting problem.  Hence source
accountability and affordability cannot by themselves supply the rank.
\end{proof}

\begin{theorem}[Fixed-atom realization-or-successor]
\label{thm:fixed-atom-realization-or-successor}
Let \(\mathfrak a\) be a joint-saturated target-active atom fixed under every
generated primal and polar operation of an effective public presentation.
Build its tree with the complete fixed signature of \(\mathfrak a\).  Exactly
one occurs: some finite level is empty and its exhaustive rejection
transcripts compose backward to a root obstruction certificate; every level
is nonempty and the fixed signature is nonordinary, producing a proper
successor; or every level is nonempty and one compatible actual approximation
path realizes \(\mathfrak a\) as an ordinary public orbit.
\end{theorem}

\begin{proof}
If a level is empty, choose the least such depth.  Apply
\cref{thm:finite-extension-obstruction} to every node of the preceding finite
level and substitute their indicator identities successively toward the
root.  The resulting finite transcript is the first alternative.  Otherwise
the finitely branching tree has nodes at every depth, so
\cref{thm:coherent-actual-branch-compactness} supplies a compatible path (and
does not require that dead side nodes have children).  The complete signature
is fixed before the tree is built, so its continuation/domain type is either
nonordinary on the whole tree or ordinary on the whole tree.  In the former
case it creates the proper successor.  In the latter,
\cref{thm:automatic-dirac-realization} gives the ordinary orbit.  Emptiness,
nonordinary type, and ordinary type are mutually exclusive.  The atom itself
is never promoted to an orbit before this test.
\end{proof}

\begin{theorem}[Two-polarity realization]
\label{thm:two-polarity-realization}
The hypothetical exclusion queue and the independent construction queue are
logically separate.  The former may schedule a complete hypothetical atom
signature, but the latter starts only from the public equation, datum,
approximation grammar, and target-blind generated records.  A singular
terminal is accepted only when an independent ordinary orbit matches the
complete signature and has cofinal contact.
\end{theorem}

\begin{proof}
Origin polarity is immutable under every constructor by
\cref{def:ordered-continuum-successors}.  Enumerate each finite independent
code universe from \(\mathsf{Pre}(\mathsf{Eq})\), run the public weak, carrier,
and total-graph checks, and refine its parent map before evaluating any target
row.  The shell, target quotient, exact provenance, continuation, and contact
coordinates are then computed as outputs of the generated record.  No
hypothetical branch ideal or contact-derived coefficient enters the code
universe or refinement map.  Compare the resulting completed signature with
the hypothetical one only after construction.
Thus no \(\mathsf{Hyp}\to\mathsf{Ind}\) conversion is used, and an accepted
match has independent origin.
\end{proof}

\begin{theorem}[Continuation-fidelity terminal test]
\label{thm:continuation-fidelity-terminal-test}
An independently realized atom is singular exactly when its ordinary
same-origin path solves the original public PDE, its continuation coordinate
is ordinary and noncontinuable, and its restrictions cross every shell in the
generated cofinal family.  An ordinary-continuable row closes as regular; a
nonordinary row is a proper successor.
\end{theorem}

\begin{proof}
Evaluate the total continuation relation on the ordinary realized path.
Its discrete status distinguishes continuation from noncontinuation, and
cofinal shell crossing is the declared target contact.  Totalization routes a
failed continuation or trace operation to its first successor; origin
polarity prevents a terminal germ from another construction entering the
cell.
\end{proof}

\begin{theorem}[Explicit singular-mechanism equivalence]
\label{thm:explicit-singular-mechanism-equivalence}
An independent realization is a singular terminal if and only if it satisfies
\textup{(A1)}--\textup{(A4)} of
\cref{def:joint-saturated-structural-atom}, \textup{(E1)}--\textup{(E6)} of
\cref{def:explicit-singular-mechanism}, and the continuation-fidelity test.
\end{theorem}

\begin{proof}
An independent singular terminal has its complete compiler transcript, which
supplies the atom, realization, and continuation rows.  Conversely those
ordinary rows give an independent public PDE orbit with cofinal contact and
noncontinuation.  The two-polarity realization theorem excludes every
contact-assumed construction.
\end{proof}

\begin{theorem}[Fair parallel support scheduler]
\label{thm:fair-parallel-support-scheduler}
There is a fair dovetailing schedule for all thirty-one support queues,
carrier-dual searches, polarity refinements, and independent realization
searches.  Every finite proof, successor, or realization prefix is processed
after finitely many master stages.
\end{theorem}

\begin{proof}
At master stage \(n\), process every tuple of support, branch depth, cylinder
depth, graph index, dual-word index, and approximation level whose code
length is at most \(n\).  Retain state between visits.  Every fixed finite
tuple is therefore visited eventually, and no nonterminating search blocks
another.
\end{proof}

\begin{theorem}[Compact finite-cover regularity certificates]
\label{thm:compact-finite-cover-regularity-certificates}
If every actual contact candidate is excluded by a finite target-local
certificate, finitely many such certificates exclude the compact contact
candidate set.
\end{theorem}

\begin{proof}
Each strict finite certificate excludes an open neighborhood in the
finite-cylinder topology.  These neighborhoods cover the compact candidate
set, so a finite subcover exists.  Concatenate its checked ordered
identities, source routes, and continuation exclusions.
\end{proof}

\begin{theorem}[Regularity-certificate completeness]
\label{thm:regularity-certificate-completeness}
For an engine instance generated by an effective public presentation,
emptiness of the actual contact set
implies that every remaining
branch is eventually removed by a finite ordered contradiction, target
nullity, continuation, class incompatibility, a finite price prefix exceeding
the public budget, or a finite realization obstruction.  Hence a finite global regularity
certificate exists.
\end{theorem}

\begin{proof}
Fix a contact signature.  If its finite actual tree had a node at every
depth, \cref{thm:coherent-actual-branch-compactness,%
thm:automatic-dirac-realization} would give an actual contacting orbit,
contrary to emptiness.  Hence it has a least empty finite level.  Exhaustive
finite checking and the indicator calculation in
\cref{thm:finite-extension-obstruction} compose to a root obstruction
transcript.  Its checked graph enclosures give a finite cylinder on which the
same transcript remains valid.  No-relaxation gap removes positive
functionals on that cylinder, while sourcewise terminalization handles every
unbounded target-amplitude route.  These cylinders cover the compact contact
candidate set; apply the finite-cover theorem.  Thus completeness here is a
consequence of \textup{(RC1)}--\textup{(RC4)}, not a separately assumed proof
language.
\end{proof}

\begin{theorem}[Singular-witness criterion and Noetherian completeness]
\label{thm:singular-witness-completeness}
Let a datum have an ordinary local germ in the public solution graph, and
apply the generated finite admissible-carrier stratification and target-aligned
branch-and-bound construction.  On an affordable cell with nonempty actual
contact set, the independent realization queue produces a recursively
checkable explicit singular mechanism precisely when a contacting orbit has
an enumerated direct-evaluation proof or a compatible actual-prefix path
generated by a total public program.  The program generates all convergence,
residual, variance, domain, provenance, and contact moduli through
\cref{thm:public-actual-path-effective-moduli}.

If the cell is generated Noetherian, the queue reaches such a finite singular
terminal after finitely many transitions.
\end{theorem}

\begin{proof}
Choose an actual contacting orbit in the cell.  Its unit shell cascade and
finite support exhaustion give a persistent nonempty source support with
unbounded canonical amplitude.  Apply sourcewise terminalization.  A charged
prefix above the generated ceiling contradicts affordability, target nullity
contradicts shell contact, and every nonzero failure of passivity, carrier
boundedness, exterior finiteness, graph fidelity, approximation, or
realization is a source-resolved successor.  Hence the route cannot terminate
before ordinary realization.

Apply approximation-grammar completeness to the orbit and its generated
signature.  A direct row gives the first certificate form.  A compatible
actual path gives a semantic orbit by automatic Dirac realization.  It gives
a recursively checkable certificate exactly when a total public program
generates the path; the public actual-path theorem then generates all listed
moduli.  Every failed approximation operation is target-null or a typed
successor by target-aligned algorithm accountability, so no approximation
failure is omitted from the semantic source calculus.  Conversely, every
accepted singular mechanism contains a direct proof or its finite total path
program, proving necessity of the criterion.

On a generated Noetherian cell, finite target-aligned terminalization reaches
a checked terminal.  Actual contact excludes the regularity terminal, while
two-polarity realization and the continuation-fidelity test identify the
remaining terminal as the stated independent singular mechanism.
\end{proof}

\begin{theorem}[Global binary semantics]
\label{thm:global-binary-semantics}
For every effective public presentation, every datum whose ordinary local
germ is emitted by its public local-solution graph, and every affordable
target-aligned cell generated by finite physical stratification,
\[
 \text{all source routes close regularly}
 \quad\Longleftrightarrow\quad \Reach_\Sigma^{\Phys}=0,
 \qquad
 \text{some independent terminal atom exists}
 \quad\Longleftrightarrow\quad \Reach_\Sigma^{\Phys}>0.
\]
\end{theorem}

\begin{proof}
Let an ordinary branch contact \(\Sigma\).  Canonical shell accountability
produces its cofinal family of unit same-origin shell currents.  Forming each
cumulative public joint record before \(q_\Sigma\), and then applying the
universal provenance morphism, selects a persistent nonempty exact-support
route with unbounded canonical amplitude.  Apply the exhaustive
target-aligned carrier stratification.  On its affordable cell, uniform
conversion, local target-aligned charging, support-transfer conservation,
and sourcewise terminalization either produce a finite prefix exceeding the
generated public ceiling or pass the complete signature to actual-prefix
realization; the first alternative closes that cell.  A target-null carrier
value closes directly.  Every remaining failure of passivity, carrier lower
boundedness, exterior-input finiteness, coherence, or domination has a least
five-source target-aligned owner and therefore enters the same sourcewise
successor and realization route.  Coherent-branch compactness, automatic
Dirac realization, and two-polarity realization identify the semantic
ordinary contacting orbit in the independent fiber.  A finite singular
certificate is supplied only under the criterion of
\cref{thm:singular-witness-completeness}.

Conversely, an independent terminal atom contains a verified public weak
orbit, ordinary continuation-status value, and cofinal canonical shell
crossings, hence is an ordinary contacting orbit.  The closed target-aligned
proof discipline and universal five-source provenance place every nonnull
contact current in one of the thirty-one routes and send every graph,
equality, carrier, or realization failure to its typed successor.  Therefore
no further contact route or terminal value exists.
\end{proof}

\begin{theorem}[Mutual exclusivity and no third terminal]
\label{thm:global-mutual-exclusivity-no-third-terminal}
No datum has both a valid regularity and singular certificate.  Moreover a
positive functional, invariant measure, moment table, zero-price law,
equality kernel, compactified state, prospective atom, graph boundary,
nonordinary realization value, infinite rank chain, or closure fixed point
is an internal compiler state, never a terminal.
\end{theorem}

\begin{proof}
A regular certificate excludes every ordinary same-origin contact cell,
whereas a singular certificate supplies a verified member of one.  The
finite verifier cannot accept both.  The transition rules respectively route
the listed objects through actual-support recovery, resaturation, consequence
generation, fidelity succession, no-emergence, or realization.
\end{proof}

\begin{theorem}[Global regularity on the public ordinary datum domain]
\label{thm:all-data-global-regularity}
Fix an effective public presentation.  If the binary engine returns
regularity on every affordable target-aligned cell generated from every datum
in \(\mathscr D_{\rm eng}(\mathfrak P_{\rm eff})\), then every ordinary
maximal branch from that public ordinary datum domain is global and remains
in the presentation's ordinary class.
\end{theorem}

\begin{proof}
The datum domain is the initial-trace projection of the native public
local-solution graph, so no local-existence premise is added.  Exhaustive
carrier and algorithm stratification places every target-bearing value in a
generated affordable cell or its typed successor.  The singular output is
absent on all of them.  Global binary semantics gives target avoidance, and
terminal coverage converts avoidance of finite noncontinuation into
continuation on every finite interval.
\end{proof}

\begin{theorem}[Effective cylinder approximation and checker soundness]
\label{thm:effective-cylinder-approximation-checker-soundness}
For every effective public presentation and every affordable target-aligned
cell it generates, each bounded observable and total graph carries nested
rational finite-cylinder inner and outer enclosures.  Exactly one occurs: the
enclosures have the composed public approximation modulus, their target
projection is zero, or their least persistent nonordinary enclosure
coordinate is a typed five-source successor.  A finite
regularity certificate checked by rational arithmetic, symbolic identities,
PSD certificates, interval bounds, graph relations, and source-origin tags
soundly excludes its claimed target cell.
\end{theorem}

\begin{proof}
Primitive constructors enumerate nested graph enclosures; structural
composition and rational truncation preserve nesting.  Shrinking radii
produce the composed modulus.  A nonshrinking radius is retained in the
expanded joint record, and no-creation plus provenance makes its nonzero
target projection the next successor.  Expanding every checked identity and
applying weak duality excludes a positive target-active realization, proving
soundness on every ordinary enclosure cell.
\end{proof}

\begin{theorem}[Effective extension and realization certificates]
\label{thm:effective-extension-realization-certificates}
At each actual-prefix node of an effective engine instance generated by its
public presentation,
parallel enumeration eventually finds either a checkable actual child or a
finite extension-obstruction certificate.  Every independent singular
terminal has a finite formal certificate or a finite program-plus-proof
object establishing its path, convergence moduli, residuals, variances,
domain rows, continuation failure, and cofinal contact.
\end{theorem}

\begin{proof}
At a node, enumerate its finite extension list and run the terminating row
checker on every entry.  An accepted entry is a child; if none is accepted,
the rejection transcripts and verified indicator table produce the finite
obstruction by \cref{thm:finite-extension-obstruction}.  Thus this decision
terminates directly and does not appeal to a semidecision procedure for an
assumed exact cone.  The same effective grammar represents the data required by
\textup{(D1)}--\textup{(D4)} and the contact test, so verifier soundness
checks any resulting realization object.  Target-aligned algorithm
accountability classifies every failed finite approximation operation.  On a
generated Noetherian contacting cell,
\cref{thm:generated-noetherian-target-terminalization,thm:singular-witness-completeness}
ensure that the finite process reaches a direct proof or a public
program-plus-proof object.
\end{proof}

\begin{definition}[Proof-carrying engine output]
\label{def:proof-carrying-engine-output}
For an accepted effective presentation and datum, a finite
\emph{regularity output} is a package
\[
 \mathsf{Regular}\bigl(\mathfrak P,u_0,\Budget,\Sigma,
       (C_A)_{\varnothing\ne A\subseteq\Channels}\bigr),
\]
where each \(C_A\) is a finite target-local exclusion transcript: a typed
ordered identity, target-null relation, continuation or incompatibility row,
budget-exceeding finite price prefix, or finite realization obstruction.
Every \(C_A\) contains the indices of its public weak and total-graph rows,
the active stratum, shell covector, cumulative joint prefix, exact word
\(A\), origin, physical-event indices, carrier owner, and graph enclosures.
A finite \emph{singularity output} is a package
\[
 \mathsf{Singular}\bigl(\mathfrak P,u_0,u_\ast,\mathfrak a,
       C_{\rm orbit},C_{\rm contact},C_{\rm closure}\bigr),
\]
where \(C_{\rm orbit}\) is an enumerated direct-orbit proof or the generated
total program-plus-proof for one independent approximation path and its
\textup{(D1)}--\textup{(D4)} moduli,
\(C_{\rm contact}\) verifies the ordinary continuation and cofinal-shell
rows, and \(C_{\rm closure}\) verifies the joint atom and all resolved
successor rows.  These three certificates carry the same public relation
indices, datum, origin, shell covectors, exact provenance, cumulative joint
frontier, event indices, and physical carrier.  The orbit program is derived
from \(\mathsf{Pre}(\mathsf{Eq})\); contact is checked only after that
independent program has been verified.

The verifier \(\mathsf{Check}_{\rm eng}\) first invokes
\(\mathsf{Check}_{\rm pres}\).  It then checks every finite transcript,
interval, graph modulus, proof-language identity, source word, origin tag,
and program certificate in the displayed package.  It also verifies that
joint formation precedes \(q_\Sigma\), that every price cites
\eqref{eq:target-aligned-charge-rule} and one canonical event index, and that
every source tag is the value of the public provenance morphism.  It rejects
malformed, incomplete-dependency, or mixed-polarity packages.
\end{definition}

\begin{theorem}[Proof-carrying output verifier]
\label{thm:proof-carrying-output-verifier}
\(\mathsf{Check}_{\rm eng}\) terminates on every finite proposed output.
Acceptance of a regularity package proves that no ordinary maximal branch
from \(u_0\) reaches \(\Sigma\).  Acceptance of a singularity package proves
that \(u_\ast\) is an independent ordinary public orbit realizing the
displayed target contact.  No accepted package has both output tags.
\end{theorem}

\begin{proof}
Presentation verification terminates by
\cref{thm:effective-canonical-presentation}.  Each remaining check is finite
rational arithmetic, a finite symbolic identity, a public graph enclosure,
or verification of a total code and its certified modulus.  The dependency
checks are the finite typing rules of
\cref{thm:engine-target-aligned-public-closure}.  The regularity
transcripts cover all thirty-one supports and are sound by the finite-cover
and checker-soundness theorems.  The singularity transcript satisfies the
explicit-mechanism equivalence.  Origin-tag checking and
\cref{thm:global-mutual-exclusivity-no-third-terminal} exclude mixed output
tags.
\end{proof}

\begin{theorem}[Public target-aligned closure of analytic premises]
\label{thm:public-target-aligned-premise-closure}
On the presentation-generated ordinary datum domain, every premise used by
the binary engine is one of the following native public values:
\begin{enumerate}[label=\textup{(PC\arabic*)}]
\item the ordinary local germ is an element of the target-blind public
  local-solution graph whose initial-trace projection defines the datum
  domain;
\item passivity, the carrier lower bound, exterior-input finiteness, the
  generated admissible-carrier budget, and finite-carrier affordability are
  values of the expanded native carrier graph;
\item graph, shell, and cylinder approximation moduli are shrinking values of
  nested public enclosures;
\item approximation coverage, finite extension, and realization rows are
  values of the saturated public instruction grammar; and
\item continuation and cofinal contact are values of the same-origin public
  continuation and shell graphs; and
\item every equation-wide estimate used downstream is first expanded into a
  finite public identity, restricted to the active target cell, and assigned
  its carrier owner; an unexpanded global norm bound is not an engine premise.
\end{enumerate}
Every unfavorable value in \textup{(PC2)}--\textup{(PC6)} is cumulatively
joined before target projection and has exactly one disposition: it is
target-null, it closes an affordable cell by a finite physical or ordered
contradiction, or its least nonzero target projection is a typed five-source
successor.  Hence no local-existence, budget, passivity, coercivity,
finite-carrier, graph-modulus, approximation, realization, convergence,
continuation, or contact certificate is carried as an additional analytic
premise by the datum.  Effective termination additionally requires the
generated Noetherian property of
\cref{def:generated-noetherian-target-cell}; target-aligned classification
alone does not imply that property.
\end{theorem}

\begin{proof}
\textup{(PC1)} is
\cref{def:declared-global-analysis-datum-class}.  Exhaustive target-aligned
carrier stratification gives \textup{(PC2)} and routes every unfavorable
carrier value.  The effective presentation theorem and effective cylinder
theorem give \textup{(PC3)}, retaining every nonshrinking enclosure as a
total-graph successor.  Approximation-grammar completeness and target-aligned
public-algorithm accountability give \textup{(PC4)} and retain every infinite
actual-prefix tree as a semantic realization value.
The public continuation graph and continuation-fidelity test give
\textup{(PC5)}.  Universal physical-quantity ownership and the public carrier
admission theorem give \textup{(PC6)}.  In every case expanded-derivation no-creation, universal
five-source provenance, recursive carrier--polar exhaustion, and
infinite-chain no-emergence give the displayed three dispositions.  These
operations take place on the presentation-generated affordable target spine,
so no ambient-space or arbitrary-\(X\) premise is introduced.
\end{proof}

\begin{theorem}[Halting and proof-carrying output on generated Noetherian cells]
\label{thm:halting-proof-carrying-output}
For every datum with an ordinary local germ emitted by an effective public
presentation, run the engine on the generated Noetherian target-aligned cells
produced by its generated finite admissible-carrier stratification.  The fair master schedule
halts with exactly one machine-verifiable
output accepted by \(\mathsf{Check}_{\rm eng}\): a regularity certificate
jointly excluding the thirty-one supports, or an independent explicit
singular mechanism with orbit, contact, and closure certificates.
\end{theorem}

\begin{proof}
By \cref{thm:public-target-aligned-premise-closure}, every analytic premise of
the schedule is a generated public value and every unfavorable value is
already an internal successor.  The generated Noetherian rank supplies finite
progress by \cref{thm:generated-noetherian-target-terminalization}.
Finite admissible-carrier stratification and thirty-one-support exhaustion enumerate
all possible target-bearing cells.  On each cell, the public algorithm
accountability theorem processes every finite constructor, approximation,
selector, total-graph value, carrier--polar descendant, and cumulative limit.
A target-null value closes the cell; a charged prefix above its generated
ceiling closes it by contradiction; and every nonzero remainder is the next
typed source successor.  The no-emergence trichotomy and fixed-atom
realization prevent an infinite execution from becoming a third outcome.

If no actual contact remains, regularity-certificate completeness gives
finite local exclusions and compactness gives a finite global package.  If
contact remains, the Noetherian singular-witness theorem gives a direct-orbit
proof or a public path program together with all effective moduli.  The fair
schedule enumerates all finite supports, cells, instruction words, and proof
rows, so it discovers the applicable finite package.
Verifier soundness accepts it, while global mutual exclusivity prevents
acceptance of both tags.  The argument uses only presentation-generated
Noetherian target-aligned structure.  Without the Noetherian rank,
algorithm accountability remains a semantic classification and this halting
conclusion does not follow.
\end{proof}

\begin{theorem}[Binary proof-carrying terminalization on generated Noetherian cells]
\label{thm:binary-proof-carrying-terminalization}
For every effective public presentation and every datum whose ordinary local
germ is emitted by the public local-solution graph, apply the native finite
admissible-carrier stratification and the target-aligned public branch-and-bound
construction.  On the resulting generated Noetherian target-aligned cells,
the compiler terminates with exactly one of
\[
 \operatorname{Regular}(C_{\rm reg})
 \qquad\text{or}\qquad
 \operatorname{Singular}(u_\ast,\mathfrak a,C_{\rm sing}).
\]
Acceptance by \(\mathsf{Check}_{\rm eng}\) proves avoidance of the complete
continuation target by every ordinary maximal branch in the regular case.  In
the singular case it proves that \(u_\ast\) is an independent ordinary solution of the public
equation with its declared data and physical ledger, and that it crosses the
cofinal target-shell family.  Every possible target-bearing contribution of a public algorithm is included
in the saturated five-source profile.  All successor, equality, recurrence,
compactification, graph-boundary, approximation, path-selection, and
prospective-atom values are internal compiler states.
\end{theorem}

\begin{proof}
Public target-aligned premise closure generates every analytic input to the
compiler.  Exhaustive target-aligned carrier stratification generates the
affordable
cells and routes every nonaffordable target-bearing failure to its least
five-source successor.  Target-aligned public-algorithm accountability then
places every algorithmic contribution in the same cumulative joint record
and saturated profile; nullity, finite over-budget closure, proper
succession, and fixed-atom realization exhaust its semantic values.  The
generated Noetherian rank makes every proper succession finite.  Thus global
binary semantics applies on every generated cell, and the halting theorem
supplies exactly one of the two output schemas.  The proof-carrying verifier
supplies their semantics, the explicit singular-mechanism equivalence gives
the singularity clause, and global mutual exclusivity excludes simultaneous
acceptance.  No theorem is invoked for an arbitrary ambient \(X\): all
quantifiers range over records, cells, and algorithms generated from the
public presentation and the verified Noetherian successor order.
\end{proof}

\begin{remark}[Four-layer engine organization]
\label{rem:four-layer-engine-organization}
The proof architecture is read in four layers.  The ownership layer consists
of the effective presentation, five-source provenance, graph-complete
closure, target quotient, and shell accountability.  The conversion layer
consists of canonical exact-support amplitude, support conversion, carrier
admission, transfer, and unresolved-progress conservation.  The
terminalization layer consists of positive-price duality, zero-price
resaturation, no-emergence, actual-prefix realization, and sourcewise
terminalization.  The effective-engine layer consists of the certificate
language, output verifier, fair scheduler, and generated-Noetherian halting
theorem.  In
particular, compact extrema, positive functionals, equality kernels, and
hypothetical atoms occur only as intermediate values in these layers.
\end{remark}

\begin{theorem}[Target-aligned public-presentation closure of the engine]
\label{thm:engine-target-aligned-public-closure}
Every state, transition, certificate, and accepted output of the binary engine
belongs to \(\mathsf{Proof}_{\Sigma}(\mathfrak P)\) and retains the canonical
public-presentation dependency record.  More precisely:
\begin{enumerate}[label=\textup{(O\arabic*)}]
\item an ownership state is the target quotient of a cumulative same-origin
  public joint record, paired with a canonical shell covector and tagged by
  the exact word supplied by \(\operatorname{Prov}\);
\item a conversion or charging transition is a restriction of
  \eqref{eq:target-local-fact-normal-form} and
  \eqref{eq:target-aligned-charge-rule} to the active stratum, with every gap,
  event index, null term, graph defect, and support-transfer output retained;
\item a terminalization or realization transition acts on the same cumulative
  joint record through prospective backward factorization, carrier--polar
  exhaustion, and the actual-prefix tree, while preserving origin polarity;
\item a regularity or singularity certificate is a finite restriction of
  these generated rows and records the public relation indices, shell
  covectors, exact supports, origins, physical owners, and graph enclosures
  used by every inference.
\end{enumerate}
Consequently the engine has no parallel source space, target test, physical
budget, realization semantics, or proof cone.  A proposed transition or
certificate missing one of the fields in \textup{(O1)}--\textup{(O4)} is
rejected as ill typed.
\end{theorem}

\begin{proof}
The public-presentation input contract and
\cref{thm:closed-target-aligned-proof-discipline} generate the ambient proof
algebra.  Canonical shell accountability, universal five-source provenance,
and the presentation-generated prospective domain give \textup{(O1)}.  The
definition of the exact-support conversion ledger, support-transfer
conservation, and no-double-charging give \textup{(O2)}.  The cumulative joint
record, exact backward factorization, recursive carrier--polar exhaustion,
and the actual finite-prefix tree give \textup{(O3)}.  Finally,
\(\mathsf{Check}_{\rm eng}\) checks precisely the finite transcripts and
metadata listed in \textup{(O4)}.  Induction over the engine transition list
therefore preserves all fields.  The only terminal constructors are the two
checked output schemas, so the conclusion also holds for accepted outputs.
\end{proof}

\begin{definition}[Continuum extension functor]
\label{def:continuum-extension-functor}
Let \(\mathfrak C\) send a finite cylinder record to its evaluation state,
a finite current to its distributional cylinder pairings, a finite target
probe to its pullback on \(X_{\rm gc}\), and a constructor word to the closure
of its typed graph.  It sends a convergent family to its sourcewise
graph-complete closure and sends stopping, restriction, signed addition,
composition, and physical-ledger addition to the corresponding closed graph
operations.
\end{definition}

\begin{theorem}[Part-I preservation under continuum completion]
\label{thm:part-i-continuum-preservation}
The functor \(\mathfrak C\) preserves the five source tags, constructor
composition, signed sums, restriction, first-contact stopping, and additive
physical ledgers.  Quotient, chart, and terminal operations append only
\(\Abs\), \(\Lift\), and \(\Bdry\), respectively.  For every cylinder record
\(R\),
\begin{equation}
 \label{eq:continuum-functor-preservation}
 \Anc(\mathfrak C R)
 =\operatorname{cl}_{\rm src}\Anc(R),
 \qquad
 \Pi_\Sigma\mathfrak C R
 =\mathfrak C(\Pi_\Sigma R).
\end{equation}
Consequently sourcewise limit closure creates no new primitive source and the
target-null kernel commutes with continuum completion.  On a finite state
space, \(\mathfrak C\) is the identity and prospective contact accountability
reduces to the Part-I first-hit identity.
\end{theorem}

\begin{proof}
Each assertion is first checked on an atomic constructor.  Native currents
retain their assigned source.  A quotient changes only provenance and appends
\(\Abs\); a chart or augmentation appends \(\Lift\); a terminal restriction
or exterior trace appends \(\Bdry\).  The graph of a composite is the
composition of the typed graphs, so induction on constructor length proves
the assertions for finite records.  Sourcewise closure is taken separately
inside each tagged graph, proving the first identity in
\eqref{eq:continuum-functor-preservation}.  Target probes are cylinder
functionals and therefore commute with pointwise graph closure, proving the
second identity.  Their common kernel is consequently closed under
\(\mathfrak C\).  Finite state spaces have no additional analytic boundary:
their evaluation image is already closed, so the construction returns the
original finite algebra and its first-hit decomposition.
\end{proof}

\begin{proposition}[Public-presentation input contract]
\label{prop:public-presentation-contract}
After \(\mathsf{Eq}\), its public datum class, and its native formal carrier
identity are fixed, every object used by the target-reachability track is
generated functorially: the test core by rational enumeration, the solution
and continuation graphs by satisfaction of the weak relation, the operation
list by the typed syntax, the topology by bounded evaluation coordinates, the
target by failed continuation, the source words by normal-form compilation,
and the physical measure by the positive Jordan part of the signed carrier
identity.  A PDE specialization therefore supplies no additional
problem-specific regularity input, proof object, favorable constant, or
target-relative lemma.
\end{proposition}

\begin{proof}
Each construction is the one stated in items \textup{(P1)}--\textup{(P8)}.
All subsequent objects are defined by graph closure, target quotient, Jordan
decomposition, normalized shell formation, compact activity lifting, or the
successor operator.  These operations consume only the preceding generated
objects.  Inspection of the dependency chain shows that none has an input
slot for any of the forbidden conclusions listed above.
\end{proof}

\begin{definition}[Target-blind public input]
\label{def:target-blind-public-input}
Let \(\mathsf{Pre}(\mathsf{Eq})\) be the signature formed before the
continuation target, target shells, normalized bad windows, source quotients,
or physical-price strata are constructed.  It contains the displayed
initial--boundary value problem, its declared local regularity class, weak
test core, optional standard local solution theorem, native algebraic
symmetries and scalings, and formal multiplier identities.  It may use a fixed universal
library of target-independent results such as distribution theory, integration
by parts, Sobolev-space definitions, Jordan decomposition, compactness of
countable products of compact metric spaces, weak compactness of probability
measures on such spaces, Archimedean ordered-algebra separation, and the
spectral theorem.

A fact is an admissible public input precisely when it is expressible in
\(\mathsf{Pre}(\mathsf{Eq})\) and remains unchanged if the continuation target
and all target observables are deleted.  In particular, the following are not
public inputs:
\begin{equation}
 \label{eq:forbidden-regularity-inputs}
 \begin{gathered}
 \text{compactness of a target-approaching sequence, an epsilon-regularity or
 continuation conclusion,}\\
 \text{a target price gap, a blow-up Liouville theorem, absence of a
 scale-soliton,}\\
 \text{vanishing of a target quotient, or any decomposition selected after
 singular shells are formed.}
 \end{gathered}
\end{equation}
All objects in \eqref{eq:forbidden-regularity-inputs} must be generated and
processed by the successor algebra.
\end{definition}

\begin{theorem}[Syntactic no-smuggling criterion]
\label{thm:syntactic-no-smuggling}
The dependency graph of the target-reachability compiler has exactly one edge
from user data: the inclusion of
\(\mathsf{Pre}(\mathsf{Eq})\) into the canonical presentation.  Every target,
source, price, compactness, equality, and recurrence object lies strictly
downstream of that edge.  Hence no target-relative conclusion can enter as a
public input.  Adding any fact in
\eqref{eq:forbidden-regularity-inputs} changes the signature and does not count
as an evaluation of the original presentation.
\end{theorem}

\begin{proof}
The target is generated from the continuation graph, the shells from the
target distance in \(X_{\rm ev}\), the source cells from normal-form syntax,
the price from the native carrier, and the recurrent equations from the
successor fixed point.  These constructions form a directed acyclic
dependency chain beginning at \(\mathsf{Pre}(\mathsf{Eq})\).  None has a
backward input slot.  Every forbidden fact mentions a downstream object and
therefore cannot be a term of the pre-target signature.
\end{proof}

\begin{theorem}[Closed target-aligned proof discipline]
\label{thm:closed-target-aligned-proof-discipline}
Every target, source, invariant, physical-price, or reachability assertion in
the continuum structural calculus is a statement in the following generated
proof algebra:
\[
 \mathsf{Proof}_{\Sigma}(\mathfrak P):=\operatorname{Cl}_{\rm fin,lim}
 \left(
   \mathsf{Pre}(\mathsf{Eq}),\;
   \text{public weak and total graphs},\;
   \text{target-shell covectors},\;
   \text{five-source constructors},\;
   \text{joint conditional increments},\;
   \text{physical carrier rows}
 \right).
\]
Here \(\operatorname{Cl}_{\rm fin,lim}\) means finite typed construction,
source-preserving graph completion, and monotone completion of finite
same-origin stopped records.  In particular:
\begin{enumerate}[label=\textup{(P\arabic*)}]
\item every premise is a public-presentation coordinate or the conclusion of
  an earlier typed construction in \(\mathsf{Proof}_{\Sigma}(\mathfrak P)\);
\item every target-visible term has a complete joint frontier before source
  projection and is assigned one nonempty support
  \(A\subseteq\{\Geom,\Caus,\Abs,\Lift,\Bdry\}\);
\item every positive target coordinate has a finite same-origin detecting
  occurrence, while a zero coordinate deletes its target-incidence edge;
\item every physical assertion is the evaluation of the unique physical
  carrier coordinate of that occurrence, with equality re-entering the same
  joint source word rather than creating a new conclusion.
\end{enumerate}
Thus no theorem in the calculus takes a whole-orbit estimate, a target-side
compactness statement, a continuation conclusion, a rigidity assertion, or a
problem-specific certificate as an input.  Such a statement can occur only
when it has first been compiled into a target-aligned coordinate of the
public presentation.
\end{theorem}

\begin{proof}
The input contract and \cref{thm:syntactic-no-smuggling} give
\textup{(P1)}.  Structural induction over the closed constructor grammar
gives \textup{(P2)}: joining precedes projection, while quotient, lift, and
boundary constructors append only their declared letters.  The finite-record
zero/positive equivalence
\eqref{eq:actual-source-zero-equivalence}--
\eqref{eq:actual-source-positive-equivalence} gives \textup{(P3)}; a
target-visible limit has a detecting finite prefix and therefore re-enters
the same source cell.  Finally,
\cref{thm:actual-source-price-equality-decision} gives \textup{(P4)}.
These are precisely the permitted proof rules, so induction on a derivation
in the displayed closure proves the claim.
\end{proof}

\begin{corollary}[Uniform proof rule for the thirty-one target supports]
\label{cor:uniform-thirty-one-support-proof-rule}
For each nonempty support \(A\subseteq\Channels\), every proof concerning
the contribution of \(A\) is obtained by restricting
\(\mathsf{Proof}_{\Sigma}(\mathfrak P)\) to its complete source word,
forming its joint detector before any coordinate is discarded, and applying
the zero/positive and carrier/equality rules of
\cref{thm:closed-target-aligned-proof-discipline}.  The resulting statement
is either target absence on that cell or an actual finite public-presentation
occurrence carrying that support and its physical charge.  Since the
nonempty supports are exactly the thirty-one nonzero subsets of
\(\Channels\), this is a proof rule for each structural source cell, not a
separate global PDE argument.
\end{corollary}

\begin{proof}
Apply \textup{(P2)}--\textup{(P4)} of
\cref{thm:closed-target-aligned-proof-discipline} to the fixed word \(A\).
The five-constructor grammar has no other nonempty word.
\end{proof}

\begin{corollary}[Textbook-only specialization interface]
\label{cor:textbook-only-specialization}
The external specialization datum ends with
\(\mathsf{Pre}(\mathsf{Eq})\).  Its mathematical content consists of the
displayed equation, its native formal carrier identity, and any cited
target-blind local theorem already attached to its standard presentation.
In particular, no later
row has an argument slot for a PDE-specific compactness proof, shell estimate,
price bound, equality analysis, continuation criterion, or prescribed value
of \(Z_n\).  The compiler itself generates the total graphs, the ordered cones,
their finite proof identities, and their complementary positive functionals.
\end{corollary}

\begin{proof}
The first sentence is \cref{def:target-blind-public-input}.  The dependency
statement follows from \cref{thm:syntactic-no-smuggling}.  Totalization and
the successor relation generate both values of every downstream relation;
the proof/model alternative is
\cref{thm:ordered-proof-or-model}.  None of these constructions
adds an input edge to the dependency graph.
\end{proof}

\begin{proposition}[Non-vacuity of the total solution graph]
\label{prop:nonvacuous-solution-graph}
For every datum, exactly one of the following occurs: the total local-solution
fiber contains at least one maximal orbit, or it contains the nonzero form
defect \(\bot_{\rm loc}\).  In particular, no reachability or regularity
functional defined from the presentation assigns the value zero because the
orbit fiber is empty.
\end{proposition}

\begin{proof}
The total graph is the disjoint union of its ordinary solution values and its
explicit failure value.  The empty ordinary fiber therefore selects
\(\bot_{\rm loc}\).  Subsequent quotients retain that record unless its target
projection is proved to vanish.
\end{proof}

\begin{figure}[tbp]
\centering
\resizebox{.96\linewidth}{!}{%
\begin{tikzpicture}[fw diagram,node distance=9mm and 13mm]
  \node[fw source,text width=3.0cm] (pde)
    {fixed equation and\\initial state};
  \node[fw provenance,text width=2.8cm,right=of pde] (orbit)
    {ordinary maximal branch\\\(u:[0,T_*)\to X\)};
  \node[fw use,text width=3.0cm,right=of orbit] (events)
    {pre-contact target events\\\(Q\in\mathfrak Q_u\)};
  \node[fw geom,text width=3.1cm,right=of events] (physical)
    {event restriction\\\(\mu_u|_Q\)};
  \node[fw terminal,text width=3.2cm,right=of physical] (reach)
    {single ceiling\\\(\mu_u(\mathfrak Q_u)\le\Budget\)};
  \draw[fw arrow] (pde)--(orbit);
  \draw[fw arrow] (orbit)--(events);
  \draw[fw arrow] (events)--(physical);
  \draw[fw arrow] (physical)--(reach);
\end{tikzpicture}%
}
\caption{One physical-budget architecture.  The equation generates one
countably additive orbit measure, while target approach selects pre-contact
events in its measurable event algebra.  A selected event pays for target
progress only after the target-local carrier factorization proves the
charging rule; the measure's total mass is bounded by \(\Budget\).}
\label{fig:physical-budget-architecture}
\end{figure}

\begin{definition}[Generated admissible carrier and budget]
\label{def:physical-budget}
Let \(u\) be an ordinary value of the total solution graph.  Enumerate the
target-blind carrier candidates generated by \(\mathcal K\).  A candidate is
\emph{admissible} on a public cell only if a finite derivation from the weak
equation verifies all of the following before any target coordinate is read:
an exact signed balance, a lower-bounded storage functional \(K_u\), a
nonnegative localizable production measure, cancellation of internal faces,
and a Jordan decomposition of genuine exterior work.  The fair compiler
creates one carrier-tagged stratum for each admitted derivation.  A proof
package works inside one such stratum and may charge an event only to its
tagged carrier, so different candidate identities are never added together
or used to charge the same occurrence twice.  Failed candidates retain their
first failed checker row as total-graph successors rather than finite-budget
values.

On an admissible ordinary branch, signed gluing gives
\begin{equation}
 \label{eq:generated-physical-carrier-identity}
 K_u(t)+\mu_u((s,t])+W_u^-((s,t])
 =K_u(s)+W_u^+((s,t]),
 \qquad 0\le s<t<T_*(u),
\end{equation}
where \(W_u^\pm\) are the Jordan parts of true exterior work and
\(\mu_u\) is the sum, with the native coefficients of the identity, of all
generated nonnegative production and balance-defect measures.  Internal
transport, chart, gauge, and interface currents have already cancelled and
do not enter \(\mu_u\).  Every surviving summand retains its five-source
ancestry.  This orbitwise construction bounds the available physical measure;
it does not yet assign any portion of that measure to target progress.  Such
an assignment is made only by
\cref{def:target-aligned-physical-charge} on an active branch.

If \(K_u\ge K_->-\infty\) and
\(W_u^+([0,T_*(u)))<\infty\), define the single orbit budget
\begin{equation}
 \label{eq:generated-physical-budget}
 \Budget(u_0,W^+)
 :=K_u(0)-K_-+W_u^+([0,T_*(u)))<\infty.
\end{equation}
Then \eqref{eq:generated-physical-carrier-identity} derives
\begin{equation}
 \label{eq:physical-control-bound}
 \mu_u(\mathfrak Q_u)\le \Budget(u_0,W^+).
\end{equation}
Failure of the carrier identity, nonnegativity, the lower bound, or finite
exterior input is the corresponding total-graph value.  Its target-visible
part enters the activity lift and ordered successor algebra.  A finite
ordered identity may exclude that value; otherwise it remains in the positive
structural model.  It is neither discarded nor misread as finite-budget
regularity.  The finite-budget reachability theorems apply only to the
ordinary branch on which \eqref{eq:physical-control-bound} has been derived.

The measure is localizable.  For pairwise disjoint
\(Q_j\in\mathcal A_u\),
\begin{equation}
 \label{eq:physical-countable-additivity}
 \sum_j\mu_u(Q_j)
 =\mu_u\!\left(\bigsqcup_jQ_j\right)
 \le \Budget.
\end{equation}
No target-dependent expenditure may be appended.  An energy, Hamiltonian,
entropy, action, commutator, or localized identity refines \(\mu_u\) only
after it has been
proved to be an algebraic decomposition of
\eqref{eq:generated-physical-carrier-identity}.  If the basic carrier assigns
finite price to a singular cascade, that cascade proceeds to the zero-cost
successor calculus; the budget is not strengthened after the fact.

For an energy-dissipative presentation this construction gives
\(K_u=\Energy(u)\) and
\(\Budget=\Energy(u_0)-E_-+W_u^+\).  Here ``finite energy budget'' means
finiteness of the cumulative nonnegative expenditure controlled by that
identity, not merely finiteness of the stored value \(\Energy(u(t))\).
Other equations may generate a different admissible Lyapunov carrier, but
the same checker rules apply and the carrier cannot be chosen in response to
the singular target.
\end{definition}

\begin{theorem}[Public carrier admission and no budget smuggling]
\label{thm:public-carrier-admission}
The carrier compiler fairly enumerates every finite multiplier derivation in
\(\mathcal K\) and decides each finite admission transcript.  Every admitted
budget is therefore a consequence of the target-blind public equation,
boundary conditions, datum, and exterior-input graph.  Conversely, an
energy, entropy, action, norm, price schedule, or coercive estimate that lacks
such a transcript cannot occur as a budget premise or authorize target
charge.  Failure to admit a proposed carrier is retained as its typed public
graph value.
\end{theorem}

\begin{proof}
Order the finite multiplier words, integration-by-parts identities, sign
decompositions, lower-bound proofs, and rational exterior-work bounds by
their public codes and dovetail their checks.  Exact symbolic equality,
positivity in the generated quadratic module, lower-bound comparison, and
finite-cylinder work bounds are finite checker operations.  An accepted
word supplies every field of \cref{def:physical-budget}, so its ceiling
follows from \eqref{eq:generated-physical-carrier-identity}.  A rejected word
has a least failed row in the corresponding total graph.  The target and its
shell covectors do not occur in this enumeration, and the proof-package type
has no field for an untraced estimate.  Hence neither target-dependent choice
nor a user-supplied favorable certificate can enter the carrier ledger.
\end{proof}

\begin{theorem}[Universal physical-quantity ownership]
\label{thm:universal-physical-quantity-ownership}
Let \(F(u)\) be any scalar, vector, tensor, measure, flux, action, entropy,
energy, enstrophy, moment, capacity, frequency, scale, or localized quantity
generated from the public equation.  Every use of \(F\) in target
accounting is evaluated through the complete expanded derivation of \(F\)
and its balance graph.  Exactly one of the following holds on a target-active
cell:
\begin{enumerate}[label=\textup{(O\arabic*)}]
\item its target projection vanishes;
\item its owner-tagged positive part is algebraically dominated by the single
  measure \(\mu_u\) through a finite identity in
  \(\mathscr D_{\mathfrak b}^{\rm alg}\);
\item its owner belongs to
  \(\mathscr D_{\mathfrak b}^{\Phys}\setminus
    \mathscr D_{\mathfrak b}^{\rm alg}\); the approximating rows and all
  target-shell errors form the cone-boundary successor, and no charge is
  declared;
\item its owner lies outside \(\mathscr D_{\mathfrak b}^{\Phys}\); finite-cylinder
  separation produces the carrier-polar successor with the same five-source
  ancestry.
\end{enumerate}
Finite, conserved, monotone, coercive, scale-invariant, or summable values of
the constituents do not imply any of these conclusions for a composite.
They are predicates inside its expanded joint record.  In particular, two
physical quantities may cancel or amplify only through their signed joint
balance, and that joint balance is processed before either positive parts or
normalization are taken.  The theorem creates no second budget.
\end{theorem}

\begin{proof}
Form \(\operatorname{Exp}(F)\), retain every intermediate and owner-tagged
restriction, and apply
\cref{thm:continuum-composition-no-creation}.  Target-nullity gives
\textup{(O1)}.  Algebraic membership of the resulting target owner gives
\textup{(O2)}.  Closure-only membership gives the total cone-boundary graph
in \textup{(O3)}.  On the complement, finite-cylinder separation and
\cref{thm:recursive-carrier-polar-exhaustion} give \textup{(O4)}.  Signed
gluing precedes the cone test, so cancellation is retained exactly.  Since
only domination by \(\mu_u\) authorizes charge, no auxiliary physical
quantity becomes another budget.
\end{proof}

\begin{example}[Energy--work budget]
On the zero-defect branch of the lower-energy, dissipation-balance, and
finite-work rows, an energy \(\Energy\), lower value \(E_-\), and nonnegative
dissipation measure \(\mu_{\mathrm{diss}}\) satisfy
\[
 \Energy(u(t))+\mu_{\mathrm{diss}}([0,t])
 \le \Energy(u_0)+W^+([0,t]),
\]
and \(W^+([0,\infty))<\infty\).  Then
\[
 \mu_{\mathrm{diss}}([0,\infty))
 \le \Energy(u_0)-E_-+W^+([0,\infty)).
\]
Taking \(\mu_u=\mu_{\mathrm{diss}}\) gives the single budget
\[
 \Budget=\Energy(u_0)-E_-+W^+([0,\infty)).
\]
If this physical measure assigns finite total price to a singular cascade,
the cascade enters the zero-price successor calculus.  No stronger measure or
second budget is introduced after the target has been inspected.
\end{example}

\begin{figure}[tbp]
\centering
\resizebox{\linewidth}{!}{%
\begin{tikzpicture}[fw diagram,node distance=9mm and 13mm]
  \node[fw source,text width=2.8cm] (initial)
    {initial physical data\\and energy level};
  \node[fw provenance,text width=2.8cm,right=of initial] (input)
    {positive forcing and\\boundary work \(W^+\)};
  \node[fw geom,text width=3.1cm,right=of input] (measure)
    {nonnegative expenditure measure\\\(\mu_u\)};
  \node[fw use,text width=3.0cm,right=of measure] (events)
    {pairwise disjoint\\target events \(Q_j\)};
  \node[fw terminal,text width=3.1cm,right=of events] (bound)
    {\(\sum_j\mu_u(Q_j)\le\Budget\)};
  \draw[fw arrow] (initial)--(measure);
  \draw[fw arrow] (input)--(measure);
  \draw[fw arrow] (measure)--(events);
  \draw[fw arrow] (events)--(bound);
\end{tikzpicture}
}%
\caption{Physical expenditure accounting.  The initial state and legitimate
positive input bound one countably additive measure.  Pairwise disjoint
target events have disjoint restrictions of that measure.  Only restrictions
admitted by the local target-aligned charging rule pay for shell progress.}
\label{fig:physical-budget}
\end{figure}

\begin{definition}[Fixed-Caus gradient-flow presentation]
\label{def:fixed-caus-presentation}
Let
\(\Vspace\hookrightarrow\Hspace\hookrightarrow\Vspace^*\) be a Gelfand
triple.  A fixed-Caus gradient-flow presentation consists of a differentiable
energy \(\Energy\), a symmetric nonnegative mobility
\(\mathbb K(u):\Vspace^*\to\Vspace\), a fixed constitutive transport
\(A_{\Caus}(u)\in\Vspace^*\), and an exterior input
\(f_{\Bdry}\in\Vspace^*\), with equation
\begin{equation}
 \label{eq:fixed-caus-flow}
 \dot u+A_{\Caus}(u)
 =-\mathbb K(u)D\Energy(u)+f_{\Bdry}.
\end{equation}
The Caus term is energy-skew when
\begin{equation}
 \label{eq:caus-energy-skew}
 \langle D\Energy(u),A_{\Caus}(u)\rangle=0.
\end{equation}
Write
\[
 \bigl\|\mathbb K(u)^{1/2}\xi\bigr\|^2
 :=\langle \xi,\mathbb K(u)\xi\rangle
\]
for \(\xi\in\Vspace^*\); this notation does not assume a bounded inverse or
a globally defined square root on \(\Vspace^*\).
Here \(\Geom\) is the mobility current
\(-\mathbb K D\Energy\), \(\Caus\) is the declared transport,
and \(\Bdry\) contains the exterior input and boundary flux.  Quotients and
symmetry reductions applied to this orbit have \(\Abs\) ancestry; charts,
blow-up rescalings, and augmented variables have \(\Lift\) ancestry.
\end{definition}

\begin{proposition}[Fixed-Caus expenditure bound]
\label{prop:fixed-caus-expenditure-bound}
Let \(u\) satisfy \cref{eq:fixed-caus-flow,eq:caus-energy-skew} on
\([0,T]\), with the chain rule for \(\Energy(u)\).  Then
\begin{equation}
 \label{eq:fixed-caus-energy-identity}
 \frac{\dd}{\dd t}\Energy(u(t))
 =-\bigl\|\mathbb K(u)^{1/2}D\Energy(u)\bigr\|^2
  +\langle D\Energy(u),f_{\Bdry}\rangle.
\end{equation}
If \(\Energy\ge E_-\), then
\begin{equation}
 \label{eq:fixed-caus-action-bound}
 \int_0^T
 \bigl\|\mathbb K(u)^{1/2}D\Energy(u)\bigr\|^2\dd t
 \le
 \Energy(u_0)-E_-
 +\int_0^T\langle D\Energy(u),f_{\Bdry}\rangle_+\dd t.
\end{equation}
For a weak solution with an energy inequality, the same statement holds with
an additional nonnegative energy-defect measure on the left.
\end{proposition}

\begin{proof}
Pair \eqref{eq:fixed-caus-flow} with \(D\Energy(u)\).  The chain rule gives
the left side of \eqref{eq:fixed-caus-energy-identity}, the Caus pairing
vanishes by \eqref{eq:caus-energy-skew}, and positivity of \(\mathbb K\)
gives the squared norm.  Integrate, discard the terminal energy above
\(E_-\), and replace exterior work by its positive part.  The weak form is
the increment formula for the energy inequality; its defect is nonnegative.
\end{proof}

\section{The five-source orbit algebra}
\label{sec:source-algebra}

\begin{definition}[Continuum source alphabet]
\label{def:source-alphabet}
The source alphabet is
\[
 \Channels=\{\Geom,\Caus,\Abs,\Lift,\Bdry\}.
\]
A raw continuum record is an orbit occurrence together with a finite ancestry
word \(\Anc(r)\in\Channels^*\).  Its source support
\(\Src(r)\subseteq\Channels\) is the set of letters appearing in that word.
The primitive and closure rules are:
\begin{enumerate}[label=\textup{(S\arabic*)}]
\item symmetric mobility, Dirichlet-form, reversible, and dissipative current
records start in \(\Geom\);
\item after the canonical native form split, the remaining oriented bulk
evolution current starts in \(\Caus\); this includes
fixed transport, reaction, nonvariational forcing, memory, and a raw bulk
current whose attempted form realization is nonordinary;
\item quotient, symmetry reduction, gauge removal, and coarse variables append
\(\Abs\);
\item charts, renormalizations, blow-up variables, auxiliary-state extensions,
and Young- or defect-variable lifts append \(\Lift\);
\item boundary traces, exterior fluxes, killing, terminal values, and tails
read at the boundary of a compactification append \(\Bdry\);
\item finite algebraic combination, restriction, and concatenation take the
union of the input source supports.
\end{enumerate}
The symbol \(\Res\) is absent from this source alphabet.  It is introduced
only as the target-null coordinate of \cref{def:target-null-residual}.
\end{definition}

\begin{remark}[Channel and source terminology]
The channels \(\Geom\) and \(\Abs\) carry intrinsic reusable structure,
while \(\Caus\), \(\Lift\), and \(\Bdry\) orient, reorganize, and read out
that structure.  A \emph{source cell} is a provenance cell in the complete
five-channel normal form; this terminology does not promote the three derived
channels to independent creators of intrinsic structure.
\end{remark}

\begin{figure}[tbp]
\centering
\resizebox{.98\linewidth}{!}{%
\begin{tikzpicture}[fw diagram,node distance=7mm and 9mm]
  \node[fw geom,text width=2.45cm] (g)
    {\(\Geom\)\\mobility, form, dissipation};
  \node[fw use,text width=2.45cm,right=of g] (c)
    {\(\Caus\)\\fixed directed transport};
  \node[fw provenance,text width=2.45cm,right=of c] (a)
    {\(\Abs\)\\quotient and gauge reduction};
  \node[fw source,text width=2.45cm,right=of a] (l)
    {\(\Lift\)\\chart and augmented state};
  \node[fw terminal,text width=2.45cm,right=of l] (b)
    {\(\Bdry\)\\trace, flux, killing, tail};
  \node[fw box,text width=3.4cm,below=15mm of a] (word)
    {composite record\\finite ancestry word in \(\Channels^*\)};
  \draw[fw arrow] (g)--(word);
  \draw[fw arrow] (c)--(word);
  \draw[fw arrow] (a)--(word);
  \draw[fw arrow] (l)--(word);
  \draw[fw arrow] (b)--(word);
\end{tikzpicture}%
}
\caption{The five continuum source roles.  The canonical native form split
routes its positive symmetric output to Geom and its complete oriented bulk
remainder to Caus.  Abs and Lift record quotient and presentation currents,
and Bdry records interface, exterior, and terminal currents.  Composite
records retain the union of their ancestries.}
\label{fig:five-continuum-sources}
\end{figure}

\begin{remark}[Forcing and constraints]
The source compiler routes a volume force according to its constitutive
role.  A fixed directed body force may be \(\Caus\), while exterior work or a
boundary supply is \(\Bdry\).  A Lagrange multiplier implementing a quotient
constraint has \(\Abs\) ancestry.  The five labels are assigned from the
declared weak equation, not from the sign of a later estimate.
\end{remark}

\begin{definition}[Operational weak-PDE signature]
\label{def:operational-weak-pde-signature}
An operational weak-PDE signature is formed downstream from the public
presentation.  It consists of a Gelfand triple
\(V\hookrightarrow H\hookrightarrow V^*\), a separating test core
\(\mathcal D\subset V\), an initial datum, the continuation target \(\Sigma\)
generated by \cref{def:singular-target}, and a weak
evolution relation
\begin{equation}
 \label{eq:operational-weak-evolution}
 \langle \dot u,\varphi\rangle_{V^*,V}
 =\mathcal A(u;\varphi)+\mathcal F_{\rm ext}(u;\varphi),
 \qquad \varphi\in\mathcal D .
\end{equation}
The native spatial relation is entered before any realization property is
asserted.  Its totalized realization relation has the following ordinary
values.
\begin{enumerate}[label=\textup{(W\arabic*)}]
\item A fixed-Caus gradient realization:
\[
 \mathcal A(u;\varphi)
 =-\langle \mathbb K(u)D\Energy(u),\varphi\rangle
  -\langle A_{\Caus}(u),\varphi\rangle ,
\]
with \(\mathbb K(u)\) symmetric and nonnegative.
\item A sector-form realization on a real \(L^2\)-space:
\[
 \mathcal E=\mathcal E^{\rm s}+\mathcal E^{\rm a},
 \qquad
 \mathcal E^{\rm s}(v,w)
 =\tfrac12\{\mathcal E(v,w)+\mathcal E(w,v)\},
 \quad
 \mathcal E^{\rm a}(v,w)
 =\tfrac12\{\mathcal E(v,w)-\mathcal E(w,v)\},
\]
where \(\mathcal E^{\rm s}\) is a Dirichlet form and
\(\mathcal E^{\rm a}\) satisfies the sector bound.
\end{enumerate}
If neither ordinary value exists, the realization relation returns
\(\bot_{\rm nat}\).  The raw fixed bulk current remains a \(\Caus\) occurrence
and the failed attempted Geom--Caus split is a form-defect record with support
\(\{\Geom,\Caus\}\).  Its target-active quotient is evaluated by the form
row.  Thus absence of a gradient, Dirichlet, positivity, or sector
realization is admitted as structure rather than excluded from the signature.
The signature also declares the symmetry, quotient, chart, auxiliary-state,
and trace operations that may be applied to
\eqref{eq:operational-weak-evolution}.  It contains no assertion of
compactness, regularity, target avoidance, positive physical price, or
continuation beyond the declared graph.
\end{definition}

\begin{definition}[Continuum source compiler]
\label{def:continuum-source-compiler}
Let \(\mathsf{Syn}(\mathfrak P)\) be the smallest typed syntax containing the
native terms of \eqref{eq:operational-weak-evolution} and closed under every
finite algebraic or differential operation appearing in its expression
graph, signed sum, composition, restriction, localization, polarization,
commutator, quotient, chart, state augmentation, trace, total-graph
completion, and sourcewise weak limit.  The source compiler
\[
 \mathsf C_\mathfrak P:
 \mathsf{Syn}(\mathfrak P)\longrightarrow
 2^{\Channels}\setminus\{\varnothing\}
\]
is defined recursively by
\begin{align}
 \mathsf C_\mathfrak P(-\mathbb K D\Energy)
   &=\{\Geom\},&
 \mathsf C_\mathfrak P(\mathcal E^{\rm s}_{\rm local/jump})
   &=\{\Geom\},\nonumber\\
 \mathsf C_\mathfrak P(-A_{\Caus})
   &=\{\Caus\},&
 \mathsf C_\mathfrak P(\mathcal E^{\rm a})
   &=\{\Caus\},\label{eq:operational-source-compiler}\\
 \mathsf C_\mathfrak P(\mathcal F_{\rm ext})
   &=\{\Bdry\},&
 \mathsf C_\mathfrak P(\bot_{\rm nat})
   &=\{\Geom,\Caus\}.\nonumber
\end{align}
Here \(A_{\Caus}\) denotes the complete oriented bulk remainder returned by
the canonical native form split; on the native failure
branch it denotes the raw bulk current.  It is not restricted to a
first-order or antisymmetric operator.
Finite sums take unions.  Quotients append \(\Abs\); charts and auxiliary
variables append \(\Lift\); traces, killing, terminal values, and
compactification boundaries append \(\Bdry\).  Restriction, concatenation,
and sourcewise weak closure retain support.  A constraint multiplier is
compiled through the operation that generates it: quotient constraints append
\(\Abs\), while an exterior constraint supply appends \(\Bdry\).
\end{definition}

\begin{definition}[Primitive dynamical-role partition]
\label{def:primitive-dynamical-role-partition}
Expand the displayed weak equation into its finite expression graph before
forming target shells.  A displayed integral, series, or nonlocal operator
is one bulk leaf unless an exact public identity expands it.  Exact
decomposition identities are applied in their fixed public enumeration, so
different valid presentations remain distinct ancestry records.  A
\emph{primitive dynamical occurrence} is a leaf
together with its first edge that contributes to a weak current pairing.
The totalized native form-split graph of
\cref{def:operational-weak-pde-signature} is then applied.  Its ordinary
value may return both a positive symmetric occurrence and an oriented
remainder; its nonordinary value returns the raw bulk current and the typed
form defect.  The following ordered role test assigns each resulting
occurrence.
\begin{enumerate}[label=\textup{(R\arabic*)}]
\item A generated positive symmetric form, mobility, reversible quadratic
  form, carré du champ, or its polarization is \(\Geom\).
\item Every oriented remainder and every raw bulk evolution relation on the
  nonordinary split branch is \(\Caus\).  This is the complement of
  \textup{(R1)} in the canonical native form split, not an
  assumption that the current is antisymmetric or has a gradient
  representation.
\item An identification of states, constraints, gauges, symmetries, or
  coarse variables appends \(\Abs\).
\item A change or enlargement of the state presentation---including a chart,
  localization state, history state, stochastic path, modulation variable,
  blow-up variable, or weak-product variable---appends \(\Lift\).
\item A spatial, temporal, exterior, killed, terminal, or compactification
  interface appends \(\Bdry\).
\end{enumerate}
Scalar coefficients, datum values, and exact storage terms carry no
independent source: they acquire the ancestry of the current-producing edge
in which they occur.  A path meeting several roles retains their union.  The
test is ordered only to make the primitive assignment disjoint; it does not
erase composite ancestry.
\end{definition}

\begin{lemma}[Exhaustion of primitive current paths]
\label{lem:primitive-current-edge-exhaustion}
Every path from a resulting primitive occurrence of a canonical physical
presentation to a weak current pairing has one of the following forms:
\begin{enumerate}[label=\textup{(E\arabic*)}]
\item an interior bulk edge, assigned by exactly one of
  \textup{(R1)} and \textup{(R2)}, followed by zero or more quotient, lift,
  or boundary edges;
\item a constraint, gauge-reaction, or identification current, assigned
  \(\Abs\), retaining every upstream bulk role;
\item a chart velocity, auxiliary-state, modulation, or presentation-defect
  current, assigned \(\Lift\), retaining every upstream role;
\item a prescribed exterior, trace, or terminal current, assigned
  \(\Bdry\), retaining every upstream role.
\end{enumerate}
Thus every primitive current occurrence has one head in \(\Channels\) and a
nonempty support containing that head.  An expression leaf with no path to a
weak current pairing contributes zero to every prospective shell current.
\end{lemma}

\begin{proof}
An edge of the weak expression graph either changes the evolving state in
the interior, identifies states, changes the state presentation, or reads an
interface value.  These are respectively the bulk, quotient, lift, and
boundary cases in
\cref{def:primitive-dynamical-role-partition}.  The cases are disjoint after
the ordered test.

For a bulk edge, apply the canonical totalized form-split graph.  On its
ordinary branch, polarization gives the \(\Geom\) occurrence and its
algebraic remainder is the oriented \(\Caus\) occurrence; either may be zero.
On its nonordinary branch the raw bulk current remains \(\Caus\), while the
failed form graph has support \(\{\Geom,\Caus\}\).  Thus no positivity,
closedness, sector, or gradient property is required for the split, and the
fixed enumeration makes the ancestry of each occurrence unique.
Quotient, lift, and boundary edges append their corresponding letters.
Their horizontal or base descendants retain the upstream head; their
vertical, interface, or defective current outputs have the head declared in
\textup{(E2)}--\textup{(E4)}.  A prescribed exterior or terminal current
begins with \(\Bdry\).

A coefficient or datum with no path through a current-producing edge is not
a summand of the weak evolution current.  It therefore contributes zero to
\(\frac{\dd}{\dd t}\Phi(u_t)\), even when it depends explicitly on space or
time.  When it enters a bulk or boundary current, it acquires that edge's
ancestry.  Exact storage contributes only through the endpoints of an
already sourced current identity and acquires that identity's ancestry.
\end{proof}

\begin{definition}[Free provenance algebra]
\label{def:free-continuum-provenance-algebra}
Let \(\mathfrak T(\mathfrak P)\) be the free totalized term algebra on the
primitive occurrences of
\cref{def:primitive-dynamical-role-partition}, with constructors
\[
 +,\ \circ,\ \operatorname{res},\ q,\ \ell,\ \operatorname{tr},\
 \operatorname{Graph},\ \operatorname{cl}_{\rm src}.
\]
They denote signed finite combination, composition, restriction, quotient,
lift/chart, trace/terminal readout, total-graph completion, and sourcewise
closure.  Its provenance semilattice is
\[
 \mathfrak P_5:=
 \bigl(2^{\Channels},\cup,\varnothing\bigr).
\]
The provenance morphism
\(\operatorname{Prov}:\mathfrak T(\mathfrak P)\to\mathfrak P_5\)
maps primitive roles to their letters, maps finite constructors to union,
appends \(\Abs,\Lift,\Bdry\) under \(q,\ell,\operatorname{tr}\), and
preserves support under restriction, totalization, and sourcewise closure.
Before \(\operatorname{Graph}\) is applied, typed elaboration inserts the
interface constructor of the least operation being totalized: for example a
weak-product state inserts \(\ell\), a failed quotient inserts \(q\), and a
trace failure inserts \(\operatorname{tr}\).  Thus Graph preserves the
already complete typed support; it cannot suppress a failure-interface
letter.
The empty value is reserved for current-free constants.  Exact cancellation
does not erase provenance; it places a nonempty sourced record in
\(\Null_\Sigma\).

Every atomic current term also has a \emph{head}
\(\operatorname{hd}(T)\in\Channels\), namely the output-current sort of the
typed term.  A native occurrence has the sort assigned by the primitive
role partition.  A composition has the sort of its outer current-producing
constructor and retains the provenance of every input.  A signed sum of
different current sorts is a total current rather than an atomic term and is
expanded into its typed summands.  Ordinary restriction graphs preserve the
head.  The horizontal quotient descendant and ordinary lift-base descendant
preserve it; quotient
intertwining/vertical/cokernel/graph coordinates have head \(\Abs\), while
lift interface/fiber/concentration/oscillation/scale/graph coordinates have
head \(\Lift\).  Every boundary-current readout, ordinary or defective, has
head \(\Bdry\); and a native form-split occurrence has head
\(\Geom,\Caus\), or \(\Bdry\) as assigned above.  A signed sum has no head until it is expanded into its
top-level atomic summands.  Always
\(\operatorname{hd}(T)\in\operatorname{Prov}(T)\).
\end{definition}

\begin{theorem}[Universal five-source provenance theorem]
\label{thm:universal-five-source-provenance}
The morphism \(\operatorname{Prov}\) is the unique provenance assignment
which agrees with the primitive dynamical-role partition and respects all
constructors of the closed continuum algebra.  More generally, let
\((S,\vee,0)\) be an idempotent join semilattice and let
\(\tau:\mathfrak T(\mathfrak P)\to S\) respect the same constructors, taking
the value \(s_c\) on a primitive role \(c\).  There is a unique join
homomorphism
\begin{equation}
 \label{eq:provenance-universal-factor}
 \overline\tau:\mathfrak P_5\to S,
 \qquad
 \overline\tau(A)=\bigvee_{c\in A}s_c,
 \qquad
 \tau=\overline\tau\circ\operatorname{Prov}.
\end{equation}
If a finite generated weak
current is written as its top-level signed sum
\(J=\sum_{i=1}^m\operatorname{ev}(T_i)\), then
\begin{equation}
 \label{eq:finite-provenance-factorization}
 q_\Sigma J
 =\sum_{c\in\Channels}q_\Sigma J^c(J)
 =\sum_{\varnothing\ne A\subseteq\Channels}q_\Sigma J_A(J),
 \qquad
 J^c(J):=\sum_{\operatorname{hd}(T_i)=c}\operatorname{ev}(T_i),
 \quad
 J_A(J):=\sum_{\operatorname{Prov}(T_i)=A}
              \operatorname{ev}(T_i).
\end{equation}
Every graph-complete or sourcewise limit retains the \emph{joint} vector of
head components and exact-support components together with the finite
factorization graph.  Thus no component is discarded before closure and no
sum of independently compactified extended values is asserted.  As a space
of provenance-tagged atomic occurrences,
\begin{equation}
 \label{eq:closed-five-source-span}
 \mathscr J_{\Sigma}^{\rm dyn}(E)
 =
 \bigsqcup_{\varnothing\ne A\subseteq\Channels}
 q_\Sigma\mathscr J_A^{\Phys,\Sat}(E) ,
\end{equation}
where \(\mathscr J_{\Sigma}^{\rm dyn}(E)\) is the target-active atomic
current-occurrence space generated by all public orbit records of physical
ceiling \(E\), and
\(\mathscr J_A^{\Phys,\Sat}(E)\) is the current evaluation of the saturated
record cell with exact support \(A\).  Analytically equal currents with
different histories remain distinct in the disjoint union.
\end{theorem}

\begin{proof}
Existence follows by recursion.  Primitive existence and disjointness are
\cref{lem:primitive-current-edge-exhaustion}.  Signed addition and
composition retain the union of the source leaves occurring in the
expression graph.  Quotient, lift, and trace constructors append their
letters.  Restriction changes no source leaf.  A total-graph value stores
the same input expression graph together with its first nonordinary output,
so it has the same ancestry, with the interface letter already appended.
Sourcewise closure is taken after the tag is fixed and therefore preserves
it.

For uniqueness, let \(P\) be another assignment with these properties.
It agrees with \(\operatorname{Prov}\) on primitive leaves.  Structural
induction on the term tree shows agreement under each finite constructor.
Total-graph and closure compatibility then gives agreement on completed
terms.  This is the universal property of the free totalized term algebra.
The same induction gives a unique head because every current-valued
constructor in the typed syntax has one output sort.  A graph defect changes
that sort only to the declared interface-failure sort.  This typing does not
alter the complete provenance support.

For the semilattice assertion, the values \(s_c\) determine a unique join
homomorphism on the free finite-subset semilattice by the displayed formula.
Constructor induction gives
\(\tau(T)=\bigvee_{c\in\operatorname{Prov}(T)}s_c\), which is precisely
\eqref{eq:provenance-universal-factor}.  Thus a sixth irreducible value would
require either a sixth primitive role or a source-creating constructor;
\cref{lem:primitive-current-edge-exhaustion} and the declared constructor
list exclude both.

To prove \eqref{eq:finite-provenance-factorization}, group the top-level
summands first by their heads and then by the nonempty support of their
complete term trees.  Each is a partition of the same finite summand list.
Evaluation is linear on weak pairings, so both grouped
identities hold before the target quotient and hence after applying
\(q_\Sigma\).  Terms with empty provenance are current-free constants and
vanish in the quotient; sourced exact cancellations retain their ancestry
and also vanish there.  Every bounded target readout is a structural-test
coordinate; it therefore commutes with graph-complete sourcewise limits.
Taking the closure of the joint finite factorization graph and retaining the
exact provenance tags proves
\eqref{eq:closed-five-source-span}.
\end{proof}

\begin{corollary}[Audit of derived and singularity-scale occurrences]
\label{cor:derived-source-audit}
Each commonly arising continuum occurrence is routed as follows:
\begin{longtable}{p{.31\linewidth}p{.58\linewidth}}
\toprule
Occurrence & Forced provenance route\\
\midrule
nonlinear product or coefficient degeneration &
ordinary output-current head with the union of its bulk-input ancestries; a
nonordinary weak-product state appends \(\Lift\), while a coefficient-domain,
chart, or interface failure receives the head and appended letter of its
least failed typed operation; every case retains the input union\\
commutator or polarized multilinear descendant &
the output head dictated by its signed weak-current role, with the union of
all input ancestries; a positive square generated by polarization is a
\(\Geom\)-headed production descendant, while a failed polarization graph is
\(\Lift\)-headed\\
reaction, nonvariational body term, or nonlocal bulk operator &
\(\Caus\) after the canonical split, with its ordinary positive symmetric
component in \(\Geom\)\\
time derivative, conservation law, or exact storage &
the shell derivative is replaced by the sourced right-hand-side current;
storage endpoints retain the ancestry of their carrier identity\\
initial datum, coefficient, parameter, or target-shell observable &
no source by itself; it acquires the ancestry of the current-producing edge
through which it affects shell progress\\
higher-order time equation &
\(\Lift\) for its first-order phase-space presentation, followed by the
\(\Geom/\Caus\) split of the resulting bulk evolution\\
memory or delay &
\(\Caus\) as a bulk relation, with \(\Lift\) appended when represented on a
history state\\
constraint multiplier, pressure, gauge, or symmetry reduction &
input ancestry together with \(\Abs\)\\
moving frame, modulation, localization state, or blow-up microscope &
input ancestry together with \(\Lift\)\\
weak defect, Reynolds stress, concentration, oscillation, or Young measure &
sourcewise limit of the nonlinear input ancestry, represented in the
\(\Lift\) normal form\\
stochastic drift, covariance, and sample path &
drift ancestry; \(\Geom\) for positive quadratic variation; \(\Lift\) for
the path augmentation\\
moving physical boundary or free interface &
\(\Lift\) for its chart/state motion and \(\Bdry\) for its oriented trace\\
shock, internal discontinuity, or topology-changing interface &
input ancestry together with \(\Bdry\) for its trace and \(\Lift\) when the
interface position is a state variable\\
nonuniqueness of weak evolution &
no new source; every branch retains the sources of its own weak current and
the total solution graph retains all branches\\
spatial infinity, exhaustion tail, killing, initial endpoint, or terminal
germ & input ancestry together with \(\Bdry\)\\
loss of closed range, trace, chain rule, covariance, compactness, or
continuation & first nonordinary total-graph value with the ancestry of the
operation whose graph failed\\
\bottomrule
\end{longtable}
The table is a consequence of the constructor theorem; it introduces no
additional routing convention.
\end{corollary}

\begin{proof}
Apply \cref{thm:universal-five-source-provenance} to the indicated expression
graph.  Positive quadratic variation is the \(\Geom\) branch of
\cref{prop:wiener-lift-source-routing}; every other row is one of the finite,
interface, total-graph, or closure constructors in the theorem.
\end{proof}

\begin{theorem}[No additional singularity source]
\label{thm:no-additional-continuum-singularity-source}
Let \(u\) be a public orbit in a finite physical stratum, and let \(W\) be
any atomic ordinary or graph-complete current occurrence in the canonical
top-level expansion of the prospective derivative of a continuation-shell
observable along \(u\).  Exactly one of
the following target-relative alternatives holds:
\begin{enumerate}[label=\textup{(N\arabic*)}]
\item \(q_\Sigma W=0\), so \(W\in\Null_\Sigma\) and the occurrence is
  residual for singular contact;
\item \(q_\Sigma W\ne0\), and the occurrence has a nonempty exact provenance
  support \(A\subseteq\Channels\) and a unique head
  \(c=\operatorname{hd}(W)\in A\), hence belongs to one of the \(31\)
  saturated source cells and exactly one of the five head normal forms.
\end{enumerate}
If \(u\) contacts the complete singular target, at least one occurrence of
type \textup{(N2)} appears on every normalized shell crossing.  Therefore
every singular contact uses target-aligned structure in
\(\Geom,\Caus,\Abs,\Lift,\Bdry\), possibly with composite support, and no
additional target-bearing source exists.
\end{theorem}

\begin{proof}
The prospective derivative is obtained by pairing the shell covector with
the public weak equation, its generated transformations, and their
sourcewise graph completions.  The totalized Stieltjes chain rule in
\cref{thm:defect-totalization-five-channel} exhausts its ordinary, jump,
Cantor, concentration, oscillation, and variation-boundary values.  By
\cref{thm:universal-five-source-provenance}, every occurrence in that
pairing belongs to a saturated cell with exact support
\(A\subseteq\Channels\), unless it is a current-free constant.  Constants
vanish under the orbit derivative.  For a nonempty source cell, either its target quotient
vanishes, giving \textup{(N1)}, or it does not, giving \textup{(N2)}.
These cases are disjoint and exhaustive.

If contact occurs, the normalized shell observable has unit endpoint
increment.  On a finite-value record, the affine shell graph
\eqref{eq:finite-affine-shell-graph} therefore has a positive summand.  On an
infinite-variation record, the closed owner relation
\eqref{eq:closed-shell-owner-floor} gives a positive atomic occurrence
directly.  Thus some occurrence is of type \textup{(N2)} without invoking an
extended signed sum.  Sourcewise completion covers weak defects and
singularity-scale limits, while total graphs cover every failed operation.
There are \(2^5-1=31\) possible nonempty supports.
Within a fixed head, the ordered refinements are exhaustive by
\cref{thm:continuum-geom-normal-form-exhaustion,%
thm:continuum-caus-normal-form-exhaustion,%
thm:continuum-abs-normal-form-exhaustion,%
thm:continuum-lift-normal-form-exhaustion,%
thm:continuum-bdry-normal-form-exhaustion}.
\end{proof}

\begin{remark}[Taxonomy versus exclusion]
\label{rem:source-exhaustion-not-regularity}
The source theorem classifies prospective singularity-forming structure; it
does not by itself bound that structure.  A contact-seeded formation-active
point may persist in any of the \(31\) cells and then gives a lower-bound
witness for the corresponding saturated source profile.  It constitutes a
singular outcome only if an independent public orbit realizes it.  Regularity
is obtained by the later target-local invariant and physical-price analysis
that upper-bounds every active source cell below the normalized contact floor.
The conclusion established here is that no target-bearing contribution lies
outside the five-source algebra.
\end{remark}

\begin{corollary}[Four-way audit of a proposed mechanism]
\label{cor:four-way-source-audit}
For a fixed public presentation, every proposed singularity mechanism is
resolved by the following exhaustive audit:
\begin{enumerate}[label=\textup{(A\arabic*)}]
\item if it occurs in the weak current graph, its constructor path gives one
  of the \(31\) source supports;
\item if it is a limit or failed operation on that graph, sourcewise
  completion or totalization gives the same inherited support;
\item if every target-shell pairing vanishes, it lies in \(\Null_\Sigma\);
\item if it has neither a public current path nor a total-graph limit of such
  paths, it is not an occurrence of the displayed evolution.
\end{enumerate}
Thus an analytically named object cannot form an additional source merely by
being singular, noncompact, nonlocal, distributional, or difficult to
estimate.
\end{corollary}

\begin{proof}
The first two cases are
\cref{thm:universal-five-source-provenance}.  The third is the definition of
\(\Null_\Sigma\).  The total solution graph and exhaustive approximation
tower contain precisely the ordinary public paths and their graph-complete
limits, so failure of both membership tests gives the fourth case.
\end{proof}

\begin{corollary}[No source created by a continuum microscope]
\label{cor:no-source-from-continuum-microscope}
Let two orbit descriptions be related by a generated quotient, chart,
localization, modulation, or compactification graph.  On its ordinary
covariance cell, their target-active currents have the same native heads and
ancestry, with only the declared \(\Abs,\Lift,\Bdry\) interface letters
appended.  On every nonordinary covariance cell, the first descent or graph
value has the corresponding interface head and retains the input ancestry.
Hence a blow-up microscope, moving frame, weak compactification, or alternate
gauge cannot create an additional singularity source.
\end{corollary}

\begin{proof}
Apply the head and provenance recursion of
\cref{def:free-continuum-provenance-algebra} to the generated transformation
graph.  Ordinary descent preserves weak current pairings and the target
quotient.  The totalized nonordinary branch is one of the quotient, lift, or
boundary coordinates specified in that recursion.
\end{proof}

\begin{theorem}[Operational native-source exhaustion]
\label{thm:operational-native-source-exhaustion}
For every operational weak-PDE signature, \(\mathsf C_\mathfrak P\) is
defined on every finite dynamically generated record and has image in the
five inherited source families.  Its evaluation commutes with the weak
equation: for every \(\varphi\in\mathcal D\),
\begin{equation}
 \label{eq:compiled-weak-current}
 \langle\dot u,\varphi\rangle
 =\sum_{c\in\Channels}\langle J^c(u),\varphi\rangle .
\end{equation}
The source word of a transformed record is the union of the native word and
the letters appended by the transformation.  There is no untyped native or
finite-constructor remainder.
\end{theorem}

\begin{proof}
The primitive paths of the displayed weak relation are exhausted by
\cref{lem:primitive-current-edge-exhaustion}.  In particular, every native
bulk term outside an ordinary positive symmetric form is routed through the
complementary \(\Caus\) branch, and every direct exterior occurrence is
\(\Bdry\).

On ordinary branch \textup{(W1)}, the displayed realization refines the
native relation into the mobility, directed-transport, and exterior terms.
Nonnegativity of \(\mathbb K\) identifies the first as \(\Geom\); the fixed
transport and exterior relation give \(\Caus\) and \(\Bdry\).
On ordinary branch \textup{(W2)}, polarization uniquely separates the symmetric and
antisymmetric forms.  The Beurling--Deny representation uniquely separates
the symmetric form into strongly local, jump, and killing terms; the first
two are \(\Geom\) and killing is \(\Bdry\).  The antisymmetric form is
\(\Caus\).

On the failure branch, retain the weak bulk current itself as a fixed
\(\Caus\) occurrence, so \eqref{eq:compiled-weak-current} still holds, and
record the failed attempted split with composite
\(\{\Geom,\Caus\}\)-ancestry.  The form row evaluates its target-active
quotient without asserting the missing realization.

The universal structural induction of
\cref{thm:universal-five-source-provenance} proves the remaining claim.  Each
constructor is covered by exactly one recursive clause in
\cref{def:continuum-source-compiler}; finite combination takes unions and
cannot erase an ancestor.  Applying the evaluation map at each induction
step gives \eqref{eq:compiled-weak-current}.  Weak closure is treated
sourcewise, so it changes the analytic representative but not the cell
label.
\end{proof}

\begin{corollary}[Failure is compiled structure]
\label{cor:failure-is-compiled-structure}
If an operation in \(\mathsf{Syn}(\mathfrak P)\) lacks a classical value on a
limit record, its totalized graph defect has the source support assigned to
the approximating graph, together with \(\Lift\) or \(\Bdry\) when the
failure occurs at a chart or terminal interface.  Consequently lack of
closed range, compactness, a classical trace, or a strong nonlinear limit
cannot create a sixth source.
\end{corollary}

\begin{proof}
Apply the compiler before graph closure, then use the sourcewise closure rule.
The graph defect is a limit of records in that tagged cell.  The chart and
terminal cases use the corresponding recursive clauses.
\end{proof}

\begin{theorem}[Sector-form source compiler]
\label{thm:sector-form-source-compiler}
Let \(X\) be locally compact and separable, let \(m\) be a Radon measure of
full support, and let \((\mathcal E^{\mathrm s},D(\mathcal E))\) be a
regular symmetric Dirichlet form on the real space \(L^2(X,m)\).  Write its
Beurling--Deny decomposition, with normalization absorbed into \(J\), as
\begin{equation}
 \label{eq:beurling-deny-source-split}
 \mathcal E^{\mathrm s}(u,v)
 =\mathcal E^{\mathrm c}(u,v)
 +\int_{X\times X\setminus\mathrm{diag}}
   (u(x)-u(y))(v(x)-v(y))\,J(\dd x,\dd y)
 +\int_Xuv\,\kappa(\dd x).
\end{equation}
Let \(\mathcal E^{\mathrm a}\) be antisymmetric on the same domain and
satisfy the sector bound
\begin{equation}
 \label{eq:sector-source-bound}
 |\mathcal E^{\mathrm a}(u,v)|
 \le C_{\mathrm{sec}}\,
   \mathcal E^{\mathrm s}_1(u,u)^{1/2}
   \mathcal E^{\mathrm s}_1(v,v)^{1/2}.
\end{equation}
Then the form
\(\mathcal E=\mathcal E^{\mathrm s}+\mathcal E^{\mathrm a}\) is closed and
sectorial on \(D(\mathcal E)\), and its source assignment is exhaustive:
\(\mathcal E^{\mathrm c}\) and the symmetric jump form are \(\Geom\),
\(\mathcal E^{\mathrm a}\) is the fixed \(\Caus\) form, and the killing
measure \(\kappa\) is \(\Bdry\).  The uniqueness of
\eqref{eq:beurling-deny-source-split} leaves no additional symmetric
Markovian primitive.  Quotients and state-space enlargements append
\(\Abs\) and \(\Lift\) as usual.
\end{theorem}

\begin{proof}
The Beurling--Deny theorem gives
\eqref{eq:beurling-deny-source-split} and uniqueness of its strongly local,
jump, and killing measures
\citep{FukushimaOshimaTakeda2011}.  Antisymmetry gives
\(\mathcal E^{\mathrm a}(u,u)=0\), while
\eqref{eq:sector-source-bound} makes it continuous in the
\(\mathcal E^{\mathrm s}_1\)-norm.  Hence \(\mathcal E\) has the same form
norm completeness as \(\mathcal E^{\mathrm s}\) and satisfies the sector
condition.  The stated routing is exactly the source assignment of
\cref{def:source-alphabet}.  A quasi-regular form has the same conclusion
after a declared quasi-homeomorphic regular realization.
\end{proof}

\begin{proposition}[Exact-drift symmetrization]
\label{prop:exact-drift-symmetrization}
Let \(M\) be a Riemannian domain, let
\(\Psi\in C^2(M)\).  On the locally-finite-weight branch let
\(\pi(\dd x)=e^{-\Psi(x)}\dd x\) is a locally finite measure.
For
\[
 L_\Psi g=\Delta g-\nabla\Psi\cdot\nabla g,
\]
totalize the integration-by-parts trace.  On its zero boundary-defect branch,
the resulting form core satisfies
\begin{equation}
 \label{eq:exact-drift-symmetric-form}
 -\langle f,L_\Psi g\rangle_{L^2(\pi)}
 =\int_M\nabla f\cdot\nabla g\,\dd\pi.
\end{equation}
Thus an exact drift is absorbed into a symmetric \(\Geom\) form by a weighted
\(\Abs\)-gauge.  If an exact/nonexact decomposition of a general drift is
proved on the declared domain, only its nonexact remainder retains independent
\(\Caus\) ancestry.  No global boundedness of \(e^{\pm\Psi}\) is inferred;
comparison with another physical measure is evaluated by its own totalized
physical-transfer row.
\end{proposition}

\begin{proof}
The weighted divergence identity
\[
 L_\Psi g=e^\Psi\operatorname{div}(e^{-\Psi}\nabla g)
\]
and the stated trace condition give
\eqref{eq:exact-drift-symmetric-form}.  Symmetry and nonnegativity of the
right side prove the source statement; changing the reference measure is the
declared gauge operation.
\end{proof}

\begin{proposition}[Wiener lift creates no additional source]
\label{prop:wiener-lift-source-routing}
Let \(H\) and \(U\) be separable Hilbert spaces and let a legal It\^o
evolution have drift \(b:H\to H\) and Hilbert--Schmidt noise coefficient
\(\sigma:H\to\mathcal L_2(U,H)\).  Totalize the generator relation on bounded
\(C^2\) cylinder functions.  Its ordinary branch is
\begin{equation}
 \label{eq:ito-generator-source-split}
 LF(u)=\langle DF(u),b(u)\rangle
 +\frac12\operatorname{Tr}
   \bigl[\sigma(u)\sigma(u)^*D^2F(u)\bigr].
\end{equation}
Then the carré du champ is
\begin{equation}
 \label{eq:ito-carre-du-champ}
 \Gamma(F)
 :=L(F^2)-2FLF
 =\|\sigma^*DF\|_U^2\ge0.
\end{equation}
Thus the diffusion covariance and predictable quadratic variation are
\(\Geom\); the drift retains its declared
\(\Geom/\Caus/\Bdry\) decomposition; and adjoining the Wiener path is
\(\Lift\).  Noise therefore creates no sixth continuum source.
\end{proposition}

\begin{proof}
Apply the second-order chain rule to \(F^2\) in
\eqref{eq:ito-generator-source-split}; the drift terms cancel and the
remaining trace is \(\|\sigma^*DF\|_U^2\).  It\^o's formula gives the same
integrand for the predictable quadratic variation of the martingale part.
The source labels follow from positivity of \(\Gamma\), the fixed drift
decomposition, and the state-space augmentation rule in
\cref{def:source-alphabet}.
\end{proof}

\section{Physical-budget-saturated sourcewise completion}
\label{sec:sourcewise-completion}

The closed continuum algebra contains every mathematical limit generated by
an ordinary finite-budget orbit.  Its topology is generated canonically by the
equation, target, physical balance, and constructor grammar.  Strong
convergence of a particular nonlinear term is subsequently read as a graph
invariant in this completion.

\begin{definition}[Structural test system]
\label{def:structural-test-system}
For each nonempty \(A\subseteq\Channels\) and derived physical ceiling \(\Budget\),
let \(\Records_A^0(\Budget)\) be the raw orbit records with source support
equal to \(A\).  The structural test system
\(\mathfrak F_A=(\phi_n^A)_{n\ge1}\) is generated recursively from the
countable coordinates \eqref{eq:canonical-evaluation-metric}.  At every finite
grammar depth it contains:
\begin{enumerate}[label=\textup{(T\arabic*)}]
\item the weak PDE pairings and source-current pairings;
\item the target shells and target fluxes;
\item every truncation of the physical-expenditure and entropy-production
coordinates generated by the fixed balance;
\item every chart, quotient, boundary, and tail coordinate generated by the
source grammar;
\item the total graphs of every finite algebraic operation generated by that
grammar, encoded by rational distance-to-graph truncations.
\end{enumerate}
An unbounded real-valued test is represented by its rational bounded
truncations.  Enumerating grammar depth and rational parameters gives one
countable family with values in \([-1,1]\).  No map or topology may be appended
after a row value is inspected.  Failure of an operation to extend as an
ordinary continuous map is recorded by the boundary of its total graph.
\end{definition}

\begin{definition}[Physical-budget-saturated closed orbit algebra]
\label{def:physical-saturated-algebra}
Let
\[
 \iota_A(r)=(\phi_n^A(r))_{n\ge1}
 \in K_A:=[-1,1]^{\mathbb N}.
\]
The space \(K_A\) is compact metrizable in the product metric.
The sourcewise saturated cell and the full physical orbit algebra are
\begin{equation}
 \label{eq:sourcewise-saturation}
 \Records_A^{\Phys,\Sat}(\Budget)
 :=\overline{\iota_A(\Records_A^0(\Budget))}^{K_A},
 \qquad
 \Records^{\Phys,\Sat}(\Budget)
 :=\bigsqcup_{\varnothing\ne A\subseteq\Channels}
 \Records_A^{\Phys,\Sat}(\Budget).
\end{equation}
A point of the closure is a mathematical limit record.  Its ancestry support
is the cell label \(A\).  A finer ancestry word may also be retained when it
is retained by the generated test system.  The closure contains all sequential
limits of ordinary orbit records seen by the generated tests.  Here \(\Src\) denotes
provenance support, not the smallest support of an analytically simplified
formula.
\end{definition}

\begin{remark}[Why the disjoint union is retained]
The same analytic denotation can arise through different source histories.
The closed algebra retains those histories as distinct records.  Forgetting
ancestry is an evaluation map from the disjoint union to the ambient analytic
space; it is not part of the source decomposition.
\end{remark}

\begin{figure}[tbp]
\centering
\resizebox{.96\linewidth}{!}{%
\begin{tikzpicture}[fw diagram,node distance=9mm and 12mm]
  \node[fw source,text width=3.0cm] (orbit)
    {ordinary orbit records\\derived ceiling \(\Budget\)};
  \node[fw provenance,text width=3.1cm,right=of orbit] (cells)
    {source cells\\\(\Records_A^0(\Budget)\)};
  \node[fw geom,text width=3.1cm,right=of cells] (tests)
    {structural test embedding\\\(\iota_A:\Records_A^0\to K_A\)};
  \node[fw use,text width=3.1cm,right=of tests] (closure)
    {cellwise mathematical closure\\\(\Records_A^{\Phys,\Sat}\)};
  \node[fw terminal,text width=3.2cm,right=of closure] (algebra)
    {closed orbit algebra\\\(\bigsqcup_A\Records_A^{\Phys,\Sat}\)};
  \draw[fw arrow] (orbit)--(cells);
  \draw[fw arrow] (cells)--(tests);
  \draw[fw arrow] (tests)--(closure);
  \draw[fw arrow] (closure)--(algebra);
\end{tikzpicture}%
}
\caption{Sourcewise physical saturation.  Closure is taken separately in
each ancestry cell before the cells are assembled, and every cell retains the
same physical ceiling \(\Budget\).}
\label{fig:sourcewise-completion}
\end{figure}

\begin{theorem}[Source-preserving continuum completion]
\label{thm:source-preserving-completion}
Every sequence \((r_i)\) of ordinary orbit records with \(\Src(r_i)=A\) that
converges in the structural test topology has its limit record in
\(\Records_A^{\Phys,\Sat}(\Budget)\).  Consequently
\begin{equation}
 \label{eq:source-preserving-limit}
 \Src\bigl(\lim_i r_i\bigr)
 =A
 =\bigcup_i\Src(r_i).
\end{equation}
If the sequence first passes through quotient or chart maps, \(\Abs\) or
\(\Lift\), respectively, is appended before taking the limit.  If the limit
is read as a boundary or tail trace, \(\Bdry\) is appended.  Thus defect
measures, oscillation variables, concentration profiles, closure directions,
blow-up germs, and tails create no source outside \(\Channels\).
\end{theorem}

\begin{proof}
By definition, \(\Records_A^{\Phys,\Sat}(\Budget)\) is the sequential closure
of the test image of the raw cell with provenance support \(A\); compact
metrizability makes this equal to its topological closure.  Hence every limit
of such a sequence belongs to that cell.  The algebraic rules of
\cref{def:source-alphabet} take unions of source supports and append the
labels of quotient, chart, and boundary operations before the closure is
formed.  The new coordinates of a limit record are limits of tests on the
same tagged sequence; they do not supply a new primitive letter.  This proves
\eqref{eq:source-preserving-limit} and each listed case.
\end{proof}

\begin{corollary}[Mixed sequences]
\label{cor:mixed-source-sequences}
Every sequence of raw records has source support in the finite set
\(2^{\Channels}\setminus\{\varnothing\}\).  It has a subsequence with constant
source support.  Each convergent such subsequence has its limit
therefore contained in an existing saturated cell.
\end{corollary}

\begin{proof}
There are only \(2^5-1\) nonempty source supports.  One support occurs
infinitely often.  Apply \cref{thm:source-preserving-completion} to that
subsequence.
\end{proof}

\begin{definition}[Source realization map]
\label{def:source-realization-map}
Let \(\mathcal P_A\) and \(\mathcal J_A\) be Hilbert spaces of potentials
and currents for a fixed source support \(A\subseteq\Channels\).  A
\emph{source realization map} is a densely defined closed operator
\[
 G_A:D(G_A)\subseteq\mathcal P_A\longrightarrow\mathcal J_A
\]
whose range consists of the currents carrying an actual potential witness in
that source cell.  Its saturated current space is
\(\overline{\Ran G_A}^{\mathcal J_A}\).  The closure has the same source
support \(A\), independently of whether every closure direction has a
potential witness.
\end{definition}

\begin{theorem}[Closed-range realization criterion]
\label{thm:closed-range-realization}
Let \(G:D(G)\subseteq\mathcal P\to\mathcal J\) be a source realization
map.  The following conditions are equivalent.
\begin{enumerate}[label=\textup{(R\arabic*)}]
\item \(\Ran G\) is closed in \(\mathcal J\).
\item There is \(\gamma>0\) such that
  \begin{equation}
   \label{eq:closed-range-poincare}
   \|G\phi\|_{\mathcal J}^{2}
   \ge\gamma\,\dist_{\mathcal P}(\phi,\ker G)^{2},
   \qquad \phi\in D(G).
  \end{equation}
\item The restriction of \(G^*G\) to \((\ker G)^\perp\) has spectrum in
  \([\gamma,\infty)\) for some \(\gamma>0\).
\item The Moore--Penrose inverse \(G^+\) is bounded on \(\Ran G\).
\end{enumerate}
When these conditions hold, every saturated current in
\(\overline{\Ran G}\) has a potential witness, unique modulo \(\ker G\).
When they fail, elements of \(\overline{\Ran G}\setminus\Ran G\) remain
limit records in the same source cell by
\cref{thm:source-preserving-completion}; failure of potential realization
does not define another source family.  In the reachability calculus this
operator theorem is activated only after restriction to the selected
target-local source cell; its complementary range value is then a polarity
successor of that cell.
\end{theorem}

\begin{proof}
The closed-range theorem for densely defined closed operators gives the
equivalence of \textup{(R1)}, \textup{(R2)}, and boundedness of \(G^+\).
On \((\ker G)^\perp\),
\(\|G\phi\|^2=\langle G^*G\phi,\phi\rangle\), so the spectral theorem makes
\textup{(R2)} equivalent to \textup{(R3)}.  Closed range gives
\(\overline{\Ran G}=\Ran G\), and the kernel is exactly the ambiguity in a
potential witness.  The final assertion is
\cref{thm:source-preserving-completion} applied to the source cell generated
by \(G\).
\end{proof}

\begin{theorem}[Target-relative gradient exhaustiveness]
\label{thm:target-relative-gradient-exhaustiveness}
Let \(G:D(G)\subseteq\mathcal P\to\mathcal J\) be a densely defined closed
realization of one inherited source cell, and let
\(P_{\overline G}\) be the orthogonal projection onto
\(\overline{\Ran G}\).  Then every \(j\in\mathcal J\) has the unique split
\begin{equation}
 \label{eq:closed-gradient-hodge-split}
 j=P_{\overline G}j+(I-P_{\overline G})j,
 \qquad
 P_{\overline G}j\in\overline{\Ran G},
 \quad (I-P_{\overline G})j\in\ker G^*.
\end{equation}
The first term belongs to the saturated cell generated by \(G\), including
when it lies in \(\overline{\Ran G}\setminus\Ran G\).  Relative to the
target family \(\Gamma_\Sigma\), the second term has the further canonical
split into its class in
\(\ker G^*/(\ker G^*\cap\Null_\Sigma)\) and its target-null part.  Hence
failure of closed range creates neither a new source nor a target-visible
residual: it creates a closure direction in the inherited cell, while every
orthogonal target-visible class remains a current of the native compiled
equation and retains its native ancestry.
\end{theorem}

\begin{proof}
The range closure is a closed Hilbert subspace, so the projection theorem
gives \eqref{eq:closed-gradient-hodge-split}; the identity
\((\Ran G)^\perp=\ker G^*\) identifies its complement.  The definition of
the sourcewise completion and
\cref{thm:source-preserving-completion} assign
\(P_{\overline G}j\) to the same tagged cell as the approximating gradients.
Quotienting the complementary space by its intersection with
\(\Null_\Sigma\) is exact.  Its kernel is target-null by definition, and a
nonzero quotient class has a nonzero target pairing.  Such a class cannot be
renamed residual; its ancestry is the ancestry of the native current from
which \(j\) was compiled by
\cref{thm:continuum-normal-form-compilation}.  These alternatives exhaust
\(\mathcal J\).
\end{proof}

\begin{corollary}[Regularized realization]
\label{cor:regularized-source-realization}
For \(j\in\mathcal J\) and \(\lambda>0\), the functional
\[
 \phi\longmapsto
 \|G\phi-j\|_{\mathcal J}^{2}+\lambda\|\phi\|_{\mathcal P}^{2}
\]
has a unique minimizer \(\phi_\lambda\in D(G)\).  The resulting current
\(G\phi_\lambda\) is a stable regularized realization, with the parameter
\(\lambda\) recording its bias.  Existence of this minimizer for every
\(\lambda>0\) does not imply any of the equivalent conditions in
\cref{thm:closed-range-realization}.
\end{corollary}

\begin{proof}
On \(D(G)\), the form
\(\langle G\phi,G\psi\rangle+\lambda\langle\phi,\psi\rangle\) is closed and
coercive in its form norm.  The functional
\(\psi\mapsto\langle j,G\psi\rangle\) is continuous in that norm, so the
Riesz representation theorem gives a unique minimizer.  The final statement
follows because adding \(\lambda I\) supplies coercivity even when the
unregularized range is not closed.
\end{proof}

\begin{figure}[tbp]
\centering
\resizebox{.97\linewidth}{!}{%
\begin{tikzpicture}[fw diagram,node distance=7mm and 9mm]
  \node[fw source,text width=2.6cm] (raw)
    {typed orbit quantity\\source support \(A\)};
  \node[fw provenance,text width=2.6cm,right=of raw] (normalize)
    {symmetry or blow-up chart\\append \(\Abs/\Lift\)};
  \node[fw geom,text width=2.4cm,above right=10mm and 10mm of normalize] (defect)
    {weak defect or\\oscillation measure};
  \node[fw use,text width=2.4cm,right=of normalize] (profile)
    {concentration or\\blow-up profile};
  \node[fw terminal,text width=2.4cm,below right=10mm and 10mm of normalize] (tail)
    {tail or boundary trace\\append \(\Bdry\)};
  \node[fw box,text width=3.0cm,right=15mm of profile] (same)
    {sourcewise closed cell\\no new primitive source};
  \draw[fw arrow] (raw)--(normalize);
  \draw[fw arrow] (normalize)--(defect);
  \draw[fw arrow] (normalize)--(profile);
  \draw[fw arrow] (normalize)--(tail);
  \draw[fw arrow] (defect)--(same);
  \draw[fw arrow] (profile)--(same);
  \draw[fw arrow] (tail)--(same);
\end{tikzpicture}%
}
\caption{Continuum limit phenomena inherit ancestry.  Normalization changes
the source word in the declared Abs/Lift manner; a terminal trace appends
Bdry.  Weakness or concentration changes the analytic realization, not the
source alphabet.}
\label{fig:limit-ancestry}
\end{figure}

\begin{remark}[Structural and evaluator topologies]
The graph-complete state space closes every typed operation together with its
defect coordinates.  A stronger PDE topology is represented by adjoining its
bounded seminorm and graph coordinates to the cylinder algebra.  Compactness
in that topology is then an ordered state-space assertion.  Failure of the
strong coordinate to be finite or closed is its typed graph value in the same
source cell.
\end{remark}

\section{Target quotient and the null residual}
\label{sec:target-null-residual}

\begin{definition}[Generated prospective target family]
\label{def:prospective-target-functionals}
Let \(\mathscr J^{\Phys,\Sat}(\Budget)\) be the current-occurrence part of
the closed orbit algebra.  The finite-current locus is replaced, when
necessary, by its Hausdorff locally convex envelope under the continuously
extended addition and scalar multiplication; the kernel of the envelope map
is retained as a tagged closure defect.  Records at an infinite
compactification endpoint are treated as boundary or tail records and are
tested by their bounded readouts, not inserted into this vector space.  If
\(\Sigma=\varnothing\), the target-shell algebra contains only constants and
\(\Gamma_\Sigma=\{0\}\).  Otherwise let
\(\rho_\Sigma(x)=d_{\rm ev}(x,\Sigma)\), where \(\Sigma\) is the complete
continuation boundary of \cref{def:singular-target}.  The \emph{target-shell
algebra} is the least countable algebra containing every rational ramp of
\(\rho_\Sigma\), its pullback by every ordinary generated covariance map, and
the rational cylindrical approximants supplied by the coordinates
\eqref{eq:canonical-evaluation-metric}.

For nonempty \(\Sigma\), the prospective target family \(\Gamma_\Sigma\) is
the closed normalized cone generated by every derived finite-cylinder shell current,
together with the killed-flux and predictable-compensator differentials when
their types occur in the fixed equation.  In particular it contains
\[
 J\longmapsto
 \langle D\widehat\phi_{r,m,I}(z),J\rangle
\]
on the ordinary branch and the corresponding Stieltjes, oscillation,
concentration, exterior, or variation-boundary pairing on every other branch.
Thus neither the probes nor their normalization are application data: they
are generated by the continuation target, approximation tower, and verified
covariance grammar.
\end{definition}

\begin{definition}[Target-null residual]
\label{def:target-null-residual}
The target-null space and target-active quotient are
\begin{equation}
 \label{eq:target-null-space}
 \Null_\Sigma
 :=\bigcap_{\Lambda\in\Gamma_\Sigma}\ker\Lambda,
 \qquad
 \mathscr J_\Sigma
 :=\mathscr J^{\Phys,\Sat}(\Budget)/\Null_\Sigma.
\end{equation}
The post-saturation residual is any closed current coordinate containing zero
and satisfying
\begin{equation}
 \label{eq:residual-is-target-null}
 \mathscr J_{\Res}^{\Phys,\Sat}\subseteq\Null_\Sigma.
\end{equation}
It is an audit coordinate for target-invisible transport, not a sixth source
family.  A current is target-visible precisely when its quotient class in
\(\mathscr J_\Sigma\) is nonzero.
\end{definition}

\begin{theorem}[Target quotient commutes with sourcewise completion]
\label{thm:target-quotient-completion}
For every source support \(A\subseteq\Channels\),
\[
 \overline{\Records_A+\Null_\Sigma}^{\,\rm gc}/\Null_\Sigma
 =
 \overline{q_\Sigma(\Records_A)}^{\,\mathscr J_\Sigma}.
\]
The kernel \(\Null_\Sigma\) is closed under sourcewise limits, signed gluing,
restriction, and every ordinary descendant map \(T\) whose covariance graph
proves \(T^*\Gamma_\Sigma\subseteq\Gamma_\Sigma\).  Consequently a
target-conditioned carrier, cohomology, moment, or pruning construction may
be performed on \(\mathscr J_\Sigma\) after ancestry has been fixed exactly on
that descent cell.  Every failed or noncovariant descent is retained instead
as the target-visible defect of \cref{thm:projection-first-saturation}.
\end{theorem}

\begin{proof}
Every element of \(\Gamma_\Sigma\) is a continuous cylinder-current
coordinate on the graph-complete space.  Its kernel is closed, and so is their
intersection.  The displayed identity is the standard quotient--closure
identity for a continuous linear quotient by a closed kernel.  Signed gluing
and restriction preserve every shell pairing by their defining graph
identities.  An ordinary descendant map preserves the common kernel when
\(T^*\Gamma_\Sigma\subseteq\Gamma_\Sigma\): for
\(J\in\Null_\Sigma\) and \(\Lambda\in\Gamma_\Sigma\),
\(\Lambda(TJ)=(T^*\Lambda)(J)=0\).  Without this covariance relation no such
claim is made; the descent defect is a separate target-active coordinate.
The displayed identity follows from the definition of the
quotient topology after saturating the source cell by the kernel.  Sourcewise
completion precedes this saturation, so no ancestry tag is lost.
\end{proof}

\begin{definition}[Compiled continuum current]
\label{def:closed-current-normal-form}
A continuum current is \emph{compiled} when it is generated from the native
terms of the fixed equation by the finite constructor rules
\textup{(S1)}--\textup{(S6)} of \cref{def:source-alphabet}.  Its evaluation
has the normal form
\begin{equation}
 \label{eq:raw-current-normal-form}
 J=J^{\Geom}+J^{\Caus}+J^{\Abs}+J^{\Lift}+J^{\Bdry}+J^{\Res},
 \qquad J^{\Res}\in\Null_\Sigma,
\end{equation}
as an equality of weak current pairings.  Empty components are allowed.
The five displayed components are the head partition
\(J^c(J)\) of \eqref{eq:finite-provenance-factorization}; within each
component every atomic occurrence retains its exact ancestry support
\(A\).  Thus the five additive current roles and the \(31\) provenance cells
are compatible classifications, not competing decompositions.
The residual is an optional regrouping in the common target kernel after the
five source components have been evaluated; the choice \(J^{\Res}=0\) is
always admissible.  It is not a constructor and carries no source letter.
\end{definition}

\begin{theorem}[Continuum normal-form compilation]
\label{thm:continuum-normal-form-compilation}
For every operational weak-PDE signature, totalize the native realization
relation and adjoin the quotient, chart, auxiliary-state, trace, forcing, and
boundary operations generated by \cref{def:source-alphabet}.  Every finite
orbit-current expression is compiled and has
\eqref{eq:raw-current-normal-form}.  A failed native realization additionally
generates its target-active form-defect class with
\(\{\Geom,\Caus\}\)-ancestry.  After sourcewise completion, every finite
target pairing has the same normal form.  Thus five-channel exhaustion does
not presume a gradient or sector realization.
\end{theorem}

\begin{proof}
On the fixed-Caus ordinary branch, positivity of \(\mathbb K\), the directed
transport, and the exterior term give \(\Geom\), \(\Caus\), and \(\Bdry\),
respectively.  On the form ordinary branch, the sector-form compiler makes the same
assignment to the symmetric mobility, antisymmetric directed, and killing
terms.  The universal induction in
\cref{thm:universal-five-source-provenance} proves the claim for every
orbit-current expression.  A quotient appends
\(\Abs\), a chart or auxiliary variable appends \(\Lift\), a trace appends
\(\Bdry\), and finite combination takes the union of the already assigned
supports.  These exhaust the transformed ordinary branches.

On the native failure branch, retain the raw fixed bulk current in \(\Caus\)
and its totalized attempted split in the composite Geom--Caus form-defect
cell, as in \cref{thm:operational-native-source-exhaustion}.  This additional
case remains inside the five-source alphabet.

For a sourcewise convergent net, each prospective target functional is a
coordinate of the structural test system.  Its value therefore passes to the
limit.  By \cref{thm:source-preserving-completion}, the limit remains in the
same tagged source cell.  The image in the quotient by \(\Null_\Sigma\) is
therefore generated by the five cell images.  Taking \(J^{\Res}=0\), or
regrouping any already identified target-null summand as \(J^{\Res}\), gives
\eqref{eq:raw-current-normal-form} without requiring a complemented kernel
and without adding a primitive source.
\end{proof}

\begin{theorem}[No post-saturation target-visible residual]
\label{thm:no-target-visible-residual}
Let \(\mathfrak P_{\Phys}\) be generated by the compiled continuum algebra
of \cref{thm:continuum-normal-form-compilation}.  Every target-visible current
record in \(\mathscr J^{\Phys,\Sat}(\Budget)\) has nonempty ancestry in
\(\Channels\), and its target-active quotient is generated by the five head
components, refined by the \(31\) possible nonempty exact-support cells of
\eqref{eq:closed-five-source-span}.  The residual source profile vanishes:
\begin{equation}
 \label{eq:zero-residual-source-profile}
 G_{\Budget,\Sigma}^{\Res,\mathrm{src},\Sat}
 :=\sup_{J\in\mathscr J_{\Res}^{\Phys,\Sat}}
   \sup_{\Lambda\in\Gamma_\Sigma}(\Lambda J)_+
 =0.
\end{equation}
If a provisional limit object pairs nontrivially with a target functional,
then \cref{thm:source-preserving-completion} places it in the existing source
cell inherited from its generating records.  Thus every dynamically generated
target-visible limit belongs to one of the five source families.
\end{theorem}

\begin{proof}
For raw occurrences, use
\eqref{eq:finite-provenance-factorization} and apply the quotient map
\(q_\Sigma:\mathscr J^{\Phys,\Sat}\to\mathscr J_\Sigma\) to
\eqref{eq:raw-current-normal-form}.  The residual term vanishes because it
lies in \(\Null_\Sigma\).  Continuity of every
\(\Lambda\in\Gamma_\Sigma\) passes the same identity to each sourcewise
limit.  The source-preserving theorem assigns every nonresidual limit term to
the closure of its existing exact-support cell.  Thus a nonzero quotient
class lies in a five-head component and one of its \(31\) provenance
refinements.  Finally, every functional in \(\Gamma_\Sigma\)
vanishes on \(\mathscr J_{\Res}^{\Phys,\Sat}\), so both nested suprema in
\eqref{eq:zero-residual-source-profile} are zero.
\end{proof}

\begin{corollary}[Five-source exhaustion and residual nullity]
\label{cor:five-source-exhaustion-residual-nullity}
Every target-visible atomic current occurrence generated by an ordinary
orbit, a sourcewise weak limit, a concentration or oscillation coordinate, a
terminal record, or a typed graph defect has a unique head
\[
 \operatorname{hd}(J)\in
 \{\Geom,\Caus,\Abs,\Lift,\Bdry\}
\]
and one nonempty exact provenance support
\[
 \operatorname{Prov}(J)=A\subseteq
 \{\Geom,\Caus,\Abs,\Lift,\Bdry\},
 \qquad \operatorname{hd}(J)\in A.
\]
A total current is the signed sum of its uniquely typed atomic occurrences;
it is not assigned an artificial common head.  Sourcewise completion retains
the joint head and exact-support vectors before the target quotient is taken.
After this saturation, every remaining residual coordinate belongs to
\(\Null_\Sigma\).  Hence the target-active quotient contains no sixth source
and no target-visible unowned remainder.
\end{corollary}

\begin{proof}
Primitive head uniqueness is
\cref{lem:primitive-current-edge-exhaustion}.  The unique extension of heads
and provenance through finite constructors, total graphs, and sourcewise
closure is \cref{thm:universal-five-source-provenance}.  Its finite
factorization graph is closed jointly, so neither head nor exact support is
chosen after passage to a limit.  Concentration, oscillation, terminal, and
graph-defect values are typed closure coordinates of their generating
records and therefore obey the same statement.  Finally,
\cref{thm:no-target-visible-residual} places the post-saturation residual in
the common kernel of all generated target functionals.
\end{proof}

\begin{theorem}[Projection-first saturation]
\label{thm:projection-first-saturation}
Let \(\mathsf G\) be the continuum constructor grammar generated by the
public equation, and let
\(q_\Sigma:\mathscr J^{\Phys,\Sat}\to\mathscr J_\Sigma\) be the quotient by
\(\Null_\Sigma\).  Totalize the descent graph of every constructor
\(T\in\mathsf G\).  On its ordinary descent cell set
\(T_\Sigma(q_\Sigma J):=q_\Sigma(TJ)\); on every other cell retain the first
descent defect
\begin{equation}
 \label{eq:target-projection-descent-defect}
 \mathcal D_{T,\Sigma}(J)
 :=q_\Sigma(TJ)-T_\Sigma(q_\Sigma J)
\end{equation}
with the ancestry of \(T\) and the target covector that detects it.  If
\(\mathsf G_\Sigma\) denotes the resulting projected grammar, then
\begin{equation}
 \label{eq:projection-first-saturation}
 q_\Sigma\!\left(\operatorname{Sat}_{\mathsf G}\mathscr J\right)
 =\operatorname{Sat}_{\mathsf G_\Sigma}
   \!\left(q_\Sigma\mathscr J\right).
\end{equation}
The identity holds sourcewise and with the physical carrier and same-origin
coordinates retained.  Thus the reachability calculus may quotient by
\(\Null_\Sigma\) before generating descendants.  A failed commutation is a
projected graph coordinate, so projection first loses no target-visible
current.
\end{theorem}

\begin{proof}
For a native current, \eqref{eq:projection-first-saturation} is the quotient
of \eqref{eq:raw-current-normal-form}.  Suppose it holds for the inputs of a
constructor.  On the ordinary descent graph, the definition of \(T_\Sigma\)
gives the identity for its output.  On the complementary graph, add
\(\mathcal D_{T,\Sigma}\); then
\(q_\Sigma TJ=T_\Sigma q_\Sigma J+\mathcal D_{T,\Sigma}(J)\) by
\eqref{eq:target-projection-descent-defect}.  Source-preserving totalization
retains the ancestry of the failed operation.  This proves both inclusions
for finite words by structural induction.

Every prospective target functional is a coordinate of the evaluation
compactification.  Hence \(q_\Sigma\) is continuous on each sourcewise total
graph, and the two finite-word inclusions pass to sourcewise completion.
Signed carrier gluing commutes with the quotient because it is performed
before Jordan decomposition.  Physical domination and same-origin relations
are retained coordinates rather than quotient relations.  This proves the
completed identity.
\end{proof}

\begin{corollary}[Source type and invariant value are independent]
\label{cor:source-type-bound-status}
Each saturated source cell has a value in every applicable fast-track
invariant: null or nonnull quotient, finite or infinite radius, positive or
zero price, zero or nonzero equality remainder, zero or positive capacity,
finite or infinite resistance, and divergent or finite cascade price.  These
values change neither the source word nor the target quotient.  Failure of a
compactness or resolvent invariant therefore returns the complementary
invariant branch; it never creates an untyped or residual source.
\end{corollary}

\begin{figure}[tbp]
\centering
\resizebox{.96\linewidth}{!}{%
\begin{tikzpicture}[fw diagram,node distance=9mm and 12mm]
  \node[fw source,text width=3.0cm] (limit)
    {orbit-generated limit\\with inherited ancestry};
  \node[fw use,text width=2.9cm,right=of limit] (test)
    {all prospective target\\pairings vanish?};
  \node[fw residual,text width=3.0cm,above right=10mm and 14mm of test] (null)
    {yes\\target-null \(\Res\)};
  \node[fw provenance,text width=3.2cm,below right=10mm and 14mm of test] (reentry)
    {no\\re-enter inherited source cell};
  \node[fw terminal,text width=3.1cm,right=16mm of reentry] (status)
    {exact complementary\\invariant branch};
  \draw[fw arrow] (limit)--(test);
  \draw[fw arrow] (test)--node[above left,font=\scriptsize]{yes}(null);
  \draw[fw arrow] (test)--node[below left,font=\scriptsize]{no}(reentry);
  \draw[fw arrow] (reentry)--(status);
\end{tikzpicture}%
}
\caption{Target-visible limits retain their inherited source cell.  The only
residual output is target-null.  A nonzero pairing fixes structural
ancestry before any attempt is made to estimate the corresponding source
profile.}
\label{fig:target-visible-reentry}
\end{figure}

\section{Prospective contact accountability}
\label{sec:prospective-accountability}

Target contact may produce a retrospective object, such as a singular
germ or a completed hitting distribution.  Source accountability is instead
read from the evolution strictly before contact.  The next three theorems give
deterministic, semigroup, and stochastic versions of the same factorization.

\begin{definition}[Canonical target shells and total shell current]
\label{def:target-shell-presentation}
If \(\Sigma=\varnothing\), set \(\tau_\Sigma=+\infty\),
\(\Gamma_\Sigma=\{0\}\), and generate no shell records.  Suppose henceforth
in this definition that \(\Sigma\ne\varnothing\).  In the compact metric
space \(X_{\rm ev}\), put
\(\rho_\Sigma(x)=d_{\rm ev}(x,\Sigma)\),
\(U_r(\Sigma)=\{\rho_\Sigma<r\}\), and, for rational \(r>0\),
\begin{equation}
 \label{eq:canonical-target-shell}
 \Phi_r(x)=
 \begin{cases}
  1,&\rho_\Sigma(x)\le r,\\
  2-\rho_\Sigma(x)/r,&r<\rho_\Sigma(x)<2r,\\
  0,&\rho_\Sigma(x)\ge2r.
 \end{cases}
\end{equation}
Then \(\Phi_r\) is continuous and \(r^{-1}\)-Lipschitz,
\(\bigcap_{r\in\mathbb Q_+}U_{2r}(\Sigma)=\Sigma\), and
\begin{equation}
 \label{eq:normalized-target-shell}
 \Phi_r=0\quad\text{on }X_{\rm ev}\setminus U_{2r}(\Sigma),
 \qquad
 \Phi_r=1\quad\text{on }U_r(\Sigma).
\end{equation}

Adjoin the bounded coordinate \(\Phi_r\) to the graph-complete algebra and
its defining distance-ramp graph to the ordered relations.  The rational
cylinder algebra separates the resulting compact state space.  By
Stone--Weierstrass, for every \(m\ge1\) there is a rational cylinder
polynomial \(\phi_{r,m}\) such that
\begin{equation}
 \label{eq:cylinder-shell-approximation}
 \|\phi_{r,m}-\Phi_r\|_\infty\le2^{-m}.
\end{equation}
Choose the polynomial first with error \(2^{-(m+2)}\).  The two strictly
positive cylinder polynomials
\(2^{-m}\pm(\phi_{r,m}-\Phi_r)\) then belong to the Archimedean quadratic
module by the Archimedean Positivstellensatz \citep{Schmudgen2017}.  The
canonical choice is the first polynomial in the fixed enumeration accompanied
by these two finite ordered identities.  For a crossing interval
\(I=[a,b]\) and \(m\) with
\(\delta_{r,m}(I):=\phi_{r,m}(u_b)-\phi_{r,m}(u_a)>0\), put
\[
 \widehat\phi_{r,m,I}
 :=\frac{\phi_{r,m}-\phi_{r,m}(u_a)}{\delta_{r,m}(I)}.
\]
The reciprocal is not inserted analytically: its total graph is generated by
the bounded coordinate \(h_{r,m,I}\) and the polynomial relations
\[
 \delta_{r,m}(I)h_{r,m,I}=1,
 \qquad
 \widehat\phi_{r,m,I}
 =h_{r,m,I}\bigl(\phi_{r,m}-\phi_{r,m}(u_a)\bigr).
\]
Since a crossing gives
\(\delta_{r,m}(I)\ge1-2^{1-m}>0\) for \(m\ge2\), this graph has one ordinary
value and a uniform bound.  The event-normalized cylinder shell therefore has
endpoint increment one on every crossing record.

On every finite approximation record, the weak evolution and the typed graph
chain rules derive a signed Stieltjes current
\begin{equation}
 \label{eq:derived-cylinder-shell-current}
 \dd\widehat\phi_{r,m,I}(u_t)
 =\sum_{c\in\Channels}\dd Q^c_{r,m,I}(t).
\end{equation}
Discrete time steps contribute their exact jump measures.  A failed product,
trace, inverse, chart, or covariance rule contributes the graph current of
its least failed typed operation to the corresponding inherited source cell.
There is no undifferentiated shell remainder.

The current coordinates are
\[
 Q^c_{r,m,I}(\psi)
 :=\int_I\psi\,\dd Q^c_{r,m,I},
 \qquad \psi\in C^1_{\mathbb Q}(I).
\]
Their sourcewise graph closure defines distributional shell currents on
\(X_{\rm gc}\).  On the bounded-variation branch they are finite signed Radon
measures with their Lebesgue, jump, and Cantor parts retained.  The provenance
disintegration additionally retains oscillation, concentration, and exterior
coordinates inherited from the approximating state records.  If their total
variation is unbounded, the normalized current and the compactified variation
coordinate
\[
 \left(\frac{Q^c}{1+\|Q^c\|_{\rm TV}},
       \frac{\|Q^c\|_{\rm TV}}{1+\|Q^c\|_{\rm TV}}\right)
\]
give the typed variation-boundary value.  The five totals are compactified
\emph{jointly}, rather than summed after separate compactification.  Fix the
enumeration of \(\Channels\), write
\[
 q=(q_c)_{c\in\Channels},\qquad q_c=Q^c_{r,m,I}(1),
 \qquad \iota(x)=\frac{x}{1+|x|},
\]
and, on every finite approximation, retain the affine shell graph
\begin{equation}
 \label{eq:finite-affine-shell-graph}
 \sum_{c\in\Channels}q_c=1,
 \qquad
 o(q):=\min\{c:q_c\ge1/5\}.
\end{equation}
The set defining \(o(q)\) is nonempty because the five signed coordinates
sum to one.  The \emph{total shell-current graph} is the closure of the joint
records
\begin{equation}
 \label{eq:total-shell-current-relation}
 \left((Q^c)_{c\in\Channels},(\iota(q_c))_{c\in\Channels},o(q)\right)
 \quad\text{subject to \eqref{eq:finite-affine-shell-graph}.}
\end{equation}
After passage to a subsequence the finite owner is constant.  Consequently
every value of the total graph has a source owner \(o\in\Channels\).  Write
\(\overline q_c:=\iota^{-1}(\lim\iota(q_c))\in[-\infty,+\infty]\), using the
continuous order-compactification inverse.  Then
\begin{equation}
 \label{eq:closed-shell-owner-floor}
 \overline q_o\in[1/5,+\infty],
 \qquad\text{equivalently}\qquad
 \iota(\overline q_o)\in[1/6,1].
\end{equation}
If all five totals are finite, the closed affine relation also gives
\(\sum_c\overline q_c=1\).  At an infinite-variation endpoint no sum of
extended signed quantities is formed; the retained owner relation
\eqref{eq:closed-shell-owner-floor} is the exact target-aligned information
which survives closure.  Thus cancellation cannot turn an
\(\infty-\infty\) expression into a fictitious residual current.

The contact time is
\(\tau_\Sigma(u)=\inf\{t:u(t)\in\Sigma\}\) for a state target and is
\(T_*(u)\) for the singular-event target when an ordinary same-origin maximal
path has no ordinary regular extension.  The latter event is physical even
when its terminal analytic graph coordinate is nonordinary.  An arbitrary
source-complete terminal record remains nonphysical until its preterminal path and
noncontinuation flag pass
\cref{thm:approximation-orbit-realization-criterion}.  A completed physical
contact record has
the following exhaustive shellwise form.  For each sufficiently small
\(r\), either the shell is entered before contact, in which case there are
strictly pre-contact times
\begin{equation}
 \label{eq:shell-crossing-times}
 0\le t_r^-<t_r^+<\tau_\Sigma(u),
 \qquad
 u(t_r^-)\notin U_{2r}(\Sigma),
 \qquad
 u(t_r^+)\in U_r(\Sigma).
\end{equation}
or it is not entered before contact, in which case the same relations hold
with \(t_r^+=\tau_\Sigma(u)\), and the endpoint increment is stored as the
\(\Bdry\)-headed jump current of the completed path.  For a source-complete terminal
record not attached to an ordinary same-origin preterminal path, the same
alternatives are statements only in \(\widehat\Sigma_{\rm sing}\); they
generate prospective currents but no physical-contact conclusion.  For a
realized singular event, the completed same-origin restrictions enter every
sufficiently small event shell before or at the terminal time.  The
metric shells separate every point outside the closed target, while contact
supplies the inner endpoint.  A target at infinity is already a boundary
point of \(X_{\rm ev}\); an additional ordinary normalization merely appends
its \(\Abs/\Lift\) ancestry.
\end{definition}

\begin{proposition}[Shell cofinality and realization fidelity]
\label{prop:shell-cofinality-realization-fidelity}
The canonical metric shells are cofinal for the closed source-complete target
\(\widehat\Sigma_{\rm sing}\).  Their restriction to a same-origin fiber of
an ordinary maximal preterminal path detects physical noncontinuation exactly
through its discrete continuation-status coordinate.  Hence shell crossing
in the source-complete graph space is sufficient for prospective source accounting,
whereas physical singular contact additionally requires either that realized
preterminal event or an ordinary realization of a state target.  No
graph-boundary current, unattached compactified germ, or shell limit supplies
such realization by itself.
\end{proposition}

\begin{proof}
Cofinality is the metric identity
\(\bigcap_{r\in\mathbb Q_+}U_{2r}(\widehat\Sigma_{\rm sing})
=\widehat\Sigma_{\rm sing}\).  On a singular-event fiber the preterminal
path and its origin are ordinary, and the discrete continuation-status
coordinate separates extendible from nonextendible maximal paths; no claim is
made that the terminal analytic graph coordinate is ordinary.  For a state
target, ordinary realization follows from
\cref{thm:approximation-orbit-realization-criterion}.  Outside these fibers the
same shells see only the source-preserving closure, so only prospective accountability
is sound.
\end{proof}

\begin{theorem}[Total normalized shell-crossing accountability]
\label{thm:deterministic-prospective-accountability}
Let \(u\) be an ordinary maximal orbit generated by the compiled continuum
algebra, with \(u(0)\notin\Sigma\), on the contact branch.  For every
sufficiently small \(r\), choose the completed crossing interval
\(I_r=[t_r^-,t_r^+]\), where \(t_r^+<\tau_\Sigma\) on a pre-contact entry and
\(t_r^+=\tau_\Sigma\) on a terminal entry, and then \(m\ge2\).
Equation \eqref{eq:cylinder-shell-approximation} gives
\(\delta_{r,m}(I_r)\ge1-2^{1-m}>0\).  Write the physical current as
\[
 J=\sum_{c\in\Channels}J^c+J^{\Res},
\]
where \(J^{\Res}\in\Null_\Sigma\).  Evaluate the derived cylinder-shell
current on \(I_r\).  Its ordinary signed source progress is
\begin{equation}
 \label{eq:shell-source-progress}
 Q_{c,r,m}(u)
 :=\int_{I_r}
 \langle D\widehat\phi_{r,m,I_r}(u_t),J_t^c\rangle\dd t.
\end{equation}
On every retained realization of this same-origin crossing branch, the shell
observable has unit endpoint increment.  If the five totals are finite, the
exact accountability identity is
\begin{equation}
 \label{eq:deterministic-prospective-factorization}
 \boxed{
 1
 =\widehat\phi_{r,m,I_r}(u(t_r^+))
  -\widehat\phi_{r,m,I_r}(u(t_r^-))
 =\sum_{c\in\Channels}Q^c_{r,m,I_r}(1).}
\end{equation}
On the ordinary branch,
\begin{equation}
 \label{eq:residual-shell-progress-zero}
 \int_{I_r}
 \langle D\widehat\phi_{r,m,I_r}(u_t),J_t^{\Res}\rangle\dd t=0.
\end{equation}
On every finite-value realization of this same-origin crossing branch,
\begin{equation}
 \label{eq:deterministic-source-extraction}
 \sum_{c\in\Channels}\bigl(Q^c_{r,m,I_r}(1)\bigr)_+\ge1,
 \qquad
 \max_{c\in\Channels}\bigl(Q^c_{r,m,I_r}(1)\bigr)_+
 \ge\frac15.
\end{equation}
On every realization, including an infinite-variation boundary value, the
joint shell graph supplies a canonical owner
\(o_{r,m,I_r}\in\Channels\) with
\begin{equation}
 \label{eq:deterministic-boundary-owner-extraction}
 \overline Q^{\,o_{r,m,I_r}}_{r,m,I_r}(1)\in[1/5,+\infty].
\end{equation}
The positive contribution may be ordinary, jump, Cantor, oscillation,
concentration, exterior, or variation-boundary current, but it always has the
ancestry of one of the five source cells.  Thus, without requiring a classical
chain rule,
\begin{equation}
 \label{eq:contact-implies-aligned-current}
 \boxed{\begin{gathered}
 u\text{ contacts }\Sigma\quad\Longrightarrow\\
 \forall r\text{ sufficiently small}\ \exists\text{ a five-channel source
 occurrence }R_r:\operatorname{Prog}_\Sigma(R_r)>0.
 \end{gathered}}
\end{equation}
Thus contact itself exposes a target-aligned occurrence in an existing source
cell at every sufficiently small shell.  It is strictly pre-contact when that
shell is entered before contact and is a \(\Bdry\)-headed endpoint occurrence
otherwise.
\end{theorem}

\begin{proof}
The cylinder approximation gives the stated positivity of
\(\delta_{r,m}(I_r)\), and event normalization gives the unit increment.
On every finite approximation,
\eqref{eq:finite-affine-shell-graph} gives the source sum.  On the ordinary
branch, insert the current normal form.  The generated shell differential
belongs to \(\Gamma_\Sigma\), so residual target-nullity proves
\eqref{eq:residual-shell-progress-zero}.  Taking positive parts of a
five-term signed sum equal to one proves
\eqref{eq:deterministic-source-extraction}.  Along any convergent net of
finite records, choose the canonical owner and then a subnet on which its
value is constant.  Closedness of \(\iota(q_o)\ge1/6\) proves
\eqref{eq:deterministic-boundary-owner-extraction} without forming an
extended signed sum.  The first-failure theorem and sourcewise closure assign
the owner its inherited nonempty support.  A terminal endpoint increment is
the completed path's boundary jump and therefore has \(\Bdry\) head.
\end{proof}

\begin{corollary}[No spontaneous deterministic contact]
\label{cor:prospective-source-before-contact}
Let \(u\) be an ordinary maximal orbit generated by the compiled continuum
algebra, with \(u(0)\notin\Sigma\).  Suppose that every ordinary and
graph-boundary derived shell current is target-null on every completed
crossing record, including a possible terminal boundary jump.  On the
ordinary approach branch this reads
\begin{equation}
 \label{eq:all-source-shell-null}
 \langle D\widehat\phi_{r,m,I}(u_t),J_t^c\rangle=0
 \quad\text{for a.e. }t<\tau_\Sigma(u),
 \quad c\in\Channels.
\end{equation}
Then \(u\) cannot contact \(\Sigma\).  Target avoidance also follows when every
derived source-current mass is nonpositive.  For a
singularity at infinity, these statements apply to the canonical boundary
shells in the generated evaluation compactification.
\end{corollary}

\begin{proof}
If contact occurred, \cref{thm:deterministic-prospective-accountability} would
produce a positive source-resolved derived current on every sufficiently small
shell.  This contradicts the stated target-nullity.  Under the stronger
ordinary zero-pairing condition, the same conclusion follows directly from
\eqref{eq:finite-affine-shell-graph}; on an infinite-variation endpoint it
follows from the retained owner floor
\eqref{eq:closed-shell-owner-floor}.
\end{proof}

\begin{corollary}[Normalized contact at infinity]
\label{cor:normalized-contact-at-infinity}
On the branch where singular contact is defined after generated normalizations
\(v_r=\mathcal S_r u\) into a state or germ compactification carrying
normalized shells \(\widetilde\Phi_r\).  If the normalized current
\(\widetilde J_r\) has the compiled normal form, including every chart derivative with
its \(\Abs/\Lift\) ancestry, then every normalized crossing interval satisfies
\[
 1
 =\sum_{c\in\Channels}
   \int_{I_r}
   \langle D\widetilde\Phi_r(v_r(t)),
           \widetilde J_{r,t}^{\,c}\rangle\dd t,
 \qquad
 \max_{c\in\Channels}
 \left(
   \int_{I_r}
   \langle D\widetilde\Phi_r(v_r(t)),
           \widetilde J_{r,t}^{\,c}\rangle\dd t
 \right)_+
 \ge\frac15.
\]
This holds when the five current totals are finite.  On an
infinite-variation value the
joint normalized shell graph instead retains an owner \(o\in\Channels\) with
compactified total \(\overline Q^{o}(1)\ge1/5\); no extended signed sum is
used.
Thus blow-up at infinity cannot occur without target-aligned structure in an
existing source cell.
\end{corollary}

\begin{proof}
Apply \cref{thm:deterministic-prospective-accountability} in the normalized
state space.  Source preservation under the legal chart is
\cref{thm:source-preserving-completion}.
\end{proof}

\begin{figure}[tbp]
\centering
\resizebox{.97\linewidth}{!}{%
\begin{tikzpicture}[fw diagram,node distance=8mm and 10mm]
  \node[fw source,text width=2.7cm] (orbit)
    {contact forces shell\\crossing \(I_r\)};
  \node[fw provenance,text width=2.7cm,right=of orbit] (shell)
    {normalized shell\\\(\Delta\Phi_r=1\)};
  \node[fw box,text width=3.0cm,right=of shell] (chain)
    {joint shell-current graph\\finite sum or closed owner};
  \node[fw residual,text width=2.5cm,above right=9mm and 10mm of chain] (residual)
    {residual pairing\\\(=0\)};
  \node[fw geom,text width=2.7cm,right=15mm of chain] (cells)
    {five typed source totals\\same-origin closure};
  \node[fw terminal,text width=3.1cm,right=of cells] (expose)
    {some \((Q_{c,r})_+\ge1/5\)};
  \draw[fw arrow] (orbit)--(shell);
  \draw[fw arrow] (shell)--(chain);
  \draw[fw arrow] (chain)--(residual);
  \draw[fw arrow] (chain)--(cells);
  \draw[fw arrow] (residual)--(cells);
  \draw[fw arrow] (cells)--(expose);
\end{tikzpicture}%
}
\caption{Normalized shell-crossing accountability.  Contact forces a unit
increase of the shell observable along the actual pre-contact current.  The
target-null residual contributes zero, so one of the five source integrals has
positive part at least \(1/5\).}
\label{fig:deterministic-prospective}
\end{figure}

\begin{theorem}[Resolvent source-response comparison]
\label{thm:resolvent-prospective-accountability}
Let \(L\) and \(L_0\) be closed killed generators on the same Banach space,
and let \(\lambda\in\rho(L)\cap\rho(L_0)\).  Write
\[
 R_L(\lambda)=(\lambda-L)^{-1},
 \qquad R_0(\lambda)=(\lambda-L_0)^{-1}.
\]
Totalize domain membership of \(R_0(\lambda)k_\Sigma\) and the generator
difference.  On their ordinary branch,
\begin{equation}
 \label{eq:generator-difference-sources}
 L-L_0=\sum_{c\in\Channels}K^c+K^{\Res}
\end{equation}
on that vector.  Then, for every continuous initial functional \(\ell_0\),
\begin{equation}
 \label{eq:resolvent-prospective-factorization}
 \begin{aligned}
 \langle\ell_0,(R_L(\lambda)-R_0(\lambda))k_\Sigma\rangle
 ={}&\sum_{c\in\Channels}
 \langle R_L(\lambda)^*\ell_0,
 K^cR_0(\lambda)k_\Sigma\rangle\\
 &+\langle R_L(\lambda)^*\ell_0,
 K^{\Res}R_0(\lambda)k_\Sigma\rangle.
 \end{aligned}
\end{equation}
If the last pairing is target-null, it vanishes.  If the real part \(p\) of
the left side is positive, some existing source cell has real pairing at
least \(p/5\).
Any domain or generator-split failure is returned by the form/realization row
and appears only as the ordinary value of that specialized row.  The
factorization is used for reachability only on the retained shell/source cell;
the five displayed summands inherit the source words of
\eqref{eq:generator-difference-sources}.
\end{theorem}

\begin{proof}
The second resolvent identity gives
\[
 R_L(\lambda)-R_0(\lambda)
 =R_L(\lambda)(L-L_0)R_0(\lambda).
\]
Pair with \(\ell_0\) and \(k_\Sigma\), then insert
\eqref{eq:generator-difference-sources}.  Target-nullity removes the residual
term.  A positive sum of five real numbers has a summand at least one fifth
of that sum.
\end{proof}

\begin{remark}[Arrival semantics]
Equation \eqref{eq:resolvent-prospective-factorization} is an operator
identity.  It becomes a difference of discounted contact responses when the
killed semigroups admit the target-flux realization of
\cref{thm:killed-resolvent-stopwatch}.  No probabilistic meaning is inferred
from the algebraic resolvent alone.
\end{remark}

\begin{figure}[tbp]
\centering
\resizebox{.98\linewidth}{!}{%
\begin{tikzpicture}[fw diagram,node distance=8mm and 10mm]
  \node[fw source,text width=2.6cm] (reference)
    {reference killed generator\\\(L_0\)};
  \node[fw provenance,text width=2.8cm,right=of reference] (difference)
    {pre-hit difference\\\(L-L_0\)};
  \node[fw geom,text width=2.8cm,right=of difference] (identity)
    {second resolvent identity\\\(R_L(L-L_0)R_0\)};
  \node[fw use,text width=2.8cm,right=of identity] (pair)
    {initial/target-flux\\cross pairing};
  \node[fw terminal,text width=3.1cm,right=of pair] (sources)
    {response difference factors\\through five source cells};
  \node[fw residual,text width=2.7cm,below=13mm of identity] (res)
    {target-null residual\\pairing \(=0\)};
  \draw[fw arrow] (reference)--(difference);
  \draw[fw arrow] (difference)--(identity);
  \draw[fw arrow] (identity)--(pair);
  \draw[fw arrow] (pair)--(sources);
  \draw[fw dashed arrow] (res)--(pair);
\end{tikzpicture}%
}
\caption{Resolvent source-response comparison.  A discounted target-response
difference is sandwiched around the pre-hit generator difference and therefore
inherits its five-source decomposition.  The no-spontaneous-contact theorem
itself comes from normalized shell crossing.}
\label{fig:resolvent-prospective}
\end{figure}

\begin{theorem}[Stochastic no-spontaneous-contact accountability]
\label{thm:stochastic-prospective-accountability}
Let \(\tau_\Sigma\) be a stopping time and
\(N_t=\mathbf 1_{\{\tau_\Sigma\le t\}}\).  Totalize the Doob--Meyer and
source-split relations.  On the branch where \(N-A\) is a martingale, the
predictable compensator on \([0,T]\) has the finite-variation
decomposition
\begin{equation}
 \label{eq:compensator-source-decomposition}
 A=\sum_{c\in\Channels}A^c+A^{\Res},
\end{equation}
where all measures are supported on \(\{t<\tau_\Sigma\}\), the signed source
terms are integrable, every term starts from zero,
\(\mathbb P\{\tau_\Sigma=0\}=0\), and
\(\mathbb E A_T^{\Res}=0\).  Then
\begin{equation}
 \label{eq:stochastic-prospective-factorization}
 \mathbb P\{\tau_\Sigma\le T\}
 =\sum_{c\in\Channels}\mathbb E A_T^c.
\end{equation}
If \(p_T:=\mathbb P\{\tau_\Sigma\le T\}>0\), one pre-contact source
compensator satisfies \(\mathbb E A_T^c\ge p_T/5\).  In particular, all five
source expectations cannot vanish when contact has positive probability.
Failure of the martingale, finite-variation, support, or integrability
relation is its stochastic total-graph defect and remains in the fluctuation
row.
\end{theorem}

\begin{proof}
Taking expectations of the integrable martingale \(N-A\), using the stated
zero initial values, gives
\(\mathbb E N_T=\mathbb E A_T\).  Insert
\eqref{eq:compensator-source-decomposition} and use the zero expectation of
the residual term.  The final assertion follows from the five-term
pigeonhole argument.
\end{proof}

\begin{figure}[tbp]
\centering
\resizebox{.95\linewidth}{!}{%
\begin{tikzpicture}[fw diagram,node distance=9mm and 13mm]
  \node[fw source,text width=2.8cm] (stop)
    {first-contact process\\\(N_t=\mathbf1_{\{\tau_\Sigma\le t\}}\)};
  \node[fw provenance,text width=2.9cm,right=of stop] (doob)
    {Doob--Meyer split\\\(N-A\) martingale};
  \node[fw geom,text width=3.0cm,right=of doob] (comp)
    {predictable pre-contact\\compensator \(A\)};
  \node[fw use,text width=3.0cm,right=of comp] (decomp)
    {five source measures\\plus target-null residual};
  \node[fw terminal,text width=3.0cm,right=of decomp] (contact)
    {positive contact exposes\\one predictable source};
  \draw[fw arrow] (stop)--(doob);
  \draw[fw arrow] (doob)--(comp);
  \draw[fw arrow] (comp)--(decomp);
  \draw[fw arrow] (decomp)--(contact);
\end{tikzpicture}%
}
\caption{Stochastic no-spontaneous-contact accountability.  First contact is
charged by its predictable pre-contact compensator.  A positive contact
probability forces a positive contribution from one of the five source
measures before the stopping time.}
\label{fig:stochastic-prospective}
\end{figure}

\begin{corollary}[Continuum prospective source exhaustion]
\label{cor:continuum-prospective-exhaustion}
Normalized deterministic contact forces a positive completed-crossing
occurrence with ancestry in \(\Channels\); it is pre-contact on the approach
branch and \(\Bdry\)-headed on a terminal jump.  Positive-probability
stochastic contact forces a predictable pre-contact occurrence.  A positive
killed-resolvent response difference has the same source-extraction property.
The post-saturation residual contribution is zero in all three interfaces.
\end{corollary}

\begin{definition}[Saturated prospective source profile]
\label{def:saturated-prospective-profile}
Fix one of the three prospective interfaces, restricted to records satisfying
its stated target-null residual condition, and fix its normalization.  Write
\(\operatorname{Prog}_\Sigma^c(R)\) for the signed \(c\)-source contribution:
\(Q_{c,r}\) for a deterministic shell, \(\mathbb E A_T^c\) for a stochastic
contact record, and the real \(c\)-pairing in
\eqref{eq:resolvent-prospective-factorization} for a resolvent comparison.
For a deterministic infinite-variation boundary record, use its joint shell
owner \(o\) and set
\begin{equation}
 \label{eq:boundary-owner-progress-convention}
 \operatorname{Prog}_\Sigma^c(R)
 :=\begin{cases}
     \overline Q^{o}(1),&c=o,\\
     0,&c\ne o.
   \end{cases}
\end{equation}
This is a target-aligned readout of the closed affine graph, not a proposed
sum of the underlying extended signed currents.  For finite additive records
set
\begin{equation}
 \label{eq:total-source-progress}
 \operatorname{Prog}_\Sigma(R)
 :=\sum_{c\in\Channels}\operatorname{Prog}_\Sigma^c(R).
\end{equation}
For a boundary-owner record, use the same notation for the sum in
\eqref{eq:boundary-owner-progress-convention}, namely
\(\operatorname{Prog}_\Sigma(R)=\overline Q^{o}(1)\ge1/5\).
Define
\begin{align}
 \mathsf T_{\Sigma,c}^{\Phys,\Sat}(\Budget)
 &:=\sup_{R\in\Records^{\Phys,\Sat}(\Budget)}
       \bigl(\operatorname{Prog}_\Sigma^c(R)\bigr)_+,
 \label{eq:channel-prospective-profile}\\
 \mathsf T_\Sigma^{\Phys,\Sat}(\Budget)
 &:=\sum_{c\in\Channels}
       \mathsf T_{\Sigma,c}^{\Phys,\Sat}(\Budget).
 \label{eq:total-prospective-profile}
\end{align}
The supremum is taken over the ordinary graph values and the owner-lifted
boundary values of the selected interface and may equal \(+\infty\).  It is
also saturated over every universal expanded composite record of
\cref{def:cumulative-joint-continuum-record}: the profile sees the joint value
before any postprocessing, normalization, averaging, or residual quotient.
Inputs outside an ordinary graph enter the corresponding totalized defect
cell, so every input and every target-visible cross term has a structural
value in an inherited source cell.  This is the direct continuum analogue of
Part~I's saturated profile over complete represented derivations.
\end{definition}

\begin{definition}[Source-support and realization-mode profiles]
\label{def:source-support-realization-profiles}
For each nonempty \(A\subseteq\Channels\), set
\[
 \operatorname{Prog}_\Sigma^A(R)
 :=\sum_{c\in A}\operatorname{Prog}_\Sigma^c(R)
   \mathbf 1_{\{\Src(R)=A\}}.
\]
Every compiled record also carries the subset of realization-mode flags
\[
 \Modes=\{\mathsf{exact},\mathsf{closure},\mathsf{harmonic},
          \mathsf{concentration},\mathsf{boundary}\}
\]
determined by its construction:
\begin{enumerate}[label=\textup{(M\arabic*)}]
\item \(\mathsf{exact}\): it has a witness in the range of its native
realization map;
\item \(\mathsf{closure}\): it lies in the sourcewise range closure, possibly
outside the range;
\item \(\mathsf{harmonic}\): its class in
\(\ker G_A^*/(\ker G_A^*\cap\Null_\Sigma)\) is nonzero;
\item \(\mathsf{concentration}\): it is read through a Young measure, defect
measure, profile, or blow-up-state lift;
\item \(\mathsf{boundary}\): it is a trace, killing, tail, terminal germ, or
graph-boundary record.
\end{enumerate}
Flags may coexist; they refine ancestry and do not form new sources.  For
bookkeeping fix the order
\[
 \mathsf{exact}<\mathsf{closure}<\mathsf{harmonic}
 <\mathsf{concentration}<\mathsf{boundary}
\]
and let \(m_*(R)\) be the maximal flag carried by \(R\).  Thus the complete
flag set is retained, while every record belongs to exactly one
\emph{primary-mode} cell.  Generated operations only append flags, so
\(m_*\) is monotone along a successor path and hence eventually constant.
For \(m\in\Modes\), define
\begin{equation}
 \label{eq:source-mode-profile}
 \mathsf T_{\Sigma,A,m}^{\Phys,\Sat}(\Budget)
 :=\sup_{\substack{R\in\Records_A^{\Phys,\Sat}(\Budget)\\m_*(R)=m}}
       \bigl(\operatorname{Prog}_\Sigma^A(R)\bigr)_+,
\end{equation}
with supremum \(0\) for an empty class.
\end{definition}

\begin{theorem}[Realization-mode exhaustion]
\label{thm:realization-mode-exhaustion}
Every target-active continuum record belongs to a nonempty source cell
\(A\subseteq\Channels\), has at least one flag in \(\Modes\), and has exactly
one primary mode.  Moreover,
\begin{equation}
 \label{eq:mode-profile-exhaustion}
 \mathsf T_\Sigma^{\Phys,\Sat}(\Budget)
 \le
 \sum_{\varnothing\ne A\subseteq\Channels}
 \sum_{m\in\Modes}
 \mathsf T_{\Sigma,A,m}^{\Phys,\Sat}(\Budget).
\end{equation}
Thus exact currents, nonclosed-range limits, harmonic classes,
concentrations, and boundary defects exhaust the analytic realization modes
of target progress without enlarging the five-source alphabet.
\end{theorem}

\begin{proof}
Native records are exact.  A sourcewise limit is in the closure mode.  The
orthogonal decomposition
\(\mathcal J_A=\overline{\Ran G_A}\oplus\ker G_A^*\), followed by the target
quotient, assigns every target-active orthogonal record to the harmonic mode.
The compiler marks every profile lift and every boundary operation with the
last two flags.  These cases cover the finite syntax and its sourcewise
completion.  Every positive source contribution is therefore bounded by at
least one term on the right of \eqref{eq:mode-profile-exhaustion}; summing
gives the result.  Multiple flags only enlarge the right side.
\end{proof}

\begin{theorem}[Source-only prospective envelope]
\label{thm:source-only-envelope}
Every record in the domain of a prospective interface satisfies
\begin{equation}
 \label{eq:source-only-envelope}
 \bigl(\operatorname{Prog}_\Sigma(R)\bigr)_+
 \le \mathsf T_\Sigma^{\Phys,\Sat}(\Budget).
\end{equation}
The envelope consists solely of the five source profiles.  Any finite upper
bounds on those profiles give a complete prospective progress bound.  A
finite-value normalized deterministic contact record satisfies
\(\operatorname{Prog}_\Sigma(R)=1\); an infinite-variation contact record has
owner progress at least \(1/5\).
\end{theorem}

\begin{proof}
On a finite additive record, each of
\cref{thm:deterministic-prospective-accountability,thm:resolvent-prospective-accountability,thm:stochastic-prospective-accountability}
has, after removal of its target-null residual pairing, the form
\[
 \operatorname{Prog}_\Sigma(R)
 =\sum_{c\in\Channels}\operatorname{Prog}_\Sigma^c(R).
\]
Take the positive part, bound every signed summand by its positive part, and
then apply
\eqref{eq:channel-prospective-profile}.
For a deterministic variation-boundary record,
\eqref{eq:boundary-owner-progress-convention} has a single nonzero source
coordinate, so the same bound follows without summing extended signed
currents.  Every other total-graph defect is evaluated in its inherited
source cell by its bounded target readout, to which the same one-coordinate
argument applies.
\end{proof}

\begin{definition}[Singularity-aligned structural support]
\label{def:singularity-aligned-support}
For a pre-contact record \(R\) in any prospective interface, define
\begin{equation}
 \label{eq:singularity-aligned-support}
 \Align_\Sigma(R)
 :=\left\{c\in\Channels:
       \operatorname{Prog}_\Sigma^c(R)>0\right\}
 \cup
 \begin{cases}
  \Src(R),&R\text{ is a total-graph defect with }
             \operatorname{Prog}_\Sigma(R)>0,\\
  \varnothing,&\text{otherwise}.
 \end{cases}
\end{equation}
Equivalently, in the deterministic interface an occurrence is aligned when
the singularity-descending covector \(D\Phi_r(u_t)\) has positive integrated
pairing with its source current on the crossing interval.
The word ``aligned'' is target-relative: it does not rename a source as
\(\Caus\), and it is evaluated only along the fixed equation.
\end{definition}

\begin{theorem}[No-spontaneous-progress principle]
\label{thm:target-alignment-necessity}
For every prospective record, including total-graph defect records,
\begin{equation}
 \label{eq:positive-progress-needs-alignment}
 \operatorname{Prog}_\Sigma(R)>0
 \quad\Longrightarrow\quad
 \Align_\Sigma(R)\ne\varnothing.
\end{equation}
For an ordinary five-term additive record one also has
\begin{equation}
 \label{eq:ordinary-five-term-margin}
 \qquad
 \max_{c\in\Channels}
   \bigl(\operatorname{Prog}_\Sigma^c(R)\bigr)_+
 \ge \frac{\operatorname{Prog}_\Sigma(R)}5.
\end{equation}
If every ordinary source occurrence and every total-graph defect represented by \(R\) lies in
\(\Null_\Sigma\), then
\begin{equation}
 \label{eq:null-structure-zero-progress}
 \operatorname{Prog}_\Sigma(R)=0,
\end{equation}
regardless of the size, oscillation, or compactness of those target-null
currents.  More generally, if every source contribution is nonpositive, then
positive target progress is impossible.  For a normalized deterministic shell
crossing, the derived Stieltjes-current identity yields a positive
contribution of at least \(1/5\) in one inherited source cell on every
ordinary or graph-boundary branch.
\end{theorem}

\begin{proof}
Equation \eqref{eq:total-source-progress} is the compiled five-term identity
on finite additive records.  A positive finite sum has a positive summand,
and its largest positive part is at least one fifth of the sum.  On a
variation-boundary record, the owner convention
\eqref{eq:boundary-owner-progress-convention} already supplies a positive
summand with inherited nonempty support; no extended sum is asserted.
Source-preserving graph closure assigns every other positive defect readout
its inherited support.  If every occurrence lies in the common target
kernel, each ordinary or bounded defect readout vanishes.  The shell
assertion is \cref{thm:deterministic-prospective-accountability}.
\end{proof}

\begin{corollary}[No-spontaneous-contact regularity]
\label{cor:target-avoidance-regularity}
Let the target have a normalized shell presentation, and let every legal
orbit be generated by the compiled algebra on its regular pre-contact interval.  If
\(\Align_\Sigma(R)=\varnothing\) for every legal shell record, then the target
is unreachable.  In particular, if all five source currents have zero class
in \(\mathscr J_\Sigma\) along every completed crossing record, their ancestry is
retained while their target contribution lies in \(\Null_\Sigma\), and contact is
impossible:
\begin{equation}
 \label{eq:null-current-direct-regularity}
 \boxed{
 q_\Sigma(J^c)=0
 \ \text{on every completed target-shell crossing for all }c\in\Channels
 \quad\Longrightarrow\quad
 \tau_\Sigma=\infty.}
\end{equation}
This conclusion requires neither compactness of the full orbit nor a positive
physical-price estimate.
\end{corollary}

\begin{proof}
If an orbit contacted \(\Sigma\),
\cref{thm:deterministic-prospective-accountability} would provide a normalized
shell record with unit progress.  By
\cref{thm:target-alignment-necessity}, its aligned support would be nonempty,
contrary to the displayed null-source condition.  Under that quotient value, every
shell source integral vanishes, so the same contradiction follows directly
from \eqref{eq:deterministic-prospective-factorization}.
\end{proof}

\begin{remark}[Why the target quotient is stronger than a full norm bound]
A classical energy estimate commonly controls every component of a current,
including directions irrelevant to the singular target.  The quotient
\(\mathscr J_\Sigma\) discards a current only after all declared prospective
target pairings vanish.  The discarded current may be large or noncompact;
no estimate of it is needed for \eqref{eq:null-structure-zero-progress}.
Only the target-aligned quotient directions proceed to resistance,
compactness, rigidity, or physical-price invariants.  The quotient strictly
reduces the regularity problem because target-null directions are absent from
all downstream estimates.
\end{remark}

\part{Physical Structural-Invariant Calculus}
\label{part:invariant-calculus}

\section{Structural invariant reduction of analytic arguments}
\label{sec:structural-invariant-reduction}

\begin{remark}[Ordinary-branch convention for analytic modules]
Every conditional identity in this part describes the ordinary value of a
total graph generated from the public signature.  Its operator, form, chart,
or multiplier must already occur in the typed syntax or in the fixed
target-independent textbook library.  The condition is never appended as
specialization data: when it does not hold, its complementary graph value is
sent to the activity lift and successor relation.  Hamiltonian, commutator,
compactness, tail, rigidity, and entropy modules can shorten a pruning
transcript, but the core contact-to-successor theorem does not require their
ordinary values.
\end{remark}

The closed source algebra separates the origin of a continuum quantity from
the estimate used to control it.  Once a source family has been placed in a
closed target-null quotient or removed from the admissible orbit family by a
proved physical exclusion, later invariants are evaluated on the remaining
reduced space.  The analytic calculation then has only three possible jobs:
close a remaining class, compare two nonnegative forms, or resolve the
equality set of a sharp comparison.

\begin{definition}[Quotient-compatible structural reduction]
\label{def:invariant-compatible-reduction}
Let \(Y\) be a Hausdorff locally convex space of currents, profiles, or
limit records, and let \(K\subseteq Y\) be a closed linear subspace already
annihilated by the downstream target and invariant maps.  Write
\[
 \pi_K:Y\longrightarrow Y/K.
\]
A downstream map, form, evaluation closure, or source transport is
\emph{compatible with the reduction} when it is constant on cosets of \(K\)
and therefore factors continuously through \(\pi_K\).  A previously excluded
closed family is removed from the admissible domain before the quotient is
formed.  Thus a reduction never discards a target-visible family without the
proved invariant identity that excludes that family.
\end{definition}

\begin{theorem}[Propagation of structural reductions]
\label{thm:structural-reduction-propagation}
Let \(K_0\subseteq K_1\subseteq Y\) be closed subspaces.  Totalize every
later map relative to both reductions.  On the zero descent-defect branch,
the canonical map is a topological linear isomorphism
\begin{equation}
 \label{eq:third-isomorphism-structural-reduction}
 (Y/K_0)/(K_1/K_0)\cong Y/K_1.
\end{equation}
On a nonzero descent-defect branch that defect is retained in its inherited
source cell.  Consequently a source direction closed at one stage may be propagated into
the next quotient and omitted from every later compatible calculation.  A
later invariant computed after reduction has the same value as the
corresponding invariant on \(Y\) after quotienting by all previously closed
directions.
\end{theorem}

\begin{proof}
The quotient map \(Y/K_0\to Y/K_1\) is continuous, open, and has kernel
\(K_1/K_0\).  The topological first isomorphism theorem gives
\eqref{eq:third-isomorphism-structural-reduction}.  Every zero-defect
downstream object factors through these quotient maps, so its reduced and
unreduced evaluations agree under the displayed isomorphism.  A failure to
factor is the nonzero descent-defect branch stipulated in the theorem.
\end{proof}

\begin{theorem}[Target-quotient coercivity]
\label{thm:target-quotient-coercivity}
Let \(Y\) be a Hilbert source space, let \(K\subseteq Y\) be a closed
target-null subspace, and let \(T:Y\to Z_\Sigma\) be bounded with
\(K\subseteq\ker T\).  Let \(q:Y\to[0,\infty]\) be an extended nonnegative
quadratic form with finite domain \(D(q)\) satisfying
\(D(q)+K\subseteq D(q)\), and define its relaxed quotient
form by
\begin{equation}
 \label{eq:relaxed-quotient-form}
 q_{\mathrm{red}}([y])
 :=\inf_{k\in K}q(y+k),
 \qquad [y]\in Y/K.
\end{equation}
If, on the required reduced domain,
\begin{equation}
 \label{eq:target-quotient-coercivity}
 \|[y]\|_{Y/K}^2\le Cq_{\mathrm{red}}([y]),
\end{equation}
then
\begin{equation}
 \label{eq:target-only-bound}
 \|Ty\|_{Z_\Sigma}^2
 \le \|\widetilde T\|^2Cq_{\mathrm{red}}([y])
 \le \|\widetilde T\|^2Cq(y),
\end{equation}
where \(\widetilde T:Y/K\to Z_\Sigma\) is the descended target map.  No
bound for the full norm \(\|y\|_Y\) follows or is needed.  When invoked by
the regularity calculus, \(Y/K\), \(q_{\rm red}\), and \(T\) are first
restricted to the contact-generated active source cell; failure of
\eqref{eq:target-quotient-coercivity} is its complementary local structural
value.
\end{theorem}

\begin{proof}
Because \(T\) annihilates \(K\), the quotient universal property gives the
bounded map \(\widetilde T\) with \(T=\widetilde T\pi_K\).  Hence
\[
 \|Ty\|\le\|\widetilde T\|\,\|[y]\|_{Y/K}.
\]
Apply \eqref{eq:target-quotient-coercivity}; the final inequality follows by
using \(k=0\) in \eqref{eq:relaxed-quotient-form}.
\end{proof}

\begin{example}[Strict gain over full-state coercivity]
\label{ex:strict-quotient-gain}
Let \(Y=K\oplus Z_\Sigma\), let \(T(k,z)=z\), and let
\(q(k,z)=\|z\|^2\).  Then
\eqref{eq:target-quotient-coercivity} holds with \(C=1\), while no inequality
\(\|(k,z)\|_Y^2\le C'q(k,z)\) can hold when \(K\ne\{0\}\).  The sequence
\((nk_0,0)\) may diverge in the full state norm but is exactly target-null.
Thus the quotient theorem proves every target estimate available from \(q\)
without controlling that divergent motion.
\end{example}

\begin{definition}[Structural evaluation invariants]
\label{def:structural-evaluation-invariants}
Work on a zero-descent-defect reduced space \(Y_{\mathrm{red}}\), with a
continuous target projection \(\Pi_\Sigma:Y_{\mathrm{red}}\to Z_\Sigma\).
The space and every family below are restricted to one propagated
same-origin target/source cell: a fixed nonempty source support, realization
mode, target-shell floor, and all earlier zero transcript coordinates.  No
extremum in this definition ranges over a source-transverse state.
The structural invariants used below have the following forms.
\begin{enumerate}[label=\textup{(I\arabic*)}]
\item For a retained target/source family \(\mathcal F\subseteq Y_{\mathrm{red}}\), its
  \emph{target closure invariant} is
  \begin{equation}
   \label{eq:target-closure-invariant}
   \mathfrak H_{\mathcal F}^{\Sigma}
   :=\overline{\operatorname{span}
      \Pi_\Sigma(\mathcal F)}^{Z_\Sigma}.
  \end{equation}
\item For nonnegative functionals \(D,N_+\) on a reduced target/source family
  \(\mathcal P\), the \emph{sharp domination radius} is
  \begin{equation}
   \label{eq:sharp-domination-radius}
   \alpha_*(N_+\mid D)
   :=\sup_{\substack{V\in\mathcal P\\D(V)>0}}
       \frac{N_+(V)}{D(V)}.
  \end{equation}
  The value is \(+\infty\) when some \(V\) has \(D(V)=0<N_+(V)\).
  For a normalized retained family \(\mathcal N\subseteq\mathcal P\), the
  \emph{physical-price gap} is
  \begin{equation}
   \label{eq:sharp-physical-price-gap}
   \delta_*(D\mid\mathcal N):=\inf_{V\in\mathcal N}D(V).
  \end{equation}
\item When \(\alpha_*<\infty\), let
  \(\widehat{\mathcal E}_{\alpha_*}\) be the closure in the canonical
  evaluation compactification of the normalized equality family
  \[
   \mathcal E_{\alpha_*}
   :=\{V\in\mathcal P:D(V)>0,\;
        N_+(V)=\alpha_*D(V)\}.
  \]
  Totalize the remainder relation generated by the equality identities.  Its
  ordinary graph image has \emph{equality--rigidity invariant}
  \begin{equation}
   \label{eq:equality-rigidity-invariant}
   \mathfrak H_{\mathrm{eq}}^{\Sigma}
   :=\overline{\operatorname{span}
     \Pi_\Sigma\mathscr R_{\mathrm{eq}}
       (\widehat{\mathcal E}_{\alpha_*})}^{Z_\Sigma}.
  \end{equation}
  and its failure and graph-boundary values enter
  \(\mathfrak H_{\rm sat}^{\Sigma}\).  For \(\alpha_*=+\infty\), the
  equality row is canonically inapplicable and the positive null-cost or
  unbounded-feeding class is already returned by the radius row.
  The equality identities may be saturation equations, conservation laws,
  modulation equations, Pohozaev identities, or Liouville constraints.
\end{enumerate}
Each invariant is computed after all zero-descent-defect source, gauge,
trace, and target-null reductions already proved have been propagated.
\end{definition}

\begin{theorem}[Structural normal form for invariant evaluation]
\label{thm:structural-invariant-normal-form}
Within the interfaces of this paper, every analytic input to target exclusion
has one of the following exact structural meanings.
\begin{enumerate}[label=\textup{(N\arabic*)}]
\item A retained target/source family is target-null on the reduced interface exactly when
  \(\mathfrak H_{\mathcal F}^{\Sigma}=0\).
\item The domination inequality \(N_+\le aD\) holds on the retained
  target/source family \(\mathcal P\)
  exactly when \(\alpha_*(N_+\mid D)\le a\).
\item The normalized target/source family \(\mathcal N\) has a uniform positive physical
  price exactly when \(\delta_*(D\mid\mathcal N)>0\).
\item A sharp equality branch has no target-visible remainder exactly when
  \(\mathfrak H_{\mathrm{eq}}^{\Sigma}=0\).
\end{enumerate}
Source closure and target-null residual calculations use \textup{(N1)};
resolvent, flux, and carrier identities generate \textup{(N2)}; physical
prices generate \textup{(N3)}; and zero-cost equations generate
\textup{(N4)}.  None is an endpoint.  Its nonclosing value supplies a
normalized witness to the activity lift, after which the successor and
viability operators continue the calculation.
\end{theorem}

\begin{proof}
Item \textup{(N1)} is the definition of the closed target projection in
\eqref{eq:target-closure-invariant}.  For \textup{(N2)}, divide by \(D(V)\)
when it is positive; the convention for a positive \(D\)-null mode gives the
remaining case.  Item \textup{(N3)} is the definition of the infimum in
\eqref{eq:sharp-physical-price-gap}.  Item \textup{(N4)} follows from
\eqref{eq:equality-rigidity-invariant}: its vanishing is precisely the
vanishing of every target projection of the equality remainder.  These are
the four structural coordinates occurring in
\cref{thm:no-target-visible-residual,thm:source-only-envelope,def:physical-price-interface,lem:compactness-rigidity-price}.
\end{proof}

\begin{definition}[Operational fast-track evaluator]
\label{def:operational-fast-track-evaluator}
For a source/mode cell \((A,m)\), propagate all previously vanished closed
subspaces and let \(Y_{A,m}^{\Sigma}\) be the resulting target-active
quotient.  Introduce the ordered row set
\begin{equation}
 \label{eq:fast-track-row-set}
 \begin{aligned}
 \mathscr R_{\rm FT}:=\{&\mathrm{form},\mathrm{grad},\mathrm{chain},
 \mathrm{bdry},\mathrm{tail},\\
 &\mathrm{target},\mathrm{feed},\mathrm{price},\mathrm{route},
 \mathrm{capacity},\\
 &\mathrm{scale},\mathrm{sat},\mathrm{fluct},\mathrm{cov},
 \mathrm{succ}\}.
 \end{aligned}
\end{equation}
Its fast-track value is the row-indexed tuple
\begin{equation}
 \label{eq:operational-fast-track-tuple}
 \mathfrak I_{A,m}^{\Sigma}
 :=\bigl(\mathfrak I_{A,m}^{\Sigma,r}\bigr)_
       {r\in\mathscr R_{\rm FT}},
\end{equation}
where the first five and covariance entries are target-active total-graph defect
quotients; \(\mathfrak H_{A,m}^{\Sigma}\) is the target closure;
\(\alpha^*\) is the feeding radius; \(\delta^*\) is the physical-price
infimum; \(\mathcal R\) is effective resistance; \(\operatorname{Cap}\) is
target capacity; \(\mathscr S\) is the pulled-back nonnegative scale-price
series; \(\mathfrak H_{\rm sat}^{\Sigma}\) is the projected saturation or
equality remainder; and \(\chi^\Sigma\) is the target-active covariance,
bracket, or fluctuation radius when the evolution is stochastic.  An
inapplicable row has the singleton value \(\{\bot_{\rm n/a}\}\) precisely
when its defining type is absent from the fixed PDE signature, as for the
fluctuation row of a deterministic equation.  Applicability is not a choice.
Every other entry is an exact quotient, graph defect,
infimum, supremum, kernel, capacity, resistance, or series value.  No
favorable analytic property occurs in the definition.
\end{definition}

\begin{definition}[Structural property compiler]
\label{def:structural-property-compiler}
For a row \(r\in\mathscr R_{\rm FT}\), let \(F_r^{\rm tot}\) be its
totalized interface restricted to the current presentation-generated
target-aligned record cell and let
\(\mathcal K_r^\Sigma\) be the closed ledger already propagated into that
cell.  Every supremum, infimum, kernel, graph value, and carrier comparison
below is taken on finite stopped records and their source-preserving closure;
there is no ambient solution-space quantifier.  The structural property
compiler \(\mathsf P_\Sigma\) assigns the
following exact invariant statements:
\begin{align}
 \mathsf P_\Sigma[\text{native realization or local existence}]
   &:=\mathfrak H_{{\rm form},A,m}^{\Sigma}=0,\label{eq:property-compiler-form}\\
 \mathsf P_\Sigma[\text{target-relative closed gradient range}]
   &:=\mathfrak H_{{\rm grad},A,m}^{\Sigma}=0,\label{eq:property-compiler-grad}\\
 \mathsf P_\Sigma[\text{target-relevant compactness of generated limits}]
   &:=\mathfrak H_{{\rm bdry},A,m}^{\Sigma}=0,\label{eq:property-compiler-compact}\\
 \mathsf P_\Sigma[\text{pre-contact chain rule}]
   &:=\mathfrak H_{{\rm chain},A,m}^{\Sigma}=0,\label{eq:property-compiler-chain}\\
 \mathsf P_\Sigma[\text{constraint-tail closure}]
   &:=\mathfrak H_{{\rm tail},A,m}^{\Sigma}=0,\label{eq:property-compiler-tail}\\
 \mathsf P_\Sigma[\text{generated operation descends}]
   &:=\mathfrak H_{{\rm cov},A,m}^{\Sigma}=0,\label{eq:property-compiler-cov}\\
 \mathsf P_\Sigma[N_+\le aD]
   &:=\alpha_{A,m}^*(N_+\mid D)\le a,\label{eq:property-compiler-radius}\\
 \mathsf P_\Sigma[\inf_{\mathcal N}D>0]
   &:=\delta_{A,m}^*(D\mid\mathcal N)>0,\label{eq:property-compiler-gap}\\
 \mathsf P_\Sigma[\text{routable target current}]
   &:=\mathcal R_{A,m}(\Sigma)<\infty,\label{eq:property-compiler-resistance}\\
 \mathsf P_\Sigma[\text{target polar}]
   &:=\operatorname{Cap}_{A,m}(\Sigma)=0,\label{eq:property-compiler-capacity}\\
 \mathsf P_\Sigma[\text{finite-budget scale exclusion}]
   &:=\mathscr S_{A,m}^{\Sigma}=+\infty,\label{eq:property-compiler-scale}\\
 \mathsf P_\Sigma[\text{sharp equality or saturation closes}]
   &:=\mathfrak H_{{\rm sat},A,m}^{\Sigma}=0,\label{eq:property-compiler-sat}\\
 \mathsf P_\Sigma[\text{fluctuation is target-null or physically routed}]
   &:=\chi_{A,m}^{\Sigma}=0\ \text{or its resistance is finite}.
   \label{eq:property-compiler-fluct}
\end{align}
The successor row is the finite-depth sequence
\(Z_{n+1,A,m}=\mathcal V_\eta(Z_{n,A,m})\); its two values are finite
pruning with its rational-cylinder transcript and cumulative joint saturation
of every retained cofinal transcript.  The latter is reprojected and is
either target-null, has a least unprocessed finite joint increment, or is a
primal--polar fixed structural atom under
\cref{thm:diagonal-target-aligned-saturation,thm:recursive-carrier-polar-exhaustion}.
Here every object is formed after quotienting by \(\mathcal K_r^\Sigma\).
Any algebraic identity derivable from the public weak equation may refine a
row into nonnegative descendants.  The compiler has no input slot for an external
compactness, coercivity, Liouville, or realization statement.
\end{definition}

\begin{definition}[Propagated physical fast-track ledger]
\label{def:propagated-physical-fast-track-ledger}
Order the rows as in \eqref{eq:fast-track-row-set}.  Begin with the closed
target-null ledger \(\mathcal K_0^\Sigma=\Null_\Sigma\).  If row \(i\) has
closed branch \(C_i\) and transport \(T_i\) into the next row, set
\begin{equation}
 \label{eq:propagated-physical-ledger}
 \mathcal K_{i+1}^\Sigma
 :=\overline{\mathcal K_i^\Sigma+T_i(C_i)}.
\end{equation}
All spaces, totalized interfaces, target projections, physical forms, and
subsequent invariants are then replaced by their descended objects on the
quotient by \(\mathcal K_{i+1}^\Sigma\).  If a map fails to descend, its
nonzero commutator with the quotient map is the next row's target-active
graph defect.  Thus quotient compatibility is tested by the construction
rather than imposed on it.
\end{definition}

\begin{theorem}[Analytic-property compilation]
\label{thm:analytic-property-compilation}
Every analytic property used by the continuum argument to decide whether a
presentation-generated source/mode cell can contribute to singular contact is equivalent to the
corresponding structural invariant statement in
\cref{def:structural-property-compiler}.  The opposite analytic property is
equivalent to the complementary invariant value and produces a tagged
activity profile and successor.  In particular, the fast track never presupposes closed
range, compactness, a chain rule, a coercive estimate, a positive price,
rigidity, polarity, finite resistance, summable tails, fluctuation control,
orbit realization, or continuation.
\end{theorem}

\begin{proof}
Existence and regularity properties are membership statements in the graph
of a partial map.  Totalization makes their negations explicit graph values,
and quotienting their closed graph boundary by \(\Null_\Sigma\) proves
\eqref{eq:property-compiler-form}--\eqref{eq:property-compiler-tail} and
\eqref{eq:property-compiler-cov}.
In the compactness row only target-visible loss of compactness survives;
target-null escape is already in \(J_{\rm res}\).  Thus the invariant tests
exactly the compactness needed downstream rather than ambient compactness of
the orbit.  The radius and gap equivalences follow
directly from the supremum and infimum definitions, including positive
\(D\)-null directions and zero-price sequences.  The spectral theorem gives
the resistance/range/harmonic trichotomy, with the harmonic part retained in
the target quotient.  Capacity zero is the definition of the polar branch;
positive capacity is its complementary target-compatible class.  A
nonnegative scale-price series either diverges or is finite.  Saturation
equations define a remainder map, so closure of the equality branch is
exactly vanishing of its target quotient.  Finally, stochastic drift,
bracket, covariance, and conditional-expectation contributions are currents
in the lifted state space; quotienting their target-null part and applying
the same radius/resistance construction gives
\eqref{eq:property-compiler-fluct}.  Every complementary value enters the
successor row.  Its finite-depth values and cumulative joint closure are
exhaustive by \cref{thm:diagonal-target-aligned-saturation}; compactness is
used only to reconstruct the compatible joint record, not to declare that
record terminal.  These cases exhaust
\eqref{eq:fast-track-row-set}.
\end{proof}

\begin{definition}[Target-relevant analytic statement grammar]
\label{def:target-relevant-statement-grammar}
Let \(\mathsf{An}_\Sigma(\mathfrak P)\) be the least grammar generated from
the fixed PDE signature by:
\begin{enumerate}[label=\textup{(A\arabic*)}]
\item membership in, or failure of, a native, limit, chain, trace, tail,
  normalization, or continuation graph;
\item vanishing or nonvanishing after the target quotient;
\item comparison of two nonnegative physical forms, including an infimum or
  sharp radius;
\item range, kernel, resistance, capacity, and harmonic statements for the
  generated operator or form;
\item convergence or divergence of a nonnegative physical-price series;
\item a saturation/equality remainder after sharp comparison;
\item covariance, bracket, compensator, or conditional-expectation statements
  in a generated stochastic lift;
\item membership in a finite-depth successor set, a cumulative joint-prefix
  state, or the fixed-atom locus of its graph-complete saturation.
\end{enumerate}
The grammar is closed under finite conjunction, sourcewise limit, pullback by
a generated chart, and descent through a generated quotient.  These closure
operations are themselves totalized.
\end{definition}

\begin{theorem}[Exhaustive fast-track compilation]
\label{thm:exhaustive-fast-track-compilation}
The property compiler extends uniquely to a total map
\begin{equation}
 \label{eq:total-fast-track-compiler}
 \mathsf P_\Sigma:
 \mathsf{An}_\Sigma(\mathfrak P)\longrightarrow
 \prod_{\substack{\varnothing\ne A\subseteq\Channels\\m\in\Modes}}
 \prod_{r\in\mathscr R_{\rm FT}}
 \operatorname{Val}\!\left(\mathfrak I_{A,m}^{\Sigma,r}\right).
\end{equation}
Every statement that can affect the contact identity or the physical
expenditure contradiction occurs in exactly one primitive row, possibly with
several source/mode tags.  Its negation is the complementary value of that
same row.  Hence no analytic property, compactness assertion, or
equation-specific estimate is an unregistered premise of the fast track.
\end{theorem}

\begin{proof}
Clauses \textup{(A1)}--\textup{(A8)} map respectively to the total-graph
rows; target row; feeding/price rows; gradient/routing/capacity rows; scale
row; saturation row; fluctuation row; and finite-depth successor row.
Source and realization tags are assigned before evaluation by
\cref{thm:operational-native-source-exhaustion,thm:realization-mode-exhaustion}.
Finite conjunction gives a product value.  Sourcewise limits remain in their
tagged closures.  Chart pullback appends \(\Abs/\Lift\), and quotient descent
either factors or produces its totalized descent defect.  Structural
induction therefore defines \(\mathsf P_\Sigma\) on the whole grammar and
leaves no unmatched constructor.

By normalized shell accountability, contact affects the argument only through
a target-active current or a totalized shell/interface defect.  By the
physical balance, exclusion affects it only through a nonnegative expenditure
comparison or series.  These are clauses \textup{(A1)}--\textup{(A6)};
stochastic and successor-path semantics are exactly
\textup{(A7)}--\textup{(A8)}.
Thus the grammar covers every input to the contact-versus-expenditure proof.
\end{proof}

\begin{theorem}[Closed-input rule]
\label{thm:closed-input-rule}
The value of \(\mathfrak I_{A,m}^{\Sigma,r}\) is determined by the fixed PDE
signature, target, physical balance, and the closures generated from them.
An auxiliary lemma may prove that value but cannot be added as data defining
it.  Any proposed object not generated by this algebra either:
\begin{enumerate}[label=\textup{(E\arabic*)}]
\item is derived from the fixed signature, in which case the source compiler
  and property compiler assign it to existing rows; or
\item changes the signature, in which case the conclusion concerns the
  enlarged equation/presentation and is not a conclusion about the original
  PDE.
\end{enumerate}
These alternatives exhaust the generated continuum algebra.
\end{theorem}

\begin{proof}
The evaluation compactification is generated by
\cref{def:structural-test-system}; the source syntax by
\cref{def:continuum-source-compiler}; and the analytic statement syntax by
\cref{def:target-relevant-statement-grammar}.  Membership in either syntax is
therefore determined before an invariant is evaluated.  Structural induction
gives \textup{(E1)}.  Failure of membership means that a new primitive map,
topology, law, or balance has been appended, which is precisely a change of
signature and gives \textup{(E2)}.
\end{proof}

\begin{theorem}[Ledger propagation and no re-estimation]
\label{thm:physical-ledger-propagation}
At each fast-track row, exactly one of the closed and complementary branches
is selected by its invariant value.  A closed branch propagated by
\eqref{eq:propagated-physical-ledger} remains zero in every later descended
target quotient.  If a downstream object does not descend, its descent
defect is retained in the appropriate existing source/mode cell.  Therefore
later rows neither assume nor re-estimate any property closed by an earlier
row.
\end{theorem}

\begin{proof}
For closed subspaces the third isomorphism theorem gives
\((Y/\mathcal K_i)/(\mathcal K_{i+1}/\mathcal K_i)\cong
Y/\mathcal K_{i+1}\).  Hence every propagated class is zero in all later
quotients.  The same statement for closed convex ledgers follows by passing
to the associated quotient seminorm.  Failure of descent is precisely a
nonzero commutator between the downstream map and the quotient map, which is
a totalized graph defect.  Sourcewise closure retains its ancestry.  The
exhaustive two-branch statement follows from
\cref{thm:analytic-property-compilation}.
\end{proof}

\begin{theorem}[Fast-track branch-to-successor completeness]
\label{thm:fast-track-branch-completeness}
For every target-active source/mode cell, the tuple
\(\mathfrak I_{A,m}^{\Sigma}\) returns one of the following exhaustive
structural branches.
\begin{enumerate}[label=\textup{(F\arabic*)}]
\item Each form, gradient, chain, boundary, tail, and covariance defect
quotient vanishes or survives as a target-visible total-graph defect with the
same source ancestry.
\item \(\mathfrak H_{A,m}^{\Sigma}=0\): the cell is target-null and closes.
Otherwise its nonzero quotient is retained.
\item \(\alpha_{A,m}^*\le1\): feeding is dominated by the generated carrier.
If \(\alpha_{A,m}^*>1\), or \(D=0<N_+\) for a member of the cell, the
maximizing sequence becomes a target-feeding activity profile.
\item \(\delta_{A,m}^*>0\): the cell has positive normalized price.
If \(\delta_{A,m}^*=0\), a zero-price sequence exists in the closure and is
retained with the same ancestry.
\item \(\mathfrak H_{{\rm eq},A,m}^{\Sigma}=0\): the sharp equality branch
closes.  Otherwise its target-visible equality remainder survives.
\item \(\operatorname{Cap}_{A,m}(\Sigma)=0\): the target is polar for that
cell.  Positive capacity supplies a normalized target-active activity profile.
\item \(\mathcal R_{A,m}(\Sigma)<\infty\): the source can be routed through
the generated geometry.  Infinite resistance supplies its harmonic component
to the activity lift.
\item For a scale cascade, the pulled-back series
\(\sum_j r_j^{\kappa_{A,m}}\delta_{A,m,j}\) either diverges, excluding that
finite-budget cascade, or is finite, yielding vanishing-price windows for the
ordered successor and recurrent-saturation compiler.
\item The saturation quotient \(\mathfrak H_{{\rm sat},A,m}^{\Sigma}\)
vanishes or retains the target-visible equality remainder.  Every Pohozaev,
Noether, virial, frequency, or entropy identity already generated by the
equation refines its nonnegative descendants; the remaining equality record becomes
a successor.
\item The fluctuation invariant \(\chi_{A,m}^{\Sigma}\) is target-null,
physically routed at finite resistance, or survives as a target-visible
covariance/bracket class in the existing lifted source cell.
\item The total weak-solution graph retains ordinary values and every
graph-boundary witness.  Both enter the successor graph; persistent boundary
values survive only through an infinite zero-price witness path.
\item An inapplicable row returns \(\{\bot_{\rm n/a}\}\); failure of a
type-present map to descend returns its graph defect and activates the same
successor operator.
\end{enumerate}
Every complementary branch is a value of the same invariant used by the
closing branch.  The tuple is not a verdict: after each row, the canonical
activity map sends every nonzero target quotient to
\(\mathcal R_\eta\), and the viability sets \(Z_n\) either remove it at finite
depth or construct its recurrent zero-price mechanism.  Before either action,
the row is converted to the target-local normal form
\eqref{eq:target-local-fact-normal-form}; hence no row acts on an
equation-wide class or charges a source-transverse positive quantity.
\end{theorem}

\begin{proof}
The graph-defect row follows from totalization.  The target, radius, gap, and
equality rows follow from
\cref{thm:structural-invariant-normal-form} and the definitions of the
closed quotient, supremum, infimum, and equality remainder.  In
the radius row, the convention \(\alpha^*=+\infty\) records a positive
\(D\)-null feeding direction.  If \(\delta^*=0\), the definition of the
infimum gives a sequence with prices tending to zero; sourcewise closure
retains its limit records.

Zero capacity is the polar branch of the associated form.  Its complement is
positive capacity and supplies an activity profile.  The spectral theorem gives the finite
resistance/range branch and the orthogonal harmonic remainder.  Finally a
nonnegative series either diverges or has finite sum.  The physical pullback
identity makes these precisely the excluded branch and the
successor-activation branch.  The latter is closed by
\cref{thm:finite-price-saturation-route}; it is not a terminal formation
class.  The saturation and fluctuation alternatives are
\cref{thm:analytic-property-compilation}.  The singleton inapplicable value
and totalized descent defect exhaust the remaining interface cases.  Applying
the activity lift and \cref{thm:finite-depth-structural-evaluator} to
each nonzero quotient proves the final constructive statement.
\end{proof}

\begin{corollary}[Estimate compression after structural derivation]
\label{cor:estimate-replacement-invariant-evaluation}
Once a graph relation, positive-form identity, target-null quotient, or
zero-set relation has been derived from the symbolic presentation, subsequent
rows use only that structural output.  Formal identities and universal
functional calculus refine the generated positive cone directly.  When they
do not remove a row, its graph-boundary, positive-null, or zero-price witnesses
are generated and passed to the successor calculus.
Closed outputs are propagated by \cref{thm:physical-ledger-propagation}.
\end{corollary}

\begin{definition}[Canonical symbolic structural transcript]
\label{def:canonical-structural-transcript}
For the fixed PDE, complete continuation target, carrier, and derived ceiling,
define
\begin{equation}
 \label{eq:canonical-structural-transcript}
 \Theta_\Sigma^{\rm syn}(\mathfrak P,\Budget)
 :=\left(
   \mathsf{Syn}(\mathfrak P),
   \mathsf{Graph}^{\rm tot}(\mathfrak P),
   \mathsf{Pos}(\mathcal K),
   \Pi_\Sigma,
   \mathsf{Cov}(\mathsf{Op}),
   \mathsf{Anc}
 \right).
\end{equation}
The entries are respectively the typed equation syntax, total graphs of its
generated operations, the cone of nonnegative terms occurring in the proved
carrier identity, the generated target morphism, the total covariance graphs,
and the event/source ancestry maps.  Exact extrema, capacities, resistances,
realization images, contact orbits, continuation verdicts, and Liouville
conclusions are not entries.  They are semantic outputs of rows constructed
from \(\Theta_\Sigma^{\rm syn}\).
\end{definition}

\begin{theorem}[Public-signature structural sufficiency]
\label{thm:estimate-erasure}
The symbolic transcript functorially determines the source cells, totalized
interfaces, target-null quotient, carrier complex, activity measures, and
successor relation used by the regularity argument:
\begin{equation}
 \label{eq:estimate-erasure-factorization}
 \mathsf{Compiler}_\Sigma
 =\overline{\mathsf{Compiler}}_\Sigma
   \circ\Theta_\Sigma^{\rm syn}.
\end{equation}
Every row is constructed together with its ordinary and complementary witness
graphs.  No exact row value or actual-contact image is consulted in this
construction.  Two presentations with isomorphic symbolic transcripts
therefore have isomorphic prospective successor and zero-cost fixed-point
algebras.  An analytic derivation can only prove a relation already present in
the transcript or refine its generated positive cone; it cannot insert a
favorable value.
\end{theorem}

\begin{proof}
Proceed by structural induction over the typed syntax and then by the fixed row
order \eqref{eq:fast-track-row-set}.  Each constructor supplies its ordinary
graph, explicit failure value, source transport, and event ancestry.  Graph
closure supplies the corresponding boundary witnesses.  The target morphism
forms the null quotient, and the positive carrier cone supplies the
expenditure coordinate.  Quotient descent either factors or produces its
commutator graph.  These operations construct both branches of every row
without deciding which branch contains a given future orbit record.

The activity and successor constructions use only the generated metric,
target shell, graph witnesses, and event ancestry.  The prospective
fixed-point construction below uses their closed images, not the image of
actual contact orbits.  This proves \eqref{eq:estimate-erasure-factorization}
and the isomorphism statement.
\end{proof}

\begin{corollary}[Closed structural input]
\label{cor:closed-structural-input}
A PDE specialization supplies only the target-blind public signature of
\cref{def:target-blind-public-input}.  The continuum algebra generates
\eqref{eq:canonical-structural-transcript}.  Every identity obtainable by
formal manipulation of that signature or by the fixed universal library is
inserted automatically as a graph relation or positive-descendant refinement.
The ordinary graph and every first-failure witness are inputs to the canonical
activity lift and successor operator of
\cref{def:ordered-continuum-successors,def:retained-successor-relation}.  Their
descending zero-cost closure is computed by
\cref{def:total-successor-saturation-algebra}.  Thus every generated branch is
reduced to a target-null class or to an explicit recurrent zero-cost mechanism;
there is no unevaluated row and no additional invariant value to be supplied.
\end{corollary}

\begin{proposition}[Positive-separator evaluation of an equality branch]
\label{prop:positive-equality-separator}
Totalize the separator relation on
\(\widehat{\mathcal E}_{\alpha_*}\).  On its zero graph-defect branch the
identity is
\begin{equation}
 \label{eq:positive-equality-separator}
 0=Q_\Sigma\!\left(
    \Pi_\Sigma\mathscr R_{\mathrm{eq}}(Z)\right)
   +\sum_{j=1}^{m}q_j(Z),
\end{equation}
where \(Q_\Sigma\ge0\) is positive definite on \(Z_\Sigma\) and every
\(q_j\ge0\).  Then
\(\mathfrak H_{\mathrm{eq}}^{\Sigma}=0\).
\end{proposition}

\begin{proof}
Every term on the right of \eqref{eq:positive-equality-separator} is
nonnegative, so each term vanishes.  Positive definiteness gives
\(\Pi_\Sigma\mathscr R_{\mathrm{eq}}(Z)=0\) for every \(Z\).  Taking the
closed linear span proves the claim.
\end{proof}

\section{Coherent carrier passivity and structural successor closure}
\label{sec:carrier-passivity-successors}

The scalar price infimum of
\cref{def:structural-evaluation-invariants} is not the terminal evaluator for
a finite-price cascade.  A finite measure on infinitely many disjoint target
windows produces a vanishing-price subsequence.  The continuum analogue of
the residual closure in the discrete algebra is to follow that subsequence
through all compact, concentrating, escaping, diffuse, gauge, and equality
outputs.  This section constructs that closure directly.

\begin{definition}[Coherent carrier complex]
\label{def:coherent-carrier-complex}
Let \(u\) be one physical orbit.  A \emph{coherent carrier complex} consists
of a locally finite rooted tree \(\mathcal T\), cells \(Q_\alpha\), and
nonnegative localization tests \(\psi_\alpha\) with the following
properties.
\begin{enumerate}[label=\textup{(C\arabic*)}]
\item Every child is a subevent of its parent and, on the region where the
  local balance is evaluated,
  \begin{equation}
   \label{eq:carrier-partition}
   \psi_\alpha=
   \sum_{\beta\in\operatorname{ch}(\alpha)}\psi_\beta.
  \end{equation}
\item All cells are restrictions, chart images, or normalized descendants of
  the same orbit.  Their transition maps satisfy the cocycle composition
  law.  Galilean, gauge, and quotient representatives are fixed compatibly on
  overlaps.
\item The weak equation gives a local carrier relation, linear in the test,
  \begin{equation}
   \label{eq:local-carrier-relation}
   K_\psi(t)+D_I(\psi)+M_I(\psi)
   =K_\psi(s)+J_I(\psi)+E_I(\psi).
  \end{equation}
  Here \(D_I,M_I\ge0\) are respectively the intrinsic production and the
  weak energy-defect measure, \(J_I\) is the signed transport current, and
  \(E_I\) is true exterior input.  An inequality is written as equality by
  adjoining its nonnegative defect to \(M_I\).
\item The normalized target observable is constant under the allowed
  transition maps.  Consequently target retention is inherited by a child
  selected from a target-retaining parent.
\end{enumerate}
The complex is generated from the shell cascade by taking the maximal
locally finite family of compatible target-retaining descendants.  A failure
of local finiteness, the partition relation, transition coherence, or gauge
compatibility is a totalized \(\Lift/\Bdry\) successor and is not an
additional admissibility condition.
\end{definition}

\begin{theorem}[Signed carrier gluing]
\label{thm:signed-carrier-gluing}
Let \(\mathcal T_0\subset\mathcal T\) be a finite rooted subtree.  Sum
\eqref{eq:local-carrier-relation} over \(\mathcal T_0\) before taking positive
parts.  Every current through an internal face cancels with the same current
with opposite orientation in the adjacent cell.  In particular, conservative
transport, pressure or constraint transport, chart translation, gauge
motion, cutoff motion, and exact scale motion contribute only through the
root, the leaves, or the true exterior boundary.  The result has the form
\begin{equation}
 \label{eq:signed-carrier-tree-balance}
 K_{\rm leaf}(\mathcal T_0)
 +\mu_{\Geom}(\mathcal T_0)
 +\mu_{\rm def}(\mathcal T_0)
 =K_{\rm root}
 +I_{\rm ext}(\mathcal T_0)
 +H_{\rm sc}(\mathcal T_0),
\end{equation}
where the first two measures on the left are nonnegative.  The term
\(H_{\rm sc}\) is the cohomology class of scale motion that is not an edge
coboundary in the chosen scale-critical carrier.  It has inherited
\(\Caus/\Lift\) ancestry.  It vanishes when the scale current admits a
coherent carrier potential.  No sum of the absolute values of the internal
currents occurs in \eqref{eq:signed-carrier-tree-balance}.
\end{theorem}

\begin{proof}
Linearity in \(\psi\) and \eqref{eq:carrier-partition} imply that the carrier,
production, defect, transport, and exterior terms of a parent are the sums of
the corresponding child terms.  Give every common face the orientation
induced by its parent.  The child has the opposite induced orientation, so
its signed flux cancels the parent flux.  Pressure-gauge constants cancel
because their weak pairing with the divergence-free or constraint-compatible
test is zero.  Transition coherence makes translation, gauge, and exact scale
terms edge coboundaries; the sum of a coboundary over a finite tree is its
root value minus its leaf values.  The only possible scale term left after
this cancellation is its class modulo the range of the edge coboundary,
which is \(H_{\rm sc}\).  The intrinsic production and weak balance defect
retain their signs and add over the cells.  This proves
\eqref{eq:signed-carrier-tree-balance}.
\end{proof}

\begin{definition}[Target-aligned physical charge]
\label{def:target-aligned-physical-charge}
Let \(Q\) be a pre-contact event selected by a shell covector \(D\Phi\).
First glue every compatible internal face and form the signed target current
on the finite-value branch,
\begin{equation}
 \label{eq:aggregate-target-current}
 q_{\Phi,Q}:=
 \sum_{c\in\{\Geom,\Caus,\Abs,\Lift,\Bdry\}}
 \int_Q\langle D\Phi,J^c\rangle .
\end{equation}
On a variation-boundary branch, replace this finite sum by the joint
shell-current owner readout \(q_{\Phi,Q}^{\rm own}:=\overline Q^o(1)\ge1/5\)
and its increasing bounded rational truncations.  Every charging relation is
closed jointly with the owner coordinate; no extended signed sum is formed.
A nonnegative measure \(P\) is an \emph{admissible charge} for this event
only if a generated local carrier identity proves
\begin{equation}
 \label{eq:target-aligned-charge-rule}
 (q_{\Phi,Q})^+
 \le \Delta K_Q+P(Q)+I_{\rm ext}^+(Q)+n_Q,
 \qquad n_Q=\langle D\Phi,N_Q\rangle=0,
 \quad N_Q\in\Null_\Sigma,
\end{equation}
with the carrier increments telescoping on same-origin concatenation and
\(P\) dominated by the public physical ledger.  Positive quantities not
appearing in such an identity are structural coordinates, not expenditure.
Positive parts of individual internal currents are never charged before
signed gluing.  On the boundary-owner branch,
\eqref{eq:target-aligned-charge-rule} denotes the compatible family of
bounded truncation inequalities; monotone passage to the owner value is made
only after the carrier domination has been established.

Equation \eqref{eq:target-aligned-charge-rule} is an eventwise
factorization, not yet a physical-price inequality.  The measure \(P\)
becomes an admissible target price only after the generated storage row
either represents the cumulative \(\Delta K_Q\) terms as a coboundary with a
finite root--leaf contribution on the selected same-origin family, or retains
every positive storage return as a target-aligned \(\Caus\)-headed
successor.  No telescoping is inferred from disjointness alone.
\end{definition}

\begin{theorem}[Local target-aligned charging rule]
\label{thm:local-target-aligned-charging-rule}
Every physical charge used by the successor algebra is obtained by
\cref{def:target-aligned-physical-charge}.  On a disjoint same-origin shell
cascade, summing \eqref{eq:target-aligned-charge-rule} cancels target-null
terms and exposes storage increments, physical production, and true exterior
input separately.  Internal carrier increments cancel only after the
intervening gaps are inserted into the same oriented carrier complex.  The
remaining root--leaf storage is evaluated by its coboundary/return row.
Consequently repeated target progress is paid by the restriction of the
public physical measure and true positive exterior input precisely on the
zero-storage-transfer cell; every nonzero exact storage transfer is itself a
target-aligned successor.  If no generated identity dominates a nonzero
target current, that current is likewise passed to the polarity stratifier;
it is not assigned zero price and is not charged to an unrelated positive
invariant.
\end{theorem}

\begin{proof}
The shell identity exposes finite \(q_{\Phi,Q}\), or its closed owner
readout, before contact.  Signed carrier
gluing gives the aggregate current and the root--leaf form
\eqref{eq:signed-carrier-tree-balance}.  Restriction to the target quotient
annihilates \(N_Q\), hence \(n_Q=0\).  Add the identities over pairwise
disjoint same-origin events and insert the oriented gap cells between
consecutive events.  The full complex telescopes, but removing the gaps
leaves exactly their signed storage-transfer current; hence disjointness
alone does not cancel it.  Countable additivity bounds the \(P\)-terms by the
single physical ledger and the exterior terms by its declared input
component.  Apply the generated coboundary graph to the remaining storage.
Its zero or root--leaf-bounded value gives the admitted physical charge,
while a nonzero signed return value gives the storage successor.  The grammar
admits no other rule for inserting a quantity into the ledger.  Failure to
derive the displayed domination has a nonzero total-graph coordinate, whose
polarity split gives the stated successor.
\end{proof}

\begin{definition}[Carrier cohomology and aggregate feeding]
\label{def:carrier-cohomology}
Let \(\mathscr W_\eta\) be the saturated closure of normalized windows whose
target observable is at least \(\eta>0\), after all generated symmetry and
gauge quotients.  Include the values of the carrier, production, exterior
input, and every totalized graph coordinate in the evaluation topology.
Let \(S\) be the successor or unit-window shift on its ordinary graph.

On a coherent ordinary branch, write \(F_{\rm sc}\) for the aggregate signed
inward scale-space feed left after \cref{thm:signed-carrier-gluing}.  Its
target-active carrier cohomology is
\begin{equation}
 \label{eq:carrier-cohomology}
 \mathfrak H_{\rm car}^{\Sigma}
 :=\overline{\Pi_\Sigma F_{\rm sc}}
 \big/
 \overline{\Pi_\Sigma\{\Psi-\Psi\circ S:\Psi\in C(\mathscr W_\eta)\}}.
\end{equation}
The quotient is totalized when \(S\), the carrier, or a transition is partial.
Its graph-boundary values retain their source ancestry.  The ordinary
carrier-deficit cocycle is
\begin{equation}
 \label{eq:carrier-excess-cocycle}
 g_{\rm car}
 :=F_{\rm sc}-D-M.
\end{equation}
For a finite block, its sum is the aggregate inward feeding minus production,
up to root--leaf storage and true exterior input.  Coboundaries disappear in
block averages and under invariant measures.  If
\([F_{\rm sc}]=0\) in \eqref{eq:carrier-cohomology}, then for every
\(\varepsilon>0\) there is a continuous \(\Psi_\varepsilon\) whose coboundary
approximates the feed class within \(\varepsilon\); no exact correction is
asserted.
\end{definition}

\begin{theorem}[Target-relative passivity]
\label{thm:target-relative-passivity}
Let \(\mathcal K_\eta\subseteq\mathscr W_\eta\) be a compact
\(S\)-invariant ordinary carrier class generated by the activity lift.  On the
ordinary signed-carrier branch, finite exterior input is already part of the
derived physical budget.  Then every
\(S\)-invariant Borel probability measure \(\nu\) on \(\mathcal K_\eta\)
satisfies
\begin{equation}
 \label{eq:invariant-passivity}
 \int_{\mathcal K_\eta}g_{\rm car}\,\dd\nu\le0.
\end{equation}
Since every invariant measure annihilates continuous coboundaries, this is
equivalent to
\begin{equation}
 \label{eq:invariant-production-passivity}
 \int F_{\rm sc}\,\dd\nu
 \le\int(D+M)\,\dd\nu.
\end{equation}
Equality in \eqref{eq:invariant-passivity} is the zero carrier-deficit branch.
It enters the zero-density successor fixed point together with its
unnormalized block ledger.  It implies vanishing of a
particular production term only when the generated consequence algebra
contains a domination identity \(p=q+r\) with that production as \(q\).
Even then, the conclusion is vanishing density on the recurrent support, not
vanishing of the cumulative production or source cocycle of the originating
path.
Every measure and current in this statement is restricted to
\(\mathcal K_\eta\), generated by one retained target/source cell; the
theorem has no consequence for invariant measures outside that cell.
\end{theorem}

\begin{proof}
Apply \eqref{eq:signed-carrier-tree-balance} to the first \(N\) consecutive
cells of a coherent path.  Internal currents cancel and the carrier terms
telescope.  The carrier is bounded on the compact evaluation class.  The exterior input has finite total mass, so its
contribution divided by \(N\) tends to zero.  Integrating the block identity
against an invariant measure cancels the endpoint terms exactly and gives
\eqref{eq:invariant-passivity}, which is
\eqref{eq:invariant-production-passivity}.  Continuous coboundaries and their
uniform closure have zero integral against every invariant measure, justifying
the cohomological quotient without producing an exact potential.
This averaging proves only the stationary density identity.  The finite-block
exterior, production, and source totals remain in the augmented occurrence
ledger and are restored before any target projection.
\end{proof}

\begin{corollary}[Totalized passivity branch]
\label{cor:totalized-passivity-branch}
For every target-retaining carrier route, exactly one of the following is
generated:
\begin{enumerate}[label=\textup{(T\arabic*)}]
\item the coherent ordinary branch of
  \cref{thm:target-relative-passivity};
\item a first failure of local finiteness, carrier partition, transition
  coherence, gauge descent, compact retention, or carrier correction,
  together with its closed-graph witness descendants;
\item a target-null output removed by the generated quotient.
\end{enumerate}
The second output enters \(\mathcal R_\eta\); the third closes.  Positive
exterior variation has \(\Bdry\) ancestry and is eligible for charge only
through the local target-aligned carrier identity.  Consequently compactness,
coherence and carrier cohomology are processed by the successor operator.
\end{corollary}

\begin{proof}
Totalize each partial relation in the order displayed in
\cref{def:coherent-carrier-complex,def:carrier-cohomology}.  If all ordinary
values occur, \cref{thm:target-relative-passivity} applies.  The first
nonordinary value has a witness sequence in its closed total graph and
therefore defines the successor in
\cref{def:total-successor-saturation-algebra}.  Positive exterior variation
enters the physical measure through
\cref{thm:local-target-aligned-charging-rule}.  The first-failure alternatives are exhaustive
because all graph coordinates belong to the compact activity lift.
\end{proof}

\begin{definition}[Carrier saturation algebra]
\label{def:carrier-saturation-algebra}
The \emph{carrier saturation algebra} is the closed target projection
\begin{equation}
 \label{eq:carrier-saturation-algebra}
 \mathfrak H_{\rm eq,car}^{\Sigma}
 :=\overline{\operatorname{span}}
 \left\{
   \Pi_\Sigma\supp\nu:
   \begin{array}{l}
   \nu\text{ is invariant on a generated compact successor core},\\[-1mm]
   \displaystyle\int g_{\rm car}\,\dd\nu=0
   \end{array}
 \right\}.
\end{equation}
The support is evaluated through all five source tags and all totalized
zero-deficit equations.  It may contain equilibria, relative equilibria, or
nonstationary recurrent conservative dynamics.  The subsequent zero-cost
fixed point, rather than an assumed rigidity statement, decides its target
projection.
\end{definition}

\begin{theorem}[Ergodic strict passivity or saturation]
\label{thm:strict-passivity-or-saturation}
Let \(\mathcal K_\eta\) be as in
\cref{thm:target-relative-passivity}, and set
\begin{equation}
 \label{eq:maximal-invariant-carrier-average}
 \beta_\eta
 :=\sup_{\nu\in\mathcal P_S(\mathcal K_\eta)}
   \int g_{\rm car}\,\dd\nu,
\end{equation}
where \(\mathcal P_S(\mathcal K_\eta)\) is the compact set of
\(S\)-invariant Borel probability measures.  Then
\(\beta_\eta\le0\), and exactly one of the following holds.
\begin{enumerate}[label=\textup{(P\arabic*)}]
\item \(\beta_\eta<0\).  There exist an integer \(N_\eta\ge1\) and
  \(\delta_\eta>0\) such that
  every \(N_\eta\)-block retained in \(\mathcal K_\eta\) satisfies
  \begin{equation}
   \label{eq:strict-block-passivity}
   -\sum_{k=0}^{N_\eta-1}g_{\rm car}(S^kw)
   \ge\delta_\eta.
  \end{equation}
\item \(\beta_\eta=0\).  The class \(\mathcal K_\eta\) carries an invariant
  maximizing measure and its
  target projection belongs to
  \(\mathfrak H_{\rm eq,car}^{\Sigma}\).
  The maximizing measure is only a support selector; the original block
  cocycle remains attached without division by \(N\).
\end{enumerate}
If \(\mathfrak H_{\rm eq,car}^{\Sigma}=0\), only \textup{(P1)} is possible.
The positive block price is therefore derived from the absence of a
target-visible equality state; no contact margin or compactness gap is inserted.
\end{theorem}

\begin{proof}
Let \(h=g_{\rm car}\).  Compactness of
\(\mathcal P_S(\mathcal K_\eta)\) and continuity of the integral give a
maximizing measure.  By \cref{thm:target-relative-passivity},
\(\beta_\eta\le0\).

Suppose \(\beta_\eta<0\) but \textup{(P1)} fails.  Then there are
\(w_n\in\mathcal K_\eta\) with block length \(N_n=n\) and total carrier
deficit smaller than \(n^{-1}\).  In particular the mean of \(h\) has
lower limit at least zero.  Form
the empirical measures
\[
 \nu_n:=\frac1{N_n}\sum_{k=0}^{N_n-1}\delta_{S^kw_n}.
\]
Compactness of probability measures on \(\mathcal K_\eta\) gives a weakly
convergent subsequence.  The endpoint identity
\[
 \int f\circ S\,\dd\nu_n-\int f\,\dd\nu_n
 =\frac{f(S^{N_n}w_n)-f(w_n)}{N_n}
\]
shows that its limit \(\nu\) is invariant.  Continuity of the evaluation
coordinates and the lower bound give \(\int h\,\dd\nu\ge0\), while
\cref{thm:target-relative-passivity} gives the reverse inequality.  Hence
\(\int h\,\dd\nu=0\), contradicting
\(\beta_\eta<0\).  Thus \textup{(P1)} holds.  If
\(\beta_\eta=0\), a maximizing invariant measure is a saturation measure and
gives \textup{(P2)}.  Conversely, an invariant saturation measure contradicts
any uniform positive block deficit after integration.  Vanishing of
\(\mathfrak H_{\rm eq,car}^{\Sigma}\) excludes a target-retaining instance of
\textup{(P2)}.
\end{proof}

\subsection{Target-aligned dual-carrier compilation}
\label{subsec:target-aligned-dual-carrier}

The preceding passivity theorem is used only after contact has selected a
source and a target covector.  This is the continuum selector cut.  It avoids
an equation-wide classification of equality states.

\begin{definition}[Active branch signature]
\label{def:active-branch-signature}
An \emph{active branch signature} is a tuple
\begin{equation}
 \label{eq:active-branch-signature}
 \mathfrak b=(A,m,\Gamma,\eta,\mathbf R),
\end{equation}
where \(A\ne\varnothing\) is a persistent source support, \(m\) is its
eventual realization mode, \(\Gamma\) is a generated target-shell family,
\(\eta>0\) is its retained source floor, and \(\mathbf R\) is the same-origin
cascade tag.  The branch ideal \(I_{\mathfrak b}\) is generated by the public
equation and by the relations already established on this branch:
source support \(A\), mode \(m\), shell pairing at least \(\eta\), chart and
gauge choices, event order, and every earlier zero graph coordinate.
No relation outside this list is inserted into \(I_{\mathfrak b}\).
\end{definition}

\begin{proposition}[Branch-local consequence rule]
\label{prop:branch-local-consequence-rule}
Let \(q=0\), \(q\ge0\), or \(q_1\le q_2\) be derivable from the public weak
equation modulo \(I_{\mathfrak b}\).  Then that relation holds on every
ordinary or zero-defect realization of \(\mathfrak b\).  Its failure on a
total-graph boundary is the next typed successor.  Consequently a relation
need only be proved on the active branch on which its antecedent has already
been generated; no assertion about states outside that branch is required.
\end{proposition}

\begin{proof}
Every generator of \(I_{\mathfrak b}\) vanishes on an ordinary realization
of the branch, so substitution in the finite weak-equation derivation proves
the relation there.  The closed total graph records the first generator that
fails under passage to a limit.  Source-preserving completion assigns that
failure its inherited source and realization mode.  These two cases exhaust
the totalized branch.
\end{proof}

\begin{definition}[Target-local structural stratum]
\label{def:target-local-structural-stratum}
For an active branch \(\mathfrak b=(A,m,\Gamma,\eta,\mathbf R)\), let
\(\mathbb G_{\mathfrak b}\) be the total evaluation graph generated by the
public weak equation, its five source currents, the shell family, and all
operations already used on the branch.  Its \emph{target-local structural
stratum} is
\begin{equation}
 \label{eq:target-local-structural-stratum}
 X_{\mathfrak b}:=
 \{z\in\mathbb G_{\mathfrak b}:I_{\mathfrak b}(z)=0,
       \ Q_{A,\Phi}(z)\ge\eta,\
       \operatorname{Origin}(z)=\mathbf R\}.
\end{equation}
All consequence, compactness, production, and equality operations below are
performed after restriction to \(X_{\mathfrak b}\).  Values outside this
stratum are irrelevant to the contact branch and are never quantified over.
\end{definition}

\begin{definition}[Generated polarity stratification]
\label{def:generated-polarity-stratification}
Fix once and for all a countable enumeration
\((F_j)_{j\ge0}\) of the scalar rational-cylinder coordinates generated by
the public equation and the target-cyclic grammar.  Vector, tensor, measure,
and operator coordinates enter through their rational scalar probes.  Replace
an unbounded real coordinate by
\(\widehat F_j=F_j/(1+|F_j|)\), and retain nonordinary values in its total
graph.  On a target-local stratum, the \(j\)-th \emph{polarity split} is the
disjoint ordered partition
\begin{equation}
 \label{eq:generated-polarity-split}
 \{\widehat F_j=0\},\qquad
 \{s\widehat F_j\in[2^{-k},2^{-k+1})\}
 \ (s\in\{-1,1\},\ k\ge1),\qquad
 \partial_{\rm gr}F_j .
\end{equation}
Endpoints are assigned to the leftmost displayed cell, and
\(\partial_{\rm gr}F_j\) is further split by the first failed public graph
operation and its inherited source word.  Intersect every cell with
\(X_{\mathfrak b}\).  First dovetail the finite identities and cylinder
moments of the parent local Archimedean cone.  A finite derivation of \(-1\)
prunes the parent cell.  Otherwise compatible moments reconstruct a positive
local functional; disintegrate it over the countable partition
\eqref{eq:generated-polarity-split} and select the least cell of positive
target mass.  That cell is a target-active successor with the same origin and
with its shell current renormalized to the universal source floor.  Thus no
semantic emptiness or exact-zero oracle is used.

Thus both possession and failure of an analytic property are structural
values.  For example, closed descent is the zero cell of its cokernel
coordinate; nonclosed descent is a signed dyadic or graph-boundary cell, not
an omitted hypothesis.
\end{definition}

\begin{definition}[Resolved-prefix rank]
\label{def:resolved-prefix-rank}
The enumeration in
\cref{def:generated-polarity-stratification} is dovetailed with the
enumerations of weak words, source descendants, quotient operations, carrier
rows, and total-graph defects.  A branch has \emph{resolved-prefix rank}
\begin{equation}
 \label{eq:resolved-prefix-rank}
 \rho(\mathfrak b):=sup\{N:\text{the polarity cell of every }F_j,
                         \ j<N,\text{ belongs to }I_{\mathfrak b}\}.
\end{equation}
When \(\rho(\mathfrak b)<\infty\) and
\(F_{\rho(\mathfrak b)}\) is split, every retained proper successor
adjoins its selected cell to the branch ideal and therefore has larger
resolved prefix.  If that split generates new operations, they enter later
in the same universal dovetailing.  The rank is consequently constructed by
the algebra; it is not presentation data and no well-foundedness premise is
required.
\end{definition}

\begin{theorem}[Target-local stratification exhaustiveness]
\label{thm:target-local-stratification-exhaustiveness}
Let \(\mathfrak b\) be a same-origin active contact branch.  At its next
unprocessed coordinate exactly one of the following occurs:
\begin{enumerate}[label=\textup{(L\arabic*)}]
\item a finite local ordered identity proves that the current target cell is
empty or target-null;
\item a generated carrier identity charges the retained cell to the finite
physical ledger;
\item a zero, signed dyadic, or graph-boundary cell has nonzero target
projection and becomes a proper successor \(\mathfrak b'\) with
\(\rho(\mathfrak b')>\rho(\mathfrak b)\).
\end{enumerate}
The alternatives exhaust ordinary values, loss of every analytic property
used by the derivation, and all limit values of the coordinate.  No global
estimate and no target-nullity statement is an input.
\end{theorem}

\begin{proof}
The total graph makes \eqref{eq:generated-polarity-split} exhaustive, and the
endpoint convention makes it disjoint.  Apply the Archimedean proof/model
alternative to the parent target cell.  Its finite \(-1\)-identity gives
\textup{(L1)}.  On the complementary cell, evaluate its target detector on
the presentation-generated reachable domain.  By
\cref{thm:reachable-source-saturation-finite-occurrence}, positivity returns
a finite stopped same-origin occurrence with its shell covector, complete
frontier, and source word.  The finite affine identity or closed owner
relation acts on that occurrence and restores the fixed floor.  If it is
dominated by a public carrier row it is \textup{(L2)}.  Otherwise its failed
polarity relation is appended to the occurrence record, producing
\textup{(L3)} and increasing the resolved prefix.  Total graph ancestry
proves the assertion for failed analytic properties.
\end{proof}

\begin{theorem}[Diagonal target-aligned saturation]
\label{thm:diagonal-target-aligned-saturation}
An infinite chain of nonempty proper successors selected from one contact
branch is evaluated by first forming its cumulative same-origin joint record,
closing it under every mixed public operation and total graph, and only then
reapplying the target projection.  Exactly one of the following occurs:
\begin{enumerate}[label=\textup{(D\arabic*)}]
\item the cumulative target projection vanishes, so all induced source
pairings with the retained shell covectors vanish along the retained path and
shell contact is impossible;
\item a least unprocessed target increment occurs on a finite joint prefix;
its polarity, complete mixed-operation word, and inherited five-channel
ancestry form the next proper successor;
\item every finite increment is source-resolved, joint saturation is fixed,
every source owner is fixed under recursive carrier-polar exhaustion, and the
totalized ordinary realization row reconstructs a joint-saturated
  primal--polar five-channel prospective source atom retaining the cofinal
  shell family and its origin polarity.
\end{enumerate}
The third conclusion is a saturated source-profile witness.  It is an
explicit singular mechanism only in the independent-origin fiber of
\cref{thm:no-recycled-contact-realization}.  A compatible
inverse limit, unbounded resolved-prefix rank, or raw ordinary path is an
internal state and never a conclusion.
\end{theorem}

\begin{proof}
At depth \(n\), join the first \(n\) successor records before forming any
residual.  The joint grammar adjoins every mixed product, polarization,
constraint descent, carrier row, and total-graph value generated from that
prefix.  For any nonzero finite positive observer law \(\nu\) on the common
measurable same-origin evaluation space, let \(\mathcal A_n\) be the sigma algebra of the depth-\(n\) joint
record and let \(P_n\) be conditional expectation onto
\(L^2(\mathcal A_n,\nu)\).  The observer algebras increase.  Their
conditional expectations of the centered shell contrast \(y_\Phi\) therefore converge strongly, with the exact
orthogonal increment decomposition
\[
 P_\infty y_\Phi=\sum_{n\ge0}d_n,
 \qquad
 d_0=P_0y_\Phi,\quad d_n=(P_n-P_{n-1})y_\Phi,
 \qquad
 \|P_\infty y_\Phi\|_2^2=\sum_{n\ge0}\|d_n\|_2^2.
\]
Consequently a target
component of the graph-complete joint limit is the sum of finite joint
increments; it cannot first appear outside all finite prefixes.

If the full projection vanishes, prospective accountability makes every
source pairing vanish on the retained path and proves \textup{(D1)}.  If a
least unprocessed increment exists, source-preserving totalization supplies
its five-channel ancestry and polarity cell, giving \textup{(D2)}.  Notice
that no lower bound on the increment is used.  If every increment is
processed, apply the joint closure once more.  Failure of fixedness generates
another mixed coordinate and returns to \textup{(D2)}.  At a fixed point the
unnormalized shell cocycle, common origin, complete source ancestry, and all
public graph rows remain present.  Before realization, apply
\cref{thm:recursive-carrier-polar-exhaustion} to every canonical owner in the
fixed enumeration.  A charge, null value, or new polar coordinate returns to
\textup{(D1)} or \textup{(D2)}.  Only simultaneous primal--polar fixedness
reaches the totalized realization row; its nonordinary value is another
proper successor and its ordinary value gives the prospective source atom in
\textup{(D3)}.  Origin-polarity preservation prevents a contact-derived
ordinary value from becoming an independent realization.  Thus rank growth by itself proves nothing and no
target-bearing structure is assigned to a residual limit.
\end{proof}

\begin{definition}[Target-local structural fact]
\label{def:target-local-structural-fact}
Fix a same-origin active branch \(\mathfrak b\), its shell covector
\(\lambda=D\Phi\), and its target-local stratum \(X_{\mathfrak b}\).  A
relation is an \emph{admissible target-local structural fact} when its
complete expanded derivation has first been universally joined and saturated,
and its restriction to \(X_{\mathfrak b}\) has the signed form
\begin{equation}
 \label{eq:target-local-fact-normal-form}
 Q_{\mathfrak b,\lambda}
 =\delta K_{\mathfrak b}
  +\sum_{c\in\Channels}R_{\mathfrak b,\lambda}^{c}
  +P_{\mathfrak b}+B_{\mathfrak b}
  +N_{\mathfrak b}+E_{\mathfrak b}.
\end{equation}
Here \(Q_{\mathfrak b,\lambda}\) is the canonical owner of the selected
target progress in that expanded joint record;
\(R^c_{\mathfrak b,\lambda}\) is the unabsorbed signed remainder with
inherited \(c\)-ancestry;
\(\delta K_{\mathfrak b}\) is a same-origin carrier coboundary;
\(P_{\mathfrak b}\ge0\) is admitted only through a carrier identity dominated
by the public physical measure; \(B_{\mathfrak b}\) is true exterior input;
\(N_{\mathfrak b}\in\Null_\Sigma\); and \(E_{\mathfrak b}\) is the ordered
list of nonordinary or failed graph coordinates.  Compatible internal
currents are added with their signs before any positive part is formed, and
each expanded term occurs exactly once: either in a signed remainder, in the
generated production, in exterior input, in the null term, or in the defect
list.

The fact may close or charge the branch only on
\(X_{\mathfrak b}\).  A nonzero component of \(E_{\mathfrak b}\), or a
nonzero uncharged component of the signed source sum, is passed to the local
polarity successor with the same origin and ancestry.
No identity involving a sum, product, localization, normalization, limit, or
other composite is admissible when applied termwise before this expanded
joint factorization.
\end{definition}

\begin{theorem}[Target-local factorization of continuum facts]
\label{thm:target-local-fact-factorization}
Every relation used by the continuum reachability or regularity calculus is
an admissible target-local structural fact after restriction to the active
branch that invoked it.  This applies to weak-equation, gradient, resolvent,
Hamiltonian, commutator, compactness, constraint, entropy, covariance,
carrier, and continuation relations.  Consequently:
\begin{enumerate}[label=\textup{(F\arabic*)}]
\item a relation on states outside \(X_{\mathfrak b}\) has no effect on that
branch;
\item a target-null term closes no nonnull source cell and contributes no
physical charge;
\item a nonnegative quantity pays for target progress only through
\cref{def:target-aligned-physical-charge}; and
\item failure of localization, descent, ancestry, signed gluing, carrier
domination, or realization is a target-local polarity successor rather than
an omitted antecedent.
\end{enumerate}
\end{theorem}

\begin{proof}
First apply universal saturated-composition accountability to the complete
expanded derivation of the invoked relation.  Pair its joint value with the
selected shell covector and restrict every
input and graph coordinate to \(X_{\mathfrak b}\).  The current normal form
\eqref{eq:raw-current-normal-form} gives the five source terms and the
target-null term.  Structural induction over restriction, quotient, chart,
trace, adjoint, polarization, and limit constructors preserves their source
words by \cref{thm:source-preserving-completion}.  The carrier constructor
separates a coboundary, generated nonnegative production, and true exterior
input; \cref{thm:signed-carrier-gluing} cancels compatible internal faces
before taking positive parts.  Totalization places the first failed operation
in \(E_{\mathfrak b}\).  These are exactly the terms of
\eqref{eq:target-local-fact-normal-form}.  The local charging rule gives
\textup{(F2)}--\textup{(F3)}, while target-local polarity stratification gives
\textup{(F1)} and \textup{(F4)}.
\end{proof}

\begin{corollary}[Activation rule for universal identities]
\label{cor:target-local-identity-activation}
An equation-wide functional-analytic identity is a library relation.  It
enters the regularity calculus only through its factorization
\eqref{eq:target-local-fact-normal-form} on a contact-generated active cell.
Thus an equation-wide operator bound, compactness statement, spectral gap, entropy
law, or coercive inequality has no independent regularity consequence; only
its target-local source components and admissible physical charge do.
\end{corollary}

\begin{corollary}[No unscoped reachability inference]
\label{cor:no-unscoped-reachability-inference}
Every derivation that changes a physical reachability profile factors through
an active branch signature \(\mathfrak b\) and a structural fact of the form
\eqref{eq:target-local-fact-normal-form}.  More precisely, its final
nonidentity step is exactly one of:
\begin{enumerate}[label=\textup{(R\arabic*)}]
\item the target-null quotient annihilates every source remainder on
  \(X_{\mathfrak b}\);
\item the signed local carrier identity charges the selected shell progress
  to the restriction of the public physical measure;
\item a nonordinary coordinate creates a same-origin polarity successor; or
\item after every cumulative source owner has passed recursive
  carrier--polar exhaustion, diagonal saturation reconstructs a
  primal--polar fixed same-origin path with actual target contact.
\end{enumerate}
Consequently no theorem about an unselected state family, no transverse
positive form, and no equation-wide analytic estimate can remove, charge, or
continue a contact branch.
\end{corollary}

\begin{proof}
The ordered reachability calculus consumes a relation only through the local
cone generated by \(I_{\mathfrak b}\), the retained shell floor, and the
factorization \eqref{eq:target-local-fact-normal-form}.  Its zero target
projection is \textup{(R1)}.  Its admitted production and exterior terms are
processed by \cref{thm:local-target-aligned-charging-rule}, giving
\textup{(R2)}.  Every nonzero defect or unabsorbed signed remainder is sent by
\cref{thm:target-local-stratification-exhaustiveness} to \textup{(R3)}.
Every cumulative owner is then processed by
\cref{thm:recursive-carrier-polar-exhaustion}: cone membership returns to
\textup{(R2)}, target-nullity returns to \textup{(R1)}, and separation
creates another instance of \textup{(R3)}.
Compatible cofinal successors are processed by
\cref{thm:diagonal-target-aligned-saturation}; only their primal--polar fixed
ordinary realization gives \textup{(R4)}.
These are the complete outgoing edges of the target-local fact compiler, so
there is no fifth inference rule with equation-wide scope.
\end{proof}

\begin{theorem}[Canonical record-scale selector]
\label{thm:canonical-record-scale-selector}
Suppose a same-origin target-retaining cascade carries scale coordinates
\(r_n>0\).  Exactly one of the following generated alternatives holds:
\begin{enumerate}[label=\textup{(S\arabic*)}]
\item \(\inf_n r_n>0\), and the cascade remains in a fixed-scale stratum;
\item \(\inf_n r_n=0\), and it has a canonical subsequence
\(n_0<n_1<\cdots\) defined by
\begin{equation}
 \label{eq:record-scale-recursion}
 n_{j+1}:=\min\{n>n_j:r_n\le\tfrac12 r_{n_j}\}.
\end{equation}
For the second alternative, interpolate \(\log r_{n_j}\) affinely along the
ordered normalized windows.  The resulting absolutely continuous chart scale
satisfies
\begin{equation}
 \label{eq:record-scale-sign}
 a:=\frac{\dot r}{r}\le0
 \quad\text{a.e.},
 \qquad r(\tau)\downarrow0.
\end{equation}
The selected centers, charts, and transitions retain the origin tag
\(\mathbf R\).  Thus \eqref{eq:record-scale-sign} is a consequence of scale
collapse, not an analytic hypothesis on the original orbit.
\end{enumerate}
\end{theorem}

\begin{proof}
If the first alternative fails, choose \(n_0\).  Since the tail infimum is
zero, the set in \eqref{eq:record-scale-recursion} is nonempty; its least
element defines \(n_{j+1}\).  The selected scales decrease by at least a
factor two and therefore tend to zero.  Affine interpolation of their
logarithms is nonincreasing, which gives
\eqref{eq:record-scale-sign}.  All coordinates are restrictions and chart
images of the original cascade, so same-origin ancestry is preserved.
\end{proof}

\begin{definition}[Source-adapted dual-carrier grammar]
\label{def:source-adapted-dual-carrier-grammar}
Fix \(\mathfrak b=(A,m,\Gamma,\eta,\mathbf R)\).  For
\(\Phi\in\Gamma\), begin with its descending covector \(D\Phi\) and generate
the smallest countable rational test module \(\mathscr W_{\mathfrak b}\)
closed under:
\begin{enumerate}[label=\textup{(D\arabic*)}]
\item adjoints of the source operators occurring in \(A\), including killed
generator and predictable-compensator pullbacks when present;
\item constraint projection, gauge quotient, chart pullback, record-scale
restriction, and multiplication by the generated rational cutoff algebra;
\item signed parent--child gluing, temporal summation by parts, carrier
coboundaries, and cancellation of target-null directions;
\item polarization and completion of squares using the native positive
\(\Geom\) form and every nonnegative defect generated by the weak equation.
\end{enumerate}
For a word \(w\in\mathscr W_{\mathfrak b}\), expansion against the weak
equation produces a \emph{dual-carrier row}
\begin{equation}
 \label{eq:source-adapted-dual-carrier-row}
 Q_{A,\Phi}
 =\delta K_w+P_w+B_w+N_w+E_w.
\end{equation}
Here \(Q_{A,\Phi}\) is the selected source pairing, \(\delta K_w\) is a
bounded carrier coboundary, \(P_w\ge0\) is generated physical production,
\(B_w\) is true exterior input, \(N_w\in\Null_\Sigma\), and \(E_w\) is the
ordered list of total-graph defects.  Paired internal currents are added with
their signs before \(P_w\) is formed.  A production row is admitted only when
the public carrier identity proves that its total mass belongs to the single
physical ledger.
\end{definition}

\begin{remark}[Constructive enumeration]
The grammar is generated from the public equation and the selected shell
covector.  At depth \(k\), enumerate weak multipliers, adjoint pulls, cutoffs,
partitions, chart words, and square completions of syntactic length at most
\(k\), then reduce them modulo \(I_{\mathfrak b}\).  This is the continuum
counterpart of saturating all prospective selectors in Part~I.  A
specialization supplies the equation and its textbook weak identities; it
does not supply a dual carrier or an equality theorem.
\end{remark}

\begin{theorem}[Target-aligned carrier duality]
\label{thm:target-aligned-carrier-duality}
For every active zero-price successor branch \(\mathfrak b\), the generated
dual-carrier hierarchy has exactly one of the following outcomes.
\begin{enumerate}[label=\textup{(C\arabic*)}]
\item A finite ordered identity shows that the active branch is target-null.
\item At a finite depth it derives a row
\eqref{eq:source-adapted-dual-carrier-row} whose graph defects are zero on the
branch and whose exterior positive variation is in the physical ledger.
Repeated target progress is then charged only to \(P_w\), up to a bounded
telescoping carrier and finite exterior input.
\item The ordered dual cone has a positive normalized functional \(L\) with
\begin{equation}
 \label{eq:dual-carrier-separating-law}
 L(I_{\mathfrak b})=0,\qquad
 L(P_w)=0\ \ (w\in\mathscr W_{\mathfrak b}),
 \qquad L(Q_{A,\Phi})\ge\eta.
\end{equation}
It is used only to select a finite dual coefficient word for the next carrier
calculation attached to the already specified owner
\(Q_{A,\Phi}\), with its mode, target, chart, origin, and complete parent
coordinates.  The functional is not inserted into the source-evaluation
domain and supplies neither a source occurrence nor an active branch.  Source
presence requires a finite same-origin pre-contact occurrence from the
presentation-generated reachable domain.
\end{enumerate}
No conclusion requires control of a source or equality state outside the
target-active branch \(\mathfrak b\).
\end{theorem}

\begin{proof}
Form the Archimedean ordered cylinder algebra generated by
\(I_{\mathfrak b}\), the nonnegative productions \(P_w\), the exterior
ledger, and the inequality \(Q_{A,\Phi}\ge\eta\).  If its quadratic module
contains \(-1\) at a finite level, it gives \textup{(C1)}.  Otherwise dovetail
the dual-carrier words.  A word whose exact telescoping identity charges the
retained crossings gives \textup{(C2)}.  If neither finite event occurs, the
ordered Hahn--Banach theorem applied to the cofinal cylinder cones gives a
positive normalized functional annihilating the branch ideal.  Since this is
a zero-price successor branch and every admitted \(P_w\) is dominated by the
physical ledger, it also annihilates every \(P_w\), while the defining shell
inequality retains the target floor.  This is \textup{(C3)}.  The separating
functional is used to extract the next finite dual coefficient word.  This
word acts only on the source owner already present in the branch and cannot
establish that an occurrence realizing that owner exists.  Weak duality
makes the alternatives disjoint after selecting the first finite closing row;
the ordered proof/model theorem makes them exhaustive.
\end{proof}

\begin{definition}[Target-local carrier cone and polar successor]
\label{def:target-local-carrier-polar-successor}
For an active signature \(\mathfrak b\), let
\(\mathscr Y_{\mathfrak b}^{\Sigma}\) be the closed locally convex space
generated by the source-owned target-current cylinder coordinates on
\(X_{\mathfrak b}\), equipped with the initial topology of the countable
family of bounded rational cylinder seminorms.  Its target-effective space is
\(\mathscr Y_{\mathfrak b}^{\Sigma}/\overline{\Null_\Sigma}\).  Let
\(\mathscr D_{\mathfrak b}^{\rm alg}\) be the algebraic convex cone generated
by finite rows of
\begin{equation}
 \label{eq:target-local-carrier-cone}
 \delta\mathscr K_{\mathfrak b}
 +\mathscr P_{\mathfrak b}^{\Phys}
 +\mathscr B_{\mathfrak b}^{\Phys}
 +\Null_\Sigma
\end{equation}
and put
\(\mathscr D_{\mathfrak b}^{\Phys}
 :=\overline{\mathscr D_{\mathfrak b}^{\rm alg}}\)
in this topology.
Here \(\delta\mathscr K_{\mathfrak b}\) is included with both signs,
\(\mathscr P_{\mathfrak b}^{\Phys}\) contains only nonnegative production
already dominated by the single public physical measure, and
\(\mathscr B_{\mathfrak b}^{\Phys}\) contains only true exterior input
already present in that ledger.  Membership in
\(\mathscr D_{\mathfrak b}^{\rm alg}\) is a finite generated target-local
charging identity.  Membership only in its closure is a limiting domination
value, not a finite proof: the total cone-closure graph retains the
approximating rows and their cylinder errors, and closes an owner only after
those errors vanish on every target shell used downstream.

If \(Q\notin\mathscr D_{\mathfrak b}^{\Phys}\), a
\emph{carrier-polar successor} is the first finite-cylinder separator in the
fixed dual enumeration whose coefficient tuple is retained through its total
graph, or the corresponding graph-boundary value,
satisfying
\begin{equation}
 \label{eq:carrier-polar-separator}
 \lambda(Q)>0,
 \qquad
 \lambda\big|_{\delta\mathscr K_{\mathfrak b}+\Null_\Sigma}=0,
 \qquad
 \lambda(P+B)\le0
 \quad
 (P+B\in\mathscr P_{\mathfrak b}^{\Phys}
              +\mathscr B_{\mathfrak b}^{\Phys}).
\end{equation}
Adjoin \(\lambda\), every adjoint pullback of \(\lambda\) through the public
source words, every polarization and square it generates, and all their
total-graph values as dual coefficient data.  Retain the canonical
source owner, unnormalized shell cocycle, and same-origin tag only when they
come from a finite presentation-generated occurrence; do not infer them from
\(\lambda\).  Denote this dual operation by
\(\mathfrak P_\Sigma(\mathfrak b,Q)\).
\end{definition}

\begin{lemma}[Finite-cylinder detection of strict carrier separation]
\label{lem:rational-carrier-separation}
If \(Q\notin\mathscr D_{\mathfrak b}^{\Phys}\), then there are a finite
cylinder projection \(\pi_F\), a real coefficient vector \(q_F\), and a
margin \(\varepsilon>0\) such that the cylinder functional
\(\lambda_F=q_F\!\cdot\!\pi_F\) satisfies
\[
 \lambda_F(Q)\ge\varepsilon,
 \qquad
 \sup_{D\in\mathscr D_{\mathfrak b}^{\Phys}}\lambda_F(D)\le0.
\]
The coefficient tuple and the verification that it is nonpositive on the
projected cone are retained through the total dual graph.  Rational
enumeration is complete only on finite rational-polyhedral projections; no
rationalization or boundary-membership claim is made in general.
\end{lemma}

\begin{proof}
The quotient by the closed target-null space is locally convex and Hausdorff.
Strict separation of the point \(Q\) from the closed convex cone gives a
continuous functional and a positive margin.  Continuity for the initial
cylinder topology makes that functional factor through a finite projection
\(\pi_F\).  This proves the displayed statement.  If the projected cone is
rational polyhedral, ordinary finite-dimensional linear programming gives a
rational separator after reducing the margin.  Without that extra property,
replacing the real coefficient vector by a rational vector need not preserve
membership in the polar cone, which is why the total dual graph is retained.
\end{proof}

\begin{definition}[Finite-cylinder carrier comparison problem]
\label{def:finite-cylinder-physical-source-profile}
Fix an active same-origin branch \(\mathfrak b\), a finite generated cylinder
\(F\), and a source-owned target current \(Q\).  Enlarge \(F\), when
necessary, so that it contains \(Q\), the order unit \(1\), the physical
ledger coordinate \(p_{\Phys}\), and every coordinate occurring in the
finite relations imposed at this stage.  Let \(A_F\) be the resulting
finite-dimensional real order-unit space modulo the projected equality
ideal and target-null subspace.  Write \(C_F\subseteq A_F\) for the closed
cone generated by the projected squares, admitted physical productions,
true exterior inputs, and signed carrier coboundaries.

For \(b\ge0\), define the normalized finite-cylinder dual test set and its
carrier upper bound by
\begin{align}
 \mathcal S_F(b)
 &:={\bigl\{L\in A_F^*:L(1)=1,\ L(C_F)\ge0,
                     \ 0\le L(p_{\Phys})\le b\bigr\}},
 \label{eq:finite-cylinder-state-set}\\
 \mathsf U_{F,Q}(b)
 &:=\sup_{L\in\mathcal S_F(b)}L(Q),
 \label{eq:finite-cylinder-source-profile}
\end{align}
with supremum \(-\infty\) when \(\mathcal S_F(b)=\varnothing\).
The cylinder is called \emph{checked} when every relation used to form
\(A_F\) and \(C_F\) carries its finite derivation from the public equation.
Thus \(\mathcal S_F(b)\) is a dual test set generated by checked relations,
not a source-evaluation domain.  Its only use is to derive a finite upper
bound or separating inequality for evaluations of actual occurrences.
\end{definition}

\begin{definition}[Finite-cylinder occurrence presence]
\label{def:finite-cylinder-occurrence-presence}
Let \(\mathcal O_F^{\rm pre}(\mathfrak b)\) be the finite same-origin
pre-contact occurrences obtained by restricting ordinary paths in the
presentation-generated reachable domain of \(\mathfrak b\) to the cylinder
\(F\).  Each \(R\in\mathcal O_F^{\rm pre}(\mathfrak b)\) defines its actual
evaluation functional \(\operatorname{ev}_R\in\mathcal S_F(b_R)\).  A source
owner \(Q\) is \emph{present above \(\eta\) on \(F\)} precisely when
\begin{equation}
 \label{eq:finite-cylinder-occurrence-presence}
 \exists R\in\mathcal O_F^{\rm pre}(\mathfrak b)
 \quad \operatorname{ev}_R(Q)\ge\eta .
\end{equation}
The occurrence, not an arbitrary member of \(\mathcal S_F(b)\), carries the
chronology, public initial datum, source ancestry, and pre-contact origin
needed by prospective accountability.
For a physical ceiling \(b\), the finite-prefix source profile is
\begin{equation}
 \label{eq:finite-cylinder-actual-source-profile}
 \mathsf T_{F,Q}^{\rm occ}(b)
 :=\sup\{\operatorname{ev}_R(Q):
 R\in\mathcal O_F^{\rm pre}(\mathfrak b),\ b_R\le b\}.
\end{equation}
This is the only source profile associated with \(F\).
\end{definition}

\begin{theorem}[Exact finite-cylinder carrier-bound duality]
\label{thm:exact-finite-cylinder-source-profile-duality}
Let \(F\) be a checked cylinder.  Then \(\mathcal S_F(b)\) is compact and,
whenever it is nonempty,
\begin{equation}
 \label{eq:finite-cylinder-profile-dual}
 \mathsf U_{F,Q}(b)
 =\inf_{\alpha\ge0,\,\beta\in\mathbb R}
 \left\{\alpha b+\beta:
 \beta 1+\alpha p_{\Phys}-Q\in C_F\right\}.
\end{equation}
If the infimum is attained only in \(\overline{C_F^{\rm alg}}=C_F\), the
approximating cone rows and their errors are retained as a cone-boundary
successor.  It is a finite charging proof precisely when the displayed
element belongs to \(C_F^{\rm alg}\).

For every threshold \(\eta\), exactly one of the following statements about
the carrier upper bound holds:
\begin{enumerate}[label=\textup{(F\arabic*)}]
\item a checked dual row gives \(\mathsf U_{F,Q}(b)<\eta\);
\item \(\mathsf U_{F,Q}(b)=\eta\), and the exposed equality face
\begin{equation}
 \label{eq:finite-cylinder-equality-face}
 \mathcal E_{F,Q}(b,\eta)
 :=\{L\in\mathcal S_F(b):L(Q)=\eta\}
\end{equation}
is used only to identify the vanishing terms of an attaining or approximating
dual row; those finite algebraic terms, rather than a member of the face, are
adjoined to the carrier calculation;
\item \(\mathsf U_{F,Q}(b)>\eta\); this numerical value supplies no source
presence statement.
\end{enumerate}
In \textup{(F2)} the equality face is processed by generated squares,
polarizations, adjoints, symmetries, and the dyadic ordered-coercivity rule.
Neither \textup{(F2)} nor \textup{(F3)} establishes source presence.  The
actual occurrence profile satisfies only the rigorously licensed inequality
\begin{equation}
 \label{eq:actual-profile-below-carrier-bound}
 \mathsf T_{F,Q}^{\rm occ}(b)\le\mathsf U_{F,Q}(b).
\end{equation}
Presence is exactly \eqref{eq:finite-cylinder-occurrence-presence}.
\end{theorem}

\begin{proof}
The order-unit normalization bounds every coordinate of \(A_F^*\); positivity
and the ledger inequality define a closed subset.  Finite dimensionality
therefore makes \(\mathcal S_F(b)\) compact.  Its support value at \(Q\) is
the primal value in \eqref{eq:finite-cylinder-source-profile}.  The polar of
the intersection of the normalized positive dual cone with
\(L(p_{\Phys})\le b\) consists exactly of the affine majorants
\(Q\le\alpha p_{\Phys}+\beta1\), with \(\alpha\ge0\).  Finite-dimensional
separation, applied to the epigraph of the support function, gives
\eqref{eq:finite-cylinder-profile-dual}; this is the standard bipolar proof
and requires neither semialgebraicity nor rational coefficients.

Compactness gives a maximizer of the dual carrier-bound problem; the
maximizer is not retained as a source-evaluation value.  Comparison of the
maximum with \(\eta\) gives the three disjoint cases.  Every actual evaluation
\(\operatorname{ev}_R\) is an admissible dual test, which proves
\eqref{eq:actual-profile-below-carrier-bound}.  In the equality case, positivity implies
that every nonnegative summand on which an attaining dual row has zero total
value vanishes separately.  For a nonattained boundary row the same
statement is expressed by its approximating rows and errors, and
no finite proof is asserted.  The listed equality constructors preserve the
same finite cylinder after adjoining their finitely many output coordinates;
\cref{thm:dyadic-ordered-coercivity-terminal-rule} applies whenever its
generated row is present.
\end{proof}

\begin{theorem}[Monotone cylinder carrier bounds]
\label{thm:monotone-cylinder-source-profile-completion}
Let \(F_1\subseteq F_2\subseteq\cdots\) be the canonical enumeration of the
checked target-relative cylinders containing \(Q,1,p_{\Phys}\), and let
restriction map \(\mathcal S_{F_{n+1}}(b)\) into
\(\mathcal S_{F_n}(b)\).  Then
\begin{equation}
 \label{eq:monotone-cylinder-profile-limit}
 \mathsf U_{F_1,Q}(b)\ge\mathsf U_{F_2,Q}(b)\ge\cdots,
 \qquad
 \mathsf U_{\infty,Q}(b):=\inf_n\mathsf U_{F_n,Q}(b).
\end{equation}
For every actual finite same-origin occurrence whose restriction is defined
on \(F_n\),
\begin{equation}
 \label{eq:occurrence-below-every-cylinder-bound}
 \operatorname{ev}_R(Q)\le\mathsf U_{F_n,Q}(b_R).
\end{equation}
Hence \(\mathsf U_{F_n,Q}(b)<\eta\) at one finite depth excludes every
actual occurrence of physical cost at most \(b\) and target value at least
\(\eta\).  If no such finite depth is found, the carrier-bound calculation
has produced no source-presence or terminal statement.  Source presence and
cofinal accumulation remain entirely in the actual occurrence domain of
\cref{def:finite-cylinder-occurrence-presence} and the prospective backward
frontier.
\end{theorem}

\begin{proof}
Adding checked relations can only shrink the dual test set, proving
monotonicity of its support bound.  The actual evaluation functional of an
occurrence belongs to every finite dual test set on which that occurrence is
defined; \eqref{eq:occurrence-below-every-cylinder-bound} follows.  The
finite exclusion statement is its contrapositive.  No inverse limit of dual
test functionals is formed.
\end{proof}

\begin{corollary}[Part-I form of continuum envelope evaluation]
\label{cor:part-I-continuum-profile-evaluation}
Let a stopped first-contact readout have target floor \(\eta>0\).  Apply the
backward frontier before terminal passage, classify every conditional joint
increment by its complete five-source ancestry, and evaluate each resulting
owner through
\cref{thm:exact-finite-cylinder-source-profile-duality,thm:monotone-cylinder-source-profile-completion}.
Then a checked carrier ceiling below \(\eta\) excludes that contact.  Source
presence, when contact is the premise, is supplied only by its finite
same-origin pre-contact occurrence and
\cref{def:finite-cylinder-occurrence-presence}.  The dual calculation neither
adds nor replaces an occurrence.
Fixedness, cone-boundary membership, an irrational separator, failure of
semialgebraic elimination, or survival through every finite cylinder creates
no additional source or terminal value.  This is the continuum counterpart
of bounding the saturated capability profile in Part~I while keeping actual
executions separate from dual bounds.
\end{corollary}

\begin{theorem}[Universal recursive amplification exhaustion]
\label{thm:recursive-carrier-polar-exhaustion}
Let \(Q\) be the canonical source owner of a nonzero target increment created
by any vertex of a universal expanded derivation on an active same-origin
branch \(\mathfrak b\).  This includes primitive, cross-term, localized,
stopped, normalized, scale-transferred, recurrent, limiting, tail, and
boundary owners.  Exactly one of the following compiler transitions occurs.
\begin{enumerate}[label=\textup{(P\arabic*)}]
\item The target projection of \(Q\) vanishes after the generated quotient,
so the owner cannot contribute to shell contact.
\item On the complementary quotient cell,
\(Q\in\mathscr D_{\mathfrak b}^{\rm alg}\).  Its finite cone identity telescopes
storage, removes target-null terms, and charges the remaining positive mass
to the public physical ledger.
\item The owner belongs to
\(\mathscr D_{\mathfrak b}^{\Phys}\setminus
\mathscr D_{\mathfrak b}^{\rm alg}\).  Its cone-closure approximation rows,
target-shell errors, and inherited owner form the next typed cone-boundary
successor.  It is not declared charged until the limiting errors vanish on
the retained shell cocycle.
\item The carrier-polar finite-prefix calculation
\(\mathfrak P_\Sigma(\mathfrak b,Q)\) is nonzero and is processed before
circulation, recurrence, compact realization, or another shell is evaluated.
\end{enumerate}
A finite value of the universal amplification graph is merely another
coordinate of the cone test; it closes the owner only when it supplies the
displayed membership identity.  An unbounded or zero-denominator value is
processed by \textup{(P3)} or \textup{(P4)}, according to its cone-closure
value, with the numerator, denominator, source word, and
unnormalized cocycle retained.  The compiler applies this transition to every
owner in the fixed enumeration of the expanded joint record before any
postprocessing or realization row.

A cofinal polar chain is joined to, but kept type-distinct from, its primal
owner chain.  Its joint comparison limit is charged, proves the primal owner
target-null, exposes a finite unprocessed primal occurrence, or is fixed under
the comparison grammar.  Fixedness supplies no source atom.  On a
contact-seeded execution, prospective source presence comes solely from the
retained finite same-origin occurrences; an independently constructed
ordinary public orbit is still required for a singular-existence verdict.
\end{theorem}

\begin{proof}
The first cell is removed before the cone test.  Algebraic membership,
closure-only membership, and nonmembership are disjoint total-graph values.
The first gives the finite charging row in \textup{(P2)}; the second gives the
retained convergence graph in \textup{(P3)}.  On their complement, geometric
Hahn--Banach separation of \(Q\) from the closed convex cone
\(\mathscr D_{\mathfrak b}^{\Phys}\), followed by
\cref{lem:rational-carrier-separation}, gives
\(\lambda\) satisfying \eqref{eq:carrier-polar-separator}; the coboundary and
null summands are linear spaces and are therefore annihilated.  The topology
of \(\mathscr Y_{\mathfrak b}^{\Sigma}\) is generated by its finite-cylinder
seminorms, so a continuous strict separator is detected at finite cylinder
depth.  Its coefficient tuple is a value of the totalized dual graph: an
ordinary value gives the separator and a nonordinary value is its typed dual
graph boundary.  Totalization retains the separator's association with the
owner; it does not turn the separator into a source occurrence.

Apply the adjoint and polarization clauses of
\cref{def:source-adapted-dual-carrier-grammar}.  A previously absent
separator or descendant is a finite proper comparison refinement.  For a
cofinal chain, retain the cumulative primal occurrence record and its dual
comparison record as distinct typed coordinates.  Apply
\cref{thm:continuum-composition-no-creation} only to the primal occurrence
coordinates: every nonzero limiting target component then occurs on a finite
same-origin prefix.  Dual fixedness supplies only a finite-prefix bound.
This proves exhaustiveness of the four cone values and prevents a positive
uncharged owner from being passed directly to realization.
Because the proof used only algebraic membership, closure, and separation for an
arbitrary generated owner, it applies verbatim to every constructor named in
the theorem.  Totalization of its amplification graph prevents division by a
vanishing physical quantity or normalization factor from deleting the
numerator.  Dovetailing the countable owner enumeration before realization
proves the universal assertion.
\end{proof}

\begin{corollary}[Carrier-admission completeness]
\label{cor:carrier-admission-completeness}
Every nonzero target owner has exactly one physical disposition: it is
target-null; it has a finite algebraic domination identity by the single
public carrier; it has only a cone-closure domination and therefore produces
its typed approximation-error successor; or it has strict carrier-cone
nonmembership and therefore produces a finite-cylinder polar separator with
all generated adjoint and polarization descendants.  No positive owner is
stored as an unpriced residual state.
\end{corollary}

\begin{proof}
These are respectively cases \textup{(P1)}--\textup{(P4)} of
\cref{thm:recursive-carrier-polar-exhaustion}.  Their disjointness is the
disjointness of target-nullity, algebraic cone membership, closure-only cone
membership, and strict nonmembership in the total cone-closure graph.
\end{proof}

\begin{definition}[Exact-support transfer ledger]
\label{def:exact-support-transfer-ledger}
Let a positive exact-support amplitude token \(\alpha_A(R)\) enter a
single finite compiler transition.  Its checked routing table partitions the
token among certificate-closed output, target-null output, immediate
same-support children, and immediate strict-support children.  Write the
nonnull amounts as
\[
 \alpha_{A,\rm closed}(R),\qquad
 \alpha_{A,C}^{\rm stay}(R)
   \quad(C\in\operatorname{Child}_A(R)),\qquad
 \alpha_{A'}^{\rm child}(R)
 \quad(A'\supsetneq A).
\]
The routing table is the target projection of the corresponding public typed
constructor graph after its inputs have been joined, and every child support
is computed by the public provenance morphism.  It includes a nonnegative
amount for every output slot and a finite symbolic identity asserting that
their sum is the input token.  It is
formed after the signed parent current has been grouped and its positive part
taken, and is retained alongside, rather than substituted for, the signed
target-current pairing of each child.  Target-null outputs receive the zero
quotient token.  A live same-support token is never silently counted as
closed.
\end{definition}

\begin{theorem}[Support-transfer conservation]
\label{thm:support-transfer-conservation}
Every successor of an exact support \(A\) either retains \(A\), or a quotient,
lift, or boundary constructor appends its declared letter and routes its
amplitude token to an exact strict superset \(A'\supsetneq A\).  For every
transfer ledger,
\begin{equation}
 \label{eq:support-transfer-amplitude-conservation}
 \alpha_A(R)
 =\alpha_{A,\rm closed}(R)
  +\sum_{C\in\operatorname{Child}_A(R)}
       \alpha_{A,C}^{\rm stay}(R)
  +\sum_{A'\supsetneq A}\alpha_{A'}^{\rm child}(R).
\end{equation}
Every nonnull child carries the parent source route and the target covector
that generated the token.  Along a branch, strict support growth occurs at
most \(5-|A|\) times.
\end{theorem}

\begin{proof}
The provenance morphism preserves support under same-word constructors and
appends only \(\Abs,\Lift,\Bdry\) at the corresponding interfaces.  Totalize
the interface before sourcewise closure and check the finite routing-table
sum.  Its same-support, strict-growth, closed, and null output slots are
disjoint and exhaustive by their tags; the null slot has zero quotient value.
This proves \eqref{eq:support-transfer-amplitude-conservation}.  The
underlying signed current identities are retained separately, so the token
identity does not assert additivity of positive parts across a signed current
decomposition.  Source-preserving closure copies the tagged token and its
covector to each live child.  Each strict transfer adjoins at least one new
letter from a five-letter alphabet, proving the final bound.
\end{proof}

\begin{definition}[Canonical physical-charge index]
\label{def:canonical-physical-charge-index}
Every finite shell occurrence carries its original pairwise-disjoint event
index \(j\) from the canonical same-origin cascade of
\cref{thm:canonical-disjoint-shell-cascade}.  Each structural descendant
copies \(j\), the public carrier coordinate, shell covector, and origin.  When the carrier
admission calculation first produces an ordinary chargeable output for this
index, the fixed compiler order marks that output as
\(\operatorname{ch}(j)\).  Every other descendant with index \(j\) carries
zero physical charge and only its structural amplitude token.  A later
return, resaturation, or cumulative join may refine the metadata of
\(\operatorname{ch}(j)\), but may not create a second marked output.
\end{definition}

\begin{theorem}[No double charging]
\label{thm:no-double-charging}
For every finite compiler subtree generated from a pairwise disjoint
same-origin shell family \((Q_j)\),
\[
 \sum_{\text{\rm charged leaves }L}p_{\Phys}(L)
 \le\mu_u\!\left(\bigsqcup_jQ_j\right)
 \le\Budget(u_0).
\]
The same assertion holds for every finite prefix of a cofinal compiler
execution.  In particular, no divergent price can be created by repeatedly
charging distinct structural descriptions of one physical occurrence.
\end{theorem}

\begin{proof}
By \cref{def:canonical-physical-charge-index}, each chargeable leaf has a
unique original index and the map from charged leaves to indices is
injective.  The charge of that leaf is the restriction of the single public
physical measure to its original event, or a disjoint subatom of that
restriction certified by the carrier row.  Summing over an injective finite
index set is bounded by countable additivity of \(\mu_u\) on the disjoint
family \((Q_j)\).  Every return or resaturation descendant retains its index,
so induction on the finite subtree proves the claim.  Applying the same
argument to increasing finite prefixes gives the cofinal assertion.
\end{proof}

\begin{definition}[Unresolved target-amplitude ledger]
\label{def:unresolved-target-amplitude-ledger}
An amplitude token entering a compiler transition is partitioned into a
certificate-closed charged token, a target-null token, and the tokens of its
ordered children.  The charged token records the target amplitude discharged
by a finite carrier, continuation, or incompatibility certificate; it is not
identified with the physical price.  The null token has value zero after
target projection.  Support-transfer and re-entry children retain their
unnormalised amplitude token, source route, origin, and shell covector.
\end{definition}

\begin{theorem}[Conservation of unresolved target progress]
\label{thm:conservation-unresolved-target-progress}
Every finite compiler transition satisfies the exact token invariant
\begin{equation}
 \label{eq:conservation-unresolved-target-progress}
 \mathsf A_{\rm in}
 =\mathsf A_{\rm charged}+\mathsf A_{\rm null}
  +\sum_{\rm children}\mathsf A_{\rm child},
 \qquad \mathsf A_{\rm null}=0\ \text{in }\mathscr J_\Sigma.
\end{equation}
At a cofinal transition the same identity holds for the cumulative joint
record after sourcewise joining and target reprojection.
\end{theorem}

\begin{proof}
For a primitive exact-support current the identity is the partition of the
positive amplitude token after its signed current has been grouped.  A
carrier certificate closes the corresponding token; target-nullity closes it
with zero quotient value; and every other disposition is a child token.
Support transfer preserves the identity by
\cref{thm:support-transfer-conservation}.  Storage, equality, and graph
re-entry preserve it because the parent token is passed to the first
contextual successor, with any certificate-closed portion removed once.
These cases exhaust the carrier-polar transition by
\cref{thm:recursive-carrier-polar-exhaustion}.

For a cofinal chain, apply the finite identity to every bounded truncation
of the cumulative token.  The joint record retains all child routes before
projection, and the truncations increase monotonically.  Monotone passage to
the cumulative join gives \eqref{eq:conservation-unresolved-target-progress}.
If a purported residual had nonzero quotient value, finite no-emergence
would expose another child; otherwise it is the null term.  Thus reprojection
does not lose unresolved target progress.
\end{proof}

\begin{corollary}[No pre-polar realization]
\label{cor:no-pre-polar-realization}
A polar functional, equality functional, or compatible dual inverse limit
never enters an ordinary-orbit realization row.  It is a finite-prefix bound
record.  Only a presentation-generated same-origin primal approximation
branch may enter realization, after each of its canonical source owners has
been processed by \cref{thm:recursive-carrier-polar-exhaustion}.  Comparison
fixedness neither supplies nor manufactures a source occurrence.
\end{corollary}

\begin{corollary}[No diffuse comparison-to-presence loophole]
\label{cor:no-diffuse-primal-polar-composition}
Let \((\mathcal J_n^{\rm pr})\) be the cumulative same-origin primal
occurrence records of a cofinal compiler execution, and let
\((\Lambda_n)\) be the associated finite-prefix polar calculations.  The target map
acts on \(\mathcal J_n^{\rm pr}\), not on \(\Lambda_n\).  If the primal limit
has nonzero target projection, then some finite same-origin primal prefix has
nonzero target projection.
Equivalently,
\begin{equation}
 \label{eq:no-diffuse-primal-polar-composition}
 q_\Sigma(\mathcal J_n^{\rm pr})=0\quad\text{for every }n
 \quad\Longrightarrow\quad
 q_\Sigma(\mathcal J_\infty^{\rm pr})=0.
\end{equation}
Hence no countable family of individually target-null or already canceled
primal occurrences can create target-aligned structure only at the limit,
and no collection of polar comparisons can create source presence at all.
If the cumulative primal target contribution is nonzero, its canonical owner
is exposed on a finite same-origin prefix and is processed again by
\cref{thm:recursive-carrier-polar-exhaustion}.  Neither division by the
number of returns nor passage to an average is permitted before this test.
\end{corollary}

\begin{proof}
Carrier--polar separators are coefficient records in the dual cylinder
algebra.  They may test a primal current but are outside the domain of the
target projection.  Apply \cref{thm:continuum-composition-no-creation} to the
cumulative primal occurrence algebra.  Its exact orthogonal finite-prefix
decomposition makes the target projection of the limit the sum of its actual
finite occurrence increments.  Thus a nonzero projection has a nonzero
finite-prefix increment.  The occurrence owner and its complete ancestry are
retained by the cumulative primal join, so recursive carrier--polar
exhaustion applies to that owner.  The final
sentence follows from the unnormalized cocycle coordinate
\(A_{\Phi,N}=N\), which is included before every recurrent normalization.
\end{proof}

\begin{definition}[Target-local equality stratum]
\label{def:target-local-equality-kernel}
The equality kernel of \(\mathfrak b\) is
\begin{equation}
 \label{eq:target-local-equality-kernel}
 \mathcal E_{\mathfrak b}
 :=\left\{z: I_{\mathfrak b}(z)=0,
                   \ P_w(z)=0\text{ for all }
                   w\in\mathscr W_{\mathfrak b}\right\}.
\end{equation}
It is restricted immediately to the retained shell cell
\(Q_{A,\Phi}\ge\eta\).  Its target-null cell is precisely
\begin{equation}
 \label{eq:target-local-null-test}
 \mathcal E_{\mathfrak b}\cap\{Q_{A,\Phi}\ge\eta\}=\varnothing.
\end{equation}
The finite-level quadratic modules are dovetailed with the generated polarity
stratification.  A derivation of \(-1\) after adjoining
\(Q_{A,\Phi}-\eta\) closes the local cell; compatible positive functionals
are used only to select further finite equality rows and are not retained in
the source-evaluation domain.  A target-retaining
equality successor is present only when a finite same-origin pre-contact
occurrence satisfies the equality rows.  Thus \eqref{eq:target-local-null-test} is an output of
local saturation, not a property to be established separately and not a
classification of all solutions of the equation.
\end{definition}

\begin{proposition}[Covariant modulation conjugation]
\label{prop:covariant-modulation-conjugation}
Suppose a dual-carrier word on an active moving-chart branch gives
\begin{equation}
 \label{eq:abstract-modulated-carrier}
 K'=-D+2aK+R,
 \qquad K,D\ge0,
\end{equation}
where \(R\) is the remaining signed source current.  Let \(\theta>0\) solve
the scalar terminal or initial-value problem
\begin{equation}
 \label{eq:abstract-covariant-scale}
 \theta'+2a\theta=-1.
\end{equation}
Then the generated carrier \(\mathcal G=\theta K\) satisfies the exact
identity
\begin{equation}
 \label{eq:abstract-conjugated-carrier}
 \mathcal G'=-K-\theta D+\theta R.
\end{equation}
No sign or integrability condition on \(a\) is used.  If \(R\) is canceled by
an adjoint word, a signed interface gluing, or a target-null quotient, then
\(K+\theta D\) is a nonnegative production generator in
\(\mathscr W_{\mathfrak b}\).
\end{proposition}

\begin{proof}
Differentiate \(\theta K\) and substitute
\cref{eq:abstract-modulated-carrier,eq:abstract-covariant-scale}:
\[
 (\theta K)'=(-1-2a\theta)K
 +\theta(-D+2aK+R)=-K-\theta D+\theta R.
\]
All operations belong to the scalar adjoint and carrier-coboundary clauses of
\cref{def:source-adapted-dual-carrier-grammar}.
\end{proof}

\begin{theorem}[Branch-local modulation stratification]
\label{thm:branch-local-modulation-reduction}
Let \(\mathfrak b\) be a record-scale active branch, so that \(a\le0\) by
\cref{thm:canonical-record-scale-selector}.  Suppose its generated
stationary equations contain
\begin{equation}
 \label{eq:branch-local-sign-identity}
 \int(P-aO)\,\dd\varpi=0,
 \qquad P,O\ge0,
\end{equation}
Then every retained stationary law is supported on
\begin{equation}
 \label{eq:branch-zero-production-stratum}
 \mathcal E_{\mathfrak b}^{P=0}
 :=\mathcal E_{\mathfrak b}\cap\{P=0\}.
\end{equation}
Restrict the local ordered cone to
\(\mathcal E_{\mathfrak b}^{P=0}\).  A finite \(-1\)-identity closes its
retained target cell; otherwise its compatible positive moments select the
next finite equality tests.  If a finite same-origin pre-contact occurrence
is present on this stratum, the shell covector, five source currents, and
every generated operation of that occurrence are restricted to
\eqref{eq:branch-zero-production-stratum}, its next polarity coordinate is
recorded, and prospective accountability is rerun there.  The positive
moments alone do not create this successor.  No property or
semantic emptiness decision for the zero set is supplied as a premise.
\end{theorem}

\begin{proof}
Both terms in \(P-aO=P+(-a)O\) are nonnegative.  Equation
\eqref{eq:branch-local-sign-identity} therefore gives \(P=0\)
\(\varpi\)-almost surely, proving
\eqref{eq:branch-zero-production-stratum}.  Apply the local Archimedean
proof/model alternative after adjoining \(P=0\) and the retained shell floor.
A finite contradiction closes the carrier test; otherwise the positive law
is used only to select another finite coefficient row.  Source presence and reselection use
only an occurrence in \(\mathcal O_F^{\rm pre}(\mathfrak b)\) satisfying the
same equality rows.  Restriction, polarity refinement, and prospective
reselection of that occurrence are generated operations in
\cref{def:generated-polarity-stratification,def:source-adapted-dual-carrier-grammar}.
\end{proof}

\begin{theorem}[Exhaustive loss-of-compactness routing]
\label{thm:target-aligned-compactness-routing}
Let a sequence in an active branch fail to converge in an ordinary topology
used by a dual-carrier row.  Its image in the generated evaluation
compactification has a subsequence with exactly one primary successor mode:
\begin{equation}
 \label{eq:compactness-successor-modes}
 \mathsf{atom},\ \mathsf{separated},\ \mathsf{nested},\
 \mathsf{diffuse},\ \mathsf{tail},\ \mathsf{rough},\
 \mathsf{gauge},\ \mathsf{constraint},\ \mathsf{equality}.
\end{equation}
The compactified value selects the next finite cylinder test; it carries no
source-presence or target-floor semantics.  The target floor remains on the
finite same-origin occurrences generating the sequence.  Atom, separated,
nested, and diffuse values are disintegrations
of oscillation or concentration and append \(\Lift\); tail values append
\(\Bdry\); gauge values append \(\Abs\); rough equation and constraint values
retain the ancestry of the failed operation.  If a scale coordinate tends to
zero, \cref{thm:canonical-record-scale-selector} supplies the nonincreasing
record chart before the next test is evaluated.  If the mode's conditional
increment vanishes on every generating occurrence, that occurrence cell
closes in \(\Null_\Sigma\).  Otherwise the first finite same-origin
occurrence detecting the mode is processed by
\cref{thm:target-aligned-carrier-duality} with the same owner.
\end{theorem}

\begin{proof}
The normalized activity measures are probabilities on the compact observer
space.  Sequential compactness and disintegration give the mutually ordered
readouts in \eqref{eq:compactness-successor-modes}; the first nonordinary
total-graph coordinate fixes the primary mode.  Completion identifies the
next coordinate test but creates no occurrence.  The prospective
finite-prefix theorem preserves the target floor on the generating
occurrence sequence.  Applying the new test there gives either zero actual
conditional increments or a first nonzero finite increment.  The latter
retains the failed owner's ancestry and re-enters the dual-carrier compiler.
\end{proof}

\begin{theorem}[Continuous selector-accountability closure]
\label{thm:continuous-selector-accountability-closure}
Let an ordinary finite-physical-budget orbit contact \(\Sigma\).  Then there
is an active branch signature \(\mathfrak b\) with
\(\eta\ge1/5\).  Along that branch, repeated contact has exactly one of the
following consequences:
\begin{enumerate}[label=\textup{(A\arabic*)}]
\item a generated dual carrier charges infinitely many disjoint crossings to
the single finite physical ledger, a contradiction;
\item a generated target-null, regular, continuable, or class-incompatible
successor contradicts retained contact;
\item target-local polarity stratification produces a proper successor with
strictly larger resolved-prefix rank;
\item the retained finite same-origin occurrences continue cofinally after
the carrier coefficient enumeration produces no further finite closing row.
The occurrences, through their unnormalized shell cocycle, witness the lower
bound in the saturated prospective source cell.
\end{enumerate}
In the fourth alternative, prospective source presence is the cofinal family
of finite pre-contact occurrences.  A singular-existence verdict still
requires an independent-origin ordinary orbit.  A stationary
law, equality kernel, ordinary path, or failed analytic property is processed
by the third alternative and is never a terminal conclusion.  Hence contact
forces a nonzero target-aligned cell of the fully saturated five-source
prospective profile.  An upper bound below its retained shell floor
contradicts contact by \cref{thm:source-only-envelope}.
\end{theorem}

\begin{proof}
Shell accountability and finite source exhaustion select a persistent
support and give the floor \(1/5\).  Eventual constancy of the realization
mode gives \(m\), while the cascade supplies \(\Gamma\) and the origin tag.
This constructs \(\mathfrak b\).  Apply
\cref{thm:target-aligned-carrier-duality}.  A finite carrier row telescopes
over the oriented complex containing both selected events and their gaps.
On the zero/root--leaf-bounded storage and finite-exterior-input cell, the
remaining term is nonnegative physical production, and infinite retained
progress contradicts the physical ledger.  A nonzero storage return or
unabsorbed exterior-input coordinate is instead the next target-aligned
successor.  Every zero target projection or terminal
regularity relation gives \textup{(A2)}.  Otherwise
\cref{thm:target-aligned-compactness-routing} produces the next active branch.
Any stationary equality law is restricted to its target-local structural
stratum and passed through
\cref{thm:target-local-stratification-exhaustiveness}.  A finite nonnull cell
gives the proper successor in \textup{(A3)}.  If proper successors continue
cofinally, form their cumulative join.  By
\cref{thm:continuum-composition-no-creation}, it is target-null, exposes a
least unprocessed finite joint increment, or is fixed under every generated
operation after all increments have been source-resolved.
The first gives \textup{(A2)}, the second returns to \textup{(A3)}, and the
third is passed through
\cref{thm:recursive-carrier-polar-exhaustion} for every canonical primal and
polar owner.  Target-nullity gives \textup{(A2)}, cone membership gives
\textup{(A1)}, and a separator is a proper successor in \textup{(A3)}.
Only simultaneous primal--polar fixedness retains the unit shell cocycle and
complete five-channel ancestry and gives \textup{(A4)}.  By
\cref{thm:no-recycled-contact-realization}, its contact-derived origin remains
hypothetical and cannot supply singular existence.  Thus no equality test,
analytic classification, or raw inverse limit remains as a premise.
\end{proof}

\begin{definition}[Canonical activity lift and ordered successor rows]
\label{def:ordered-continuum-successors}
Let \(\mathfrak C_\eta\) be the closure, in the countable product of the
generated evaluation compactification, of all indexed normalized shell
cascades
\[
 \mathbf R=((R_n,Q_n,p_n))_{n\ge1},
 \qquad \operatorname{Prog}_\Sigma(R_n)\ge\eta,
\]
that satisfy the total weak-equation, shell, and event-order graphs.  On an
ordinary contact value, the \(Q_n\) are the canonical pairwise disjoint
physical events, \(p_n=\mu_u(Q_n)\), and the progress is one.  Here \(p_n\)
is the raw event-expenditure coordinate.  It becomes an admissible price for
shell progress only when the target-local factorization satisfies
\eqref{eq:target-aligned-charge-rule}; otherwise the carrier-domination graph
value is retained in \(R_n\) as a successor coordinate.  At a graph
boundary, \(p_n\) is the lower-semicontinuous event-expenditure coordinate.
Boundary points of
\(\mathfrak C_\eta\) retain the sequences that approximate them.  The compact
cascade coordinate \(\mathbf R\) is an \emph{origin tag}; every descendant
below retains it, so successor chains cannot splice different orbits or
cascades.

The tag also carries an origin polarity
\begin{equation}
 \label{eq:origin-polarity}
 \operatorname{pol}_{\rm org}(\mathbf R)
 \in\{\mathsf{Hyp},\mathsf{Ind}\}.
\end{equation}
A record extracted after assuming target contact has polarity
\(\mathsf{Hyp}\).  A record has polarity \(\mathsf{Ind}\) only when its
public orbit is generated from the declared PDE data without using target
contact, noncontinuation, or any descendant of either statement as an input.
Every restriction, limit, quotient, chart, adjoint, polarization, return, and
joint-saturation constructor preserves origin polarity.  In particular,
there is no constructor \(\mathsf{Hyp}\to\mathsf{Ind}\).

Within the successor algebra, ``zero price'' abbreviates zero restriction of
the public physical measure on a retained target-aligned event.  It does not
assert a target-progress domination inequality.  When that inequality is
ordinary and admitted, zero price is also zero admissible target charge; when
it fails, its total-graph coordinate remains part of the same zero-expenditure
successor and is processed before any exclusion inference.

The generated compact observer space \(\Omega_{\rm ev}\) is the product of
the state, spacetime, scale, source, chart, constraint, origin, and total-graph
coordinates occurring in the symbolic transcript.

On the bounded-variation branch, let \(q_n^c\) be the finite signed Radon
measure containing all ordinary, jump, Cantor, oscillation, concentration,
and exterior parts of the derived \(c\)-source shell current.  After pushing
them to \(\Omega_{\rm ev}\), set
\begin{equation}
 \label{eq:canonical-activity-measure}
 a_n:=\sum_{c\in\Channels}(q_n^c)^+,
 \qquad
 \nu_n:=\frac{a_n}{a_n(\Omega_{\rm ev})}.
\end{equation}
The target floor and \eqref{eq:finite-affine-shell-graph} imply
\(a_n(\Omega_{\rm ev})\ge\eta\), and an ordinary normalized contact cascade
has the stronger lower bound one.  If a current has no finite Radon value, its
normalized variation-boundary value and its owner from
\eqref{eq:closed-shell-owner-floor} are placed at the corresponding typed
boundary coordinate; positive rational truncations of the owner current give
the probability-valued lift without summing the remaining extended signed
coordinates.
The lifted record consists of
\begin{equation}
 \label{eq:canonical-activity-lift}
 \widehat R_n=
 \left(\mathbf R,R_n,\nu_n,
       \frac{a_n(\Omega_{\rm ev})}{1+a_n(\Omega_{\rm ev})},
       \frac{p_n}{1+p_n},n\right).
\end{equation}

Run the following total graphs in order on this lift:
\begin{equation}
 \label{eq:ordered-successor-rows}
 \begin{gathered}
 \text{form/constraint realization}
 \longrightarrow \text{chart/gauge/scale coherence},\\
 \text{activity-measure lift}
 \longrightarrow \text{atomic/non-atomic/boundary disintegration},\\
 \text{state/scale disintegration}
 \longrightarrow \text{Young, profile-tree, and tail coordinates},\\
 \text{carrier cohomology}
 \longrightarrow \text{strict passivity/saturation}
 \longrightarrow \text{zero-cost recurrent closure}.
 \end{gathered}
\end{equation}
The grammar includes the cofinal tail shift
\((R_n,Q_n,p_n)\mapsto(R_{n+1},Q_{n+1},p_{n+1})\).  Hence a target-retaining zero-price
witness sequence always has a successor; an analytic failure is carried by
its total-graph coordinate rather than by an artificial failure self-loop.
The familiar realization modes
\begin{equation}
 \label{eq:successor-realization-modes}
 \mathsf{atom},\quad\mathsf{separated},\quad\mathsf{nested},\quad
 \mathsf{diffuse},\quad\mathsf{tail},\quad\mathsf{rough},\quad
 \mathsf{gauge},\quad\mathsf{constraint},\quad\mathsf{equality}.
\end{equation}
are canonical readouts of \(\nu_n\) and its coordinate disintegrations; they
are not asserted to be mutually exclusive alternatives.  Every ordinary or
graph-boundary output retains its five-channel ancestry, original event index,
and target floor.  Hence a failed analytic relation is immediately
another lifted successor input, not an unevaluated row.
\end{definition}

\begin{theorem}[Origin-polarity preservation and no recycled contact]
\label{thm:no-recycled-contact-realization}
The generated continuum algebra preserves
\(\operatorname{pol}_{\rm org}\).  Consequently a contact-selected compiler
execution and every one of its finite or cofinal descendants have polarity
\(\mathsf{Hyp}\).  Such a record may prove a lower bound on the saturated
prospective source profile, but it cannot prove existence of a singular
orbit.  A singular-existence conclusion requires an ordinary public orbit
with polarity \(\mathsf{Ind}\).
\end{theorem}

\begin{proof}
The primitive contact selector is tagged \(\mathsf{Hyp}\).  Every generated
constructor copies the complete origin cascade into its output, and
cumulative joining is defined only for equal origin tags.  Restriction,
quotient, chart, adjoint, polarization, return, closure, totalization, and
joint saturation therefore preserve both the origin and its polarity.
Structural induction over the grammar proves preservation on every finite
descendant.  At a cofinal stage the origin coordinate is constant and belongs
to a discrete closed factor of the joint graph, so it also survives the
cumulative limit.

An independently generated orbit begins in the public-data fiber tagged
\(\mathsf{Ind}\).  Since the grammar contains no origin-changing constructor,
the hypothetical and independent fibers are disjoint.  The target-profile
lower bound uses only the retained shell cocycle and remains valid in the
hypothetical fiber; singular existence is an existential statement in the
independent fiber.
\end{proof}

\begin{theorem}[Compact activity-profile exhaustion]
\label{thm:target-complete-concentration-exhaustion}
Every sequence of lifted records \(\widehat R_n\) has a convergent subsequence
in
\begin{equation}
 \label{eq:activity-lift-compact-space}
 \widehat X_\Sigma
 :=\mathfrak C_\eta\times X_{\rm ev}\times\mathcal P(\Omega_{\rm ev})
   \times[0,1]^2\times\overline{\mathbb N},
\end{equation}
where \(\mathcal P(\Omega_{\rm ev})\) carries the narrow topology and
\(\overline{\mathbb N}=\mathbb N\cup\{\infty\}\).  Its limit retains the
target floor (unit progress for a contact origin) and source ancestry.  The limiting activity law has the
canonical decomposition
\begin{equation}
 \label{eq:canonical-activity-decomposition}
 \nu=\sum_{j\ge1}m_j\delta_{z_j}
      +\nu_{\rm na},
 \qquad m_j>0,
\end{equation}
and its projections to the compactified scale, spatial-boundary, state-law,
and total-graph coordinates record, respectively, nested or separated
profiles, tails, Young oscillation, and rough/gauge/constraint failures.
These coordinates exhaust the limit because they are coordinates of the
single probability law \(\nu\); no separate compactness input or
remainder class is required.
\end{theorem}

\begin{proof}
The spaces \(\mathfrak C_\eta\), \(X_{\rm ev}\), and
\(\Omega_{\rm ev}\) are compact metric.  Hence
\(\mathcal P(\Omega_{\rm ev})\) is compact metric, and so is the product in
\eqref{eq:activity-lift-compact-space}.  Sequential compactness gives a
convergent subsequence.  Unit target progress and the finitely many source
supports are closed coordinates; pass first to a constant-support subsequence.

Every probability measure has a unique at most countable atomic part and a
nonatomic remainder, which proves \eqref{eq:canonical-activity-decomposition}.
Coordinate projection and disintegration record all spatial, scale, state-law,
and graph-boundary behavior.  Boundary escape is mass on the compactification
boundary, rather than failure of tightness.  Since every generated observation
is one of these countable coordinates, equality of all coordinates is equality
of lifted records.  Thus no unrecorded limit output remains.
\end{proof}

\begin{definition}[Witness-indexed retained successor relation]
\label{def:retained-successor-relation}
Let \(\mathfrak A_\eta\subset\widehat X_\Sigma\) be the closed set of lifted
records with target progress at least \(\eta>0\).  An elementary successor
witness is a tuple
\begin{equation}
 \label{eq:elementary-successor-witness}
 (w,w',\mathbf R,\sigma,O),
\end{equation}
where \(\mathbf R=((R_n,Q_n,p_n))_n\in\mathfrak C_\eta\),
\(\sigma:\mathbb N\to\mathbb N\) is strictly increasing, and \(O\) is one
of the ordered total-graph operations in
\eqref{eq:ordered-successor-rows}.  Enumerate these operations and compactify
their index by \(\overline{\mathbb N}\); an operation-index boundary is
represented by the corresponding total-graph limit.  The two endpoints must be limits along
the \emph{same} selector:
\begin{equation}
 \label{eq:same-origin-successor-limits}
 w=\lim_{k\to\infty}\widehat R_{\sigma(k)},
 \qquad
 w'=\lim_{k\to\infty}O^{\rm tot}(\widehat R_{\sigma(k)}),
\end{equation}
and both retain the origin coordinate \(\mathbf R\).  Let
\(\mathfrak W_\eta^0\) be the set of these tuples.  Compactify the selector by
its graph in \(\overline{\mathbb N}^{\mathbb N}\), retain the compactified
operation index, take the closure
\(\mathfrak W_\eta=\overline{\mathfrak W_\eta^0}\), and define
\begin{equation}
 \label{eq:closed-retained-successor}
 \mathcal R_\eta:=\pi_{w,w'}(\mathfrak W_\eta)
 \subseteq\mathfrak A_\eta\times\mathfrak A_\eta.
\end{equation}
The witness space is compact, so \(\mathcal R_\eta\) is closed and every edge
has a nonempty compact witness fiber.  Every witness carries the common
origin, approximating sequence, source tags, unit target coordinate, and
original event indices.  A graph-boundary value is therefore another edge
limit.  It is never promoted to an actual orbit state without its witness
sequence, and edges from different origin cascades cannot be concatenated.
\end{definition}

\begin{lemma}[Diagonal event ancestry and no double counting]
\label{lem:successor-no-double-counting}
Let a finite or infinite chain of successors have one ordinary contact origin
\(\mathbf R=((R_n,Q_n,p_n))_n\).  For every \(k\) one can choose an original
index \(N_k\) such that \(N_{k+1}>N_k\) and the \(N_k\)-th original record,
after the first \(k\) ordered operations in the witness chain, approximates
the \(k\)-th successor coordinate within \(2^{-k}\).  Consequently the physical
events \(Q_{N_k}\) are pairwise disjoint and
\begin{equation}
 \label{eq:successor-no-double-counting}
 \sum_k\mu_u(Q_{N_k})\le\mu_u\!\left(\bigsqcup_nQ_n\right)
 \le\Budget.
\end{equation}
\end{lemma}

\begin{proof}
Choose a witness from the compact fiber of every edge.  At depth \(k\), use
the same-origin limits \eqref{eq:same-origin-successor-limits} and a diagonal
approximation of the first \(k\) witness tuples by elementary witnesses.
Their selectors have cofinal image in the one origin sequence.  Choose
\(N_k>N_{k-1}\) sufficiently far out that the composed depth-
\(k\) graph value is within \(2^{-k}\) of \(w_k\).  No single event is used to
approximate two different depths.  The events therefore belong to distinct
members of the original disjoint shell family, and countable additivity gives
\eqref{eq:successor-no-double-counting}.
\end{proof}

\begin{theorem}[Successor recurrence]
\label{thm:continuum-successor-recurrence}
The relation \(\mathcal R_\eta\) is closed in the compact metric lifted
topology.  If there is an infinite chain
\begin{equation}
 \label{eq:infinite-successor-chain}
 w_0\mathcal R_\eta w_1\mathcal R_\eta w_2\mathcal R_\eta\cdots,
\end{equation}
then the path space of all such chains is compact and invariant under the
left shift.  The orbit closure of the given path carries a shift-invariant
Borel probability measure.  Every point in the support of its coordinate
projection is a limit generated from one retained cascade origin with a
diagonal witness sequence,
retains the target floor, and retains the union of source ancestries.  No claim
of pointwise realization by a new PDE orbit is used.
\end{theorem}

\begin{proof}
The witness bundle is compact and coordinate projection is continuous, so
\(\mathcal R_\eta\) is compact and hence closed.  The set
\(\mathfrak A_\eta\) is compact in the lifted topology,
so \(\mathfrak A_\eta^{\mathbb N}\) is compact.  The path space is the
intersection, over all \(j\), of the closed conditions
\((w_j,w_{j+1})\in\mathcal R_\eta\); it is therefore compact and is nonempty
by \eqref{eq:infinite-successor-chain}.  Empirical measures of the shifted
given path have a weakly convergent subsequence, and the endpoint calculation
proves shift invariance.  Apply
\cref{lem:successor-no-double-counting} to obtain the diagonal original-event
witnesses.  Sourcewise closure preserves ancestry and the target floor.
\end{proof}

\begin{definition}[Target-aligned ordered path calculus]
\label{def:target-aligned-ordered-path-calculus}
Fix an active branch
\(\mathfrak b=(A,m,\Gamma,\eta,\mathbf R)\).  Let
\(\mathfrak A_{\mathfrak b}\) be the closed lifted stratum with that source
support, realization mode, shell family and floor, origin tag, and all prior
zero transcript coordinates.  Let
\(\mathcal B_{\mathfrak b}^{\mathbb Q}\) be the unital rational cylinder
algebra generated by finitely supported path coordinates on
\(\mathfrak A_{\mathfrak b}^{\mathbb N}\), including:
\begin{enumerate}[label=\textup{(\alph*)}]
\item the graph-complete equation, carrier, operation-defect, and source
coordinates at every rational path time;
\item the closed successor-witness coordinates for every adjacent pair;
\item the bounded price coordinate \(\bar p=p/(1+p)\);
\item the normalized target pairing and the positive sourcewise shell masses;
\item the shift pullback \(S^*\).
\end{enumerate}
Generators annihilated by every element of \(\Gamma_\Sigma\) are removed
before \(\mathcal B_{\mathfrak b}^{\mathbb Q}\) is formed.

Let \(I_{\mathfrak b}\) be the ideal generated by the weak equation, signed carrier,
typed graph, same-origin, successor, zero-density, target-floor, and
source-support equalities together with the complete active-branch signature.
The zero-density generator acts only on the normalized support coordinate.
The unnormalized occurrence cocycle, physical expenditure, and source currents
are independent coordinates and no element of \(I_{\mathfrak b}\) identifies
them with zero.
Let \(M_{\mathfrak b}\) be the quadratic module
generated by squares, normalized coordinate bounds, nonnegative production
and defect masses, and the positive target-alignment coordinates.  Finally,
let
\[
 \mathcal D_S
 :=\operatorname{span}_{\mathbb Q}
 \{f\circ S-f:f\in\mathcal B_{\mathfrak b}^{\mathbb Q}\}
\]
and define the ordered proof cone
\begin{equation}
 \label{eq:target-aligned-proof-cone}
 \mathcal K_{\mathfrak b}:=
 I_{\mathfrak b}+M_{\mathfrak b}+\mathcal D_S.
\end{equation}
A finite ordered derivation is an explicit equality
\[
 -1=i+m+d,\qquad
 i\in I_{\mathfrak b},\quad m\in M_{\mathfrak b},\quad d\in\mathcal D_S,
\]
written with rational cylinder polynomials.
\end{definition}

\begin{theorem}[Ordered proof or stationary structural model]
\label{thm:ordered-proof-or-model}
For every reachable active branch \(\mathfrak b\), exactly one of the following holds.
\begin{enumerate}[label=\textup{(O\arabic*)}]
\item \(-1\in\mathcal K_{\mathfrak b}\).  Then a finite ordered derivation proves
that no stationary target-aligned zero-density successor support in \(X_{\mathfrak b}\)
exists.
\item There is a linear functional
\(L:\mathcal B_{\mathfrak b}^{\mathbb Q}\to\mathbb R\) such that
\begin{equation}
 \label{eq:ordered-model-functional}
 L(1)=1,\qquad L(M_{\mathfrak b})\ge0,\qquad
 L(I_{\mathfrak b})=L(\mathcal D_S)=0.
\end{equation}
Its continuous extension is integration against a shift-invariant Borel
probability law on the graph-complete path space.  That law has zero physical
expenditure density and satisfies the full signature \(\mathfrak b\), including target
pairing at least \(\eta\), source support \(A\), mode \(m\), origin
\(\mathbf R\), and all same-origin successor coordinates.  It is returned
jointly with the unnormalized multiplicity ledger inherited from \(\mathbf R\).
\end{enumerate}
The alternatives are exclusive.  Every assertion in \textup{(O1)} is checked
by finite rational algebra; every object in \textup{(O2)} is reconstructed
from the public graph relations.
\end{theorem}

\begin{proof}
Pass first to the rational vector-space quotient by
\(I_{\mathfrak b}+\mathcal D_S\).  The normalized coordinate bounds make the image
of \(M_{\mathfrak b}\) an Archimedean order cone with order unit \([1]\).  If
\(-[1]\) is not in this cone, the order-unit separation theorem gives a
positive linear functional with value one on \([1]\).  Pulling it back gives
\eqref{eq:ordered-model-functional}.  Positivity makes the functional bounded
in the order-unit norm.  The Archimedean representation theorem for quadratic
modules, followed by the Riesz representation theorem, therefore represents
its continuous extension by a Borel probability measure supported on the
joint positivity set of the graph-complete path relations
\citep{Schmudgen2017}.  Annihilation of
\(\mathcal D_S\) is precisely shift invariance.  Annihilation of
\(I_{\mathfrak b}\) imposes the equation, successor, same-origin, zero-density,
target-floor, source, mode, origin, and prior-zero relations.  This proves
\textup{(O2)}.  Since the cumulative coordinates are not generators of this
zero-density ideal, separation cannot erase them.

Conversely, a functional satisfying \eqref{eq:ordered-model-functional} sends
every \(i+m+d\in\mathcal K_{\mathfrak b}\) to a nonnegative number.  It therefore
cannot coexist with \(-1\in\mathcal K_{\mathfrak b}\).  If no such functional
exists, order-unit separation gives \(-1\in\mathcal K_{\mathfrak b}\), proving
exhaustiveness.
An element of the algebraic sum is, by definition, a finite expression, so
the derivation uses finitely many rational relations.
\end{proof}

\begin{definition}[Finite cylinder-moment hierarchy]
\label{def:finite-cylinder-moment-hierarchy}
Fix an active branch \(\mathfrak b\) and enumerations of its generators and
proof rules.  At level \(\ell=(n,d,N)\), retain the first \(n\) generators, monomials of degree at
most \(2d\), and successor coordinates through time \(N\).  The primal cone
\(\mathsf P_{\mathfrak b}^{\ell}\) consists of truncated functionals whose moment
and localizing matrices for \(M_{\mathfrak b}\) are positive semidefinite and which
annihilate the retained ideal and shift-difference rows.  The dual cone
\(\mathsf D_{\mathfrak b}^{\ell}\) consists of rational identities
\begin{equation}
 \label{eq:finite-dual-derivation}
 -1=i_\ell+\sum_j s_j^2m_j+d_\ell
\end{equation}
using only the retained rows.  Exact rational matrix arithmetic and rational
interval enclosures check every dual identity.
\end{definition}

\begin{theorem}[Proof-producing hierarchy completeness]
\label{thm:proof-producing-hierarchy-completeness}
For a fixed reachable active branch \(\mathfrak b\):
\begin{enumerate}[label=\textup{(\roman*)}]
\item every feasible primal truncation gives valid lower moment bounds and
every dual derivation is a sound exclusion of the source cell;
\item if \(-1\in\mathcal K_{\mathfrak b}\), then
\(\mathsf D_{\mathfrak b}^{\ell}\) contains a derivation at some finite level;
\item if every finite level is primal feasible, a diagonal subsequence of
normalized moment tables converges to a functional satisfying
\eqref{eq:ordered-model-functional}.
\end{enumerate}
Thus semantic membership and compact-set inclusion are absent from the
pruning interface.
\end{theorem}

\begin{proof}
Weak duality follows by applying a feasible positive truncated functional to
\eqref{eq:finite-dual-derivation}.  A finite expression for
\(-1\in\mathcal K_{\mathfrak b}\) contains finitely many generators, a finite
degree, and a finite shift depth, so it occurs at some level, proving
\textup{(ii)}.  For \textup{(iii)}, the Archimedean bounds make each scalar
moment sequence bounded.  Repeated diagonal extraction gives a value on every
cylinder monomial.  Positivity of every finite moment and localizing matrix,
and annihilation of every retained relation at all sufficiently large levels,
pass to the limit.  The resulting normalized positive functional satisfies
\eqref{eq:ordered-model-functional}; apply
\cref{thm:ordered-proof-or-model}.
\end{proof}

\begin{proposition}[Sound scope of the ordered hierarchy]
\label{prop:sound-scope-ordered-hierarchy}
The ordered hierarchy has the following exact logical contract.
\begin{enumerate}[label=\textup{(S\arabic*)}]
\item A displayed finite identity \(-1\in\mathcal K_{\mathfrak b}\) is a
  checked exclusion of its stated finite target cell.
\item Feasibility of every finite relaxation reconstructs a positive
  functional on the completed ordered cylinder algebra.  It does not by
  itself reconstruct a PDE orbit, an analytic continuation germ, or a
  singular solution.
\item Nonpolynomial maps, unbounded operators, distributional products,
  measures, and limiting relations are represented by their bounded total
  graphs.  Finite semialgebraic elimination applies only to the finite
  polynomial moment relaxation of those graph coordinates.
\item Strict inequalities may be removed by an Archimedean ordered identity.
  A zero-margin equality cell remains a positive-functional successor unless
  a finite ideal/quadratic-module identity removes it or the
  approximation--orbit criterion realizes it.
\item Empty decreasing compact semantic sets imply an empty finite level, but
  a regularity proof uses that level only with the recorded finite deletion
  identities.  Semantic emptiness is not a proof object.
\item A primal--polar fixed point and a cofinal cumulative successor are
  internal graph-complete states.  They become physical terminals only
  through \cref{thm:approximation-orbit-realization-criterion}; otherwise
  their next nonordinary, cone-boundary, or target-visible coordinate remains
  in the successor algebra.
\end{enumerate}
Consequently the hierarchy is proof-producing and sound, but failure to find
a finite identity is never interpreted as a physical counterexample.  No
relative-completeness claim for all true PDE inequalities or real-radical
equalities is used anywhere in a reachability conclusion.
\end{proposition}

\begin{proof}
Items \textup{(S1)}--\textup{(S2)} are
\cref{thm:ordered-proof-or-model,thm:proof-producing-hierarchy-completeness}.
Item \textup{(S3)} follows from total graph completion and the definition of
the finite moment levels.  For \textup{(S4)}, strict Archimedean positivity
has a finite quadratic-module representation, whereas equality is imposed in
the ideal and its complementary functional is retained.  Nested compactness
proves the semantic assertion in \textup{(S5)}; checking the concatenated
finite identities supplies the only permitted proof step.  Item \textup{(S6)}
is the Dirac realization gate and the no-compactification-promotion
corollary.  These implications establish the contract without assuming a
complete calculus for the full analytic graph language.
\end{proof}

\begin{definition}[Prospective zero-density successor fixed point]
\label{def:total-successor-saturation-algebra}
Fix a reachable active branch \(\mathfrak b\).  In this definition and the
following pruning results, \(\mathfrak A_\eta\), \(\mathcal R_\eta\),
\(Z_n\), and \(Z_\infty^\Sigma\) denote their restrictions to the closed
same-origin stratum \(\mathfrak A_{\mathfrak b}\); the branch index is
suppressed only to keep the formulas readable.  Let
\(p:\mathfrak A_{\mathfrak b}\to[0,\infty]\) be the lower-semicontinuous relaxation
of the original-event price coordinate.  For a closed
\(K\subseteq\mathfrak A_{\mathfrak b}\), define the viability operator
\begin{equation}
 \label{eq:zero-cost-viability-operator}
 \mathcal V_\eta(K)
 :=\left\{w\in K:p(w)=0,
       \ \mathcal R_\eta(w)\cap K\ne\varnothing\right\}.
\end{equation}
It maps closed sets to closed sets, because it is the projection of the compact
set \(\mathcal R_\eta\cap(K\times K)\cap(\{p=0\}\times K)\).
Starting with the target-active zero-density support set
\(Z_0=\mathfrak A_\eta\cap\{p=0\}\), set
\begin{equation}
 \label{eq:finite-depth-zero-cost-closure}
 Z_{n+1}=\mathcal V_\eta(Z_n),
 \qquad
 Z_\infty^\Sigma=\bigcap_{n=0}^{\infty}Z_n.
\end{equation}
Thus \(Z_n\) consists exactly of the support records from which a target-retaining
zero-density successor path of length \(n\) can be generated.  Every such
record remains paired with its finite-prefix multiplicity ledger.  Compactness and
closedness give
\begin{equation}
 \label{eq:omega-viability-fixed-point}
 \mathcal V_\eta(Z_\infty^\Sigma)=Z_\infty^\Sigma.
\end{equation}
Indeed, for \(w\in Z_\infty^\Sigma\), choose
\(w_n\in Z_n\) with \(w\mathcal R_\eta w_n\); a convergent subsequence has a
limit in every \(Z_n\), and the closed successor graph retains the edge.
Consequently the countable finite-depth iteration is the complete topological
description of viability.  It is not, by itself, a proof procedure.
If \(Z_\infty^\Sigma=\varnothing\), nested compactness implies
\(Z_N=\varnothing\) for some finite \(N\).  This semantic fact becomes a
checkable exclusion only when the ordered hierarchy below produces an
explicit identity \(-1\in\mathcal K_{\mathfrak b}\) on each active component.  If
every \(Z_n\) is nonempty, their intersection contains an infinite zero-density
support path; this statement does not identify the accumulated cocycle with
its normalized support law.

Let \(\mathsf{Path}(Z_\infty^\Sigma)\) be the compact path space of
\(\mathcal R_\eta\) restricted to this fixed point.  Let
\(\mathcal P_{\mathcal R_\eta}^{0,\Sigma}\) be all shift-invariant Borel
probability laws on this path space.  The prospective successor-saturation
record associated with \(\boldsymbol\nu\) is its joint law
\(\operatorname{Aug}(\boldsymbol\nu)\) with the prefix functions
\[
 \Gamma_N=\sum_{j<N}\delta_{e_j},\qquad
 A_{\Phi,N}=\sum_{j<N}a_\Phi(e_j),\qquad
 \Pi_N=\sum_{j<N}p(e_j),\qquad
 Q_N^c=\sum_{j<N}q^c(e_j).
\]
These are random finite measures or scalar coordinates on the path space;
their joint law is determined by \(\boldsymbol\nu\) and is not extra data.
The prospective successor-saturation
algebra is
\begin{equation}
 \label{eq:total-successor-saturation-algebra}
 \mathfrak H_{\rm succ}^{\Sigma}
 :=
 \overline{\operatorname{span}}
 \left\{
  \bigl(\Pi_\Sigma(\pi_0)_\#\boldsymbol\nu,
        \operatorname{Aug}(\boldsymbol\nu)\bigr):
  \boldsymbol\nu\in\mathcal P_{\mathcal R_\eta}^{0,\Sigma}
 \right\}.
\end{equation}
Every object in this definition is constructed prospectively from
\(\Theta_\Sigma^{\rm syn}\).  Neither actual contact orbits nor exact invariant
values occur.  A total-graph failure contributes only when it belongs to an
infinite target-retaining zero-density support path; naming the failure
supplies no self-loop.  The augmented ledger is part of every generator, so
the closed span is taken in the product of the support and cumulative-cocycle
spaces rather than in the support space alone.
\end{definition}

\begin{definition}[Finite ordered pruning transcript]
\label{def:finite-pruning-transcript}
Put
\begin{equation}
 \label{eq:elimination-rank-sets}
 E_n:=\mathfrak A_\eta\setminus Z_n.
\end{equation}
A \emph{finite ordered pruning transcript} for active branch \(\mathfrak b\) is
the complete finite expression
\begin{equation}
 \label{eq:finite-pruning-node}
 -1=i+\sum_{j=1}^{M}s_j^2m_j+d,
 \qquad
 i\in I_{\mathfrak b},\quad m_j\in M_{\mathfrak b}^{\rm gen},
 \quad d\in\mathcal D_S,
\end{equation}
including the rational coefficients and the indices of every equation,
carrier, graph, source, target-floor, same-origin, successor, positivity, and
shift relation used.  Here \(M_{\mathfrak b}^{\rm gen}\) is the enumerated set of
quadratic-module generators.  Verification consists only of expanding
\eqref{eq:finite-pruning-node}, checking the cited ideal and shift rows, and
checking that each multiplier is a square.  The transcript is accepted only
after this finite calculation succeeds.  Its existence is never inferred
from semantic emptiness of \(Z_n\).  Among valid transcripts use the first in
the fixed enumeration.
\end{definition}

\begin{theorem}[Ordered pruning or same-origin structural model]
\label{thm:canonical-finite-pruning-transcript}
The elimination sets are open and satisfy
\begin{equation}
 \label{eq:elimination-rank-recursion}
 E_{n+1}
 =E_n\cup
 \left\{w\in Z_n:
   \mathcal R_\eta(w)\cap Z_n=\varnothing\right\}.
\end{equation}
The recursion is the topological shadow of the ordered hierarchy.  Exactly one
of the following proof-producing alternatives holds for every reachable
active branch \(\mathfrak b\).
\begin{enumerate}[label=\textup{(P\arabic*)}]
\item At some finite moment level, the dual hierarchy produces
\(-1\in\mathcal K_{\mathfrak b}\).  The identity itself is the finite proof-valid
pruning transcript and excludes the cell.
\item \(-1\notin\mathcal K_{\mathfrak b}\).  Order-unit separation gives the
positive functional of \cref{thm:ordered-proof-or-model}, hence an infinite
same-origin zero-price path law with target floor \(\eta\).  Equivalently, a
compatible family of all finite moment restrictions reconstructs this law.
\end{enumerate}
Thus every pruning assertion is verified by the ordered proof kernel.  A
presently unsuccessful finite search is neither an exclusion nor a model; the
model branch is the mathematical alternative \(-1\notin\mathcal K_{\mathfrak b}\).
It is successor data: zero-price resaturation must still restrict its kernel,
reselect its prospective source, and totalize its realization graph.
\end{theorem}

\begin{proof}
The set \(E_0\) is open because \(Z_0\) is closed.  If \(E_n\) is open, then
\begin{equation*}
 \left\{w:\mathcal R_\eta(w)\cap Z_n\ne\varnothing\right\}
 =\pi_1\!\left(
   \mathcal R_\eta\cap(\mathfrak A_\eta\times Z_n)
  \right)
\end{equation*}
is compact and hence closed.  This proves
\eqref{eq:elimination-rank-recursion} and openness of \(E_{n+1}\).

Apply \cref{thm:proof-producing-hierarchy-completeness}.  In its dual branch,
the finite identity is exactly \eqref{eq:finite-pruning-node}; applying any
putative positive stationary functional gives \(-1\ge0\), so the cell is
excluded.  In the other algebraic branch,
\cref{thm:ordered-proof-or-model} supplies a normalized positive functional
annihilating the successor and shift-difference relations and reconstructs its
stationary path law.
Same-origin coordinates belong to the ideal and hence hold almost surely.
The alternatives are exclusive by weak duality.
\end{proof}

\begin{definition}[Generated zero-density support dynamics]
\label{def:generated-saturation-equations}
Let \(\boldsymbol\nu\) be a stationary target-retaining successor law with
zero physical-expenditure density, paired with the unnormalized occurrence
ledger \((\Gamma_N,A_{\Phi,N},\Pi_N,(Q_N^c)_c)_N\).  The density equality is
an identity of the support selector; it imposes no zero relation on
\(\Pi_N\) or \(Q_N^c\).  Let
\(\mathsf{Conseq}(\mathfrak P)\) be the least countable consequence algebra
containing the public weak equation and native carrier identity and closed
under rational linear combination, localization by the generated test core,
integration by parts already encoded in the weak formulation, polarization,
adjoint formation, quotient descent, chart pullback, Jordan decomposition,
completion of squares, composition of total graphs, and sourcewise graph
closure.  If a product, chain, trace, or covariance operation lacks an
ordinary value, the algebra inserts its total-graph defect rather than the
desired identity.  This grammar is enumerated from
\(\mathsf{Pre}(\mathsf{Eq})\) and has no additional input slot.

Let \(\mathscr P_{\rm gen}\) be the countable family of all generated
nonnegative functions \(q\) for which
\(\mathsf{Conseq}(\mathfrak P)\) contains a domination identity
\begin{equation}
 \label{eq:generated-nonnegative-completion}
 p=q+r_q,
 \qquad q\ge0,\quad r_q\ge0.
\end{equation}
Include \(q=p\) and \(r_p=0\).  This definition takes the union of all valid
positive descendants and does not require a choice of a maximal or unique
decomposition.  The family is generated by the fixed equation and carrier:
mobility squares, weak energy defects, positive commutator squares,
Fisher-information or Hessian squares, constraint defects, and target-ray
defects occur only when their defining identity is already a relation in
\(\mathsf{Conseq}(\mathfrak P)\).  No extra positive term is inserted.

The \emph{saturation equations} are
\begin{equation}
 \label{eq:generated-saturation-equations}
 q=0\quad(q\in\mathscr P_{\rm gen}),
 \qquad
 \dot u=J_{\Geom}+J_{\Caus}+J_{\Abs}+J_{\Lift}+J_{\Bdry},
\end{equation}
together with the totalized constitutive, constraint, chart, and covariance
relations.  No tangency to a symmetry orbit is asserted.  Equilibria, relative
equilibria, periodic or recurrent conservative orbits, scale breathers, and
persistent graph-boundary dynamics are retained whenever they solve these
relations.  If
\(\mathcal F_{\rm gen}\) is the generated finite or countable family of
partial relations used on the current branch, let
\(\partial_\infty F\) denote the recurrent graph-boundary values obtained
from infinite witness-successor paths for \(F\), rather than the entire graph
boundary.  Define the totalized saturation set
\begin{equation}
 \label{eq:totalized-saturation-set}
 \mathscr Z_{\rm sat}^{\rm tot}
 :=
 \left(
  \bigcap_{q\in\mathscr P_{\rm gen}}\ker q
  \cap\{\dot u=J_{\rm tot}\}
 \right)
 \sqcup
 \bigsqcup_{F\in\mathcal F_{\rm gen}}\partial_\infty F .
\end{equation}
Thus a failed relation enters the saturation set only through a recurrent
zero-density graph-boundary path retaining its multiplicity ledger.  The definition imposes only identities already
derived from the equation and does not convert stationarity of a law into
stationarity of a state.
\end{definition}

\begin{theorem}[Saturation-equation compiler]
\label{thm:saturation-equation-compiler}
Every generator of \(\mathfrak H_{\rm succ}^{\Sigma}\) is supported on the
totalized solution set of \eqref{eq:generated-saturation-equations} and on the
greatest viable set \(Z_\infty^\Sigma\).  Conversely, every stationary law on
the prospective path space supported on this solution set generates a class
in \(\mathfrak H_{\rm succ}^{\Sigma}\).  Hence
\begin{equation}
 \label{eq:saturation-algebra-kernel-form}
 \mathfrak H_{\rm succ}^{\Sigma}
 =
 \overline{\operatorname{span}}\,
 \Pi_\Sigma
 \left(
 Z_\infty^\Sigma
 \cap\mathscr Z_{\rm sat}^{\rm tot}
 \right),
\end{equation}
where the right side denotes the coordinate supports of stationary path laws
on the displayed set.  It is generated entirely by algebraic kernels, the
closed successor graph, source quotients, and the descending viability
operator.
\end{theorem}

\begin{proof}
For a stationary zero-density selector law,
\(\int p\,\dd\boldsymbol\nu=0\), so \(p=0\) almost surely.  For every
\(q\in\mathscr P_{\rm gen}\),
\eqref{eq:generated-nonnegative-completion} gives \(0\le q\le p\), hence
\(q=0\) almost surely on the ordinary branch.  The full weak
current identity remains one of the total graph coordinates, so its ordinary
support solves \(\dot u=J_{\rm tot}\) without any symmetry-tangency inference.
On a graph-boundary branch, membership in the path fixed point supplies the
recurrent value in \(\partial_\infty F\).  This proves the first inclusion in
\eqref{eq:saturation-algebra-kernel-form}.  Conversely, a stationary path law
in the displayed ordinary joint kernel has \(p=\mathcal P_0=0\); a recurrent
graph-boundary law has the same zero price by membership in
\(Z_\infty^\Sigma\).  These are density statements on the selected support.
The accompanying cumulative coordinates retain the exact identities
\(\Gamma_N=A_{\Phi,N}\gamma_N\),
\(\Pi_N=A_{\Phi,N}\int p\,\dd\gamma_N\), and
\(Q_N^c=A_{\Phi,N}\int q^c\,\dd\gamma_N\); they are not set to zero.
Both are generators of the successor-saturation algebra.
Taking target projections and closed spans proves equality.
\end{proof}

\begin{corollary}[Complete zero-cost structural reduction]
\label{cor:algebraic-equality-reduction}
The finite-depth recursion \eqref{eq:finite-depth-zero-cost-closure} has as its
intersection the largest target-retaining zero-cost successor set.  Its recurrent projection
satisfies the exact alternative
\begin{equation}
 \label{eq:zero-cost-structural-alternative}
 \mathfrak H_{\rm succ}^{\Sigma}=0
 \quad\Longleftrightarrow\quad
 \Pi_\Sigma\supp\boldsymbol\nu=0
 \text{ for every stationary law on }Z_\infty^\Sigma.
\end{equation}
At each successor stage the compiler applies the five-source quotient to the
full zero-cost weak equation and to every recurrent graph-boundary coordinate.
A target-null output is removed at the next stage.  Any nonzero fixed-point
output is a positive structural model of recurrent conservative,
relative-equilibrium, profile, or graph-boundary dynamics satisfying the
relations retained so far.  It is not asserted to be an actual singular PDE
orbit and is not terminal.  It is passed to the zero-price resaturation
operator of \cref{def:zero-price-resaturation-successor}; only the exhaustive
alternatives of \cref{thm:no-terminal-zero-price-model} may terminate that
route.  No compactness gap or separate rigidity datum appears in this
reduction.  A checked identity \(-1\in\mathcal K_{\mathfrak b}\) closes a finite
stage; the complementary positive functional supplies the next kernel and
prospective-source calculation.
\end{corollary}

\begin{theorem}[Zero-density successor support with multiplicity retention]
\label{thm:zero-mean-successor-reduction}
Let an infinite witness-indexed successor chain originate from the disjoint
shell sequence of a finite-expenditure contact orbit.  Then every coordinate
of its diagonal path has relaxed original-event expenditure zero, and its path
closure carries a stationary successor support law with zero physical
expenditure density.  The law is paired with the original cumulative records,
which satisfy
\[
 A_{\Phi,N}=N,\qquad
 \Pi_N=\sum_{j<N}p_j,\qquad
 \sum_cQ_N^c=N.
\]
Consequently its augmented target projection lies in
\(\mathfrak H_{\rm succ}^{\Sigma}\).  No cumulative physical or structural
coordinate is declared zero.
\end{theorem}

\begin{proof}
By countable additivity, the original event expenditures
\(p_n=\mu_u(Q_n)\) are summable and hence tend to zero.  Every successor uses a
subsequence of these same indexed events.  The diagonal selection in
\cref{lem:successor-no-double-counting} therefore approximates each path
coordinate by records whose expenditures tend to zero.  Lower semicontinuity
and nonnegativity give zero expenditure at every path coordinate.  Whenever
that coordinate is used as target price, admissibility follows separately
from \cref{thm:local-target-aligned-charging-rule}; otherwise its failed
domination relation remains in the successor graph.

\Cref{thm:continuum-successor-recurrence} gives a compact shift-invariant path
closure.  A weak limit of its empirical path measures is shift invariant and
has zero price density.  The empirical normalization is stored jointly with
the original occurrence ledger and the exact identities
\(\Gamma_N=A_{\Phi,N}\gamma_N\) and
\(\Pi_N=A_{\Phi,N}\int p\,\dd\gamma_N\) restore multiplicity before target
projection.  Target retention and source ancestry are closed
coordinates, so the augmented law is one of the generators in
\eqref{eq:total-successor-saturation-algebra}.
\end{proof}

\begin{theorem}[Estimate-free finite-budget contact routing]
\label{thm:finite-price-saturation-route}
Let a contact orbit with finite physical expenditure generate infinitely many
pairwise disjoint target-retaining windows \(I_j\).  The canonical activity
lift and cofinal tail shift generate an infinite path in
\(Z_\infty^\Sigma\), and
\begin{equation}
 \label{eq:finite-price-to-fixed-point}
 \Pi_\Sigma Z_\infty^\Sigma\ne0,
 \qquad
 \mathfrak H_{\rm succ}^{\Sigma}\ne0.
\end{equation}
Thus finite price is completely reduced to an explicit recurrent zero-density
support with its full cumulative occurrence, physical, and source ledger; it
is not a terminal scalar branch or an outstanding analytic row.
No lower price, compactness estimate in the original state topology,
coercivity constant, or Liouville input is used: the only numerical fact is
countable additivity of the equation-generated physical measure.
\end{theorem}

\begin{proof}
Countable additivity gives
\(\sum_j\mu_u(I_j)<\infty\), hence after a subsequence
\(\mu_u(I_j)\to0\).  Apply the compact activity lift.  Its limit has relaxed
price zero, unit endpoint shell increment, the closed source-owner floor, and
inherited source support.  Apply each
ordered total graph to cofinal tails.  Ordinary and boundary values are both
coordinates of the compact lift, and the tail shift supplies a successor at
every stage.  This gives an infinite target-retaining density-zero support
path while retaining the unnormalized prefix ledger.

The descending viability recursion retains exactly the points admitting all
finite continuations of such a path; compactness supplies an infinite path in
their intersection.  Hence the path lies in \(Z_\infty^\Sigma\).
\Cref{thm:zero-mean-successor-reduction} supplies a stationary zero-density
support law with multiplicity restored.
Its unit shell coordinate and positive owner coordinate make its target
projection nonzero, proving
\eqref{eq:finite-price-to-fixed-point}.
\end{proof}

\begin{theorem}[Ordered prospective regularity alternative]
\label{thm:structural-bruteforce-regularity}
For a fixed canonical presentation, run the total local-solution and carrier
graphs, source compiler, canonical shells, signed carrier compiler,
source-adapted dual-carrier grammar, activity lift, and zero-cost successor
fixed point.  Every target-visible local-solution or carrier defect is
inserted into the same successor graph.  Define the generated regularity mechanism
algebra
\begin{equation}
 \label{eq:terminal-structural-algebra}
 \mathfrak H_{\rm out}^{\Sigma}
 :=\overline{\operatorname{span}}
  \{\Pi_\Sigma\mathfrak a:\mathfrak a\text{ is a joint-saturated
      prospective source atom}\},
\end{equation}
where explicit realization is understood in the sense of
\cref{def:explicit-singular-mechanism}.  Put \(\eta_0=1/5\).  Exactly
one of the following algebraic alternatives holds:
\begin{enumerate}[label=\textup{(R\arabic*)}]
\item every algebraically reachable active branch is eliminated by a local
  target-null, continuation, class-incompatibility, or divergent-charge
  relation produced by polarity refinement and resaturation; every ordinary
  maximal finite-expenditure branch then remains inside the declared
  continuation graph and avoids the complete target;
\item zero-price resaturation reconstructs a prospective source atom
  \(\mathfrak a\).  It supplies the closed reduced weak equation, source
  ancestry, vanishing productions, complete typed values for every auxiliary
  graph and their descendants, physical and origin ledgers, and a retained
  source-profile floor at least \(\eta_0\).
\end{enumerate}
Independently, if the public local-solution or carrier graph has no ordinary
value, its explicit first-failure graph coordinate is retained.  This is an
existence or carrier-realization output, not a regularity conclusion, and the
regularity alternative is applied only to ordinary fibers.
Every actual finite-expenditure contact orbit constructs a hypothetical
output of type \textup{(R2)} with its own diagonal event witnesses;
stationary positive laws
and equality kernels encountered on the way are nonterminal resaturation
states.  Consequently
\begin{equation}
 \label{eq:terminal-algebra-regularity}
 \boxed{
 \begin{gathered}
 \mathsf T_\Sigma^{\Phys,\Sat}(\Budget)<\eta_0
 \\
 \Longrightarrow
 \text{ every ordinary maximal finite-carrier branch avoids the declared target.}
 \end{gathered}}
\end{equation}
This is the continuum prospective-profile bound corresponding to Part~I's
universal accountability inequality.  A closing branch carries
the finite transcript \eqref{eq:finite-pruning-node}; the complementary branch
is iterated by \cref{thm:no-terminal-zero-price-model} until it closes or
produces the source-profile data listed in \textup{(R2)}.  Singular existence
is asserted only by an independently realized member satisfying
\cref{def:explicit-singular-mechanism}.
\end{theorem}

\begin{proof}
Insert ordinary values of the total local-solution and carrier graphs into the
activity lift.  The form value \(\bot_{\rm loc}\) is retained separately.  On
an ordinary finite-expenditure branch, suppose a maximal
orbit contacts the complete target.  Canonical shells give the disjoint unit
increment cascade and its source-owner floor.  A divergent generated lower-price series contradicts
countable additivity; the finite-price tail instead enters the successor
graph.  \Cref{thm:finite-price-saturation-route} constructs a stationary law
in \(Z_\infty^\Sigma\) with nonzero target coordinate and original shell-event
witnesses.  \Cref{thm:continuous-selector-accountability-closure} either
closes its active branch or reprocesses its typed equality successor.
\Cref{thm:no-terminal-zero-price-model} turns every nonclosing compatible
route into the source-profile witness \textup{(R2)}.  Its origin remains
hypothetical by \cref{thm:no-recycled-contact-realization}.  Otherwise each finite local elimination in
\textup{(R1)} contradicts the retained target coordinate; diagonal
saturation supplies the same conclusion for a cofinal target-null route.
If the saturated profile is below \(\eta_0\), \textup{(R2)} is impossible;
terminal coverage then keeps every ordinary maximal branch inside the
declared continuation graph.  This
proves \eqref{eq:terminal-algebra-regularity}.
\end{proof}

\begin{proposition}[Covariant entropy as a carrier correction]
\label{prop:covariant-entropy-carrier-correction}
Let a covariant entropy cocycle \(\Entropy\) satisfy the gluing law of
\cref{thm:entropy-cocycle-gluing}.  If its local balance has the form
\begin{equation}
 \label{eq:entropy-carrier-correction}
 \Entropy(Sw)-\Entropy(w)
 =F_{\rm sc}(w)-D(w)-M(w)-R(w),
 \qquad R\ge0,
\end{equation}
then \(\Psi=-\Entropy\) is the canonical carrier correction of
\cref{def:carrier-cohomology}, and
\begin{equation}
 \label{eq:entropy-passivity-deficit}
 D+M-F_{\rm sc}+\Psi-\Psi\circ S=R\ge0.
\end{equation}
Every invariant saturation measure is supported on \(R=0\).  It is supported
on \(D=M=0\) only when the generated carrier transcript supplies a
nonnegative domination identities expressing the zero-price deficit with
\(D\) and \(M\) as positive descendants.  Thus a Perelman-type formula refines the carrier cohomology and
zero-price equations.  Its equality models are exactly the ordinary or
recurrent graph-boundary solutions of those generated equations; relative
equilibria are one possible subclass, not an automatic conclusion.
\end{proposition}

\begin{proof}
Rearrange \eqref{eq:entropy-carrier-correction} to obtain
\eqref{eq:entropy-passivity-deficit}.  The entropy cocycle glues on one
coherent orbit, so its increments telescope and define a carrier coboundary.
For an invariant saturation measure the nonnegative right side has zero
integral and therefore \(R=0\) almost surely.  Any further vanishing follows
only by applying the generated positive-descendant rule of
\cref{def:generated-saturation-equations}; the full weak equation remains in
the support coordinates.
\end{proof}

\begin{corollary}[Gradient-kernel equality reduction]
\label{cor:gradient-kernel-equality-reduction}
Let \(\Omega\) be connected and let the ordinary saturation equations on a
normalized successor core contain
\begin{equation}
 \label{eq:local-gradient-kernel-production}
 D_m(V):=\int_{K_m}|\nabla V|^2\,\dd y=0,
 \qquad m\ge1,
\end{equation}
for an increasing compact exhaustion \(K_m\Subset\Omega\).  Then every
ordinary zero-expenditure stationary law is supported on spatially constant
states.  If the generated gauge quotient sends spatial constants to
\(\Null_\Sigma\), its ordinary equality contribution to
\(\mathfrak H_{\rm succ}^{\Sigma}\) vanishes.  Every nonordinary value of the
gradient, exhaustion, connected-core, or gauge relation remains a
witness-successor and is reduced by \(\widehat{\mathcal R}_\eta\).
\end{corollary}

\begin{proof}
For a zero-expenditure stationary law \(\boldsymbol\nu\), nonnegativity and
the saturation-equation compiler give
\(\int D_m(V)\,\dd\boldsymbol\nu=0\) for every \(m\).  Hence
\(D_m(V)=0\) for \(\boldsymbol\nu\)-almost every state, simultaneously for
all \(m\).  The exhaustion gives \(\nabla V=0\) almost everywhere on
\(\Omega\).  The weak fundamental theorem on connected domains makes \(V\)
spatially constant.  The target projection then vanishes by the stated gauge
quotient.  Totalization of each relation gives the last assertion.
\end{proof}

\section{Covariant action and the physical budget}
\label{sec:covariant-action}

The fixed-Caus formulation separates the current selected by the PDE from
the metric used to measure irreversible expenditure.  This separation is
physical: the action below is the quantity appearing in the energy identity.
The explicit inverse notation in this section is used only when the mobility
has the Hilbert realization stated next.  In a general Gelfand triple, the
same identities are read at the level of the closed mobility form.
These action identities remain library relations until a shell covector
selects an active source cell.  Their reachability content is exactly the
target-local five-source factorization of
\cref{thm:target-local-fact-factorization}; only the resulting
carrier-dominated production may enter the physical ledger.

\begin{definition}[Mobility current space]
\label{def:mobility-current-space}
Fix a state \(u\), a real Hilbert tangent space \(\Hspace_u\), and a bounded
self-adjoint nonnegative mobility
\(\mathbb K(u)\in\mathcal L(\Hspace_u)\).  Let
\(\mathbb K(u)^{-1/2}\) denote the Moore--Penrose inverse on
\((\ker\mathbb K(u))^\perp\).  The mobility current space is the
Cameron--Martin space
\[
 \mathfrak J_u:=\Ran\mathbb K(u)^{1/2}
\]
with inner product and norm
\[
 (J_1,J_2)_{u,*}
 :=\bigl\langle\mathbb K(u)^{-1/2}J_1,
                  \mathbb K(u)^{-1/2}J_2\bigr\rangle_{\Hspace_u},
 \qquad
 \|J\|_{u,*}^2
 :=\bigl\|\mathbb K(u)^{-1/2}J\bigr\|_{\Hspace_u}^2.
\]
Equivalently, this is the completion of \(\Ran\mathbb K(u)\) in the
displayed norm.  Covectors are represented in \(\Hspace_u\), with components
in \(\ker\mathbb K(u)\) identified when they act on currents.
For a velocity \(v\), put
\[
 W(u,v):=v+A_{\Caus}(u)-f_{\Bdry}(u).
\]
The dissipation Lagrangian is
\begin{equation}
 \label{eq:dissipation-lagrangian}
 \mathscr L(u,v)
 :=\begin{cases}
 \frac12\|W(u,v)\|_{u,*}^2,&W(u,v)\in\mathfrak J_u,\\
 +\infty,&\text{otherwise}.
 \end{cases}
\end{equation}
Its physical action on an interval \(I\) is
\(\Action_I(u)=\int_I2\mathscr L(u,\dot u)\dd t\).
\end{definition}

\begin{definition}[Covariant dissipation Hamiltonian]
\label{def:covariant-dissipation-hamiltonian}
For \(p\in\Hspace_u\), define
\begin{equation}
 \label{eq:covariant-dissipation-hamiltonian}
 \mathscr H_{\mathrm{cov}}(u,p)
 :=-\langle p,A_{\Caus}(u)-f_{\Bdry}(u)\rangle
   +\frac12\|\mathbb K(u)^{1/2}p\|^2.
\end{equation}
This is the Legendre--Fenchel dual of
\(v\mapsto\mathscr L(u,v)\).
\end{definition}

\begin{proposition}[Fenchel equality and the selected current]
\label{prop:fenchel-selected-current}
Whenever the displayed quantities are finite,
\begin{equation}
 \label{eq:fenchel-current-square}
 \mathscr L(u,v)+\mathscr H_{\mathrm{cov}}(u,p)-\langle p,v\rangle
 =\frac12\|W(u,v)-\mathbb K(u)p\|_{u,*}^2\ge0.
\end{equation}
Equality holds precisely when
\(W(u,v)=\mathbb K(u)p\) in the current quotient.  For the fixed-Caus
gradient flow \eqref{eq:fixed-caus-flow}, the selected velocity has
\[
 W(u,\dot u)=-\mathbb K(u)D\Energy(u),
 \]
so equality holds at \(p=-D\Energy(u)\), modulo the mobility kernel.
\end{proposition}

\begin{proof}
Substitute \eqref{eq:dissipation-lagrangian} and
\eqref{eq:covariant-dissipation-hamiltonian}; since
\(v=W-A_{\Caus}+f_{\Bdry}\), the linear terms combine to
\(-\langle p,W\rangle\).  Completing the square in the
\(\mathbb K^{-1}\)-metric gives \eqref{eq:fenchel-current-square}.  The
equality and selected-current statements follow directly.
\end{proof}

\begin{definition}[Physical carrier balance]
\label{def:physical-carrier-balance}
Let \(u:[0,T_*)\to X\) be an ordinary maximal branch.  Expand the native
signed carrier graph of \cref{def:physical-budget} and take the Jordan
decomposition of its signed currents.  Denote the resulting scalar by
\(\mathcal K_u\), the gross positive production by \(\mathfrak D_u\), the
returned internal feeding by \(\Phi_u^+\), and true positive exterior input
by \(R_u^+\).  Its ordinary graph value is, for \(0\le s<t<T_*\),
\begin{equation}
 \label{eq:physical-carrier-balance}
 \mathcal K_u(t)+\mathfrak D_u((s,t])
 \le
 \mathcal K_u(s)+\Phi_u^+((s,t])+R_u^+((s,t]).
\end{equation}
Here all four quantities are generated from the one signed identity.  The
generated passivity radius is the least extended-real \(\beta\) for which the
measure domination
\begin{equation}
 \label{eq:physical-passivity-domination}
 \Phi_u^+\le\beta\,\mathfrak D_u.
\end{equation}
holds.  Whether \(\beta\le1\), whether \(\mathcal K_u\) has a lower bound,
and whether the exterior input is finite are values of generated relations,
not additional data.  All terms are expressed in the physical clock.  A normalized-clock identity
is routed through the totalized clock-comparison relation in
\cref{prop:physical-clock-comparison}.
\end{definition}

\begin{theorem}[Carrier compiler on one admitted stratum]
\label{thm:carrier-physical-budget}
If the generated values satisfy \(\beta\le1\),
\(\inf_t\mathcal K_u(t)=K_->-\infty\), and
\(R_u^+([0,T_*))<\infty\), then
\begin{equation}
 \label{eq:carrier-expenditure-measure}
 \mu_u:=\mathfrak D_u-\Phi_u^+
\end{equation}
is a nonnegative countably additive physical expenditure measure and
\begin{equation}
 \label{eq:carrier-physical-budget}
 \mu_u([0,T_*))
 \le
 \mathcal K_u(0)-K_-+R_u^+([0,T_*))
 =:\Budget<\infty.
\end{equation}
If \(\beta<1\), then
\begin{equation}
 \label{eq:strict-passivity-coercivity}
 \mu_u\ge(1-\beta)\mathfrak D_u.
\end{equation}
For \(\beta=1\), \(\mu_u\) is the exact flux-deficit measure.  These are
values checked on the carrier-tagged stratum, not hypotheses supplied with
the datum.  If any displayed
relation has the opposite value, its total-graph witness is sent directly to
the activity lift and successor recursion.  No alternative expenditure or
independent budget is appended later.  This theorem constructs the available
orbitwise physical measure.  A target event spends that measure only after
its source-resolved shell current satisfies the local charging rule
\eqref{eq:target-aligned-charge-rule}.
\end{theorem}

\begin{proof}
The domination \(\Phi_u^+\le\mathfrak D_u\) makes the difference in
\eqref{eq:carrier-expenditure-measure} a nonnegative Radon measure.  Rearrange
\eqref{eq:physical-carrier-balance} with \(s=0\), use
\(\mathcal K_u(t)\ge K_-\), and let \(t\uparrow T_*\).  Monotone convergence
gives \eqref{eq:carrier-physical-budget}.  Finally,
\(\Phi_u^+\le\beta\mathfrak D_u\) gives
\eqref{eq:strict-passivity-coercivity}.
\end{proof}

\begin{definition}[Presentation-generated ordinary datum domain]
\label{def:declared-global-analysis-datum-class}
Fix an accepted effective presentation
\(\mathfrak P_{\rm eff}\).  Its presentation-generated ordinary datum domain
\(\mathscr D_{\rm eng}(\mathfrak P_{\rm eff})\) consists of the data \(u_0\)
which occur as initial traces of ordinary elements of the target-blind public
local-solution graph on some rational interval
\([0,\tau]\), \(\tau>0\).  That graph is generated from
\(\mathsf{Pre}(\mathsf{Eq})\) and includes the weak equation, boundary rows,
and declared ordinary coordinates.  Membership is therefore the source
projection of a native public graph, not a certificate carried by the datum
and not a target-relative admission test.  The native carrier graph, its lower-bound graph, exterior-input
graph, passivity graph, and all of their totalized failure values remain
generated coordinates of the accepted public presentation.  Finite budget,
passivity, and carrier coercivity are compiler outputs, rather than datum
admission certificates.
\end{definition}

\begin{theorem}[Ordinary-fiber existence and exhaustive target-aligned carrier stratification]
\label{thm:ordinary-fiber-global-finite-stratum-coverage}
Let \(u_0\in\mathscr D_{\rm eng}(\mathfrak P_{\rm eff})\).  Then the ordinary
local-solution fiber over \(u_0\) is nonempty.  On every ordinary maximal
branch \(u\) from \(u_0\), including every branch when uniqueness fails, and
for every candidate enumerated by the target-blind public multiplier grammar,
the expanded carrier graph and every target-aligned owner generated from it
have exactly one of the following dispositions:
\begin{enumerate}[label=\textup{(F\arabic*)}]
\item the generated passivity, lower-carrier, and exterior-input rows produce
  one carrier-tagged nonnegative physical measure
  \(\mu_u=\mathfrak D_u-\Phi_u^+\) and a finite public ceiling
  \begin{equation}
   \label{eq:global-finite-stratum-coverage}
   \mu_u(\mathfrak Q_u^{\rm aff})
   \le \mu_u([0,T_*(u)))
   \le \mathcal K_u(0)-K_-(u)+R_u^+([0,T_*(u)))
   =:\Budget(u)<\infty,
  \end{equation}
  where \(\mathfrak Q_u^{\rm aff}\) consists precisely of the pre-contact
  events whose source-resolved owners satisfy the local target-aligned
  charging rule;
\item the complete cumulative same-origin target projection of the failed
  carrier rows vanishes;
\item a least nonzero target-aligned failure owner is generated, with its
  canonical shell covector, exact five-source provenance, origin, event
  index, joint prefix, and carrier-polar disposition.  It is the next typed
  structural successor.
\end{enumerate}
The three cells are exhaustive for every enumerated carrier.  Thus every
target-bearing failure of
passivity, carrier lower boundedness, exterior-input finiteness, carrier
coherence, or algebraic domination is processed by the same public
target-aligned successor algebra as every other possible cause of contact.
No such value is an admission failure, an unpriced residual, or a regularity
conclusion.
\end{theorem}

\begin{proof}
By definition of \(\mathscr D_{\rm eng}(\mathfrak P_{\rm eff})\), \(u_0\)
is the initial trace of an ordinary element of the native target-blind
local-solution graph.  Extend that element to a maximal branch in the same
public graph; if extensions are
nonunique, retain every maximal extension.  Fix any such branch.  The
public carrier compiler of \cref{thm:public-carrier-admission} expands the
enumerated signed identity, joins all
same-origin inputs, cancels internal currents, and totalizes the passivity,
lower-bound, exterior-input, and coherence graphs.  If all their ordinary
finite values occur, \cref{thm:carrier-physical-budget} gives the tagged
measure and
ceiling in \textup{(F1)}.  The local charging rule restricts its use to
\(\mathfrak Q_u^{\rm aff}\), and countable additivity gives
\eqref{eq:global-finite-stratum-coverage}.

Otherwise retain the complete expanded derivation of the first failed graph
coordinate and form its cumulative joint record before target projection.
The universal five-source provenance theorem assigns every nonzero projected
owner its exact ancestry.  A zero projection gives \textup{(F2)}.  On the
complement, universal physical-quantity ownership and recursive
carrier--polar exhaustion give algebraic carrier admission, a cone-boundary
successor, or a finite-cylinder polar successor.  An algebraically admitted
finite part joins the affordable ledger; every other nonzero value is the
typed successor in \textup{(F3)}.  The continuum composition no-creation
theorem exposes any nonzero cumulative target component at a finite joint
prefix.  Hence no limit value supplies a fourth cell.  The construction is
uniform over the total solution fiber and therefore applies to every maximal
branch.
\end{proof}

\begin{definition}[Reduced passivity radius]
\label{def:reduced-passivity-radius}
Let \(\mathcal P_{\mathrm{car}}^{\mathrm{red}}\) be a normalized carrier
profile space obtained after zero-descent-defect source reductions.  Let
\(D(V)\ge0\) and \(\Phi_+(V)\ge0\) be the descended gross-expenditure and
feeding-flux densities in the same physical units.  The reduced passivity
radius is
\begin{equation}
 \label{eq:reduced-passivity-radius}
 \beta_*^{\mathrm{red}}
 :=\alpha_*(\Phi_+\mid D)
 =\sup_{\substack{V\in\mathcal P_{\mathrm{car}}^{\mathrm{red}}\\D(V)>0}}
   \frac{\Phi_+(V)}{D(V)},
\end{equation}
with the null-mode convention of
\cref{def:structural-evaluation-invariants}.  A bound
\(\beta_*^{\mathrm{red}}\le1\) supplies
\eqref{eq:physical-passivity-domination} whenever the carrier identity is
local in the normalized profiles and its physical-clock pullback is proved.
\end{definition}

\begin{proposition}[Quadratic passivity is a spectral invariant]
\label{prop:quadratic-passivity-radius}
On the quadratic-form realization branch of the reduced Hilbert profile space,
\[
 D(V)=\langle AV,V\rangle,
 \qquad
 \Phi_+(V)=\max\{\langle BV,V\rangle,0\},
\]
where \(A\ge0\), \(B=B^*\), all \(D\)-null target-visible feeding modes have
been separated, and \(B\) is bounded in the \(A\)-form norm.  Let
\(S_{\mathrm{flux}}=A^{-1/2}BA^{-1/2}\) be the descended self-adjoint form
operator on the completion modulo \(\ker A\).  Then
\begin{equation}
 \label{eq:quadratic-passivity-spectrum}
 \beta_*^{\mathrm{red}}
 =\max\{\sup\sigma(S_{\mathrm{flux}}),0\}.
\end{equation}
Thus every multiplier or casewise proof of \(\Phi_+\le D\) is, on this
reduced quadratic interface, an upper bound for one spectral radius.
\end{proposition}

\begin{proof}
The \(A\)-form completion is a Hilbert space after quotienting by
\(\ker A\).  The form induced by \(B\) is represented there by
\(S_{\mathrm{flux}}\).  Rayleigh--Ritz identifies the supremum of
\(\langle BV,V\rangle/D(V)\) with
\(\sup\sigma(S_{\mathrm{flux}})\); taking the positive part gives
\eqref{eq:quadratic-passivity-spectrum}.
\end{proof}

\begin{figure}[tbp]
\centering
\resizebox{.98\linewidth}{!}{%
\begin{tikzpicture}[fw diagram,node distance=8mm and 10mm]
  \node[fw geom,text width=2.7cm] (gradient)
    {fixed-Caus energy\\and action identity};
  \node[fw provenance,text width=2.7cm,below=8mm of gradient] (hypo)
    {twisted Geom--Caus\\commutator carrier};
  \node[fw source,text width=2.7cm,below=8mm of hypo] (mourre)
    {positive-commutator\\virial carrier};
  \node[fw use,text width=3.0cm,right=15mm of hypo] (balance)
    {one carrier balance\\\(\mathcal K+\mathfrak D\le
      \mathcal K_0+\Phi^++R^+\)};
  \node[fw provenance,text width=2.8cm,right=of balance] (radius)
    {reduced passivity radius\\\(\beta_*^{\mathrm{red}}\le1\)};
  \node[fw terminal,text width=3.0cm,right=of radius] (budget)
    {single physical measure\\\(\mu_u=\mathfrak D-\Phi^+\le\Budget\)};
  \draw[fw arrow] (gradient)--(balance);
  \draw[fw arrow] (hypo)--(balance);
  \draw[fw arrow] (mourre)--(balance);
  \draw[fw arrow] (balance)--(radius);
  \draw[fw arrow] (radius)--(budget);
\end{tikzpicture}%
}
\caption{Carrier-to-budget compiler.  Gradient-flow, hypocoercive, and
positive-commutator identities are different ways to expose the same
structural object: a lower-bounded carrier whose feeding flux is dominated by
gross expenditure.  The output is the sole physical budget measure.}
\label{fig:carrier-physical-budget}
\end{figure}

\section{Hamiltonian meanings and chart covariance}
\label{sec:hamiltonian-meanings}

The constructions in this section are universal operator identities, not
equation-wide regularity assertions.  In a reachability argument they are
activated only after restriction to a contact-generated stratum
\(X_{\mathfrak b}\), where
\cref{thm:target-local-fact-factorization} separates their
\(\Geom/\Caus/\Abs/\Lift/\Bdry\) components, target-null part, carrier
coboundary, and total-graph defects.  Only that factorized local value may
enter a physical charge or close a branch.

\begin{definition}[Geometric spectral Hamiltonian]
\label{def:geometric-hamiltonian}
Let \(\mathfrak q\) be a densely defined closed nonnegative symmetric form on
a real Hilbert space \(\Hspace\).  Its geometric Hamiltonian is the unique
self-adjoint operator \(H\ge0\) satisfying
\[
 \mathfrak q(u,v)=\langle H^{1/2}u,H^{1/2}v\rangle,
 \qquad u,v\in D(\mathfrak q).
\]
This Hamiltonian records only the \(\Geom\) form; the fixed antisymmetric
current remains in \(\Caus\).
\end{definition}

\begin{proposition}[Separation of Hamiltonian meanings]
\label{prop:hamiltonian-separation}
The following constructions require distinct data.
\begin{enumerate}[label=\textup{(\roman*)}]
\item A closed symmetric form determines the spectral Hamiltonian \(H\) of
\cref{def:geometric-hamiltonian}.
\item A canonical conservative Hamiltonian vector field requires a Poisson
tensor \(\mathbb J\), a domain on which the Jacobi identity holds, and a
functional \(\mathcal H\); the field is \(\mathbb J D\mathcal H\).
\item A positive mobility and a covariant connection determine the Fenchel
Hamiltonian \(\mathscr H_{\mathrm{cov}}\) of
\cref{def:covariant-dissipation-hamiltonian}.
\end{enumerate}
An antisymmetric Caus form alone supplies directed transport but does not
provide the Poisson or Jacobi data in item \textup{(ii)}.
\end{proposition}

\begin{proof}
The representation theorem for closed nonnegative forms gives item
\textup{(i)} \citep[Chapter~VI]{Kato1995}.  Item \textup{(ii)} is the
definition of a Poisson Hamiltonian system; antisymmetry is only one of the
Poisson requirements.  The square completion in
\cref{prop:fenchel-selected-current} proves item \textup{(iii)}.
\end{proof}

\begin{figure}[tbp]
\centering
\resizebox{.97\linewidth}{!}{%
\begin{tikzpicture}[fw diagram,node distance=8mm and 11mm]
  \node[fw geom,text width=2.8cm] (form)
    {closed symmetric form\\\(\mathfrak q\)};
  \node[fw terminal,text width=2.9cm,right=of form] (spectral)
    {spectral Hamiltonian\\\(H\ge0\)};
  \node[fw source,text width=2.8cm,below=13mm of form] (poisson)
    {Poisson tensor and\\Jacobi-domain data};
  \node[fw terminal,text width=2.9cm,right=of poisson] (canonical)
    {conservative Hamiltonian\\vector field};
  \node[fw provenance,text width=2.8cm,below=13mm of poisson] (mobility)
    {mobility plus\\covariant connection};
  \node[fw terminal,text width=2.9cm,right=of mobility] (fenchel)
    {dissipation Hamiltonian\\\(\mathscr H_{\rm cov}\)};
  \node[fw residual,text width=3.0cm,right=17mm of canonical] (skew)
    {antisymmetric Caus alone\\directed transport only};
  \draw[fw arrow] (form)--(spectral);
  \draw[fw arrow] (poisson)--(canonical);
  \draw[fw arrow] (mobility)--(fenchel);
  \draw[fw dashed arrow] (skew)--(canonical);
\end{tikzpicture}%
}
\caption{Three noninterchangeable Hamiltonian constructions.  The dashed
line marks a missing implication: skew transport alone does not establish a
Poisson structure.}
\label{fig:three-hamiltonian-meanings}
\end{figure}

\begin{theorem}[Same-current covariance]
\label{thm:same-current-covariance}
Let \(v=\Phi_t(u)\) be a generated \(C^1\) chart relation.  Totalize its
diffeomorphism and tangent-map relations.  On the branch where
\(T_t=D\Phi_t(u):\Hspace_u\to\Hspace_v\) is a bounded linear isomorphism
with bounded inverse, with every expression on the right evaluated at
\(u=\Phi_t^{-1}(v)\), define
\begin{equation}
 \label{eq:transported-covariant-data}
 \widetilde{\mathbb K}=T_t\mathbb KT_t^*,
 \qquad
 \widetilde A_{\Caus}=T_tA_{\Caus},
 \qquad
 A_{\Lift}^{\Phi}=-\partial_t\Phi_t,
 \qquad
 \widetilde f_{\Bdry}=T_tf_{\Bdry}.
\end{equation}
Then
\begin{equation}
 \label{eq:same-current-covariance}
 \dot v+\widetilde A_{\Caus}(v)+A_{\Lift}^{\Phi}(v)
       -\widetilde f_{\Bdry}(v)
 =T_t\bigl(\dot u+A_{\Caus}(u)-f_{\Bdry}(u)\bigr).
\end{equation}
The map \(T_t\) is an isometry between the corresponding mobility-current
spaces equipped with their transported metrics, so the Lagrangian and
physical action are invariant.  The chart has
\(\Lift\) ancestry, and quotienting any declared symmetry additionally has
\(\Abs\) ancestry.
\end{theorem}

\begin{proof}
Differentiate \(v=\Phi_t(u)\) and substitute
\eqref{eq:transported-covariant-data}; the \(\partial_t\Phi_t\) terms cancel,
giving \eqref{eq:same-current-covariance}.  For \(\widetilde J=T_tJ\), the
variational identity
\[
 \|J\|_{\mathbb K^{-1}}^2
 =\sup_{p\in\Hspace_u}
   \{2\langle p,J\rangle-\langle p,\mathbb Kp\rangle\}
\]
and the bijective change of covector
\(p=T_t^*\widetilde p\) give
\(\|\widetilde J\|_{\widetilde{\mathbb K}^{-1}}
  =\|J\|_{\mathbb K^{-1}}\).
Integrating proves action invariance.  The source labels follow from
\cref{def:source-alphabet}.
\end{proof}

\begin{proposition}[Chart covariance of the dissipation Hamiltonian]
\label{prop:covariant-hamiltonian-chart}
On the ordinary covariance-graph branch of
\cref{thm:same-current-covariance}, let
\(p=T_t^*\widetilde p\), and define the transformed Hamiltonian using the
total connection \(\widetilde A_{\Caus}+A_{\Lift}^{\Phi}\).  Then
\begin{equation}
 \label{eq:hamiltonian-chart-shift}
 \widetilde{\mathscr H}_{\mathrm{cov}}(t,v,\widetilde p)
 =\mathscr H_{\mathrm{cov}}(u,p)
  +\langle\widetilde p,\partial_t\Phi_t(u)\rangle.
\end{equation}
The additional term is the time-dependent Lift connection.  It is not a new
physical Caus term.
\end{proposition}

\begin{proof}
Insert \eqref{eq:transported-covariant-data} into
\eqref{eq:covariant-dissipation-hamiltonian}.  The quadratic term is invariant
because \(\widetilde{\mathbb K}=T_t\mathbb KT_t^*\), while
\[
-\langle\widetilde p,
 T_t(A_{\Caus}-f_{\Bdry})-\partial_t\Phi_t\rangle
 =-\langle p,A_{\Caus}-f_{\Bdry}\rangle
 +\langle\widetilde p,\partial_t\Phi_t\rangle.
\]
\end{proof}

\begin{remark}[A change of physical clock is additional data]
\label{rem:chart-clock-separation}
\Cref{thm:same-current-covariance} keeps the physical time variable fixed.
If a blow-up chart also introduces \(\tau=\theta(t)\), the current,
mobility, and action density must be transformed with the exact Jacobian
\(\dd t/\dd\tau\).  Action invariance, or comparison with the original
physical measure, is not automatic; the required one-sided comparison is
the content of \cref{prop:physical-clock-comparison}.
\end{remark}

\begin{figure}[tbp]
\centering
\resizebox{.98\linewidth}{!}{%
\begin{tikzpicture}[fw diagram,node distance=8mm and 10mm]
  \node[fw source,text width=2.5cm] (u) {state \(u\)};
  \node[fw provenance,text width=3.0cm,right=of u] (chart)
    {legal chart\\\(v=\Phi_t(u)\), \(T_t=D\Phi_t\)};
  \node[fw source,text width=2.5cm,right=of chart] (v) {state \(v\)};
  \node[fw geom,text width=2.9cm,below=14mm of u] (current)
    {physical current and\\mobility \(J,\mathbb K\)};
  \node[fw geom,text width=3.0cm,below=14mm of v] (pushed)
    {push-forward current\\\(T_tJ,T_t\mathbb KT_t^*\)};
  \node[fw use,text width=2.8cm,below=13mm of chart] (action)
    {same action and\\target aiming};
  \node[fw provenance,text width=3.0cm,right=of v] (connection)
    {Lift connection\\\(-\partial_t\Phi_t\)};
  \node[fw terminal,text width=3.0cm,below=of connection] (shift)
    {Hamiltonian shift\\\(+\langle\widetilde p,\partial_t\Phi_t\rangle\)};
  \draw[fw arrow] (u)--(chart);
  \draw[fw arrow] (chart)--(v);
  \draw[fw arrow] (current)--(action);
  \draw[fw arrow] (action)--(pushed);
  \draw[fw arrow] (v)--(pushed);
  \draw[fw dashed arrow,bend left=22] (chart) to (connection);
  \draw[fw arrow] (connection)--(shift);
\end{tikzpicture}%
}
\caption{Same-current chart covariance.  A chart transports the physical
current and metric; its time dependence is a Lift connection with the exact
Hamiltonian shift shown.}
\label{fig:same-current-chart}
\end{figure}

\section{Killed resolvents, resistance, and target realization}
\label{sec:killed-resolvents}

Every resolvent, capacity, and resistance relation below is a library
relation in the sense of
\cref{cor:target-local-identity-activation}.  Its regularity use is confined
to the active target/source cell selected by contact.  A failed domain,
closedness, killing, flux, or range relation is therefore a typed local
successor; it is not an equation-wide obstruction and supplies no physical
charge by itself.

\begin{definition}[Killed geometric Hamiltonian]
\label{def:killed-hamiltonian}
Let \(\mathfrak q\) be a densely defined closed nonnegative symmetric
Markovian form and let \(\Sigma\) be quasi-closed for a regular or
quasi-regular realization.  The killed form is the restriction
\[
 D(\mathfrak q_\Sigma)
 =\{u\in D(\mathfrak q):\widetilde u=0\text{ q.e. on }\Sigma\},
 \qquad
 \mathfrak q_\Sigma=\mathfrak q|_{D(\mathfrak q_\Sigma)},
\]
and its self-adjoint operator, when the restricted form is closed, is
\(H_\Sigma\ge0\).  Adding an admitted nonnegative killing potential instead
defines a \emph{soft}-killed form.  It is a separate construction unless
convergence to the hard restriction is proved.
\end{definition}

\begin{proposition}[Existence of the killed Hamiltonian]
\label{prop:killed-hamiltonian-exists}
If \(D(\mathfrak q_\Sigma)\) is closed in the form norm, then
\(\mathfrak q_\Sigma\) is closed and determines a unique self-adjoint
\(H_\Sigma\ge0\).
\end{proposition}

\begin{proof}
A form-norm closed subspace of the domain of a closed form is complete in that
norm.  The representation theorem for closed nonnegative forms gives
\(H_\Sigma\).
\end{proof}

\begin{proposition}[Committor is the canonical binary-contact quotient]
\label{prop:committor-contact-quotient}
Let \(A\) and \(B\) be disjoint quasi-closed sets for an associated right
process, and let
\[
 h(x):=\mathbb P_x\{\tau_A<\tau_B\}
\]
be its quasi-continuous committor.  Define \(q_h(x)=h(x)\) outside the
exceptional set.  If a quotient \(q:X\to Q\) preserves this binary contact
law in the precise sense that \(h=\bar h\circ q\) quasi-everywhere for some
\(\bar h:Q\to[0,1]\), then
\begin{equation}
 \label{eq:committor-quotient-factorization}
 q_h=\bar h\circ q
 \qquad\text{quasi-everywhere}.
\end{equation}
Hence \(q_h\) is the coarsest quotient preserving the scalar event
\(\{\tau_A<\tau_B\}\): every other such quotient is finer than its level-set
partition.  This assertion concerns the binary hitting probability, not the
full hitting-time distribution.
\end{proposition}

\begin{proof}
The preservation branch is exactly
\eqref{eq:committor-quotient-factorization}.  Therefore
\(q(x)=q(y)\) implies \(h(x)=h(y)\), so every fiber of \(q\) lies in a fiber
of \(q_h\).  Conversely, \(q_h\) itself preserves the event by the defining
identity for \(h\).
\end{proof}

\begin{theorem}[Killed-resolvent stopwatch]
\label{thm:killed-resolvent-stopwatch}
Let \(L_\Sigma\) generate a killed semigroup \(P_t^\Sigma\) with growth
bound \(\omega\).  Let \(\ell_0\) be an initial functional and let
\(k_\Sigma\) be a terminal-flux vector such that
\begin{equation}
 \label{eq:killed-flux-law}
 \mathbb P_{\ell_0}\{\tau_\Sigma\in\dd t,\tau_\Sigma<\infty\}
 =\langle\ell_0,P_t^\Sigma k_\Sigma\rangle\dd t
\end{equation}
as finite measures.  Then, for \(\lambda>\omega\),
\begin{equation}
 \label{eq:killed-arrival-resolvent}
 \mathbb E_{\ell_0}
 \bigl[e^{-\lambda\tau_\Sigma};\tau_\Sigma<\infty\bigr]
 =\langle\ell_0,(\lambda-L_\Sigma)^{-1}k_\Sigma\rangle.
\end{equation}
Without \eqref{eq:killed-flux-law}, the right side remains a cross-resolvent
pairing and has no automatic first-contact interpretation.
\end{theorem}

\begin{proof}
Integrate \eqref{eq:killed-flux-law} against \(e^{-\lambda t}\).  The Bochner
Laplace transform of the semigroup is
\(\int_0^\infty e^{-\lambda t}P_t^\Sigma\dd t
=(\lambda-L_\Sigma)^{-1}\), which proves the identity.
\end{proof}

\begin{theorem}[Resolvent Fenchel dual and resistance]
\label{thm:resolvent-resistance}
Let \(K\ge0\) be self-adjoint with closed form \(\mathfrak q\), and let
\(e\in\Hspace\).  For \(\lambda>0\),
\begin{equation}
 \label{eq:resolvent-fenchel}
 \langle e,(K+\lambda)^{-1}e\rangle
 =\sup_{v\in D(\mathfrak q)}
 \left\{2\operatorname{Re}\langle e,v\rangle
 -\mathfrak q(v,v)-\lambda\|v\|^2\right\}.
\end{equation}
If \(P_{\ker K}e=0\), then monotone convergence gives
\begin{equation}
 \label{eq:resistance-limit}
 \lim_{\lambda\downarrow0}\langle e,(K+\lambda)^{-1}e\rangle
 =\int_{(0,\infty)}s^{-1}\dd\nu_e(s)
 =:\mathcal R_K(e)\in[0,\infty].
\end{equation}
Finiteness of \(\mathcal R_K(e)\) is equivalent to boundedness of
\(v\mapsto\langle e,v\rangle\) on the energy space with norm
\(\mathfrak q(v,v)^{1/2}\); in that case it has a unique weak energy
representer modulo \(\ker K\).
\end{theorem}

\begin{proof}
For \(R_\lambda=(K+\lambda)^{-1}\), completing the square in the form norm
of \(K+\lambda\) shows that the supremum is attained at
\(v=R_\lambda e\) and equals \(\langle e,R_\lambda e\rangle\).
The spectral theorem gives
\(\langle e,R_\lambda e\rangle
=\int(\lambda+s)^{-1}\dd\nu_e(s)\).  Monotone convergence proves
\eqref{eq:resistance-limit}.  The last equivalence is the Riesz theorem on the
completion of \(D(\mathfrak q)/\ker K\).
\end{proof}

\begin{theorem}[Harmonic defect equals resistance error]
\label{thm:harmonic-defect-resistance}
Let \(K\ge0\) be self-adjoint with form \(\mathfrak q\), and let
\[
 \mathcal E_K:=\overline{D(\mathfrak q)/\ker K}^{\mathfrak q}
\]
be its energy space.  For \(\ell\in\mathcal E_K^*\), totalize the weak
solution relation \(\mathfrak q(x,\phi)=\ell(\phi)\) modulo \(\ker K\).  On
its ordinary branch let \(x\) denote the solution and let
\(v\in D(\mathfrak q)\) be a trial state with residual
\begin{equation}
 \label{eq:harmonic-defect-functional}
 e_v(\phi):=\mathfrak q(v,\phi)-\ell(\phi).
\end{equation}
Then
\begin{equation}
 \label{eq:harmonic-defect-resistance}
 \|e_v\|_{\mathcal E_K^*}^2
 =\|[v-x]\|_{\mathcal E_K}^2
 =\mathfrak q(v-x,v-x).
\end{equation}
If \(e_v\) is represented by a vector \(e\perp\ker K\) in \(\Hspace\), then
\begin{equation}
 \label{eq:vector-defect-resistance}
 \|e_v\|_{\mathcal E_K^*}^2
 =\langle e,K^+e\rangle
 =\mathcal R_K(e).
\end{equation}
Thus an approximate committor, elliptic corrector, or resolvent state is
controlled by one resistance invariant of its harmonic defect; no
\(\ker K\) component of the trial state enters the error.  When that kernel
is a closed gauge or target-null space, the corresponding discarded
motion requires no estimate.
\end{theorem}

\begin{proof}
The weak equation and \eqref{eq:harmonic-defect-functional} give
\[
 e_v(\phi)=\mathfrak q(v-x,\phi).
\]
The Riesz map of \(\mathcal E_K\) is an isometry onto its dual, proving
\eqref{eq:harmonic-defect-resistance}.  When the residual is represented by
\(e\), the spectral theorem identifies its squared energy-dual norm with
\(\int_{(0,\infty)}s^{-1}\dd\nu_e(s)\), which is
\eqref{eq:vector-defect-resistance}; the middle pairing is understood as
this spectral quadratic-form value.
\end{proof}

\begin{corollary}[Defect-aligned drift bound]
\label{cor:defect-aligned-drift-bound}
If \(e\perp\ker K\) and \(\mathcal R_K(e)<\infty\), then for every
\(w\in D(\mathfrak q)\),
\begin{equation}
 \label{eq:defect-aligned-drift-bound}
 |\langle e,w\rangle|
 \le\mathcal R_K(e)^{1/2}\mathfrak q(w,w)^{1/2}.
\end{equation}
Only the component of \(w\) paired with the target-selected routable defect
enters a reachability charge, and it does so through the local carrier rule.
\end{corollary}

\begin{proof}
Use the energy representer from
\cref{thm:resolvent-resistance} and apply Cauchy--Schwarz in
\(\mathcal E_K\).
\end{proof}

\begin{proposition}[Routable and zero-mode resistance split]
\label{prop:resistance-zero-mode-split}
For arbitrary \(e\in\Hspace\), write
\[
 e=e_0+e_\perp,
 \qquad
 e_0=P_{\ker K}e,
 \qquad
 e_\perp=P_{(\ker K)^\perp}e.
\]
Then
\begin{equation}
 \label{eq:resistance-zero-mode-split}
 \langle e,(K+\lambda)^{-1}e\rangle
 =\frac{\|e_0\|^2}{\lambda}
  +\int_{(0,\infty)}\frac{1}{s+\lambda}\dd\nu_{e_\perp}(s).
\end{equation}
Consequently the zero-discount resistance is finite precisely when
\(e_0=0\) and
\(\int_{(0,\infty)}s^{-1}\dd\nu_{e_\perp}(s)<\infty\).  A nonzero
target-visible zero mode is an inherited Geom/Caus zero-frequency occurrence,
not a residual coordinate.  It is removed only when it belongs to a declared
gauge or to \(\Null_\Sigma\); otherwise a separate carrier or price
row evaluates it.
\end{proposition}

\begin{proof}
The spectral projection of \(K\) at \(0\) contributes
\(\lambda^{-1}\|e_0\|^2\); the remaining spectral measure is supported on
\((0,\infty)\), which gives \eqref{eq:resistance-zero-mode-split}.  Monotone
convergence yields the finiteness criterion.  The source statement follows
from \cref{thm:source-preserving-completion,thm:no-target-visible-residual}.
\end{proof}

\begin{proposition}[Capacity and committor realization]
\label{prop:capacity-committor-realization}
Let \(A\) and \(B\) be disjoint quasi-closed sets for a regular or
quasi-regular Dirichlet form.  The condenser capacity
\[
 \Capacity(A,B)
 =\inf\{\mathfrak q(v,v):v\in D(\mathfrak q),\ \widetilde v=1
 \text{ on }A,\ \widetilde v=0\text{ on }B\}
\]
is an analytic Geom invariant.  On the zero realization-defect branch, if the minimizer is unique and the
associated right process realizes the weak Dirichlet problem, then its
quasi-continuous version is the committor
\(x\mapsto\mathbb P_x\{\tau_A<\tau_B\}\) quasi-everywhere.  The probabilistic
conclusion is the ordinary value of the right-process realization row;
positive capacity alone is only its prospective value
\citep{FukushimaOshimaTakeda2011,MaRoeckner1992}.
\end{proposition}

\begin{theorem}[Polar-target capacity criterion]
\label{thm:zero-capacity-polarity}
Let \(\mathfrak q\) be a symmetric quasi-regular Dirichlet form with
associated right process, and let \(\Sigma\) be quasi-closed.  Write
\[
 \Capacity_{\mathfrak q}(\Sigma)
 :=\inf\{\mathfrak q(v,v)+\|v\|_2^2:
      v\in D(\mathfrak q),\ \widetilde v\ge1
      \text{ q.e. on a neighborhood of }\Sigma\}.
\]
If \(\Capacity_{\mathfrak q}(\Sigma)=0\), then \(\Sigma\) is polar:
\begin{equation}
 \label{eq:zero-capacity-polarity}
 \mathbb P_x\{\tau_\Sigma<\infty\}=0
 \qquad\text{for quasi-every }x.
\end{equation}
This is a direct Geom nonreachability criterion for the Markov realization
and requires no cascade argument.  Positive capacity selects the complementary row
but supplies no hitting lower bound without a starting-law or nonpolarity
statement.
\end{theorem}

\begin{proof}
For a quasi-regular Dirichlet form, sets of zero \(1\)-capacity are properly
exceptional up to a capacity-zero enlargement.  The associated right process
avoids every properly exceptional set for quasi-every starting point, which
gives \eqref{eq:zero-capacity-polarity}
\citep{FukushimaOshimaTakeda2011,MaRoeckner1992}.
\end{proof}

\begin{proposition}[Quotient pullback of zero capacity]
\label{prop:quotient-capacity-pullback}
Let \(q:X\to Q\) identify a target
\(\Sigma=q^{-1}(\Sigma_Q)\) quasi-everywhere, and let
\(\mathfrak q_X,\mathfrak q_Q\) be symmetric quasi-regular Dirichlet forms.
Totalize the lift relation which, for every admissible quotient potential \(f\)
with
\(\widetilde f\ge1\) near \(\Sigma_Q\), there is a lifted potential \(\psi_f\)
with \(\widetilde\psi_f\ge1\) near \(\Sigma\) and
\begin{equation}
 \label{eq:quotient-capacity-lift}
 (\mathfrak q_X)_1(\psi_f,\psi_f)
 \le C(\mathfrak q_Q)_1(f,f)+\varepsilon(f).
\end{equation}
If a minimizing family for \(\Capacity_{\mathfrak q_Q}(\Sigma_Q)\) can be
chosen with \(\varepsilon(f)\to0\), then
\begin{equation}
 \label{eq:quotient-capacity-bound}
 \Capacity_{\mathfrak q_X}(\Sigma)
 \le C\Capacity_{\mathfrak q_Q}(\Sigma_Q).
\end{equation}
In particular, zero quotient capacity pulls back to zero capacity on \(X\).
\end{proposition}

\begin{proof}
Every \(\psi_f\) is an admissible competitor for the capacity of \(\Sigma\).
Apply \eqref{eq:quotient-capacity-lift} along the stated minimizing family and
take the liminf.
\end{proof}

\begin{figure}[tbp]
\centering
\resizebox{.98\linewidth}{!}{%
\begin{tikzpicture}[fw diagram,node distance=8mm and 10mm]
  \node[fw source,text width=2.7cm] (initial)
    {initial functional\\\(\ell_0\)};
  \node[fw geom,text width=2.8cm,right=of initial] (killed)
    {killed semigroup\\\(P_t^\Sigma\)};
  \node[fw terminal,text width=2.7cm,right=of killed] (flux)
    {terminal flux\\\(k_\Sigma\)};
  \node[fw provenance,text width=2.8cm,below=13mm of killed] (laplace)
    {Laplace transform\\cross resolvent};
  \node[fw use,text width=3.0cm,right=of laplace] (fenchel)
    {diagonal specialization\\Fenchel dual};
  \node[fw geom,text width=3.0cm,right=of fenchel] (resistance)
    {zero-discount limit\\effective resistance};
  \draw[fw arrow] (initial)--(killed);
  \draw[fw arrow] (killed)--(flux);
  \draw[fw arrow] (killed)--(laplace);
  \draw[fw arrow] (flux)--(laplace);
  \draw[fw dashed arrow] (laplace)--(fenchel);
  \draw[fw arrow] (fenchel)--(resistance);
\end{tikzpicture}%
}
\caption{Killed arrival and resistance are related but distinct.  Arrival is
a cross pairing that requires a flux realization.  The Fenchel and resistance
identities concern a diagonal self-adjoint specialization.}
\label{fig:killed-resolvent-resistance}
\end{figure}

\section{Full-generator response and physical aiming}
\label{sec:full-generator-aiming}

Full-generator formulas enter the reachability calculus only through the
same target-local activation rule.  The input and output projections below
are those generated by the active branch signature, and every differentiated
or factorized term retains its five-source ancestry.  Amplification outside
that target-cyclic cell is transverse to the branch under evaluation.

\begin{definition}[Target-active full resolvent]
\label{def:target-active-resolvent}
Let \(L=-H+A_{\Caus}+B_{\Bdry}\) be a closed generator with a declared
resolvent point \(\lambda\).  For closed mathematical input and target-output
subspaces \(\mathscr I\) and \(\mathscr O_\Sigma\), with bounded projections
\(P_{\mathscr I}\) and \(P_{\mathscr O}\), define
\begin{equation}
 \label{eq:target-active-full-resolvent}
 \Psi_\Sigma(\lambda;L)
 :=\bigl\|P_{\mathscr O}(\lambda-L)^{-1}P_{\mathscr I}\bigr\|.
\end{equation}
This quantity measures nonnormal amplification visible to the fixed target
interface.  It evaluates the existing Geom--Caus--Bdry generator,
not an additional source.
\end{definition}

\begin{proposition}[Nonnormal source invariance]
\label{prop:nonnormal-source-invariance}
Let \(L_0\) be closed on a Banach space and let every \(K^c\) be bounded.
Put \(L_\theta=L_0+\sum_c\theta_cK^c\) on \(D(L_0)\).  On the branch
\(\lambda\in\rho(L_\theta)\) locally in \(\theta\), this is the open
ordinary value of the totalized resolvent relation; on it the resolvent is
norm differentiable and
\begin{equation}
 \label{eq:resolvent-source-derivative}
 \partial_{\theta_c}(\lambda-L_\theta)^{-1}
 =(\lambda-L_\theta)^{-1}K^c(\lambda-L_\theta)^{-1}.
\end{equation}
Thus every first variation of target-active response is factored through the
same source term \(K^c\).  Large pseudospectral response can expose a large
Geom--Caus interaction, but it creates no new nonnormal source.
\end{proposition}

\begin{proof}
Differentiate
\((\lambda-L_\theta)(\lambda-L_\theta)^{-1}=I\) on the common domain and
solve for the derivative.  Bounded input/output compression preserves the
factorization.
\end{proof}

\begin{figure}[tbp]
\centering
\resizebox{.97\linewidth}{!}{%
\begin{tikzpicture}[fw diagram,node distance=8mm and 10mm]
  \node[fw geom,text width=2.5cm] (h)
    {Geom part\\\(-H\)};
  \node[fw use,text width=2.5cm,below=12mm of h] (a)
    {fixed Caus part\\\(A_{\Caus}\)};
  \node[fw box,text width=3.0cm,right=16mm of h,yshift=-6mm] (full)
    {full nonnormal resolvent\\\((\lambda-L)^{-1}\)};
  \node[fw source,text width=2.7cm,right=of full,yshift=10mm] (input)
    {input space\\\(\mathscr I\)};
  \node[fw terminal,text width=2.7cm,right=of full,yshift=-10mm] (output)
    {target readout\\\(\mathscr O_\Sigma\)};
  \node[fw provenance,text width=3.0cm,right=18mm of full] (response)
    {target-active response\\\(\Psi_\Sigma(\lambda;L)\)};
  \node[fw residual,text width=3.0cm,right=of response] (meaning)
    {bound on existing cells\\no new source label};
  \draw[fw arrow] (h)--(full);
  \draw[fw arrow] (a)--(full);
  \draw[fw arrow] (input)--(response);
  \draw[fw arrow] (full)--(response);
  \draw[fw arrow] (output)--(response);
  \draw[fw arrow] (response)--(meaning);
\end{tikzpicture}%
}
\caption{Full nonnormal response retains both Geom and fixed Caus.  Input and
target compression select a response coefficient; they do not change the
operator's source ancestry.}
\label{fig:full-nonnormal-response}
\end{figure}

\begin{definition}[Target-gradient ray]
\label{def:target-gradient-ray}
Let \(F\) be a differentiable target shell.  Totalize membership of its mobility gradient
\[
 G_F:=\mathbb K(u)DF(u)
\]
in \(\mathfrak J_u\).  On the membership branch, for a physical current \(J\in\mathfrak J_u\),
define its positive target power and aiming contribution by
\begin{equation}
 \label{eq:target-power-alignment}
 P_F(J):=(J,G_F)_{u,*},
 \qquad
 \operatorname{Align}_F(J)
 :=\frac{(P_F(J)_+)^2}{\|G_F\|_{u,*}^2},
\end{equation}
with value zero when \(G_F=0\).  By the Cameron--Martin pairing,
\[
 P_F(J)=\langle DF(u),J\rangle
\]
for every \(J\in\mathfrak J_u\); components of \(DF\) in
\(\ker\mathbb K(u)\) annihilate that current space.
\end{definition}

\begin{theorem}[Projection onto the target ray]
\label{thm:target-ray-projection}
For every \(J,G\in\mathfrak J_u\),
\begin{equation}
 \label{eq:target-ray-projection}
 \|J\|_{u,*}^2
 =\frac{\bigl((J,G)_{u,*}\bigr)_+^2}{\|G\|_{u,*}^2}
 +\dist_{u,*}^2(J,\mathbb R_+G),
\end{equation}
with the first term zero when \(G=0\).  Hence
\begin{equation}
 \label{eq:pointwise-action-aiming}
 (P_F(J)_+)^2\le\|J\|_{u,*}^2\|G_F\|_{u,*}^2.
\end{equation}
\end{theorem}

\begin{proof}
The orthogonal projection of \(J\) onto the closed ray \(\mathbb R_+G\) has
coefficient \(\bigl((J,G)_{u,*}\bigr)_+/\|G\|_{u,*}^2\).  Pythagoras gives
\eqref{eq:target-ray-projection}; dropping the nonnegative distance gives
\eqref{eq:pointwise-action-aiming}.
\end{proof}

\begin{proposition}[Integrated action--aiming bound]
\label{prop:action-aiming-bound}
Let \(I\) be a pre-contact interval.  On the finite ordinary branch of the
two totalized squared-norm integrals below,
\begin{equation}
 \label{eq:integrated-action-aiming}
 \int_I P_F(J(t))_+\dd t
 \le
 \left(\int_I\|J(t)\|_{u(t),*}^2\dd t\right)^{1/2}
 \left(\int_I\|G_F(u(t))\|_{u(t),*}^2\dd t\right)^{1/2}.
\end{equation}
For the selected Geom current of \cref{def:fixed-caus-presentation}, the
square of the first factor is the physical action and is bounded by the
right-hand side of \eqref{eq:fixed-caus-action-bound}.
\end{proposition}

\begin{proof}
Apply \eqref{eq:pointwise-action-aiming} and then Cauchy--Schwarz in time.
The final assertion is \cref{prop:fixed-caus-expenditure-bound}.
\end{proof}

\begin{figure}[tbp]
\centering
\resizebox{.97\linewidth}{!}{%
\begin{tikzpicture}[fw diagram,node distance=8mm and 10mm]
  \node[fw source,text width=2.7cm] (current)
    {physical current\\\(J\)};
  \node[fw provenance,text width=2.8cm,below=12mm of current] (shell)
    {mathematical target shell\\\(F\)};
  \node[fw geom,text width=2.8cm,right=of shell] (gradient)
    {mobility gradient ray\\\(\mathbb R_+G_F\)};
  \node[fw box,text width=3.0cm,right=17mm of current] (projection)
    {orthogonal projection of \(J\)\\onto \(\mathbb R_+G_F\)};
  \node[fw use,text width=2.8cm,above right=9mm and 12mm of projection] (aligned)
    {target-aligned\\action};
  \node[fw residual,text width=2.8cm,below right=9mm and 12mm of projection] (defect)
    {transverse ray\\defect};
  \node[fw terminal,text width=3.0cm,right=38mm of projection] (identity)
    {physical action\\= aligned + transverse};
  \draw[fw arrow] (current)--(projection);
  \draw[fw arrow] (shell)--(gradient);
  \draw[fw arrow] (gradient)--(projection);
  \draw[fw arrow] (projection)--(aligned);
  \draw[fw arrow] (projection)--(defect);
  \draw[fw arrow] (aligned)--(identity);
  \draw[fw arrow] (defect)--(identity);
\end{tikzpicture}%
}
\caption{Physical aiming geometry.  The target shell is a mathematical test
of the already selected current.  Projection splits the physical action into
target-aligned and transverse expenditure.}
\label{fig:physical-aiming}
\end{figure}

\section{Degenerate geometry and commutator-generated carriers}
\label{sec:degenerate-geom-routing}

Direct resistance controls only the positive spectral subspace of the Geom
operator.  Fixed Caus transport may nevertheless move a Geom-null direction
into a dissipative direction.  The correct reduced invariant is then a
positive commutator Gramian together with a carrier identity realizing that
Gramian along the actual evolution.
These relations are evaluated only on the contact-generated
\(\Geom\)--\(\Caus\) cell.  A commutator square becomes physical expenditure
only after its target-local factorization proves domination by the public
carrier; every complementary graph value remains a same-origin successor.

\begin{definition}[Degenerate Geom--Caus core]
\label{def:degenerate-geom-caus-core}
Let \(\Hspace\) be a complex Hilbert space and
\[
 L=-S+A,
 \qquad S=S^*\ge0,
 \qquad A^*=-A,
\]
on a common invariant core \(\mathcal D\).  Let \(\mathcal G\subseteq\ker S\)
be a closed \(A\)-invariant gauge subspace, and write
\(P_{\mathcal G}^\perp\) for its orthogonal complement.  Put
\(C_0=S^{1/2}\).  Totalize the iterated commutator relations on
\(\mathcal D\); their ordinary graph branch sets
\[
 C_{k+1}=[A,C_k]\quad\text{on }\mathcal D.
\]
For fixed \(K\) and \(c_k>0\), totalize the quadratic form
\begin{equation}
 \label{eq:commutator-gramian-form}
 q_K(v):=\sum_{k=0}^{K}c_k\|C_kv\|^2.
\end{equation}
On its closable graph branch its positive self-adjoint operator is denoted by \(Q_K\), and
its reduced bracket gap is
\begin{equation}
 \label{eq:reduced-bracket-gap}
 \gamma_K
 :=\inf_{\substack{v\in D(q_K)\cap\mathcal G^\perp\\\|v\|=1}}q_K(v).
\end{equation}
Failure of a domain, closure, or gauge relation is its totalized graph defect;
algebraic commutator span alone does not select the closing branch.
\end{definition}

\begin{proposition}[Bracket resistance invariant]
\label{prop:bracket-resistance-invariant}
If \(\gamma_K>0\), then
\(Q_K\ge\gamma_K I\) on \(\mathcal G^\perp\), and for every
\(e\in\mathcal G^\perp\),
\begin{equation}
 \label{eq:bracket-resistance-bound}
 \mathcal R_K^{\mathrm{br}}(e)
 :=\langle e,Q_K^{-1}e\rangle
 \le\gamma_K^{-1}\|e\|^2.
\end{equation}
The quotient kernel
\[
 \mathfrak h_K:=\mathcal G^\perp\cap\ker Q_K
\]
is the structural rank deficit when \(\gamma_K=0\).  Every target-visible
element of \(\mathfrak h_K\) retains its Geom--Caus ancestry and remains an
existing source occurrence for the routing, carrier, and price rows.
\end{proposition}

\begin{proof}
The first assertion is the definition of the form lower bound in
\eqref{eq:reduced-bracket-gap}; functional calculus gives
\eqref{eq:bracket-resistance-bound}.  The kernel is the intersection of the
kernels of the closed commutator forms in
\eqref{eq:commutator-gramian-form}.  Source preservation and the target-null
residual theorem give the final assertion.
\end{proof}

\begin{definition}[Twisted-carrier realization evaluator]
\label{def:twisted-carrier-invariant}
A bracket Gramian \(Q_K\) has a twisted-carrier realization when there is a
bounded self-adjoint operator \(M\) on \(\mathcal G^\perp\), constants
\(0<m_-\le m_+<\infty\) and \(c_M>0\), and a common evolution core on which
\begin{align}
 m_-\|v\|^2&\le\langle Mv,v\rangle\le m_+\|v\|^2,
 \label{eq:twisted-carrier-equivalence}\\
 2\operatorname{Re}\langle Mv,Lv\rangle
 &\le-c_M q_K(v).
 \label{eq:twisted-carrier-lyapunov}
\end{align}
For a forced equation, the positive part of
\(2\operatorname{Re}\langle Mu,f_{\Bdry}\rangle\dd t\) is the candidate
\(\Bdry\)-input current.  It enters \(R_u^+\) only when the public carrier
identity identifies it as true exterior input; otherwise its target-local
projection is a boundary successor.
\end{definition}

\begin{theorem}[Hypocoercive carrier produces physical expenditure]
\label{thm:hypocoercive-carrier-budget}
On the zero graph-defect branch of
\cref{def:twisted-carrier-invariant}, the evolution and carrier chain-rule
relations for an ordinary maximal branch have ordinary values
\(\dot u=Lu+f_{\Bdry}\) and \(\dd\langle Mu,u\rangle/\dd t\).  On this branch,
\[
 \mathcal K_u(t):=\langle Mu(t),u(t)\rangle,
 \qquad
 \mathfrak D_u(\dd t):=c_Mq_K(u(t))\dd t,
 \qquad
 \Phi_u^+=0
\]
satisfy the carrier balance of
\cref{def:physical-carrier-balance}.  Hence
\begin{equation}
 \label{eq:hypocoercive-physical-budget}
 \mu_u([0,T_*))
 \le \langle Mu_0,u_0\rangle+R_u^+([0,T_*))
 \le m_+\|u_0\|^2+R_u^+([0,T_*))
\end{equation}
after subtracting any known lower carrier level.  If \(\gamma_K>0\), this
expenditure controls every nongauge direction in \(\mathcal G^\perp\).
The commutator calculation is therefore a structural compiler for a physical
budget in a partially dissipative system \citep{Villani2009}.  Its
regularity use is restricted to the active target/source cell selected by the
shell current; nongauge directions outside that cell are not charged to its
target progress.
\end{theorem}

\begin{proof}
Differentiate \(\langle Mu,u\rangle\) and use
\eqref{eq:twisted-carrier-lyapunov}; the positive forcing pairing gives the
declared input measure.  Apply
\cref{thm:carrier-physical-budget}.  When \(\gamma_K>0\),
\(q_K(u)\ge\gamma_K\|P_{\mathcal G}^\perp u\|^2\).
\end{proof}

\begin{definition}[Positive-commutator carrier]
\label{def:positive-commutator-carrier}
Let the evolution be \(\dot u=-iHu\), with \(H=H^*\), and let
\(P=\mathbf 1_\Delta(H)\) be a reducing spectral projection.  A self-adjoint
conjugate operator \(C\) supplies a positive-commutator carrier when the form
\[
 B:=i[H,C]
\]
is well defined and nonnegative on a common invariant core in \(P\Hspace\),
and the expectation \(\langle u(t),Cu(t)\rangle\) has a finite upper bound
\(C_+\) along the orbit.  The strict Mourre branch has
\(PBP\ge\theta P\) modulo the declared gauge for some \(\theta>0\).
\end{definition}

\begin{theorem}[Conservative carrier produces physical expenditure]
\label{thm:positive-commutator-budget}
Under \cref{def:positive-commutator-carrier}, let \(u_0\in P\Hspace\).  On
the zero chain-defect branch, the totalized commutator relation gives
\begin{equation}
 \label{eq:positive-commutator-identity}
 \frac{\dd}{\dd t}\langle u(t),Cu(t)\rangle
 =\langle u(t),Bu(t)\rangle.
\end{equation}
Consequently
\begin{equation}
 \label{eq:positive-commutator-physical-budget}
 \mu_u(Q):=\int_Q\langle u(t),Bu(t)\rangle\dd t,
 \qquad
 \mu_u([0,T_*))
 \le C_+-\langle u_0,Cu_0\rangle.
\end{equation}
When \(PBP\ge\theta P\), this budget controls
\(\theta\int_0^{T_*}\|Pu(t)\|^2\dd t\).  Thus a virial, Morawetz, or Mourre
identity contributes to the same physical-budget calculus whenever its
carrier variation is finite \citep{JensenMourrePerry1984}.  In the
reachability calculus, the contribution is restricted to the retained
target-local spectral/source cell and is charged only through
\cref{thm:local-target-aligned-charging-rule}.
\end{theorem}

\begin{proof}
Differentiate the expectation on the common core to obtain
\eqref{eq:positive-commutator-identity}.  Nonnegativity of \(B\) makes the
right side a measure density.  Integrate and use the upper carrier bound.
The strict estimate follows from \(PBP\ge\theta P\).
On the complementary branch the distributional difference between the two
sides of \eqref{eq:positive-commutator-identity} is the retained
\(\mathfrak H_{\rm chain}^{\Sigma}\)-class, so the theorem does not presume
the chain rule in the fast track.
\end{proof}

\begin{corollary}[Complete routing alternatives for a zero mode]
\label{cor:zero-mode-routing-alternatives}
Let \(e_0\) be the zero mode in
\cref{prop:resistance-zero-mode-split}.  Direct resistance controls
\(e_\perp\).  A target-visible \(e_0\) may be controlled by a
twisted-carrier evaluator or a positive-commutator carrier, with its price
eligible for the same \(\mu_u\) only through its target-local carrier
factorization.  The complementary values of both evaluators select the
harmonic nonroutable branch of
\cref{thm:fast-track-branch-completeness}; it remains a Geom--Caus source
class and does not move to \(\Res\).
\end{corollary}

\section{Compactness and rigidity as a price compiler}
\label{sec:compactness-price}

Compactness is not an additional continuum source and is not part of the
definition of the residual.  Its role is quantitative: it converts a
qualitative zero-price rigidity statement into a uniform positive price on a
closed class of target-approaching windows.  The topology used in this
conversion must be generated from seminorms or graph coordinates derived
from the fixed PDE.  The following lemma is one method for evaluating
\(\delta_*\); its tuple is not input to the fast track.  It is invoked only
on the target/source stratum that produced the window family.  Each failed
clause is totalized into the boundary, chain, or price row of that same
stratum.

\begin{definition}[Canonical compactness--rigidity row]
\label{def:typed-compactness-invariant}
Fix a reachable active branch \(\mathfrak b\) at physical ceiling
\(\Budget\), and let \(\mathcal W_{\mathfrak b}\) be the normalized
pre-contact windows generated by the equation's \(\Abs/\Lift\) grammar that
have the origin, source support, realization mode, shell floor, and earlier
zero coordinates recorded in \(\mathfrak b\).  Inside the closure of its
target-local stratum \(X_{\mathfrak b}\), define
\begin{equation}
 \label{eq:canonical-window-closures}
 K_{\mathfrak b}:=\overline{\mathcal W_{\mathfrak b}},
 \qquad
 B_{\Sigma,\mathfrak b}:=\overline{\mathcal W_{\mathfrak b,\Sigma}},
\end{equation}
where \(\mathcal W_{\mathfrak b,\Sigma}\) consists of those same-origin
windows carrying the selected normalized target-shell crossing.  Thus
\(K_{\mathfrak b}\) is compact and \(B_{\Sigma,\mathfrak b}\) is closed by
construction; neither is an analytic premise and neither contains a
source-transverse window.

On the ordinary admitted value of the local charging graph, after the
storage-transfer row has its zero/root--leaf-bounded value, let \(d_0\) be
the nonnegative window expenditure obtained by restricting the fixed
physical measure to that target-active event.  A failed charging graph or a
nonzero storage return does not enter \(d_0\); it is the corresponding price
or \(\Caus\)-storage successor.  The canonical lower-semicontinuous
relaxation of the admitted \(d_0\) is
\begin{equation}
 \label{eq:canonical-lsc-price}
 \overline d(w):=\sup_{U\ni w}
  \inf\{d_0(v):v\in U\cap\mathcal W_{\mathfrak b}\}
 \in[0,\infty],
\end{equation}
where \(U\) ranges over neighborhoods of \(w\), and the infimum of an empty
set is \(+\infty\).  Equivalently, this is the largest lower-semicontinuous
functional whose restriction to \(\mathcal W_{\mathfrak b}\) is bounded above by
\(d_0\).
The compactness--rigidity value is the exact extended number
\begin{equation}
 \label{eq:compactness-rigidity-gap}
 \delta_{\mathrm{rig}}(\mathfrak b)
 :=\inf_{w\in B_{\Sigma,\mathfrak b}}\overline d(w),
\end{equation}
with value \(+\infty\) when \(B_{\Sigma,\mathfrak b}=\varnothing\).  The branches
\(\delta_{\mathrm{rig}}>0\) and \(\delta_{\mathrm{rig}}=0\) are exhaustive;
the latter is the closed zero-price target-approach class.  Every topology,
normalization, target class, and price in this row is therefore generated
from the fixed PDE, target, and physical balance.
\end{definition}

\begin{definition}[Canonical physical-transfer radius]
\label{def:physical-price-interface}
For each generated event \(Q\), write \(w_Q=\mathcal N_Qu\) and let \(\mu_u\)
be the orbit's physical expenditure measure.  Define
\begin{equation}
 \label{eq:canonical-physical-transfer-radius}
 c_{\mathrm{pay}}^*(\mathfrak b)
 :=\inf_{\substack{Q:\ w_Q\in B_{\Sigma,\mathfrak b}\\
                         0<\overline d(w_Q)<\infty}}
       \frac{\mu_u(Q)}{\overline d(w_Q)}\in[0,\infty],
\end{equation}
using \(+\infty\) for an empty indexing family.  A window with
\(\overline d(w_Q)=\infty\) and finite \(\mu_u(Q)\), or a failure of the
normalization/pullback relation, is retained by the price graph-defect row.
For every ordinary finite-price window,
\begin{equation}
 \label{eq:physical-price-domination}
 \mu_u(Q)\ge
 c_{\mathrm{pay}}^*(\mathfrak b)\,\overline d(w_Q).
\end{equation}
Consequently the physical-transfer question has the exact complementary
values \(c_{\mathrm{pay}}^*>0\) and \(c_{\mathrm{pay}}^*=0\), together with
its typed graph-defect value.  No comparison constant is supplied as data.
\end{definition}

\begin{lemma}[Compactness--rigidity row evaluation]
\label{lem:compactness-rigidity-price}
The value in \eqref{eq:compactness-rigidity-gap} is positive exactly when the
canonical closed target-approach class contains no sequence with relaxed physical
price tending to zero.  If \(B_{\Sigma,\mathfrak b}\ne\varnothing\), the infimum is
attained.
\end{lemma}

\begin{proof}
The set \(B_{\Sigma,\mathfrak b}\) is compact and \(\overline d\) is lower
semicontinuous by relaxation.  Hence the infimum is attained.  Its value is
zero exactly when a minimizing sequence has price tending to zero; compactness
then supplies a convergent subsequence in \(B_{\Sigma,\mathfrak b}\), and lower semicontinuity
identifies the zero-price branch.  The converse follows from the constant
sequence at a zero-price minimizer.
\end{proof}

\begin{theorem}[Uniform physical price branch]
\label{thm:uniform-physical-price}
On the invariant branch
\(c_{\mathrm{pay}}^*(\mathfrak b)>0\) and
\(\delta_{\mathrm{rig}}(\mathfrak b)>0\), every ordinary generated event \(Q\)
whose normalized window belongs to \(B_{\Sigma,\mathfrak b}\) satisfies
\begin{equation}
 \label{eq:uniform-physical-event-price}
 \mu_u(Q)\ge \delta_{\Phys}(\mathfrak b),
 \qquad
 \delta_{\Phys}(\mathfrak b)
 :=c_{\mathrm{pay}}^*(\mathfrak b)
   \delta_{\mathrm{rig}}(\mathfrak b)>0.
\end{equation}
On either zero branch, the corresponding zero-transfer or zero-price sequence
is retained in the same source/mode cell.
\end{theorem}

\begin{proof}
Apply \eqref{eq:physical-price-domination} and then
\eqref{eq:compactness-rigidity-gap}.  The complementary statements follow
from the two infimum definitions and sourcewise closure.
\end{proof}

\begin{theorem}[Borderline passivity price compiler]
\label{thm:borderline-passivity-price}
On the invariant branch \(\beta_*^{\mathrm{red}}=1\) in
\cref{def:reduced-passivity-radius}, define the normalized carrier deficit
\begin{equation}
 \label{eq:normalized-carrier-deficit}
 \Delta_{\mathrm{car}}(V):=D(V)-\Phi_+(V)\ge0.
\end{equation}
Let \(\overline\Delta_{\mathrm{car}}\) be its canonical relaxation
\eqref{eq:canonical-lsc-price} on the generated compact sets
\eqref{eq:canonical-window-closures}, and define
\begin{equation}
 \label{eq:borderline-equality-rigidity}
 \delta_{\mathrm{car}}^*(\mathfrak b)
 :=\inf_{V\in B_{\Sigma,\mathfrak b}}\overline\Delta_{\mathrm{car}}(V).
\end{equation}
Define the clock-transfer radius by the same infimum construction as
\eqref{eq:canonical-physical-transfer-radius}, with
\(\overline d=\overline\Delta_{\mathrm{car}}\); call it
\(c_{\mathrm{clock}}^*(\mathfrak b)\).  On the positive values of both rows,
physical pullback gives
\begin{equation}
 \label{eq:borderline-deficit-domination}
 \mu_u(Q)\ge c_{\mathrm{clock}}^*
 \overline\Delta_{\mathrm{car}}(\mathcal N_Qu),
\end{equation}
and every selected target-bad event has price
\begin{equation}
 \label{eq:borderline-uniform-price}
 \mu_u(Q)\ge
 c_{\mathrm{clock}}^*\delta_{\mathrm{car}}^*(\mathfrak b)>0.
\end{equation}
Thus the critical value \(\beta_*^{\mathrm{red}}=1\) is closed by an
equality--rigidity invariant on the reduced saturation class.  If either
infimum is zero, its minimizing sequence is the complementary price class.
\end{theorem}

\begin{proof}
The canonical relaxation is lower semicontinuous on the compact generated
class, so \cref{lem:compactness-rigidity-price} evaluates
\(\delta_{\mathrm{car}}^*\).  The transfer-radius definition gives
\eqref{eq:borderline-deficit-domination}; multiplication yields
\eqref{eq:borderline-uniform-price}.
\end{proof}

\begin{corollary}[Separator closure of the borderline branch]
\label{cor:separator-borderline-passivity}
On the zero separator-defect branch, the target projection of every
zero-deficit profile vanishes by
\cref{prop:positive-equality-separator}.  Hence the canonical target-bad
closure contains no zero of \(\overline\Delta_{\mathrm{car}}\), so compactness
gives \(\delta_{\mathrm{car}}^*>0\).  The positive clock-transfer branch of
\cref{thm:borderline-passivity-price} then gives the uniform event price; its
zero branch retains the exact zero-transfer sequence.
\end{corollary}

\begin{proposition}[A standard local compactness compiler]
\label{prop:aubin-lions-compiler}
Let \(X_0,X,X_1\) be Banach spaces with
\(X_0\Subset X\hookrightarrow X_1\), the first embedding compact and the
second continuous.  If \(1\le p<\infty\), \(1<q\le\infty\), and a family
of windows is bounded in
\[
 L^p(0,1;X_0)\cap W^{1,q}(0,1;X_1),
\]
then it is relatively compact in \(L^p(0,1;X)\).  Consequently, its closure
agrees with the corresponding compact branch of \(K_{\mathfrak b}\).  The target
and price rows are still evaluated by the canonical target closure and
lower-semicontinuous relaxation; the lemma supplies no separate closedness or
price premise.
\end{proposition}

\begin{proof}
The relative compactness is the Aubin--Lions lemma in the stated range
\citep{Aubin1963,Lions1969,Simon1986}.  The final assertion follows by mapping
that relatively compact family into the generated evaluation closure.
\end{proof}

\begin{remark}[Quotient compactness row]
If a compact group acts continuously on a compact normalized set, its
quotient is compact.  For a noncompact symmetry group, fixing a gauge only
defines a slice.  The generated closure of that slice and the totalized gauge
graph select the compact and graph-defect branches, respectively.  In either
case the quotient operation has \(\Abs\)-ancestry and does not alter the
source of the limiting physical record.
\end{remark}

\begin{remark}[What failure of compactness means]
\label{rem:compactness-invariant}
A nonordinary value of a normalization, target-persistence,
lower-semicontinuity, or physical-pullback graph is pushed into the compact
activity lift with the relevant
\(\Geom/\Caus/\Abs/\Lift/\Bdry\) ancestry.  A zero value of either sharp
infimum supplies its minimizing sequence to the same lift.  Finite-depth
viability then removes the descendants or constructs their recurrent
zero-price path.  No target-visible element is transferred to \(\Res\), and
this optional price accelerator creates no new obligation in the core proof.
\end{remark}

\begin{figure}[tbp]
\centering
\resizebox{.98\linewidth}{!}{%
\begin{tikzpicture}[fw diagram,node distance=8mm and 10mm]
  \node[fw source,text width=2.7cm] (family)
    {target-aligned\\source family};
  \node[fw provenance,text width=2.8cm,right=of family] (quotient)
    {propagated compatible\\quotient \(Y_{\mathrm{red}}\)};
  \node[fw residual,text width=2.8cm,above right=9mm and 11mm of quotient] (closure)
    {closure invariant\\\(\mathfrak H_{\mathcal F}^{\Sigma}=0\)};
  \node[fw geom,text width=2.8cm,right=of quotient] (radius)
    {sharp radius or gap\\\(\alpha_*,\,\delta_*\)};
  \node[fw use,text width=2.8cm,below right=9mm and 11mm of quotient] (equality)
    {critical equality\\separator};
  \node[fw terminal,text width=3.1cm,right=16mm of radius] (ledger)
    {null/divergent closure\\or zero-price successors};
  \draw[fw arrow] (family)--(quotient);
  \draw[fw arrow] (quotient)--(closure);
  \draw[fw arrow] (quotient)--(radius);
  \draw[fw arrow] (quotient)--(equality);
  \draw[fw arrow] (closure)--(ledger);
  \draw[fw arrow] (radius)--(ledger);
  \draw[fw arrow] (equality)--(ledger);
\end{tikzpicture}%
}
\caption{Identity-to-successor compiler.  Once closed null directions are
propagated out, generated identities expose a target closure, domination
radius, physical-price gap, or critical equality relation.  Every nonclosing
value is normalized and passed to the successor operator.  Each evaluator returns both its closing and
complementary structural relation to the ordered algebra.}
\label{fig:estimate-invariant-compiler}
\end{figure}

\FloatBarrier

\section{Constraint tails as finite structural obstructions}
\label{sec:constraint-tails}

Nonlocal constraints often enter a local evolution through an elliptic
multiplier, potential, or pressure.  A direct estimate of the whole nonlocal
field obscures the smaller structural question: which part is fixed by the
local source, which part is a gauge, which part decays from the exterior, and
which finite-dimensional harmonic mode remains visible to the target?  The
following interface makes that separation explicit.
It is formed only after restriction to the active branch.  Its local,
harmonic, gauge, and exterior coordinates are the
\(\Geom/\Caus/\Abs/\Lift/\Bdry\) factorization of that branch's constraint
current; no norm of the full constraint field has an independent
reachability consequence.

\begin{definition}[Elliptic constraint--tail interface]
\label{def:elliptic-constraint-tail}
Let \(U_0\Subset U\) be open sets, let

\[
 \mathcal C p=S(V)\quad\text{in }U
\]

be a linear elliptic constraint in a declared weak sense, and let
\(\mathcal G_{U_0}\) be its gauge space on \(U_0\).  The compiler forms the
product of the following five total graphs.  Its ordinary value consists of:
\begin{enumerate}[label=\textup{(T\arabic*)}]
\item a local solution operator \(\mathcal L_U\) satisfying
  \(\mathcal C\mathcal L_US(V)=S(V)\) on \(U\);
\item a decomposition on \(U_0\),
  \begin{equation}
   \label{eq:constraint-local-tail-split}
   p|_{U_0}=p_{\mathrm{loc}}+h,
   \qquad p_{\mathrm{loc}}:=\mathcal L_US(V)|_{U_0},
   \qquad \mathcal Ch=0\ \text{in }U_0;
  \end{equation}
\item a Banach space \(\mathcal H(U_0)\) of homogeneous solutions with a
  continuous splitting
  \begin{equation}
   \label{eq:homogeneous-tail-splitting}
   \mathcal H(U_0)=\mathcal H_{\mathrm{low}}
      \oplus\mathcal H_{\mathrm{dec}},
  \end{equation}
  where \(\mathcal H_{\mathrm{low}}\) is finite dimensional;
\item a dyadic exterior expansion
  \(P_{\mathrm{dec}}h=\sum_{k\ge0}h_k\) convergent in
  \(\mathcal H_{\mathrm{dec}}\), with
  \begin{equation}
   \label{eq:tail-summability-invariant}
   \mathfrak s_{\mathrm{tail}}(h)
   :=\sum_{k\ge0}\|h_k\|_{\mathcal H_{\mathrm{dec}}}<\infty;
  \end{equation}
\item a linear gauge realization
  \(\Gamma:\mathcal G_{U_0}\to\mathcal H_{\mathrm{low}}\) with closed
  range.
\end{enumerate}
The residual low-mode obstruction is the cokernel class
\begin{equation}
 \label{eq:tail-cokernel-invariant}
 \mathfrak o_{\mathrm{tail}}(p)
 :=[P_{\mathrm{low}}h]
 \in \mathcal H_{\mathrm{low}}/\Ran\Gamma.
\end{equation}
The local term retains the ancestry of \(S(V)\); gauge removal appends
\(\Abs\), and the exterior expansion and its limiting trace append
\(\Bdry\).
None of \textup{(T1)}--\textup{(T5)} is additional specialization data.
Candidate operators, splittings, and dyadic pieces are enumerated from the
typed weak constraint and the fixed universal elliptic library.  The first
relation without an ordinary value is retained as its total-graph successor
with the displayed ancestry.
\end{definition}

\begin{theorem}[Constraint--tail structural normal form]
\label{thm:constraint-tail-normal-form}
On the ordinary value of the elliptic constraint--tail product graph, every
field has, on \(U_0\), the
decomposition
\begin{equation}
 \label{eq:constraint-tail-normal-form}
 p=p_{\mathrm{loc}}+\Gamma g+h_{\mathrm{dec}}+h_{\mathrm{obs}},
\end{equation}
where \(g\in\mathcal G_{U_0}\),
\(h_{\mathrm{dec}}\in\mathcal H_{\mathrm{dec}}\), and
\(h_{\mathrm{obs}}\) is any representative of the class
\(\mathfrak o_{\mathrm{tail}}(p)\).  The first three terms are closed by,
respectively, the local source law, a propagated quotient, and the
summability invariant \eqref{eq:tail-summability-invariant}.  Consequently:
\begin{enumerate}[label=\textup{(C\arabic*)}]
\item if \(\mathfrak o_{\mathrm{tail}}(p)=0\), the entire constraint field
  closes using only its inherited source, \(\Abs\), and \(\Bdry\);
\item if \(\mathfrak o_{\mathrm{tail}}(p)\ne0\) and its target projection
  vanishes, its remaining class lies in \(\Null_\Sigma\);
\item if \(\mathfrak o_{\mathrm{tail}}(p)\ne0\) and is target-visible, it is
  a finite-dimensional \(\Bdry/\Abs\)-ancestry occurrence evaluated by the
  carrier and price rows.  It is not a new residual source.
\end{enumerate}
Thus the nonlocal constraint is reduced internally to the local source graph,
the summability graph, and a finite-dimensional target-projected cokernel.
Every nonordinary value is already a witness-indexed successor rather than an
additional analytic task.
\end{theorem}

\begin{proof}
Apply the continuous projections in
\eqref{eq:homogeneous-tail-splitting} to \(h\).  Closed range and finite
dimensionality give
\(
 \mathcal H_{\mathrm{low}}
 =\Ran\Gamma\oplus H_{\mathrm{obs}}
\)
for a finite-dimensional complement \(H_{\mathrm{obs}}\).  This gives
\eqref{eq:constraint-tail-normal-form}; changing the complement changes
\(h_{\mathrm{obs}}\) only by a gauge term and hence preserves the cokernel
class.  Absolute convergence in
\eqref{eq:tail-summability-invariant} closes the decaying series.  The three
routing conclusions follow from quotient propagation,
\cref{thm:source-preserving-completion}, and the definition of
\(\Null_\Sigma\).
\end{proof}

\begin{proposition}[Constraint tails enter the physical budget only through a
carrier law]
\label{prop:constraint-tail-carrier-law}
On the ordinary constraint-tail branch, let the signed tail contribution be
\(\Psi_{\mathrm{tail}}\) to a localized carrier balance.  It is charged to
the physical expenditure measure only if the equation supplies a measure
inequality of one of the forms
\begin{equation}
 \label{eq:constraint-tail-physical-domination}
 (\Psi_{\mathrm{tail}})_+\le\Phi_u^+,
 \qquad\text{or}\qquad
 (\Psi_{\mathrm{tail}})_+\le R_u^+,
\end{equation}
with the right side appearing in
\eqref{eq:physical-carrier-balance}.  Tail summability alone proves analytic
closure but supplies no physical price and no budget bound.  In a target
argument these inequalities are first restricted to the active shell/source
cell and then passed through \eqref{eq:target-aligned-charge-rule}.
\end{proposition}

\begin{proof}
The carrier compiler uses only the measure terms displayed in
\eqref{eq:physical-carrier-balance}.  Either inequality in
\eqref{eq:constraint-tail-physical-domination} places the positive tail
contribution in that balance.  Without such a comparison,
\eqref{eq:tail-summability-invariant} is a norm statement on a homogeneous
field and has no implication for \(\mu_u\).
\end{proof}

\section{Covariant entropy cocycles and physical production}
\label{sec:entropy-cocycles}

An entropy enters the reachability argument through the physical production
that it controls along one orbit.  Local entropy formulas in changing
blow-up charts
must therefore be joined by an explicit cocycle law.  Independent
windowwise monotonicity does not telescope.  Perelman's entropy formula is a
classical model of covariant monotonicity with a rigid equality case
\citep{Perelman2002}.  The formulation below derives the required
equation-independent cocycle contracts directly from a local balance and its
reset relation.
For reachability, each local balance and reset is first compiled by
\cref{thm:target-local-fact-factorization}.  Entropy production is therefore
not a second budget and is not charged across unrelated windows: it
contributes to the sole physical measure only on the same-origin
target/source cell on which the public carrier proves domination.

\begin{definition}[Physical entropy cocycle]
\label{def:physical-entropy-cocycle}
Let \(I_j=[a_j,b_j]\), \(j\ge1\), be ordered normalized windows obtained from
pairwise disjoint physical events \(Q_j\).  A physical entropy cocycle on
this family consists of
\begin{enumerate}[label=\textup{(E\arabic*)}]
\item scalar primitives \(\Entropy_j:I_j\to\mathbb R\);
\item nonnegative Radon production measures \(\Production_j\) and
  nonnegative error measures \(r_j\) on \(I_j\);
\item reset errors \(e_j\ge0\) and a constant \(c_{\Entropy}>0\);
\item constants \(M_+<\infty\) and \(M_->-\infty\);
\end{enumerate}
such that, for \(a_j\le s<t\le b_j\),
\begin{align}
 \Entropy_j(t)+c_{\Entropy}\Production_j((s,t])
 &\le \Entropy_j(s)+r_j((s,t]),
 \label{eq:local-entropy-production}\\
 \Entropy_{j+1}(a_{j+1})
 &\le \Entropy_j(b_j)+e_j,
 \label{eq:entropy-reset-comparison}\\
 \Entropy_1(a_1)&\le M_+,
 &\Entropy_j(t)&\ge M_-.
 \label{eq:entropy-two-sided-bound}
\end{align}
The measures \(\Production_j\) are coordinate expressions of one physical
production law on the events \(Q_j\); chart Jacobians and time changes are
included in that law.  The half-open convention in
\eqref{eq:local-entropy-production} prevents endpoint atoms from being
counted twice.
\end{definition}

\begin{theorem}[Entropy-cocycle gluing]
\label{thm:entropy-cocycle-gluing}
For every \(N\ge1\), a physical entropy cocycle satisfies
\begin{equation}
 \label{eq:entropy-cocycle-bound}
 c_{\Entropy}\sum_{j=1}^{N}\Production_j(I_j)
 \le
 M_+-M_-
 +\sum_{j=1}^{N}r_j(I_j)
 +\sum_{j=1}^{N-1}e_j.
\end{equation}
If the two error series are finite, then
\(\sum_j\Production_j(I_j)<\infty\).  If in addition every selected bad
window satisfies \(\Production_j(I_j)\ge\delta>0\), their number is at most
\begin{equation}
 \label{eq:entropy-window-count}
 \frac{M_+-M_-+\sum_jr_j(I_j)+\sum_je_j}
      {c_{\Entropy}\delta}.
\end{equation}
\end{theorem}

\begin{proof}
Apply \eqref{eq:local-entropy-production} on every whole window.  Summing
gives
\[
 c_{\Entropy}\sum_{j=1}^{N}\Production_j(I_j)
 \le
 \sum_{j=1}^{N}
   \bigl(\Entropy_j(a_j)-\Entropy_j(b_j)\bigr)
 +\sum_{j=1}^{N}r_j(I_j).
\]
The reset inequalities telescope the first sum from above:
\[
 \sum_{j=1}^{N}
   \bigl(\Entropy_j(a_j)-\Entropy_j(b_j)\bigr)
 \le
 \Entropy_1(a_1)-\Entropy_N(b_N)+\sum_{j=1}^{N-1}e_j.
\]
Use \eqref{eq:entropy-two-sided-bound}.  This proves
\eqref{eq:entropy-cocycle-bound}.  Monotone convergence gives the countable
statement, and \(N\delta\le\sum_{j=1}^{N}\Production_j(I_j)\) gives
\eqref{eq:entropy-window-count}.
\end{proof}

\begin{definition}[Perelman-type covariant invariant row]
\label{cert:perelman-compiler}
A Perelman-type construction is evaluated by the totalized tuple
\(\mathfrak I_{\rm Per}\) whose six coordinates are the graph defects of the
following relations for the fixed PDE and its generated window atlas:
\begin{enumerate}[label=\textup{(P\arabic*)}]
\item Each local primitive is a quotient-invariant scalar, typically a
  localized energy or variance plus an adjoint entropy and explicit
  pressure, drift, modulation, and boundary corrections.
\item Differentiation along the \emph{fixed} PDE gives
  \eqref{eq:local-entropy-production}; all unsigned remainders are placed in
  the displayed error measure, not suppressed formally.
\item The target-visible quotient of production by its nonnegative physical
  floor \(\mathfrak d\) is recorded by the price-gap row.
\item The primitives have the lower bound in
  \eqref{eq:entropy-two-sided-bound}.  Its sign must match the orientation
  of \eqref{eq:local-entropy-production}.
\item Transport along the same physical orbit, chart pullback, and any change
  of a finite-horizon adjoint datum satisfy
  \eqref{eq:entropy-reset-comparison}, with summable reset errors.
\item Production either defines the expenditure measure \(\mu_u\), as in
  \cref{cor:entropy-physical-budget}, or is dominated by the declared
  \(\mu_u\) through an explicit change-of-variables inequality of the form
  \eqref{eq:physical-price-domination}.
\end{enumerate}
The entropy-cocycle branch is the common zero set of these six target-active
graph defects.  Any nonzero coordinate is retained with its inherited source
ancestry and sent to the form, chain, price, tail, or boundary row indicated
by the failed relation.  Thus the construction neither presumes the entropy
scalar nor accepts it as additional favorable data.  These coordinates
isolate the production-and-cocycle mechanism of Perelman's entropy formula
\citep{Perelman2002}.
\end{definition}

\begin{proposition}[Weighted zero-production rigidity]
\label{prop:weighted-zero-production-rigidity}
Let \(D\subset\mathbb R^d\) be connected and open, let
\(v\in L^2(I;W^{1,2}_{\mathrm{loc}}(D))\), and let
\(m\ge m_*>0\) almost everywhere on \(D\times I\).  If a nonnegative
production functional obeys
\begin{equation}
 \label{eq:weighted-dirichlet-floor}
 \mathfrak d(w)
 \ge c_0\int_I\int_Dm(x,t)|\nabla v(x,t)|^2\dd x\dd t
\end{equation}
with \(c_0>0\), then \(\mathfrak d(w)=0\) implies that
\(v(\cdot,t)\) is spatially constant on \(D\) for almost every \(t\).
Consequently, any closed bad class carrying a strictly positive
mean-subtracted oscillation is disjoint from \(\{\mathfrak d=0\}\).  The
regularity compiler invokes this conclusion only on the target-local stratum
whose selected production is \(\mathfrak d\); it gives no exclusion for an
unrelated bad class or source cell.
\end{proposition}

\begin{proof}
Nonnegativity and \eqref{eq:weighted-dirichlet-floor} imply
\(\nabla v=0\) almost everywhere.  A Sobolev function with zero weak gradient
on a connected domain is almost everywhere constant.  Its mean-subtracted
oscillation is therefore zero, which excludes membership in the stated bad
class.
\end{proof}

\begin{proposition}[Physical-clock comparison]
\label{prop:physical-clock-comparison}
Let \(\mu\) be a physical measure on an event.  On the absolutely-continuous
clock branch, let a normalized production measure have density \(w\ge0\):
\(\dd\nu=w\dd\mu\).  On a measurable set \(Q\),
\begin{align}
 w\le C,\quad 0<C<\infty,\quad\text{a.e. on }Q
 &\Longrightarrow \mu(Q)\ge C^{-1}\nu(Q),
 \label{eq:clock-upper-comparison}\\
 w\ge c>0\quad\text{a.e. on }Q
 &\Longrightarrow \mu(Q)\le c^{-1}\nu(Q).
 \label{eq:clock-lower-comparison}
\end{align}
Thus a lower bound for normalized production yields a lower physical price
only from an upper bound such as \eqref{eq:clock-upper-comparison}; the
opposite density bound has the opposite implication.
\end{proposition}

\begin{proof}
Integrate the pointwise bounds in \(\nu(Q)=\int_Qw\dd\mu\).
\end{proof}

\begin{corollary}[Entropy as the physical expenditure measure]
\label{cor:entropy-physical-budget}
On the finite-error branch of \cref{def:physical-entropy-cocycle} and
\eqref{eq:entropy-cocycle-bound}, on the atomic
\(\sigma\)-algebra generated by the disjoint events \(Q_j\), define
\[
 \nu_u\!\left(\bigsqcup_{j\in S}Q_j\right)
 :=\sum_{j\in S}\Production_j(I_j),
 \qquad S\subseteq\mathbb N.
\]
Then \(\nu_u\) is a finite, countably additive physical measure on that event
family.  It coincides with the restriction of \(\mu_u\) in
\cref{def:physical-budget} precisely when the public carrier identity
identifies this entropy production as physical expenditure.  Otherwise it is
a generated nonnegative structural coordinate and cannot replace or enlarge
\(\mu_u\).  Even in the coincident case, a shell event is charged only through
its target-local factorization \eqref{eq:target-local-fact-normal-form}.
\end{corollary}

\begin{proof}
Countable additivity follows directly from the nonnegative series defining
\(\nu_u\); finiteness is \cref{thm:entropy-cocycle-gluing}.
\end{proof}

\begin{remark}[Equality models are row-local]
A Gaussian may be the equality state of an adjoint Fisher-information term,
while a constant field is the equality state of a weighted Dirichlet floor.
These statements concern different nonnegative terms.  A valid rigidity
argument uses the equality set of the particular floor that is both compactly
stable and linked to target badness; it does not merge the two equality
models.
\end{remark}

\begin{figure}[tbp]
\centering
\resizebox{.98\linewidth}{!}{%
\begin{tikzpicture}[fw diagram,node distance=8mm and 10mm]
  \node[fw source,text width=2.7cm] (orbit)
    {one physical orbit};
  \node[fw provenance,text width=2.7cm,right=of orbit] (atlas)
    {ordered chart windows\\\(I_j\)};
  \node[fw geom,text width=2.65cm,above right=8mm and 11mm of atlas] (local)
    {local production\\\(\Production_j\)};
  \node[fw provenance,text width=2.65cm,below right=8mm and 11mm of atlas] (reset)
    {reset comparison\\\(+e_j\)};
  \node[fw use,text width=2.8cm,right=36mm of atlas] (glue)
    {one scalar cocycle\\and one physical measure};
  \node[fw terminal,text width=2.8cm,right=of glue] (finite)
    {finite total production\\and finitely many \(\delta\)-windows};
  \draw[fw arrow] (orbit)--(atlas);
  \draw[fw arrow] (atlas)--(local);
  \draw[fw arrow] (atlas)--(reset);
  \draw[fw arrow] (local)--(glue);
  \draw[fw arrow] (reset)--(glue);
  \draw[fw arrow] (glue)--(finite);
\end{tikzpicture}%
}
\caption{Covariant entropy-cocycle accounting.  Local production estimates become
orbitwise additive only after chart and adjoint-reset comparisons identify
them along the same-origin physical orbit.  The output is a finite physical production measure that
can serve as the orbit expenditure \(\mu_u\).}
\label{fig:physical-entropy-cocycle}
\end{figure}

\FloatBarrier


\section{Wick-rotated compensated quotients and structural-action learning}
\label{sec:wick-compensated-learning}

This section acts on the general nonreversible compensated-quotient record
already defined in Part~I.  On its spectral-\(\Abs\) branch, the Wick map
replaces the real quotient evolution by unitary spectral evolution and fixes
the represented real \(\Geom\) operator, the selected \(\Caus\) operator,
their authority record, and every saturated charge.  The resulting generator
is nonnormal in general, while its quotient response has a positive
self-adjoint learning kernel.

Thus the only new datum is the choice to apply the Wick map.  All source
typing, quotient and fiber maps, geometric forms, causal provenance,
authority restrictions, compensation inputs, and resource coordinates are
those of the input Part~I record.  The transformed response stability,
response kernel, memory kernel, action, and learning rules below are
consequences evaluated on that same record.
Its representation, confidence, authority, and model-selection coordinates
belong to the inherited Part~I computational ledger; they never augment the
continuum physical budget.  Conversely, any balance in this section used for
PDE reachability must first pass through
\cref{thm:target-local-fact-factorization,thm:local-target-aligned-charging-rule}
on the contact-generated source cell.  The learning bounds themselves make
no continuum regularity claim.

\subsection{Spectral quotient and real geometric complement}
\label{subsec:wick-block-model}

\begin{definition}[Wick transform of an inherited compensated record]
\label{def:wick-compensated-quotient}
Let \(\mathscr C_I\) be a dynamically qualified Part~I master compensated
quotient--geometry record in its represented spectral-\(\Abs\) realization.
Write \(\mathcal Q_{\mathbb R}\) and \(\mathcal Y_{\mathbb R}\) for its
quotient and complementary real Hilbert spaces.  In the notation of that
record, \(H_{\Abs}\) is the nonnegative self-adjoint quotient operator,
\(G\) is the real nonnegative geometric operator, and
\[
 K_C=\begin{pmatrix}0&C^*\\-C&A_{\mathcal Y}\end{pmatrix}
\]
is the actual selected skew \(\Caus\) block with its existing authority and
provenance certificate.  These are fields of \(\mathscr C_I\), not new
structural data.

Let \(\mathcal Q=\mathcal Q_{\mathbb R}\otimes_{\mathbb R}\mathbb C\) and
\(\mathcal Y=\mathcal Y_{\mathbb R}\otimes_{\mathbb R}\mathbb C\), and
extend every field of \(\mathscr C_I\) complex-linearly.  The Wick transform
acts only on the quotient evolution.  The complex extension of \(G\)
commutes with conjugation and restricts to the original operator on
\(\mathcal Y_{\mathbb R}\), so the \(\Geom\) generator remains the
complexification of the same real geometric operator (although a coupled
trajectory need not remain in \(\mathcal Y_{\mathbb R}\)):
\begin{equation}
 \label{eq:wick-record-map}
 \mathfrak W_{\Abs}:
 -H_{\Abs}\oplus(-G)+K_C
 \longmapsto
 -iH_{\Abs}\oplus(-G)+K_C.
\end{equation}
Equivalently, on \(\mathcal H_W:=\mathcal Q\oplus\mathcal Y\), its
transformed generator is
\begin{equation}
 \label{eq:wick-block-generator}
 \mathcal L_W(C)
 :=\begin{pmatrix}
      -iH_{\Abs}&C^*\\
      -C&-G+A_{\mathcal Y}
    \end{pmatrix}
 =-\begin{pmatrix}0&0\\0&G\end{pmatrix}
  +\begin{pmatrix}-iH_{\Abs}&C^*\\-C&A_{\mathcal Y}\end{pmatrix}.
\end{equation}
The \emph{positive self-energy normal form} is the inherited subbranch
\(A_{\mathcal Y}=0\); the general inherited branch retains the full
nonnormal resolvent.  The spectral calculus gives
\begin{equation}
 \label{eq:wick-spectral-calculus}
 e^{-\tau H_{\Abs}}
 =\int_{[0,\infty)}e^{-\tau\omega}\,\dd E_{H_{\Abs}}(\omega),
 \qquad
 e^{-itH_{\Abs}}
 =\int_{[0,\infty)}e^{-it\omega}\,\dd E_{H_{\Abs}}(\omega).
\end{equation}
The two spectral integrals are ordinary for every vector in their displayed
time domains.  Their identification by the notation \(\tau=it\) is evaluated
by the Part~I analytic-continuation relation: it is an ordinary equality on
the record's continuation domain and otherwise carries that relation's
nonordinary boundary value.  No continuation domain is added to structural
admissibility.
\end{definition}

The map \(\mathfrak W_{\Abs}\) preserves the source cell and the complete
record.  In particular, \(H_{\Abs}\), \(-iH_{\Abs}\), and their spectral
projectors have the same \(\Abs\) ancestry; \(G,C,A_{\mathcal Y}\), the
authority envelope, and their costs are unchanged.  Equation
\eqref{eq:wick-spectral-calculus} is spectral functional calculus, rather
than a formal substitution in a learned trajectory.

\begin{proposition}[Exact inheritance under the Wick map]
\label{prop:wick-exact-inheritance}
Let \(\mathscr C_I^W=\mathfrak W_{\Abs}(\mathscr C_I)\).  Under the canonical
scalar-extension embeddings,
\begin{equation}
 \label{eq:wick-inherited-record-identities}
 \Src(\mathscr C_I^W)=\Src(\mathscr C_I),\qquad
 \Anc(\mathscr C_I^W)=\Anc(\mathscr C_I),\qquad
 \mathfrak A(\mathscr C_I^W)=\mathfrak A(\mathscr C_I),\qquad
 \mathbf c_{\rm str}(\mathscr C_I^W)=\mathbf c_{\rm str}(\mathscr C_I).
\end{equation}
The represented \(\Geom\) form, selected \(\Caus\) current, action metric,
and quotient/fiber transports agree with their real restrictions in
\(\mathscr C_I\).  Consequently the source, structural-admissibility, and
authority gates have the same truth value before and after the map.  The
inherited compensation data remain fixed, while every response or dynamic-
stability quantity containing the quotient time generator is reevaluated
with \(-H_{\Abs}\) replaced by \(-iH_{\Abs}\).
\end{proposition}

\begin{proof}
Complexification is an isometric scalar extension and commutes with sums,
adjoints, compositions, projections, and restrictions to the original real
spaces.  Equation \eqref{eq:wick-record-map} fixes every block except the
quotient diagonal.  The source, ancestry, authority, structural cost,
current, metric, and transport fields therefore remain identical.  Their
algebraic and authority gates are preserved.  A resolvent, self-energy,
memory, stability factor, or learning kernel contains the transformed
diagonal and is therefore recomputed as a derived Wick response rather than
copied from the input record.
\end{proof}

\begin{table}[tbp]
\centering
\small
\begin{tabular}{@{}p{.23\linewidth}p{.31\linewidth}p{.36\linewidth}@{}}
\toprule
Wick quantity & Part~I source & Exhaustive evaluation\\
\midrule
Generator realization
& dynamic and sectorial realization field
& bounded or ordinary closed-form value; otherwise typed realization defect\\
Resolvent and self-energy
& full-generator dynamic record
& Hille--Yosida value on \(\Re z>0\); totalized graph value elsewhere\\
Response dimension
& finite evaluation carrier and effective-dimension machinery
& finite carrier trace; population extended trace in \([0,\infty]\)\\
Spectral gap and commutation
& represented operators and spectral calculus
& computed gap/status; modal formula or exact singular-value formula\\
Legal bridge and action
& Caus authority envelope and current metric
& distance to the legal-current closure; empty family gives \(+\infty\)\\
Identification
& represented dynamics chart and score candidates
& derived design range/kernel; full or partial identification bound\\
Statistical law
& i.i.d., mixing, sequential, martingale, and channel cells
& recordwise minimum ordinary radius; every unavailable candidate equals
  \(+\infty\)\\
RL conversion
& safe-action, behavior/deployment, and actor records
& extended critic, converter, support, and policy coordinates\\
Model choice
& exact source cells, charged nets, Rademacher union, and PAC--Bayes mixture
& finite-stage saturated profile and Pareto strict inverse; empty/nonattained
  limit is retained and typed\\
Continuum energy
& operator-range and approximation records
& Moore--Penrose quadratic form and extended suprema; no closed-range or
compactness gate\\
\bottomrule
\end{tabular}
\caption{Inheritance and totalization ledger for the Wick specialization.
Every row is an evaluation of Part~I data or a consequence of it.}
\label{tab:wick-inheritance-totalization}
\end{table}
\FloatBarrier

\begin{proposition}[No-new-premise closure of the Wick specialization]
\label{prop:wick-no-new-premise-closure}
The input domain of \(\mathfrak W_{\Abs}\) is exactly the represented,
dynamically qualified spectral-\(\Abs\) cell of the Part~I master record.
Within that inherited domain, the Wick chapter adds no structural, analytic,
statistical, or model-selection premise.  More precisely:
\begin{enumerate}[label=\textup{(\roman*)}]
\item bounded, sectorial, unbounded-bridge, authority, source, channel,
dependence, control, and model-family alternatives are values of existing
Part~I record fields;
\item spectral gaps, commutators, design excitation, operator errors, source
residuals, converter constants, policy divergences, and optimization errors
are computed coordinates in \([0,\infty]\), not admission conditions;
\item a missing analytic or statistical certificate has its Part~I
nonordinary value.  Every dependent upper-error or risk candidate has value
\(+\infty\), every dependent lower-usefulness candidate has value zero, and
a derived coordinate without an ordered extension retains its typed
nonordinary value;
\item an empty legal, zero-action, or risk-feasible family produces its typed
authority, representation, confidence, or budget status; and
\item every finite-rate formula is the restriction of a general extended
bound to the named Part~I rate cell, while all complementary cells retain the
general bound.
\end{enumerate}
Consequently no favorable branch is used to define admissibility for this
specialization.  The only new operation is
\(-H_{\Abs}\mapsto-iH_{\Abs}\), after which all quantities are derived or
totalized.
\end{proposition}

\begin{proof}
Exact inheritance is \cref{prop:wick-exact-inheritance}.  Part~I already
partitions represented records by realization, authority, source, channel,
dependence, and control status and supplies charged finite stages for family
selection.  Evaluation of a field gives either its ordinary value or its
existing nonordinary value.  Extended-real arithmetic, the conventions
\(\inf\varnothing=+\infty\) and \(\sup\varnothing=0\), and typed empty-family
statuses make every construction total.  The propositions below derive the
new response, action, loss, and transfer coordinates from those fields and
make each finite specialization an explicit record cell.
\end{proof}

\begin{proposition}[Generation and physical budget]
\label{prop:wick-generation-budget}
On the bounded realization branch of the inherited Part~I record,
\(\mathcal L_W(C)\) generates a contraction semigroup and a
solution \(u=(x,y)\) of \(\dot u=\mathcal L_W(C)u\) satisfies
\begin{align}
 \frac{\dd}{\dd t}\|x\|^2
 &=2\operatorname{Re}\langle Cx,y\rangle,
 \label{eq:wick-quotient-balance}\\
 \frac{\dd}{\dd t}\|y\|^2
 &=-2\operatorname{Re}\langle Cx,y\rangle
   -2\|G^{1/2}y\|^2,
 \label{eq:wick-geometry-balance}\\
 \frac{\dd}{\dd t}\bigl(\|x\|^2+\|y\|^2\bigr)
 &=-2\|G^{1/2}y\|^2.
 \label{eq:wick-total-energy-balance}
\end{align}
Consequently the Wick-rotated \(\Abs\) block and the authorized \(\Caus\)
block transfer carrier between the two sectors without expenditure.  The
real \(\Geom\) block supplies the unique nonnegative expenditure in this
model.
\end{proposition}

\begin{proof}
The second matrix in \eqref{eq:wick-block-generator} is skew-adjoint.
Taking real parts in the two component equations gives
\eqref{eq:wick-quotient-balance} and
\eqref{eq:wick-geometry-balance}; their transfer terms cancel.  Thus
\(\mathcal L_W(C)\) is dissipative.  The bounded perturbation theorem, or
direct finite-dimensional exponentiation, gives the contraction semigroup.
\end{proof}

\begin{proposition}[Closed-operator and sectorial-complement realization]
\label{prop:wick-sectorial-realization}
On the closed sectorial branch of the inherited Part~I record, the
complexified complement form defines its maximal dissipative operator
\(-G+A_{\mathcal Y}\).  On its inherited bounded-bridge cell, the block operator
\(\mathcal L_W(C)\) on
\(D(H_{\Abs})\oplus D(-G+A_{\mathcal Y})\) is maximal dissipative.  The same
conclusion holds on the ordinary unbounded-bridge realization cell produced
by the Part~I sectorial Caus-authority interface.  Its complementary
nonordinary cell retains the typed realization defect.  Identity
\eqref{eq:wick-total-energy-balance} holds in integrated form,
\begin{equation}
 \label{eq:wick-integrated-energy-balance}
 \|u(t)\|^2+2\int_0^t\|G^{1/2}y(s)\|^2\dd s=\|u(0)\|^2,
\end{equation}
on the equality branch, and with \(\le\) for weak limits carrying a
nonnegative energy defect.
\end{proposition}

\begin{proof}
The complement sectorial form gives \(-G+A_{\mathcal Y}\), while
\(-iH_{\Abs}\) generates a unitary group.  Their direct sum is maximal
dissipative.  The off-diagonal bridge is bounded and skew-adjoint, so the
bounded dissipative perturbation theorem preserves maximal dissipativity.
The ordinary unbounded-bridge conclusion is exactly the closed-operator
perturbation conclusion already certified by the Part~I sectorial interface
\citep{Kato1995,Ouhabaz2005}.  Approximation by vectors in a common core and
lower semicontinuity give the integrated identity or its defect version.
\end{proof}

\subsection{Lumpability defect, self-energy, and memory}
\label{subsec:wick-feshbach-memory}

Let \(U:\mathcal Q\to\mathcal H_W\) be \(Ux=(x,0)\), let
\(P=U^*\), and let \(P_\Delta=\mathbf1_\Delta(H_{\Abs})\) be a declared
spectral band.

\begin{proposition}[Band-lumpability defect]
\label{prop:wick-band-lumpability}
The compressed quotient generator is
\(P\mathcal L_WU=-iH_{\Abs}\), and its full-operator defect on \(\Delta\)
is
\begin{equation}
 \label{eq:wick-band-defect}
 E_{W,\Delta}
 :=(I-UP)\mathcal L_WUP_\Delta
 =\binom{0}{-CP_\Delta}.
\end{equation}
Hence the spectral band is exactly lumpable if and only if
\(CP_\Delta=0\).  Its intrinsic quasi-lumpability size is
\(\varepsilon_{W,\Delta}=\|CP_\Delta\|\), supplemented in an empirical
record by the separately measured model and compression residuals.  Under
the canonical complexification isometry this is exactly the norm of the
corresponding cross-sector defect in \(\mathscr C_I\); Wick rotation changes
the quotient diagonal and introduces no new lumpability hypothesis.
\end{proposition}

\begin{proof}
Apply the two block rows of \(\mathcal L_W\) to \((x,0)\), compress the
first component, and subtract its lift.  The off-diagonal block is fixed by
\eqref{eq:wick-record-map}, so the resulting vector is the complexification
of the inherited Part~I defect.  Complexification is isometric, which proves
the norm statement.
\end{proof}

\begin{theorem}[Feshbach response and geometric self-energy]
\label{thm:wick-feshbach-response}
The Part~I realization field is exhaustive.  On either ordinary realization
cell of \cref{prop:wick-generation-budget,prop:wick-sectorial-realization},
every \(z\in\mathbb C\) with \(\operatorname{Re}z>0\) satisfies, by maximal
dissipativity,
\begin{equation}
 \label{eq:wick-derived-resolvent-bounds}
 \bigl\|(z+G-A_{\mathcal Y})^{-1}\bigr\|
 \le\frac1{\operatorname{Re}z},
 \qquad
 \bigl\|(z-\mathcal L_W(C))^{-1}\bigr\|
 \le\frac1{\operatorname{Re}z}.
\end{equation}
Define
\begin{equation}
 \label{eq:wick-self-energy}
 \Sigma_C(z)
 :=C^*(z+G-A_{\mathcal Y})^{-1}C,
 \qquad
 F_W(z):=z+iH_{\Abs}+\Sigma_C(z).
\end{equation}
On the ordinary inherited unbounded-bridge cell, these expressions denote the
closed forms and associated operators supplied by that Part~I realization
certificate.
Then \(F_W(z)\) is invertible and
\begin{equation}
 \label{eq:wick-compressed-resolvent}
 P(z-\mathcal L_W(C))^{-1}U=F_W(z)^{-1}.
\end{equation}
In the positive self-energy normal form, \(z=\lambda>0\) gives
\begin{equation}
 \label{eq:wick-positive-self-energy}
 \Sigma_C(\lambda)=C^*(\lambda+G)^{-1}C\ge0.
\end{equation}
Thus real geometric relaxation compensates the imaginary quotient defect by
a positive frequency-dependent self-energy.  For
\(\operatorname{Re}z\le0\), the same block formula is the ordinary value of
the existing totalized resolvent relation whenever that value exists.  A
pole, domain failure, or noninvertible Schur complement is its nonordinary
graph value.  Hence every \(z\) is accounted for, while the learning kernel
uses the automatically ordinary half-line \(\lambda>0\) on the ordinary
realization cells.  On the complementary nonordinary realization cell, the
self-energy, Schur complement, and learning kernel inherit the nonordinary
graph value and every response-dependent upper certificate equals
\(+\infty\); the record is retained with its realization status.
\end{theorem}

\begin{proof}
Both the complement generator \(-G+A_{\mathcal Y}\) and the full Wick
generator are maximal dissipative by
\cref{prop:wick-generation-budget,prop:wick-sectorial-realization}.  The
right half-plane therefore belongs to both resolvent sets and the
Hille--Yosida estimate gives \eqref{eq:wick-derived-resolvent-bounds}.
The block operator \(z-\mathcal L_W\) is
\[
 \begin{pmatrix}
  z+iH_{\Abs}&-C^*\\ C&z+G-A_{\mathcal Y}
 \end{pmatrix}.
\]
The lower-right block and the full block are invertible.  Block Gaussian
elimination therefore shows that their Schur complement \(F_W(z)\) is
invertible and gives \eqref{eq:wick-compressed-resolvent}.  On the ordinary
inherited unbounded-bridge cell the same factorization is first taken on the
common form core and then closed by its Part~I realization theorem.  When
\(A_{\mathcal Y}=0\) and \(z=\lambda>0\), the middle resolvent in
\eqref{eq:wick-positive-self-energy} is positive self-adjoint.  Outside the
right half-plane, totalization records exactly the failure of one of these
ordinary elimination steps.
\end{proof}

\begin{corollary}[Mori--Zwanzig memory equation]
\label{cor:wick-memory-equation}
On the bounded-bridge cell, for classical initial data
\((x_0,y_0)\), the quotient coordinate satisfies
\begin{equation}
 \label{eq:wick-memory-equation}
 \dot x(t)
 =-iH_{\Abs}x(t)
  -\int_0^tK_C(t-s)x(s)\dd s
  +C^*e^{t(-G+A_{\mathcal Y})}y_0,
 \qquad
 K_C(t):=C^*e^{t(-G+A_{\mathcal Y})}C.
\end{equation}
For every \(\operatorname{Re}z>0\), the Laplace transform of \(K_C\) is
\(\Sigma_C(z)\).  Equations
\eqref{eq:wick-compressed-resolvent} and
\eqref{eq:wick-memory-equation} are respectively the Feshbach and
Mori--Zwanzig forms of the same quotient elimination
\citep{Feshbach1958,Zwanzig1960,Mori1965}.
On the ordinary unbounded-bridge cell, the same display is interpreted in the
mild input--output or closed-form sense supplied by the Part~I realization
record, first on its common core and then by closure.  A datum outside that
recorded admissibility domain has the inherited nonordinary evolution value;
no additional observation-domain premise is imposed here.
\end{corollary}

\begin{proof}
Variation of constants in the \(y\)-equation gives
\[
 y(t)=e^{t(-G+A_{\mathcal Y})}y_0
 -\int_0^te^{(t-s)(-G+A_{\mathcal Y})}Cx(s)\dd s.
\]
Substitute this identity in \(\dot x=-iH_{\Abs}x+C^*y\).  Laplace
transformation gives \eqref{eq:wick-self-energy}.  On the ordinary
unbounded-bridge cell, the Part~I common-core factorization in
\cref{thm:wick-feshbach-response} gives the same identity, and closure gives
the recorded mild input--output relation.
\end{proof}

\begin{definition}[Dark quotient space]
\label{def:wick-dark-space}
On the bounded-bridge cell, the dark quotient space is
\begin{equation}
 \label{eq:wick-dark-space}
 \mathcal D_{\mathrm{dark}}
 :=\{x\in\mathcal Q:
       Ce^{-itH_{\Abs}}x=0\text{ for every }t\in\mathbb R\}.
\end{equation}
It is the largest \(H_{\Abs}\)-reducing subspace contained in \(\ker C\).
On the ordinary unbounded-bridge cell, \(Ce^{-itH_{\Abs}}\) denotes the
closed observation orbit supplied by the Part~I realization record and the
same intersection is taken on its recorded common domain.  A domain failure
retains the typed realization value instead of defining a smaller dark
space.  Thus the dark-space coordinate is totalized by the same realization
partition as the Feshbach response.
In dimension \(d_Q<\infty\),
\begin{equation}
 \label{eq:wick-finite-dark-space}
 \mathcal D_{\mathrm{dark}}
 =\bigcap_{k=0}^{d_Q-1}\ker(CH_{\Abs}^k).
\end{equation}
The target-visible component of \(\mathcal D_{\mathrm{dark}}\) is a
\(\Caus/\Geom\) reach obstruction.  It remains available to a direct
spectral readout unless that readout is separately excluded.
\end{definition}

\begin{proposition}[Totalized imaginary-axis and bracket alternatives]
\label{prop:wick-dark-bracket-alternatives}
Define the derived geometric gap
\[
 g_*:=\inf\sigma(G)\in[0,\infty).
\]
Every imaginary-axis eigenpair \(\mathcal L_Wu=i\beta u\),
\(u=(x,y)\), satisfies
\begin{equation}
 \label{eq:wick-imaginary-kernel-branch}
 y\in\ker G,
 \qquad
 -iH_{\Abs}x+C^*y=i\beta x,
 \qquad
 -Cx+A_{\mathcal Y}y=i\beta y.
\end{equation}
On the exhaustive cell \(g_*>0\), this reduces to
\(u=(x,0)\), \(Cx=0\), and \(H_{\Abs}x=-\beta x\).  Consequently that
cell has no imaginary-axis eigenvalue exactly when the corresponding
\(H_{\Abs}\)-eigenspaces have zero intersection with \(\ker C\).  On the
complementary cell \(g_*=0\), equation
\eqref{eq:wick-imaginary-kernel-branch} retains the geometric kernel rather
than discarding it.  In both cells the commutator Gramian of
\cref{def:degenerate-geom-caus-core}, applied to
\[
 S=\begin{pmatrix}0&0\\0&G\end{pmatrix},
 \qquad
 A=\begin{pmatrix}-iH_{\Abs}&C^*\\-C&A_{\mathcal Y}\end{pmatrix},
\]
has its Part~I extended reduced gap.  A positive value supplies the existing
twisted-carrier estimate; the zero value returns its target-visible bracket
kernel as the typed dark/rank-defect branch.  These alternatives exhaust the
record and introduce no coercivity or bracket assumption.
\end{proposition}

\begin{proof}
Taking real parts of
\(\langle u,\mathcal L_Wu\rangle=i\beta\|u\|^2\) and using
\eqref{eq:wick-total-energy-balance} gives
\(\langle y,Gy\rangle=0\).  Since \(G\ge0\), this is equivalent to
\(G^{1/2}y=0\), hence to \(y\in\ker G\), and the two block equations give
\eqref{eq:wick-imaginary-kernel-branch}.  The derived inequality
\(G\ge g_*I\) makes \(y=0\) on the cell \(g_*>0\).  The remaining claims
are the two values of the Part~I bracket-resistance and twisted-carrier
alternative for the displayed \(S+A\) decomposition.
\end{proof}

\subsection{Authority-constrained structural action}
\label{subsec:wick-structural-action}

The bridge in \eqref{eq:wick-block-generator} is an operator.  Its action
record uses the actual edge or current representation induced by that same
operator.  This prevents a favorable target current from replacing an
imposed PDE current.

\begin{definition}[Wick evaluation of the inherited bridge action]
\label{def:wick-bridge-action}
Fix a pre-contact state of \(\mathscr C_I\).  Let
\((\mathfrak J,\langle\cdot,\cdot\rangle_{\mathbb M^{-1}})\), the
target-normalized potential \(D\), the actual authorized bridge current
\(J_C\), and its authority envelope \(\mathfrak A\) be the corresponding
Part~I fields.  Put
\[
 J_D:=-\mathbb M\nabla D,
 \]
and define
\begin{align}
 \mathcal A_{\Caus}(J_C)
 &:=\|J_C\|_{\mathbb M^{-1}}^2,
 &
 \mathcal E_{\mathrm{aim}}(D)
 &:=\|J_D\|_{\mathbb M^{-1}}^2,
 \label{eq:wick-action-aiming}\\
 \mathcal P_{\Caus}(J_C,D)
 &:=\langle J_C,J_D\rangle_{\mathbb M^{-1}},
 &
 \mathscr W_{\mathfrak A}(C;D)
 &:=\frac12\mathcal A_{\Caus}(J_C)
   +\frac12\mathcal E_{\mathrm{aim}}(D)
   -\mathcal P_{\Caus}(J_C,D).
 \label{eq:wick-zero-waste-action}
\end{align}
All fields and transports in this display remain those of \(\mathscr C_I\).
The Wick map changes none of them; it only permits the displayed action to be
evaluated jointly with the transformed spectral response.
\end{definition}

\begin{theorem}[Inner least-action bridge coordinate]
\label{thm:wick-least-action-law}
Every inherited bridge action evaluated after the Wick transform satisfies
\begin{equation}
 \label{eq:wick-completed-action-square}
 \mathscr W_{\mathfrak A}(C;D)
 =\frac12\|J_C-J_D\|_{\mathbb M^{-1}}^2\ge0.
\end{equation}
Let \(\mathfrak J_{\mathfrak A}\) be the possibly empty set of represented
currents induced by legal bridges, and use
\(\operatorname{dist}(J_D,\varnothing):=+\infty\).  Then
\begin{equation}
 \label{eq:wick-action-distance}
 \inf_{C\in\mathfrak A}\mathscr W_{\mathfrak A}(C;D)
 =\frac12\operatorname{dist}_{\mathbb M^{-1}}
   \bigl(J_D,\overline{\mathfrak J_{\mathfrak A}}\bigr)^2.
\end{equation}
\begin{enumerate}[label=\textup{(\roman*)}]
\item On the inherited closed-convex authority cell, the infimum is attained
  at the metric projection
  \(J_C=\Pi_{\mathfrak J_{\mathfrak A}}J_D\).  On every other nonempty cell,
  the Part~I search transcript returns an \(\varepsilon\)-minimizer and
  charges its cover and optimization error.  The empty cell returns the
  authority-obstruction value \(+\infty\).
\item Under free authority, the infimum is zero precisely when
  \(J_D\in\overline{\mathfrak J_{\mathfrak A}}\), and zero waste is attained
  precisely when \(J_D\in\mathfrak J_{\mathfrak A}\).  Thus density and exact
  representability are not conflated.
\item Under hybrid or RL authority, the same distance is evaluated on the
  declared control family.  Its convex, nonconvex, nonattained, and empty
  cells are handled by the preceding exhaustive rule.
\item Under imposed PDE authority, \(\mathfrak J_{\mathfrak A}\) is a
  singleton and the action evaluates its fixed current.
\end{enumerate}
For a declared positive target power \(p\), every legal current with
\(\mathcal P_{\Caus}(J_C,D)\ge p\) obeys
\begin{equation}
 \label{eq:wick-power-action-lower-bound}
 \mathcal A_{\Caus}(J_C)
 \ge\frac{p^2}{\mathcal E_{\mathrm{aim}}(D)}
\end{equation}
on the cell \(\mathcal E_{\mathrm{aim}}(D)>0\).  Equality holds exactly for
\(J_C=(p/\mathcal E_{\mathrm{aim}}(D))J_D\) when this current is authorized.
On the complementary cell \(\mathcal E_{\mathrm{aim}}(D)=0\), no current can
satisfy \(\mathcal P_{\Caus}(J_C,D)\ge p\); the power request therefore
returns the authority/target obstruction rather than a missing denominator.
This minimization evaluates the action coordinate
\(\mathbf c_{\rm act}(\mathscr C_W)\).  Zero action is one favorable
coordinate; structural optimality is the subsequent saturated-profile
comparison of \cref{def:wick-saturated-optimality}.
\end{theorem}

\begin{proof}
Expand the square in \eqref{eq:wick-completed-action-square}; taking the
infimum gives \eqref{eq:wick-action-distance} because distance is unchanged
by closure.  The Hilbert projection theorem gives the closed-convex cell,
while the definition of an infimum supplies \(\varepsilon\)-minimizers on
every other nonempty cell.  The authority statements are its free,
control-family, and singleton specializations.  Cauchy--Schwarz
gives
\(p^2\le\mathcal A_{\Caus}(J_C)\mathcal E_{\mathrm{aim}}(D)\), with
the displayed equality condition; when the second factor is zero, the
left-hand side cannot be positive.
\end{proof}

\begin{corollary}[Authority monotonicity]
\label{cor:wick-authority-monotonicity}
The Part~I authority order compares envelopes only at fixed quotient,
geometry, target, and cost convention.  Every ordered pair
\(\mathfrak A_1\preceq\mathfrak A_2\), equivalently
\(\mathfrak J_{\mathfrak A_1}\subseteq
\mathfrak J_{\mathfrak A_2}\), satisfies
\[
 \inf_{C\in\mathfrak A_2}\mathscr W_{\mathfrak A_2}(C;D)
 \le
 \inf_{C\in\mathfrak A_1}\mathscr W_{\mathfrak A_1}(C;D).
\]
The corresponding saturated attainable profile is ordered in the opposite
direction:
imposed \(\le\) hybrid \(\le\) free whenever the represented families are
nested with no larger implementation cost.
\end{corollary}

\subsection{Positive response kernel and quantitative transfer}
\label{subsec:wick-learning-kernel}

The quotient response \(F_W(\lambda)^{-1}\) need not be normal.  Statistical
learning therefore uses a positive transform of the full response rather
than a diagonalization of \(F_W\).

\begin{definition}[Wick response kernel]
\label{def:wick-response-kernel}
Fix \(\lambda>0\).  The complement and full resolvents, and the
invertibility of \(F_W(\lambda)\), are consequences of
\cref{thm:wick-feshbach-response} on every ordinary realization, both in the
positive-self-energy and general complement cells.  There define
\begin{equation}
 \label{eq:wick-response-kernel}
 T_{W,\lambda}
 :=\bigl(I+F_W(\lambda)^*F_W(\lambda)\bigr)^{-1},
 \qquad
 A_{W,\lambda}:=T_{W,\lambda}^{-1}
 =I+F_W(\lambda)^*F_W(\lambda).
\end{equation}
The operator \(T_{W,\lambda}\) is a positive contraction.  On a nonordinary
realization, \(T_{W,\lambda}\) is the unavailable response-kernel value and
all response-dependent upper bounds equal \(+\infty\).  For every ordinary
kernel, Part~I supplies each finite evaluation carrier \(\mathscr E\)---an empirical sample,
mixing-block sample, or sequential path/tree---has its represented evaluation
map \(E_{\mathscr E}\), with the normalization of its Part~I theorem, and
finite positive response matrix
\begin{equation}
 \label{eq:wick-empirical-response-compression}
 T_{W,\lambda}^{\mathscr E}
 :=E_{\mathscr E}T_{W,\lambda}E_{\mathscr E}^*,
 \qquad
 \operatorname{Tr}_{\mathscr E}T_{W,\lambda}
 :=\operatorname{Tr}T_{W,\lambda}^{\mathscr E}<\infty.
\end{equation}
This is the trace used by the Part~I empirical Rademacher bounds.  It is
finite because \(\mathscr E\) is finite, independently of every compactness
property of the ambient response.  Its empirical ridge effective dimension
is exactly the Part~I Gram-matrix quantity
\begin{equation}
 \label{eq:wick-empirical-effective-dimension}
 \widehat N_{W,\mathscr E}(\lambda,\rho)
 :=\operatorname{Tr}\!\left[
 T_{W,\lambda}^{\mathscr E}
 (T_{W,\lambda}^{\mathscr E}+\rho I)^{-1}\right]
 \le \dim\mathscr E.
\end{equation}

For ridge parameter \(\rho>0\), define the population effective dimension as
the extended positive trace
\begin{align}
 N_W(\lambda,\rho)
 &:=\operatorname{Tr}_{\rm ext}\!\left[
  T_{W,\lambda}(T_{W,\lambda}+\rho I)^{-1}\right]
 \notag\\
 &:=\sup_{P:\,P=P^*=P^2,\ \operatorname{rank}P<\infty}
 \operatorname{Tr}\!\left[
 P T_{W,\lambda}(T_{W,\lambda}+\rho I)^{-1}P\right]
 \in[0,\infty],
 \label{eq:wick-effective-dimension}\\
 m_W(\lambda,\rho;f)
 &:=\frac{\|T_{W,\lambda}(T_{W,\lambda}+\rho I)^{-1}f\|^2}
 {\|f\|^2}.
 \label{eq:wick-target-capture}
\end{align}
The capture ratio is defined to be zero at \(f=0\), as in the Part~I
target-capture convention.
Every represented nested dense Galerkin family \(P_n\uparrow I\) satisfies
\begin{equation}
 \label{eq:wick-effective-dimension-monotone-limit}
 N_W(\lambda,\rho)
 =\sup_n\operatorname{Tr}\!\left[
 P_nT_{W,\lambda}(T_{W,\lambda}+\rho I)^{-1}P_n\right].
\end{equation}
Thus finiteness of the population coordinate is a derived value.  When the
supremum is infinite, the effective-dimension candidate equals \(+\infty\)
and the Part~I recordwise certificate minimum uses another valid same-loss
bound or returns its existing capacity/budget obstruction.  This operation
forms the risk coordinate of the record; the saturated profile retains the
structural comparison.
\end{definition}

\begin{proposition}[Part~I trace accounting without an ambient trace gate]
\label{prop:wick-totalized-trace-accounting}
Put
\(R_{W,\rho}=T_{W,\lambda}(T_{W,\lambda}+\rho I)^{-1}\).
For every represented nested finite-rank family \(P_n\uparrow I\),
\begin{equation}
 \label{eq:wick-totalized-trace-limit}
 \operatorname{Tr}_{\rm ext}R_{W,\rho}
 =\lim_{n\to\infty}\operatorname{Tr}(P_nR_{W,\rho}P_n)
 \in[0,\infty].
\end{equation}
The value is finite exactly when \(R_{W,\rho}\) is trace class; trace-class
behavior is therefore a conclusion about this coordinate, not an
admissibility hypothesis.  Every empirical Rademacher, block, and sequential
trace remains finite on its own evaluation carrier even when the value in
\eqref{eq:wick-totalized-trace-limit} is infinite.
\end{proposition}

\begin{proof}
Choose an orthonormal basis adapted to the nested ranges of \(P_n\).  Since
\(R_{W,\rho}\ge0\), the finite traces are the increasing partial sums of
\(\sum_j\langle e_j,R_{W,\rho}e_j\rangle\).  Their limit is the extended
positive trace, independently of the adapted basis.  A positive bounded
operator is trace class precisely when this sum is finite.  The final
statement follows because every evaluation compression in
\eqref{eq:wick-empirical-response-compression} is a finite positive matrix;
\eqref{eq:wick-empirical-effective-dimension} additionally bounds its ridge
mode count by the carrier dimension.
\end{proof}

\begin{proposition}[Exhaustive modal and nonnormal response formulas]
\label{prop:wick-response-formulas}
On the positive-self-energy cell \(A_{\mathcal Y}=0\), put
\(B=\lambda I+\Sigma_C(\lambda)=B^*\).  Evaluate the derived spectral-
commutation status: its ordinary commuting value means that the spectral
measures of \(H_{\Abs}\) and \(B\) commute, and its complementary value means
that they do not.  On the commuting cell, their joint spectral representation
has scalar response
weight
\begin{equation}
 \label{eq:wick-modal-weights}
 \eta^W(\omega,\sigma)
 =\frac{1}{1+(\lambda+\sigma)^2+\omega^2};
\end{equation}
in a joint pure-point cell this reads
\(H_{\Abs}\phi_j=\omega_j\phi_j\),
\(\Sigma_C(\lambda)\phi_j=\sigma_j\phi_j\), and
\(\eta_j^W=\eta^W(\omega_j,\sigma_j)\).  On the noncommuting cell,
\begin{equation}
 \label{eq:wick-nonnormal-precision}
 A_{W,\lambda}
 =I+B^2+H_{\Abs}^2+i[B,H_{\Abs}],
\end{equation}
as a quadratic-form identity on the inherited common form domain, and the
singular values of \(F_W(\lambda)\), rather than separate eigenvalues,
determine \(T_{W,\lambda}\).

On the general cell \(A_{\mathcal Y}\ne0\), set
\(B=\lambda I+\Sigma_C(\lambda)\), which need not be self-adjoint.  The
always-valid precision identity is
\begin{equation}
 \label{eq:wick-general-nonnormal-precision}
 A_{W,\lambda}
 =I+B^*B+H_{\Abs}^2+i(B^*H_{\Abs}-H_{\Abs}B).
\end{equation}
Thus every record has an exact singular-value formula; the modal formula is
an additional ordinary value derived only on the recorded commuting
positive-self-energy cell, and its absence removes no record.
\end{proposition}

\begin{proof}
The joint spectral theorem applied on the commuting cell gives
\eqref{eq:wick-modal-weights}.  On the positive-self-energy cell,
\((B+iH_{\Abs})^*(B+iH_{\Abs})=(B-iH_{\Abs})(B+iH_{\Abs})\), which gives
\eqref{eq:wick-nonnormal-precision}.  Without self-adjointness of \(B\), the
same product expands instead to
\(B^*B+H_{\Abs}^2+i(B^*H_{\Abs}-H_{\Abs}B)\), proving
\eqref{eq:wick-general-nonnormal-precision}.
\end{proof}

\begin{theorem}[Master Wick learning transfer]
\label{thm:wick-master-learning-transfer}
For each complete Wick record, evaluate every Part~I learning-certificate
candidate already carried by that record.  In each ordinary candidate,
substitute
\begin{equation}
 \label{eq:wick-learning-dictionary}
 T_k=T_{W,\lambda},\qquad
 A=A_{W,\lambda},\qquad
 E_q=E_{W,\Delta},
\end{equation}
and retain its sample law, loss, source cell, confidence allocation, and
authority record.  Every conclusion of that Part~I candidate then holds with
\(N_{\mathrm{eff}}(\rho)=N_W(\lambda,\rho)\) and capture
\(m_k(\rho)=m_W(\lambda,\rho;f_*)\).  A candidate whose Part~I record fields
are nonordinary has value \(+\infty\).  The recordwise minimum over valid
same-loss candidates forms the Part~I statistical-radius coordinate, and the
typed failure mechanism is unchanged.  Across complete records, the inherited
saturated profile remains the comparison rule.

In particular, for the energy ball
\(\mathcal F_W(R)=\{f:\langle f,A_{W,\lambda}f\rangle\le R^2\}\),
let \(S\) be the declared \(m\)-point evaluation carrier and use its Part~I
compression \(T_{W,\lambda}^S\) from
\eqref{eq:wick-empirical-response-compression}.  Then
\begin{equation}
 \label{eq:wick-rademacher}
 \widehat{\mathfrak R}_S(\mathcal F_W(R))
 \le\frac{R}{m}\sqrt{\operatorname{Tr}_S T_{W,\lambda}},
\end{equation}
and, on the Part~I \([0,1]\)-valued \(L_\ell\)-Lipschitz loss cell, the
corresponding candidate satisfies with probability at least \(1-\delta\),
\begin{equation}
 \label{eq:wick-fixed-operator-risk}
 L(f)\le L_S(f)
 +\frac{2L_\ell R}{m}\sqrt{\operatorname{Tr}_S T_{W,\lambda}}
 +3\sqrt{\frac{\log(2/\delta)}{2m}}
 +\varepsilon_{\mathrm{task}}.
\end{equation}
On the named Part~I rate cell
\(f_*=T_{W,\lambda}^{r}g\), \(0<r\le1\), and
\(\eta_j^W\asymp j^{-2\beta}\), \(\beta>1/2\), its certified
kernel-ridge oracle inequality gives
\begin{equation}
 \label{eq:wick-krr-rate}
 \|\widehat f_{\rho_m}-f_*\|_{L^2}^2
 \lesssim m^{-4\beta r/(4\beta r+1)},
 \qquad
 \rho_m\asymp m^{-1/(2r+1/(2\beta))}.
\end{equation}
The constants and probability event are exactly those of the fixed-kernel
oracle certificate; \eqref{eq:wick-krr-rate} is asserted only for the named
eigenvalue and source class.
\end{theorem}

\begin{proof}
The operator \(T_{W,\lambda}\) is positive and self-adjoint, so the algebraic
gate of the Part~I kernel and energy-ball candidates is derived rather than
assumed.
Its empirical compression \(T_{W,\lambda}^S\) is a finite positive matrix,
so the Part~I ellipsoid support calculation gives
\eqref{eq:wick-rademacher}, and
symmetrization and contraction give
\eqref{eq:wick-fixed-operator-risk}.  The source residual is
\(O(\rho^{2r})\), while
\(N_W(\lambda,\rho)\asymp\rho^{-1/(2\beta)}\); balancing variance and bias
gives \eqref{eq:wick-krr-rate}.  This named eigenvalue class makes the
extended effective dimension finite as a conclusion.  The remaining
transfers are substitutions in positive-operator statements and preserve
their recorded cells and units.
\end{proof}

\begin{proposition}[Learned-operator stability]
\label{prop:wick-learned-operator-stability}
On every finite Part~I evaluation carrier, let two represented
positive-normal-form records have the derived error coordinates
\[
 \varepsilon_H:=\|H-\widehat H\|,
 \qquad
 \varepsilon_C:=\|C-\widehat C\|,
 \qquad
 \varepsilon_G:=\|G-\widehat G\|.
\]
Put
\(r=\max\{\|(\lambda+G)^{-1}\|,
\|(\lambda+\widehat G)^{-1}\|\}\) and
\(c=\max\{\|C\|,\|\widehat C\|\}\).  Then
\begin{align}
 \|F_W-\widehat F_W\|
 &\le\varepsilon_H+2cr\varepsilon_C+c^2r^2\varepsilon_G
 =:\varepsilon_F,
 \label{eq:wick-feshbach-perturbation}\\
 \|T_{W,\lambda}-\widehat T_{W,\lambda}\|
 &\le(\|F_W\|+\|\widehat F_W\|)\varepsilon_F.
 \label{eq:wick-kernel-perturbation}
\end{align}
Here \(r\le\lambda^{-1}\) is derived from positivity, so every finite-carrier
quantity in the display is finite.  On a continuum record, the same formulas
are evaluated with its Part~I form-to-resolvent error coordinates in
\([0,\infty]\); a nonordinary comparison gives the valid value \(+\infty\)
rather than an admissibility restriction.
The general nonnormal branch has the same formula with
\(\varepsilon_G\) replaced by the derived same-form-norm distance between
\((G-A_{\mathcal Y})-(\widehat G-\widehat A_{\mathcal Y})\) and with both
complement resolvent norms retained.
\end{proposition}

\begin{proof}
Add and subtract the two mixed self-energy terms and use
\(R-\widehat R=R(\widehat G-G)\widehat R\) for
\(R=(\lambda+G)^{-1}\).  This proves
\eqref{eq:wick-feshbach-perturbation}.  Apply the resolvent identity to
\((I+F^*F)^{-1}\) and use
\(\|F^*F-\widehat F^*\widehat F\|
\le(\|F\|+\|\widehat F\|)\|F-\widehat F\|\).
\end{proof}

\subsection{Statistically calibrated training objectives}
\label{subsec:wick-training-objectives}

The response kernel determines the regularizer, but it does not by itself
specify how the block generator is identified or how a selectable bridge is
trained.  The three inherited authority/data cells therefore evaluate three
objectives: a dynamics loss on target-free data, a response-kernel loss on
supervised data, and, for controlled systems, a Bellman/policy loss on the
represented deployment law.  Their samples and confidence allocations are
fields of the Part~I record.

\begin{definition}[Structure-preserving dynamics loss]
\label{def:wick-dynamics-training-loss}
On each represented finite-dimensional chart supplied by the Part~I learner,
quotienting by its recorded carrier gauge gives identifiable coordinates. Let
\(Z_j=(x_j,y_j,\dot x_j,\dot y_j)\), \(1\le j\le n_{\rm dyn}\),
be dynamics-only observations, and let
\(W_{x,j}\) and \(W_{y,j}\) be represented positive precision matrices.
For
\(\vartheta=(H,G,A_{\mathcal Y},C)\) satisfying
\[
 H=H^*\succeq0,\qquad G=G^*\succeq0,\qquad
 A_{\mathcal Y}^*=-A_{\mathcal Y},
\]
define
\begin{equation}
 \label{eq:wick-dynamics-training-loss}
 \mathcal L_{\rm dyn}(\vartheta)
 :=\frac1{2n_{\rm dyn}}\sum_{j=1}^{n_{\rm dyn}}
 \left(
 \|\dot x_j+iHx_j-C^*y_j\|_{W_{x,j}}^2
 +\|\dot y_j+Cx_j+Gy_j-A_{\mathcal Y}y_j\|_{W_{y,j}}^2
 \right).
\end{equation}
This is the generalized least-squares loss for the two block equations.  On
the inherited conditional-Gaussian channel cell with inverse covariance
\(W_{x,j}\oplus W_{y,j}\), it is also the negative log-likelihood up to a
parameter-independent constant.

When derivatives are unavailable, raw transitions use the likelihood loss
\begin{equation}
 \label{eq:wick-transition-training-loss}
 \mathcal L_{\rm tr}(\vartheta)
 :=\frac1{2n_{\rm dyn}}\sum_j
 \|u_{j+1}-e^{\Delta_j\mathcal L_W(\vartheta)}u_j\|_{\Omega_j^{-1}}^2
 +\frac1{2n_{\rm dyn}}\sum_j\log\det\Omega_j,
\end{equation}
with the conditional mean and covariance replaced by the represented
transition law when it is non-Gaussian.  A discretization defect, derivative
estimation error, or misspecified noise covariance is retained as a separate
model coordinate.  Neither loss uses target labels or rewards.
The fitted parameters and \(\varepsilon_{\rm opt}\) are inner training
coordinates in the sense of \cref{def:wick-saturated-optimality}.  Comparison
between complete records is deferred to the saturated profile.
\end{definition}

\begin{proposition}[Identification and response error]
\label{prop:wick-dynamics-identification}
In the linear derivative chart of
\cref{def:wick-dynamics-training-loss}, write
\[
 r_j(\vartheta_*+h)=\xi_j+\mathcal D_jh,
 \qquad
 Q_n=\frac1{n_{\rm dyn}}\sum_j
 \mathcal D_j^*W_j\mathcal D_j,
 \qquad
 s_n=\frac1{n_{\rm dyn}}\sum_j\mathcal D_j^*W_j\xi_j,
\]
where \(W_j=W_{x,j}\oplus W_{y,j}\), and where \(\vartheta_*\) is any
represented comparator; its residual \(\xi_j\) is retained, so realizability
is not presumed.  Define from the observed design
\[
 \kappa_{\rm id}:=\lambda_{\min}(Q_n),\qquad
 P_{\rm id}:=\mathbf1_{(0,\infty)}(Q_n),
\]
and let \(\kappa_{\rm id}^+\) be the least positive eigenvalue of \(Q_n\),
with \(\kappa_{\rm id}^+=+\infty\) when \(Q_n=0\).  Every represented
\(\varepsilon_{\rm opt}\)-minimizer compared with \(\vartheta_*\) obeys
\begin{equation}
 \label{eq:wick-identification-bound}
 \|P_{\rm id}(\widehat\vartheta-\vartheta_*)\|
 \le
 \frac{2\|s_n\|}{\kappa_{\rm id}^+}
 +\sqrt{\frac{2\varepsilon_{\rm opt}}{\kappa_{\rm id}^+}}.
\end{equation}
On the derived cell \(\kappa_{\rm id}>0\), \(P_{\rm id}=I\) and this is the
full-parameter bound.  On the cell \(\kappa_{\rm id}=0\),
\(\ker Q_n\) is exactly the unidentifiable design/gauge direction and remains
an obstruction coordinate; no excitation hypothesis is imposed.

One Part~I score-concentration candidate is the conditional sub-Gaussian
cell.  On that recorded cell, in a \(p_{\rm id}\)-dimensional chart, the
design certificate gives, with probability at least
\(1-\delta_{\rm id}\),
\begin{equation}
 \label{eq:wick-identification-concentration}
 \|s_n\|
 \le c_{\rm id}\sigma_{\rm id}
 \sqrt{\frac{p_{\rm id}+\log(1/\delta_{\rm id})}{n_{\rm dyn}}}.
\end{equation}
The constant \(c_{\rm id}\), the score norm, and the sampling/dependence
certificate are fields of the record.  Off that cell, the Part~I minimum uses
its mixing, martingale, or finite-description same-loss candidate; this is a
recordwise radius coordinate, while the saturated profile compares records.
An unavailable
score radius equals \(+\infty\).  Componentwise application of
\cref{prop:wick-learned-operator-stability} converts
\eqref{eq:wick-identification-bound} into certified bounds for
\(F_W\), \(T_{W,\lambda}\), effective dimension, capture, and every
operator-Lipschitz training coordinate.
\end{proposition}

\begin{proof}
The loss difference is a quadratic with Hessian \(Q_n\) and linear score
\(s_n\).  Moreover \(s_n\perp\ker Q_n\), because every vector in
\(\ker Q_n\) is killed by each weighted design map.  Writing
\(h=\widehat\vartheta-\vartheta_*\), comparison with the represented
comparator and the optimization tolerance gives
\(\langle h,Q_nh\rangle/2
 \le \|s_n\|\|P_{\rm id}h\|+\varepsilon_{\rm opt}\).
The spectral inequality
\(Q_n\succeq\kappa_{\rm id}^+P_{\rm id}\) and the quadratic formula imply
\eqref{eq:wick-identification-bound}.  The kernel component does not enter
the loss and is therefore exactly unidentifiable.  The concentration display
is the Part~I sub-Gaussian score candidate on its named cell.  The last
assertion follows
from \eqref{eq:wick-feshbach-perturbation}--
\eqref{eq:wick-kernel-perturbation} and the response-coordinate Lipschitz
bounds in \cref{cor:wick-quantitative-ledger}.
\end{proof}

\begin{definition}[Unrestricted-Caus spectral training loss]
\label{def:wick-free-caus-training-loss}
For a legal bridge \(C\), let
\[
 \mathcal H_{W,C}:=\operatorname{ran}T_{W,\lambda}(C)^{1/2},
 \qquad
 \|f\|_{W,C}^2
 :=\|T_{W,\lambda}(C)^{\dagger/2}f\|^2.
\]
On a finite represented score space this norm is
\(\langle f,A_{W,\lambda}(C)f\rangle\).  Given a target-training split
\(S_{\rm tr}=\{(X_i,Y_i)\}_{i=1}^{m_{\rm tr}}\), define the spectral
kernel-ridge estimator
\begin{equation}
 \label{eq:wick-free-inner-loss}
 \widehat f_{C,\rho}
 :=\operatorname*{arg\,min}_{f\in\mathcal H_{W,C}}
 \left\{
 \frac1{2m_{\rm tr}}\sum_{i=1}^{m_{\rm tr}}|f(X_i)-Y_i|^2
 +\frac\rho2\|f\|_{W,C}^2
 \right\}.
\end{equation}
Thus the correct spectral penalty is the closed quadratic form of the
\emph{inverse} response, equivalently \(A_{W,\lambda}\) on its form domain,
rather than \(T_{W,\lambda}\) itself.
On the ordinary Part~I evaluation cell, the represented evaluation map is
bounded on \(\mathcal H_{W,C}\); the positive ridge term therefore makes
\eqref{eq:wick-free-inner-loss} coercive in its response norm and gives its
unique minimum.  On a finite score carrier this is the positive-definite
normal equation in \eqref{eq:wick-empirical-krr-normal-equation}.  An
unavailable evaluation map or nonordinary form retains its inherited value
and makes the corresponding upper-risk candidate \(+\infty\).

Let \(\mathfrak A_{\rm dyn}\) be the dynamics confidence set obtained from
an independent target-free split.  For a designable control component it is
the full declared bridge family; for a native or imposed component it is the
identified confidence set.  On an independent validation split define
\begin{equation}
 \label{eq:wick-free-outer-loss}
 \boxed{
 \widehat{\mathcal J}_{\rm free}(C,\rho)
 =\widehat R_{\rm val}(\widehat f_{C,\rho})
 +\lambda_{\mathscr W}\mathscr W_{\mathfrak A}(C;D)
 +\operatorname{pen}_{\rm sel}(C,\rho)
 +\iota_{\mathfrak A_{\rm dyn}}(C),
 }
\end{equation}
where \(\iota\) is zero on the confidence set and \(+\infty\) outside it.
The selection penalty is a valid bridge/kernel selection deviation obtained
from a split sample, finite union, PAC--Bayes mixture, or uniform cover.
Operator estimation and optimization are recorded separately, so no radius
is charged twice.  Under
unrestricted authority the zero-waste fiber
\begin{equation}
 \label{eq:wick-zero-waste-fiber}
 \mathfrak A_D^0:=\{C\in\mathfrak A_{\rm dyn}:J_C=J_D\}
\end{equation}
is evaluated rather than presumed nonempty.  It is nonempty exactly when the
represented legal family realizes \(J_D\).  On that cell, exact least-action
training minimizes \eqref{eq:wick-free-outer-loss} over
\(\mathfrak A_D^0\).  On the empty cell, the distance
\eqref{eq:wick-action-distance} is the irreducible representation/authority
gap.  The penalized form trades statistical risk against physical action
when such a trade is part of the declared task.
The kernel-ridge fit and bridge search are inner candidate-generating
optimizations.  Their risk, action, cover charge, and optimization error enter
the complete certificate; the saturated profile and its strict inverse select
among the resulting structural records according to
\cref{def:wick-saturated-optimality}.
\end{definition}

\begin{theorem}[Oracle bound for unrestricted-Caus learning]
\label{thm:wick-free-caus-oracle}
Consider the Part~I split-validation candidate of a complete record.  Its
fields are a charged finite \(\varrho\)-net
\(\mathfrak A_{\varrho}\) of the declared legal-bridge and ridge family
\(\mathfrak A_{\rm dyn}\times\mathcal P\), of cardinality
\(N_{\varrho}\), the recorded loss range \(L_{\rm loss}\), the independent
validation carrier of size \(m_{\rm val}\), and the simultaneous fixed-kernel
oracle values
\begin{equation}
 \label{eq:wick-free-candidate-bound}
 \mathcal B_C(\rho)
 :=C_{\rm var}\frac{\sigma^2N_W^C(\lambda,\rho)}{m_{\rm tr}}
 +C_{\rm bias}B_C(\rho;f_*)
 +\varepsilon_{\rm task}
 +L_{\rm op}\varepsilon_T(C),
\end{equation}
where \(N_W^C\) has the extended meaning of
\eqref{eq:wick-effective-dimension}.  Thus \(\mathcal B_C(\rho)\) is a
well-defined value in \([0,\infty]\); a divergent mode count removes this
candidate from the recordwise finite oracle minimum rather than excluding the
model from the framework.  The bias term is
\[
 B_C(\rho;f_*)
 =\|[I-T_C(T_C+\rho I)^{-1}]f_*\|_{L^2}^2.
\]
Let
\begin{equation}
 \label{eq:wick-free-selection-radius}
 r_{\rm val}
 :=L_{\rm loss}
 \sqrt{\frac{\log(2N_{\varrho}/\delta_{\rm val})}{2m_{\rm val}}}.
\end{equation}
Take \(\operatorname{pen}_{\rm sel}=r_{\rm val}\) on this net; a
candidate-dependent valid radius gives the analogous oracle bound with that
radius inside the infimum.  Let \(L_{\rm sel}\in[0,\infty]\) be the Part~I
net-transport modulus of candidate risk plus action, and let
\(\varepsilon_{\rm out}\in[0,\infty]\) be the recorded optimization gap of
the returned pair \((\widehat C,\widehat\rho)\).  On the candidate's
simultaneous event,
\begin{align}
 \label{eq:wick-free-oracle-bound}
 &R(\widehat f_{\widehat C,\widehat\rho})-R(f_*)
 +\lambda_{\mathscr W}\mathscr W_{\mathfrak A}(\widehat C;D)
 \notag\\
 &\qquad\le
 \inf_{C\in\mathfrak A_{\rm dyn},\,\rho\in\mathcal P}
 \left\{\mathcal B_C(\rho)
 +\lambda_{\mathscr W}\mathscr W_{\mathfrak A}(C;D)\right\}
 +2r_{\rm val}+L_{\rm sel}\varrho+\varepsilon_{\rm out}.
\end{align}
An absent split, loss-range, simultaneous-oracle, or finite net-transport
field sets this candidate to \(+\infty\); the inherited Rademacher and
PAC--Bayes candidates remain in the Part~I same-loss certificate minimum.
The choice \(\mathcal P=(0,\infty)\) is permitted only on a Part~I record
carrying a charged cover and transport modulus for that continuum; otherwise
\(\mathcal P\) is exactly its represented finite ridge grid.
That minimum supplies the risk coordinate subsequently restricted in the
saturated profile.  On the cell
\(\mathfrak A_D^0\ne\varnothing\), restriction to that fiber removes the
action term.  Within its named eigenvalue/source rate cell
\(\eta_j^W(C)\asymp j^{-2\beta_C}\) and
\(f_*=T_C^{r_C}g_C\), and whose charged ridge family contains the balancing
scale or a net point with the displayed transport error, optimization of
\(\rho\) gives the certified inner rate
\(m_{\rm tr}^{-4\beta_Cr_C/(4\beta_Cr_C+1)}\) at that fixed record,
plus the displayed bridge-selection and operator errors.  Outside that rate
cell, \eqref{eq:wick-free-oracle-bound} remains the quantitative conclusion.
\end{theorem}

\begin{proof}
Hoeffding's inequality and a union bound give a validation deviation at most
\(r_{\rm val}\) for every net point.  Insert the validation risk on both
sides of the approximate-minimizer inequality.  The two deviations,
net-transport term, and optimization error give
\eqref{eq:wick-free-oracle-bound}; then substitute the simultaneous
candidate oracle certificates.  The zero-waste statement is
\cref{thm:wick-least-action-law}.  The rate follows by balancing
\(\rho^{2r_C}\) with
\(m_{\rm tr}^{-1}\rho^{-1/(2\beta_C)}\).
\end{proof}

\begin{definition}[Action-safe Wick actor--critic loss]
\label{def:wick-rl-training-loss}
Let \(P_{\mathcal A}\) project onto the represented legal-action span and
let \(P_{\rm safe}\) project onto its failure-killed safe subspace.  Define
\begin{equation}
 \label{eq:wick-action-safe-kernel}
 T_C^{\mathcal A,{\rm safe}}
 :=\left.P_{\rm safe}P_{\mathcal A}T_{W,\lambda}(C)
 P_{\mathcal A}P_{\rm safe}\right|_{\operatorname{ran}P_{\rm safe}},
 \qquad
 N_C^{\mathcal A,{\rm safe}}(\rho)
 :=\operatorname{Tr}_{\rm ext}\!\left[
 T_C^{\mathcal A,{\rm safe}}
 (T_C^{\mathcal A,{\rm safe}}+\rho I)^{-1}\right].
\end{equation}
The trace has the extended Part~I meaning of
\eqref{eq:wick-effective-dimension}; every empirical critic calculation uses
its finite evaluation compression from
\eqref{eq:wick-empirical-response-compression}.
The Part~I control record supplies a real-valued observable-score map.  Thus
every critic \(q\) below ranges over its inherited real score slice inside the
complexified response space.  Hermitian response energies are restricted to
that slice, while maxima, exponentials, rewards, and policy probabilities are
real-valued.  This is the existing control-score type, not an additional
regularity condition on the Wick response.
For an executed bridge action \(a\mapsto C(h,a)\), put
\[
 w_D(h,a):=\frac12
 \|J_{C(h,a)}-J_D(h)\|_{\mathbb M^{-1}}^2.
\]
On the action-charged authority cell, its stage reward is
\(r_\alpha(h,a)=r(h,a)-\alpha w_D(h,a)\).  On the hard-action-constraint
cell, the legal action set carries that constraint and \(\alpha=0\); the
same action is not charged a second time.

At fitted Bellman iteration \(k\), freeze a copy \(q_k^-\) of the preceding
critic.  A delayed target network retains the lag
\(\varepsilon_{{\rm lag},k}:=\|q_k^--\widehat q_k\|_\infty\).  Set
\begin{equation}
 \label{eq:wick-fitted-bellman-target}
 Y_i^{(k)}
 =r_\alpha(h_i,a_i)
 +\chi_i^{\rm cont}\gamma
  \max_{a'\in\mathcal A_{\rm safe}(h_i')}
 q_k^-(h_i',a').
\end{equation}
Here \(\chi_i^{\rm cont}\in\{0,1\}\) is the terminal-continuation field of
the complete Part~I dynamical record.  Before taking the maximum, that record
either certifies a nonempty safe successor row or returns its existing
authority/terminal status; no maximum over an empty set is inserted into the
loss.
For represented behavior-to-deployment weights \(0\le\omega_i\le\bar\omega\),
the critic loss is
\begin{equation}
 \label{eq:wick-rl-critic-loss}
 \mathcal L_{\rm crit}^{(k)}(q;C)
 :=\frac1{2m_k}\sum_{i=1}^{m_k}\omega_i
 (q(h_i,a_i)-Y_i^{(k)})^2
 +\frac\rho2
 \|(T_C^{\mathcal A,{\rm safe}})^{\dagger/2}q\|^2.
\end{equation}
Thus critic regularization again uses the inverse positive response.  The
frozen target makes \eqref{eq:wick-rl-critic-loss} a supervised regression
loss and avoids the double-sampling defect of a same-sample squared Bellman
residual.

For the neutral action row \(\nu(a\mid h)\) carried by the Part~I control
record, define on its represented support
\begin{equation}
\label{eq:wick-rl-gibbs-policy}
\begin{aligned}
 \pi_q^\tau(a\mid h)
 &=\frac{\nu(a\mid h)e^{q(h,a)/\tau}}
 {\sum_{a'\in\mathcal A_{\rm safe}(h)}
   \nu(a'\mid h)e^{q(h,a')/\tau}},
 &&a\in\mathcal A_{\rm safe}(h),\\
 \mathcal L_{\rm actor}(\pi;q)
 &=\mathbb E_{d_{\rm dep}}
 D_{\rm KL}(\pi(\cdot\mid h)\Vert\pi_q^\tau(\cdot\mid h)).
\end{aligned}
\end{equation}
The policy is zero outside \(\mathcal A_{\rm safe}(h)\).  The displayed
denominator is positive on the ordinary Part~I support cell.  If the neutral
row has zero total safe mass, the Gibbs row and actor certificate take their
typed support-obstruction values.  A zero mass on any required safe action is
retained by the support coordinate \(\nu_{\min}\) in
\cref{thm:wick-rl-loss-calibration}.
The canonical alternating objective is
\begin{equation}
 \label{eq:wick-rl-composite-loss}
 \boxed{
 \mathcal L_{W,{\rm RL}}^{(k)}(q,\pi,C)
 =\mathcal L_{\rm crit}^{(k)}(q;C)
 +\lambda_\pi\mathcal L_{\rm actor}(\pi;q)
 +\lambda_{\rm dyn}\mathcal L_{\rm dyn}(\vartheta)
 +\lambda_{\rm num}\mathcal L_{\rm num}(q,\pi,C).
 }
\end{equation}
The critic is fitted first with \(q_k^-\) fixed; the actor is then fitted to
the new critic.  Importance, occupancy, filter, numerical, bridge-selection,
and target-network costs belong to the complete RL record.
The critic, actor, and bridge updates are inner optimization transcripts.
Their completed record enters the saturated profile, which supplies the
structural order across policies, responses, and authority cells.
\end{definition}

\begin{theorem}[RL critic and policy calibration]
\label{thm:wick-rl-loss-calibration}
Work on the inherited discounted or contractive Part~I control cell, with
recorded modulus \(0\le\gamma<1\).  A finite-horizon or noncontractive
control record retains its own Part~I dynamic-programming candidate and does
not use the denominators below.
For each fitted critic step, let \(e_k^2\in[0,\infty]\) be the recordwise
minimum of the applicable simultaneous Part~I critic-oracle candidates.  It
is the critic-radius coordinate of that fixed RL record; structural comparison
uses the saturated profile.  Its Wick
effective-dimension candidate is
\begin{equation}
 \label{eq:wick-rl-critic-oracle}
 \begin{aligned}
 \mathcal E_{k,{\rm KRR}}
 &:=\frac{C_{\rm var}\sigma_{{\rm TD},k}^2
 N_{C_k}^{\mathcal A,{\rm safe}}(\rho_k)}{m_k}
 +C_{\rm bias}B_{C_k}^{\mathcal A,{\rm safe}}(\rho_k)\\
 &\quad+\varepsilon_{{\rm shift},k}
 +\varepsilon_{{\rm model},k}
 +\varepsilon_{{\rm sel},k}
 +\varepsilon_{{\rm opt},k},
 \qquad e_k^2\le \mathcal E_{k,{\rm KRR}}.
 \end{aligned}
\end{equation}
where the safe source residual includes
\(\|(I-P_{\rm safe})Q_*\|^2\).  An unavailable candidate has value
\(+\infty\), so this definition excludes no record.

Let \(\mathcal K_{\infty,k}\) be the set of same-law conversion constants
certified by the Part~I record and define the totalized conversion error
\begin{equation}
 \label{eq:wick-rl-sup-converter}
 \varepsilon_{\infty,k}
 :=\inf_{\kappa\in\mathcal K_{\infty,k}}\kappa e_k
 \in[0,\infty],
 \qquad \inf\varnothing:=+\infty.
\end{equation}
Every ordinary member of \(\mathcal K_{\infty,k}\) certifies
\(\|\widehat q_{k+1}-\mathcal T_{B,\alpha}q_k^-\|_\infty
\le\kappa e_k\).  In the extended nonnegative reals, the
restricted regularized Bellman target therefore satisfies
\begin{equation}
 \label{eq:wick-rl-fqi-recursion}
 \Gamma_{Q,K}:=\|\widehat q_K-Q_{B,\alpha}^*\|_\infty
 \le\gamma^K\Gamma_{Q,0}
 +\sum_{k=0}^{K-1}\gamma^{K-1-k}
 \bigl(\varepsilon_{\infty,k}
       +\gamma\varepsilon_{{\rm lag},k}\bigr).
\end{equation}
Define two further derived record coordinates
\[
 \nu_{\min}:=\inf_{h,a\in\mathcal A_{\rm safe}(h)}\nu(a\mid h),
 \qquad
 \kappa_\pi:=\sup_hD_{\rm KL}(\widehat\pi(\cdot\mid h)
 \Vert\pi_{\widehat q_K}^\tau(\cdot\mid h))\in[0,\infty],
\]
and use \(\log(1/0):=+\infty\).  The actor's extended one-step action
certificate is
\begin{equation}
 \label{eq:wick-rl-action-defect}
 d_{\rm act}
 :=2\Gamma_{Q,K}
 +\tau\log\frac1{\nu_{\min}}
 +\operatorname{Osc}(\widehat q_K)\sqrt{\frac{\kappa_\pi}{2}},
\end{equation}
and
\begin{equation}
 \label{eq:wick-rl-policy-oracle}
 \|V_{B,\alpha}^*-V_{\alpha}^{\widehat\pi}\|_\infty
 \le\frac{d_{\rm act}}{1-\gamma}.
\end{equation}
These displays are finite exactly on the Part~I same-law/support/actor cells
that make all their coordinates finite.  The value for the unrestricted
control problem additionally includes the
represented budget, safe-action, and observation/partial-information gaps.
When the same-law conversion set is empty,
\(\varepsilon_{\infty,k}=+\infty\):
the \(L^2(d_{\rm dep})\) critic certificate remains valid, whereas the
sup-norm Bellman and policy certificates take the honest value \(+\infty\).
\end{theorem}

\begin{proof}
The fixed-target critic is kernel ridge regression with kernel
\(T_C^{\mathcal A,{\rm safe}}\), so its effective-dimension oracle
certificate is \eqref{eq:wick-rl-critic-oracle}.  Apply the converter,
insert the target-network lag, and iterate the \(\gamma\)-contraction to obtain
\eqref{eq:wick-rl-fqi-recursion}.  The Gibbs variational identity bounds its
soft-greedy defect by \(\tau\log(1/\nu_{\min})\); Pinsker's inequality gives
the final actor term in \eqref{eq:wick-rl-action-defect}, and replacement of
the true critic by \(\widehat q_K\) costs \(2\Gamma_{Q,K}\).  The discounted
performance-difference identity sums the uniform one-step defect and gives
\eqref{eq:wick-rl-policy-oracle}.
\end{proof}

\begin{corollary}[Action-safe RL rate cell]
\label{cor:wick-rl-learning-rate}
On the named Part~I rate cell in which the nonzero eigenvalues of
\(T_C^{\mathcal A,{\rm safe}}\) obey
\(\eta_j\asymp j^{-2\beta}\), \(\beta>1/2\), the Bellman-optimal critic of
that fixed RL record has source order \(0<r\le1\), the safe approximation residual is
zero, the charged ridge family contains the balancing scale or a net point
with no larger transport error, and the shift, model, selection, and
optimization coordinates have no larger order, the general oracle value
specializes to
\begin{equation}
 \label{eq:wick-rl-rate}
 e_k^2=O\!\left(m_k^{-4\beta r/(4\beta r+1)}\right),
 \qquad
 e_k=O\!\left(m_k^{-2\beta r/(4\beta r+1)}\right).
\end{equation}
On the subcell carrying uniformly bounded conversion constants
\(\kappa_{\infty,k}\in\mathcal K_{\infty,k}\),
the fitted-iteration residual has converged, \(\tau\to0\), and the derived
\(\kappa_\pi\to0\), the policy gap has order
\begin{equation}
 \label{eq:wick-rl-policy-rate}
 O\!\left(
 \frac{m^{-2\beta r/(4\beta r+1)}}{(1-\gamma)^2}
 \right),
\end{equation}
in addition to the exact budget, safe-action, and observation approximation
gaps.  Every complementary cell retains the extended oracle
\eqref{eq:wick-rl-fqi-recursion}--\eqref{eq:wick-rl-policy-oracle}; it is not
excluded.  A direct sup-norm critic oracle may replace the square-root
conversion and yields its own certified rate.  Bellman optimality here is an
inner control coordinate of the fixed RL record; the risk-qualified saturated
profile supplies the structural optimum across RL records.
\end{corollary}

\begin{corollary}[Specialized quantitative ledger]
\label{cor:wick-quantitative-ledger}
The Part~I confidence allocator constructs one simultaneous event for every
ordinary candidate below.  On that event the following bounds specialize to
a Wick certificate as indicated; every nonordinary candidate is assigned
\(+\infty\) and remains in the record.
\begin{enumerate}[label=\textup{(\roman*)}]
\item On each nonempty represented Part~I stage, let
\(\{\mathcal H_\alpha^W\}_{\alpha\in\mathfrak A_W^{\rm src}}\) be its exact
source cells.  For the recorded \([-1,1]\)-loss candidate, source selection
obeys
\begin{equation}
 \label{eq:wick-source-cell-union}
 \widehat{\mathfrak R}_S
 \left(\bigcup_\alpha\mathcal H_\alpha^W\right)
 \le\max_\alpha\widehat{\mathfrak R}_S(\mathcal H_\alpha^W)
 +\sqrt{\frac{2\log|\mathfrak A_W^{\rm src}|}{m}}.
\end{equation}
For a cell-mixture prior and posterior,
\begin{equation}
 \label{eq:wick-source-cell-pac-bayes}
 \operatorname{KL}(Q\Vert P)
 =\operatorname{KL}(q\Vert w)
 +\sum_\alpha q_\alpha\operatorname{KL}(Q_\alpha\Vert P_\alpha).
\end{equation}
Thus the selection of a Wick source, bridge family, or authority cell is
charged once by a union, mixture, split-sample, or uniform-cover certificate.
\item For every fixed bridge and ridge scale, the complete kernel-ridge
oracle candidate is
\begin{equation}
 \label{eq:wick-krr-oracle}
 \|\widehat f_{C,\rho}-f_*\|_{L^2}^2
 \le C_{\rm var}\frac{\sigma^2N_W^C(\lambda,\rho)}m
 +C_{\rm bias}B_C(\rho;f_*).
\end{equation}
The source rate \eqref{eq:wick-krr-rate}, a minimax lower rate, or a
condition-number lower capacity bound transfers only on its declared
eigenvalue/source class.  In particular, on an \(m\)-point score space,
\begin{equation}
 \label{eq:wick-condition-capacity-lower}
 \operatorname{TaskCap}_{A_W}(y)
 \ge\frac{\|y\|_2}{\sqrt{m\,\kappa(A_W)}}.
\end{equation}
Dark, target-null, and authority-excluded components remain exact
approximation floors rather than statistical terms.
\item On the inherited moving-bridge/feature cell, a terminal kernel
\(T_{\widehat\theta}\) selected from its charged \(\varrho\)-net has
feature-learning certificate
\begin{align}
 \label{eq:wick-feature-learning-bound}
 \operatorname{ExcessRisk}
 \le{}&C_{\rm var}\frac{\sigma^2
 N_W^{\widehat\theta}(\lambda,\rho)}m
 +C_{\rm bias}B_{\widehat\theta}(\rho;f_*)\notag\\
 &+r_{\rm sel}+L_{\rm sel}\varrho
 +\varepsilon_{\rm est}+\varepsilon_{\rm opt},
\end{align}
where, for an independent bounded-loss audit of size \(m_a\), one may take
\[
 r_{\rm sel}=L_{\rm loss}
 \sqrt{\frac{\log(2\mathcal N(\Theta,\varrho)/\delta)}{2m_a}}.
\]
The fixed-response/lazy regime is the specialization with no kernel movement;
a neural or nonlinear readout adds its own represented margin or covering
certificate and is combined only in common risk units.
\item A learned family \(\{T_{W,\lambda}(\theta)\}\) uses
\eqref{eq:wick-kernel-perturbation} in the learned-operator covering bound.
Since \(A_{W,\lambda}\ge I\), its selection radius is
\begin{equation}
 \label{eq:wick-learned-rademacher}
 \inf_{\varrho>0}\left[
 \frac Rm\sqrt{\sup_\theta\operatorname{Tr}_S T_\theta}
 +R\sqrt{\frac{L_T\varrho}{m}}
 +\frac Rm\sqrt{2\log\mathcal N(\Theta,\varrho)}
 \right].
\end{equation}
Here \(\varrho\) is the charged cover resolution, not the ridge scale
\(\rho\).
\item On the represented positive defect/bracket range \(R_{E,W}\), the
proper Gaussian structural prior has precision
\(\mathbf M_{E,W}=P_{E,W}A_{W,\lambda}P_{E,W}|_{R_{E,W}}\).  Its PAC--Bayes
radius is
\begin{equation}
 \label{eq:wick-pac-bayes}
 \sqrt{\frac{\mathrm{KL}_{E,W}^{\mathrm{str}}(Q)
 +\log(2\sqrt m/\delta)}{2(m-1)}}.
\end{equation}
This candidate is ordinary for \(m\ge2\); at \(m<2\) it equals
\(+\infty\).
Null and dark directions retained by the model receive a separate proper
label-independent prior and KL charge.
\item The Part~I resistance-query candidate uses
\(\langle E_{W,\Delta},Q_{K,S}^{+}E_{W,\Delta}\rangle\).  Rayleigh
monotonicity gives monotonicity under grounding.  The existing submodularity
coordinate either certifies the \(1-1/e\) greedy batch factor or assigns that
factor-candidate \(+\infty\); the resistance query itself remains available.
\item Put \(\varepsilon_T:=\|\widehat T-T\|\), let \(g_*\ge0\) be the
derived target-band separation, and define
\[
 d_{\rm DK}:=
 \begin{cases}
  2\varepsilon_T/g_*,&g_*>2\varepsilon_T,\\
  1,&g_*\le2\varepsilon_T.
 \end{cases}
\]
Davis--Kahan gives the first value, while the universal distance bound for
orthogonal projectors gives the second
\citep{DavisKahan1970}.  Thus the projector discrepancy is always bounded by
\(d_{\rm DK}\), without a spectral-gap assumption.  Smoothed bands use the
recordwise minimum of this value and their represented operator-Lipschitz
value.  The
result propagates to capture, effective dimension, action-safe capacity, and
every declared Lipschitz certificate coordinate.
\item On the homogeneous symmetric label-noise cell \(\eta<1/2\), every
clean-risk radius for a flip-symmetric bounded loss is divided by
\(\chi_\eta:=1-2\eta\).  For squared loss the debiased
Wick estimator satisfies
\begin{equation}
 \label{eq:wick-debiased-krr}
 \|\widehat f_{C,\rho}-f_*\|_{L^2}^2
 \le C_{\rm var}\frac{v_\eta^2N_W^C(\lambda,\rho)}
 {\chi_\eta^2m}
 +C_{\rm bias}B_C(\rho;f_*).
\end{equation}
Wick source ancestry and operator-selection costs are unchanged.  General
reward, transition, action, and observation corruption enters through its
typed same-norm channel converter, not through the binary factor.
\item On the Part~I stationary \(\beta\)-mixing cell of length
\(N\), block length \(a\), and
\(\mu_a=\lfloor N/(2a)\rfloor\), put
\(\delta_a=\delta-2(\mu_a-1)\beta(a)\).  For every block length with
\(\mu_a\ge1\) and \(\delta_a>0\), the represented block response kernel
\(T_{W,a}\) gives
\begin{equation}
 \label{eq:wick-mixing-block-bound}
 |R(f)-\widehat R_N(f)|
 \le\frac{2R_A}{\mu_a}\sqrt{\operatorname{Tr}_{\rm blk(a)}T_{W,a}}
 +3\sqrt{\frac{\log(4/\delta_a)}{2\mu_a}}
 +\frac{2a}{N}.
\end{equation}
Here \(\operatorname{Tr}_{\rm blk(a)}\) is the finite compression to the
\(\mu_a\)-block evaluation carrier.  Block lengths with \(\mu_a=0\) or
\(\delta_a\le0\) have value \(+\infty\), so the recordwise minimum over all
represented block lengths is a valid bound coordinate.  For a positive
represented horizon \(T\), the Part~I predictable-path cell uses
\begin{equation}
 \label{eq:wick-sequential-radius}
 \mathfrak R_T^{\rm seq}
 \le R_A\mathfrak q_T(A_W)
 \le\frac{R_A}{T}\sqrt{\operatorname{Tr}_{\rm seq}T_{W,\lambda}}
\end{equation}
where \(\operatorname{Tr}_{\rm seq}\) is the Part~I supremum of the finite
\(T\)-step path/tree compression traces.  For a fixed response this gives the
displayed bound; a selectable response uses the learned-operator sequential
cover, and the Bayesian branch uses the martingale PAC--Bayes term
\begin{equation}
 \label{eq:wick-martingale-pac-bayes}
 \frac{\mathrm{KL}_{E,W}^{\mathrm{str}}(Q)+\log(1/\delta)}{\xi T}
 +\frac{\xi}{8}.
\end{equation}
for its represented \(\xi>0\); other coefficient values are nonordinary.
The online regret bound is
\(2L_\ell T\sup_{\mathbf z}\mathfrak R_T^{\mathrm{seq}}\).  Each complementary
dependence cell uses its own Part~I candidate or the value \(+\infty\); a
dark space alone supplies no dependence coordinate.
\item Define the acquisition coordinate
\(\Gamma_{\rm acq}:=\sup_{h,a}|\widehat g-g|\in[0,\infty]\).  The recorded
\(\varepsilon_{\rm opt}\)-maximizer loses at most
\begin{equation}
 \label{eq:wick-active-score-loss}
 \sup_ag(h,a)-g(h,\widehat a(h))
 \le2\Gamma_{\rm acq}+\varepsilon_{\rm opt}.
\end{equation}
This maximization is the operational action choice within a fixed acquisition
record.  Its error and action cost enter that record, after which the
saturated profile compares complete acquisition structures.
Over a horizon \(H\), the action-value loss is the survival-weighted sum of
these stepwise terms.  Resistance-greedy selection obtains the \(1-1/e\)
batch factor on the positive submodularity-certificate cell in the earlier
querying item; on its complement that factor-candidate is \(+\infty\).
\item For RL authority, compress \(T_{W,\lambda}\) first to the represented
action span and then to the safe subspace.  Effective-dimension monotonicity
is unconditional.  On the named Part~I safe source-order cell, the exact
source residual has the displayed upper bound:
\begin{align}
 \label{eq:wick-safe-effective-dimension-order}
 N_W^{\mathcal A,{\rm safe}}(\rho)
 &\le N_W^{\mathcal A}(\rho)\le N_W(\lambda,\rho),\\
 \label{eq:wick-safe-source-residual}
 B_W^{\mathcal A,{\rm safe}}(\rho;f_*)
 &\le C_r\|g_{\rm safe}\|^2\rho^{2r}
 +\|(I-P_{\rm safe})f_*\|^2.
\end{align}
The critic and actor objectives are
\eqref{eq:wick-rl-critic-loss} and
\eqref{eq:wick-rl-gibbs-policy}.  Define the extended coordinates
\(\varepsilon_M\) as the same-norm model defect and \(\varepsilon_B\) as the
learned Bellman residual.  On the Part~I discounted Bellman/contraction cell,
the Bellman candidate is
\begin{align}
 \|V^*-v\|_\infty
 &\le\frac{\varepsilon_B+\varepsilon_M}{1-\gamma},
 \label{eq:wick-bellman-value}\\
 \|V^*-V^{\widehat\pi_v}\|_\infty
 &\le\frac{2\gamma(\varepsilon_B+\varepsilon_M)}{(1-\gamma)^2}
 +\frac{2\varepsilon_M}{1-\gamma}.
 \label{eq:wick-bellman-policy}
\end{align}
Here \(0<\gamma<1\) is the declared discount or contraction modulus.  The
distance from \(J_D\) to the legal policy-current family is the irreducible
authority/action gap; it cannot be removed by additional value samples.
The complementary control-status cell assigns these Bellman candidates
\(+\infty\) and retains its typed control defect.
\item For a complete target-qualified Wick record, source-resolved
statistical, information, structural, and dynamical candidates combine only
after conversion to the same deployed target advantage \(g_W\):
\begin{equation}
 \label{eq:wick-source-resolved-target-minimum}
 g_W\le\min\left\{
 1,\ U_W^{\rm stat\to tgt},\
 \beta_W^++\sqrt{\frac{\log2}{2}J_W^{\rm src}},\
 eH_B\sum_a\Lambda_W^{(a)},\
 U_W^{\rm dyn\to tgt}
 \right\}.
\end{equation}
Here \(\Lambda_W^{(a)}\) is the contextual charge assigned to source
\(a\); it is distinct from the generator \(\mathcal L_W\).
For a finite sample-independent output codebook of logarithmic size
\(L_{\rm out}\), the fallback deviation candidate is
\begin{equation}
 \label{eq:wick-finite-description-radius}
 2\sqrt{\frac{2L_{\rm out}}{m_{\rm eff}}}
 +3\sqrt{\frac{\log(2/\delta_{\rm eff})}{2m_{\rm eff}}}
 +\varepsilon_{\rm dep}.
\end{equation}
\item Galerkin or finite-element records use
\eqref{eq:wick-feshbach-perturbation} and
\eqref{eq:wick-kernel-perturbation} after their form-to-resolvent error has
been certified.  Every finite Galerkin carrier uses its finite Part~I trace.
Along a represented nested dense family, the population effective dimension
is the monotone extended value in
\eqref{eq:wick-effective-dimension-monotone-limit}.  Operator approximation
errors on the \(d_n\)-dimensional carrier are charged by
\eqref{eq:wick-effective-dimension-interval}.  A represented tail estimate
may supply a finite upper envelope; without one, the upper value is
\(+\infty\).  No compactness or trace-class premise is inserted.
\item On every finite empirical band, define
\(\varepsilon_H:=\|\widehat H-H\|\),
\(\varepsilon_C:=\|\widehat C-C\|\), its derived separation
\(g_\Delta\ge0\), and
\[
 d_\Delta:=
 \begin{cases}
  2\varepsilon_H/g_\Delta,&g_\Delta>2\varepsilon_H,\\
  1,&g_\Delta\le2\varepsilon_H.
 \end{cases}
\]
With
\(c=\max\{\|C\|,\|\widehat C\|\}\),
\begin{equation}
 \label{eq:wick-empirical-defect-interval}
 \bigl|\|CP_\Delta\|-\|\widehat C\widehat P_\Delta\|\bigr|
 \le\varepsilon_C+c d_\Delta.
\end{equation}
Thus zero is excluded from this interval precisely when the data certify a
nonzero quasi-lumpability defect.  On every common certified reduced
subspace, define
\(\varepsilon_Q:=\|\widehat Q_K-Q_K\|\); Weyl's inequality bounds the
reduced-gap discrepancy by \(\varepsilon_Q\).  In quotient dimension
\(d_Q\), putting
\(\varepsilon_T=\|\widehat T-T\|\) gives
\begin{align}
 |\widehat N_W(\lambda,\rho)-N_W(\lambda,\rho)|
 &\le\frac{d_Q}{\rho}\varepsilon_T,
 \label{eq:wick-effective-dimension-interval}\\
 |\widehat m_W(\lambda,\rho;f)-m_W(\lambda,\rho;f)|
 &\le\frac{2}{\rho}\varepsilon_T
 \label{eq:wick-capture-interval}
\end{align}
for fixed normalized \(f\).  Data-dependent readouts add their separately
certified sample-splitting or uniform-selection radius.
\end{enumerate}
Each item retains the approximation floor, target-readout miss, source-cell
selection cost, model-channel perturbations, confidence allocation, and
deployment conversion of its Part~I antecedent.  An absent antecedent is the
nonordinary value \(+\infty\), never an omitted record or hidden assumption.
\end{corollary}

\begin{table}[tbp]
\centering
\small
\begin{tabular}{@{}p{.22\linewidth}p{.36\linewidth}p{.32\linewidth}@{}}
\toprule
Part~I quantitative family & Wick specialization & Record field / nonordinary output\\
\midrule
Identification and operator stability
& \eqref{eq:wick-identification-bound},
  \eqref{eq:wick-feshbach-perturbation},
  \eqref{eq:wick-kernel-perturbation}
& quotient gauge; derived design range/kernel; extended same-norm error\\
Fixed supervised learning
& \eqref{eq:wick-rademacher}, \eqref{eq:wick-krr-oracle},
  \eqref{eq:wick-krr-rate}
& positive response is derived; named source/eigenvalue cell or general oracle\\
Bridge or feature selection
& \eqref{eq:wick-free-oracle-bound},
  \eqref{eq:wick-feature-learning-bound},
  \eqref{eq:wick-learned-rademacher}
& inherited split/cover candidate; absent candidate equals \(+\infty\)\\
Bayesian and noisy learning
& \eqref{eq:wick-pac-bayes}, \eqref{eq:wick-debiased-krr},
  \eqref{eq:wick-source-cell-pac-bayes}
& inherited prior/channel cell; absent candidate equals \(+\infty\)\\
Dependent and online learning
& \eqref{eq:wick-mixing-block-bound},
  \eqref{eq:wick-sequential-radius},
  \eqref{eq:wick-martingale-pac-bayes}
& inherited dependence cell; other cells use their own candidate\\
Active learning
& resistance querying and \eqref{eq:wick-active-score-loss}
& extended resistance; submodularity coordinate controls only batch factor\\
Controlled/RL learning
& \eqref{eq:wick-rl-critic-oracle}--
  \eqref{eq:wick-rl-policy-oracle},
  \eqref{eq:wick-bellman-value}--\eqref{eq:wick-bellman-policy}
& safe/action and converter coordinates, each totalized in \([0,\infty]\)\\
Saturated-profile working frontier
& \eqref{eq:wick-minimum-lift-energy},
  \eqref{eq:wick-certified-simplicity-bound},
  \eqref{eq:wick-rl-minimal-working-bound}
& risk-restricted saturated profile and strict Pareto inverse; inner minima
  remain record coordinates; empty feasible set is typed\\
Deployment and continuum
& \eqref{eq:wick-source-resolved-target-minimum}, empirical intervals,
  resolvent/trace transfer
& extended converter, allocated simultaneous event, finite-carrier or
  extended-trace accounting\\
\bottomrule
\end{tabular}
\caption{Exhaustive location of the quantitative learning families used by
the Wick specialization.  Every row is evaluated on one complete selected
generator.  A dependent upper-error or risk field returns \(+\infty\), a
dependent lower-usefulness field returns zero, and every other unavailable
coordinate retains its typed nonordinary value; none is replaced by an
omitted theorem or numerical surrogate.}
\label{tab:wick-quantitative-coverage}
\end{table}

\subsection{Part~I model choice and the saturated certified frontier}
\label{subsec:wick-observable-minimality}

No new simplicity functional is introduced.  The model family, selection
penalties, and comparison order are inherited from Part~I.  Source choice is
charged by its exact-cell Rademacher union term or by the cell-mixture
PAC--Bayes divergence.  A label-dependent operator choice is charged by the
learned-operator cover, and coefficient selection is charged by the declared
finite coefficient family.  At fixed geometry, readout size is chosen by the
Part~I target-subspace coverage criterion.  Across complete records,
simplicity is the Pareto order of the existing saturated resource and action
vector \(\mathbf c\), with optimality interpreted throughout by
\cref{def:wick-saturated-optimality}.

\begin{proposition}[Inner minimum-response-energy coordinate]
\label{prop:wick-quotient-energy-minimality}
Let \(R_C:\mathcal Z_C\to\mathcal Y_C\) be a bounded operator between
Hilbert spaces and let \(T_C=R_CR_C^*\).  Define the totalized response
energy
\begin{equation}
 \label{eq:wick-totalized-response-energy}
 \mathcal E_{T_C}(f):=
 \begin{cases}
  \|T_C^{\dagger/2}f\|^2,&f\in\operatorname{ran}T_C^{1/2},\\
  +\infty,&f\notin\operatorname{ran}T_C^{1/2}.
 \end{cases}
\end{equation}
Then \(\operatorname{ran}T_C^{1/2}=\operatorname{ran}R_C\) and, for every
\(f\in\mathcal Y_C\),
\begin{equation}
 \label{eq:wick-minimum-lift-energy}
 \mathcal E_{T_C}(f)
 =\inf_{z\in\mathcal Z_C:\,R_Cz=f}\|z\|_{\mathcal Z_C}^2,
 \qquad \inf\varnothing:=+\infty.
\end{equation}
For \(f\in\operatorname{ran}R_C\), the infimum is attained uniquely in
\((\ker R_C)^\perp\), at
\begin{equation}
 \label{eq:wick-minimum-lift}
 z_f=R_C^\dagger f.
\end{equation}
In particular, take \(R_C=T_{W,\lambda}(C)^{1/2}\) on its represented
support.  Then the inverse-response penalty in
\eqref{eq:wick-free-inner-loss} is exactly the least squared response effort
among all latent lifts producing \(f\).  The same statement holds with
\(T_C^{\mathcal A,{\rm safe}}\) for the RL critic.  Nonclosed continuum
ranges are included: unattainable targets receive \(+\infty\), while every
attainable target has the minimum lift above.
This minimum is the response-energy coordinate of a fixed complete record.
It enters the resource/action ledger used by the saturated profile and does
not compare different response kernels or source cells.
\end{proposition}

\begin{proof}
The polar decomposition \(R_C=T_C^{1/2}V_C\) gives the operator-range
identity \(\operatorname{ran}R_C=\operatorname{ran}T_C^{1/2}\).  For an
attainable \(f\), the affine fiber \(R_C^{-1}\{f\}\) is a nonempty closed
set of the form \(z_f+\ker R_C\).  Its unique element perpendicular to
\(\ker R_C\) is the Moore--Penrose solution \(z_f=R_C^\dagger f\).
Functional calculus and the polar decomposition give
\(\|z_f\|=\|T_C^{\dagger/2}f\|\).  Pythagoras then yields
\(\|z_f+k\|^2=\|z_f\|^2+\|k\|^2\) for every \(k\in\ker R_C\).  For an
unattainable \(f\), both sides of \eqref{eq:wick-minimum-lift-energy} equal
\(+\infty\) by definition.
\end{proof}

\begin{definition}[Inherited Wick family-selection certificate]
\label{def:wick-inherited-family-selection-certificate}
Fix a resource--action cell \((B,a)\).  Let
\(\mathfrak A_W\) be the exact Part~I source cells represented in that cell.
For \(\alpha\in\mathfrak A_W\), let \(\Theta_\alpha\) be its
declared finite-coverable Wick generator/bridge family, let \(\Lambda\) be
the declared finite set of training-coefficient vectors, let \(\mathfrak K\)
be the finite capacity grid used in Part~I, and let
\(\mathcal H_{\alpha,\theta,\lambda,K}^W\) be the resulting readout or
policy-loss class at capacity \(K\).  Let
\(\mathfrak I_W\subseteq\mathfrak A_W\times\Lambda\times\mathfrak K\)
be the represented family indices.  Put
\begin{equation}
 \label{eq:wick-inherited-model-union}
 \mathcal H_W(B,a)
 :=\bigcup_{(\alpha,\lambda,K)\in\mathfrak I_W}
   \bigcup_{\theta\in\Theta_\alpha}
   \mathcal H_{\alpha,\theta,\lambda,K}^W.
\end{equation}
When \(\mathfrak I_W=\varnothing\), the union is empty, every upper
certificate below is defined to be \(+\infty\), and the learner returns the
typed source/authority/budget obstruction.  The remaining formulas are
evaluated on each nonempty finite Part~I stage; this is a case distinction,
not a restriction on the underlying family.
Every union index, net resolution, terminal kernel, coefficient vector,
training transcript, and deployment map is part of the complete record.
Equations \eqref{eq:wick-inherited-source-union-radius} and
\eqref{eq:wick-inherited-mixture-kl} are respectively the simultaneous
specializations of
\cref{eq:wick-source-cell-union,eq:wick-learned-rademacher} and
\cref{eq:wick-source-cell-pac-bayes,eq:wick-pac-bayes}.

For the Rademacher branch, define the within-family learned-response radius
\begin{equation}
 \label{eq:wick-inherited-learned-rademacher-radius}
 \mathfrak r_{\alpha,\lambda,K}^W
 :=\inf_{\varrho>0}\left[
 \frac Rm\sqrt{\sup_{\theta\in\Theta_\alpha}
                    \operatorname{Tr}_S T_\theta}
 +R\sqrt{\frac{L_{T,\alpha}\varrho}{m}}
 +\frac Rm\sqrt{2\log\mathcal N(\Theta_\alpha,\varrho)}
 \right]
\end{equation}
The supremum and covering number are extended values; no attainment or
compactness is used.  The complete exact-family union radius is
\begin{equation}
 \label{eq:wick-inherited-source-union-radius}
 \mathfrak r_{\rm fam}^W
 :=\max_{(\alpha,\lambda,K)\in\mathfrak I_W}
       \mathfrak r_{\alpha,\lambda,K}^W
 +\sqrt{\frac{2\log|\mathfrak I_W|}{m}}.
\end{equation}
Thus, for a \([0,1]\)-valued \(L_\ell\)-Lipschitz loss, the inherited
Rademacher upper certificate is
\begin{equation}
 \label{eq:wick-inherited-rademacher-ucb}
 U_{\rm Rad}^W(h)
 :=\widehat R_S(h)+2L_\ell\mathfrak r_{\rm fam}^W
 +3\sqrt{\frac{\log(2/\delta_{\rm Rad})}{2m}}
 +\varepsilon_{\rm task}+\varepsilon_{\rm op}
 +\varepsilon_{\rm dep}.
\end{equation}
Here \(L_{T,\alpha}\) is supplied by
\cref{prop:wick-learned-operator-stability}; a fixed operator uses the
fixed-response radius in \eqref{eq:wick-fixed-operator-risk}.

For the Bayesian branch, let \(\mathfrak I_{W,\rm PB}\) be the finite
hierarchical index of exact source cell, represented operator-net point, and
capacity.  Use exactly the Part~I mixture prior
\(P=\sum_{\iota\in\mathfrak I_{W,\rm PB}}w_\iota P_\iota\), fixed
independently of the labeled sample, and posterior
\(Q_\lambda=\sum_\iota q_{\lambda,\iota}Q_{\lambda,\iota}\).  Put
\begin{align}
 \label{eq:wick-inherited-mixture-kl}
 \operatorname{KL}_W(Q_\lambda\Vert P)
 &:={\rm KL}(q_\lambda\Vert w)
 +\sum_{\iota\in\mathfrak I_{W,\rm PB}}q_{\lambda,\iota}
   {\rm KL}(Q_{\lambda,\iota}\Vert P_\iota),\\
 \label{eq:wick-inherited-pac-bayes-ucb}
 U_{\rm PB}^W(Q_\lambda)
 &:=\widehat R_S(Q_\lambda)
 +\sqrt{\frac{\operatorname{KL}_W(Q_\lambda\Vert P)
 +\log(2\sqrt m|\Lambda|/\delta_{\rm PB})}{2(m-1)}}
 +\varepsilon_{\rm task}+\varepsilon_{\rm op}+\varepsilon_{\rm dep}.
\end{align}
The PAC--Bayes expression is an ordinary candidate for \(m\ge2\).  At
\(m<2\), or when its prior/posterior fields are unavailable, it has value
\(+\infty\) and does not affect the recordwise same-loss minimum.
A label-dependent \(\theta\)-selection is either an index in the fixed
hierarchical prior or is paid by the learned-operator cover; it is not treated
as a label-independent geometry.  A split-validation certificate such as
\eqref{eq:wick-free-oracle-bound} is an alternative candidate for the same
selection step; its radius is not added again to a union or mixture candidate.
Let \(U_W(M)\) be the minimum of
\eqref{eq:wick-inherited-rademacher-ucb},
\eqref{eq:wick-inherited-pac-bayes-ucb}, and the other applicable candidates
in \cref{cor:wick-quantitative-ledger}, after all have been converted to the
same population loss on one confidence event.  An unavailable candidate has
value \(+\infty\).  Let \(\Gamma_W(M)\in[0,\infty]\) denote the smallest applicable
two-sided Part~I radius on that same event, including the identical family
index, cover, channel, operator, and deployment charges.  Thus
\(U_W(M)\le\mathcal R(M)+2\Gamma_W(M)\) in the extended reals; absence of a
two-sided comparator certificate is exactly \(\Gamma_W(M)=+\infty\).
Both minima are recordwise statistical certificate coordinates.  They
implement the same-loss Rademacher/PAC--Bayes choice of Part~I; the saturated
profile below, rather than either scalar radius, orders complete records.
\end{definition}

\begin{definition}[Saturated-profile semantics of Wick optimality]
\label{def:wick-saturated-optimality}
The terms \emph{optimal}, \emph{minimal}, and \emph{simplest} retain the
Part~I hierarchy.  An optimization performed with a complete record fixed---
including least-action projection, minimum response lift, ridge fitting,
critic fitting, or minimum full-coverage readout---is an \emph{inner
optimization}.  Its value, optimizer transcript, and optimization error are
coordinates of that record.  Only the saturated profile orders two different
structural records.

Structural optimization is defined only by the inherited saturated profile,
instantiated for Wick records in \eqref{eq:wick-saturated-profile}, and its
Pareto-ordered strict inverse \eqref{eq:wick-strict-inverse}.  Let
\(\ell\) denote a fixed Part~I population-law, estimand, norm, and loss-scale
identifier after every declared converter has been applied.  Let
\(\mathscr R_{W,\ell}^{(n)}(B,a)\) denote that typed partition of the charged
finite stage produced by the Part~I learner, and suppress \(\ell\) from the
notation below by writing \(\mathscr R_W^{(n)}(B,a)\).  Records with different
identifiers define separate profiles and are never compared through a
statistical minimum.  Let \(\underline V_W(M)\) be the simultaneous Part~I
lower confidence bound for \(V_M^W\).  At requested population-risk accuracy
\(\epsilon\), the computable saturated profile and strict inverse within the
fixed typed partition are
\begin{align}
 \label{eq:wick-certified-risk-profile}
 \widehat G_{Q,W,\mathfrak A,n}^{\mathrm{sat},\le\epsilon}(B,a)
 &:=\sup_{\substack{M\in\mathscr R_W^{(n)}(B,a)\\
                    U_W(M)\le\epsilon}}
       \underline V_W(M),\\
 \label{eq:wick-certified-risk-inverse}
 \widehat{\mathfrak P}_{Q,W,\mathfrak A,n}^{>,\le\epsilon}(\theta)
 &:=\{(B,a):
 \widehat G_{Q,W,\mathfrak A,n}^{\mathrm{sat},\le\epsilon}(B,a)>\theta\}.
\end{align}
Thus structural optimality at fixed resources is the saturated supremum of
complete visibility.  A record is called an optimizer only on a cell where it
attains that supremum; otherwise the Part~I search transcript retains its
\(\varepsilon\)-maximizers and nonattainment status.  A simplest working
record is a Pareto-minimal point of the appropriate strict inverse.  The
population risk-qualified profile is
\eqref{eq:wick-risk-qualified-saturated-profile};
\cref{prop:wick-risk-qualified-saturation-inheritance} proves that it is a
nested restriction of the unrestricted Part~I profile.  Every earlier or later
inner minimizer is used only to evaluate the visibility, risk, action, or
resource coordinates that enter these profiles.
\end{definition}

\begin{proposition}[Part~I fixed-response capacity coordinate]
\label{prop:wick-minimum-capacity-readout}
Fix one Wick response and let \(P_W\) be the orthogonal projector onto its
represented response support.  For the declared target-score subspace
\(\mathcal Y_{\rm tgt}\), put
\begin{equation}
 \label{eq:wick-intrinsic-target-subspace}
 \mathcal V_W^*:=\overline{\operatorname{ran}
 (P_WP_{\mathcal Y_{\rm tgt}})},
 \qquad d_W^*:=\dim\mathcal V_W^*\in\mathbb N\cup\{\infty\}.
\end{equation}
For a reachable readout subspace \(\mathcal R_K\) with
\(\dim\mathcal R_K\le K\), define
\[
 \mathcal C_W(y):=\frac{\|P_{\mathcal V_W^*}y\|^2}{\|y\|^2},
 \qquad
 \mathcal E_W(\mathcal R_K;y)
 :=\frac{\|P_{\mathcal R_K}P_{\mathcal V_W^*}y\|^2}{\|y\|^2}.
\]
As in Part~I, both ratios are defined to be zero at \(y=0\).
Then
\begin{equation}
 \label{eq:wick-coverage-gap}
 \mathcal C_W(y)-\mathcal E_W(\mathcal R_K;y)
 =\frac{\|(I-P_{\mathcal R_K})P_{\mathcal V_W^*}y\|^2}{\|y\|^2}.
\end{equation}
Full structural coverage is equivalent to
\begin{equation}
 \label{eq:wick-projection-coverage}
 P_{\mathcal R_K}P_{\mathcal V_W^*}=P_{\mathcal V_W^*}.
\end{equation}
It forces \(K\ge d_W^*\).  On the finite-dimensional cell, any represented
family containing \(\mathcal V_W^*\) attains it at \(K=d_W^*\); on the
cell \(d_W^*=\infty\), no finite-capacity readout has full coverage.  Hence
\(d_W^*\) is the
minimum readout capacity that covers the complete target-active response
subspace, exactly as in the Part~I model-family criterion.  It is an inner
minimum at fixed response.  Across responses and source cells, capacity is a
component of \(\mathbf c\), and selection is governed by
\cref{def:wick-saturated-optimality}.
\end{proposition}

\begin{proof}
The orthogonal decomposition
\(v=P_{\mathcal R_K}v+(I-P_{\mathcal R_K})v\), with
\(v=P_{\mathcal V_W^*}y\), proves
\eqref{eq:wick-coverage-gap}.  Identity
\eqref{eq:wick-projection-coverage} holds precisely when
\(\mathcal V_W^*\subseteq\mathcal R_K\), which implies
\(d_W^*\le\dim\mathcal R_K\le K\).  On the finite-dimensional cell, taking
\(\mathcal R_{d_W^*}=\mathcal V_W^*\) proves attainability on exactly the
represented reachable-family cell containing that subspace.  Infinite
\(d_W^*\) rules out finite \(K\) directly.
\end{proof}

\begin{theorem}[Part~I oracle selection and the simplest certified Wick record]
\label{thm:wick-simplest-certified-model}
Use the Part~I confidence allocation
\(\delta_{\rm Rad}+\delta_{\rm PB}+\delta_{\rm other}\le\delta\) in
\cref{def:wick-inherited-family-selection-certificate}.  At represented
stage \(n\), let \(\mathscr R_W^{(n)}(B,a)\) be the charged finite family of
complete Wick records produced by Part~I, and let \(\mathcal R(M)\) denote
their common population loss.  For a requested risk accuracy \(\epsilon\)
and the inherited visibility threshold \(\theta\), define
\begin{equation}
 \label{eq:wick-certified-complexity-selector}
 \mathscr F_W(\epsilon,\theta;B,a)
 :=\{M\in\mathscr R_W^{(n)}(B,a):
       U_W(M)\le\epsilon,\ \underline V_W(M)>\theta\}.
\end{equation}
The selected Pareto set is
\begin{equation}
 \widehat{\mathscr M}_W(\epsilon,\theta;B,a)
 :=\operatorname{Min}_{\preceq\mathbf c}
   \mathscr F_W(\epsilon,\theta;B,a).
\end{equation}
Here \(\operatorname{Min}_{\preceq\mathbf c}\) denotes all Pareto-minimal
records in the complete resource--action order already used by the saturated
profile.  These recordwise frontiers are attached to feasible grid points of
the certified risk-qualified strict inverse
\eqref{eq:wick-certified-risk-inverse}; its full represented grid boundary is identified
in \textup{(v)}.  Then, with probability at least
\(1-\delta\):
\begin{enumerate}[label=\textup{(\roman*)}]
\item every
\(\widehat M\in\widehat{\mathscr M}_W(\epsilon,\theta;B,a)\) satisfies
\begin{equation}
 \label{eq:wick-certified-working-bound}
 \mathcal R(\widehat M)\le\epsilon,
 \qquad V^W_{\widehat M}>\theta;
\end{equation}
\item for every returned \(\widehat M\), no certified
\((\epsilon,\theta)\)-working record in
\(\mathscr R_W^{(n)}(B,a)\) has
\(\mathbf c(M)\prec\mathbf c(\widehat M)\);
\item every comparator \(M_0\in\mathscr R_W^{(n)}(B,a)\) whose extended
two-sided radius and lower visibility satisfy
\begin{equation}
 \label{eq:wick-certified-comparator-margin}
 \mathcal R(M_0)+2\Gamma_W(M_0)\le\epsilon,
 \qquad \underline V_W(M_0)>\theta,
\end{equation}
belongs to \(\mathscr F_W(\epsilon,\theta;B,a)\), and some
\(\widehat M\in\widehat{\mathscr M}_W(\epsilon,\theta;B,a)\) obeys
\begin{equation}
 \label{eq:wick-certified-simplicity-bound}
 \mathbf c(\widehat M)\preceq\mathbf c(M_0);
\end{equation}
\item when \(\mathscr R_W^{(n)}(B,a)\ne\varnothing\), an
\(\varepsilon_{\rm opt}\)-minimizer of \(U_W\) satisfies the Part~I oracle
inequality
\begin{equation}
 \label{eq:wick-penalized-minimal-oracle}
 \mathcal R(\widehat M_{\rm or})
 \le\inf_{M\in\mathscr R_W^{(n)}(B,a)}
 \{\mathcal R(M)+2\Gamma_W(M)\}+\varepsilon_{\rm opt}.
\end{equation}
This is the Part~I risk-oracle comparison at a fixed represented stage.  It
supplies the risk restriction in
\eqref{eq:wick-certified-risk-profile}; structural selection is its saturated
visibility supremum and Pareto strict inverse.
\item across the Part~I resource--action grid, the learner reports the
returned record frontiers at precisely the Pareto-minimal grid points of the
certified strict region
\[
 \widehat{\mathfrak P}_{Q,W,\mathfrak A,n}^{>,\le\epsilon}(\theta).
\]
\end{enumerate}
If the feasible set is empty, the output is the existing insufficient-
confidence, target-miss, authority, or budget status.  If several records
remain statistically unresolved and Pareto-minimal, the learner returns their
certified frontier; no uniqueness criterion is added.  The directed sequence
of stages carries the Part~I net remainder.  Its infimum and upward
feasibility region are therefore defined without compactness; Pareto-minimal
continuum records are reported exactly on stages or limit cells where they
exist, and nonattainment remains in the search transcript.
\end{theorem}

\begin{proof}
The learned-response Rademacher bound and exact-source union bound give
\eqref{eq:wick-inherited-rademacher-ucb} uniformly over
\eqref{eq:wick-inherited-model-union}.  The PAC--Bayes chain rule gives
\eqref{eq:wick-inherited-mixture-kl}; applying the structural PAC--Bayes
bound to every \(\lambda\in\Lambda\) at confidence
\(\delta_{\rm PB}/|\Lambda|\) gives
\eqref{eq:wick-inherited-pac-bayes-ucb}.  The confidence allocation and union
bound put all applicable candidates on one event.  Their recordwise minimum
is valid because they bound the same population loss after the declared
converters.

On that event,
\(\mathcal R(M)\le U_W(M)\) and
\(\underline V_W(M)\le V_M^W\), proving \textup{(i)}.  Statement
\textup{(ii)} is the definition of a Pareto-minimal element of the finite
feasible set.  A two-sided radius gives
\(U_W(M_0)\le\mathcal R(M_0)+2\Gamma_W(M_0)\), so
\eqref{eq:wick-certified-comparator-margin} places \(M_0\) in the feasible
set.  Every member of a finite partially ordered set dominates a minimal
member, proving \textup{(iii)}.  Finally,
\[
 \mathcal R(\widehat M_{\rm or})
 \le U_W(\widehat M_{\rm or})
 \le\inf_MU_W(M)+\varepsilon_{\rm opt}
 \le\inf_M\{\mathcal R(M)+2\Gamma_W(M)\}+\varepsilon_{\rm opt},
\]
which proves \textup{(iv)}.  By
\eqref{eq:wick-certified-risk-profile}--
\eqref{eq:wick-certified-risk-inverse}, a grid point belongs to the certified
strict inverse exactly when it contains a record in
\(\mathscr F_W(\epsilon,\theta;B,a)\).  The grid-level reporting rule selects
exactly the Pareto-minimal such points and returns the recordwise Pareto
frontier there, proving \textup{(v)}.  Net stability errors are already components of
\(\Gamma_W\) and \(\mathbf c\).
\end{proof}

\begin{corollary}[Saturated-profile minimum working models for free Caus and RL]
\label{cor:wick-free-rl-minimal-models}
For unrestricted Caus, take the candidates produced by
\eqref{eq:wick-free-inner-loss} over the exact source/operator/coefficient
family in \eqref{eq:wick-inherited-model-union}.  Select them by the
Rademacher certificate \eqref{eq:wick-inherited-rademacher-ucb}, the
PAC--Bayes certificate \eqref{eq:wick-inherited-pac-bayes-ucb}, or their
valid same-loss minimum to form \(U_W\), and then apply the
risk-qualified saturated profile at the inherited threshold \(\theta\).  Then
\eqref{eq:wick-certified-working-bound}--
\eqref{eq:wick-certified-simplicity-bound} hold, and the response-energy
coordinate is the exact minimum lift from
\cref{prop:wick-quotient-energy-minimality}.  On the zero-waste fiber
\(\mathfrak A_D^0\), its action coordinate is zero.  At fixed response, the
smallest full-coverage readout has \(K=d_W^*\) by
\cref{prop:wick-minimum-capacity-readout}; across responses, the inherited
cover or mixture charge enters \(\mathbf c\), and the risk-qualified saturated
profile decides the tradeoff.

For RL, let \(g_{\rm auth},g_{\rm obs}\in[0,\infty]\) be the derived
budget/safe-set and observation gaps and define the extended policy certificate
\begin{equation}
 \label{eq:wick-rl-working-certificate}
 U_{\rm RL}(M)
 :=\frac{2\Gamma_{Q,K}
 +\tau\log(1/\nu_{\min})
 +\operatorname{Osc}(\widehat q_K)\sqrt{\kappa_\pi/2}}
 {1-\gamma}
 +g_{\rm auth}+g_{\rm obs}.
\end{equation}
Fix the Part~I environment, value norm, reward convention, and deployment-law
identifier.  Over the resulting typed partition of represented RL candidates,
with the nonordinary converter/support values totalized as in
\cref{thm:wick-rl-loss-calibration}, let
\[
 \widehat{\mathscr M}_{\epsilon,\theta}^{\rm RL}
 :=\operatorname{Min}_{\preceq\mathbf c}
 \{R:U_{\rm RL}(R)\le\epsilon,\ \underline V_W(R)>\theta\},
\]
where every critic, bridge, coefficient, and policy-family selection term in
\(\Gamma_{Q,K}\) is the corresponding inherited Rademacher or PAC--Bayes
charge.  Every
\(\widehat R\in\widehat{\mathscr M}_{\epsilon,\theta}^{\rm RL}\)
satisfies
\begin{equation}
 \label{eq:wick-rl-minimal-working-bound}
 \|V_\alpha^*-V_\alpha^{\widehat\pi}\|_\infty\le\epsilon,
 \qquad V^W_{\widehat R}>\theta,
 \qquad
 \nexists R:\ U_{\rm RL}(R)\le\epsilon,\quad
 \underline V_W(R)>\theta,
 \quad\mathbf c(R)\prec\mathbf c(\widehat R).
\end{equation}
Thus the output is precisely the Part~I Pareto-minimal certified safe-policy
frontier in the risk-qualified strict inverse at the requested value accuracy
and visibility threshold.
\end{corollary}

\begin{proof}
The unrestricted statement is
\cref{thm:wick-simplest-certified-model} with the exact energy identity from
\cref{prop:wick-quotient-energy-minimality}; zero action on
\(\mathfrak A_D^0\) follows from \cref{thm:wick-least-action-law}.  For RL,
\cref{thm:wick-rl-loss-calibration} and the two declared approximation gaps
give
\(\|V_\alpha^*-V_\alpha^{\pi_M}\|_\infty\le U_{\rm RL}(M)\).
Feasibility therefore proves the first inequality in
\eqref{eq:wick-rl-minimal-working-bound}; the simultaneous lower visibility
bound proves its second conclusion, and Pareto minimality in the saturated
strict inverse proves the last statement.
\end{proof}

\subsection{Certified learner and primal--dual profile}
\label{subsec:wick-certified-learner}

\begin{definition}[Complete Wick certificate]
\label{def:complete-wick-certificate}
A complete represented Wick certificate is an extension
\[
 \mathscr C_W=
 (\mathscr C_I,\mathfrak W_{\Abs},
 E_{W,\Delta},\Sigma_C,F_W,T_{W,\lambda},
 \mathscr S,\mathscr T,\mathscr L,\mathbf c),
\]
where \(\mathscr C_I\) is the complete selected Part~I compensated
quotient--geometry record.  In particular, it already contains
\(U,P_\Delta,H_{\Abs},G,A_{\mathcal Y},C,\mathfrak A,D,J_C\), their source
and authority provenance, and their structural resource coordinates.  The
map \(\mathfrak W_{\Abs}\) is \eqref{eq:wick-record-map}; it adds no source
or resource field.  The derived ledger \(\mathscr S\) contains the
self-energy, dark-space, bracket,
quasi-lumpability, and continuum-transfer certificates;
\(\mathscr T\) contains the dynamics, free-Caus, or RL loss, its data split,
optimization error, selected bridge/kernel, and behavior-to-deployment
comparison, together with the exact source cell, declared operator and
coefficient families, Rademacher cover or PAC--Bayes mixture charge,
target-subspace coverage identity, and certified risk tolerance;
\(\mathscr L\) contains the target capture, statistical confidence, source,
noise, and deployment records; and \(\mathbf c\) is the inherited complete
Pareto resource vector together with the evaluated Caus action.  Every
derived object is computed from the Wick image of the same selected Part~I
record.  The optimization entries in \(\mathscr T\) are inner transcripts;
only the saturated profile below compares complete records.  Its Wick
visibility is the
existing master-certificate visibility evaluated with the full response and
authority data; write it \(V^W_{\mathscr C_W}\).
Its dynamic stability coordinate is the existing general-resolvent factor
\[
 S_W(\lambda)=\max\{0,1-\Delta_W(\lambda)\},
\]
where \(\Delta_W\) is the represented target-functional discrepancy between
the true compressed response and \(F_W(\lambda)^{-1}\).  It vanishes for the
exact block model of \cref{thm:wick-feshbach-response}; empirical and
continuum model residuals enter through the resolvent identity.

For an authority envelope \(\mathfrak A\), resource ceiling \(B\), and
action ceiling \(a\), let \(\mathcal C_{Q,W}^{\mathfrak A}(B,a)\) be the
represented family of complete certificates satisfying both ceilings and
having a nonzero quotient defect.  The zero-defect boundary belongs to the
existing exact-quotient cell.  Put
\begin{equation}
 \label{eq:wick-saturated-profile}
 G_{Q,W,\mathfrak A}^{\mathrm{sat}}(B,a)
 :=\sup_{\mathscr C_W\in\mathcal C_{Q,W}^{\mathfrak A}(B,a)}
 V^W_{\mathscr C_W},
 \qquad \sup\varnothing:=0.
\end{equation}
For a threshold \(\theta\), its strict inverse is the upward-closed
resource--action feasibility region
\begin{equation}
 \label{eq:wick-strict-inverse}
 \mathfrak P_{Q,W,\mathfrak A}^{>}(\theta)
 :=\{(B,a):G_{Q,W,\mathfrak A}^{\mathrm{sat}}(B,a)>\theta\}
 =\bigcup_{\substack{\mathscr C_W:\,
                V^W_{\mathscr C_W}>\theta}}
   \{(B,a):\mathbf c_{\rm res}(\mathscr C_W)\preceq B,\quad
              \mathbf c_{\rm act}(\mathscr C_W)\le a\}.
\end{equation}
For a requested population-risk accuracy \(\epsilon\in[0,\infty]\), fix one
Part~I population-law, estimand, norm, and loss-scale identifier.  After all
declared converters have been charged, let
\(\mathcal C_{Q,W,\ell}^{\mathfrak A}(B,a)\) denote the corresponding typed
partition of \(\mathcal C_{Q,W}^{\mathfrak A}(B,a)\); suppress \(\ell\) in
the display below.  The same Part~I saturation restricted to working records
is
\begin{equation}
 \label{eq:wick-risk-qualified-saturated-profile}
 G_{Q,W,\mathfrak A}^{\mathrm{sat},\le\epsilon}(B,a)
 :=\sup_{\substack{\mathscr C_W\in
                    \mathcal C_{Q,W}^{\mathfrak A}(B,a)\\
                    \mathcal R(\mathscr C_W)\le\epsilon}}
       V^W_{\mathscr C_W},
 \qquad \sup\varnothing:=0,
\end{equation}
and its strict inverse is
\begin{equation}
 \label{eq:wick-risk-qualified-strict-inverse}
 \mathfrak P_{Q,W,\mathfrak A}^{>,\le\epsilon}(\theta)
 :=\{(B,a):
 G_{Q,W,\mathfrak A}^{\mathrm{sat},\le\epsilon}(B,a)>\theta\}.
\end{equation}
Its Pareto-minimal elements are reported when they exist.  In particular,
they exist on the finite resource--action grid used by the learner below;
the region itself, rather than a possibly empty set of minimal elements, is
the exact generalized inverse on an arbitrary ordered resource space.
\end{definition}

\begin{proposition}[Risk qualification is a restriction of Part~I saturation]
\label{prop:wick-risk-qualified-saturation-inheritance}
For \(0\le\epsilon_1\le\epsilon_2\le+\infty\),
\begin{align}
 G_{Q,W,\mathfrak A}^{\mathrm{sat},\le\epsilon_1}(B,a)
 &\le
 G_{Q,W,\mathfrak A}^{\mathrm{sat},\le\epsilon_2}(B,a)
 \le G_{Q,W,\mathfrak A}^{\mathrm{sat}}(B,a),
 \label{eq:wick-risk-profile-nesting}\\
 \mathfrak P_{Q,W,\mathfrak A}^{>,\le\epsilon_1}(\theta)
 &\subseteq
 \mathfrak P_{Q,W,\mathfrak A}^{>,\le\epsilon_2}(\theta)
 \subseteq
 \mathfrak P_{Q,W,\mathfrak A}^{>}(\theta).
 \label{eq:wick-risk-inverse-nesting}
\end{align}
At \(\epsilon=+\infty\), both rightmost inclusions are equalities.  The same
identities hold at every represented stage after replacing \(G\) and
\(\mathfrak P\) by the certified quantities in
\eqref{eq:wick-certified-risk-profile}--
\eqref{eq:wick-certified-risk-inverse} and taking the unrestricted certified
profile to be the same supremum without the condition \(U_W(M)\le\epsilon\).
Consequently risk accuracy restricts the record family inside the inherited
saturated optimization; it introduces no second structural order.
\end{proposition}

\begin{proof}
The record sets in each supremum are nested because
\(\{\mathcal R\le\epsilon_1\}\subseteq
  \{\mathcal R\le\epsilon_2\}\subseteq\mathcal C_{Q,W}^{\mathfrak A}(B,a)\).
Suprema preserve inclusion, proving \eqref{eq:wick-risk-profile-nesting}.
Taking strict superlevel sets proves \eqref{eq:wick-risk-inverse-nesting}.
Every extended risk satisfies \(\mathcal R\le+\infty\), which gives equality
at \(\epsilon=+\infty\).  Replacing \(\mathcal R\) by the simultaneous upper
certificate \(U_W\) gives the identical finite-stage proof.
\end{proof}

\begin{theorem}[Bayesian comparison for the Wick saturated learner]
\label{thm:wick-bayesian-comparison}
Fix a represented stage \(n\), a common confidence event, and one population
loss.  For a complete Wick record \(M\), let \(U_{\rm PB}^W(M)\) denote the
PAC--Bayes certificate in \eqref{eq:wick-inherited-pac-bayes-ucb}, evaluated
on its carried posterior and assigned \(+\infty\) when that candidate is
unavailable.  Define the PAC--Bayes-only saturated lower profile and strict
inverse by
\begin{align}
 \widehat G_{Q,W,\mathfrak A,n}^{\rm PB,\le\epsilon}(B,a)
 &:=%
 \sup_{\substack{M\in\mathscr R_W^{(n)}(B,a)\\
                  U_{\rm PB}^W(M)\le\epsilon}}
       \underline V_W(M),
 \qquad \sup\varnothing:=0,
 \label{eq:wick-pb-only-saturated-profile}\\
 \widehat{\mathfrak P}_{Q,W,\mathfrak A,n}^{\rm PB,>,\le\epsilon}(\theta)
 &:=%
 \{(B,a):
 \widehat G_{Q,W,\mathfrak A,n}^{\rm PB,\le\epsilon}(B,a)>\theta\}.
 \label{eq:wick-pb-only-strict-inverse}
\end{align}
Then
\begin{align}
 U_W(M)&\le U_{\rm PB}^W(M),
 \label{eq:wick-pb-pointwise-domination}\\
 \widehat G_{Q,W,\mathfrak A,n}^{\rm PB,\le\epsilon}(B,a)
 &\le
 \widehat G_{Q,W,\mathfrak A,n}^{\mathrm{sat},\le\epsilon}(B,a),
 \label{eq:wick-pb-profile-domination}\\
 \widehat{\mathfrak P}_{Q,W,\mathfrak A,n}^{\rm PB,>,\le\epsilon}(\theta)
 &\subseteq
 \widehat{\mathfrak P}_{Q,W,\mathfrak A,n}^{>,\le\epsilon}(\theta).
 \label{eq:wick-pb-inverse-domination}
\end{align}
The profile inequality is strict at \((B,a)\) whenever some
\(M\in\mathscr R_W^{(n)}(B,a)\) satisfies
\begin{equation}
 U_W(M)\le\epsilon<U_{\rm PB}^W(M),
 \qquad
 \underline V_W(M)>
 \widehat G_{Q,W,\mathfrak A,n}^{\rm PB,\le\epsilon}(B,a).
 \label{eq:wick-strict-pb-improvement-cell}
\end{equation}

Let \(\mathfrak B(\mathscr C_W)\) be the universal Bayesianization of
Part~I applied to the active execution carried by a complete Wick record.
Embed the Wick resource vector in the complete Bayesian resource lattice by
filling the belief, likelihood, normalization, and belief-memory slots with
zero.  The Part~I compiler and forgetting map give
\begin{equation}
 \mathcal R(\mathfrak B(\mathscr C_W))=\mathcal R(\mathscr C_W),
 \qquad
 V(\mathfrak B(\mathscr C_W))=V^W_{\mathscr C_W},
 \qquad
 \mathbf c(\mathscr C_W)\preceq
 \mathbf c(\mathfrak B(\mathscr C_W)),
 \label{eq:wick-direct-versus-bayesianized-record}
\end{equation}
and the Caus action coordinate is unchanged.  The resource inequality is
strict precisely when the canonical Bayesian annotation has a positive
belief-specific resource slot.  Hence the direct Wick realization reaches
the same risk--visibility point no later in the saturated Pareto inverse than
its universal Bayesian annotation.

For the full record classes, the inherited Part~I comparison remains
\begin{equation}
 \operatorname{BayesArr}_{I,H}^{\rm sat}(B)
 \le \operatorname{ActArr}_{I,H}^{\rm sat}(B)
 \le
 \operatorname{BayesArr}_{I,H}^{\rm sat}
 \bigl(\Psi_{\rm uBayes}(B,H,n)\bigr).
 \label{eq:wick-inherited-bayesianization-comparison}
\end{equation}
Thus Bayesian optimality is universal after the Part~I transfer ceiling,
whereas \eqref{eq:wick-pb-profile-domination} is a same-family finite-sample
domination and \eqref{eq:wick-direct-versus-bayesianized-record} is a
same-execution resource comparison.  The Wick class satisfies
\begin{equation}
 \label{eq:wick-profile-inclusion}
 \mathcal C_{Q,W}^{\mathfrak A}(B,a)
 \subseteq\mathcal C_{Q,\mathrm{nr}}^{\mathfrak A}(B),
 \qquad
 G_{Q,W,\mathfrak A}^{\mathrm{sat}}(B,a)
 \le G_{Q,\mathrm{nr},\mathfrak A}^{\mathrm{sat}}(B).
\end{equation}
These relations give the complete comparison: pointwise certificate
domination on the common Wick family, a resource comparison for the same
execution, and transfer-ceiling equivalence for the universal classes.
\end{theorem}

\begin{proof}
By definition, \(U_W(M)\) is the recordwise minimum of every applicable
same-loss certificate, including \(U_{\rm PB}^W(M)\).  This proves
\eqref{eq:wick-pb-pointwise-domination}.  Consequently
\[
 \{M:U_{\rm PB}^W(M)\le\epsilon\}
 \subseteq \{M:U_W(M)\le\epsilon\}
\]
inside the same charged family \(\mathscr R_W^{(n)}(B,a)\).  Taking suprema
of the same lower visibility proves
\eqref{eq:wick-pb-profile-domination}; taking strict superlevel sets proves
\eqref{eq:wick-pb-inverse-domination}.  The record in
\eqref{eq:wick-strict-pb-improvement-cell} lies only in the second feasible
set and has lower visibility above the first supremum, proving strictness.

The universal Bayesianization theorem of Part~I retains every query,
transition, action, continuation, stopping, and deployment kernel of the
original active record and appends a constant represented belief, unit
likelihood and normalizer, and identity update.  It therefore preserves the
execution law, deployed population risk, and target visibility.  Its
belief-specific costs are nonnegative.  The Part~I forgetting compiler
deletes exactly those slots with no increase in cost.  Together with the exact
Wick inheritance in \cref{prop:wick-exact-inheritance}, this proves
\eqref{eq:wick-direct-versus-bayesianized-record} and its strictness
criterion.  The two profile inequalities in
\eqref{eq:wick-inherited-bayesianization-comparison} are respectively the
Part~I forgetting and universal Bayesianization inequalities.  Finally, the
Wick record inclusion in the nonreversible family gives the stated
specialized-family relation.
\end{proof}

\begin{definition}[Two-stage Wick structural learner]
\label{def:wick-structural-learner}
The learner receives the candidate records produced by the Part~I structural
learner from a represented empirical generator or closed-form approximation,
with their confidence radii, finite-coverable quotient, geometry, and
authority families, resolvent and ridge grids, and total resource ceiling.
It later receives a represented target/reward sample.  The Wick stage neither
retypes nor enlarges those candidate families.
Raw trajectories enter through a represented generator/form estimator and
its dependence certificate.  An exact PDE operator is the zero-estimation-
error specialization of the same interface.

\begin{enumerate}[label=\textup{(L\arabic*)}]
\item On dynamics-only data, enumerate the inherited charged finite net.  For
every Part~I candidate, minimize \eqref{eq:wick-dynamics-training-loss}, or
the raw-path likelihood \eqref{eq:wick-transition-training-loss}, within
that record's existing structural and authority constraints.  Fit its
represented parameters and residual, then apply
\(\mathfrak W_{\Abs}\) to the selected quotient block.  Certify
\(E_{W,\Delta}\), the derived complement-resolvent bounds, self-energy,
dark dimension, bracket gap, the derived design range/kernel, and the
identification/perturbation
bounds of
\cref{prop:wick-learned-operator-stability,prop:wick-dynamics-identification}.
\item Freeze the quotient/geometry selection or retain a simultaneous
uniform cover.  On split target/control data, choose the target band and
readout.  Under unrestricted authority, minimize
\eqref{eq:wick-free-inner-loss} and
\eqref{eq:wick-free-outer-loss}, restricting to the zero-waste fiber when
least action is a hard requirement.  Under RL authority, alternate the
frozen-target critic and Gibbs actor steps in
\eqref{eq:wick-rl-composite-loss}.  In a hybrid envelope, optimize only the
declared control part.  Every selected bridge current is certified by
\cref{thm:wick-least-action-law}; an imposed bridge remains fixed.  For each
fixed response, compute \(\mathcal V_W^*\), set \(K=d_W^*\) whenever the
reachable family verifies \eqref{eq:wick-projection-coverage}, and otherwise
retain the exact gap \eqref{eq:wick-coverage-gap}.
\item Compute the lower-confidence master visibility and every applicable
entry of \cref{cor:wick-quantitative-ledger}.  At fixed \((B,a)\), return an
\(\varepsilon\)-maximizer of the saturated profile.  Ties within the
certified tolerance use the inherited Part~I Pareto order on the complete
resource--action vector; the declared candidate order only enumerates
unresolved equals.  When population-risk accuracy \(\epsilon\)
and visibility threshold \(\theta\) are requested, construct
\(U_W\) from the inherited Rademacher union/cover or PAC--Bayes
mixture/coefficient certificate in
\cref{def:wick-inherited-family-selection-certificate}, and apply
\eqref{eq:wick-certified-complexity-selector} to the certified
risk-qualified saturated profile; on the RL branch use
\eqref{eq:wick-rl-working-certificate}.  An empty certified feasible set is
reported by its typed failure status.
\item For a requested threshold \(\theta\), enumerate the inherited
budget/action grid, return the certified strict region
\eqref{eq:wick-strict-inverse}, or
\eqref{eq:wick-certified-risk-inverse} when risk-qualified, and report its
Pareto-minimal grid points.
\end{enumerate}
The output is a certificate and one of the following typed statuses:
certified profile, target spectral miss, uncompensated model defect,
target-visible dark mode, authority obstruction, insufficient statistical
confidence, continuum-transfer defect, or budget obstruction.
\end{definition}

\begin{theorem}[Correctness and termination of the Wick learner]
\label{thm:wick-learner-correctness}
At every represented Part~I stage, the candidate set is the charged finite
net constructed by that learner, and its perturbation remainder is the
extended value supplied by \cref{prop:wick-learned-operator-stability}.
Consequently the Wick stage terminates by finite enumeration.  The Part~I
confidence allocation and union/PAC--Bayes theorem construct a simultaneous
event of probability at least \(1-\delta\); an unavailable radius is
\(+\infty\).  On that event the budget-first output has certified
visibility at least
\begin{equation}
 \label{eq:wick-learner-oracle}
 \left[
 \sup_{\mathscr C_W\in\mathcal C_{Q,W}^{\mathfrak A}(B,a)}
 V^W_{\mathscr C_W}
 -\varepsilon-\varepsilon_{\mathrm{net}}-\varepsilon_{\mathrm{stat}}
 \right]_+,
\end{equation}
where the last two terms are those derived extended perturbation and
confidence radii and \([r]_+:=\max\{0,r\}\), with an infinite subtracted
radius giving zero.  Its threshold-first output contains every truly feasible strict
budget/action point outside these radii and excludes every truly infeasible
point outside them; intersecting intervals are reported as insufficient
confidence.

On the unrestricted-Caus branch, the returned predictor satisfies the
extended oracle inequality \eqref{eq:wick-free-oracle-bound}.  On the RL
branch, every critic iteration satisfies \eqref{eq:wick-rl-critic-oracle},
and the output satisfies the extended inequalities
\eqref{eq:wick-rl-fqi-recursion} and
\eqref{eq:wick-rl-policy-oracle}.  A nonordinary converter or actor coordinate
makes only the corresponding bound infinite and returns the typed confidence,
shift, or policy-conversion status.
For the minimum-working-model option, the empty Part~I-certified feasible set
returns its typed status; every record returned on the complementary cell
satisfies
\eqref{eq:wick-certified-working-bound}--
\eqref{eq:wick-certified-simplicity-bound}; the RL specialization satisfies
\eqref{eq:wick-rl-minimal-working-bound}.  These records are exactly the
recordwise frontiers reported at Pareto-minimal points of the certified
risk-qualified strict inverse.  At fixed response, full structural
coverage uses the minimum capacity \(K=d_W^*\) whenever
\eqref{eq:wick-projection-coverage} is verified; this capacity remains an
inner coordinate of the saturated-profile comparison.

Moreover, \cref{thm:wick-bayesian-comparison} gives
\eqref{eq:wick-profile-inclusion}, and
\begin{equation}
 \label{eq:wick-generalized-inverse}
 G_{Q,W,\mathfrak A}^{\mathrm{sat}}(B,a)>\theta
 \quad\Longleftrightarrow\quad
 (B,a)\in\mathfrak P_{Q,W,\mathfrak A}^{>}(\theta).
\end{equation}
For every \(\epsilon\), the risk-qualified specialization likewise satisfies
\begin{equation}
 \label{eq:wick-risk-generalized-inverse}
 G_{Q,W,\mathfrak A}^{\mathrm{sat},\le\epsilon}(B,a)>\theta
 \quad\Longleftrightarrow\quad
 (B,a)\in
 \mathfrak P_{Q,W,\mathfrak A}^{>,\le\epsilon}(\theta).
\end{equation}
\end{theorem}

\begin{proof}
Finite enumeration at each Part~I stage and the declared tie rule give
termination.  The allocated confidence event and the extended net Lipschitz
bounds compare every population certificate
with its evaluated representative; maximizing the simultaneous lower bound
gives \eqref{eq:wick-learner-oracle}.  The interval rule proves the inverse
claim.  Sample splitting and uniform selection invoke
\cref{thm:wick-free-caus-oracle} on the free branch; fitted Bellman
contraction and Gibbs calibration invoke
\cref{thm:wick-rl-loss-calibration} on the RL branch.  The model-family
selector invokes
\cref{prop:wick-minimum-capacity-readout,thm:wick-simplest-certified-model,cor:wick-free-rl-minimal-models}.
Every record in
\(\mathcal C_{Q,W}^{\mathfrak A}\) has a nonzero
quotient defect evaluated through the full nonreversible resolvent and
therefore is a record in the existing Part~I \(Q,\mathrm{nr}\) family,
proving
\eqref{eq:wick-profile-inclusion}.
Finally, \eqref{eq:wick-generalized-inverse} is the generalized-inverse
identity for a monotone saturated profile and its strict feasibility region;
the same argument proves \eqref{eq:wick-risk-generalized-inverse} after
restriction to the risk-qualified record family.
\end{proof}

\begin{remark}[Totalized structural multipliers]
\label{rem:wick-structural-multipliers}
The inherited optimization-status field separates the finite-dimensional
candidate problems into exhaustive cells.  On the convex
constraint-qualified cell, the KKT multipliers of the resource, action,
capacity, and confidence constraints are shadow prices for those declared
structural resources.  The other cells retain their nonconvexity or failed-
qualification status and make no multiplier claim.  Identification of an
ordinary multiplier with a measured physical coupling constant is a distinct
representation question, not a structural assumption or a conclusion of the
present construction.
\end{remark}

\subsection{Empirical implementation by differentiable linear algebra}
\label{subsec:wick-empirical-implementation}

The preceding learner is directly implementable on each finite chart in the
Part~I stage.  This subsection specifies the tensor representation, the
differentiable solves, and the separation between inner training and outer
structural selection.  Every branch uses the complete source-resolved
learning ledger of Part~I.  The RL branch additionally uses its complete
dynamical learning record; when a candidate is active or structural
Bayesian, it also uses the complete neural implementation record.  Thus no
record type is silently enlarged.  In particular, the optimizer output is
one represented Wick record; backpropagation does not replace the saturated
profile by a scalar training objective.

\begin{definition}[Finite empirical Wick interface]
\label{def:wick-finite-empirical-interface}
Fix one candidate specification \(s\) in the charged finite stage
\(\mathscr R_W^{(n)}(B,a)\).  Its Part~I record supplies the following data.
\begin{enumerate}[label=\textup{(\roman*)}]
\item The quotient and complement charts
      \(\mathbb R^{d_Q}\) and \(\mathbb R^{d_Y}\), their gauges, evaluation
      maps, precision convention, and source and authority fields.
\item A bridge chart
      \(\operatorname{Dec}_{\mathfrak A,s}:\Phi_{C,s}\to
      \mathbb R^{d_Y\times d_Q}\).  Its image is the bridge family carried by
      \(s\).  On the unrestricted branch it is the full declared bridge
      chart; on the RL branch it is the inherited legal-action compiler.
\item Positive grids for \(\lambda,\rho,\tau\), the loss weights, the target
      evaluation maps, the legal and safe action masks, and the complete
      confidence allocation.  The record also carries its population
      estimand, probability-law, norm, and loss identifiers and every
      declared converter between them.  These are the grids and types already
      charged in the Part~I candidate, rather than new continuous
      hyperparameter families.
\item Identifiers and the declared overlap relation for the dynamics,
      target-training, validation, RL, and audit carriers used by that
      candidate.  Independence is required exactly when the selected Part~I
      split-sample certificate uses it.  A shared carrier instead uses its
      simultaneous union, cover, PAC--Bayes, sequential, or martingale
      certificate.  A carrier may be absent when its associated certificate
      is not invoked.
\end{enumerate}
The dynamics carrier contains one of the two records
\[
 Z_j=(x_j,y_j,\dot x_j,\dot y_j,W_{x,j},W_{y,j}),
 \qquad
 Z_j^{\rm tr}=(u_j,u_{j+1},\Delta_j,\Omega_j),
\]
used in \eqref{eq:wick-dynamics-training-loss} and
\eqref{eq:wick-transition-training-loss}, respectively.  A supervised carrier stores
\((\Phi_{\rm tr},y_{\rm tr})\) and
\((\Phi_{\rm val},y_{\rm val})\), where each \(\Phi\) is the represented
evaluation map from quotient scores to observations.  An RL carrier stores
\((h_i,a_i,r_i,h_i',\omega_i)\), its terminal indicator, and its legal and
safe masks.  When the record contains an encoder, raw observations first pass
through that inherited map; its reconstruction and deployment errors remain
separate certificate coordinates.

The terminal empirical object is
\begin{equation}
 \label{eq:wick-empirical-record}
 \widehat{\mathscr C}_{W,s}
 =\bigl(s,\widehat\vartheta_s,\widehat f_s,\widehat q_s,
        \widehat\pi_s,\mathsf{Opt}_s,\mathsf{Audit}_s,
        \mathbf c_s,\mathsf{Status}_s\bigr).
\end{equation}
Here absent supervised or RL fields are null slots.  The optimizer transcript
\(\mathsf{Opt}_s\) stores the initialization, iterates, stochastic gradients,
moments, clipping, stopping rule, terminal weights, and the represented
optimization gap.  The audit field stores split identifiers, all evaluated
residuals and intervals, the same-loss risk candidates, lower master
visibility, precision and solve residuals, and the confidence allocation.
The status belongs to the typed list in
\cref{def:wick-structural-learner}, augmented only by the numerical status
already present in the Part~I implementation record.
\end{definition}

\begin{definition}[Structure-preserving tensor compiler]
\label{def:wick-tensor-compiler}
Let \(L_H,L_G,K\) be unconstrained real trainable matrices of the appropriate
sizes and let \(\phi_C\in\Phi_{C,s}\).  Define
\begin{equation}
 \label{eq:wick-tensor-parameterization}
 H=L_HL_H^{\mathsf T},\qquad
 G=L_GL_G^{\mathsf T},\qquad
 A_{\mathcal Y}=\frac12(K-K^{\mathsf T}),\qquad
 C=\operatorname{Dec}_{\mathfrak A,s}(\phi_C).
\end{equation}
This displayed parameterization is the dense trainable cell.  If \(s\)
carries additional sparsity, tying, locality, or coefficient restrictions,
the factors range only over its inherited chart and are passed through its
exact structure compiler.  A finite or nondifferentiable Part~I chart is
enumerated and then held fixed during the differentiable subproblem.  Thus
\eqref{eq:wick-tensor-parameterization} does not enlarge the candidate family.
The displayed square factors parameterize the full dense positive-semidefinite
cones.  A rectangular factor is used only when the inherited candidate
already carries the corresponding rank ceiling.  Nonuniqueness of the
factors is an optimizer-coordinate gauge, whose storage and optimization are
charged in the record; all identification and perturbation bounds concern
the compiled operators \(H,G,A_{\mathcal Y},C\), not a choice of factor.
The real matrices are converted to the declared complex precision only when
the Wick generator or its response is evaluated.  For \(\lambda>0\), put
\begin{align}
 B_C&=\lambda I_{d_Y}+G-A_{\mathcal Y},
 &R_C&=\operatorname{solve}(B_C,C),
 \label{eq:wick-empirical-complement-solve}\\
 \Sigma_C&=C^*R_C,
 &F_C&=\lambda I_{d_Q}+iH+\Sigma_C,
 \label{eq:wick-empirical-self-energy}\\
 A_C&=I_{d_Q}+F_C^*F_C,
 &\operatorname{Resp}_C(V)&=\operatorname{solve}(A_C,V).
 \label{eq:wick-empirical-response-solve}
\end{align}
Thus \(\operatorname{Resp}_C(V)=T_{W,\lambda}(C)V\), without materializing a
matrix inverse.  For an evaluation map \(E\), its finite response Gram matrix
is
\begin{equation}
 \label{eq:wick-empirical-gram-solve}
 K_{C,E}=E\operatorname{Resp}_C(E^*).
\end{equation}
If \(E\) is represented by a redundant row system, its Part~I range
compression is applied before the solve.  A Moore--Penrose calculation uses
the same declared finite-precision range decomposition and stores the range
residual.  No eigenvalue is clipped and no diagonal jitter is added unless
that modification and its induced operator error are fields of \(s\).

For a coefficient vector \(f\in\mathbb C^{d_Q}\), the finite supervised
inner problem is
\begin{equation}
 \label{eq:wick-empirical-krr-normal-equation}
 \widehat f_{C,\rho}
 =\operatorname{solve}\!\left(
   \frac1{m_{\rm tr}}\Phi_{\rm tr}^*\Phi_{\rm tr}+\rho A_C,
   \frac1{m_{\rm tr}}\Phi_{\rm tr}^*y_{\rm tr}
  \right).
\end{equation}
This is the normal equation of
\eqref{eq:wick-free-inner-loss}.  Kernel form
\((K_{C,E}+m_{\rm tr}\rho I)\alpha=y_{\rm tr}\) is equivalent on the
represented evaluation range.  The primal form makes the inverse-response
penalty \(f^*A_Cf\) explicit.
\end{definition}

\begin{proposition}[Faithfulness and differentiability of the tensor compiler]
\label{prop:wick-tensor-compiler-faithfulness}
For every finite empirical specification \(s\), the compiler of
\cref{def:wick-tensor-compiler} has the following properties.
\begin{enumerate}[label=\textup{(\roman*)}]
\item It satisfies \(H=H^*\succeq0\), \(G=G^*\succeq0\),
      \(A_{\mathcal Y}^*=-A_{\mathcal Y}\), and
      \(C\in\operatorname{im}\operatorname{Dec}_{\mathfrak A,s}\) at every
      optimizer iterate.
\item The matrix \(B_C\) is invertible and
      \(\|B_C^{-1}\|\le\lambda^{-1}\).  The matrix \(A_C\) is Hermitian
      positive definite, \(A_C\succeq I\), and
      \(\|A_C^{-1}\|\le1\).  Hence
      \eqref{eq:wick-empirical-complement-solve}--
      \eqref{eq:wick-empirical-krr-normal-equation} have unique ordinary
      values in exact arithmetic for every \(\lambda,\rho>0\).
\item The compiled block generator is exactly
      \eqref{eq:wick-block-generator} and satisfies
      \eqref{eq:wick-total-energy-balance}.  Its computed self-energy and
      response are exactly the finite-chart restrictions of
      \eqref{eq:wick-self-energy} and \eqref{eq:wick-response-kernel}.
\item The tensor core is differentiable in \(H,G,A_{\mathcal Y},C\) at every
      ordinary finite-precision evaluation point.  Its composition with an
      inherited parameter decoder is differentiated on the decoder's
      represented differentiable cell; a finite or nondifferentiable decoder
      is handled by the inherited enumeration/search transcript.  For
      \(X=B^{-1}D\), its differential is
      \begin{equation}
       \label{eq:wick-differential-linear-solve}
       \dd X=B^{-1}(\dd D-(\dd B)X).
      \end{equation}
      For the transition loss, the differential of the matrix exponential is
      the Fr\'echet formula
      \begin{equation}
       \label{eq:wick-differential-matrix-exponential}
       D\exp(M)[\dot M]
       =\int_0^1e^{(1-t)M}\dot M e^{tM}\,\dd t.
      \end{equation}
      Reverse automatic differentiation through matrix multiplication,
      linear solves, and the matrix exponential therefore evaluates the
      derivatives of the mathematical losses on the differentiable cell.
      Since the trainable parameters are real and the terminal loss is real,
      this is the real pullback of the complex intermediate calculation.
\end{enumerate}
A failed finite-precision solve, a nonfinite loss, or a residual exceeding the
declared tolerance produces the numerical-status branch of
\eqref{eq:wick-empirical-record}.  Every dependent upper-error, risk, or
policy-gap candidate receives \(+\infty\), while every dependent lower
visibility candidate receives zero.  A raw derived coordinate that has no
ordered extended value retains its typed nonordinary status.  The record
remains in the finite stage.
\end{proposition}

\begin{proof}
The first assertion follows from
\eqref{eq:wick-tensor-parameterization}, the exact structured compiler when
present, and the range of the inherited authority decoder.  For
\(v\in\mathbb C^{d_Y}\),
\[
 \operatorname{Re}\langle v,B_Cv\rangle
 =\lambda\|v\|^2+\langle v,Gv\rangle\ge\lambda\|v\|^2.
\]
Finite-dimensional accretivity gives injectivity, hence invertibility, and
the displayed inverse bound.  Moreover, for \(w\in\mathbb C^{d_Q}\),
\(\langle w,A_Cw\rangle=\|w\|^2+\|F_Cw\|^2\), proving the claims about
\(A_C\).  Substitution in the block generator proves the energy identity,
and block elimination gives the two response identities.  Differentiating
\(BX=D\) proves \eqref{eq:wick-differential-linear-solve}; the standard
variation-of-constants identity proves
\eqref{eq:wick-differential-matrix-exponential}.  Composition and the chain
rule give the automatic-differentiation assertion on the represented
differentiable cell; the other chart types retain their Part~I search rule.
The last statement is
the extended-value convention of
\cref{prop:wick-no-new-premise-closure} applied to the numerical field of the
complete implementation record.
\end{proof}

\paragraph{Dynamics fit.}
The following PyTorch-style pseudocode describes the inner dynamics step.
\begin{quote}
\small\ttfamily
for spec in part1\_stage:\\
\quad model = WickModule(spec)\\
\quad for batch in spec.dynamics\_loader:\\
\qquad H = Lh @ Lh.T;\quad G = Lg @ Lg.T\\
\qquad Ay = (K - K.T) / 2;\quad C = spec.decode\_caus(phi\_c)\\
\qquad Lw = wick\_generator(cplx(H), cplx(G), cplx(Ay), cplx(C))\\
\qquad L = derivative\_loss(batch, H, G, Ay, C)\\
\qquad \# or: pred = expm\_action(dt * Lw, batch.u, tol=spec.solve\_tol)\\
\qquad \# \phantom{or:} L = transition\_loss(batch, pred)\\
\qquad optimizer.zero\_grad();\quad L.backward();\quad optimizer.step()\\
\quad store\_optimizer\_transcript(model, optimizer)
\end{quote}
The two branches evaluate
\eqref{eq:wick-dynamics-training-loss} and
\eqref{eq:wick-transition-training-loss}, respectively.  The terminal
design matrix \(Q_n\), score \(s_n\), range projector \(P_{\rm id}\), kernel,
and optimization gap are then evaluated from the stored trajectory; they are
not training assumptions.
The call \texttt{expm\_action} computes
\(e^{\Delta\mathcal L_W}u\) without forming the exponential.  A dense Pad\'e
exponential remains available on a small chart; a large sparse or matrix-free
chart uses the represented Krylov exponential-action backend and stores its
forward residual or error enclosure.

\paragraph{Unrestricted-Caus fit.}
For every charged \((s,\lambda,\rho)\), response and readout fitting use the
same bridge parameter.  Gradients of the validation term pass through
\eqref{eq:wick-empirical-krr-normal-equation}.
\begin{quote}
\small\ttfamily
for spec, lam, rho in charged\_free\_candidates:\\
\quad for outer\_step in spec.bridge\_steps:\\
\qquad C = spec.decode\_caus(phi\_c)\\
\qquad B = lam * Iy + G - Ay;\quad R = linalg.solve(B, C)\\
\qquad F = lam * Iq + 1j * H + C.mH @ R\\
\qquad A = Iq + F.mH @ F\\
\qquad M = Phi\_tr.mH @ Phi\_tr / m\_tr + rho * A\\
\qquad f = linalg.solve(M, Phi\_tr.mH @ y\_tr / m\_tr)\\
\qquad action = squared\_current\_distance(C, target\_current, spec.metric) / 2\\
\qquad penalty = spec.selection\_penalty(C, rho)\\
\qquad loss = validation\_risk(f) + action\_weight * action + grad\_part(penalty)\\
\qquad bridge\_optimizer.zero\_grad();\quad loss.backward()\\
\qquad bridge\_optimizer.step()\\
\quad detach\_and\_audit(spec, C, f)
\end{quote}
This is the dense small-chart backend.  The scalable backend replaces
\texttt{B}, \texttt{F}, \texttt{A}, and \texttt{M} by the operator actions in
\cref{def:wick-matrix-free-backend} and replaces each \texttt{linalg.solve} by
its residual-controlled Krylov solve.  It never forms
\(\Phi_{\rm tr}^*\Phi_{\rm tr}\), \(F_C^*F_C\), or an inverse.  The dense and
matrix-free paths solve the same normal equation and differ only by their
recorded numerical error.
Its corresponding implementation is:
\begin{quote}
\small\ttfamily
ops = compile\_linear\_operators(spec, parameters)\\
if spec.augmented\_backend:\\
\quad f, z, solve\_info = solve\_augmented\_kkt(ops, Phi\_tr, y\_tr, rho)\\
else:\\
\quad rhs = Phi\_tr.rmatvec(y\_tr) / m\_tr\\
\quad f, solve\_info = krylov\_solve(ops.apply\_M(rho), rhs, spec.solver)\\
loss = validation\_risk(Phi\_val.matvec(f)) + action\_term + selection\_term\\
hypergradient = implicit\_adjoint(loss, solve\_info, ops)\\
update\_bridge(hypergradient); store(solve\_info)
\end{quote}
The data-map handles stream their matrix--vector and adjoint products, so this
code does not retain the complete design on the accelerator.
The selection penalty and the indicator of
\(\mathfrak A_{\rm dyn}\) in \eqref{eq:wick-free-outer-loss} are certificate
and parameter-domain terms.  They are added to the displayed loss when they
vary differentiably within the candidate.  Otherwise the Part~I finite net or
search transcript compares the detached terminal candidates using their full
outer value, including the penalty; \texttt{grad\_part} is then constant
during that inner step.  If zero action is a hard task
requirement, Part~I first evaluates the represented fiber
\(\mathfrak A_D^0\).  On its nonempty cell the decoder is restricted to that
fiber and the action is identically zero.  On its empty cell the candidate
receives the distance in \eqref{eq:wick-action-distance}; training does not
replace that value by a penalty surrogate.  If the task declares a
risk--action tradeoff, the displayed action weight is the declared
\(\lambda_{\mathscr W}\) of \eqref{eq:wick-free-outer-loss}.

\paragraph{Action-safe RL fit.}
Let \(\mathscr E_{\rm RL}\) be the complete finite critic-evaluation carrier,
with the weights and normalization of its selected Part~I statistical
certificate.  Its represented safe row map
\(E_{\mathscr E}^{\mathcal A,{\rm safe}}:\mathcal Q\to
\mathbb C^{r_{\mathscr E}}\) is the full-carrier evaluation map restricted
to the legal, failure-killed action rows.  Apply the inherited exact row-range
compression when those rows are redundant.  The finite safe response is
\begin{equation}
 \label{eq:wick-empirical-safe-response}
 K_{C,\mathscr E}^{\rm safe}
 =E_{\mathscr E}^{\mathcal A,{\rm safe}}
   \operatorname{Resp}_C
   \bigl((E_{\mathscr E}^{\mathcal A,{\rm safe}})^*\bigr).
\end{equation}
This is exactly the finite compression in
\eqref{eq:wick-empirical-response-compression} applied to
\eqref{eq:wick-action-safe-kernel}.  It is positive definite after the
declared row-range compression: for \(z\ne0\),
\[
 z^*K_{C,\mathscr E}^{\rm safe}z
 =\left\langle
 (E_{\mathscr E}^{\mathcal A,{\rm safe}})^*z,
 T_{W,\lambda}(C)
 (E_{\mathscr E}^{\mathcal A,{\rm safe}})^*z
 \right\rangle>0.
\]
Without compression it is positive semidefinite and the already declared
Moore--Penrose range calculation is used.  This Gram representation is used
for a dual solve only when its row dimension is the cheaper represented
coordinate.  The default large-carrier representation is primal:
\[
 q_f=E_{\mathscr E}^{\mathcal A,{\rm safe}}f,\qquad
 \|q_f\|_{T_C^{\mathcal A,{\rm safe}}}^2=f^*A_Cf
 \quad\text{for }f\in
 (\ker E_{\mathscr E}^{\mathcal A,{\rm safe}})^{\perp_{A_C}},
\]
with the inherited minimum-energy range gauge imposed.  This is the same minimum-lift
identity as \cref{prop:wick-quotient-energy-minimality}.  It permits minibatch
evaluation of the data term while retaining one global response penalty; no
minibatch Gram inverse is introduced.  The full Gram matrix is therefore an
audit or small-dual object, not a training requirement.  One fitted iteration
is:
\begin{quote}
\small\ttfamily
f\_target = detach(f)\\
require(nonempty\_safe\_rows, status=authority\_obstruction)\\
with no\_grad():\\
\quad next\_q = real\_score(E\_next\_safe @ f\_target)\\
\quad next\_safe\_value = masked\_max(next\_q, safe\_mask)\\
\quad target = reward - alpha * bridge\_action\\
\qquad\qquad + (1 - terminal) * gamma * next\_safe\_value\\
f, critic\_info = solve\_weighted\_response\_krr(\\
\quad E\_taken, target, importance\_weights, rho, backend=spec.critic\_backend)\\
with no\_grad():\quad gibbs = safe\_gibbs(real\_score(E\_actor @ f), neutral\_row, tau)\\
actor\_loss = lambda\_pi * mean(kl(actor.safe\_row, gibbs))\\
actor\_optimizer.zero\_grad();\quad actor\_loss.backward()\\
actor\_optimizer.step()\\
if spec.differentiable\_model\_cell:\\
\quad model\_loss = frozen\_critic\_component(detach(f), C)\\
\qquad\qquad + lambda\_dyn * dynamics\_loss(model, dynamics\_batch)\\
\qquad\qquad + lambda\_num * numerical\_loss(model, numerical\_batch)\\
\quad model\_optimizer.zero\_grad();\quad model\_loss.backward()\\
\quad model\_optimizer.step()\\
store\_lag\_support\_shift\_and\_optimizer\_errors()
\end{quote}
Here \(f\) is a quotient-score coefficient vector or a batched represented
coefficient field supplied by the inherited critic chart; every displayed
\(E\)-map has the corresponding batch axes.  The safe mask is applied before both the target maximum and Gibbs
normalization.  The target tensor and the Gibbs target are detached in their
respective alternating steps.  The represented masked reduction is defined
only after the nonempty-safe-row check and therefore introduces no artificial
infinite logit.  The bridge action is charged either in the
reward or as a hard legal-action constraint, exactly as in
\cref{def:wick-rl-training-loss}; it is never counted twice.  When the bridge
or dynamics block is trainable, its permitted update uses
\eqref{eq:wick-rl-composite-loss} after the critic and actor steps, with the
critic and actor arguments frozen as specified there, and stores the induced
target-network lag and model-selection error.  Every component and displayed
weight is retained in the terminal objective.  Backpropagation is used only
on its represented differentiable slice; masks, validation rules, finite
decoders, and nondifferentiable audit coordinates are held fixed or searched
by the inherited finite transcript and are evaluated again in the detached
terminal record.
For a fixed response, \texttt{solve\_weighted\_response\_krr} uses the dense,
primal matrix-free, dual, or augmented backend to solve the convex critic
problem directly.  A nonlinear inherited critic chart instead retains its
declared WAdam transcript and the same inverse-response penalty, but it does
not receive the exact KRR-minimizer claim.  The target, actor, and bridge
updates remain alternating, so no backward graph is retained across fitted
Bellman iterations.

\begin{definition}[Detached audit and saturated empirical selector]
\label{def:wick-detached-empirical-selector}
After every terminal inner fit, detach all terminal tensors and evaluate the
complete record on its assigned audit carriers.  The audit performs the
following operations in order.
\begin{enumerate}[label=\textup{(A\arabic*)}]
\item Evaluate the generator residual, solve residuals, lumpability defect,
      dark and bracket coordinates, self-energy, response, empirical
      effective dimension, target capture, action, and target-subspace
      coverage gap.
\item Evaluate every applicable Part~I statistical candidate after conversion
      to the same population loss, and set \(U_W\) to their recordwise
      minimum as in
      \cref{def:wick-inherited-family-selection-certificate}.  The
      Rademacher union/cover, PAC--Bayes mixture, split-validation, mixing,
      sequential, and channel candidates retain their own applicability and
      confidence fields.  All candidates and every data-dependent terminal
      selection are covered by one output-selection-independent simultaneous
      event allocated before selection.  A confidence interval chosen only
      after observing the terminal output is not a certificate.  An
      unavailable candidate equals \(+\infty\).  Records with different
      population laws, estimands, norms, or unconverted loss units remain in
      different typed partitions and their upper bounds are never minimized
      against one another.
\item Evaluate the existing Part~I master visibility with the Wick stability,
      full response, actual authorized bridge, and same-record interval
      fields:
      \begin{equation}
       \label{eq:wick-empirical-master-visibility}
       V_M^W
       =\overline A_q S_W(\lambda)M_{\mathfrak G_q}
        C_{\mathfrak G_q}
        \frac{R_{\mathfrak G_q}}{1+\rho_{\mathfrak G_q}}.
      \end{equation}
      Apply the Part~I monotone interval evaluator to obtain
      \(\underline V_W\).  A missing qualification or unauthorized current
      gives lower visibility zero, as in the master certificate.
\item Store the complete resource vector, Caus action, risk and visibility
      intervals, confidence allocation, and typed status.  The computation,
      architecture, parameter, training, optimization, precision, storage,
      deployment, evaluation, and upper-audit costs are stored in the
      corresponding slots of the Part~I complete learner--deployment ledger.
      On an active or structural-Bayesian candidate these slots are exactly
      the complete neural implementation cost; on an RL candidate they also
      populate the complete dynamical learning record.
\end{enumerate}
For each fixed supervised loss type \(\ell\), let
\(\mathscr R_{W,\ell}^{(n)}(B,a)\) be its same-law, same-estimand, same-norm
partition.  At fixed \((B,a,\ell)\), return the \(\varepsilon\)-maximizers
of \eqref{eq:wick-certified-risk-profile} restricted to that partition.
Across its charged resource--action grid, return all Pareto-minimal points of
\eqref{eq:wick-certified-risk-inverse}.  The RL value-performance certificate
is already the separately calibrated \(U_{\rm RL}\) of
\eqref{eq:wick-rl-working-certificate}.  Thus a requested
\((\epsilon,\theta)\) has the two typed returns
\begin{subequations}
 \label{eq:wick-empirical-pareto-return}
 \begin{align}
 \widehat{\mathscr M}_{W,\ell}^{\rm free}
 &:={\operatorname{Min}}_{\preceq\mathbf c}
 \{M\in\mathscr R_{W,\ell}^{(n)}(B,a):
       U_W(M)\le\epsilon,\ \underline V_W(M)>\theta\},
 \label{eq:wick-empirical-pareto-return-free}\\
 \widehat{\mathscr M}_{W}^{\rm RL}
 &:={\operatorname{Min}}_{\preceq\mathbf c}
 \{M\in\mathscr R_{W,\rm RL}^{(n)}(B,a):
       U_{\rm RL}(M)\le\epsilon,\ \underline V_W(M)>\theta\}.
 \label{eq:wick-empirical-pareto-return-rl}
 \end{align}
\end{subequations}
The two sets may be identified only when a converter already carried by
Part~I places both guarantees on the same population law, estimand, norm, and
loss scale.  In that case the converted records form one named partition;
the conversion error is charged before the recordwise minimum.
An empty set returns its inherited typed status.
\end{definition}

The outer computation can be expressed without differentiating through the
structural order:
\begin{quote}
\small\ttfamily
records = []\\
for spec in part1\_stage:\\
\quad terminal = fit\_inner\_wick\_record(spec)\\
\quad records.append(certify\_detached(terminal, spec.audit\_splits))\\
for loss\_id, same\_loss\_records in typed\_partitions(records):\\
\quad for B, action\_ceiling in charged\_grid(same\_loss\_records):\\
\qquad feasible = [r for r in same\_loss\_records\\
\qquad\quad if resource\_leq(r.cost, B) and r.action <= action\_ceiling\\
\qquad\quad and working\_ucb(r) <= epsilon\\
\qquad\quad and r.visibility\_lcb > theta]\\
\qquad report(loss\_id, pareto\_minimal(\\
\qquad\qquad feasible, key=complete\_resource\_action))
\end{quote}
Here \texttt{working\_ucb} dispatches to \(U_W\) on a supervised partition
and to \(U_{\rm RL}\) on an RL value-performance partition.  The partition
key contains the population law, estimand, norm, and loss scale, so the code
cannot compare certificates with different meanings.
No weighted sum of risk, visibility, capacity, action, and computation occurs
in this outer loop.  A weighted loss may generate a record only when that
tradeoff is declared inside the task.  The record is still compared by the
same saturated profile and Pareto order as every other Part~I candidate.

\begin{table}[tbp]
\centering
\small
\begin{tabular}{@{}p{.20\linewidth}p{.35\linewidth}p{.35\linewidth}@{}}
\toprule
Empirical output & Evaluated Wick coordinate & Quantitative conclusion\\
\midrule
Dynamics fit
& \(Q_n,s_n,P_{\rm id},\ker Q_n\), generator and compression residuals
& \eqref{eq:wick-identification-bound},
  \eqref{eq:wick-feshbach-perturbation}, and
  \eqref{eq:wick-kernel-perturbation}\\
Response audit
& \(E_{W,\Delta},\Sigma_C,F_C,T_C\), dark/bracket status, empirical traces,
  capture and coverage
& \eqref{eq:wick-empirical-defect-interval}--
  \eqref{eq:wick-capture-interval}, with the full alternatives in
  \cref{cor:wick-quantitative-ledger}\\
Unrestricted Caus
& inverse-response ridge fit, actual action, bridge cover and validation
  radius
& \eqref{eq:wick-free-oracle-bound}, including operator, selection, net, and
  optimization errors\\
RL
& safe effective dimension, critic residual, converter, target lag, policy
  KL, authority and observation gaps
& \eqref{eq:wick-rl-critic-oracle},
  \eqref{eq:wick-rl-fqi-recursion},
  \eqref{eq:wick-rl-policy-oracle}, and
  \eqref{eq:wick-rl-working-certificate}\\
Statistical audit
& Rademacher union/cover, PAC--Bayes mixture, dependence and channel
  candidates in one typed same-loss partition
& their recordwise minimum \(U_W\), the source-resolved minimum
  \eqref{eq:wick-source-resolved-target-minimum}, and
  \cref{thm:wick-bayesian-comparison}\\
Structural return
& \(\underline V_W\), typed \(U_W\) or \(U_{\rm RL}\),
  \(\mathbf c_{\rm res},\mathbf c_{\rm act}\)
& \eqref{eq:wick-certified-risk-profile},
  \eqref{eq:wick-certified-risk-inverse}, and
  \eqref{eq:wick-empirical-pareto-return-free}--
  \eqref{eq:wick-empirical-pareto-return-rl}\\
\bottomrule
\end{tabular}
\caption{Empirical outputs and their quantitative certificate coordinates.
The table is an execution map for the exhaustive ledger in
\cref{cor:wick-quantitative-ledger}; it introduces no additional bound.}
\label{tab:wick-empirical-quantitative-ledger}
\end{table}

\begin{theorem}[Certified empirical realization of the Wick learner]
\label{thm:wick-certified-empirical-realization}
Apply the preceding procedure to the charged finite Part~I stage.  Then:
\begin{enumerate}[label=\textup{(\roman*)}]
\item every ordinary terminal execution compiles to a complete
      source-resolved learning ledger and a complete Wick certificate.  An
      RL execution also compiles to a complete dynamical learning record,
      and an active or structural-Bayesian execution compiles to a complete
      neural implementation record.  Every nonordinary execution compiles to
      the corresponding typed records;
\item the dynamics and unrestricted-Caus tensor losses equal
      \eqref{eq:wick-dynamics-training-loss} or
      \eqref{eq:wick-transition-training-loss},
      \eqref{eq:wick-free-inner-loss}--
      \eqref{eq:wick-free-outer-loss}, respectively, on their represented
      finite carriers.  The alternating RL steps are exactly the critic,
      actor, dynamics, and numerical components of
      \eqref{eq:wick-rl-composite-loss}, with their displayed weights and
      frozen arguments;
\item an Adam or block-preconditioned update has the represented form
      \begin{equation}
       \label{eq:wick-empirical-wadam-update}
       \theta_{k+1}
       =\theta_k-\alpha_kW_k^{-1}\widehat m_k,
      \end{equation}
      and is a diagonal or block instance of the Part~I certified WAdam
      record.  Its stationarity and global-gap conclusions are exactly those
      of the Part~I WAdam descent theorem when that record certifies its
      alignment, second-moment, smoothness, and, for the global conclusion,
      PL coordinates.  Otherwise the actual terminal optimization gap
      remains an extended record coordinate.  Even on the global-gap cell,
      that objective gap enters the saturated error only through the
      objective-to-action, objective-to-risk, or objective-to-visibility
      converter already carried by Part~I; without an ordinary converter it
      is only an optimization coordinate;
\item on the simultaneous Part~I confidence event, the returned free-Caus
      records satisfy \eqref{eq:wick-free-oracle-bound}, and the returned RL
      records satisfy \eqref{eq:wick-rl-critic-oracle}--
      \eqref{eq:wick-rl-policy-oracle}; and
\item within each typed same-loss partition, the outer return is exactly the
      Part~I Pareto-minimal certified risk-qualified Wick frontier.  The
      supervised return satisfies
      \eqref{eq:wick-certified-working-bound}--
      \eqref{eq:wick-certified-simplicity-bound}; the RL return satisfies
      \eqref{eq:wick-rl-minimal-working-bound}.  Relative to a PAC--Bayes-only
      supervised selector, its feasible profile and strict inverse obey
      \eqref{eq:wick-pb-profile-domination} and
      \eqref{eq:wick-pb-inverse-domination}; and
\item the empirical compiler adds no admission premise.  Its charts, laws,
      losses, converters, masks, confidence event, costs, and search order
      are fields of the Part~I record; the Wick generator, solves, losses,
      derivatives, and audit coordinates are finite compositions or
      totalized evaluations of those fields.
\end{enumerate}
Consequently the empirical use of backpropagation changes only the method
used to produce and record inner candidates.  The meaning of optimal,
minimal, and simplest remains the saturated-profile meaning of
\cref{def:wick-saturated-optimality}.
\end{theorem}

\begin{proof}
The tensor compiler proposition proves structural fidelity and identifies
the finite response with the response used in the theoretical losses.
Expansion of the residual and ridge quadratic forms proves the loss
identities.  The full-carrier safe compression in
\eqref{eq:wick-empirical-safe-response} is
\eqref{eq:wick-empirical-response-compression} applied to
\eqref{eq:wick-action-safe-kernel}; its principal minibatch compressions
therefore retain the same response and normalization.  Detaching the target
produces the fitted supervised critic problem in
\eqref{eq:wick-rl-critic-loss}, and the actor step is
\eqref{eq:wick-rl-gibbs-policy}.

The stored optimizer state and update have the form
\eqref{eq:wick-empirical-wadam-update}; the Part~I WAdam conclusions follow
precisely on the cells where its derived certificate coordinates are
ordinary.  Part~I's calibration step, rather than the optimizer theorem,
then determines whether its objective gap contributes to an action, risk, or
visibility bound.  The detached audit uses the same data splits, preallocated
simultaneous confidence event, source index, authority record, response, and
loss as the terminal model.  It therefore constructs the quantitative bounds of
\cref{cor:wick-quantitative-ledger,thm:wick-free-caus-oracle,thm:wick-rl-loss-calibration}.
The supervised selector
\eqref{eq:wick-empirical-pareto-return-free} is the selector in
\cref{thm:wick-simplest-certified-model}; the free and RL conclusions follow
from \cref{cor:wick-free-rl-minimal-models}, with the latter using
\eqref{eq:wick-empirical-pareto-return-rl}.  The PAC--Bayes comparison is
\cref{thm:wick-bayesian-comparison}.  Typed partitions prevent invalid
cross-loss minima.  Exact inheritance and no-new-premise closure are
\cref{prop:wick-exact-inheritance,prop:wick-no-new-premise-closure}; finite
linear algebra merely evaluates their derived coordinates.  Extended values
and typed statuses give the complementary numerical and statistical cells.
\end{proof}

\subsection{Executable refinement contract}
\label{subsec:wick-executable-refinement}

An implementation must refine one charged finite Part~I stage.  It is not a
procedure for constructing an uncharged continuum family.  The following
contract fixes the remaining programming choices and makes the return value of
every execution branch explicit.

\begin{definition}[Typed program objects]
\label{def:wick-typed-program-objects}
For each candidate \(s\in\mathscr R_W^{(n)}(B,a)\), the program stores six
typed objects.  \texttt{WickSpec} is immutable; every other object becomes an
immutable serialized value when it is attached to the terminal record.
\begin{enumerate}[label=\textup{(\roman*)}]
\item \texttt{WickSpec} stores the stage and candidate identifiers, dimensions,
      real and complex dtypes, device, authority decoder, all represented
      grids, loss and law identifiers, confidence allocation, carrier
      identifiers and overlap relation, dense or matrix-free backend, linear
      solver, restart rule, preconditioner, primal/dual choice, trace-audit
      method, precision schedule, numerical tolerances, memory ceiling, maximum
      iteration counts, and tie rule.
\item \texttt{WickParameters} stores the real trainable tensors
      \(L_H,L_G,K,\phi_C\).  Any inherited encoder, readout, critic, or actor
      parameters occupy separately named fields.
\item \texttt{WickOperators} stores either the compiled tensors
      \(H,G,A_{\mathcal Y},C,B_C,R_C,\Sigma_C,F_C,A_C\) or their exact
      \texttt{LinearOperator} actions, together with dtypes, shapes, adjoint
      actions, backend identifiers, and no optimizer state.
\item \texttt{CarrierBundle} stores the represented dynamics, training,
      validation, RL, and audit carriers together with their identifiers,
      masks, weights, terminal fields, and normalization constants.
\item \texttt{FitTranscript} stores the initialization, random seeds,
      minibatch order, optimizer states, iterates, stopping reason, terminal
      parameters, objective values, and optimization-gap candidate.
\item \texttt{AuditRecord} stores the detached residuals, statistical
      candidates, confidence event identifier, lower visibility, resource and
      action vectors, and typed status.  A successful audit constructs
      \eqref{eq:wick-empirical-record}; a failed audit constructs the same
      record with its typed failure status.
\end{enumerate}
The immutable specification is hashed and its identifier is copied into every
other object.  Thus a tensor, loss, or certificate produced under one
specification cannot be attached to another candidate without a type failure.
\end{definition}

The dense finite-chart shapes are fixed by the following table.  A structured
decoder may use smaller parameter tensors, but its output has the displayed
shape.
\begin{table}[htbp]
\centering
\small
\begin{tabular}{@{}lll@{}}
\toprule
Object & Shape and dtype & Mathematical value\\
\midrule
\(L_H,H,F_C,A_C\) & \(d_Q\times d_Q\), real for \(L_H,H\), complex otherwise
  & \(H=L_HL_H^{\mathsf T}\), \eqref{eq:wick-empirical-self-energy}--
    \eqref{eq:wick-empirical-response-solve}\\
\(L_G,G,K,A_{\mathcal Y},B_C\) & \(d_Y\times d_Y\), real before compilation
  & \(G=L_GL_G^{\mathsf T}\),
    \(A_{\mathcal Y}=(K-K^{\mathsf T})/2\)\\
\(C,R_C\) & \(d_Y\times d_Q\), real \(C\), complex \(R_C\)
  & authority-decoded bridge and \(B_C^{-1}C\)\\
\(\Sigma_C\) & \(d_Q\times d_Q\), complex
  & \(C^*R_C\)\\
\(E\) & \(r\times d_Q\), complex
  & represented evaluation map\\
\(K_{C,E}\) & \(r\times r\), complex Hermitian
  & \(E A_C^{-1}E^*\)\\
\(\Phi_{\rm tr},y_{\rm tr}\) & \(m_{\rm tr}\times d_Q\), \(m_{\rm tr}\)
  & supervised design and target\\
\(q_{\rm safe}\) & \(r_{\mathscr E}\), real
  & inherited observable-score slice\\
\bottomrule
\end{tabular}
\caption{Code-level shape contract.  The complex dtype has the same precision
as the declared real dtype.  In particular, float64 is paired with complex128.
Every adjoint in the response calculation is a conjugate transpose.}
\label{tab:wick-code-shape-contract}
\end{table}
In a matrix-free record, each table entry is a shape-checked operator handle
with \texttt{matvec}, \texttt{rmatvec}, and, when available,
\texttt{matmat}; the handle need not own a dense tensor of that shape.
Training and RL carriers are streamed in declared batches.  The audit may make
additional deterministic passes over an auditable carrier, but it does not
require that carrier to reside in accelerator memory.

The code boundary from the complex response space to supervised real outputs
is the observation map carried by \(s\).  A target may instead be represented
as a complex observation, in which case the loss is the modulus-squared loss
in \eqref{eq:wick-free-inner-loss}.  The RL observation map has real range by
\cref{def:wick-rl-training-loss}; the program applies a maximum or exponential
only after that map.  Calling \texttt{real} on an arbitrary complex score is
not an admissible replacement for this typed observation map.

\begin{definition}[Typed numerical interface]
\label{def:wick-pure-numerical-interface}
The numerical core exposes the following interfaces; each output retains the
input specification identifier.
\begin{quote}
\small\ttfamily
compile(spec, parameters) -> WickOperators | TypedFailure\\
response(spec, operators, evaluation\_map) -> ResponseHandle | TypedFailure\\
fit\_free(spec, operators, carriers) -> FitTranscript | TypedFailure\\
fit\_rl(spec, operators, carriers) -> FitTranscript | TypedFailure\\
audit(spec, operators, transcript, carriers) -> AuditRecord\\
select(stage, audit\_records, epsilon, theta) -> ParetoFrontier | TypedFailure
\end{quote}
Here \texttt{compile} evaluates
\eqref{eq:wick-tensor-parameterization} and
\eqref{eq:wick-empirical-complement-solve}--
\eqref{eq:wick-empirical-response-solve}; \texttt{response} returns a dense
Gram handle or a matrix-free operator handle and evaluates
\eqref{eq:wick-empirical-gram-solve} or
\eqref{eq:wick-empirical-safe-response}; \texttt{fit\_free} evaluates
\eqref{eq:wick-free-inner-loss}--\eqref{eq:wick-free-outer-loss}; and
\texttt{fit\_rl} evaluates the frozen critic, actor, dynamics, and numerical
blocks of \eqref{eq:wick-rl-composite-loss}.  The audit and selector are
\cref{def:wick-detached-empirical-selector}.  These functions do not infer a
grid, authority envelope, loss converter, safe mask, confidence split, or
stopping rule from the observed objective values.
The compiler, response, audit, and selector are pure functions of their stated
inputs.  The two fitting interfaces are deterministic when the stored seed,
minibatch order, optimizer backend, and deterministic-backend flag are fixed;
otherwise their complete realized transcript is part of the return value.

The runtime data-access validator checks every requested carrier identifier
against the overlap relation in \texttt{WickSpec}.  A split-sample candidate
requires the declared disjoint carriers.  A shared carrier is accepted only
for the corresponding simultaneous, uniform, sequential, martingale, or
PAC--Bayes candidate.  The confidence allocation and candidate index are
sealed before the first target-bearing carrier is read.
\end{definition}

\begin{definition}[Dense, nested matrix-free, and augmented backends]
\label{def:wick-matrix-free-backend}
The backend tag of \texttt{WickSpec} selects one of three representations of the
same finite-chart operators.
\begin{enumerate}[label=\textup{(\roman*)}]
\item The dense backend materializes the matrices in
      \eqref{eq:wick-empirical-complement-solve}--
      \eqref{eq:wick-empirical-krr-normal-equation}, uses factorizations with
      multiple right-hand sides, and caches a factorization only while every
      operator parameter in its cache key is unchanged.
\item The matrix-free backend stores the actions of
      \(L_H,L_G,K,C,C^*,\Phi,\Phi^*\).  With all adjoints taken in the declared
      complex inner product, it applies
      \begin{align}
       \mathsf Hx&=L_H(L_H^*x),&
       \mathsf Gy&=L_G(L_G^*y),&
       \mathsf A_{\mathcal Y}y&=\tfrac12(Ky-K^*y),\notag\\
       \mathsf B y&=\lambda y+\mathsf Gy-\mathsf A_{\mathcal Y}y,&
       \mathsf B^* y&=\lambda y+\mathsf Gy+\mathsf A_{\mathcal Y}y,
       \label{eq:wick-matrix-free-basic-actions}\\
       \mathsf F x&=\lambda x+i\mathsf Hx+
         C^*\operatorname{Solve}_{\mathsf B}(Cx),\notag\\
       \mathsf F^*x&=\lambda x-i\mathsf Hx+
         C^*\operatorname{Solve}_{\mathsf B^*}(Cx),\notag\\
       \mathsf A x&=x+\mathsf F^*(\mathsf Fx),\notag\\
       \mathsf M_\rho x
       &=\frac1{m_{\rm tr}}\Phi^*(\Phi x)+\rho\mathsf A x.
       \label{eq:wick-matrix-free-response-actions}
      \end{align}
      In exact arithmetic the solve symbols in these definitions denote the
      unique solutions.  At runtime they are residual-controlled iterative
      solves and their errors are included in every enclosing operator
      application.  Because an adaptively terminated inner solve makes the
      outer action inexact and iteration-dependent, the nested backend uses a
      flexible outer Krylov method (or fixes the inner tolerance and initial
      state so that the represented action is stationary).  In both cases the
      terminal candidate is accepted only through the independent residual
      audit below.
\item The augmented matrix-free backend eliminates nested solves.  Introduce
      \(z\in\mathbb C^{d_Y}\) and impose
      \(\mathsf Bz=Cf\).  Then the primal free-Caus problem is exactly
      \begin{equation}
       \min_{f,z:\,\mathsf Bz=Cf}
       \left\{
        \frac{\|\Phi f-y\|^2}{2m_{\rm tr}}
        +\frac{\rho}{2}
          \left(\|f\|^2+
          \|(\lambda I+i\mathsf H)f+C^*z\|^2\right)
       \right\}.
       \label{eq:wick-augmented-krr}
      \end{equation}
      Its block KKT operator uses only
      \(\mathsf H,\mathsf B,C,C^*,\Phi,\Phi^*\) actions.  The RL critic has the
      identical form with its safe evaluation map and weighted target.  A
      residual-controlled block solver therefore reaches the same minimizer
      without an inner solve inside every outer matrix--vector product.
\end{enumerate}
The matrix-free code never forms
\(\Phi^*\Phi\), \(F_C^*F_C\), \(K_{C,E}\), or a matrix inverse.  It uses a
block Krylov solve when several right-hand sides share an operator.  A
preconditioner and warm start may change iteration count but not the
mathematical output: their identifiers and construction costs are stored, and
acceptance always evaluates the residual of the original unpreconditioned
equation.

The free-Caus fit solves the smaller represented primal or dual system.  The
primal system is \(\mathsf M_\rho f=\Phi^*y/m_{\rm tr}\); the dual system is
\((K_{C,\Phi}+m_{\rm tr}\rho I)\alpha=y\).  The same rule is used for the RL
critic.  A primal critic stores \(f\) in the minimum-\(A_C\)-energy gauge and
evaluates only the minibatch rows \(E_{\rm batch}f\).  A dual critic stores
\(\alpha\) and evaluates \(K_{\rm batch,train}\alpha\).  The choice and its
conversion residual are fields of the candidate.
The nested and augmented formulations are equal because the constraint in
\eqref{eq:wick-augmented-krr} has the unique solution
\(z=\mathsf B^{-1}Cf\), after which its second squared norm is
\(\|\mathsf Ff\|^2\).
\end{definition}

The numerical method is selected from the proved operator structure, not from
a generic ``Krylov'' flag.  A conforming solver dispatch is:
\begin{center}
\small
\begin{tabular}{@{}p{.20\linewidth}p{.27\linewidth}p{.39\linewidth}@{}}
\toprule
System & Derived structure & Large-chart method and safeguard\\
\midrule
\(\mathsf Bz=v\), \(\mathsf B^*z=v\)
& accretive, generally non-Hermitian
& restarted GMRES; flexible GMRES when the preconditioner or inner accuracy
  varies; audit the original unpreconditioned residual\\
\(\mathsf M_\rho f=b\)
& Hermitian positive definite in exact arithmetic, with
  \(\mathsf M_\rho\succeq\rho I\)
& preconditioned conjugate gradients only for a stationary Hermitian action;
  otherwise flexible GMRES or the augmented formulation\\
\(K+m\rho I\)
& Hermitian positive definite on the compressed row range
& preconditioned or block conjugate gradients, used only when the compressed row
  dimension is favorable\\
augmented KKT system
& Hermitian indefinite under exact adjoint pairing
& MINRES with a Hermitian positive preconditioner; GMRES when finite-precision
  actions or preconditioning do not preserve that pairing\\
\(e^{\Delta\mathcal L_W}u\)
& contractive but generally nonnormal
& adaptive Arnoldi exponential action with time subdivision; a Lanczos-only
  path is not used unless the represented generator is Hermitian\\
\bottomrule
\end{tabular}
\end{center}
The default complement preconditioner uses the positive part
\(\lambda I+G\); response and KKT preconditioners may reuse inherited sparse,
domain-decomposition, or multilevel actions.  A randomized or learned
preconditioner is permitted only as an accelerator: its construction and
state are charged, and acceptance still uses the original operator residual.
Thus a poor preconditioner can increase recorded work or cause a typed
numerical-budget return, but it cannot change the accepted equation.

\begin{proposition}[A posteriori certificate for exponential actions]
\label{prop:wick-expm-action-certificate}
Let \(u(t)=e^{t\mathcal L_W}u_0\) on an inherited ordinary semigroup cell, and
let \(\widetilde u\) be the piecewise differentiable path returned by an
Arnoldi exponential-action computation with defect
\(d(t)=\dot{\widetilde u}(t)-\mathcal L_W\widetilde u(t)\).  Then
\begin{equation}
 \|u(T)-\widetilde u(T)\|
 \le \|u_0-\widetilde u(0)\|+
     \int_0^T\|d(t)\|\,\dd t.
 \label{eq:wick-expm-action-defect-bound}
\end{equation}
Consequently adaptive time subdivision can accept an exponential action from
an outward-rounded quadrature enclosure of its Arnoldi defect, without
forming either \(\mathcal L_W\) or its exponential.  A missing enclosure
returns the numerical-status value and cannot enter a finite certificate.
\end{proposition}

\begin{proof}
Variation of constants gives
\[
 \widetilde u(T)-u(T)
 =e^{T\mathcal L_W}(\widetilde u(0)-u_0)
  +\int_0^T e^{(T-t)\mathcal L_W}d(t)\,\dd t.
\]
The contraction bound derived in
\cref{prop:wick-generation-budget,prop:wick-sectorial-realization}, followed
by the triangle inequality, proves
\eqref{eq:wick-expm-action-defect-bound}.
\end{proof}

\begin{proposition}[Matrix-free adjoint for transition training]
\label{prop:wick-expm-action-adjoint}
Let \(J(u(T),\theta)\) be a real terminal transition loss on the ordinary
cell of \cref{prop:wick-expm-action-certificate}, with
\(\mathcal L_W=\mathcal L_W(\theta)\).  Define
\[
 p(T)=\partial_uJ,\qquad
 p(t)=e^{(T-t)\mathcal L_W^*}p(T).
\]
Then
\begin{equation}
 D_\theta J[\dot\theta]
 =\partial_\theta J[\dot\theta]
  +\operatorname{Re}\int_0^T
    \left\langle p(t),
      D_\theta\mathcal L_W[\dot\theta]u(t)\right\rangle\,\dd t.
 \label{eq:wick-expm-action-adjoint}
\end{equation}
Thus the large-chart backward pass uses actions of
\(\mathcal L_W\), \(\mathcal L_W^*\), and its parameter-directional action;
it need not form a dense Fr\'echet derivative.  The implementation stores a
declared checkpoint schedule and recomputes forward segments during the
adjoint sweep, so peak activation memory is bounded by the checkpoint window.

For a computed forward path \(\widetilde u\) and adjoint path
\(\widetilde p\), let \(e_u(t)\) and \(e_p(t)\) be their defect-based error
enclosures and let
\(\|D_\theta\mathcal L_W[\dot\theta]\|\le L_{\mathcal L}(\dot\theta)\)
be the factored derivative-action enclosure.  Apart from the outward-rounded
quadrature error, substitution in
\eqref{eq:wick-expm-action-adjoint} has error at most
\begin{equation}
 \int_0^T L_{\mathcal L}(\dot\theta)
 \left[e_p(t)(\|\widetilde u(t)\|+e_u(t))
       +\|\widetilde p(t)\|e_u(t)\right]\,\dd t.
 \label{eq:wick-expm-adjoint-error}
\end{equation}
If any enclosure is unavailable, this gradient-error coordinate is
\(+\infty\), exactly as for an inexact linear-solve gradient.
\end{proposition}

\begin{proof}
Differentiate \(\dot u=\mathcal L_W(\theta)u\), pair the tangent equation
with the backward adjoint, and integrate by parts to obtain
\eqref{eq:wick-expm-action-adjoint}.  Subtract the approximate and exact
integrands, add and subtract
\(\langle\widetilde p,D_\theta\mathcal L_W u\rangle\), and apply
Cauchy--Schwarz and
\(\|u(t)\|\le\|\widetilde u(t)\|+e_u(t)\).  Integration and addition of the
quadrature enclosure give \eqref{eq:wick-expm-adjoint-error}.
\end{proof}

\begin{proposition}[Implicit differentiation of inexact solves]
\label{prop:wick-implicit-solve-gradient}
Let \(S_\theta x_\theta=b_\theta\) be any ordinary solve in
\cref{def:wick-matrix-free-backend}, and let \(J(x_\theta,\theta)\) be a real
terminal scalar.  If \(z_\theta\) solves
\[
 S_\theta^*z_\theta=\partial_xJ,
\]
then the real parameter differential is
\begin{equation}
 D_\theta J[\dot\theta]
 =\partial_\theta J[\dot\theta]
  +\operatorname{Re}\langle z_\theta,
       D_\theta b_\theta[\dot\theta]
       -D_\theta S_\theta[\dot\theta]x_\theta\rangle.
 \label{eq:wick-implicit-solve-gradient}
\end{equation}
Hence reverse differentiation requires an adjoint solve but does not store the
forward Krylov iterates.  The implementation stores the primal and adjoint
residuals.  Their validated inverse bounds give forward and adjoint state-error
enclosures exactly as in \eqref{eq:wick-code-forward-error}; these are
propagated to the gradient-error coordinate whenever a Part~I optimizer
certificate uses the computed gradient.  More explicitly, write
\(e_x=\|\widetilde x-x\|\) and
\(e_z=\|\widetilde z-z\|\) for those enclosures, and suppose that a parameter
direction has certified bounds
\(\|D_\theta b[\dot\theta]\|\le L_b(\dot\theta)\) and
\(\|D_\theta S[\dot\theta]\|\le L_S(\dot\theta)\).  The gradient returned by
substituting \((\widetilde x,\widetilde z)\) in
\eqref{eq:wick-implicit-solve-gradient} then has directional error at most
\begin{equation}
 e_z\{L_b(\dot\theta)+L_S(\dot\theta)(\|\widetilde x\|+e_x)\}
 +\|\widetilde z\|L_S(\dot\theta)e_x,
 \label{eq:wick-implicit-gradient-error}
\end{equation}
in addition to the certified error in the explicit derivative
\(\partial_\theta J\).  The derivative-action bounds are computed by the same
factored automatic-differentiation graph as the operator actions and stored in
the optimizer record.  If any required bound is unavailable, the corresponding
gradient-error coordinate is \(+\infty\), so no certificate that needs this
gradient is issued.  A training run may use a looser inner tolerance, but the
detached terminal audit is evaluated at the declared certificate tolerance.
\end{proposition}

\begin{proof}
Differentiate \(S_\theta x_\theta=b_\theta\) to obtain
\(S_\theta\dot x=\dot b-\dot Sx\).  Pair this equation with the adjoint
solution and add the explicit parameter derivative of \(J\), which gives
\eqref{eq:wick-implicit-solve-gradient}.  The residual-error statement follows
by applying the validated inverse bound to the primal and adjoint residuals.
For the last assertion, subtract the approximate bilinear term from the exact
one and write the difference as
\[
 \operatorname{Re}\langle \widetilde z-z,\dot b-\dot Sx\rangle
 -\operatorname{Re}\langle\widetilde z,\dot S(\widetilde x-x)\rangle.
\]
Cauchy--Schwarz, \(\|x\|\le\|\widetilde x\|+e_x\), and the two derivative
bounds give \eqref{eq:wick-implicit-gradient-error}.
\end{proof}

\begin{proposition}[Work and memory of the scalable backend]
\label{prop:wick-scalable-complexity}
Let \(c_H,c_G,c_{A_{\mathcal Y}},c_C,c_{C^*},c_\Phi,c_{\Phi^*}\) be the
recorded costs of one operator action, let \(k_B,k_{B^*},k_M,k_{\exp}\) be the
actual terminal iteration counts, and put
\[
 c_B=c_G+c_{A_{\mathcal Y}}+O(d_Y),\qquad
 c_{B^*}=c_G+c_{A_{\mathcal Y}}+O(d_Y),\qquad
 c_F=c_H+c_C+c_{C^*}+k_Bc_B+O(d_Q),
\]
with \(c_{F^*}\) defined using \(k_{B^*}\).  One application of
\(\mathsf M_\rho\) costs
\begin{equation}
 c_M=c_\Phi+c_{\Phi^*}+c_F+c_{F^*}+O(d_Q),
 \label{eq:wick-matrix-free-matvec-cost}
\end{equation}
and one primal ridge solve costs \(O(k_Mc_M)\).  A transition prediction costs
\(O(k_{\exp}c_{\mathcal L})\), where \(c_{\mathcal L}\) is one block-generator
action.  If \(\chi_{\rm ckpt}\ge1\) is the measured recomputation factor of
the declared adjoint checkpoint schedule, its reverse pass costs
\(O(\chi_{\rm ckpt}k_{\exp}(c_{\mathcal L}+c_{\mathcal L^*}))\) and stores only
the checkpoint states and the current Arnoldi window.  If \(k_{\rm aug}\) is
the augmented KKT iteration count, its solve
cost is
\begin{equation}
 O\!\left(k_{\rm aug}
 [c_H+c_B+c_{B^*}+c_C+c_{C^*}+c_\Phi+c_{\Phi^*}+d_Q+d_Y]\right),
 \label{eq:wick-augmented-solve-cost}
\end{equation}
with no nested iteration multiplier.  If the outer, complement, and augmented
restart widths are \(r_M,r_B,r_{\rm aug}\), restarted nested storage is
\(O(d_Qr_M+d_Yr_B)\) per right-hand-side block and augmented storage is
\(O((d_Q+2d_Y)r_{\rm aug})\), in addition to the structured operator storage.

For comparison, a generic dense backend stores
\(O(d_Q^2+d_Y^2+d_Qd_Y)\) scalars and has cubic factorization cost in the
active primal or dual dimension.  Therefore the implementation selects the
dense backend only when its measured factorization and memory estimate fit the
declared budget; otherwise it uses the matrix-free backend.  These are
accounting identities in the complete resource record, not asymptotic
assumptions about a candidate.
\end{proposition}

\begin{proof}
Equations \eqref{eq:wick-matrix-free-basic-actions}--
\eqref{eq:wick-matrix-free-response-actions} contain one \(\mathsf B\) or
\(\mathsf B^*\) solve in each \(\mathsf F\) or \(\mathsf F^*\) action and one
of each in an \(\mathsf A\) action.  Adding the two data-map actions gives
\eqref{eq:wick-matrix-free-matvec-cost}; multiplication by the recorded
iteration counts gives the nested work bounds.  The KKT action of
\eqref{eq:wick-augmented-krr} contains one of each displayed primitive action,
which gives \eqref{eq:wick-augmented-solve-cost}.  Restarted Krylov methods retain only
their declared basis window.  The transition reverse-pass count follows from
\eqref{eq:wick-expm-action-adjoint}: every recomputed forward or adjoint
substep uses one generator action, and the checkpoint schedule records the
number of such recomputations.  Standard dense storage and factorization
counts give the comparison.
\end{proof}

\begin{definition}[Scalable audit and execution schedule]
\label{def:wick-scalable-audit-schedule}
The implementation uses the following hierarchy without changing the returned
record family.
\begin{enumerate}[label=\textup{(\roman*)}]
\item A finite response effective dimension always has the deterministic
      certificate
      \begin{equation}
       0\le
       \operatorname{Tr}\!\left[K(K+\rho I)^{-1}\right]
       \le\operatorname{rank}K\le\min\{r,d_Q\}.
       \label{eq:wick-cheap-effective-dimension-bound}
      \end{equation}
      This bound requires no Gram materialization.  A small problem may compute
      the trace exactly.  A large problem may use a represented stochastic
      trace or Lanczos-quadrature interval only when its random seed, probe
      count, spectral enclosure, failure probability, and computation cost are
      fields of the Part~I numerical certificate.  Its failure probability is
      included in the simultaneous confidence allocation.  An unvalidated
      randomized estimate is a diagnostic and cannot replace
      \eqref{eq:wick-cheap-effective-dimension-bound} in a risk certificate.
\item Capture, response, defect, and target-band quantities are evaluated by
      operator actions on the target vectors actually used by the certificate.
      A matrix-free eigenvalue, singular-value, projector, dark-space, or
      bracket calculation returns a validated interval from its residual and
      spectral enclosure.  If such an interval is unavailable, the associated
      upper candidate is \(+\infty\) and lower visibility is zero; the
      implementation does not request a full eigendecomposition solely to
      manufacture an ordinary value.
\item Candidates are parallelized only across sealed specifications or
      data-independent operator actions.  Carrier identifiers and random-number
      streams are disjoint or shared exactly as declared by the confidence
      record.  Parallel execution therefore changes wall-clock cost but not
      the statistical event.
\item Factorizations, preconditioners, Krylov subspaces, and warm starts are
      reused across right-hand sides, nearby charged ridge scales, and
      successive bridge iterates.  Every cache key contains the specification
      hash, operator-parameter version, dtype, device, and relevant scale.
      Updating \(H,G,A_{\mathcal Y}\), or \(C\) invalidates every response cache
      depending on it.  Reuse across changed operators is a warm start, not an
      exact cache hit, and is followed by a fresh residual check.
\item Training may use mixed precision and loose adaptive inner tolerances.
      The terminal parameters are recompiled and audited in the declared
      certificate precision.  The discrepancy between the training and audit
      operators, together with every inexact-gradient residual, is stored in
      the optimization and operator-error coordinates.
\item The stage scheduler maintains lower and upper intervals for risk,
      visibility, resources, and action.  It may stop or prune a candidate only
      when those intervals prove infeasibility or Pareto domination by an
      already audited record.  Ambiguous candidates are refined.  Thus
      successive halving, early stopping, branch-and-bound, and asynchronous
      execution are permitted only in this certified form.
\item Training checkpoints contain the specification hash, parameter and cache
      versions, optimizer state, data-stream position, random-number states,
      and current audit intervals.  Resumption therefore continues the same
      transcript rather than creating an undeclared candidate.
\end{enumerate}
The scheduler first applies algebraic and authority checks, then the cheap
bounds in \eqref{eq:wick-cheap-effective-dimension-bound}, then inner fitting,
and finally high-precision certification of the unresolved Pareto set.  This
orders work from cheapest to most expensive while preserving the exhaustive
Part~I frontier.
\end{definition}

\begin{proposition}[Soundness of certified screening]
\label{prop:wick-certified-screening}
At a fixed typed stage, suppose a partial candidate record \(M\) has certified
intervals
\[
 R(M)\in[\underline R_M,\overline R_M],\qquad
 V(M)\in[\underline V_M,\overline V_M],
\]
and lower resource/action vector \(\underline{\mathbf c}_M\).  It may be
discarded for a request \((\epsilon,\theta)\) if either
\(\underline R_M>\epsilon\) or \(\overline V_M\le\theta\), or if an audited
feasible record \(M_0\) satisfies
\[
 \mathbf c(M_0)\preceq\underline{\mathbf c}_M
 \quad\text{and at least one coordinate is strict}.
\]
Such screening leaves the Pareto-minimal feasible frontier unchanged.
\end{proposition}

\begin{proof}
The first two conditions prove that \(M\) cannot satisfy the risk or visibility
constraint.  In the third case every completion of \(M\) has resource/action
vector at least \(\underline{\mathbf c}_M\) and is therefore dominated by the
feasible record \(M_0\).  None can be a Pareto-minimal feasible point.
\end{proof}

\begin{definition}[Residual acceptance and finite-precision transfer]
\label{def:wick-residual-acceptance}
Let tildes denote the tensors returned by the numerical linear algebra.  The
audit stores certified upper enclosures of the absolute residuals
\begin{align}
 r_B&:=\|B_C\widetilde R_C-C\|,
 &r_A(V)&:=\|A_C\widetilde X-V\|,
 \label{eq:wick-code-solve-residuals}\\
 r_{\rm con}&:=\|\mathsf B\widetilde z-C\widetilde f\|,\notag\\
 r_f&:=\left\|
 \left(\frac{\Phi_{\rm tr}^*\Phi_{\rm tr}}{m_{\rm tr}}+\rho A_C\right)
 \widetilde f-\frac{\Phi_{\rm tr}^*y_{\rm tr}}{m_{\rm tr}}
 \right\|,
 \label{eq:wick-code-krr-residual}\\
 r_{\exp}(T)&:=\|u_0-\widetilde u(0)\|
  +\overline{\int_0^T\|\dot{\widetilde u}(t)
      -\mathcal L_W\widetilde u(t)\|\,\dd t},\notag\\
 r_{\rm herm}&:=\|A_C-A_C^*\|,
 &r_{\rm gram}&:=\|K_{C,E}-K_{C,E}^*\|.
 \label{eq:wick-code-structure-residuals}
\end{align}
Such an enclosure may be computed by directed rounding, interval arithmetic,
or a proved backend error bound.  A residual evaluated only in working
precision is stored as a diagnostic and cannot by itself certify acceptance.
The audit also stores scale-normalized versions of these residuals for
diagnostics.
Acceptance uses the declared absolute or mixed absolute-relative tolerances in
\texttt{WickSpec}; no hard-coded tolerance enters the mathematical record.
Put \(\operatorname{Herm}(M)=(M+M^*)/2\).  The structural audit additionally
evaluates
\begin{equation}
 \begin{aligned}
 r_H&=\|H-H^*\|,\qquad
 r_G=\|G-G^*\|,\qquad
 r_{\mathcal Y}=\|A_{\mathcal Y}+A_{\mathcal Y}^*\|,\\
 &\lambda_{\min}(\operatorname{Herm}H),\qquad
 \lambda_{\min}(\operatorname{Herm}G),\qquad
 \lambda_{\min}(\operatorname{Herm}A_C).
 \end{aligned}
 \label{eq:wick-code-compiler-audit}
\end{equation}
Let
\[
 M_f:=\frac{\Phi_{\rm tr}^*\Phi_{\rm tr}}{m_{\rm tr}}+\rho A_C.
\]
The structural action identities or a validated eigensolver produce the record
fields
\begin{equation}
 \underline\alpha_B\le
 \lambda_{\min}(\operatorname{Herm}B_C),\qquad
 \underline\alpha_A\le
 \lambda_{\min}(\operatorname{Herm}A_C),\qquad
 \underline\alpha_f\le
 \lambda_{\min}(\operatorname{Herm}M_f).
 \label{eq:wick-code-coercivity-enclosures}
\end{equation}
For the matrix-free backend, the stored floating entries are interpreted as
exact dyadic coefficients of the factor actions in
\eqref{eq:wick-matrix-free-basic-actions}--
\eqref{eq:wick-matrix-free-response-actions}; those identities give
\(\underline\alpha_B=\lambda\), \(\underline\alpha_A=1\), and
\(\underline\alpha_f=\rho\).  Rounding in their evaluation is included in the
residual enclosures.  A backend that instead materializes rounded matrix
products includes their rounding envelopes before forming the lower bounds.
It may use sharper validated values.
Every nonfinite value, shape or dtype mismatch, failed safe-row check, or
residual outside tolerance produces \texttt{TypedFailure}.  The same rule
applies when one of the three lower bounds in
\eqref{eq:wick-code-coercivity-enclosures} is unavailable or nonpositive.
Jitter, clipping,
projection, or eigenvalue truncation is permitted only when the modified
operator and its perturbation radius are stored as fields of a new candidate.

The validated coercivity bounds give the a posteriori estimates
\begin{equation}
 \|\widetilde R_C-B_C^{-1}C\|
 \le\frac{r_B}{\underline\alpha_B},
 \qquad
 \|\widetilde X-A_C^{-1}V\|
 \le\frac{r_A(V)}{\underline\alpha_A},
 \qquad
 \|\widetilde f-\widehat f_{C,\rho}\|
 \le\frac{r_f}{\underline\alpha_f}.
 \label{eq:wick-code-forward-error}
\end{equation}
The audit propagates these quantities through
\eqref{eq:wick-feshbach-perturbation}--
\eqref{eq:wick-kernel-perturbation} and stores the result as numerical operator
and optimization error.  Thus numerical acceptance never converts a residual
into an exact identity.
An augmented solve is accepted only after the audit evaluates both
\(r_{\rm con}\) and the reduced normal-equation residual \(r_f\).  Its terminal
predictor therefore receives the same \(r_f/\underline\alpha_f\) enclosure as
the nested or dense solve.  An implicit derivative through the indefinite KKT
system uses a validated smallest-singular-value bound for that system or
recomputes the adjoint through the reduced positive-definite equation.
\end{definition}

\begin{definition}[Finite executable learner]
\label{def:wick-finite-executable-learner}
The reference program executes the following finite outer algorithm.  Its
enumeration and selection control flow is deterministic under the declared tie
rule; any permitted backend nondeterminism is confined to a fitting transcript.
\begin{enumerate}[label=\textup{(E\arabic*)}]
\item Load one charged stage and validate every immutable specification,
      carrier identifier, overlap declaration, grid, converter, and confidence
      allocation before fitting.  Run the declared target-free or synthetic
      backend pilot and apply \eqref{eq:wick-backend-dispatch}.
\item For each \(s\), train or enumerate the dynamics candidate according to
      its decoder type and its declared finite stopping rule.  Compile the
      terminal generator and evaluate the identification and response-error
      coordinates.
\item For each charged \((\lambda,\rho,\tau)\) and coefficient vector, call
      \texttt{response}.  Dispatch to \texttt{fit\_free}, \texttt{fit\_rl}, or
      the imposed-PDE evaluation branch according to the authority tag.
\item Stop every optimization loop at its declared maximum iteration count or
      stopping predicate and store the actual stopping reason and terminal
      gap.  A divergent or interrupted loop returns a terminal transcript with
      its optimization status.
\item Detach the terminal tensors and call \texttt{audit}.  The audit recomputes
      all structural objects from the terminal parameters rather than trusting
      cached training values.  It then evaluates every predeclared applicable
      Part~I bound on the common confidence event.
\item Partition accepted and nonordinary records by population law, estimand,
      norm, and loss identifier.  Within each partition call \texttt{select}
      and return the risk-qualified saturated Pareto frontier
      \eqref{eq:wick-empirical-pareto-return}.
\end{enumerate}
All enumerated sets and iteration limits are finite at stage \(n\), so this
algorithm terminates.  Candidate failures remain in the finite audit ledger;
only a malformed stage that cannot be parsed produces a stage-level failure.
Write \(\mathsf C_W=\mathsf{CompleteWickRecord}\).  Its total return type is
\begin{equation}
 \mathsf{WickReturn}
 :=\bigl[
      \mathsf{FiniteLedger}(\mathsf C_W)
      \times\mathsf{ParetoFrontier}(\mathsf C_W)
    \bigr]
   \ \sqcup\ \mathsf{TypedStageFailure}.
 \label{eq:wick-code-total-return-type}
\end{equation}
The directed Part~I stage iterator is external to this finite call and retains
its existing limit and nonattainment semantics.
\end{definition}

\begin{theorem}[Program refinement and certificate soundness]
\label{thm:wick-program-refinement}
Let an execution conform to
\cref{def:wick-typed-program-objects,%
def:wick-pure-numerical-interface,%
def:wick-matrix-free-backend,%
def:wick-scalable-audit-schedule,%
def:wick-residual-acceptance,%
def:wick-finite-executable-learner}.  Then:
\begin{enumerate}[label=\textup{(\roman*)}]
\item in exact arithmetic, every successful numerical-core output equals the
      finite-chart mathematical object named by its interface;
\item in finite precision, every accepted output has the a posteriori errors
      in \eqref{eq:wick-code-forward-error}, and every dependent certificate
      includes their propagated operator or optimization error; a transition
      output additionally has the enclosure
      \eqref{eq:wick-expm-action-defect-bound};
\item every returned ordinary record satisfies
      \cref{thm:wick-certified-empirical-realization}, while every rejected or
      unavailable branch returns the corresponding totalized record;
\item the outer return is the Part~I Pareto-minimal risk-qualified saturated
      frontier within each typed partition;
\item dense and matrix-free ordinary executions have the same
      exact-arithmetic target, implicit differentiation has
      \eqref{eq:wick-implicit-solve-gradient}, transition backpropagation has
      \eqref{eq:wick-expm-action-adjoint}, scalable work and memory obey
      \eqref{eq:wick-matrix-free-matvec-cost}, and certified screening leaves
      the returned frontier unchanged.
\end{enumerate}
Consequently a conforming implementation is a refinement of the mathematical
learner: it may change optimization method, batching, device, or linear-algebra
backend only through fields recorded in the transcript and numerical audit.
It may not change the structural selection rule or the meaning of optimality.
\end{theorem}

\begin{proof}
The tensor identities and exact-arithmetic solve properties are
\cref{prop:wick-tensor-compiler-faithfulness}.  The free and RL dispatches use
the same tensors, losses, frozen arguments, masks, and normalizations as
\cref{def:wick-free-caus-training-loss,def:wick-rl-training-loss}; hence their
successful exact-arithmetic outputs are the stated finite candidate objects.
The operator actions in \cref{def:wick-matrix-free-backend} are the factored
forms of the same dense operators, so both backends have the same exact target.
Implicit differentiation and its residual transfer are
\cref{prop:wick-implicit-solve-gradient}; the work and memory statement is
\cref{prop:wick-scalable-complexity}; transition-action certification is
\cref{prop:wick-expm-action-certificate}; and its scalable reverse pass is
\cref{prop:wick-expm-action-adjoint}.
For finite precision, subtract the exact equations from
\eqref{eq:wick-code-solve-residuals}--\eqref{eq:wick-code-krr-residual} and use
the three validated inverse bounds from
\eqref{eq:wick-code-coercivity-enclosures}.  This proves
\eqref{eq:wick-code-forward-error}.  The perturbation bounds then produce the
stored response errors.

The sealed specification and data-access validator preserve the candidate
index, authority, loss type, sample law, and simultaneous confidence event.
The detached audit therefore satisfies the hypotheses of
\cref{thm:wick-certified-empirical-realization}.  Its totalized failure rules
are \cref{prop:wick-no-new-premise-closure}.  Finally, steps \textup{(E5)} and
\textup{(E6)} are exactly
\cref{def:wick-detached-empirical-selector}; the saturated-profile conclusion
follows from \cref{thm:wick-simplest-certified-model,cor:wick-free-rl-minimal-models}.
The pruning conclusion is \cref{prop:wick-certified-screening}.
\end{proof}

\begin{table}[htbp]
\centering
\small
\begin{tabular}{@{}p{.28\linewidth}p{.62\linewidth}@{}}
\toprule
Required automated check & Acceptance criterion\\
\midrule
Shape, dtype, and authority property tests
& Every tensor satisfies \cref{tab:wick-code-shape-contract}; decoded bridges
  lie in the represented authority image.\\
Compiler identity tests
& Symmetry, positivity, skew symmetry, energy cancellation, and residuals in
  \eqref{eq:wick-code-compiler-audit} meet the declared tolerances; the
  outward-validated lower bounds in
  \eqref{eq:wick-code-coercivity-enclosures} are positive.\\
Transition-action tests
& On small charts, dense and Arnoldi actions agree within
  \eqref{eq:wick-expm-action-defect-bound}; adjoint gradients agree with
  directional derivatives within \eqref{eq:wick-expm-adjoint-error}; a large
  fixture respects the declared checkpoint-memory ceiling.\\
Block-elimination tests
& Direct full-block solves agree with the compressed Feshbach solve within the
  propagated residual bound.\\
Kernel tests
& Gram matrices are Hermitian positive semidefinite; range-compressed matrices
  are positive definite; empirical effective dimensions lie in
  \([0,r]\).\\
Matrix-free and augmented tests
& Every primitive action passes an adjoint dot-product test; dense, nested,
  and augmented predictions agree within the propagated residual enclosure;
  the augmented constraint and reduced normal-equation residuals both pass.\\
Resource-regression tests
& A large sparse fixture completes under the declared peak-memory ceiling;
  allocation tracing confirms that it never materializes
  \(F_C\), \(F_C^*F_C\), \(\Phi^*\Phi\), or a full response Gram matrix.\\
Free-Caus loss tests
& Primal and dual KRR predictions agree on the represented evaluation range;
  automatic gradients agree with directional derivatives of the full outer
  loss within the declared derivative-test tolerance.\\
RL semantic tests
& Terminal rows remove continuation, unsafe rows receive zero policy mass,
  empty safe rows return the authority status, targets are detached during
  critic fitting, and the full-carrier response is reused across minibatches.\\
Selection tests
& Synthetic records with known partial order return exactly all and only the
  Pareto-minimal feasible records, separately for every typed loss partition.\\
Failure-injection tests
& Nonfinite solves, excessive residuals, missing converters, unsupported
  carrier overlaps, zero safe support, and unavailable certificates return the
  specified typed status with upper bound \(+\infty\) and lower visibility
  zero where ordered.\\
Reproducibility tests
& The specification hash, seed, carrier identifiers, minibatch order, terminal
  tensors, and audit output are stored; deterministic backends reproduce the
  declared tolerance class.\\
\bottomrule
\end{tabular}
\caption{Minimum verification suite for a conforming implementation.  These
tests verify program refinement and failure polarity.  Statistical coverage is
provided by the selected Part~I certificate on its allocated event, not by a
software test.}
\label{tab:wick-code-verification-suite}
\end{table}

The definitions above determine a direct implementation order: first implement
the immutable records and compiler, then the solve backends, the
branch-specific fits, the detached audit, and finally the typed saturated
selector.  Each layer has a local mathematical oracle in
\cref{tab:wick-code-verification-suite}; no end-to-end training run is required
to diagnose a failed algebraic layer.

\paragraph{Reference implementation plan.}
The implementation is divided into seven independently testable modules.
\begin{enumerate}[label=\textup{(P\arabic*)}]
\item \emph{Schemas and persistence.}  Implement immutable
      \texttt{WickSpec}, versioned parameter and carrier handles, transcripts,
      audit records, typed failures, serialization, and specification hashing.
      This module contains no tensor algebra.
\item \emph{Operator compiler.}  Implement dense and \texttt{LinearOperator}
      actions, exact adjoint actions, authority decoders, streaming evaluation
      maps, generator actions, and exponential actions.  Property tests verify
      structure and adjointness before a solver is introduced.
\item \emph{Numerical solvers.}  Implement dense factorizations, restarted and
      block Krylov solvers, the augmented KKT backend, preconditioner and cache
      versioning, residual enclosures, and implicit adjoint differentiation.
      Dense and matrix-free outputs are compared on small random and scalar
      instances.
\item \emph{Learning branches.}  Implement dynamics fitting, free-Caus
      bilevel fitting, action-safe fitted critics, actor fitting, and imposed
      PDE evaluation using only the preceding typed interfaces.  No branch
      accesses an audit carrier through its training API.
\item \emph{Audit and bounds.}  Recompile detached terminal parameters, run
      high-precision residual checks, compute exact or certified interval
      response quantities, evaluate all applicable Part~I same-loss
      candidates, and construct one complete record per candidate.
\item \emph{Saturated selector.}  Partition records by their law, estimand,
      norm, and loss identifier; implement resource/action dominance,
      certified screening, strict inverses, and frontier serialization.
\item \emph{Parallel orchestration.}  Schedule sealed candidates as independent
      tasks, stream datasets, checkpoint transcripts, collect audit records,
      and run the selector after task completion.  The coordinator never
      modifies a candidate specification after target-bearing data have been
      read.
\end{enumerate}
Each module is released only after its rows of
\cref{tab:wick-code-verification-suite} pass.  The scalar model of
\cref{ex:scalar-wick-quotient} is the first integration fixture; dense versus
matrix-free equality is the second; a sparse large-chart instance that cannot
materialize \(F_C\) or \(K_{C,E}\) is the required scalability fixture.

\begin{definition}[Cost-aware backend dispatch]
\label{def:wick-cost-aware-dispatch}
Let \(\mathcal B_s\) be the dense, nested matrix-free, augmented, primal, and
dual backend combinations implemented for specification \(s\).  Before
target-bearing fitting, a target-free pilot records projected work
\(\widehat W_b\), peak memory \(\widehat M_b\), communication
\(\widehat C_b\), and whether backend \(b\) reached the pilot residual
tolerance.  Let \(\preceq_s^{\rm impl}\) be the declared total tie order on
the projected work--memory--communication triples, refining their Pareto
order.  The dispatcher chooses
\begin{equation}
 \widehat b
 =\operatorname*{min}_{\preceq_s^{\rm impl}}
 \left\{b\in\mathcal B_s:
 \widehat M_b\le M_s,\ b\text{ passed the pilot}\right\}
 \label{eq:wick-backend-dispatch}
\end{equation}
An empty feasible backend set returns the numerical-budget status.  Pilot cost and the selected backend
are stored in the resource record.
\end{definition}

\begin{proposition}[Backend dispatch does not alter structural selection]
\label{prop:wick-backend-dispatch-soundness}
Every backend in \(\mathcal B_s\) targets the same dense finite-chart equations.
Consequently \eqref{eq:wick-backend-dispatch} is an inner implementation choice
for the fixed candidate \(s\).  Residual acceptance makes its terminal
operator, prediction, and gradient discrepancies record coordinates.
Structural comparison remains the risk-qualified saturated profile applied to
the completed records.
\end{proposition}

\begin{proof}
Dense, nested, augmented, primal, and dual equality follows from
\cref{def:wick-matrix-free-backend} and the representer identity.  The
finite-precision differences are bounded by
\cref{def:wick-residual-acceptance}.  Hence backend dispatch changes measured
resource and numerical-error coordinates but neither the candidate family nor
the outer order.
\end{proof}

\FloatBarrier
\begin{figure}[htbp]
\centering
\resizebox{.99\linewidth}{!}{%
\begin{tikzpicture}[fw diagram,node distance=7mm and 9mm]
  \node[fw provenance,text width=2.8cm] (stage)
    {charged Part~I stage\\charts, grids, authority};
  \node[fw use,text width=2.7cm,right=of stage] (fit)
    {inner backpropagation\\dynamics, free Caus, or RL};
  \node[fw geom,text width=2.8cm,right=of fit] (solve)
    {cost-dispatched solves\\dense, matrix-free, or augmented};
  \node[fw box,text width=2.5cm,right=of solve] (detach)
    {detached terminal\\record and audit};
  \node[fw source,text width=2.7cm,right=of detach] (minimum)
    {typed working bound\\same-loss or RL conversion};
  \node[fw terminal,text width=2.8cm,right=of minimum] (profile)
    {risk-qualified saturated\\profile and Pareto inverse};
  \node[fw residual,text width=4.0cm,below=11mm of detach] (failure)
    {solve residual, nonfinite value, missing certificate,\\or empty authority fiber};
  \draw[fw arrow] (stage)--(fit);
  \draw[fw arrow] (fit)--(solve);
  \draw[fw arrow] (solve)--(detach);
  \draw[fw arrow] (detach)--(minimum);
  \draw[fw arrow] (minimum)--(profile);
  \draw[fw dashed arrow] (solve)--(failure);
  \draw[fw dashed arrow] (failure)-| node[pos=.2,below,font=\scriptsize]
    {typed value} (minimum.south);
\end{tikzpicture}%
}
\caption{Empirical realization of the Wick learner.  Differentiation is
confined to the inner candidate fit.  The detached record is audited under
the inherited confidence allocation, all applicable same-loss bounds are
combined by their valid minimum, the RL bound is separately calibrated into
value units, and complete records are ordered only within their typed
risk-qualified saturated profile and Pareto strict inverse.}
\label{fig:wick-empirical-training-pipeline}
\end{figure}
\FloatBarrier

\subsection{A solvable two-sector model}
\label{subsec:wick-solvable-model}

\begin{example}[Scalar Wick quotient]
\label{ex:scalar-wick-quotient}
Let \(\mathcal Q=\mathcal Y=\mathbb C\),
\(H_{\Abs}=\omega\ge0\), \(G=g>0\),
\(A_{\mathcal Y}=0\), and \(C=c\in\mathbb R\).  Then
\begin{equation}
 \label{eq:scalar-wick-generator}
 \mathcal L_W(c)=
 \begin{pmatrix}-i\omega&c\\-c&-g\end{pmatrix},
 \qquad
 \frac{\dd}{\dd t}(|x|^2+|y|^2)=-2g|y|^2.
\end{equation}
The quotient defect is \(|c|\), the dark space is all of \(\mathbb C\) when
\(c=0\) and zero otherwise, and
\begin{equation}
 \label{eq:scalar-wick-response}
 \Sigma_c(z)=\frac{c^2}{z+g},
 \qquad
 F_c(z)=z+i\omega+\frac{c^2}{z+g},
 \qquad
 \eta_c^W(\lambda)
 =\frac{1}{1+\left(\lambda+\dfrac{c^2}{\lambda+g}\right)^2+\omega^2}.
\end{equation}
The poles are the roots of
\begin{equation}
 \label{eq:scalar-wick-poles}
 (z+i\omega)(z+g)+c^2=0.
\end{equation}
Take the represented current coordinate to be \(J_c=c\), the mobility to be
\(a>0\), and \(J_D=-a\nabla D\).  Then
\[
 \mathscr W(c;D)=\frac{(c-J_D)^2}{2a}.
\]
Free authority selects \(c=J_D\); a restricted interval
\([c_-,c_+]\) selects the Euclidean projection of \(J_D\) onto that
interval; and an imposed PDE coefficient \(c=c_0\) retains the fixed waste
\((c_0-J_D)^2/(2a)\).  The same \(c\) occurs in the defect, self-energy,
response weight, and action record.
\end{example}

\begin{figure}[tbp]
\centering
\resizebox{.99\linewidth}{!}{%
\begin{tikzpicture}[fw diagram,node distance=7mm and 9mm]
  \node[fw provenance,text width=3.4cm] (record)
    {Part~I compensated record\\\(\Abs,\Geom,\Caus,\mathfrak A,\mathbf c\)};
  \node[fw use,text width=2.8cm,right=of record] (wick)
    {Wick on \(\Abs\) only\\\(-H_{\Abs}\mapsto-iH_{\Abs}\)};
  \node[fw geom,text width=2.8cm,right=of wick] (fixed)
    {fixed inherited fields\\\(G,C,A_{\mathcal Y},\mathfrak A\)};
  \node[fw box,text width=3.0cm,above right=8mm and 10mm of wick] (response)
    {derived ordinary response\\or typed nonordinary value};
  \node[fw box,text width=3.0cm,below right=8mm and 10mm of wick] (action)
    {inner least-action coordinate\\\(\mathscr W_{\mathfrak A}\)};
  \node[fw source,text width=2.8cm,right=15mm of response] (kernel)
    {positive kernel \(T_{W,\lambda}\)\\or \(+\infty\) candidate};
  \node[fw use,text width=3.2cm,right=15mm of action] (cert)
    {Part~I finite-carrier bounds,\\union/mixture and allocated confidence};
  \node[fw provenance,text width=3.1cm,right=14mm of kernel] (train)
    {inner inverse-response fit\\free KRR or safe RL};
  \node[fw terminal,text width=3.0cm,right=15mm of train,yshift=-9mm] (learner)
    {risk-qualified saturated profile\\and Pareto strict inverse};
  \draw[fw arrow] (record)--(wick);
  \draw[fw arrow] (wick)--(fixed);
  \draw[fw arrow] (wick)--(response);
  \draw[fw arrow] (fixed)--(response);
  \draw[fw arrow] (record)--(action);
  \draw[fw arrow] (response)--(kernel);
  \draw[fw arrow] (action)--(cert);
  \draw[fw arrow] (kernel)--(train);
  \draw[fw arrow] (train)--(learner);
  \draw[fw arrow] (cert)--(learner);
\end{tikzpicture}%
}
\caption{Wick learning as a derived operation on a Part~I record.  The Wick
map changes only the spectral \(\Abs\) evolution.  The real \(\Geom\), actual
authorized \(\Caus\) bridge, authority envelope, provenance, and charges are
fixed inherited fields.  On every ordinary Part~I realization, Feshbach
elimination produces the self-energy and positive response kernel for every
\(\lambda>0\), as a consequence of maximal dissipativity; a nonordinary
realization returns its typed value and \(+\infty\) response candidate.  The
unchanged bridge supplies the action record.
Inverse-response KRR or action-safe RL trains the applicable authority branch.
The Part~I Rademacher union/cover or PAC--Bayes mixture/coefficient rule then
uses finite evaluation-carrier traces to form the risk coordinate.  The
risk-qualified saturated usefulness profile then returns the Pareto-minimal
points of its strict resource--action inverse at the inherited visibility
threshold.}
\label{fig:wick-compensated-learning}
\end{figure}

\begin{remark}[Relation to familiar physics]
\label{rem:wick-physics-scope}
The exact correspondences proved here are spectral Wick rotation,
Feshbach--Schur elimination, Mori--Zwanzig memory, conservative coupling,
and real geometric dissipation.  They explain why the resulting structural
presentation has the same operator shapes as effective Hamiltonians,
system--bath models, and Euclidean-to-spectral descriptions.  Quantum
measurement postulates, complete positivity of density-matrix evolution,
gauge-field reconstruction, and completeness of a finite set of physical
constants require additional axioms and are outside these theorems.
\end{remark}

\FloatBarrier


\part{Structural Reachability and Regularity}
\label{part:regularity}

\section{Singular targets and cascade compilers}
\label{sec:singular-cascade}

The algebra and normalized shell identity first evaluate whether
target-aligned structure exists.  Its absence excludes contact directly.
On the surviving aligned branch, contact itself generates the canonical
infinite disjoint shell family, and
\eqref{eq:canonical-source-price-schedule} generates its sharp lower-price
series in the single physical measure.  Divergence excludes the source cell;
finite total price generates zero-price tails and activates the compact
successor relation.  Shell accountability, cascade extraction, physical
pricing, and ordered proof/model evaluation are consecutive constructive
operations.

\begin{definition}[Ordinary continuation target and source-complete shell closure]
\label{def:singular-target}
Let \(\mathsf{Cont}^{\rm tot}\) be the total continuation relation of
\cref{def:physical-presentation}.  For a finite maximal orbit, its terminal
germ is the tuple of all limits of the generated evaluation coordinates along
times increasing to \(T_*\), together with the origin tag and the cofinal
same-orbit restriction sequence; compact metrizability supplies such germ
subsequences in the source-complete graph space.  Let
\(\mathscr B_{\rm cont}^{\rm evt}\) be the set of terminal events consisting
of an ordinary same-origin maximal preterminal path, its finite endpoint
time, the value ``no ordinary regular extension'' in
\(\mathsf{Cont}^{\rm tot}\), and its complete cofinal restriction record.
The terminal analytic germ is retained as an ordinary or nonordinary total
graph coordinate of that event; the event does not require such a germ to be
an ordinary state in the declared regularity class.  The ordinary-extension, failure, and
graph-boundary status is retained as a finite discrete coproduct coordinate
of \(X_{\rm ev}\).  Define
\begin{equation}
 \label{eq:complete-continuation-target}
 \Sigma_{\rm sing}:=\Sigma_{\rm sing}^{\rm ord}
 :=\mathscr B_{\rm cont}^{\rm evt},
 \qquad
 \widehat\Sigma_{\rm sing}:=
 \overline{\mathscr B_{\rm cont}^{\rm evt}}^{\,X_{\rm ev}^{\rm fail}},
 \qquad
 X_{\rm ev}^{\rm fail}:=
 \{x:\operatorname{ContStatus}(x)\ne\mathsf{ordinary}\}.
\end{equation}
The first set is the physical singular-event target; the second is its
source-preserving shell closure.  Physical contact means either same-origin
preterminal approach to every shell of the event target or arrival of the
ordinary preterminal path at its recorded noncontinuation event.  In the second case the completed
restriction convention in \textup{(P4)} supplies the first form as well; a
genuinely discontinuous target interface outside that convention would have
a \(\Bdry\)-headed terminal jump.
Every finite maximal noncontinuation of an ordinary path is therefore contact
with the event target and generates a cofinal source-complete shell record.  This
does not promote its terminal analytic coordinate to an ordinary state.  An
arbitrary compactified boundary point is not itself physical contact: its
preterminal path and noncontinuation event must pass
\cref{thm:approximation-orbit-realization-criterion}.  The
origin coordinate prevents a terminal germ of one orbit from manufacturing
contact for another, while the discrete continuation-status coordinate
prevents an extendible germ from entering the target merely as an untagged
analytic limit.  An empty local-solution fiber remains the separate form
defect \(\bot_{\rm loc}\); it is not target avoidance.
\end{definition}

\begin{theorem}[Complete continuation target and terminal coverage]
\label{prop:terminal-coverage}
Let \(u:[0,T)\to X_{\rm reg}\) be an ordinary same-origin branch in the
declared class.  Exactly one of the following semantic alternatives holds.
\begin{enumerate}[label=\textup{(C\arabic*)}]
\item \(T=\infty\).
\item \(T<\infty\) and \(u\) has an ordinary extension in the declared
class beyond \(T\).
\item \(T<\infty\) and the total continuation relation has the same-origin
noncontinuation value at \(T\).
\end{enumerate}
If \(u\) is maximal, \textup{(C2)} is impossible; hence every finite maximal
branch lies in \textup{(C3)}.  The event in \textup{(C3)} belongs to
\(\Sigma_{\rm sing}^{\rm ord}\), and its same-origin terminal restrictions
cross a cofinal source-complete shell family for
\(\widehat\Sigma_{\rm sing}\).  Conversely, a point of the graph-complete
closure is physical contact only when its record contains one ordinary
same-origin preterminal path and that path's noncontinuation coordinate.
The terminal analytic coordinate represents an ordinary germ only after it
passes \cref{thm:approximation-orbit-realization-criterion}.

The ordinary-extension, noncontinuation, and graph-boundary tags are disjoint
coordinates.  A graph-boundary tag is an internal totalization value and is
not a fourth semantic outcome for an ordinary branch.
\end{theorem}

\begin{proof}
The cases \(T=\infty\) and \(T<\infty\) are disjoint.  At finite \(T\), the
proposition asserting the existence of an ordinary extension has a definite
truth value.  The total continuation relation records the extension in the
positive case and its noncontinuation value in the negative case.  Hence
\textup{(C2)} and \textup{(C3)} are disjoint and exhaustive.  An ordinary
extension of a finite maximal branch contradicts maximality.

By \cref{def:singular-target}, a same-origin finite noncontinuation event is
exactly an element of \(\mathscr B_{\rm cont}^{\rm evt}\).  The completed
restriction convention in \textup{(P4)} copies its origin and status to every
cofinal restriction.  The metric shell construction therefore supplies a
cofinal source-complete shell record converging to
\(\widehat\Sigma_{\rm sing}\).  Conversely, the definition of physical
contact requires those two discrete rows; closure of analytic coordinates
alone supplies neither of them.  Finally, compact metrizability provides
convergent analytic coordinates, while
\cref{thm:approximation-orbit-realization-criterion} gives the gate for
promoting such coordinates to an ordinary analytic germ.  The typed
graph-boundary value therefore remains an internal successor.
\end{proof}

\begin{definition}[Singularity-formation support]
\label{def:singularity-formation-support}
For an ordinary maximal orbit \(u\), let
\(\mathscr D_\Sigma(u)\) be a family of pre-contact shell records
\[
 R_{r,m}=\bigl(I_r,\widehat\phi_{r,m,I_r},
              (Q^c_{r,m,I_r})_{c\in\Channels}\bigr)
\]
attached to that same orbit.  The family is \emph{shell-complete} when contact
implies that it contains such a record for every sufficiently small \(r\).
Every record generated by contact then satisfies
\begin{equation}
 \label{eq:contact-unit-shell-progress}
 \Delta\widehat\phi_{r,m,I_r}=1,
 \qquad
 \begin{cases}
  \sum_{c\in\Channels}Q^c_{r,m,I_r}(1)=1,
     &\text{on a finite-value record},\\
  \overline Q^{o}_{r,m,I_r}(1)\ge1/5,
     &\text{on a variation-boundary record}.
 \end{cases}
\end{equation}
The singularity-formation support is
\begin{equation}
 \label{eq:singularity-formation-support}
 \mathfrak S_\Sigma(u)
 :=\bigcup_{R\in\mathscr D_\Sigma(u)}\Align_\Sigma(R)
 \subseteq\Channels.
\end{equation}
It answers the target-relative question: which inherited source currents
actually have positive pairing with the covectors descending toward the
singular set?  It does not ask which terms of the PDE are large in a
target-independent norm.
\end{definition}

\begin{proposition}[A singularity requires aligned structure]
\label{prop:singularity-requires-aligned-structure}
If \(u\) is generated by the compiled continuum algebra and
\(\mathscr D_\Sigma(u)\) is shell-complete,
then
\begin{equation}
 \label{eq:contact-implies-formation-support}
 u\text{ contacts }\Sigma
 \quad\Longrightarrow\quad
 \mathfrak S_\Sigma(u)\ne\varnothing.
\end{equation}
Consequently \(\mathfrak S_\Sigma(u)=\varnothing\) implies target avoidance.
This implication is independent of the size or compactness of the full
orbit.
\end{proposition}

\begin{proof}
Contact supplies normalized records satisfying
\eqref{eq:contact-unit-shell-progress}.  By
\cref{thm:deterministic-prospective-accountability,thm:target-alignment-necessity},
each has a nonempty aligned support, contained in the union
\eqref{eq:singularity-formation-support}.  The contrapositive gives target
avoidance.
\end{proof}

\begin{theorem}[Canonical disjoint shell cascade]
\label{thm:canonical-disjoint-shell-cascade}
Let a compiled ordinary maximal branch start outside \(\Sigma\) and contact
\(\Sigma\) through a cofinal pre-contact sequence at its first contact or
maximal time, and let
\((U_r,\Phi_r)\) be a target-shell presentation.  There are scales
\(r_j\downarrow0\) and pairwise disjoint, chronologically ordered,
strictly pre-contact intervals \(I_j=[a_j,b_j]\) such that
\begin{equation}
 \label{eq:canonical-shell-cascade}
 u(a_j)\notin U_{2r_j}(\Sigma),\qquad
 u(b_j)\in U_{r_j}(\Sigma),\qquad
 \Phi_{r_j}(u(b_j))-\Phi_{r_j}(u(a_j))=1.
\end{equation}
Every interval therefore has a nonempty aligned support and a source
occurrence with positive target progress.  On every ordinary or boundary
current graph, one of the five source-resolved contributions has positive
part at least \(1/5\).  No
compactness, price gap, or contact margin is used in this extraction.
\end{theorem}

\begin{proof}
Set \(b_0=0\) and \(r_0=1\).  Since every strictly pre-contact state is outside
\(\Sigma=\bigcap_{r>0}U_{2r}(\Sigma)\), after \(b_{j-1}\) choose a regular
time \(a_j>b_{j-1}\) and then a scale
\(r_j<\min\{r_{j-1}/4,1/j\}\) for which
\(u(a_j)\notin U_{2r_j}(\Sigma)\).  Contact supplies a later strictly
pre-contact time at which the orbit belongs to \(U_{r_j}(\Sigma)\); choose
one as \(b_j\).  This inductively gives \(b_{j-1}<a_j<b_j\), hence pairwise
disjoint intervals.  No continuity premise is needed: the two membership
choices give \(\Phi_{r_j}(u(a_j))=0\) and
\(\Phi_{r_j}(u(b_j))=1\) directly.  Apply
\cref{thm:deterministic-prospective-accountability} on each interval to get
the final assertions.
\end{proof}

\begin{definition}[Generated renormalization-scale row]
\label{def:renormalization-scale-interface}
Let \(z_*\) be a contact point or germ.  The normalization grammar is the
least family containing the identity and the space, time, amplitude, gauge,
and quotient transformations for which the fixed PDE's totalized covariance
graphs have ordinary values.  Each failed covariance value is a \(\Lift\)- or
\(\Abs\)-tagged graph defect.  Applying every ordinary grammar value to the
canonical shell cascade produces scales \(r_j\downarrow0\), physical
pre-contact events \(Q_j=Q(z_*,r_j)\), and records

\[
 V_j=\mathcal S_{z_*,r_j}u\in X_{\mathrm{eval}},
\]

where every \(\mathcal S_{z,r}\) has its generated \(\Abs/\Lift\) ancestry.
Pass to the least infinite constant source-support and realization-mode
subsequence in the fixed enumeration; it exists because both alphabets are
finite.  Let \(\mathfrak b\) denote the resulting same-origin active branch
signature.  All closures and prices in this row are restricted to
\(X_{\mathfrak b}\).
The row returns:
\begin{enumerate}[label=\textup{(R\arabic*)}]
\item the geometric-separation value obtained, after the canonical dyadic
  subsequence, from
  \(r_{j+1}\le\vartheta r_j\) for some \(0<\vartheta<1\) and the events
  \(Q_j\) are pairwise disjoint;
\item the totalized target-persistence inclusion
  \(\overline{\{V_j\}}\subseteq B_{\Sigma,\mathfrak b}\), where
  \(\mathfrak b\) is the active branch extracted from this same-origin
  cascade; the complementary value is a target-persistence graph defect;
\item the total-graph value of the physical pullback relation
  \begin{equation}
   \label{eq:physical-scale-pullback}
   \mu_u(Q_j)\ge
   c_{\mathrm{sc}}^*r_j^{\kappa}
   \overline d_{\mathrm{sc}}(V_j),
  \end{equation}
\end{enumerate}
Here \(c_{\mathrm{sc}}^*\) is the sharp transfer infimum, \(\kappa\) is the
exponent returned by the covariance graph, and
\(\overline d_{\mathrm{sc}}\) is the canonical lower-semicontinuous
relaxation.  If no scalar exponent represents the pullback, the row returns
the full scale-dependent weight sequence or its graph defect.  Thus scale
covariance is evaluated rather than postulated.
\end{definition}

\begin{theorem}[Scale-cascade price compiler]
\label{thm:scale-cascade-price-compiler}
Evaluate the totalized scale interface on its generated active branch
\(\mathfrak b\).  Its normalized price value is
\begin{equation}
 \label{eq:scale-normalized-gap}
 \delta_{\mathrm{sc}}
 :=\inf_{V\in B_{\Sigma,\mathfrak b}}\overline d_{\mathrm{sc}}(V)
 \in[0,\infty].
\end{equation}
For the (j)-th generated scale stratum, define its sharp physical lower
price exactly as in \eqref{eq:canonical-source-price-schedule} and denote it
by \(\underline p_{\mathrm{sc},j}\).  The exact scale-price value is
\begin{equation}
 \label{eq:divergent-scale-weight}
 \mathscr S_\Sigma
 :=\sum_{j=1}^{\infty}\underline p_{\mathrm{sc},j}
 \in[0,\infty].
\end{equation}
On the ordinary scalar-pullback branch, whenever the displayed factors are
finite, \eqref{eq:physical-scale-pullback} gives
\begin{equation}
 \label{eq:scale-price-power-lower-bound}
 \underline p_{\mathrm{sc},j}
 \ge c_{\mathrm{sc}}^*r_j^\kappa\delta_{\mathrm{sc}}.
\end{equation}
The branch \(\mathscr S_\Sigma=+\infty\) gives
\begin{equation}
 \label{eq:scale-price-series}
 \Budget
 \ge\mu_u\!\left(\bigsqcup_{j=1}^{\infty}Q_j\right)
 =\sum_{j=1}^{\infty}\mu_u(Q_j)
 \ge\sum_{j=1}^{\infty}\underline p_{\mathrm{sc},j}
 =\mathscr S_\Sigma
 =\infty,
\end{equation}
contradicting a finite physical budget.  The stronger condition
\(a_*:=\inf_jr_j^\kappa>0\) gives the uniform event price
\(\mu_u(Q_j)\ge c_{\mathrm{sc}}^*a_*\delta_{\mathrm{sc}}>0\) on the
positive transfer and price branches.
On the positive transfer and normalized-price branches, after restricting to
\(r_j\le1\), divergence is automatic when \(\kappa\le0\), and the
scale-critical case \(\kappa=0\) transfers the normalized gap without loss.
When \(\kappa>0\), the exact criterion is divergence of the generated
lower-price series; a normalized gap alone supplies no such conclusion.  The
complementary value \(\mathscr S_\Sigma<\infty\)
produces vanishing-price windows in its inherited cell and is routed by
\cref{thm:finite-price-saturation-route}.  In particular, a positive
dissipative pullback exponent cannot be converted into a scale-independent price: the scale
motion is retained in the signed carrier cohomology and then tested by the
successor-saturation algebra.
\end{theorem}

\begin{proof}
The definition of each sharp lower price gives
\(\mu_u(Q_j)\ge\underline p_{\mathrm{sc},j}\).  Sum over the pairwise
disjoint events to obtain \eqref{eq:scale-price-series}; countable additivity
and the finite physical budget give the contradiction.  On the scalar
pullback branch, insert \eqref{eq:scale-normalized-gap} into
\eqref{eq:physical-scale-pullback} to obtain
\eqref{eq:scale-price-power-lower-bound}.  If \(0<r_j\le1\), then
\(r_j^\kappa\ge1\) for \(\kappa\le0\).  For \(\kappa>0\), the individual
weights tend to zero, so their summability must be decided from the actual
scale sequence.
\end{proof}

\begin{proposition}[Canonical blow-up invariant image]
\label{prop:blowup-structural-data}
Every contact cascade and every ordinary value of the generated normalization
grammar has the canonical invariant image
\begin{equation}
 \label{eq:blowup-invariant-tuple}
 \left(
  \vartheta,
  \overline{\{V_j\}},
  \mathfrak H_{\{V_j\}}^\Sigma,
  \kappa,
  \delta_{\mathrm{sc}}
 \right).
\end{equation}
Here \(\vartheta\le1/4\) is produced by the canonical shell subsequence, the
closure is taken in the generated evaluation compactification,
\(\mathfrak H_{\{V_j\}}^\Sigma\) records whether the cascade remains
target-visible, \(\kappa\) or its full weight-sequence replacement is the
physical pullback value, and \(\delta_{\mathrm{sc}}\) is the canonical price
infimum.  A failed normalization or pullback occupies the corresponding
\(\Lift/\Abs\) graph-defect cell.  Thus every blow-up outcome has an image;
none of the five coordinates is additional problem data.
\end{proposition}

\begin{definition}[Normalized cascade refinement]
\label{def:contact-cascade-refinement}
For a contact orbit, combine the canonical intervals from
\cref{thm:canonical-disjoint-shell-cascade} with every ordinary value of the
generated normalization grammar, then pass to the least infinite constant
source-support and realization-mode subsequence.  Let \(\mathfrak b\) be its
same-origin active branch signature.  The resulting normalized cascade
refinement consists of
\begin{enumerate}[label=\textup{(C\arabic*)}]
\item pairwise disjoint measurable pre-contact events
  \(Q_1,Q_2,\ldots\) in the domain of the same physical expenditure measure
  \(\mu_u\);
\item generated normalizations \(\mathcal N_{Q_j}\), carrying their
  \(\Abs/\Lift\) ancestry, such that
  \(w_j:=\mathcal N_{Q_j}u\in B_{\Sigma,\mathfrak b}\), for the persistent
  active branch \(\mathfrak b\) selected from this same-origin cascade;
\item for every \(j\), a normalized shell-crossing record \(R_j\) carried by
  \(w_j\) for which
  \begin{equation}
   \label{eq:cascade-unit-shell-accountability}
   \Delta\widehat\phi(R_j)=1,
   \qquad
   \operatorname{Prog}_\Sigma(R_j)\ge\frac15.
  \end{equation}
  On a finite-value record the latter quantity is the five-term sum, equals
  one, and has a component at least \(1/5\); on a variation-boundary record it
  is the retained owner progress from
  \eqref{eq:boundary-owner-progress-convention}.
\end{enumerate}
By \cref{thm:no-target-visible-residual} and the applicable prospective
shell-accountability theorem, every occurrence in clause \textup{(C3)} has a
nonempty aligned source support
\(A_j:=\Align_\Sigma(R_j)\subseteq\Channels\).  Every nonordinary
normalization value is already a typed graph-defect record.  The refinement
asserts no positive price; it generates the family on which the price row is
evaluated.
\end{definition}

\begin{proposition}[Persistent channel extraction]
\label{prop:persistent-channel-extraction}
For every cascade generated by \cref{def:contact-cascade-refinement}, some
nonempty support \(A\subseteq\Channels\) occurs for infinitely many indices.
In particular, at least one of the five elementary source labels belongs to
the ancestry of infinitely many target-approaching windows.
\end{proposition}

\begin{proof}
There are only \(2^5-1\) nonempty supports.  The infinite pigeonhole
principle gives an infinite constant-support subsequence.  Any member of that
support has the second asserted property.
\end{proof}

\begin{definition}[Target-aligned source ledger]
\label{def:target-aligned-source-ledger}
Fix an ordinary maximal branch \(u\) and a normalized cascade family
\((Q_j,R_j)_{j\ge1}\).  For every nonempty \(A\subseteq\Channels\), let
\(\mathcal F_A(u)\) be the family of source-current occurrences extracted
from those \(R_j\) for which \(\Align_\Sigma(R_j)=A\).  The source ledger is
generated as follows.  For \(j\in I_A:=\{j:\Align_\Sigma(R_j)=A\}\), let
the compiler first evaluates the total graph of the target-aligned charging
relation \eqref{eq:target-aligned-charge-rule}.  On its ordinary admitted
value, let \(\mathcal Q_{A,j}^{\rm pay}\) be the family of all generated
ordinary events under the same physical ceiling and in the same propagated
source/mode, shell/scale, and carrier-factorization stratum as \(Q_j\),
normalized to the same unit shell increment and source-owner floor.  Equality of stratum means
equality of every transcript entry preceding and including the charging row,
so the equivalence relation is canonical.  Define the sharp admissible
physical lower price
\begin{equation}
 \label{eq:canonical-source-price-schedule}
 \underline p_{A,j}
 :=\inf\{\mu_v(Q): (v,Q)\in\mathcal Q_{A,j}^{\rm pay}\}\in[0,\infty],
\end{equation}
with \(+\infty\) for an empty stratum.  The stratum contains \((u,Q_j)\), so
\(\mu_u(Q_j)\ge\underline p_{A,j}\).  Each source row has exactly one of the
following exhaustive entries:
\begin{enumerate}[label=\textup{(L\arabic*)}]
\item a \emph{null entry},
  \begin{equation}
   \label{eq:source-ledger-null-entry}
   \mathfrak H_{\mathcal F_A(u)}^\Sigma=0;
  \end{equation}
\item a \emph{divergent-price entry}, when the target quotient is nonzero and
  \begin{equation}
   \label{eq:source-ledger-divergent-entry}
   |I_A|=\infty\Longrightarrow
   \sum_{j\in I_A}\underline p_{A,j}=\infty;
  \end{equation}
\item a \emph{finite-price successor entry}, when the target quotient is
nonzero and either the charging relation has a nonordinary or failed graph
value, or its ordinary admitted value satisfies
  \begin{equation}
   \label{eq:source-ledger-formation-entry}
   \sum_{j\in I_A}\underline p_{A,j}<\infty.
  \end{equation}
\end{enumerate}
The target quotient, the totalized charging relation, and the dichotomy for a
nonnegative series make the list exhaustive.  A charging defect retains the
same source, shell covector, and event origin and enters the successor algebra;
it is never assigned price \(+\infty\).  Quotient, symmetry, conservation, carrier, resistance,
compactness--rigidity, equality, scale, and entropy calculations are
evaluators of these exact entries.  A uniform positive lower price is a
special divergent-price value.  Since \(|\Channels|=5\), the ledger has at
most \(2^5-1=31\) nonempty rows even when the cascade contains infinitely
many scales.
\end{definition}

\begin{theorem}[Physical expenditure exclusion]
\label{thm:physical-expenditure-exclusion}
Let \(\mu\) be a nonnegative countably additive measure with
\(\mu(\Omega)\le \Budget<\infty\).  Let \(Q_j\subseteq\Omega\) be pairwise
disjoint measurable sets and let \(\delta_j\ge0\).  If
\begin{equation}
 \label{eq:divergent-physical-price-schedule}
 \mu(Q_j)\ge\delta_j
 \quad\text{for every }j,
 \qquad
 \sum_{j=1}^{\infty}\delta_j=\infty,
\end{equation}
then such a family cannot exist.  The uniform-price criterion
\(\delta_j\ge\delta>0\) is a special case.
\end{theorem}

\begin{proof}
For every \(N\), finite additivity and monotonicity give
\[
 \Budget\ge
 \mu\!\left(\bigsqcup_{j=1}^{N}Q_j\right)
 =\sum_{j=1}^{N}\mu(Q_j)
 \ge \sum_{j=1}^{N}\delta_j.
\]
The partial sums on the right are unbounded by
\eqref{eq:divergent-physical-price-schedule}, contradicting the fixed upper
bound \(\Budget\).
\end{proof}

\begin{corollary}[Shell progress implies expenditure exclusion]
\label{cor:shell-price-exclusion}
On one ordinary same-origin active contact branch \(\mathfrak b\), let
\(R_j\) be the normalized shell records on pairwise disjoint pre-contact
events \(Q_j\).  On the ordinary admitted value of the target-local charging
row,
\begin{equation}
 \label{eq:unit-progress-physical-price}
 \Delta\widehat\phi(R_j)=1,
 \qquad \operatorname{Prog}_\Sigma(R_j)\ge1/5
 \quad\Longrightarrow\quad
 \mu_u(Q_j)\ge\delta_j\ge0.
\end{equation}
If \(\sum_j\delta_j=\infty\), then \(u\) cannot contact \(\Sigma\).
\end{corollary}

\begin{proof}
Contact supplies the unit shell increments and source-owner floor by
\cref{thm:deterministic-prospective-accountability}.  The price implication
is a structural fact on \(X_{\mathfrak b}\) by
\cref{thm:target-local-fact-factorization,thm:local-target-aligned-charging-rule}.
It gives \eqref{eq:divergent-physical-price-schedule}, and
\cref{thm:physical-expenditure-exclusion} gives the contradiction.
\end{proof}

\begin{theorem}[Finite structural-ledger exclusion]
\label{thm:finite-structural-ledger-exclusion}
Let \(u\) have the canonical contact cascade of
\cref{thm:canonical-disjoint-shell-cascade}.  On the fast-track value where
every persistent target-aligned support is either zero in the target quotient
or has a divergent physical-price series, \(u\)
does not contact \(\Sigma\).
\end{theorem}

\begin{proof}
On the contact branch, \cref{prop:persistent-channel-extraction} gives a
nonempty \(A\subseteq\Channels\) for which
\(I_A:=\{j:\Align_\Sigma(R_j)=A\}\) is infinite.  If \(A\) has a null entry,
then every target projection of the corresponding source-current occurrences
vanishes.  This contradicts \(A=\Align_\Sigma(R_j)\ne\varnothing\).  If
\(A\) has a divergent-price entry, the pairwise disjoint events
\((Q_j)_{j\in I_A}\) carry the divergent lower-price schedule
  \((\underline p_{A,j})_{j\in I_A}\), contradicting
\cref{thm:physical-expenditure-exclusion}.  The ledger alternatives are
exhaustive, so contact is impossible.
\end{proof}

\begin{corollary}[Estimate-to-invariant reduction]
\label{cor:estimate-to-invariant-reduction}
In the setting of \cref{thm:finite-structural-ledger-exclusion}, the canonical
transcript evaluates each persistent source support as null,
divergent-price, or finite-price successor.  By
\cref{thm:estimate-erasure}, no source-transverse estimate of the full orbit remains in
the verdict after these row values are known.  In particular, arbitrarily
large target-null motion is absent from the reduced calculation.  The last
value is not a verdict: it activates
\cref{thm:finite-price-saturation-route}.
\end{corollary}

\begin{corollary}[Entropy exclusion is physical expenditure exclusion]
\label{cor:entropy-expenditure-exclusion}
On the entropy-cocycle, finite-error, carrier-domination, and
positive-production rows of one same-origin active contact branch
\(\mathfrak b\),
\(\Production_j(I_j)\ge\delta>0\) on every selected window and contact is
impossible.
\end{corollary}

\begin{proof}
By \cref{cor:entropy-physical-budget}, the glued production is a finite
physical measure on this branch.  Its use as the price of the selected shell
events is admitted by
\cref{thm:local-target-aligned-charging-rule}.  Apply
\cref{thm:physical-expenditure-exclusion} to its restrictions to the disjoint
same-origin events \(Q_j\).
\end{proof}

\begin{figure}[tbp]
\centering
\resizebox{.98\linewidth}{!}{%
\begin{tikzpicture}[fw diagram,node distance=7mm and 9mm]
  \node[fw source,text width=2.5cm] (contact)
    {putative target contact};
  \node[fw use,text width=2.6cm,right=of contact] (descent)
    {normalized shell\\unit progress};
  \node[fw provenance,text width=2.7cm,right=of descent] (aligned)
    {nonempty singularity-\\formation support};
  \node[fw geom,text width=2.8cm,right=of aligned] (current)
    {target-cyclic spine\\\(\mathfrak J_{\Sigma,\mathfrak b}\)};
  \node[fw use,text width=2.8cm,right=of current] (queue)
    {zero, return, graph\\internal queue};
  \node[fw terminal,text width=2.8cm,above right=8mm and 10mm of queue] (divergent)
    {divergent physical\\charge};
  \node[fw residual,text width=2.7cm,right=11mm of queue] (null)
    {target-null\\retained path};
  \node[fw terminal,text width=2.9cm,below right=8mm and 10mm of queue] (witness)
    {contact-derived joint\\source occurrence};
  \draw[fw arrow] (contact)--(descent);
  \draw[fw arrow] (descent)--(aligned);
  \draw[fw arrow] (aligned)--(current);
  \draw[fw arrow] (current)--(queue);
  \draw[fw arrow] (queue)--(divergent);
  \draw[fw arrow] (queue)--(null);
  \draw[fw arrow] (queue)--(witness);
\end{tikzpicture}%
}
\caption{Compressed singularity-formation accounting.  Contact exposes one
target-aligned source cell.  Projection-first saturation and signed gluing
produce the target-cyclic feeding--storage spine.  Zero-kernel restrictions, signed-dyadic return maps,
and graph-boundary successors are internal queue transitions.  On the
contact-seeded branch the queue has three outputs: divergent physical charge,
target-nullity, or a contact-derived finite joint source occurrence.  Its
ordinary same-origin realization is evaluated in the corresponding reachable
contact cell.}
\label{fig:singular-cascade-expenditure}
\end{figure}

\section{The continuous structural regularity theorem}
\label{sec:structural-regularity-theorem}

The five-channel conclusion is derived before any regularity estimate.  The
only equation-dependent step in target exclusion is the physical evaluation
of the target-aligned class exposed by contact.

\begin{definition}[Totalization of a continuum interface]
\label{def:totalized-continuum-interface}
Let \(F:D(F)\subset X_F\to Y_F\) be any partial interface used by the
continuum argument: a constitutive realization, a gradient potential map, a
target-shell derivative, a physical-price comparison, or a continuation map.
Adjoin a failure value \(\bot_F\) and define
\[
 F^{\rm tot}(x)=
 \begin{cases}
 F(x),&x\in D(F),\\
 \bot_F,&x\notin D(F).
 \end{cases}
\]
Embed the graph of this total map in the structural test compactification and
set
\begin{equation}
 \label{eq:interface-graph-boundary}
 \partial F
 :=\overline{\operatorname{graph}F^{\rm tot}}
   \setminus\operatorname{graph}F.
\end{equation}
Thus \(\partial F\) contains both limiting graph defects and the explicit
failure records \((x,\bot_F)\).  The \emph{totalized interface record} of an
input is its graph value when the input lies in \(D(F)\), and otherwise its
class in the closed defect cell generated by \(\partial F\).  The defect cell
retains the ancestry of
the approximating graph records; failure at a chart, weak-limit, or terminal
interface appends \(\Lift\) or \(\Bdry\) according to
\cref{def:source-alphabet}.  Its target-null part is quotiented by
\(\Null_\Sigma\); its nonzero target quotient is a target-visible defect
class in the same five-channel algebra.
\end{definition}

\begin{definition}[Total continuum interface algebra]
\label{def:continuum-obstruction-algebra}
The target-relative continuum interface algebra is the finite coproduct
\begin{equation}
 \label{eq:continuum-obstruction-algebra}
 \mathfrak H_{\rm cont}^\Sigma
 :=\mathfrak H_{\rm form}^\Sigma
 \sqcup\mathfrak H_{\rm grad}^\Sigma
 \sqcup\mathfrak H_{\rm chain}^\Sigma
 \sqcup\mathfrak H_{\rm price}^\Sigma
 \sqcup\mathfrak H_{\rm cov}^\Sigma .
\end{equation}
Here the five cells are respectively the target-active quotients of:
\begin{enumerate}[label=\textup{(D\arabic*)}]
\item the graph boundary of the local-solution map and of the native
constitutive or sector-form realization;
\item \(\overline{\Ran G}\setminus\Ran G\), together with the target-active
part of \(\ker G^*\), after the exact target-null quotient;
\item the source-resolved derived Stieltjes shell currents of
\eqref{eq:derived-cylinder-shell-current}, including their ordinary,
Lebesgue, jump, Cantor, oscillation, concentration, exterior, and
variation-boundary graph coordinates;
\item the total carrier graph, including failure to produce the finite physical
measure dictated by the equation, and the zero-price limit of every
finite-expenditure shell cascade;
\item the generated chart, gauge, scale, and quotient-covariance graphs,
including every commutator defect created when an operation fails to descend.
\end{enumerate}
Every cell is source-tagged before its target-null quotient is taken.  A
missing property is therefore an input to the activity lift and successor
operator, not a terminal verdict.  Regular continuation is not a sixth
interface cell: by \cref{def:singular-target,prop:terminal-coverage}, a finite
maximal branch without an ordinary continuation generates a cofinal source-complete
shell record and is processed by the same shell and successor machinery.  Its
compactified germ is physical target contact only after the ordinary
realization gate.
\end{definition}

\begin{theorem}[Defect totalization and five-channel closure]
\label{thm:defect-totalization-five-channel}
For every fixed continuum evolution, target, and physical carrier, each of
the five interfaces in \cref{def:continuum-obstruction-algebra} has an
ordinary total-graph value or a graph-boundary value.  Every target-visible
value has ancestry in \(\Geom,\Caus,\Abs,\Lift,\Bdry\); every
target-invisible value lies in \(J_{\rm res}\).  Every target-visible value,
ordinary or defective, is converted into an activity measure and a
witness-indexed successor.  Hence lack of closed range, a classical chain
rule, positive price, or covariance is not a stopping point: it is
a concrete successor mechanism in the same closed algebra.  Finite
noncontinuation is target contact, not an unprocessed remainder.
\end{theorem}

\begin{proof}
For each partial map, the total value is either an ordinary graph value or the
failure value \(\bot_F\).  Closing the total graph adds its limiting defects,
so \eqref{eq:interface-graph-boundary} exhausts both kinds of failure.
Sourcewise closure and
\cref{thm:source-preserving-completion} retain the ancestry of every
approximating record.  The gradient cell is exhausted by
\cref{thm:target-relative-gradient-exhaustiveness}.  For the shell interface,
the finite approximation grammar gives the exact polynomial Stieltjes
identity \eqref{eq:derived-cylinder-shell-current}.  Sourcewise graph closure
passes the affine identity to every finite Radon value.  At a
variation-boundary value, the joint shell graph instead preserves the
constant canonical owner and its closed target floor
\eqref{eq:closed-shell-owner-floor}; no undefined extended signed sum is
used.  Sourcewise closure preserves the ancestry of every current and there
is no untyped chain-rule remainder.  A finite-price cascade enters
\cref{def:ordered-continuum-successors}; a carrier-graph failure is lifted by
the same total-graph coordinate.  Covariance and descent failures retain the
tags of the operation and commutator that generated them.  Finally quotient
each cell by its intersection with \(\Null_\Sigma\).  The kernel is residual.
Every nonzero quotient value enters \(\mathcal R_\eta\), and
\cref{def:total-successor-saturation-algebra} iterates it until it is
eliminated or lies on an infinite zero-price path.  Terminal coverage and the
realization gate handle finite noncontinuation without promoting a source-complete
germ.  This proves all assertions.
\end{proof}

\begin{theorem}[No favorable analytic premise in constructive closure]
\label{thm:no-favorable-analytic-premise}
Fix only the canonical weak-PDE signature, its complete continuation target,
and its equation-generated physical carrier.  The source compiler, sourcewise
closure, target quotient, total interface graphs, activity lift, successor
relation, and finite-depth viability sets are defined without requiring any of
the following properties:
\begin{equation}
 \label{eq:no-favorable-premise-list}
 \begin{gathered}
 \text{a gradient/sector realization, local continuation, compactness,}\\
 \text{continuation of the full branch, or closed gradient range,}\\
 \text{a classical chain rule, smooth shells, coercivity, a positive price
 gap, rigidity, polarity,}\\
 \text{finite resistance, summable tails, orbit realization, or
 continuation.}
 \end{gathered}
\end{equation}
Each listed property is tested by a total graph generated from the equation.
Both its ordinary and boundary values are supplied to the same activity lift,
successor relation, and finite-depth viability recursion.  A branch therefore
cannot terminate at the sentence that a property fails: it is removed by the
target-null quotient, charged by the physical carrier, or retained as an
explicit zero-price recurrent mechanism.  Consequently none of the listed
properties is a premise of the totalization, contact-accountability, or
constructive regularity theorems.
\end{theorem}

\begin{proof}
The compiler is structural induction on the fixed equation syntax.
Sourcewise closure is defined separately in each tagged cell.  Quotienting
by the common target kernel is exact by definition.  Totalization replaces
every partial interface by the disjoint union of its graph values, explicit
failure values, and graph-boundary values.  The realization modes partition
the resulting ordinary and defect records.  The positive target-current part
of each record defines its activity measure, whose normalized lift defines
\(\mathcal R_\eta\).  Iterating
\(Z_{n+1}=\mathcal V_\eta(Z_n)\) removes every record without arbitrarily long
zero-price descendants.  Compactness turns survival at every finite depth
into an infinite path and a stationary path law.  Thus both sides of every
analytic relation undergo an explicit structural operation, and no favorable
side of a dichotomy is selected.
\end{proof}

\begin{definition}[Physical reachability profile]
\label{def:physical-reachability-profile}
For a finite number \(\Budget\) arising from the carrier identity, let
\(\Orbit^{\rm ord}_{\le\Budget}\) be the ordinary maximal solution branches
whose derived expenditure satisfies
\(\mu_u(\mathfrak Q_u)\le\Budget\).  When this set is nonempty, define
\begin{equation}
 \label{eq:physical-reachability-profile}
 \Reach^{\Phys}_\Sigma(\Budget)
 :=\sup_{u\in\Orbit^{\rm ord}_{\le\Budget}}
   \mathbf 1_{\{u\text{ contacts }\Sigma\}}.
\end{equation}
There is no empty-supremum convention.  If the ordinary solution stratum is
empty, the total local-solution or carrier graph returns its explicit defect
record.  Thus \(\Reach^{\Phys}_\Sigma(\Budget)=0\) means that ordinary
maximal branches exist and none contacts the target.  The ceiling is not a
second datum: it is only a stratum of the one physical expenditure generated
by the equation.
\end{definition}

\begin{definition}[Constructive zero-density mechanism cells]
\label{def:realized-singularity-formation-cells}
Fix a derived finite physical stratum \(\Budget\).  For nonempty
\(A\subseteq\Channels\) and \(m\in\Modes\), let
\(Z_{\infty,A,m}^{\Sigma}(\Budget)\) be the part of
\(Z_\infty^\Sigma\) whose inherited support is exactly \(A\) and whose
primary mode is exactly \(m\).  Full ancestry and all nonprimary flags remain
coordinates of the record.  Define
\begin{align}
 \mathfrak F_{A,m}^{\Sigma}(\Budget)
 &:=
 \overline{\operatorname{span}}
 \left\{
  \bigl(\Pi_\Sigma(\pi_0)_\#\boldsymbol\nu,
        \operatorname{Aug}(\boldsymbol\nu)\bigr):
  \boldsymbol\nu\text{ stationary on }
  \mathsf{Path}(Z_{\infty,A,m}^{\Sigma}(\Budget))
 \right\},\label{eq:realized-formation-cell}\\
 \mathfrak F_{\Sigma}^{\Phys}(\Budget)
 &:=
 \bigoplus_{\substack{\varnothing\ne A\subseteq\Channels\\m\in\Modes}}
 \mathfrak F_{A,m}^{\Sigma}(\Budget).\label{eq:total-formation-algebra}
\end{align}
These cells are not defined from contact orbits.  They are the finite
source/mode decomposition of the prospectively generated zero-density fixed
point.  The direct sum is provenance-tagged: analytically equal
representatives with different ancestry remain distinct.  Every generator
carries its stationary path law, full weak equation or recurrent graph
boundary, target pairing, and diagonal cascade-origin witnesses.  It also
carries its unnormalized occurrence, expenditure, and source-current
prefixes.  Equality of normalized support laws therefore never identifies
two different cumulative histories.
\end{definition}

\begin{definition}[Terminal boundary formation cell]
\label{def:terminal-boundary-formation-cell}
Let \(\mathfrak F_{\Sigma}^{\rm term}(\Budget)\) be the closed
provenance-tagged span of the target projections of all completed
continuation-failure records in the derived physical stratum whose shell
entry occurs only at the terminal endpoint.  Each generator retains the
ordinary same-origin weak-equation path, its terminal continuation-status
value, the \(\Bdry\)-headed jump current, the joint shell owner, and every
auxiliary total-graph descendant.  This cell is generated prospectively from
the total continuation graph; it is not defined by choosing an observed
singular orbit.  It is disjoint from the recurrent zero-density cells because
it carries the terminal-entry flag rather than a cofinal pre-contact cascade.
\end{definition}

\begin{theorem}[Constructive source decomposition of zero-density dynamics]
\label{thm:constructive-zero-price-source-decomposition}
For every derived finite physical stratum \(\Budget\),
\begin{equation}
 \label{eq:reachability-formation-equivalence}
 \boxed{
 \mathfrak H_{\rm succ}^{\Sigma}(\Budget)
 =\mathfrak F_{\Sigma}^{\Phys}(\Budget).}
\end{equation}
The right side has at most \(31\times5\) provenance-tagged cells.  Every
finite-expenditure approach-contact orbit produces a nonzero generator in
one of them.  Every terminal-entry contact instead produces a nonzero
generator of \(\mathfrak F_{\Sigma}^{\rm term}(\Budget)\).
Consequently, whenever the ordinary stratum is nonempty,
\begin{equation}
 \label{eq:zero-mechanism-implies-avoidance}
 \mathfrak F_{\Sigma}^{\Phys}(\Budget)=0
 \quad\text{and}\quad
 \mathfrak F_{\Sigma}^{\rm term}(\Budget)=0
 \quad\Longrightarrow\quad
 \Reach_\Sigma^{\Phys}(\Budget)=0.
\end{equation}
A nonzero cell is not a label awaiting a further estimate: its generator is an
augmented stationary zero-density successor law carrying the weak equation, recurrent
graph defects, target pairing, source ancestry, and cascade-origin witness
sequence, together with the exact cumulative five-source identity; an
ordinary contact origin makes these actual physical events.
\end{theorem}

\begin{proof}
Along a successor path the inherited support can only increase.  Because the
source alphabet is finite, it is eventually constant; stationarity makes it
constant almost surely.  The primary mode is likewise monotone in a finite
ordered set and is stationary only when constant almost surely.  Decomposing
every stationary path law along these invariant finite
parts proves \eqref{eq:reachability-formation-equivalence}.

If an ordinary branch has approach contact with \(\Sigma\), the total shell
graph gives the disjoint unit-increment cascade without a chain-rule premise.
\Cref{thm:finite-price-saturation-route} converts that same cascade into a
stationary zero-density law with nonzero target projection and diagonal
origin-event witnesses.  The finite tag argument places a nonzero component
in one \((A,m)\)-cell.  If contact first occurs at the terminal endpoint,
\cref{thm:deterministic-prospective-accountability} gives its positive
\(\Bdry\)-headed owner, which is a generator of
\(\mathfrak F_{\Sigma}^{\rm term}\).  The contrapositive proves
\eqref{eq:zero-mechanism-implies-avoidance}; the stated data are supplied by
\cref{thm:saturation-equation-compiler,lem:successor-no-double-counting}.
\end{proof}

\begin{theorem}[Finite-depth structural evaluator]
\label{thm:finite-depth-structural-evaluator}
For every source/mode cell, start with its generated target-active zero-price
records and form
\begin{equation}
 \label{eq:finite-depth-cell-recursion}
 Z_{0,A,m}\supseteq Z_{1,A,m}\supseteq\cdots,
 \qquad
 Z_{n+1,A,m}=\mathcal V_\eta(Z_{n,A,m}).
\end{equation}
Exactly one of the following constructive outcomes occurs:
\begin{enumerate}[label=\textup{(E\arabic*)}]
\item at some finite cylinder-moment level the dual proof cone contains
\(-1\); the finite ordered identity produces the pruning transcript of
\cref{thm:canonical-finite-pruning-transcript} and closes the cell;
\item every finite moment level is feasible; diagonal reconstruction produces
an infinite successor path law and a nonzero generator of
\(\mathfrak F_{A,m}^{\Sigma}(\Budget)\).
\end{enumerate}
Applying the generated zero-cost equations to outcome \textup{(E2)} removes
each target-null component and recursively expands each ordinary or
graph-boundary component.  Thus every nonzero row value undergoes a further
structural operation.
\end{theorem}

\begin{proof}
Apply \cref{thm:proof-producing-hierarchy-completeness} to the path algebra
restricted to cell \((A,m)\).  Its dual branch is \textup{(E1)} and its primal
projective-limit branch is \textup{(E2)}.  The ordered ideal contains the
target floor, zero price, same-origin successors, and source/mode coordinates,
so reconstruction retains all of them.  The final reduction is
\cref{thm:saturation-equation-compiler} applied to every support component.
\end{proof}

\begin{theorem}[End-to-end constructive closure]
\label{thm:end-to-end-fast-track-completeness}
For a derived finite physical stratum \(\Budget\), apply source compilation,
the target quotient, total interface graphs, shell activity lifting,
witness-indexed successors, finite-depth viability pruning, and the generated
zero-cost equations, in that order.  The output is
\begin{equation}
 \label{eq:end-to-end-fast-track-output}
 \boxed{
 \mathfrak H_{\rm out}^{\Sigma}(\Budget)
 =\mathfrak F_{\Sigma}^{\rm term}(\Budget)
  \oplus\mathfrak H_{\rm succ}^{\Sigma}(\Budget)
 =\mathfrak F_{\Sigma}^{\rm term}(\Budget)
  \oplus\bigoplus_{A,m}\mathfrak F_{A,m}^{\Sigma}(\Budget).}
\end{equation}
On every ordinary local finite-carrier branch,
\begin{equation}
 \label{eq:end-to-end-regularity-implication}
 \mathfrak F_{\Sigma}^{\rm term}(\Budget)=0,
 \qquad
 \bigoplus_{A,m}\mathfrak F_{A,m}^{\Sigma}(\Budget)=0
 \quad\Longrightarrow\quad
 \Reach_\Sigma^{\Phys}(\Budget)=0.
\end{equation}
The compiler never consults an actual-contact image.  A contact orbit itself
constructs a nonzero cell.  Conversely, a cell that survives the evaluator
comes with an infinite zero-price path, stationary law, full generated
saturation equations, source support, target pairing, and diagonal origin
events.  A finite exclusion carries its rational ordered derivation.  A
surviving cell carries the positive-functional model of all enumerated
relations; coherent approximation-origin paths are retained separately from
source-complete records.
\end{theorem}

\begin{proof}
Totalization inserts every target-visible local-solution or carrier boundary
witness into the same successor graph.  On an ordinary finite-carrier branch,
source and mode tags partition the zero-price path laws and
\cref{thm:constructive-zero-price-source-decomposition} identifies their
direct sum with \(\mathfrak H_{\rm succ}^{\Sigma}\), proving
\eqref{eq:end-to-end-fast-track-output} on the approach component; the
terminal continuation graph supplies its disjoint \(\Bdry\) summand.  If
approach contact occurs, shell accountability, disjoint cascade extraction,
and \cref{thm:finite-price-saturation-route} construct a nonzero stationary
path law in one cell.  If terminal-entry contact occurs, its closed shell
owner gives a nonzero terminal generator.  Their joint contrapositive proves
\eqref{eq:end-to-end-regularity-implication}.  The finite-depth evaluator and
saturation-equation compiler supply the listed data for every surviving
generator, without an actual-contact filter.
\end{proof}

\begin{theorem}[Unconditional continuum contact mechanism]
\label{thm:unconditional-continuum-contact-alternative}
Fix the canonical physical continuum presentation and close its orbit algebra
sourcewise.  Every target-active current, total-graph defect, and
target-visible limit has ancestry in
\(\Geom,\Caus,\Abs,\Lift,\Bdry\), while
\begin{equation}
 \label{eq:zero-saturated-residual-source}
 G_{\Budget,\Sigma}^{\Res,\mathrm{src},\Sat}=0.
\end{equation}
Every contact orbit has exactly one of two source-complete contact forms:
\begin{enumerate}[label=\textup{(C\arabic*)}]
\item cofinal approach contact supplies a canonical sequence of pairwise disjoint
pre-contact shell intervals \(I_j\) with unit shell increment; after passage
to a subsequence, one nonempty support \(A\subseteq\Channels\) occurs on
every \(I_j\) and has a positive owner there;
\item otherwise a cofinal sequence of terminal-entry shells supplies
completed shell records whose
positive owner is \(\Bdry\)-headed and belongs to
\(\mathfrak F_{\Sigma}^{\rm term}(\Budget)\).
\end{enumerate}
The alternatives are distinguished by the generated terminal-entry flag and
are therefore disjoint.  For the complete continuation target of
\cref{def:singular-target}, the completed same-origin restriction convention
in \textup{(P4)} forces \textup{(C1)}; \textup{(C2)} is retained for other
discontinuous or killed targets.  This statement
does not require a constitutive realization, classical chain rule, or price
gap: every failed relation contributes through its derived typed graph current
with the ancestry of the failed operation.
\end{theorem}

\begin{proof}
The continuum normal-form compiler and source-preserving completion give the
five-channel statement for ordinary values.  Derived graph currents give the
same statement for every boundary value.  The target quotient gives
\eqref{eq:zero-saturated-residual-source}: by construction residual means
target-null.  On the approach branch, the canonical disjoint shell-cascade
theorem supplies the intervals.  On each interval the derived
shell observable has increment one.  On a finite current record,
\eqref{eq:finite-affine-shell-graph} selects a positive source total; on a
variation-boundary record, \eqref{eq:closed-shell-owner-floor} selects its
positive owner.  The infinite pigeonhole principle selects a persistent
nonempty support.  If a shell is first entered at the terminal endpoint,
\cref{thm:deterministic-prospective-accountability} assigns the endpoint
jump a \(\Bdry\) head and the same owner floor.  The terminal-entry flag
separates this cell from the approach cascade.
\end{proof}

\begin{theorem}[Generated carrier-price regularity routing]
\label{thm:physical-saturated-accountability}
In the setting of
\cref{thm:unconditional-continuum-contact-alternative}, restrict to the
derived finite admissible-carrier stratum
\(\mu_u(\mathfrak Q_u)\le\Budget<\infty\) of ordinary maximal branches
beginning outside \(\Sigma\).  On every cofinal approach-contact branch, for
each persistent support use the canonical
sharp prices of \eqref{eq:canonical-source-price-schedule}.  Then the
framework-generated price sequence has the following exhaustive routing:
\begin{enumerate}[label=\textup{(P\arabic*)}]
\item if
\begin{equation}
 \label{eq:persistent-source-physical-price}
 \sum_{j=1}^{\infty}\underline p_{A,j}=\infty.
\end{equation}
then that support cannot sustain contact;
\item if the series is finite, the corresponding shell tails enter the
zero-density ordered path algebra with their unnormalized prefix ledger.  A checked identity
\(-1\in\mathcal K_{\mathfrak b}\) excludes the active cell; otherwise the order-unit
alternative supplies an augmented stationary zero-density law in
\(\mathfrak F_{A,m}^{\Sigma}(\Budget)\).
\end{enumerate}
If every approach-contact persistent support is removed either by the
divergent-price argument or by a finite ordered identity, and the generated
terminal cell \(\mathfrak F_{\Sigma}^{\rm term}(\Budget)\) is target-null,
continuable, or class-incompatible, then
\begin{equation}
 \label{eq:physical-target-unreachable}
 \boxed{\Reach^{\Phys}_\Sigma(\Budget)=0.}
\end{equation}
Thus finite carrier expenditure closes every divergent-price contact route.
A finite lower-price series is not assumed impossible and is never the end of
the target-reachability track: it is retained as a target-aligned zero-price
successor with its complete five-source ancestry.
\end{theorem}

\begin{proof}
If an ordinary maximal branch had cofinal approach contact with \(\Sigma\),
\cref{thm:unconditional-continuum-contact-alternative} would supply pairwise
disjoint intervals and a persistent support \(A\).  On their selected
subsequence, the definition of the sharp price and countable additivity give
\[
 \Budget\ge\sum_j\mu_u(I_j)
 \ge\sum_j\underline p_{A,j}=\infty,
\]
a contradiction in case \textup{(P1)}.  In the finite-series case, use the
actual prices as the price coordinate in the activity lift.  Cofinal tails
have price tending to zero, so
\cref{thm:finite-price-saturation-route,thm:ordered-proof-or-model}
give precisely the two outcomes in \textup{(P2)}.  If every persistent support
is eliminated, approach contact has no possible support, contradicting shell
accountability.  A remaining contact would have terminal-entry form, whose
\(\Bdry\)-headed owner belongs to
\(\mathfrak F_{\Sigma}^{\rm term}\); the stated terminal-cell condition
excludes it.  This proves \eqref{eq:physical-target-unreachable}.
\end{proof}

\begin{remark}[Orbitwise source ledgers]
The divergent-price branch is evaluated only for the source support that
persists along a putative singular cascade.  Its canonical lower-price
schedule may depend on the orbit stratum and may decay with scale.  A uniform
positive price is stronger and yields uniform window counts; neither is a
global estimate supplied to the framework.
\Cref{thm:physical-saturated-accountability} uses only divergence of this
generated series; five-channel exhaustion, the approach cascade, and the
separate terminal \(\Bdry\) cell were proved earlier.
\end{remark}

\begin{theorem}[Total constructive continuum prospective evaluator]
\label{thm:total-continuum-regularity-alternative}
Let \(u_0\) be a datum in the public state class of the fixed equation.
Totalize the local-solution and carrier graphs and run the constructive closure
of \cref{thm:end-to-end-fast-track-completeness} on every ordinary maximal
branch.  Put \(\eta_0=1/5\).  Exactly one of the following algebraic outcomes
holds:
\begin{enumerate}[label=\textup{(\roman*)}]
\item every contact-generated same-origin branch reaches a local
target-null, continuation, class-incompatibility, or divergent
physical-charge cell; then no ordinary maximal branch from \(u_0\) with
finite physical expenditure contacts \(\Sigma_{\rm sing}\).  Terminal
coverage consequently supplies continuation for every such ordinary branch;
\item diagonal target-aligned saturation reconstructs a prospective source
atom, including its closed reduced five-channel equation, physical ledger,
origin tag, and retained cofinal shell floor.
\end{enumerate}
Finite noncontinuation of an ordinary maximal path is contact with the
complete singular-event target and generates its same-origin cofinal
source-complete shell record.  The realization gate is required only before an
arbitrary closure point, or the terminal analytic coordinate itself, is
interpreted as an ordinary PDE object.  Finite noncontinuation never produces a
separate continuation class.  An empty ordinary
local-solution fiber is retained as the form value \(\bot_{\rm loc}\) and is
not converted into either regularity or singularity.  Positive stationary
functionals, equality kernels, and typed defects are intermediate successor
states and do not constitute a third outcome.
\end{theorem}

\begin{proof}
On an ordinary branch, finite maximal time without continuation gives the
same-origin singular event and source-complete terminal record of
\cref{prop:terminal-coverage}.  Its preterminal orbit already realizes
physical contact; a nonordinary terminal analytic coordinate remains a proper
graph successor and is never promoted to a regularity-class state.  Contact
then has one of the two forms in
\cref{thm:unconditional-continuum-contact-alternative}.  Cofinal approach
gives disjoint unit-increment shells and a persistent five-channel support.
A divergent sharp price contradicts the finite carrier, while a finite-price
cascade enters the zero-price successor algebra.  The normalized shell graph
gives a source owner at least \(1/5\), so no additional target-floor parameter
is needed.  Terminal entry instead enters
\(\mathfrak F_{\Sigma}^{\rm term}\) through its \(\Bdry\)-headed owner; its
total continuation graph either closes locally or yields the
\(\Bdry\)-headed source-profile witness in \textup{(ii)}.  Every finite local closing cell contradicts the
retained target floor or the finite physical ledger, so if all generated
contact branches close, contact is impossible.  Terminal coverage then turns
that target-local conclusion into continuation of each ordinary branch.
Otherwise, on an approach branch, the ordered proof/model alternative produces the initial
target-retaining compiler state.  Apply
\cref{thm:three-output-target-aligned-current-compiler}.  Its charge and null
outputs close the contact branch; its prospective source output is precisely
\textup{(ii)}.  Nonvacuity of the public graph follows from
\cref{prop:nonvacuous-solution-graph,thm:approximation-realization-first-failure}.
\end{proof}

\begin{theorem}[Target-aligned physical fast track]
\label{thm:target-aligned-physical-fast-track}
Fix a generated finite physical stratum and let
\(\Gamma_\Sigma\), \(\mathscr J_\Sigma\), and the sourcewise shell currents be
as above.  For each reachable active branch \(\mathfrak b\), let
\(\mathcal K_{\mathfrak b}\) be the ordered proof cone generated by
\cref{def:target-aligned-ordered-path-calculus,%
def:source-adapted-dual-carrier-grammar}.  Then
\begin{equation}
 \label{eq:fast-track-proof-model-alternative}
 \boxed{
 \Reach_\Sigma^{\Phys}(\Budget)=0
 \quad\text{or}\quad
 \mathfrak S_\Sigma(\Budget)\ne\varnothing,}
\end{equation}
where \(\mathfrak S_\Sigma(\Budget)\) is the set of joint-saturated
prospective source atoms generated by zero-price resaturation, retaining
their origin polarity.  A member with contact-derived origin is a source-profile
witness; only an independently generated member passing
\cref{thm:approximation-orbit-realization-criterion} is an explicit singular
mechanism.  At an intermediate finite
level, failure of an ordered identity constructs a normalized stationary
positive-functional successor satisfying
\begin{equation}
 \label{eq:fast-track-kernel}
 \Pi_\Sigma
 \left(
  \ker Q_{\rm phys}\cap
  \ker I_{\mathfrak b}\cap
  \bigcap_{w\in\mathscr W_{\mathfrak b}}\{P_w=0\}
 \right)\ne0
\end{equation}
for some reachable \(\mathfrak b\).  The successor is reprocessed by
\cref{def:target-aligned-recurrent-closure-compiler}; support, compactness,
passivity, fidelity, and realization values are internal queue coordinates.
By \cref{thm:three-output-target-aligned-current-compiler}, the successor is
never a third mathematical output.
\end{theorem}

\begin{proof}
Suppose contact occurs.  The derived shell-current theorem gives infinitely
many disjoint unit-increment events, each with a finite or boundary-owner
source floor.  Finite source exhaustion gives a
persistent nonempty support \(A\), and event normalization gives the universal
target floor \(1/5\).  If the physical prices have
divergent sum, countable additivity contradicts the finite physical budget.
Otherwise a subsequence has price tending to zero and the same-origin
successor construction gives an augmented stationary target-aligned
zero-density law in
an active branch \(\mathfrak b\).  This law defines the positive functional in
\cref{thm:ordered-proof-or-model}; weak duality forbids
\(-1\in\mathcal K_{\mathfrak b}\).  Run the target-aligned recurrent closure
compiler on this state.  Its divergent-charge and target-null outputs
contradict contact; its remaining output is the hypothetical prospective atom
in \(\mathfrak S_\Sigma(\Budget)\).  This proves
\eqref{eq:fast-track-proof-model-alternative}.

Conversely, failure of the ordered derivation gives the positive functional
by \cref{thm:ordered-proof-or-model}.  Its ideal relations are the weak PDE,
source-adapted zero-production, successor, and stationarity equations, while
its target floor makes the target projection nonzero.  This is
\eqref{eq:fast-track-kernel}.  Approximation realization is exactly
\cref{thm:approximation-realization-first-failure}; it distinguishes a
coherent generalized solution path from its next typed successor.
Applying the three-output current compiler to every typed or ordinary outcome
gives avoidance or the prospective source-profile output after the finite
physical ledger removes the divergent-charge cell.  The latter becomes a
singular-existence output only through the independent ordinary realization
criterion.
\end{proof}

\begin{corollary}[Generated PDE prospective-accountability output]
\label{cor:pde-target-reachability}
Let \(\mathfrak U_{\rm pub}\) be the datum class in the public PDE
presentation.  Derive the physical expenditure from the carrier identity and
run the source, signed-carrier, witness-successor, viability, and saturation
compilers on every resulting finite stratum.  Put \(\eta_0=1/5\).  Then
\begin{equation}
 \label{eq:continuation-contact-criterion}
 \boxed{
 \Reach_{\Sigma_{\rm sing}}^{\Phys}>0
 \quad\Longrightarrow\quad
 \mathsf T_{\Sigma_{\rm sing}}^{\Phys,\Sat}(\Budget)\ge\eta_0 .}
\end{equation}
The contact-generated output includes the cumulative prospective atom, its nonzero
five-channel source projection, unnormalized shell cocycle, closed reduced
equation, totalized realization rows, physical and origin ledgers, and the
cofinal shell floor.  Its origin polarity is \(\mathsf{Hyp}\), so it is a
lower-bound witness for the saturated source profile and not evidence that a
singular orbit exists.  Positive
functionals, compatible equality successors, and ordinary paths lacking this
joint atom are intermediate resaturation states.  Conversely,
\(\mathsf T_{\Sigma_{\rm sing}}^{\Phys,\Sat}(\Budget)<\eta_0\) proves
\(\Reach_{\Sigma_{\rm sing}}^{\Phys}=0\).  A singular-existence conclusion
is available only from an independently generated ordinary public orbit
satisfying \cref{def:explicit-singular-mechanism}.  None of these values is a
premise supplied with the PDE presentation.
\end{corollary}

\begin{proof}
Apply \cref{thm:target-aligned-physical-fast-track} at the derived finite
ceiling of each branch.  The finite ordered hierarchy initializes the local
queue, and \cref{thm:three-output-target-aligned-current-compiler} evaluates
all its descendants.  Under contact, divergent charge contradicts the finite
physical ledger and target-nullity contradicts the unit shell increment.
The remaining cumulative joint queue is reprojected by
\cref{thm:continuum-composition-no-creation}; its finite-prefix successors
continue the compiler, and its primal--polar fixed point retains source floor
at least \(\eta_0\).  Origin-polarity preservation makes this a hypothetical
source-profile witness.  Taking the contrapositive yields the regularity
criterion.  Independent realization is governed separately by
\cref{thm:no-recycled-contact-realization,def:explicit-singular-mechanism}.
\end{proof}

\begin{corollary}[Quantitative finite-cascade bound]
\label{cor:finite-cascade-bound}
Fix an active branch \(\mathfrak b\).  If the physical expenditure satisfies
\(\mu_u(\mathfrak Q_u)\le\Budget\) and each event in a disjoint same-origin
family has normalized window in \(B_{\Sigma,\mathfrak b}\) and admissible
target-local price at least \(\delta_{\Phys}(\mathfrak b)>0\), then that family has cardinality
at most
\begin{equation}
 \label{eq:finite-cascade-bound}
 \left\lfloor
 \frac{\Budget}{\delta_{\Phys}(\mathfrak b)}
 \right\rfloor.
\end{equation}
\end{corollary}

\begin{proof}
The proof of \cref{thm:physical-expenditure-exclusion} gives
\(N\delta_{\Phys}(\mathfrak b)\le\Budget\) for every finite subfamily.
Eligibility of each price is exactly
\cref{thm:local-target-aligned-charging-rule}; events outside
\(B_{\Sigma,\mathfrak b}\) do not enter this count.
\end{proof}

\begin{figure}[tbp]
\centering
\resizebox{\linewidth}{!}{%
\begin{tikzpicture}[fw diagram,x=1cm,y=1cm]
  \node[fw source,text width=2.45cm] (equation) at (0,0)
    {fixed PDE and\\physical carrier};
  \node[fw provenance,text width=2.5cm] (algebra) at (3.1,0)
    {closed five-source orbit algebra};
  \node[fw use,text width=2.5cm] (quotient) at (6.2,0)
    {totalized target-active\\contact interfaces};
  \node[fw residual,text width=2.55cm] (null) at (9.6,3.0)
    {no aligned structure\\unit crossing impossible};
  \node[fw use,text width=2.55cm] (defect) at (9.6,.8)
    {source-resolved\\graph current};
  \node[fw geom,text width=2.55cm] (ledger) at (9.6,-2.0)
    {active branch signature\\and record-scale chart};
  \node[fw terminal,text width=2.7cm] (nullout) at (13.8,3.0)
    {target unreachable\\at the finite ceiling};
  \node[fw use,text width=2.8cm] (defectout) at (13.8,.8)
    {witness-net\\successor expansion};
  \node[fw terminal,text width=2.8cm] (divout) at (13.8,-1.4)
    {dual carrier\\physical production};
  \node[fw use,text width=2.8cm] (finiteout) at (13.8,-3.4)
    {typed compactness\\successor};
  \node[fw geom,text width=2.8cm] (sat) at (17.3,-3.4)
    {local polarity refinement\\and diagonal saturation};
  \node[fw use,text width=2.8cm] (joint) at (20.8,-3.4)
    {cumulative joint record\\reproject all prefixes};
  \node[fw terminal,text width=2.8cm] (atom) at (24.4,-3.4)
    {saturated prospective\\source-profile atom};
  \draw[fw arrow] (equation)--(algebra);
  \draw[fw arrow] (algebra)--(quotient);
  \draw[fw arrow] (quotient)--(null);
  \draw[fw arrow] (quotient)--(defect);
  \draw[fw arrow] (quotient)--(ledger);
  \draw[fw arrow] (null)--(nullout);
  \draw[fw arrow] (defect)--(defectout);
  \draw[fw arrow] (defectout) to[bend left=16] (sat);
  \draw[fw arrow] (ledger)--(divout);
  \draw[fw arrow] (ledger)--(finiteout);
  \draw[fw arrow] (divout) to[bend left=15] (sat);
  \draw[fw arrow] (finiteout)--(sat);
  \draw[fw dashed arrow] (sat) to[bend left=24] (ledger);
  \draw[fw arrow] (sat)--(joint);
  \draw[fw dashed arrow] (joint) to[bend left=22] (ledger);
  \draw[fw arrow] (joint)--(atom);
\end{tikzpicture}%
}
\caption{Constructive singularity evaluator.  Only the target-active quotient
is analyzed.  If it contains no singularity-aligned structure, normalized
shell contact is impossible.  Otherwise contact fixes an active branch
signature and, when needed, a record-scale chart.  The source-adapted dual
carrier either charges target progress to physical production or creates a
typed compactness successor.  The next unprocessed coordinate is split only
on the retained target cell; every nonnull polarity cell re-enters the same
compiler through its cumulative join.  A cofinal refinement is saturated and
reprojected as one joint record.  A finite joint-prefix component returns to
the five-source compiler exactly when it is still unprocessed; all
sub-threshold processed increments remain in the cumulative sum.  Every
cumulative owner is then charged, annihilated, or expanded by recursive
carrier--polar exhaustion.  A primal--polar fixed atom is rejoined with its
unnormalized shell cocycle and complete source ancestry, then evaluated as an
ordinary same-origin reachable contact cell.}
\label{fig:structural-regularity-compiler}
\end{figure}

\FloatBarrier

\section{Structural evaluator coverage and failure semantics}
\label{sec:invariant-hierarchy}

The following table is the coverage proof for the fast track.  Every
classical analytic question is assigned a generated invariant and every
opposite value activates a structural operation.  No row accepts a desired
analytic conclusion as input or returns a bare label.

\begin{center}
\begin{tabular}{p{.25\linewidth}p{.33\linewidth}p{.32\linewidth}}
\toprule
Analytic question & Generated structural invariant & Successor action\\
\midrule
Does a native, chain, trace, or covariance relation hold? & target-active boundary
of its total graph & activity lift of the typed graph witness\\
Does a reduction descend through later operations? & commutator of the
operation with the quotient map & typed descent defect\\
Is the gradient range closed where the target sees it? &
\(\mathfrak H_{\rm grad}^{\Sigma}\) & inherited closure or harmonic class\\
Is a tail structurally closed? & \(\mathfrak H_{\rm tail}^{\Sigma}\) &
target-visible tail cokernel\\
Is target feeding dominated? & \(\alpha_*\) & supercritical or positive
null-cost feeding profile\\
Is normalized target progress physically priced? & \(\delta_*\) & zero-price
sequence in the same source cell\\
Can the source be routed? & \(\mathcal R(\Sigma)\) & target-visible harmonic
nonroutable class\\
Does a Wick quotient admit compensated target response? &
\(F_W(z)\), \(T_{W,\lambda}\), dark space, and bracket gap &
certified response/action profile or a typed dark, authority, or transfer
obstruction\\
Is the target polar? & \(\operatorname{Cap}(\Sigma)\) & positive-capacity
target-compatible class\\
Does a scale cascade exceed finite expenditure? & \(\mathscr S^\Sigma\) &
finite-price successor chain and recurrent saturation row\\
Do internal signed currents create net target feed? & carrier cohomology
\([F_{\rm sc}-D-M]\) & nonzero scale-feed cohomology successor\\
Does a failed interface persist dynamically? & recurrent witness boundary
\(\partial_\infty F\) & target-null loss or recurrent failure mechanism\\
Does a sharp branch close? & \(\mathfrak H_{\rm sat}^{\Sigma}\) & nonzero
saturation/equality remainder\\
Are stochastic fluctuations invisible or routed? & \(\chi^\Sigma\) and
fluctuation resistance & target-visible covariance/bracket class\\
Does a weak-solution graph value persist? & recurrent witness boundary
\(\partial_\infty F\) & finite-depth removal or recurrent weak-equation path\\
Can zero-price target progress persist? &
ordered cones \(\mathcal K_{\mathfrak b}\) and finite moment levels &
finite rational derivation or a reconstructed stationary law\\
\bottomrule
\end{tabular}
\end{center}

\begin{remark}[Exact input closure]
\label{rem:invariant-evaluation}
The fast-track data are constructed from the public equation, generated target,
source grammar, and native physical carrier.  A specialization cannot append
compactness, coercivity, rigidity, nonpolarity, continuation, or a desired
invariant value.  Each relation and its complement are graph values.  The
complementary value is converted into a typed successor on the surviving
target-active quotient, and its descendants are processed by finite-depth
viability within the same generated recursion.
\end{remark}

\begin{definition}[Target-local invariant evaluation packet]
\label{def:target-local-invariant-evaluation-packet}
Let \(I\) be any invariant in the fast-track table and let
\(\mathfrak b\) be an active same-origin branch.  Its evaluation packet is
\begin{equation}
 \label{eq:target-local-invariant-evaluation-packet}
 \mathfrak E_I(\mathfrak b)
 :=\bigl(
 X_{\mathfrak b},\overline{\operatorname{Graph}(I)},
 q_\Sigma I,Q_I,\Anc(I),p_{\Phys},
 (F_{I,n},C_{I,n})_{n\ge1},
 I_I^{=},\mathfrak G_I^{\rm real}
 \bigr).
\end{equation}
Here \(\overline{\operatorname{Graph}(I)}\) is the total graph of \(I\), including its
ordinary and first-failure values; \(q_\Sigma I\) is its target quotient;
\(Q_I\) is its signed target-current owner after compatible interfaces have
been glued; \(\Anc(I)\subseteq\Channels^*\) is its complete constructor
ancestry; \(p_{\Phys}\) is the restriction of the single public physical
carrier; \((F_{I,n},C_{I,n})\) is the canonical cofinal family of checked
finite cylinders and carrier cones; \(I_I^{=}\) is the equality ideal
generated on zero physical valuation; and \(\mathfrak G_I^{\rm real}\) is
the same-origin ordinary realization graph.  Every entry is generated from
the public presentation, the selected shell covector, and earlier branch
relations.  No invariant value is part of the packet data.
\end{definition}

\begin{definition}[Invariant evaluation protocol]
\label{def:invariant-evaluation-protocol}
For an initial packet \(\mathfrak E_{I_0}(\mathfrak b_0)\), maintain a FIFO
queue \(\mathcal Q_n\) and an unnormalized cumulative same-origin joint record
\(\mathcal J_n\).  Initially
\(\mathcal Q_0=(\mathfrak E_{I_0}(\mathfrak b_0))\) and
\(\mathcal J_0=\mathfrak E_{I_0}(\mathfrak b_0)\).  On removing the first
packet \(\mathfrak E_I(\mathfrak b)\), perform these operations in order:
\begin{enumerate}[label=\textup{(P\arabic*)}]
\item project through \(q_\Sigma\), deleting a zero projection by its checked
  target-nullity row;
\item compute the next finite-cylinder physical valuation and source profile,
  deleting the descendant stratum only when a checked profile ceiling lies
  below the retained shell threshold;
\item on threshold equality, adjoin every vanishing production, square,
  polarization, adjoint, symmetry, and dyadic coercivity row, reproject, and
  append every nonzero descendant;
\item on a strict dual or positive-profile value, join primal owner and dual
  covector, pull the covector through the complete backward frontier, and
  append every nonzero conditional constructor increment;
\item on a nonordinary value, append the least failed total-graph coordinate
  with its owner, target covector, physical denominator, occurrence cocycle,
  origin, and preceding joint frontier;
\item when the finite queue is empty, evaluate the realization fiber of the
  cumulative joint record.  A finite incompatible fiber gives its checked
  prefix row; a coherent zero-defect branch enters the ordinary realization
  gate; and projective branch escape is appended as a \(\Lift\)-owned
  realization-fiber packet.
\end{enumerate}
Every appended packet is joined without normalization into
\(\mathcal J_{n+1}\).  At a cofinal execution form
\(\mathcal J_\infty=\bigvee_n\mathcal J_n\) before projection.  A nonzero
finite contextual increment re-enters \textup{(P3)}--\textup{(P5)} at its
least occurrence.  If all finite structural increments have been processed,
evaluate the complete joint realization fiber: finite incompatibility is
checked on its prefix, a coherent branch enters the ordinary gate, and branch
escape re-enters \textup{(P5)} through its total graph.  Profile witnesses,
equality faces, polar separators, graph
boundaries, and compatible inverse-limit functionals are protocol states,
never stopping values.
\end{definition}

\begin{theorem}[Unconditional local evaluator for every generated invariant]
\label{thm:unconditional-target-local-invariant-evaluator}
For every packet \(\mathfrak E_I(\mathfrak b)\), the continuum compiler
constructs the extended physical valuation
\begin{equation}
 \label{eq:invariant-saturated-valuation}
 \operatorname{val}^{\Sat}_I(Q_I)
 =\sup_n\inf\{C\ge0:Q_{I,n}\in D_{I,n}+Cp_{I,n}\}
\end{equation}
and the saturated source profile
\begin{equation}
 \label{eq:invariant-saturated-profile}
 \mathsf T_I^{\Sat}(E)
 =\inf_n\sup\{L(Q_{I,n}):L\in\mathcal S_{I,n}(E)\}.
\end{equation}
At each rational target threshold \(\eta>0\), these values drive the
exhaustive transitions of
\cref{def:invariant-evaluation-protocol}:
\begin{enumerate}[label=\textup{(I\arabic*)}]
\item a finite checked carrier or nullity row proves
  \(\mathsf T_{I,n}(E)<\eta\) on some cylinder and hence subcriticality on
  the entire descendant stratum;
\item the threshold equality face is generated, its vanishing productions
  and squares are adjoined to \(I_I^{=}\), and target projection is repeated;
\item a positive target-active value supplies a finite contextual dual
  increment or a compatible cofinal functional, retaining \(Q_I\), its
  unnormalized occurrence cocycle, physical denominator, and complete
  five-source ancestry;
\item a nonordinary coefficient supplies its least failed graph coordinate
  with the same data and is evaluated as the next invariant packet;
\item an ordinary zero-defect same-origin path supplies the corresponding
  physical value of \(I\) and alone can activate the contact realization
  row.
\end{enumerate}
The first alternative is an unconditional subcriticality proof on its
generated stratum.  In the other alternatives the exact non-subcritical
structure is an evaluated five-source record and is immediately returned to
the protocol.  None of \textup{(I2)}--\textup{(I4)} is an output.  Repeated
evaluation uses the cumulative joint record, so a
finite or infinite family of individually subthreshold coordinates cannot
lose its joint target projection.
\end{theorem}

\begin{proof}
Projection-first saturation gives \(q_\Sigma I\) and removes only its
target-null kernel.  Source exhaustion assigns every remaining atomic and
mixed current to its constructor ancestry before any estimate.  On a finite
cylinder, \cref{thm:finite-prefix-source-profile-carrier-duality} gives
\eqref{eq:invariant-saturated-valuation}, while
\cref{thm:exact-finite-cylinder-source-profile-duality} gives the finite
version of \eqref{eq:invariant-saturated-profile}.  Cofinal monotonicity is
\cref{thm:monotone-saturated-physical-valuation,thm:monotone-cylinder-source-profile-completion}.
The equality alternative is processed by
\cref{thm:saturation-equation-compiler,thm:dyadic-ordered-coercivity-terminal-rule}.
The total-graph alternative is
\cref{thm:nonordinary-shell-current-first-failure-evaluator}; its output is a
new packet because the failed operation, covector, owner, carrier, and origin
are all retained.  Ordinary realization is governed by
\cref{thm:approximation-orbit-realization-criterion}.  Finally,
\cref{thm:exact-prospective-continuum-backward-factorization,thm:cofinal-prospective-frontier-closure}
gives the cumulative joint assertion.  These constructions use only packet
entries and therefore introduce no analytic premise.
\end{proof}

\begin{definition}[Realization-fiber invariant]
\label{def:realization-fiber-invariant}
For a cumulative joint packet \(\mathcal J\), let
\(\mathscr F_n(\mathcal J)\) be the set of depth-\(n\) public approximation
prefixes having the same origin and initial datum as \(\mathcal J\), matching
its first \(n\) cylinder coordinates within \(2^{-n}\), and having its first
\(n\) normalized public residuals at most \(2^{-n}\).  Restriction gives the
rooted inverse graph
\begin{equation}
 \label{eq:realization-fiber-inverse-graph}
 \mathscr F(\mathcal J)
 :=\varprojlim_n\mathscr F_n(\mathcal J).
\end{equation}
The realization-fiber invariant is the total graph with values
\begin{enumerate}[label=\textup{(R\arabic*)}]
\item the least \(n\) for which \(\mathscr F_n(\mathcal J)=\varnothing\),
  together with the finite cylinder relations defining that fiber;
\item one element of \(\mathscr F(\mathcal J)\), with all residual-tail,
  origin, carrier, and continuation coordinates;
\item the projective branch-escape value
  \(\mathscr F_n(\mathcal J)\ne\varnothing\) for every \(n\) but
  \(\mathscr F(\mathcal J)=\varnothing\), retaining the escaping prefix net
  as one unnormalized cumulative graph packet; this value has no distinguished
  failed finite coordinate.
\end{enumerate}
The third value is a \(\Lift\)-owned target-local invariant.  It is not an
ordinary path and is inserted into the invariant protocol.
\end{definition}

\begin{theorem}[Exact realization-fiber evaluation]
\label{thm:exact-realization-fiber-evaluation}
Every cumulative joint packet has exactly one realization-fiber value from
\cref{def:realization-fiber-invariant}.  In value \textup{(R1)}, the least
empty prefix is retained together with its entire finite relation set.  A
finite algebraic separation identity closes it as checked incompatibility;
until such an identity is generated, the total emptiness graph is itself the
next finite-prefix packet and is not treated as a conclusion.  In value
\textup{(R2)}, residual-tail
evaluation gives either a least nonzero public defect or a zero-defect
compatible branch; the latter enters
\cref{thm:dirac-realization-gate}.  Value \textup{(R3)} re-enters
\cref{def:invariant-evaluation-protocol} with its complete target pairing and
five-source ancestry.  Hence absence of an approximation branch is itself an
evaluated structural invariant and never an existence conclusion.
\end{theorem}

\begin{proof}
Either some finite fiber is empty or all are nonempty.  In the latter case
the inverse limit is empty or nonempty, giving the three disjoint values.
Finite-cylinder separation gives a closing identity whenever the generated
finite convex cylinder relaxation is empty.  If the relaxation is nonempty
while the exact prefix fiber is empty, their discrepancy is a nonzero
finite realization-graph coordinate and is re-enqueued; thus emptiness is
evaluated through the generated discrepancy coordinate.
For \textup{(R2)}, the monotone residual-tail coordinates in
\cref{def:exhaustive-approximation-tower} expose the least nonzero defect; if
none exists, the branch is zero defect and the Dirac gate applies.  In
\textup{(R3)}, totalization of the restriction graph records the escaping
net before passage to its cumulative graph boundary.  Projection-first
saturation either removes its target pairing or assigns its first nonzero
conditional increment to the constructor that generated the failed
realization map.  The least-occurrence theorem is invoked only if such a
finite increment exists; otherwise the complete unnormalized branch-escape
packet remains the next \(\Lift\)-owned invariant.  This is exactly the
packet required by the protocol.
\end{proof}

\begin{theorem}[Closure and conservation of the invariant protocol]
\label{thm:invariant-evaluation-protocol-closure}
Every finite or cofinal execution of
\cref{def:invariant-evaluation-protocol} preserves the monotone ledger
\begin{equation}
 \label{eq:invariant-protocol-monotone-ledger}
 (\text{origin},\text{shell cocycle},\text{physical carrier},
   \text{complete joint frontier},\text{five-source ancestry}).
\end{equation}
Every live state has a prescribed next operation.  The only stopping values
are:
\begin{enumerate}[label=\textup{(T\arabic*)}]
\item a checked finite target-nullity or subcritical carrier derivation;
\item an ordinary zero-defect same-origin path, whose continuation flag
  decides target contact.
\end{enumerate}
Thus target-visible residuals, fixed profiles, equality faces, separators,
nonordinary graphs, and source-complete graph records are nonterminal protocol states.
\end{theorem}

\begin{proof}
Projection changes only the target-null quotient.  Finite-cylinder valuation
and equality restriction preserve origin and ancestry, while signed gluing
preserves the unnormalized shell cocycle and public carrier.  Backward
factorization assigns a nonzero dual pairing to a finite contextual
constructor increment, so \textup{(P4)} preserves the full preceding joint
frontier.  Source-preserving totalization gives the same conclusion for
\textup{(P5)}.  Induction proves
\eqref{eq:invariant-protocol-monotone-ledger}.

For a cofinal execution,
\cref{thm:cofinal-prospective-frontier-closure,thm:continuum-composition-no-creation}
shows that every unprocessed nonzero target component has a least finite
contextual occurrence and therefore returns to a finite transition.  On the
complementary fully processed joint record,
\cref{thm:exact-realization-fiber-evaluation} gives finite fiber evaluation,
a zero-defect branch entering the Dirac gate, or a branch-escape packet
entering transition \textup{(P5)}.  Hence every nonterminal state has the
successor prescribed in the definition, and only \textup{(T1)}--\textup{(T2)}
stop.
\end{proof}

\begin{center}
\begin{tabular}{p{.22\linewidth}p{.35\linewidth}p{.31\linewidth}}
\toprule
Protocol state & Mandatory next operation & Permitted closure\\
\midrule
target-null packet & checked quotient deletion & finite nullity\\
finite subcritical profile & replay dual carrier row & finite subcriticality\\
equality face & enlarge equality ideal and reproject & none\\
dual/profile witness & backward pullback and constructor split & none\\
nonordinary graph value & least-failure packet & none\\
cofinal cumulative join & finite contextual re-entry or realization graph & none\\
complete detector upper profiles & cofinal zero derivation & structural absence\\
uniform positive detector floor & compatible functional and backward finite occurrence & structural presence; then realization\\
coincident valuation cuts & rational-cut comparison & exact extended-real value\\
separated valuation cuts & target-project the cut-gap owner & structural presence; then realization\\
zero-defect ordinary path & continuation/contact readout & ordinary physical value\\
\bottomrule
\end{tabular}
\end{center}

\begin{theorem}[Invariant-by-invariant fast-track evaluation]
\label{thm:fast-track-invariant-by-invariant-evaluation}
The protocol of
\cref{def:invariant-evaluation-protocol,thm:unconditional-target-local-invariant-evaluator,thm:invariant-evaluation-protocol-closure}
applies to every
row of the fast-track table as follows.
\begin{enumerate}[label=\textup{(F\arabic*)}]
\item Native, chain, trace, and covariance relations use their total-graph
  boundary as \(I\); the least failed relation is the next \(\Lift\)-owned
  activity packet.
\item Quotient descent uses the commutator with the quotient map; its zero
  cell descends, and every nonzero cell is an \(\Abs\)-owned contextual
  increment.
\item Gradient closure uses \(\mathfrak H_{\rm grad}^{\Sigma}\); its zero
  class is the closed-range cell and its nonzero class is retained with
  \(\Geom/\Caus\) gradient ancestry.
\item Constraint-tail closure uses
  \(\mathfrak H_{\rm tail}^{\Sigma}\); signed interface gluing removes exact
  tails and the surviving cokernel is a \(\Bdry\)-owned packet.
\item Feeding domination uses \(\alpha_*\); finite carrier domination is
  subcritical, while supercritical or null-cost feeding retains the aggregate
  signed current and its source join.
\item Physical pricing uses \(\delta_*\); a positive price is charged to
  \(p_{\Phys}\), and a zero-price sequence enters equality saturation with
  its unnormalized numerator.
\item Routing uses \(\mathcal R(\Sigma)\); a routed value inherits its
  destination owner, and a nonroutable harmonic value remains in the
  target quotient with the owner of the failed routing constructor.
\item Wick response uses \(F_W(z)\), \(T_{W,\lambda}\), the dark space, and
  the bracket gap; response and action are physically valued, while dark,
  authority, and transfer defects retain their generated graph owners.
\item Target polarity uses \(\operatorname{Cap}(\Sigma)\); zero capacity is
  target-null and positive capacity is a target-compatible \(\Geom\) packet.
\item Scale expenditure uses \(\mathscr S^\Sigma\); every scale step retains
  its pulled-back physical weight, and recurrent saturation keeps the full
  multiplicity cocycle.
\item Signed internal feeding uses the carrier-cohomology class
  \([F_{\rm sc}-D-M]\); exact classes telescope, while nonzero classes are
  \(\Lift/\Caus\)-owned scale-feed packets.
\item Persistent interface failure uses \(\partial_\infty F\); loss of the
  target pairing is null, while a nonzero recurrent boundary is evaluated
  with its interface owner.
\item Sharp equality uses \(\mathfrak H_{\rm sat}^{\Sigma}\); the zero class
  enters the generated equality ideal and a nonzero remainder returns to the
  source profile.
\item Stochastic visibility uses \(\chi^\Sigma\) and fluctuation resistance;
  zero covariance is target-null and a positive covariance or bracket class
  has \(\Geom/\Caus\) ancestry.
\item Weak-solution persistence uses the total weak-equation boundary;
  finite-depth removal is checked by an equation identity and a recurrent
  boundary retains its ordinary/nonordinary realization coordinate.
\item Zero-price progress uses \(\mathcal K_{\mathfrak b}\) and its finite
  moment levels; finite infeasibility gives a checked nullity identity, while
  compatible levels produce one joined positive functional that immediately
  enters equality and ordinary-realization evaluation.
\end{enumerate}
For every row, the complementary value is therefore a concrete packet for the
same protocol and is enqueued immediately.  None of the sixteen checks
requests compactness, coercivity,
rigidity, a critical norm, or the desired invariant value from the
specialization.
\end{theorem}

\begin{proof}
For \textup{(F1)} and \textup{(F15)}, apply defect totalization and
\cref{thm:approximation-realization-first-failure}.  Quotient functoriality
and target-relative gradient exhaustion give \textup{(F2)}--\textup{(F3)}.
\Cref{thm:constraint-tail-normal-form,prop:constraint-tail-carrier-law} gives
\textup{(F4)}.  Carrier duality and the physical-price interface give
\textup{(F5)}--\textup{(F6)}.  Zero-mode routing and the Wick total graphs give
\textup{(F7)}--\textup{(F8)}.  Capacity polarity gives \textup{(F9)}.
\Cref{thm:scale-cascade-price-compiler,def:carrier-cohomology} gives
\textup{(F10)}--\textup{(F11)}.  Recurrent boundary totalization and the
saturation-equation compiler give \textup{(F12)}--\textup{(F13)}.  The
stochastic gradient/covariance source decomposition gives \textup{(F14)}.
Finally, ordered proof-or-model evaluation and the Dirac realization gate
give \textup{(F16)}.  Each cited construction produces precisely the packet
entries in \eqref{eq:target-local-invariant-evaluation-packet}; applying
\cref{thm:unconditional-target-local-invariant-evaluator} proves every row.
\end{proof}

\begin{remark}[Separation of invariant roles]
The logical interfaces enforce ten distinctions.
\begin{enumerate}[label=\textup{(N\arabic*)}]
\item An imposed \(\Caus\) term is transported by a chart and is not
  optimized toward the target.  A free or hybrid authority envelope may
  optimize only its declared represented control slot.
\item A killed spectral resolvent is mapped to arrival only through the
  totalized terminal-flux realization row; a nonzero realization defect is
  retained.
\item Finiteness of a basic energy does not imply a positive cost at every
  singular scale; the validity or failure of
  \eqref{eq:physical-price-domination} is the physical-price row.
\item Local entropy monotonicity does not yield an orbitwise production bound without the
  reset relation; its failure is the corresponding cocycle graph defect.
\item Preservation of a nonlinear badness predicate under weak limits is the
  boundary-defect value in the canonical evaluation closure.
\item Failure of an interface creates its totalized graph-defect record.
  Target visibility makes that defect part of the singularity-formation
  algebra; target invisibility places it in \(J_{\rm res}\).
\item Large or noncompact motion in \(\Null_\Sigma\) requires no downstream
  estimate because it cannot increase a declared prospective target
  functional.
\item A quotient is propagated through the zero descent-defect branch;
  failure to factor is retained as a target-active descent defect.
\item Summability of a constraint tail closes the analytic tail but supplies
  a physical price only through a carrier comparison such as
  \eqref{eq:constraint-tail-physical-domination}.
\item A positive normalized scale gap excludes an infinite cascade when the
  pulled-back lower prices have the divergent total mass required by
  \eqref{eq:divergent-scale-weight}; a nonvanishing infimum is a stronger
  sufficient condition.
\end{enumerate}
\end{remark}

\section{A reusable PDE instantiation protocol}
\label{sec:pde-instantiation}

For a particular deterministic PDE, the specialization datum is only its
public textbook presentation from
\cref{prop:public-presentation-contract,cor:textbook-only-specialization}.
Recording that datum requires no new mathematical result.  The compiler then
performs the following generated operations.
\begin{enumerate}[label=\textup{(I\arabic*)}]
\item Read the displayed equation, datum class, domain, boundary conditions,
  weak test core, declared regularity class, native transformations, and
  formal carrier multiplier.  Enumerate the exhaustive rational approximation
  grammar, its tested residuals, and its signed carrier records.  Generate the
  total solution graph, continuation graph, and complete singular target.
\item Form the graph-complete ordered cylinder algebra.  Every nonlinear
  operation carries its input, output, oscillation, concentration, tail, and
  first-failure coordinates.  Ordinary states embed as evaluation states, and
  zero-defect approximation paths satisfy the public weak equation.
\item Expand the formal carrier multiplier against the weak equation, glue
  common faces with their signs, and take the Jordan decomposition.  This
  generates the one measure \(\mu_u\) and its derived ceiling.  Every
  nonordinary value is inserted into the activity lift; no replacement
  measure is requested.
\item Run the normal-form source compiler.  It labels native mobility or
  reversible form \(\Geom\), fixed directed transport \(\Caus\), quotient
  operations \(\Abs\), generated charts and augmented variables \(\Lift\),
  and terminal or exterior fluxes \(\Bdry\).  Construct the primitive
  expression graph, apply its canonical form split, and propagate both the
  unique current head of each atomic component and complete provenance
  support through every constructor.  The universal provenance theorem then verifies that the
  resulting \(31\) support cells contain every public current path.
\item Form the sourcewise mathematical closure of every dynamically generated
  weak limit, concentration, oscillation, tail, defect, and blow-up germ.
\item Generate finite-cylinder target shells and their source-resolved
  Stieltjes currents.  Contact supplies unit shell progress, the canonical
  disjoint cascade, a persistent five-source support, its detecting covector,
  and the same-origin active signature.  The signature also contains the
  canonical record-scale chart when the selected scales collapse.  Directions
  annihilated by the shell covector are removed immediately in
  \(\Null_\Sigma\).  Apply projection-first saturation; every nonzero descent
  or normalization defect retains the ancestry of its failed operation.
\item On the resulting target spine, generate the forward ancestry module and
  the backward target-cyclic covector module.  Compute their finite incidence
  cores and pass to the minimal target-active realization
  \(\mathsf M_{\Sigma,\mathfrak b}\).  Delete every equation, constraint,
  chart, defect, tail, and carrier coordinate outside a complete path from a
  source occurrence to the shell covector.  On the retained realization,
  restrict the public equation, constraints, gradient structure, native
  carrier, charts, and total graphs.  Glue compatible currents with their signs and form the single aggregate row
  \(\mathfrak J_{\Sigma,\mathfrak b}\).  This step simultaneously generates
  the retained target-cyclic covectors, nonnegative productions, constraint tails,
  commutator squares, compactness coordinates, feeding cohomology, and exact
  storage transfer.  Only carrier-dominated production can spend the public
  physical measure.
\item On the minimal realization, apply the five canonical normal-form
  projections: \cref{sec:continuum-geom-normal-form},
  \cref{sec:continuum-caus-normal-form},
  \cref{sec:continuum-abs-normal-form},
  \cref{sec:continuum-lift-normal-form}, and
  \cref{sec:continuum-bdry-normal-form}.  Compute the geometric
  range/kernel/capacity tuple; the directed cohomology, storage, passivity,
  circulation, and compressed-resolvent tuple; the quotient
  descent/vertical/cokernel/target-descent tuple; the lift
  interface/fiber/concentration/oscillation/scale tuple; and the boundary
  exterior/tail/terminal tuple.  Enumerate the \(31\) possible nonempty source supports
  and form their nested invariant images
  \(\mathfrak F_{A,n}^{\Sigma}(E)\).  This step is generated from the public
  graphs and target covectors; no value of an invariant is supplied.
\item Run the target-aligned recurrent closure compiler of
  \cref{def:target-aligned-recurrent-closure-compiler}.  Its internal queue
  computes the recurrent circulation price and processes zero kernels,
  signed-dyadic return laws, and first graph-boundary values in one fixed
  order.  Positive price returns a generated carrier potential; zero price
  returns the next circulation support.  By
  \cref{thm:three-output-target-aligned-current-compiler}, the output is a
  divergent physical charge, target-nullity on the retained path, or a finite
  contact-derived joint source occurrence with unchanged origin polarity.
  That occurrence is rejoined with its complete same-origin history and
  evaluated in the reachable contact cell.  Hence no
  separate compactness, equality, tail, pressure, scale, or rigidity stage is
  part of the specialization protocol.  At each finite depth, the incidence
  fast track precedes the ordered moment calculation, so no source-transverse
  coordinate enters that calculation.
\end{enumerate}

This protocol is equation-agnostic.  Different continuum equations generate
different row values in steps \textup{(I2)}, \textup{(I5)}, and
\textup{(I7)}--\textup{(I9)}; none supplies desired values as additional
input.  The closed source algebra, local target restriction, charging rule,
and failure semantics do not change.  Universal analytic identities used in
any step are activated only through
\cref{thm:target-local-fact-factorization}.

\subsection{Target-cyclic degeneracy compilation}
\label{subsec:target-cyclic-degeneracy}

The equality case of a physical carrier law must itself be evaluated by the
closed algebra.  It is not an auxiliary rigidity hypothesis.  The following
construction extracts all consequences that the public equation forces from
a target covector and from vanishing physical production.

\begin{definition}[Target-cyclic covector module]
\label{def:target-cyclic-covector-module}
Fix an active branch \(\mathfrak b\), a selected target shell \(\Phi\), and
the five-channel current
\[
 J=J^{\Geom}+J^{\Caus}+J^{\Abs}+J^{\Lift}+J^{\Bdry}.
\]
Let \(\mathscr C^0_{\mathfrak b}\) be the rational localization module
generated by \(D\Phi\) on the retained cylinder algebra.  Inductively,
\(\mathscr C^{n+1}_{\mathfrak b}\) is obtained from
\(\mathscr C^n_{\mathfrak b}\) by adjoining every ordinary value of
\begin{enumerate}[label=\textup{(C\arabic*)}]
\item the adjoint transports \((DJ^c)^*\lambda\), the covector Lie
derivatives \(\mathcal L_{J^c}\lambda\), and all iterated commutators of
these operations, for \(c\in\{\Geom,\Caus,\Abs,\Lift,\Bdry\}\);
\item multiplication by rational cylinder cutoffs, restriction to record
scales, pullback by generated charts, and descent through the generated
constraint, gauge, and boundary quotients;
\item polarization of every multilinear operation in the public weak
equation and of every defect form generated by its total graph; and
\item the source-adapted carrier and square-completion operations of
\cref{def:source-adapted-dual-carrier-grammar}.
\end{enumerate}
Here \(DJ^c\) denotes the total graph of rational difference quotients on
the cylinder core.  Its ordinary value is the derivative when that derivative
exists; no differentiability of the weak current is presumed.
Every operation is totalized.  Its first nonordinary value is routed with its
source ancestry and the covector that exposed it; it is not adjoined as an
ordinary word.  Define
\begin{equation}
 \label{eq:target-cyclic-module}
 \mathscr C^\infty_{\mathfrak b}
 :=\overline{\bigcup_{n\geq0}\mathscr C^n_{\mathfrak b}}^{\rm cyl}.
\end{equation}
This is the least covector module forced by the public dynamics and the
target shell.  In particular, no multiplier, compactness statement, or
rigidity assertion is supplied as specialization data.
\end{definition}

\begin{theorem}[Target-cyclic consequence completeness]
\label{thm:target-cyclic-consequence-completeness}
The module \(\mathscr C^\infty_{\mathfrak b}\) contains every finite
target-relevant consequence derivable from the selected shell covector by the
declared target-cyclic grammar: adjoint transport, Lie differentiation,
commutation, localization, quotient or chart pullback, boundary pullback,
polarization, square completion, carrier coboundary formation, and
total-graph completion.  A missing ordinary value is retained as the typed
successor of its first failed operation.  This is completeness for the
declared public consequence grammar.
\end{theorem}

\begin{proof}
Every finite derivation has a finite syntax tree.  Its leaves lie in
\(\mathscr C^0_{\mathfrak b}\), and each internal node is one constructor
listed in \textup{(C1)}--\textup{(C4)} of
\cref{def:target-cyclic-covector-module}.  Induction on tree depth places
every ordinary node in some \(\mathscr C^n_{\mathfrak b}\).  Totalization
supplies the typed first-failure value at every nonordinary node, with its
source ancestry and detecting covector.  Cylinder closure then gives the
stated module.  No claim concerns analytic consequences outside the declared
grammar.
\end{proof}

\begin{definition}[Generated production family and joint kernel]
\label{def:generated-production-joint-kernel}
For \(\lambda\in\mathscr C^\infty_{\mathfrak b}\), expand the exact weak
identity obtained by pairing \(\lambda\) with the public equation.  Glue
signed internal currents before taking positive parts.  The
\emph{target-cyclic production family} \(\mathscr P_{\mathfrak b}\) consists
of every nonnegative term produced algebraically by that expansion:
mobility squares, Dirichlet and commutator forms, positive graph defects,
boundary production, and their polarized positive localizations.  A member
of \(\mathscr P_{\mathfrak b}\) is charged to the physical ledger only when
the same expansion proves its domination by the public carrier measure.
Otherwise it is retained as an equality-kernel coordinate.  Put
\begin{equation}
 \label{eq:target-cyclic-joint-kernel}
 X^n_{\mathfrak b}
 :=X_{\mathfrak b}\cap
   \bigcap_{\substack{P\in\mathscr P_{\mathfrak b}\cr
                      \operatorname{depth}(P)\leq n}}
   \{P=0\},
 \qquad
 X^\infty_{\mathfrak b}:=\bigcap_{n\geq0}X^n_{\mathfrak b}.
\end{equation}
The zero set is interpreted in the total ordered graph, so failure of lower
semicontinuity, closedness, trace descent, or constraint descent is itself a
nonzero typed coordinate with inherited source.  Thus a missing analytic
property enlarges the generated structure; it is never presumed away.
\end{definition}

\begin{theorem}[Zero-production kernel soundness]
\label{thm:zero-production-kernel-soundness}
Let \(\gamma\) be a zero-price recurrent support selector and let
\(P\in\mathscr P_{\mathfrak b}\) be a generated nonnegative production with
\(0\le P\le p_{\Phys}\) on its carrier-admission row.  Then
\[
 \int P\,\dd\gamma=0,
 \qquad P=0\quad\gamma\text{-almost surely}.
\]
The resaturation restriction to
\(X^\infty_{\mathfrak b}\) of
\eqref{eq:target-cyclic-joint-kernel} adjoins only these pointwise production
relations.  It retains the unnormalised occurrence, expenditure, and
source-current coordinates
\((\Gamma_N,\Pi_N,(Q_N^c)_c)_N\) of the selected path.
\end{theorem}

\begin{proof}
Zero price gives \(\int p_{\Phys}\,\dd\gamma=0\).  Positivity and the
domination \(P\le p_{\Phys}\) give \(\int P\,\dd\gamma=0\), hence
\(P=0\) almost surely.  The kernel definition adds the resulting zero
relations to the ordered branch ideal.  By
\cref{thm:circulation-multiplicity-restoration}, the normalized selector is
stored jointly with the unnormalised occurrence ledger rather than replacing
it.  Thus kernel restriction changes neither \(\Pi_N\) nor any \(Q_N^c\).
\end{proof}

\begin{definition}[Forward structural ancestry module]
\label{def:forward-structural-ancestry-module}
Fix an active signature \(\mathfrak b\).  Let
\(\mathscr R^0_{\mathfrak b}\) consist of the five projected source-current
occurrences in its unit shell identity, with their origin and carrier rows.
Inductively, \(\mathscr R^{n+1}_{\mathfrak b}\) is obtained from
\(\mathscr R^n_{\mathfrak b}\) by adjoining the ordinary outputs and first
nonordinary total-graph values of every public operation that accepts a word
already present: weak evolution, localization, multilinear polarization,
constraint and gauge descent, chart transport, boundary trace, signed
interface gluing, carrier formation, return, and same-origin succession.  Put
\begin{equation}
 \label{eq:forward-structural-ancestry-module}
 \mathscr R^\infty_{\mathfrak b}
 :=\overline{\bigcup_{n\ge0}\mathscr R^n_{\mathfrak b}}^{\rm cyl}.
\end{equation}
Each word retains its complete source ancestry.  This is the forward closure
of what the public equation can generate from the source occurrence; it is
formed without a norm bound or compactness premise.
\end{definition}

\begin{definition}[Target-incidence core and minimal realization]
\label{def:target-incidence-minimal-realization}
For finite depths \(m,n\), form the bipartite total-incidence graph
\(\mathsf G^{m,n}_{\Sigma,\mathfrak b}\).  Its left vertices are the words of
\(\mathscr R^m_{\mathfrak b}\), its right vertices are the covectors of
\(\mathscr C^n_{\mathfrak b}\), and it contains:
\begin{enumerate}[label=\textup{(I\arabic*)}]
\item an ordinary edge \((R,\lambda)\) for every nonzero signed-dyadic cell
of the pairing \(\langle\lambda,R\rangle\);
\item a typed edge for the first nonordinary value of that pairing or of an
operation needed to form it; and
\item the directed constructor edges recording forward ancestry and backward
adjoint transport; and
\item one typed frontier vertex for each well-typed public constructor schema
whose next instance has not yet been expanded, with edges from its admitted
input types to its possible output and adjoint types.
\end{enumerate}
Delete every vertex that lies on no directed path from one of the five source
roots to the root target covector \(D\Phi\), allowing paths through frontier
vertices, and repeat deletion until stable.  A frontier vertex is replaced by
its ordinary, zero, signed-dyadic, and graph-boundary values when its instance
enters the enumeration.
Denote the surviving finite graph by
\(\mathsf M^{m,n}_{\Sigma,\mathfrak b}\).  Its diagonal closure
\begin{equation}
 \label{eq:minimal-target-active-realization}
 \mathsf M_{\Sigma,\mathfrak b}
 :=\overline{\bigcup_{k\ge0}
       \mathsf M^{k,k}_{\Sigma,\mathfrak b}}^{\rm gc}
\end{equation}
is the \emph{minimal target-active realization}; frontier vertices are absent
from this diagonal closure because every fixed constructor instance is
eventually expanded.  A pairing value zero in the
ordinary graph is deleted; failure of the zero pairing to descend through a
constructor is its typed incidence edge and survives precisely when it lies
on a source--target path.
\end{definition}

\begin{theorem}[Target-incidence compression]
\label{thm:target-incidence-compression}
For every active signature \(\mathfrak b\):
\begin{enumerate}[label=\textup{(M\arabic*)}]
\item every contribution to a shell crossing, every target-visible limit
defect, and every target-visible first failure belongs to
\(\mathsf M_{\Sigma,\mathfrak b}\);
\item deletion of the complementary forward words and backward covectors
does not change the shell-current identity, the admitted physical price, the
aggregate feeding--storage row, the recurrent circulation price, the
prospective source profile, or the set of independently realized singular
mechanisms;
\item the reduction is minimal: every retained nonfrontier vertex occurs on a
source--target constructor path or on a typed prefix that can extend through
the current frontier, while removing it deletes at least one algebraically
possible target-current word at that depth; every spurious frontier path is
removed at the first depth that resolves its failed extension;
\item if \(\mathsf M_{\Sigma,\mathfrak b}=\varnothing\), then the signature
is target-null.  If it is nonempty, each of its terminal source--target paths
has one of the five source ancestries, and the sum of their signed pairings is
the unit shell progress.
\end{enumerate}
Thus structural analysis of singular formation may be performed entirely on
\(\mathsf M_{\Sigma,\mathfrak b}\).
\end{theorem}

\begin{proof}
Start with the exact unit shell-current identity.  Every term has a source
root in \(\mathscr R^0_{\mathfrak b}\) and the target root \(D\Phi\).  Apply
one public constructor.  On its ordinary graph, the chain rule or weak
duality moves the constructor forward on the current or backward by its
adjoint on the covector, preserving the pairing.  On a nonordinary graph
value, source-preserving totalization inserts the failed operation as the
typed incidence edge.  Induction over constructor depth proves
\textup{(M1)}.

A deleted vertex has no complete source--target path.  Its ordinary pairing
with every target-cyclic descendant is zero, and every failed descent that
could change this conclusion would itself be a typed path and hence would
survive.  Removing all such vertices therefore changes neither the shell
identity nor any object generated from its target-visible terms.  Physical
price is admitted only through those terms, signed feeding and storage are
their glued images, circulation uses their return edges, and diagonal
realization uses their same-origin paths.  This proves \textup{(M2)}.

Finite graph pruning retains exactly the vertices on source--target paths or
on their well-typed not-yet-instantiated constructor prefixes.  Exhaustive enumeration
resolves every fixed frontier instance, proving \textup{(M3)}.  Compact
diagonal closure preserves the realized paths and all typed boundary values.
Empty core annihilates every source pairing,
so prospective accountability forbids shell crossing.  On a nonempty core,
group the retained root paths by their inherited source labels and recover
the five-source shell identity, proving \textup{(M4)}.
\end{proof}

\begin{theorem}[Commutation of the continuum compiler with minimal
target realization]
\label{thm:minimal-realization-compiler-commutation}
Let \(\mathsf P_{\mathsf M}\) denote restriction to
\(\mathsf M_{\Sigma,\mathfrak b}\), and let
\(\lvert\mathsf M_{\Sigma,\mathfrak b}\rvert\) denote the closed cylinder
state support on which all retained incidence paths are realized.  Target quotient, sourcewise completion,
generated production, signed carrier gluing, joint-kernel restriction,
first-return formation, circulation duality, circulation elimination, and
finite target branch-and-bound commute with \(\mathsf P_{\mathsf M}\).
In particular,
\begin{align}
 \delta^{\rm cyc}_{\Sigma,\mathfrak b}
 &=\delta^{\rm cyc}_{\Sigma,\mathfrak b}
       \big|_{\mathsf M_{\Sigma,\mathfrak b}},
 \label{eq:minimal-realization-cycle-price}\\
 X_\infty
 &=X_\infty\cap
   \lvert\mathsf M_{\Sigma,\mathfrak b}\rvert,
 \label{eq:minimal-realization-singular-core}\\
 \mathscr U_n
 &=\mathsf P_{\mathsf M_n}\mathscr U_n,
 \label{eq:minimal-realization-finite-cover}
\end{align}
where \(\mathsf M_n\) is the finite incidence core at the coordinates used at
depth \(n\).  Consequently the finite evaluator never generates a moment,
multiplier, defect, tail, or graph coordinate outside a complete
source--target path.
\end{theorem}

\begin{proof}
Each listed operation is generated functorially from pairings and public
constructor edges.  By \cref{thm:target-incidence-compression}, its
target-visible inputs and every failure of its descent lie in the incidence
core.  Restricting before or after the operation therefore gives the same
ordinary values and the same typed boundary values.  The physical objective
and target-return normalization are unchanged, proving
\eqref{eq:minimal-realization-cycle-price}.  Return, origin, and realization
edges are among the constructor paths, giving
\eqref{eq:minimal-realization-singular-core}.  At finite depth, graph pruning
is a finite preliminary step and the remaining moment rows are identical,
which proves \eqref{eq:minimal-realization-finite-cover}.
\end{proof}

\begin{corollary}[Finite incidence fast track]
\label{cor:finite-incidence-fast-track}
At cylinder depth \(n\), the minimal target-active realization is computed
before any ordered moment problem by two finite graph searches: forward from
the five source roots and backward from \(D\Phi\).  Their intersection is
the conservative incidence core
\(\mathsf M^{n,n}_{\Sigma,\mathfrak b}\), including its finitely many typed
constructor-frontier vertices.  If the finite incidence graph has
\(V_n\) vertices and \(E_n\) constructor or pairing edges, this pruning uses
\(O(V_n+E_n)\) graph operations.  The depth-
\(n\) moment program then contains only coordinates and localizing rows
indexed by the intersection.  Empty intersection closes the finite target
cell by the zero-pairing identities attached to the deleted edges; nonempty
intersection supplies the exact list of currently generated target-aligned
rows to price or refine, together with the frontier schemas that must be
expanded before any further deletion.
\end{corollary}

\begin{proof}
A vertex lies on a source--target path exactly when it is both forward
reachable from a source root and backward reachable from the target root.
Breadth-first or depth-first search computes these two sets in linear graph
time.  The commutation theorem permits deletion before production and moment
generation.  Frontier vertices prevent deletion of a prefix that a later
well-typed constructor can connect.  The stored zero or nonordinary pairing
value on each resolved boundary edge supplies the stated local identity or
typed successor.
\end{proof}

\begin{figure}[tbp]
\centering
\resizebox{.98\linewidth}{!}{%
\begin{tikzpicture}[fw diagram,node distance=8mm and 11mm]
 \node[fw provenance,text width=3.0cm] (sources)
   {five source roots\\forward ancestry};
 \node[fw use,text width=3.0cm,right=of sources] (forward)
   {public constructor graph\\\(\mathscr R^n_{\mathfrak b}\)};
 \node[fw provenance,text width=3.0cm,below=of sources] (target)
   {singular shell\\backward adjoint closure};
 \node[fw use,text width=3.0cm,right=of target] (backward)
   {target-cyclic module\\\(\mathscr C^n_{\mathfrak b}\)};
 \node[fw provenance,text width=3.1cm,right=14mm of forward,
       yshift=-7mm] (intersection)
   {source--target incidence core\\
    \(\mathsf M^n_{\Sigma,\mathfrak b}\)};
 \node[fw use,text width=3.0cm,right=of intersection] (compiler)
   {production, carrier,\\circulation, refinement};
 \node[fw terminal,text width=2.8cm,above right=7mm and 8mm of compiler] (empty)
   {empty incidence/core\\target-null};
 \node[fw use,text width=2.9cm,below right=7mm and 8mm of compiler] (path)
   {nested retained prefixes\\join and reproject};
 \draw[fw arrow] (sources)--(forward);
 \draw[fw arrow] (target)--(backward);
 \draw[fw arrow] (forward)--(intersection);
 \draw[fw arrow] (backward)--(intersection);
 \draw[fw arrow] (intersection)--(compiler);
 \draw[fw arrow] (compiler)--(empty);
 \draw[fw arrow] (compiler)--(path);
\end{tikzpicture}
}
\caption{Bidirectional target-incidence fast track.  Forward closure records
what the equation can generate from the five source occurrences.  Backward
closure records what the singular shell can detect through adjoint and
total-graph descendants.  Only their incidence core is passed to physical
pricing, recurrent circulation, and finite refinement.  A nested retained
prefix is a live joint-saturation state, not a terminal.}
\label{fig:target-incidence-fast-track}
\end{figure}

\begin{definition}[Covariant production valuation]
\label{def:covariant-production-valuation}
Let \(r\) be a generated record scale.  For every carrier term \(K\),
production \(P\), and target current \(Q\), the compiler obtains their scale
weights by substituting the public scaling chart into the displayed weak
identity:
\[
 K=r^{\kappa_K}\widetilde K,
 \qquad P=r^{\kappa_P}\widetilde P,
 \qquad Q=r^{\kappa_Q}\widetilde Q.
\]
These are valuations of equation terms, not analytic estimates.  Terms of
equal weight are first combined with their signs.  If the resulting carrier
law has critical weight, it is used directly.  If a positive power of \(r\)
would erase normalized production as \(r\downarrow0\), the scale velocity
\(a=\dot r/r\) is retained and
\cref{prop:covariant-modulation-conjugation} is applied.  Exact scale
coboundaries are absorbed into the covariant carrier; nonexact scale
cohomology remains an explicit \(\Caus\)- or \(\Lift\)-current.  A failure of
the chart, valuation, conjugation, or coboundary identity is the corresponding
total-graph successor.  In particular, the compiler never infers
\(\widetilde P=0\) from \(r^{\kappa_P}\widetilde P\to0\).
\end{definition}

\begin{definition}[Defect-polarization compiler]
\label{def:defect-polarization-compiler}
Suppose a public \(d\)-linear operation \(B\) has total weak graph coordinate
\[
 B(u_\nu,\ldots,u_\nu)\rightharpoonup
 B(u,\ldots,u)+\mathcal R_B.
\]
The compiler reconstructs the symmetrized multilinear defect from its
rational diagonal values by the polarization identity.  When an ordered
nonsymmetric component is used by the weak equation, the rational grammar
also adjoins all mixed-input graph coordinates; these determine that
component separately.  It then applies the following exclusive rules.
\begin{enumerate}[label=\textup{(D\arabic*)}]
\item A positive combination is admitted only when positivity follows from
the public positive cone, an exact square, or lower semicontinuity of an
already positive carrier form.  It is added to
\(\mathscr P_{\mathfrak b}\).
\item A signed defect is retained as a signed current and glued across all
compatible faces before its Jordan decomposition.
\item Failure of positivity, polarization, gluing, localization, or descent
is retained as the first nonordinary graph value with the ancestry of \(B\).
\end{enumerate}
Thus oscillation and concentration are generated structural coordinates of
the weak graph.  They cannot be discarded as errors and cannot be promoted
to physical expenditure without a carrier identity.
\end{definition}

\begin{proposition}[Functorial constraint and boundary descent]
\label{prop:target-cyclic-functorial-descent}
The target-cyclic construction commutes with every ordinary constraint,
gauge, chart, trace, and boundary quotient generated by the public
presentation.  More precisely, if \(q\) is such an ordinary quotient, then
\[
 q^*\mathscr C^n_{q\mathfrak b}
 \subseteq\mathscr C^n_{\mathfrak b},
 \qquad
 q(X^n_{\mathfrak b})\subseteq X^n_{q\mathfrak b}
\]
on its ordinary graph domain.  Equality holds only on coordinates for which
the generated lifting graph is ordinary and onto.  A missing reverse
inclusion is therefore recorded by its cokernel coordinate.  Outside the
ordinary domain, the least failed commutation, tail, trace, chart, or gauge
coordinate is routed as a target-null value, a physically charged value, or
the next target-cyclic successor.  No surjectivity, closed range, compact
trace, or tail decay is assumed.
\end{proposition}

\begin{proof}
The assertion holds at depth zero because the target shell is pulled back
through the same quotient.  If it holds at depth \(n\), naturality of
pullback gives the claim for adjoint transport and Lie derivatives; the
multilinear polarization identity gives it for weak defects; signed gluing
gives it for interfaces; and the public covariance identity gives it for
record-scale charts.  These are exactly the generators in
\cref{def:target-cyclic-covector-module}.  Induction proves the ordinary
inclusions.  An ordinary onto lifting graph supplies the reverse inclusions.
Totalization supplies a value when naturality or lifting fails, and the
source-preserving graph closure retains the failed operation and the target
covector that detected it.  The three stated routes are then the exhaustive
source-accountability alternatives.
\end{proof}

\begin{theorem}[Target-cyclic joint-kernel descent]
\label{thm:target-cyclic-joint-kernel-descent}
Let \(\mathfrak b\) be a same-origin target-active branch obtained from a
legal orbit with finite physical ledger.  Run
\cref{def:target-cyclic-covector-module,def:generated-production-joint-kernel,def:covariant-production-valuation,def:defect-polarization-compiler}
in the fixed enumeration of rational cylinders and public graph operations.
Exactly one of the following occurs, with priority given to the first finite
event in that enumeration.
\begin{enumerate}[label=\textup{(K\arabic*)}]
\item A generated covariant carrier law charges a disjoint cofinal family of
shell crossings by a nonnegative physical measure whose required total mass
diverges.  This contradicts the finite physical ledger.
\item On the retained same-origin path support in
\(X^\infty_{\mathfrak b}\), every source current selected by the generated
target-cyclic covectors is target-null:
\[
 \langle\lambda,J^c\rangle=0
 \quad
 (\lambda\in\mathscr C^\infty_{\mathfrak b}\text{ along that support},
 c\in\{\Geom,\Caus,\Abs,\Lift,\Bdry\}).
\]
Hence that retained path cannot cross the selected shell.  No assertion is
made about the rest of \(X^\infty_{\mathfrak b}\).
\item A production, defect, commutator, scale-cohomology, constraint,
boundary, chart, or total-graph coordinate is nonzero.  Its least such value,
together with its zero, signed dyadic, or graph-boundary polarity, detecting
covector, and inherited origin, is the next proper target-active successor;
the child is first joined to the complete parent transcript and reprojected.
\item All finite coordinates of the cumulative joint record are resolved,
compatible, and ordinary, and joint saturation is a fixed point.  The fixed
record retains a nonzero unnormalized shell cocycle and complete inherited
five-channel ancestry.  Its totalized realization row is ordinary and gives
a same-origin path crossing every cofinal target shell.  Thus it is a
joint-saturated prospective source atom in the sense of
\cref{def:joint-saturated-structural-atom}, with its original origin polarity,
not an inverse-limit or
resolved-prefix output.
\end{enumerate}
No compactness, coercivity, Liouville property, closed-range statement, or
equality classification occurs among the hypotheses.
\end{theorem}

\begin{proof}
Contact gives unit shell increment across each normalized shell.  By
prospective accountability, each completed crossing has a finite positive
source component or a positive boundary owner.  Starting from that shell covector, differentiate the pairing
through the public equation.  The chain rule gives the adjoint transports and
Lie derivatives in \textup{(C1)}; localization, quotient descent, weak
passage, and carrier completion give \textup{(C2)}--\textup{(C4)}.  Hence
every term exposed by any finite iteration is either an exact signed current,
a generated nonnegative production, or a total-graph value.  The signed
terms are glued before their positive parts are formed, and
\cref{def:covariant-production-valuation} prevents collapsing scales from
erasing a production.  This proves that the enumeration omits no descendant
of the initial target pairing.

If a carrier domination charges the cofinal disjoint crossings with divergent
total mass, \textup{(K1)} holds.  Restrict next to the actual same-origin path
support.  If every generated target pairing vanishes there, the shell
derivative vanishes along that path, so the fundamental theorem of calculus
forbids its shell crossing and \textup{(K2)} holds.  No global annihilation
of the joint kernel is used.  Otherwise choose the least nonzero or
nonordinary coordinate of the cumulative joint record and apply its generated polarity split.
Source-preserving totalization,
\cref{prop:target-cyclic-functorial-descent,thm:target-local-stratification-exhaustiveness}
give \textup{(K3)} without an analytic exclusion assumption.  Joining the
child to the whole preceding transcript before reprojection retains every
mixed operation through which individually target-null coordinates could
become jointly target-aligned.

It remains to consider the case in which every finite coordinate is resolved
and compatible in its induced graph.  Form the cumulative joint records of
\cref{def:cumulative-joint-continuum-record} before taking a limit.  By
\cref{thm:continuum-composition-no-creation}, an unprocessed target component
of their graph-complete limit already has a least nonzero finite joint
increment; that prefix is a successor of type \textup{(K3)}.  If the full
joint target projection vanishes, the joint limit belongs to \textup{(K2)}.
The only remaining case, after every increment has acquired its complete
five-source ancestry, is a fixed point of the
joint, graph, source, carrier, and target-cyclic closures.  Its unnormalized
shell cocycle cannot vanish, because each completed crossing contributes one
before any return normalization.  Totalize its same-origin realization row.
A nonordinary value is again a finite successor of type \textup{(K3)}; an
ordinary contacting value is precisely the structural atom in
\textup{(K4)}.  Hence neither increasing rank nor inverse-limit existence is
a terminal argument.  The fixed priority rule makes the four outputs
disjoint.
\end{proof}

\begin{corollary}[Generated mobility-annihilator test]
\label{cor:generated-mobility-annihilator-test}
If the restriction of the target-cyclic module to the retained same-origin
path support annihilates all five induced currents there, then that path
cannot contact the target.  Otherwise the least nonzero restricted pairing,
together with its generated polarity cell, is the initial covector of the
next successor.  Thus the covector
needed to continue the equality analysis is generated by the equation and
the target; it is never an external multiplier or a user-supplied rigidity
certificate.
\end{corollary}

\subsection{Target-relative resaturation of zero-price dynamics}
\label{subsec:zero-price-resaturation}

A zero-price equality law is not a terminal object.  It is a new public
dynamical stratum on which prospective accountability is run again.  This is
the continuum counterpart of forming the saturated closure before taking the
residual in Part~I.

\begin{definition}[Zero-price resaturation successor]
\label{def:zero-price-resaturation-successor}
Let
\(\mathfrak b_n=(A_n,m_n,\Gamma_n,\eta_n,\mathbf R)\) be a reachable active
branch with ordered ideal \(I_n\), dual-carrier module \(\mathscr W_n\), and
zero-price recurrent law \(L_n\).  Retain its cumulative joint record
\(\mathcal J_n\) and unnormalized shell cocycle.  Its
\emph{resaturation successor} is
constructed by the following generated operations.
\begin{enumerate}[label=\textup{(R\arabic*)}]
\item Starting from the selected shell covector, generate
\(\mathscr C^\infty_{\mathfrak b_n}\) and
\(\mathscr P_{\mathfrak b_n}\) by
\cref{def:target-cyclic-covector-module,def:generated-production-joint-kernel}.
Adjoin every production annihilated by \(L_n\), and form the closed ordered
kernel
\begin{equation}
 \label{eq:zero-price-kernel-restriction}
 X_{n+1}:=X_n\cap\bigcap_{P\in\mathscr P_{\mathfrak b_n}}\{P=0\}.
\end{equation}
\item Restrict the public weak equation, its constraints, carrier, charts,
and total graphs to \(X_{n+1}\).  Recompute the five source projections there;
no old source tag or origin coordinate is discarded.
\item Apply the prospective shell identity on the restricted pre-contact
path.  If contact remains, select a persistent source support and target
covector with the normalized floor \(1/5\).
\item Apply the covariant production valuation and defect-polarization
compiler.  Adjoin every resulting adjoint, quotient, localization,
signed-gluing, covariance, commutator, defect, and square-completion relation
to \(I_n\), and adjoin its detecting covector to \(\mathscr W_n\).  Exact
scale coboundaries are absorbed and nonexact scale cohomology is retained as
a source current.
\item Totalize every operation used above.  Its first nonordinary value is a
typed successor.  If every value is ordinary, the restriction gives the next
zero-price branch \(\mathfrak b_{n+1}\).
\item Apply the generated polarity split to the next unprocessed coordinate
on the retained same-origin target stratum.  Prune a cell only by its finite
local ordered identity; every positive-moment cell appends its polarity
relation to the ideal.  A normalized return law may select its support; the
unnormalized target current is unchanged before any physical charge is made.
\item Expand the complete derivation of the selected child, join it with the
expanded records in \(\mathcal J_0,\ldots,\mathcal J_n\), generate every
mixed operation and polarization, saturate, and reapply the five-channel
target projection.  Retain every physical owner and normalization
denominator.  A
least unprocessed target increment is the next source-resolved successor even
when each of its individual input coordinates was target-null.
\end{enumerate}
The successor is \emph{proper} when either its ordered ideal or its observable
module strictly increases after reduction by the preceding relations.  Every
proper successor has strictly larger resolved-prefix rank.  An infinite rank
chain is evaluated by its cumulative joint record under
\cref{thm:continuum-composition-no-creation}; rank growth is not a terminal
argument.
\end{definition}

\begin{proposition}[Monotone resaturation]
\label{prop:monotone-zero-price-resaturation}
Along a zero-price resaturation chain,
\begin{align}
 \label{eq:monotone-zero-price-resaturation}
 I_0&\subseteq I_1\subseteq\cdots,
 &\mathscr W_0&\subseteq\mathscr W_1\subseteq\cdots,
 &X_0&\supseteq X_1\supseteq\cdots,\notag\\
 \mathcal A_0&\subseteq\mathcal A_1\subseteq\cdots,
 &\rho(\mathfrak b_0)&<\rho(\mathfrak b_1)<\cdots.
\end{align}
Every retained branch has the same origin tag, and every completed shell
prefix retains \(A_{\Phi,N}=N\).
At a limit stage use
\begin{equation}
 \label{eq:limit-zero-price-resaturation}
 I_\infty:=\overline{\bigcup_n I_n}^{\rm ord},
 \qquad
 \mathscr W_\infty:=\overline{\bigcup_n\mathscr W_n}^{\rm cyl},
 \qquad
 X_\infty:=\bigcap_nX_n,
 \qquad
 \mathcal J_\infty
 :=\operatorname{Sat}_{\mathsf G_\Sigma}
   \overline{\bigcup_n\mathcal J_n}^{\rm gc}.
\end{equation}
These closures create no source: every limit relation retains the ancestry of
the finite words that generate it, and every target component of
\(\mathcal J_\infty\) is the limit of finite joint-prefix components.
\end{proposition}

\begin{proof}
Kernel restriction only adds zero relations, while dual-carrier compilation
only adds words and their consequences.  The polarity step appends the next
resolved coordinate, and the joint step enlarges the cumulative sigma
algebra, proving the inclusions and the rank inequality.  The shell selector
is applied to the same pre-contact path, so the origin tag is unchanged and
each completed shell contributes one to the unnormalized cocycle.  The closures
in \eqref{eq:limit-zero-price-resaturation} are the inverse-limit closures of
the countable rational cylinder grammar.  Source-preserving completion
and \cref{thm:continuum-composition-no-creation} therefore give the final
assertion.
\end{proof}

\begin{theorem}[Proper-successor monotonicity]
\label{thm:proper-successor-monotonicity}
For a proper same-origin successor define
\[
 \Xi(\mathfrak b)
 :=\bigl(\lvert A(\mathfrak b)\rvert,\rho(\mathfrak b),
          n_{\rm graph}(\mathfrak b),n_{\rm polar}(\mathfrak b),
          n_{\rm carrier}(\mathfrak b)\bigr)
\]
in lexicographic order, where the last three coordinates count the processed
ordered graph, polarity, and carrier--polar prefixes.  Every proper
successor strictly increases \(\Xi\).  Support growth is monotone and has at
most the five source letters; on fixed support, resolved-prefix rank
increases, and processed owner and graph indices never decrease.
\end{theorem}

\begin{proof}
Source ancestry is joined by union, so a newly visible source letter enlarges
the first coordinate and no letter can later be removed.  Otherwise a
polarity successor appends its least unresolved coordinate, a graph successor
appends its least failed relation, and a carrier--polar successor appends a
new separator or descendant.  The cumulative join retains all previous
prefixes, so none of these indices can decrease.  The definition of
properness and \cref{prop:monotone-zero-price-resaturation} then give the
strict lexicographic increase.  This is a canonical monotonicity measure,
not a finite-termination assertion: an infinite chain is evaluated through
its cumulative joint record.
\end{proof}

\subsection{Recurrent passivity and support transfer}
\label{subsec:terminal-passivity-support}

The resolved-prefix rank detects a first failed operation, but rank growth
alone does not evaluate a current which repeatedly reappears after
renormalization.  The carrier must be followed through the same-origin
successor chain.  The following construction performs that step before a
limit path is declared to be a singular mechanism.

\begin{definition}[Target-local marked carrier complex]
\label{def:target-local-marked-carrier-complex}
Fix a same-origin active branch \(\mathfrak b\) and its compact retained path
space.  A marked carrier cell is a tuple
\[
 C=(w,w^+,\omega,K^-,K^+,P,H,E),
\]
where \(w\mathcal R_\eta w^+\), \(\omega\) is the retained successor witness,
\(K^\pm\) are the endpoint values of one generated carrier, \(P\ge0\) is the
sum of the generated nonnegative productions in that cell, \(H\) is the
complete signed interior, chart, scale, constraint, and interface current
after compatible faces have been glued, and \(E\) is the true exterior or
physical-input current.  Every coordinate is restricted to
\(X_{\mathfrak b}\) and carries the detecting shell covector.  Its total graph
contains the exact relation
\begin{equation}
 \label{eq:target-local-marked-carrier}
 P=K^- -K^+ +H+E.
\end{equation}
Let \(\mathscr B_{\mathfrak b}\) be the closed linear space generated by
shift coboundaries \(F\circ S-F\), paired internal-face currents, exact chart
and gauge currents, and currents annihilated by the selected target quotient.
The \emph{target-local feeding class} is
\begin{equation}
 \label{eq:target-local-feeding-cohomology}
 [H]_{\mathfrak b}
 \in
 \mathscr H_{\mathfrak b}
 :=\overline{\operatorname{span}\{H\}}/\mathscr B_{\mathfrak b}.
\end{equation}
Exterior input is not put in this quotient.  Its ordinary carrier-dominated
part is charged by \eqref{eq:target-aligned-charge-rule}; every other value is
its \(\Bdry\)- or carrier-domination successor.
\end{definition}

\begin{proposition}[Stationary carrier reduction]
\label{prop:stationary-carrier-reduction}
Let \(\varpi\) be a shift-invariant probability law on a concatenated marked
carrier path in one active branch.  On the ordinary finite-integral value,
\begin{equation}
 \label{eq:stationary-carrier-cohomology}
 \int P\,\dd\varpi
 =\int [H]_{\mathfrak b}\,\dd\varpi
  +\int E\,\dd\varpi.
\end{equation}
Here the first integral on the right means evaluation of the cohomology class
by the stationary law.  If an endpoint, concatenation, integrability,
quotient, or exterior coordinate is nonordinary, its least total-graph value
is the corresponding proper successor.
\end{proposition}

\begin{proof}
Sum \eqref{eq:target-local-marked-carrier} over a finite consecutive block.
Concatenation cancels every common endpoint carrier.  Divide by the block
length and pass to the empirical path law.  Bounded endpoint storage vanishes,
as does every shift coboundary.  Paired faces and exact chart or gauge currents
already vanish in the quotient.  This gives
\eqref{eq:stationary-carrier-cohomology}.  Each operation used in this
calculation is totalized, so its first nonordinary value gives the stated
successor instead of an identity.
The division constructs a recurrent support law only.  The unnormalized sums
of \(P\), \(H\), \(E\), and all five source currents remain coordinates of
the same structural record and are not consequences of
\eqref{eq:stationary-carrier-cohomology}.
\end{proof}

\begin{theorem}[Target-aligned passivity--support alternative]
\label{thm:target-aligned-passivity-support}
Let an infinite zero-price successor path retain the active signature
\(\mathfrak b\), and form its stationary marked carrier law.  Exactly one of
the following events occurs in the fixed total-graph order.
\begin{enumerate}[label=\textup{(V\arabic*)}]
\item A true exterior or input current is admitted by the target-local
carrier graph.  Its restrictions to the diagonal original events are charged
to the public physical measure; divergent required mass contradicts the
finite ledger, and finite mass enters the zero-input tail successor.
\item A marked-carrier, gluing, quotient, support-transfer, or realization
coordinate is nonordinary and gives its typed proper successor.
\item The ordinary feeding class has nonzero positive stationary value.  The
normalized positive measure
\begin{equation}
 \label{eq:positive-feeding-support-law}
 \dd\lambda_H
 :=\frac{([H]_{\mathfrak b})_+\,\dd\varpi}
         {\int([H]_{\mathfrak b})_+\,\dd\varpi}
\end{equation}
has a target-visible support projection.  Its first atomic, separated,
nested, diffuse, tail, rough, gauge, constraint, or recurrent projection is
the next same-origin target-active return record, retaining the source word
of the current that finances \(P\).  It is a proper successor when it resolves
a new coordinate; otherwise its positive-measure signed-dyadic cell defines
the induced recurrent-current return map.
\item The ordinary exterior value and feeding cohomology vanish.  Then
\begin{equation}
 \label{eq:terminal-passivity-zero-production}
 \int P\,\dd\varpi=0,
 \qquad P=0\quad\varpi\text{-almost surely},
\end{equation}
and the path law is restricted to the target-cyclic joint kernel generated by
these zero productions.
\end{enumerate}
Thus lack of passivity is itself a target-aligned structural current.  It
cannot be retained as an unevaluated equality law.
\end{theorem}

\begin{proof}
Apply \cref{prop:stationary-carrier-reduction}.  The exterior coordinate is
processed first by the local charging rule, giving \textup{(V1)} or its
successor.  Any nonordinary operation gives \textup{(V2)}.  On the remaining
ordinary zero-exterior cell, nonnegativity of \(P\) gives
\[
 \int[H]_{\mathfrak b}\,\dd\varpi=\int P\,\dd\varpi\ge0.
\]
If the value is positive, \eqref{eq:positive-feeding-support-law} is a
probability measure on the already compact activity observer space.  Apply
\cref{thm:target-complete-concentration-exhaustion}.  If every support
projection were target-null, the target-local functional defining
\([H]_{\mathfrak b}\) would integrate to zero.  Hence one projection is
target-visible and supplies \textup{(V3)} with its same-origin witness.
On an equivalent support cell, Poincare recurrence for the least positive
signed-dyadic set supplies the induced return map, so equivalence does not
discard the current.
If the value is zero, \eqref{eq:stationary-carrier-cohomology} gives
\(\int P\,\dd\varpi=0\).  Since \(P\ge0\), it vanishes almost surely, proving
\textup{(V4)}.
\end{proof}

\begin{theorem}[Target-visible fidelity defects generate carrier successors]
\label{thm:fidelity-defect-support-successor}
Let \(F\) be a public equation, constitutive, nonlinear, constraint, trace,
chart, or realization operation, and let \(\mathcal D_F\) be a nonzero value
of its total-graph defect on an active branch \(\mathfrak b\).  For a
target-cyclic covector \(\lambda\), exactly one of the following occurs.
\begin{enumerate}[label=\textup{(D\arabic*)}]
\item \(\langle\lambda,\mathcal D_F\rangle=0\); this occurrence of the defect
lies in the target kernel.
\item The normalized variation of
\(\langle\lambda,\mathcal D_F\rangle\), localized by the rational cylinder
core, produces an atomic, separated, nested, diffuse, tail, or graph-boundary
support record with the same origin and the ancestry of \(F\).  A support
record inequivalent to the parent is a proper target-active successor.
\item Every positive support record is equivalent to the parent in the
resolved coordinates.  Its event-count law is recurrent.  The defect pairing
is then retained, with its sign, as either a nonzero component of
\(\Pi_{\Sigma,\mathfrak b}[H]_{\mathfrak b}\) or an exact carrier-storage
transfer in \(\mathscr B_{\mathfrak b}\).  Both values are evaluated by
\cref{thm:target-local-passivity-fast-track}; an exact transfer is not
identified with a target-null current.
\end{enumerate}
Consequently no target-visible fidelity defect can disappear at a recurrent
limit: it is a proper support successor, a recurrent feeding current, or a
recurrent storage-transfer current.
\end{theorem}

\begin{proof}
Apply the scalar distribution
\(T=\langle\lambda,\mathcal D_F\rangle\) to the enumerated rational local test
core.  Either every value is zero, which is \textup{(D1)}, or there is a
least test with nonzero value.  The total variation of its rational
localizations is a nonzero positive finite measure on each finite cylinder;
if its variation is infinite, use the normalized-variation boundary
coordinate.  Normalize and apply compact activity-profile exhaustion.
Source-preserving completion retains the ancestry of \(F\), and the
same-origin selector retains the original event indices.  This gives
\textup{(D2)}.

If one support record is inequivalent to the parent, its witness is a proper
successor.  Otherwise take empirical laws of the equivalent support records.
Successor recurrence gives a shift-invariant law.  Pairing the public weak
equation with the transported localizations of \(\lambda\) places \(T\) in
the complete signed current \(H\); paired internal occurrences cancel before
the quotient.  If its class is nonzero, its target projection is retained as
a feeding coordinate.  If its class is zero, the definition of the closed
coboundary space supplies bounded cylinder storage functions whose
coboundaries converge to the pairing.  Their finite-block sums are endpoint
transfers.  The support law retains those endpoints and therefore routes the
value as a carrier-storage coordinate.  Only annihilation by every
target-cyclic covector gives \textup{(D1)}.  This proves \textup{(D3)} without
the false implication that nonzero variation has nonzero stationary mean.
\end{proof}

\begin{theorem}[Fidelity-defect successor]
\label{thm:fidelity-defect-successor}
Let a public equation, chain, constraint, chart, trace, covariance, carrier,
or continuation relation have a target-visible defect on an active branch.
Its first failed finite total-graph coordinate determines a nonzero
target-cyclic pairing, a source owner, and an actual finite same-origin
ancestor occurrence.  Appending that coordinate and its zero, signed-dyadic,
or graph-boundary polarity cell gives a proper successor.  If every
target-cyclic pairing of the defect is zero, the defect is deleted only in
\(\Null_\Sigma\).
\end{theorem}

\begin{proof}
Enumerate the graph residuals and choose the least nonordinary value or the
least nonzero target pairing.  The total graph retains both its input record
and its finite approximating sequence, while
\cref{thm:fidelity-defect-support-successor} supplies its source-preserving
support or carrier route.  On a positive bounded graph detector,
\cref{thm:actual-support-return} recovers an actual finite same-origin
ancestor occurrence.  The selected graph coordinate was unresolved in the
parent, so adjoining it strictly enlarges the ordered ideal or the
observable graph prefix and is proper.  If no target-cyclic pairing is
nonzero, the definition of the target quotient gives precisely
\(\Null_\Sigma\); it does not erase a target-visible residual.
\end{proof}

\begin{definition}[Target-active recurrent component]
\label{def:target-active-recurrent-component}
Let \(\Omega_{\mathfrak b}\) be the compact path space of same-origin
successor chains and let \(S\) be the left shift.  A \emph{target-active
recurrent component} is the support of an ergodic \(S\)-invariant probability
law \(\varpi\) obtained from a contact-generated chain and satisfying the
retained shell floor.  Every generated scalar coordinate has a fixed zero,
signed-dyadic, or graph-boundary value on the corresponding measurable cell.
A target-visible value is retained even when its support successor is
equivalent in the previously resolved coordinates.  In that case it is a
recurrent current coordinate and remains an active operation of the compiler.

Such a component exists for every infinite chain: take a weak limit of its
empirical shift laws and then an ergodic component with nonzero target mass.
The empirical law selects recurrent support only.  It is stored jointly with
the cumulative record and unnormalized cocycle of
\cref{def:cumulative-joint-continuum-record}; target projection is recomputed
on that join.  This construction asserts recurrence of structural records
only.  The source record retains the same ordinary stopped history from which
it was generated; its target status is determined only after rejoining that
complete history and evaluating the reachable contact cell.
\end{definition}

\begin{theorem}[Target-relative recurrent stratification]
\label{thm:target-relative-terminal-stratification}
Every infinite same-origin zero-price resaturation route has exactly one of
the following constructive outputs.
\begin{enumerate}[label=\textup{(T\arabic*)}]
\item a finite target-null, continuation, class-incompatibility, or divergent
physical-charge relation;
\item a proper target-active successor produced by a compactness mode, a
fidelity defect, or a previously unprocessed feeding or storage current;
\item a shift-invariant equality law satisfying
\eqref{eq:terminal-passivity-zero-production}, followed by the
target-cyclic mobility-annihilator test;
\item a target-active recurrent component.  On it, every nonzero aggregate
feeding or storage current is selected by a signed-dyadic return set and
becomes the next recurrent-current successor.  If all such currents vanish,
the component lies in the generated zero-production joint kernel.  A cofinal
  chain of recurrent-current successors is accumulated in
  \(\mathcal J_\infty\) and reprojected.  It is target-null, has a nonzero
  unprocessed finite joint or polar increment, or is a primal--polar fixed
  joint-saturated structural atom.  Its hypothetical ordinary realization is
  a source-profile witness; only its independent ordinary realization is an
  explicit singular mechanism.
\end{enumerate}
A nonzero public-equation, constitutive, constraint, boundary, fidelity,
support-transfer, feeding, or storage coordinate is therefore an operation
on \textup{(T4)}, never a terminal name.  In
\textup{(T3)}, a target-null joint kernel closes; a nonzero target-cyclic
pairing is the next successor.  Consequently an unending classification
chain is replaced by a recurrent current system with a definite structural
evolution.
\end{theorem}

\begin{proof}
Apply compact activity-profile exhaustion to the diagonal same-origin path.
An atomic, separated, nested, diffuse, tail, rough, gauge, constraint, or
equality value is a coordinate of its single activity law.  Process its first
target-visible value.  The carrier values are then exhausted by
\cref{thm:target-aligned-passivity-support}.  Cases \textup{(V1)}--
\textup{(V2)} give \textup{(T1)}--\textup{(T2)}; case \textup{(V3)} gives a
proper support successor when it resolves a new coordinate and otherwise
gives the recurrent-current return system in \textup{(T4)}; case
\textup{(V4)} gives \textup{(T3)}.

Repeat on every proper successor.  Strict resolved-prefix growth and compact
diagonal saturation give a same-origin joint limit record.  By
\cref{thm:continuum-successor-recurrence} it carries an invariant law; select
an ergodic component of nonzero target mass.  Apply
\cref{thm:fidelity-defect-support-successor} to every retained public-equation
defect.  A target-visible defect returns to \textup{(T2)} or becomes a
recurrent feeding or storage coordinate; every remaining defect is
target-null.

For a nonzero recurrent coordinate choose the least signed-dyadic set
\(A=\{sF\in[2^{-k},2^{-k+1})\}\) of positive \(\varpi\)-measure.  Poincare
recurrence makes the first-return map \(S_A\) finite almost everywhere, and
the normalized restriction \(\varpi(A)^{-1}\varpi|_A\) is \(S_A\)-invariant.
This law selects the recurrent-current support in \textup{(T4)}.  Rejoin it
with the complete same-origin witnesses, source ancestry, and unnormalized
shell cocycle before projection.  If the join resolves a new coordinate it is
\textup{(T2)}.  If equivalent successors recur indefinitely, form their
cumulative joint record and apply
\cref{thm:continuum-composition-no-creation}: it is target-null, has a finite
joint-prefix successor carrying the least unprocessed increment, or has all
primal increments source-resolved.  In the last case apply
\cref{thm:recursive-carrier-polar-exhaustion}; a separator is the next
successor, while only simultaneous primal--polar fixedness gives a structural
atom.  Rejoining the last cell with its complete same-origin stopped history
gives the actual contact-cell test in \textup{(T4)}.  A nonordinary value is
another typed successor.  If no nonzero coordinate
remains, stationary carrier reduction gives \textup{(T3)}.
This proves exhaustiveness without asserting that an increasing refinement
chain has a sink.
\end{proof}

\begin{corollary}[Local criterion excluding a recurrent singular component]
\label{cor:terminal-atom-regularity-criterion}
Suppose the generated passivity row of every reachable active signature has
target-null zero-production kernel, and every recurrent feeding, fidelity,
storage, or support-transfer current is either charged by the public physical
carrier or has target-null return law.  Then no target-active recurrent
component contacts the continuation target.  Hence every finite-physical-
expenditure ordinary branch is regular and continuable.  These statements
are outputs of the local ordered cones, return maps, and total graphs of the
active signatures; they are not specialization data.
\end{corollary}

\begin{proof}
A contact branch produces an infinite zero-price route after the divergent-
price cases have been removed.  Apply
\cref{thm:target-relative-terminal-stratification}.  The first three outputs
close or continue.  In the recurrent output, every nonzero current is charged
or has target-null return law by hypothesis.  The remaining component lies in
the zero-production kernel, which is target-null by the generated local
ordered identity, contradicting its shell floor.
\end{proof}

\begin{theorem}[Target-local passivity fast track]
\label{thm:target-local-passivity-fast-track}
For a reachable active signature \(\mathfrak b\), after exterior input has
been processed, the full recurrent carrier analysis reduces to one aggregate
target-current row,
\begin{equation}
 \label{eq:passivity-fast-track-class}
 \mathfrak J_{\Sigma,\mathfrak b}
 :=\bigl(\Pi_{\Sigma,\mathfrak b}[H]_{\mathfrak b},
         \Pi_{\Sigma,\mathfrak b}H^{\rm stor}_{\mathfrak b}\bigr),
\end{equation}
where \(H^{\rm stor}_{\mathfrak b}\) is the endpoint-transfer coordinate of
the closed coboundary graph.  Its total ordered graph has the following
exhaustive evaluation.
\begin{enumerate}[label=\textup{(F\arabic*)}]
\item A finite identity in \(\mathcal K_{\mathfrak b}\) proves
\(\Pi_{\Sigma,\mathfrak b}[H]_{\mathfrak b}=0\).  Stationary carrier
reduction forces the generated production to vanish.  A nonzero storage
coordinate is then selected by its signed-dyadic return set; if storage also
vanishes, the target-cyclic joint kernel is evaluated on that same cell.
\item The class has a nonzero signed-dyadic value.  Its positive support law
\eqref{eq:positive-feeding-support-law}, source word, detecting covector, and
same-origin witness define a positive-density return-map record.  The compiler applies the
five-source normal form and prospective accountability directly to this
record; it does not separately estimate the individual pressure, transport,
chart, cutoff, or scale terms from which the class was assembled.
\item The projection, gluing, support, or realization graph is nonordinary.
Its first graph-boundary value, with the same data, is the next active record.
\end{enumerate}
Thus a specialization evaluates one aggregate target-aligned row.  Exact
storage and nonexact feeding are distinguished, but neither is discarded.
Atomic, separated, nested, diffuse, and tail modes occur only as support
coordinates of the selected return law and introduce no additional proof
stage.
\end{theorem}

\begin{proof}
Signed carrier gluing constructs \([H]_{\mathfrak b}\) and its storage graph
before a positive part is taken.  Apply the target-local factorization and
the generated polarity split to \eqref{eq:passivity-fast-track-class}.  The
zero cohomology cell is \textup{(F1)} by
\cref{prop:stationary-carrier-reduction,thm:target-aligned-passivity-support}.
On that cell the exact component is a block endpoint transfer and is split by
the same signed-dyadic return construction.  Positive and negative dyadic
feeding cells retain the signed current; reverse the orientation on a
negative cell and apply the positive support construction, giving
\textup{(F2)}.  Poincare recurrence supplies the induced return law.
Totalization gives \textup{(F3)}.  Compact activity exhaustion is applied to
the single support measure in \textup{(F1)}, \textup{(F2)}, or
\textup{(F3)}, which proves the final assertion.
\end{proof}

\begin{theorem}[Actual-support return]
\label{thm:actual-support-return}
Let a positive recurrent feeding or storage support selector be produced by
\cref{thm:target-local-passivity-fast-track}.  Then it yields an increasing
sequence of actual finite same-origin occurrence records, rather than only a
probability law.  Each selected return carries its original event index,
finite public weak-equation record, complete source ancestry, canonical
target amplitude, physical valuation, return time, and predecessor and
successor records on the same orbit.
\end{theorem}

\begin{proof}
The recurrent selector is supported on the graph-complete closure of the
presentation-generated finite occurrence records.  A positive signed-dyadic
cell contains a rational cylinder neighborhood of positive selector mass.
Density and
\cref{thm:finite-contextual-occurrence-amplitude-retention} supply an actual
finite same-origin occurrence in that neighborhood with retained positive
amplitude.  Poincare recurrence supplies infinitely many returns to the
positive cell.  Choose them inductively with strictly increasing original
event indices and apply the same finite-occurrence recovery at each return.
The original index, predecessor, successor, and return-time coordinates are
part of the recurrent closure graph, while multiplicity restoration reattaches
the unnormalised occurrence and physical ledgers before target projection.
\end{proof}

\begin{definition}[Target-aligned recurrent closure compiler]
\label{def:target-aligned-recurrent-closure-compiler}
For a reachable active signature \(\mathfrak b\), a compiler state is
\begin{equation}
 \label{eq:target-aligned-compiler-state}
 \mathfrak q=(X,I,\mathscr C,\mathscr P,\varpi,
              \mathfrak J_{\Sigma,\mathfrak b},
              \delta^{\rm cyc}_{\Sigma,\mathfrak b},\mathbf R,
              \mathcal J_{\le k},\Gamma_{\le k},A_{\Phi,\le k},
              \Pi_{\le k},(Q^c_{\Phi,\le k})_{c\in\Channels}),
\end{equation}
where \(X\) is the current same-origin target cell, \(I\) its ordered ideal,
\(\mathscr C\) the target-cyclic covector module, \(\mathscr P\) the generated
production family, \(\varpi\) its recurrent law when present,
\(\mathfrak J_{\Sigma,\mathfrak b}\) the aggregate feeding--storage row, and
\(\delta^{\rm cyc}_{\Sigma,\mathfrak b}\) its sharp recurrent circulation
price when the return graph has been formed, and \(\mathbf R\) the
original-event witnesses.  The final entries are the cumulative joint record,
occurrence multiplicity, unnormalized shell cocycle, cumulative physical
expenditure, and cumulative five-source currents from
\cref{def:cumulative-joint-continuum-record,def:target-current-circulation-problem}.
Before enqueueing a state,
join it with every preceding state on the same execution, saturate the joint
record, and then restrict all coordinates to the minimal target-active
realization \(\mathsf M_{\Sigma,\mathfrak b}\).  Initialize the queue with the
algebraically reachable shell state.  Process its least generated coordinate by the
following fixed transition order:
\begin{enumerate}[label=\textup{(Q\arabic*)}]
\item Dovetail the finite identities in the local cone and the dual
circulation-price hierarchy.  A verified
\(-1\)-identity removes the cell.  A carrier-domination identity records the
restriction of the public physical measure; a cofinal family with divergent
total charge removes the cell.  A positive recurrent circulation price
produces the dual carrier potential and removes the recurrent cell by
infinite return charge.
\item On the zero cell of
\(\Pi_{\Sigma,\mathfrak b}[H]_{\mathfrak b}\), adjoin every forced
nonnegative production to the zero-production joint kernel.  Apply
prospective accountability on the restricted same-origin path.  Vanishing of
all target-cyclic pairings removes the cell in \(\Null_\Sigma\); otherwise
the least nonzero pairing is appended to the queue.
\item On the zero-density circulation supplied by the primal hierarchy, select
a signed-dyadic feeding or storage cell of positive recurrent mass, form its
first-return map and normalized induced law, and append the sign,
dyadic level, source ancestry, detecting covector, and return-time graph.
The induced law selects support only.  Rejoin it with
\((\mathcal J_{\le k},\Gamma_{\le k},A_{\Phi,\le k},\Pi_{\le k},
(Q^c_{\Phi,\le k})_c)\) and impose the multiplicity-restoration identities
before target projection.  No zero density is copied into a cumulative
zero relation.  The
resulting joint state is the next queue element even when its previously
resolved scalar coordinates agree with those of its parent.
\item On a graph-boundary value, append the first failed operation, its total
graph, source ancestry, target covector, and same-origin support law.  Apply
the same transition order to that typed state.
\end{enumerate}
Only descendants generated from the selected shell covector are enqueued.
Every other coordinate is removed by the target quotient.  Compatible signed
currents are glued before \textup{(Q2)}--\textup{(Q3)}, so the queue contains
aggregate target currents rather than separate pressure, transport, chart,
cutoff, scale, and interface estimates.  At an infinite execution, the queue
forms \(\mathcal J_\infty\) and applies
\cref{thm:continuum-composition-no-creation}; it never promotes the compatible
inverse limit to an output before this joint reprojection.
\end{definition}

\begin{definition}[Target-cyclic structural spine]
\label{def:target-cyclic-structural-spine}
The \emph{target-cyclic structural spine} of an active signature
\(\mathfrak b\) is the least closed subgraph
\begin{equation}
 \label{eq:target-cyclic-structural-spine}
 \Spine_{\Sigma,\mathfrak b}
 \subseteq \operatorname{Sat}_{\mathsf G_\Sigma}
                  (q_\Sigma\mathscr J)
\end{equation}
that contains the five projected source currents paired with the selected
shell covector and is closed under target-cyclic adjoint transport, projected
total graphs, signed carrier gluing, zero-production restriction,
signed-dyadic first return, same-origin succession, and diagonal realization.
The physical valuation on the spine is the restriction of the public carrier
measure.  No coordinate with zero target projection belongs to the spine.

Equivalently, \(\Spine_{\Sigma,\mathfrak b}\) is generated by starting from
the unit shell-current identity and repeatedly adjoining only the next
coordinate selected by \textup{(Q2)}--\textup{(Q4)}.  It is therefore
enumerated from the public presentation, target shells, and physical carrier;
it is not an additional specialization object.
\end{definition}

\begin{proposition}[Incidence presentation of the target spine]
\label{prop:incidence-presents-target-spine}
Let \(\operatorname{pr}_{\mathscr R}\) forget the covector and incidence-edge
coordinates while retaining current ancestry, return, origin, and total-graph
values.  Then
\begin{equation}
 \label{eq:incidence-presents-target-spine}
 \operatorname{pr}_{\mathscr R}
 \mathsf M_{\Sigma,\mathfrak b}
 =\Spine_{\Sigma,\mathfrak b}.
\end{equation}
Thus the minimal realization is the bidirectional, finitely prunable
presentation of the target-cyclic structural spine.
\end{proposition}

\begin{proof}
Every retained incidence path starts at a projected source occurrence and is
closed under the target-cyclic adjoint, total-graph, gluing, kernel, return,
origin, and realization operations that define the spine.  Its current
projection is therefore contained in the spine.  Conversely, structural
induction over the least spine closure attaches to every spine word the
target-cyclic covector and constructor path that generated it.  This gives a
vertex of the incidence realization with that current projection.  The two
inclusions prove \eqref{eq:incidence-presents-target-spine}.
\end{proof}

\begin{proposition}[Spine restriction of the recurrent compiler]
\label{prop:spine-restriction-recurrent-compiler}
Every nonclosing state of the target-aligned recurrent closure compiler lies
in \(\Spine_{\Sigma,\mathfrak b}\).  Conversely, every finite spine word is a
finite compiler state with the same ordered ideal, source ancestry, physical
valuation, and origin witness.  Hence restricting the compiler to the spine
changes none of its outputs.
\end{proposition}

\begin{proof}
The initial shell-current state is a spine generator.  Transition
\textup{(Q2)} applies a target-cyclic adjoint or joint-kernel restriction;
\textup{(Q3)} applies a first-return constructor; and \textup{(Q4)} applies a
projected total graph.  These are precisely the closure operations in
\cref{def:target-cyclic-structural-spine}.  Induction places every queued
state in the spine.  Conversely, structural induction over the least closed
subgraph expresses every finite spine word as a finite composition of those
transitions.  Projection-first saturation supplies every descent defect, and
source-preserving completion retains its ancestry.  Carrier and origin
coordinates are unchanged by \cref{thm:projection-first-saturation}.
\end{proof}

\begin{definition}[Cumulative joint continuum record]
\label{def:cumulative-joint-continuum-record}
Let
\(\mathfrak q_0\to\mathfrak q_1\to\cdots\)
be any same-origin sequence of structural records generated by the public
continuum grammar.  Its depth-
\(n\) \emph{cumulative joint record} is
\begin{equation}
 \label{eq:cumulative-joint-continuum-record}
 \mathcal J_n
 :=\operatorname{Sat}_{\mathsf G_\Sigma}
   \operatorname{Join}(\mathfrak q_0,\ldots,\mathfrak q_n).
\end{equation}
The join contains every retained coordinate of the prefix, every mixed
product and polarization generated by the public operations, every ordinary
or total-graph relation needed to evaluate those coordinates jointly, the
union of their five-source ancestries, their common origin, and the complete
shell-current history.  Joining is an accounting operation on one physical
history.  It introduces no physical expenditure and carries the restriction
of the original physical measure.

For a nonzero finite positive observer law \(\nu\) on the common measurable
same-origin evaluation space on which all \(\mathcal J_n\) are defined, let
\(\mathcal A_n:=\sigma(\mathcal J_n)\), let
\(P_n:L^2(\nu)\to L^2(\mathcal A_n,\nu)\) be conditional expectation, and
put
\begin{equation}
 \label{eq:cumulative-joint-limit-algebra}
 \mathcal A_\infty:=\sigma\!\left(\bigcup_{n\ge0}\mathcal A_n\right),
 \qquad
 \mathcal J_\infty
 :=\operatorname{Sat}_{\mathsf G_\Sigma}
   \overline{\bigcup_{n\ge0}\mathcal J_n}^{\,\rm gc}.
\end{equation}
For a target shell \(\Phi\), let
\(a_\Phi(e_j):=\Phi(u_j^{\rm out})-\Phi(u_j^{\rm in})=1\) be the full
normalized progress of the \(j\)-th completed shell crossing.  The joint record retains the
\emph{unnormalized target cocycle}
\begin{equation}
 \label{eq:unnormalized-target-cocycle}
 \Gamma_N:=\sum_{j<N}\delta_{e_j},
 \qquad
 A_{\Phi,N}:=\sum_{j<N}a_{\Phi}(e_j),
 \qquad
 \Pi_N:=\sum_{j<N}p_{\Phys}(e_j),
 \qquad
 A_{\Phi,N}=\sum_{c\in\Channels}Q^c_{\Phi,N}(1).
\end{equation}
The normalized first-return law used to locate recurrent support is stored as
a separate coordinate; it never replaces
\(\Gamma_N\), \(A_{\Phi,N}\), \(\Pi_N\), or the source-resolved currents in
\eqref{eq:unnormalized-target-cocycle}.

More generally, every generated continuum object \(q\) carries its
\emph{expanded derivation} \(\operatorname{Exp}(q)\): the finite directed
acyclic graph containing all primitive inputs, every intermediate value, and
the constructor at each vertex.  The constructor list includes signed and
nonlinear combination, restriction, localization, stopping, conditioning,
chart change, scaling, quotient, trace, adjoint, polarization, weak passage
to a total-graph value, and finite evaluation of a directed limit.  If a
vertex uses physical quantities \(\mu^1,\ldots,\mu^k\), the expanded record
retains their owner-tagged restrictions and the exact domination or balance
graph invoked at that vertex.  It never replaces them by an unsigned sum.
For a quotient \(N/D\), empirical average, rescaled chart, or conditional
law, the record retains \((N,D)\), its zero-denominator boundary cell, and
the unnormalized cocycles from which the normalized value was formed.

For a same-origin family \((q_j)\), its \emph{universal cumulative record}
is the saturation of the join of
\(\operatorname{Exp}(q_0),\ldots,\operatorname{Exp}(q_n)\), including every
mixed constructor generated by the public grammar.  Thus
\eqref{eq:cumulative-joint-continuum-record} is not a construction reserved
for shell currents or recurrent scale charts; it is the mandatory evaluation
interface for every composite continuum quantity.
\end{definition}

\begin{definition}[Prospective continuum backward frontier]
\label{def:prospective-continuum-backward-frontier}
Fix one completed target-shell crossing, beginning at its last entrance into
the outer shell before contact, and stop its public evolution at the first
subsequent entrance into the inner shell.  A path that exits the outer shell
first is a failed attempt and is not included in this completed-crossing
record.
Let \(\rho\) be the finite expanded derivation of the stopped shell increment
\[
 q_\rho:=\Phi_r(u_{\rm out})-\Phi_r(u_{\rm in})=1,
\]
formed only from values available on this pre-contact interval.
A \emph{backward frontier} is an ordered finite set of vertices
\(\Gamma=(v_1,\ldots,v_k)\) meeting every directed path from a primitive
state-dependent input of \(\operatorname{Exp}(q_\rho)\) to its terminal
vertex.  Every \(v_\ell\) is a proper ancestor of the terminal shell label:
it is an input, state, current occurrence, or integrated current on the
half-open pre-contact history, rather than the completed contact flag.  Write
\[
 J_{\Gamma,\ell}:=(f_{v_1},\ldots,f_{v_\ell}),
 \qquad
 \mathcal A_{\Gamma,\ell}:=\sigma(J_{\Gamma,\ell}),
 \qquad 0\le \ell\le k,
\]
where \(J_{\Gamma,0}\) is constant.  The frontier record contains the
ancestry of all \(v_\ell\), their joint interface, and the complete
downstream subgraph from \(\Gamma\) to \(q_\rho\).  A conditional extension
is the pair \((J_{\Gamma,\ell-1},f_{v_\ell})\), rather than the isolated
coordinate \(f_{v_\ell}\).

For a nonzero finite positive law \(\nu\) on the measurable carrier of this record,
normalize \(\nu\) to a probability law and let \(P_{\Gamma,\ell}\) be
conditional expectation onto
\(L^2(\mathcal A_{\Gamma,\ell},\nu)\).  For every bounded centered
target-shell contrast \(y_\Phi\), define
\begin{equation}
 \label{eq:continuum-backward-conditional-increment}
 \Delta_{\Gamma,\ell}(y_\Phi)
 :=\bigl\|
 (P_{\Gamma,\ell}-P_{\Gamma,\ell-1})y_\Phi
 \bigr\|_{L^2(\nu)}^2.
\end{equation}
When \(y_\Phi\ne0\), its normalized conditional increment is
\begin{equation}
 \label{eq:continuum-backward-normalized-conditional-increment}
 \widehat\Delta_{\Gamma,\ell}(y_\Phi)
 :=\frac{\Delta_{\Gamma,\ell}(y_\Phi)}
          {\|y_\Phi\|_{L^2(\nu)}^2}.
\end{equation}
This definition uses no compact embedding, critical-norm bound, convergence
modulus, or perturbative estimate.
For an actual crossing there is a canonical law requiring no analytic
choice.  Choose the frontier first at the source-resolved integrated-current
outputs in the exact shell identity, and then expand their complete
ancestries.  Let \(\rho^\circ\) be the corresponding zero-current comparison
record: it retains the same public coefficients and initial endpoint and
replaces those current outputs by the additive zero of their current spaces.
This is the algebraic zero object for the exact current identity, not an
asserted solution of the equation.  Hence \(q_{\rho^\circ}=0\).  On the
two-record space
\begin{equation}
 \label{eq:canonical-two-history-law}
 \Omega_\rho:=\{\rho,\rho^\circ\},
 \qquad
 \nu_\rho:=\tfrac12(\delta_\rho+\delta_{\rho^\circ}),
 \qquad
 y_{\Phi,\rho}:=q-\tfrac12,
\end{equation}
one has \(\|y_{\Phi,\rho}\|_2^2=1/4\).  Every frontier coordinate is
evaluated separately on the two records; no infeasible hybrid tuple is
introduced.
\end{definition}

\begin{theorem}[Exact prospective backward factorization]
\label{thm:exact-prospective-continuum-backward-factorization}
For every frontier in
\cref{def:prospective-continuum-backward-frontier}, contraction of the finite
downstream graph produces a measurable map \(h_{\rho,\Gamma}\) satisfying
\begin{equation}
 \label{eq:continuum-backward-factorization}
 q_\rho=h_{\rho,\Gamma}\circ J_{\Gamma,k}.
\end{equation}
Consequently, if \(P_{q_\rho}\) denotes conditional expectation onto
\(\sigma(q_\rho)\), then
\begin{align}
 \bigl\|P_{\Gamma,k}y_\Phi\bigr\|_2^2
 &=\bigl\|P_{q_\rho}y_\Phi\bigr\|_2^2
   +\bigl\|(P_{\Gamma,k}-P_{q_\rho})y_\Phi\bigr\|_2^2,
 \label{eq:continuum-backward-pythagoras}\\
 \bigl\|P_{\Gamma,k}y_\Phi\bigr\|_2^2
 &=\sum_{\ell=1}^{k}\Delta_{\Gamma,\ell}(y_\Phi).
 \label{eq:continuum-backward-increment-sum}
\end{align}
In particular,
\begin{equation}
 \label{eq:continuum-load-bearing-contextual-extension}
 \bigl\|P_{q_\rho}y_\Phi\bigr\|_2^2>0
 \quad\Longrightarrow\quad
 \Delta_{\Gamma,\ell}(y_\Phi)>0
 \quad\hbox{for some }\ell.
\end{equation}
For the canonical source-current frontier and two-record law
\eqref{eq:canonical-two-history-law}, the terminal shell coordinate satisfies
\(P_{q_\rho}y_{\Phi,\rho}=y_{\Phi,\rho}\), and therefore
\begin{equation}
 \label{eq:canonical-crossing-conditional-charge}
 \sum_{\ell=1}^{k}
 \Delta_{\Gamma,\ell}(y_{\Phi,\rho})
 =\bigl\|P_{\Gamma,k}y_{\Phi,\rho}\bigr\|_2^2
 =\frac14,
 \qquad
 \sum_{\ell=1}^{k}
 \widehat\Delta_{\Gamma,\ell}(y_{\Phi,\rho})=1.
\end{equation}
Thus every actual crossing, without an observer-law hypothesis, exposes a
positive finite contextual source extension.  Further backward expansion
retains that source output together with its complete constructor ancestry;
it does not require a second comparison record.
The complete contextual extension
\((J_{\Gamma,\ell-1},f_{v_\ell})\) in
\eqref{eq:continuum-load-bearing-contextual-extension} is re-entered into the
five-source projection before any residual quotient.  Its ancestry is the
union of the ancestries of its inputs.  Thus a nonlinear or conditional
synergy may carry target alignment even when every input is target-null in
isolation, but that synergy is a finite source-owned constructor occurrence.
\end{theorem}

\begin{proof}
Every path from a primitive state-dependent input to the terminal vertex
meets \(\Gamma\).  Evaluating the vertices below the cut in topological order
therefore gives \(h_{\rho,\Gamma}\) and proves
\eqref{eq:continuum-backward-factorization}.  Hence
\(\sigma(q_\rho)\subseteq\mathcal A_{\Gamma,k}\).  Conditional expectations
onto the nested closed subspaces satisfy
\(P_{q_\rho}P_{\Gamma,k}=P_{\Gamma,k}P_{q_\rho}=P_{q_\rho}\), which gives
\eqref{eq:continuum-backward-pythagoras}.  The spaces
\(L^2(\mathcal A_{\Gamma,\ell},\nu)\) are nested, so their successive
projection differences are pairwise orthogonal.  Since the shell contrast is
centered, \(P_{\Gamma,0}y_\Phi=0\); Pythagoras gives
\eqref{eq:continuum-backward-increment-sum}.  A positive left side forces a
positive summand.  On the canonical two-history space, \(q_\rho\) takes the
values one and zero and \(y_{\Phi,\rho}=q_\rho-1/2\), so conditioning on
\(q_\rho\) fixes the contrast.  This proves
\eqref{eq:canonical-crossing-conditional-charge}.  Expanding the
load-bearing constructor and applying the
source-preserving normal form assigns the contextual input word to
\(\Geom,\Caus,\Abs,\Lift\), or \(\Bdry\), with the union of its input
ancestries.  Projection is applied to this joint word, so individual
target-nullity cannot delete the mixed increment.
\end{proof}

\begin{theorem}[Cofinal prospective-frontier closure]
\label{thm:cofinal-prospective-frontier-closure}
Let \((e_j)_{j\ge0}\) be any same-origin sequence of completed shell
crossings, and flatten the expanded backward frontiers of the first \(N\)
crossings into the cumulative record \(\mathcal J_N\) of
\cref{def:cumulative-joint-continuum-record}.  Let
\(\mathcal A_N=\sigma(\mathcal J_N)\) and let \(P_N\) be conditional
expectation on any finite positive law carried by the measurable joint
record.  Then
\begin{equation}
 \label{eq:cofinal-prospective-frontier-orthogonal-sum}
 P_\infty y_\Phi=\sum_{N\ge0}(P_N-P_{N-1})y_\Phi,
 \qquad
 \|P_\infty y_\Phi\|_2^2
 =\sum_{N\ge0}\|(P_N-P_{N-1})y_\Phi\|_2^2,
\end{equation}
with \(P_{-1}=0\).  Therefore a target-visible cumulative record has a
target-visible finite joint prefix.  Applying
\cref{thm:exact-prospective-continuum-backward-factorization} inside the
least such prefix exposes a finite contextual extension with inherited
five-channel ancestry.

For the actual crossings, let \(\nu_{\rho_j}\) and
\(y_{\Phi,\rho_j}\) be their canonical two-record laws and contrasts.  Form
the finite orthogonal occurrence sum
\begin{equation}
 \label{eq:finite-occurrence-frontier-hilbert-sum}
 \mathcal H_N:=\bigoplus_{j<N}L^2(\nu_{\rho_j}),
 \qquad
 \mathbf Y_N:=\bigoplus_{j<N}y_{\Phi,\rho_j}.
\end{equation}
Order the pairs \((j,\ell)\) lexicographically and reveal the corresponding
frontier coordinates in that order.  The resulting conditional-expectation
subspaces of \(\mathcal H_N\) are nested, and their exact orthogonal
increments give
\begin{equation}
 \label{eq:finite-occurrence-frontier-exact-accounting}
 \|\mathbf Y_N\|_{\mathcal H_N}^2
 =\frac N4
 =\sum_{j<N}\sum_{\ell}
     \Delta_{\Gamma_j,\ell}(y_{\Phi,\rho_j}),
 \qquad
 N=\sum_{j<N}\sum_{\ell}
     \widehat\Delta_{\Gamma_j,\ell}(y_{\Phi,\rho_j}).
\end{equation}

The same join retains the unnormalized identities
\begin{equation}
 \label{eq:cofinal-prospective-frontier-cocycle}
 A_{\Phi,N}=N,
 \qquad
 A_{\Phi,N}=\sum_{c\in\Channels}Q^c_{\Phi,N}(1).
\end{equation}
Thus passage to a subsequence, empirical law, recurrent average, renormalized
chart, weak limit, or compactification cannot replace the cumulative target
cocycle by its average.  If every finite contextual increment is
target-null, then \(P_\infty y_\Phi=0\).  Independently of any auxiliary
observer law, \eqref{eq:finite-occurrence-frontier-exact-accounting} retains
the complete accumulated target alignment of the first \(N\) crossings.
If contact persists, positive finite contextual source increments are thus
present and remain attached to the cumulative physical ledger.
These alternatives require neither terminal compactness nor perturbative
critical-norm bounds.
\end{theorem}

\begin{proof}
The sigma algebras \((\mathcal A_N)\) increase.  Hilbert-space martingale
convergence gives strong convergence of \(P_Ny_\Phi\) to
\(P_\infty y_\Phi\); orthogonality of successive differences gives
\eqref{eq:cofinal-prospective-frontier-orthogonal-sum}.  A nonzero sum has a
nonzero finite partial sum, and the least prefix at which its squared norm
increases contains a positive conditional increment.  The preceding theorem
supplies its finite contextual owner and ancestry.  Each completed normalized
crossing contributes exactly one to the shell cocycle before normalization.
Joining occurrences is additive and does not identify them, proving
\eqref{eq:cofinal-prospective-frontier-cocycle}.  Summing
\eqref{eq:canonical-crossing-conditional-charge} over the first \(N\)
crossings in the orthogonal sum
\eqref{eq:finite-occurrence-frontier-hilbert-sum} gives
\eqref{eq:finite-occurrence-frontier-exact-accounting}.  The common joint
projection detects cross-history synergy, while the occurrence direct sum
and unnormalized cocycle retain multiplicity and physical ownership.
\end{proof}

\begin{corollary}[Backward-saturated residual nullity]
\label{cor:backward-saturated-residual-nullity}
Fix a physical ceiling and saturate the continuum algebra over the expanded
backward frontiers of all legal same-origin finite stopped histories below
that ceiling, all finite contextual joins, and their directed cumulative
joins.  After applying the five source projections to this saturated class,
the remaining prospective residual is target-null:
\begin{equation}
 \label{eq:backward-saturated-residual-nullity}
 G^{J_{\rm res},{\rm src,sat}}_{\Sigma,\Phys}=0.
\end{equation}
Indeed, any residual with nonzero target pairing determines a positive
conditional increment on a finite expanded joint prefix.  Its detecting
contextual word is then an element of the same saturation and has inherited
\(\Geom/\Caus/\Abs/\Lift/\Bdry\) ancestry, contradicting its placement in
the post-saturation residual.
\end{corollary}

\begin{proof}
Suppose a residual class has nonzero pairing with a bounded shell contrast.
For a finite class, apply
\cref{thm:exact-prospective-continuum-backward-factorization}; for a directed
class, first apply \cref{thm:cofinal-prospective-frontier-closure}.  In either
case a finite contextual extension has positive conditional increment.  Its
expanded constructor graph belongs to the declared saturation.  Apply
\cref{thm:continuum-normal-form-compilation,thm:no-target-visible-residual} to
that graph.  Its target-active quotient is generated by the five head
components with inherited exact-support ancestry, while the residual lies in
\(\Null_\Sigma\).  Hence every residual pairing vanishes, which is
\eqref{eq:backward-saturated-residual-nullity}.
\end{proof}

\begin{corollary}[No terminal analytic prerequisite]
\label{cor:no-terminal-analytic-prerequisite}
Prospective source extraction for singular contact is performed on finite
stopped shell histories and closed by
\cref{thm:cofinal-prospective-frontier-closure}.  It requires neither a
bounded local critical norm, compactness of terminal windows, convergence to
an ordinary blow-up profile, nor perturbative control in a critical space.
When one of those properties holds, it supplies an additional generated
coordinate of the same frontier.  When it fails, its typed total-graph value
retains the same inherited five-source ancestry; failure cannot delete the
finite target increment or create a source family.
\end{corollary}

\begin{definition}[Same-origin occurrence space and physical comparison measure]
\label{def:same-origin-occurrence-measure}
Let \((e_j)_{j\in J}\) be the completed target-shell crossings of one
physical orbit, with the intervening gaps inserted before any current is
restricted.  Give the countable occurrence space
\(\Omega_{\rm occ}:=J\) its full sigma algebra and define
\begin{equation}
 \label{eq:occurrence-progress-and-price-measures}
 \Gamma_\Sigma:=\sum_{j\in J}a_j\delta_j,
 \qquad
 \mu_{\Phys}:=\sum_{j\in J}p_j\delta_j,
 \qquad
 a_j:=\Phi_j(u_j^{\rm out})-\Phi_j(u_j^{\rm in})=1.
\end{equation}
Here \(p_j\) is the restriction to the physical image of the complete
nonnegative public carrier measure admitted by
\cref{def:target-aligned-physical-charge}; it is not normalized dissipation
in the moving chart.  If several crossings overlap in physical space-time,
replace them by the disjoint atoms of their finite incidence algebra and
retain the crossing multiplicity on each atom.  Thus
\(\mu_{\Phys}(\Omega_{\rm occ})\) is bounded by the single physical ledger,
whereas \(\Gamma_\Sigma([0,N))=N\).

For each source \(c\in\Channels\), signed gluing and universal expanded
derivation give a signed occurrence current \(Q^c\) with
\begin{equation}
 \label{eq:source-resolved-occurrence-identity}
 \Gamma_\Sigma=\sum_{c\in\Channels}Q^c
\end{equation}
on every finite subset.  The equality, the five owners, the physical measure,
and the multiplicity measure are joined before taking a directed limit.
Chart changes merely relabel the atoms; they do not rescale either measure.
\end{definition}

\begin{theorem}[Exact compositional accounting against the physical ledger]
\label{thm:exact-compositional-physical-accounting}
On the same-origin occurrence space of
\cref{def:same-origin-occurrence-measure}, the positive target-progress
measure has the unique Lebesgue decomposition
\begin{equation}
 \label{eq:target-progress-lebesgue-decomposition}
 \Gamma_\Sigma=f_\Sigma\mu_{\Phys}+\Gamma_\Sigma^{\perp},
 \qquad
 f_\Sigma\ge0,
 \qquad
 \Gamma_\Sigma^{\perp}\perp\mu_{\Phys}.
\end{equation}
It is computed atomwise by
\begin{equation}
 \label{eq:atomic-target-amplification}
 f_\Sigma(j)=p_j^{-1}\quad\hbox{when }p_j>0,
 \qquad
 \Gamma_\Sigma^{\perp}
   =\sum_{p_j=0}\delta_j.
\end{equation}
On each occurrence atom, apply the finite-dimensional positive/negative split
to the five signed values in
\eqref{eq:source-resolved-occurrence-identity}.  Atomwise directed joining
then assigns every nonzero part of \(f_\Sigma\mu_{\Phys}\) and
\(\Gamma_\Sigma^{\perp}\) to at least one of
\(\Geom,\Caus,\Abs,\Lift,\Bdry\), with its complete mixed ancestry.  In
particular the singular measure in
\eqref{eq:target-progress-lebesgue-decomposition} is not
\(J_{\rm res}\).

If the orbit has finite physical budget and completes infinitely many
crossings, then necessarily
\begin{equation}
 \label{eq:finite-budget-amplification-alternative}
 \operatorname*{ess\,sup}_{\mu_{\Phys}}f_\Sigma=+\infty
 \quad\hbox{or}\quad
 \Gamma_\Sigma^{\perp}(\Omega_{\rm occ})>0.
\end{equation}
Equivalently, if the expanded five-source algebra proves, on the actual
same-origin occurrence record,
\begin{equation}
 \label{eq:compositional-regularity-test}
 \Gamma_\Sigma^{\perp}=0,
 \qquad
 f_\Sigma\le C<\infty
 \quad\mu_{\Phys}\text{-a.e.},
\end{equation}
then the number of completed crossings satisfies the sharp bound
\begin{equation}
 \label{eq:physical-crossing-bound}
 \#J=\Gamma_\Sigma(\Omega_{\rm occ})
 \le C\mu_{\Phys}(\Omega_{\rm occ})<\infty,
\end{equation}
and contact with a target requiring a cofinal shell family is impossible.
Conversely, every finite-budget cofinal contact produces one of the two
source-owned prospective amplification values in
\eqref{eq:finite-budget-amplification-alternative}.  They enter the
zero-production, carrier--polar, and realization transitions; neither value
is a singular-existence terminal by itself.
\end{theorem}

\begin{proof}
The physical images of the incidence atoms are disjoint, so restriction of
the public carrier gives the countably additive measure \(\mu_{\Phys}\).
The unnormalized shell cocycle gives the countably additive counting measure
\(\Gamma_\Sigma\).  The Lebesgue decomposition theorem gives
\eqref{eq:target-progress-lebesgue-decomposition}; on a countable atomic
space its values are exactly \eqref{eq:atomic-target-amplification}.

For every finite prefix, signed gluing and prospective accountability give
\eqref{eq:source-resolved-occurrence-identity}.  Expand all mixed
constructors and split the five real source values on each occurrence atom
only after this equality has been formed.  Because the occurrence space is
countably atomic, source-preserving passage to the directed join then assigns
each nonzero positive aggregate atom to the union of the five-source
ancestries of its finite detecting word.  No countably additive extension of
an individual signed \(Q^c\) is used.  Singularity with respect to
\(\mu_{\Phys}\) changes the comparison class, not the structural owner.

Suppose that \(\mu_{\Phys}(\Omega_{\rm occ})<\infty\),
\(\Gamma_\Sigma^{\perp}=0\), and \(f_\Sigma\le C\).  Integration of
\eqref{eq:target-progress-lebesgue-decomposition} gives
\[
 \#J=\int f_\Sigma\,\dd\mu_{\Phys}
 \le C\mu_{\Phys}(\Omega_{\rm occ})<\infty.
\]
Its contrapositive is \eqref{eq:finite-budget-amplification-alternative}.
No compactness, continuation criterion, scale summability, or lower price per
normalized window enters the argument.
\end{proof}

\begin{definition}[Physical target-transition capacity]
\label{def:physical-target-transition-capacity}
For a physical ceiling \(E<\infty\), let
\(\mathscr O_{\Phys}(E)\) be the same-origin orbit records whose public
physical carrier has mass at most \(E\).  Define
\begin{align}
 \label{eq:physical-target-transition-capacity}
 \mathsf t^{\Phys}_\Sigma(E)
 &:=
 \sup_{R\in\mathscr O_{\Phys}(E)}
 \Gamma_{\Sigma,R}(\Omega_{{\rm occ},R}),\\
 \label{eq:physical-target-amplification-envelope}
 \mathsf a^{\Phys}_\Sigma(E)
 &:=
 \sup_{R\in\mathscr O_{\Phys}(E)}
 \sup_{A\in\mathcal A_{{\rm occ},R}^{\rm fin}}
 \frac{\Gamma_{\Sigma,R}(A)}{\mu_{\Phys,R}(A)}.
\end{align}
Here \(\mathcal A_{{\rm occ},R}^{\rm fin}\) is the algebra of finite
unions of generated occurrence atoms.  The quotient is \(+\infty\) when
its numerator is positive and its denominator is zero.  Both suprema are
values of the saturated target-local algebra: they are taken only after the
five-source join, same-origin restriction, and unnormalized multiplicity
record have been formed.
\end{definition}

\begin{theorem}[Target-capacity duality]
\label{thm:physical-target-capacity-duality}
For every \(E<\infty\),
\begin{equation}
 \label{eq:physical-target-capacity-bound}
 \mathsf t^{\Phys}_\Sigma(E)
 \le E\,\mathsf a^{\Phys}_\Sigma(E),
\end{equation}
with the usual extended-real convention.  More precisely, for a fixed
same-origin record \(R\), the following are equivalent:
\begin{enumerate}[label=\textup{(D\arabic*)}]
\item there is \(C<\infty\) such that
  \(\Gamma_{\Sigma,R}(A)\le C\mu_{\Phys,R}(A)\) for every generated
  finite occurrence set \(A\);
\item the inequality extends to the occurrence sigma algebra and
  \(\Gamma_{\Sigma,R}^{\perp}=0\),
  \(\lVert f_{\Sigma,R}\rVert_{L^\infty(\mu_{\Phys,R})}\le C\);
\item after every finite generated reordering of occurrence atoms, every
  finite prefix satisfies \(A_{\Phi,N}\le C\Pi_N\).
\end{enumerate}
The least admissible \(C\) is
\begin{equation}
 \label{eq:least-physical-amplification-constant}
 C_R=
 \operatorname*{ess\,sup}_{\mu_{\Phys,R}}f_{\Sigma,R}
 \quad\text{if }\Gamma_{\Sigma,R}^{\perp}=0,
 \qquad C_R=+\infty\quad\text{otherwise}.
\end{equation}
Consequently a finite physical target-transition capacity is not an axiom:
for each fixed occurrence record it is exactly the bounded-amplification
value of the five-source join.  A cofinal contact at finite physical budget
forces \(C_R=+\infty\) on that contact-derived record and thereby exposes an
unbounded-amplification or physical-singular five-source successor.  This
recordwise conclusion is not a computation of reachability from the initial
datum and is not a singular-existence statement.
\end{theorem}

\begin{proof}
Finite unions of atoms form an algebra generating the occurrence sigma
algebra.  If \textup{(D1)} holds, monotone approximation extends the measure
inequality to that sigma algebra.  The Lebesgue decomposition
\eqref{eq:target-progress-lebesgue-decomposition} then has no singular part,
and its Radon--Nikodym derivative is at most \(C\); this is \textup{(D2)}.
The converse follows by integration.  Taking \(A\) to be the first \(N\)
occurrences proves \textup{(D3)}.  Conversely, the target program may order
any finite generated occurrence set first while retaining the intervening
gaps; hence \textup{(D3)} for every generated ordering gives \textup{(D1)}.
The minimality property of the Radon--Nikodym derivative gives
\eqref{eq:least-physical-amplification-constant}.  Finally, for each \(R\),
\[
 \Gamma_{\Sigma,R}(\Omega_{{\rm occ},R})
 \le C_R\mu_{\Phys,R}(\Omega_{{\rm occ},R}),
\]
and taking suprema proves \eqref{eq:physical-target-capacity-bound}.
\end{proof}

\begin{definition}[Finite-prefix physical valuation]
\label{def:finite-prefix-physical-valuation}
Fix a same-origin occurrence record and a finite occurrence cylinder \(F\).
Let \(Y_F\) be the finite-dimensional real space of all source-owned current
coordinates detected on \(F\), after quotienting the target-null subspace.
Write \(Q_F\in Y_F\) for a target-current owner and
\(p_F\in Y_F\) for the restriction of the complete public physical measure,
including admitted production and exterior input with their public weights.
Let \(D_F\subset Y_F\) be the closed convex cone generated by the restrictions
to \(F\) of
\[
 \delta\mathscr K+\Null_\Sigma ,
\]
where the two orientations of every storage coboundary are included, internal
currents are glued with their signs before taking a positive part, and every
production or exterior term has first been mapped into \(p_F\) by its exact
public carrier row.  A term for which that map is nonordinary remains a graph
coordinate and is not inserted into \(D_F\).  Define the extended physical valuation
\begin{equation}
 \label{eq:finite-prefix-physical-valuation}
 \operatorname{val}_F(Q_F)
 :=\inf\{C\ge0:Q_F\in D_F+C p_F\}.
\end{equation}
The infimum of the empty set is \(+\infty\).  Thus a finite value is a
finite-prefix carrier domination, while an infinite value records failure of
that domination on the same prefix; neither value is a terminal verdict.
\end{definition}

\begin{theorem}[Exact finite-prefix source-profile/carrier duality]
\label{thm:finite-prefix-source-profile-carrier-duality}
For every \(F\) and \(Q_F\) as in
\cref{def:finite-prefix-physical-valuation},
\begin{equation}
 \label{eq:finite-prefix-carrier-duality}
 \operatorname{val}_F(Q_F)
 =\sup\left\{
   \lambda(Q_F):
   \lambda\in Y_F^*,\quad
   \lambda(D_F)\le0,\quad
   \lambda(p_F)\le1
 \right\}.
\end{equation}
The equality is interpreted in the extended nonnegative reals.  If the
right-hand side is finite and the primal set is nonempty, the infimum is
attained.  If it is infinite, then for every \(M<\infty\) a finite-prefix
dual word satisfies
\begin{equation}
 \label{eq:finite-prefix-unbounded-dual-witness}
 \lambda_M(D_F)\le0,
 \qquad \lambda_M(p_F)\le1,
 \qquad \lambda_M(Q_F)>M.
\end{equation}
Each \(\lambda_M\) is joined to \(Q_F\), pulled backward through the complete
joint frontier, and assigned the ancestry of the constructor that first
creates its nonzero conditional target increment.  Hence dual separation
creates no additional source family.
\end{theorem}

\begin{proof}
Pass to the quotient by the lineality space of \(D_F\).  The epigraph
\[
 \{(Q,C):C\ge0,\ Q\in D_F+C p_F\}
\]
is a closed convex cone in the finite-dimensional space
\(Y_F\times\mathbb R\).  Separation of \((Q_F,C)\) from this cone, followed
by normalization of the coefficient of \(C\), gives the lower bound in
\eqref{eq:finite-prefix-carrier-duality}; membership gives the reverse bound.
Closedness gives attainment whenever the feasible set is nonempty and the
value is finite.  If the gauge is infinite, separation at levels
\(M=1,2,\ldots\) gives \eqref{eq:finite-prefix-unbounded-dual-witness}.
The backward-frontier factorization and conditional-increment identity apply
to the finite joint record \((Q_F,\lambda_M)\).  Constructor-level charging
therefore assigns every nonzero dual readout to one of the five inherited
owners.
\end{proof}

\begin{theorem}[Monotone saturated physical valuation]
\label{thm:monotone-saturated-physical-valuation}
Let \(F_1\subset F_2\subset\cdots\) be the canonical exhaustive sequence of
finite occurrence cylinders, including every intervening gap, source word,
storage endpoint, exterior-input coordinate, and total-graph value generated
up to the corresponding depth.  Then
\begin{equation}
 \label{eq:saturated-physical-valuation}
 \operatorname{val}^{\Sat}(Q)
 :=\sup_n\operatorname{val}_{F_n}(Q_{F_n})
\end{equation}
is independent of the chosen cofinal enumeration and is the least
\(C\in[0,+\infty]\) for which every finite generated restriction admits a
domination by \(D_{F_n}+Cp_{F_n}\).  Equivalently,
\begin{equation}
 \label{eq:saturated-profile-dual}
 \operatorname{val}^{\Sat}(Q)
 =\sup_n\ \sup_{\substack{\lambda(D_{F_n})\le0\\
                           \lambda(p_{F_n})\le1}}
        \lambda(Q_{F_n}).
\end{equation}
A finite value charges every finite target prefix at one common physical
rate.  An infinite value supplies a cofinal family of finite target-aligned
dual increments with their unnormalized occurrence cocycle.  It is processed
by amplification and equality saturation before any averaging or limiting
normalization.
\end{theorem}

\begin{proof}
Restriction of a domination on \(F_{n+1}\) is a domination on \(F_n\), so the
finite valuations are nondecreasing.  Two cofinal enumerations dominate each
other termwise after passing to subsequences, which proves independence.
Taking the least common constant over all finite restrictions gives
\eqref{eq:saturated-physical-valuation}; applying
\cref{thm:finite-prefix-source-profile-carrier-duality} at every depth gives
\eqref{eq:saturated-profile-dual}.  The final assertions retain, rather than
divide by, the cumulative occurrence coordinate.
\end{proof}

\begin{theorem}[Variation-safe finite-prefix current principle]
\label{thm:variation-safe-finite-prefix-current-principle}
For each source \(c\in\Channels\), let \(Q_F^c\) denote its signed current on
a finite occurrence cylinder \(F\).  All source identities and cancellations
are performed at this level.  Exactly one of the following holds along the
canonical exhaustion:
\begin{enumerate}[label=\textup{(J\arabic*)}]
\item
\(\sup_n |Q_{F_n}^c|(F_n)<\infty\).  Then the compatible finite-prefix
currents extend uniquely to a finite signed measure \(Q^c\) on the occurrence
sigma algebra, and finite-prefix signed gluing passes to that measure.
\item
\(\sup_n |Q_{F_n}^c|(F_n)=\infty\).  The increasing variation coordinate,
together with the first finite prefix on which each rational level is
crossed, is the source-owned variation-boundary successor.  No global signed
measure or cancellation across that boundary is used.
\end{enumerate}
In both cases the positive contact measure
\(\Gamma_\Sigma\) remains countably additive.  Consequently every use of a
global Jordan or Lebesgue decomposition in the target calculus applies to
\(\Gamma_\Sigma\) or to a source current satisfying \textup{(J1)}; otherwise
the calculus remains finite-prefix local.
\end{theorem}

\begin{proof}
Compatibility assigns one weight to each occurrence atom.  Under
\textup{(J1)}, the absolute sum of those weights is finite, so their series
defines the unique finite signed measure on the full sigma algebra and agrees
with every finite restriction.  If the variations are unbounded, the
least prefix crossing each rational level is defined and preserves the
source owner.  The two alternatives are exhaustive.  Countable additivity of
\(\Gamma_\Sigma\) follows independently from its occurrence-counting
definition in \cref{def:same-origin-occurrence-measure}.
\end{proof}

\begin{corollary}[Zero physical denominator is an evaluated source cell]
\label{cor:zero-physical-denominator-evaluated-cell}
Let \(A\) be a generated occurrence atom with
\(\mu_{\Phys}(A)=0\) and \(\Gamma_\Sigma(A)>0\).  Then a source
\(c\in\Channels\) satisfies \(Q^c(A)>0\).  On the one-atom cylinder
\(F=\{A\}\), exactly one of the following finite computations occurs:
the target quotient of \(Q_F^c\) is zero; a finite carrier row in
\(D_F+C p_F\) charges it; or the dual value
\(\operatorname{val}_F(Q_F^c)=+\infty\) generates the normalized
zero-production tangent of
\cref{thm:amplification-zero-production-tangent}.  Hence singularity of
\(\Gamma_\Sigma\) with respect to the physical measure is an evaluated
five-source cell, not a singularity verdict and not a residual source.
\end{corollary}

\begin{proof}
The source-resolved finite-set identity gives the source \(c\).  Apply
\cref{thm:finite-prefix-source-profile-carrier-duality}.  Since \(p_F=0\),
every finite positive target value outside \(D_F\) has unbounded gauge.  Its
dual witnesses retain the same finite owner and enter target normalization.
\end{proof}

\begin{theorem}[Nonordinary shell-current first-failure evaluator]
\label{thm:nonordinary-shell-current-first-failure-evaluator}
Consider a finite stopped shell history and totalize, in constructor order,
every product, trace, pressure solve, chart map, scale map, and Stieltjes
pairing used by its shell-current identity.  If the ordinary identity fails,
let \(v\) be the least failed graph.  Join to \(v\) its target covector,
physical denominator, shell multiplicity, complete preceding joint frontier,
and source owner.  Then its target-visible part has exactly one of the
following values:
\begin{enumerate}[label=\textup{(G\arabic*)}]
\item its conditional target increment is zero, so it lies in
\(\Null_\Sigma\);
\item its finite-prefix physical valuation is finite and the corresponding
carrier row charges it;
\item its valuation is infinite and a finite dual witness enters the
source-owned amplification/equality evaluator;
\item a descendant graph used by that evaluator is nonordinary, in which
case the least such descendant, with the enlarged joint frontier, is the
next graph record.
\end{enumerate}
No nonordinary value is removed before joint target projection, and no graph
value is promoted to an ordinary orbit or contact event.
\end{theorem}

\begin{proof}
The total graphs form a finite ordered coproduct on the stopped history, so a
least failed graph exists.  Exact backward factorization through its complete
joint frontier and orthogonal conditional increments decide \textup{(G1)}.
On the complementary target quotient apply
\cref{thm:finite-prefix-source-profile-carrier-duality}; its finite and
infinite values give \textup{(G2)} and \textup{(G3)}.  Totalize every graph
used in that finite computation.  An ordinary value completes the stated
row, while the least failed descendant gives \textup{(G4)}.  Source
preservation under graph completion assigns the ancestry of \(v\) throughout.
The realization gate is separate, so none of these graph values implies
physical existence.
\end{proof}

\begin{theorem}[Prospective source-amplification extraction]
\label{thm:prospective-source-amplification-extraction}
Let a same-origin pre-contact record have the source-resolved occurrence
identity
\[
 \Gamma_\Sigma=\sum_{c\in\Channels}Q^c
\]
on every finite occurrence set and finite physical measure
\(\mu_{\Phys}\).  No global signed-measure extension of the individual
\(Q^c\) is required.  The following alternatives hold.
\begin{enumerate}[label=\textup{(X\arabic*)}]
\item if \(\operatorname*{ess\,sup}f_\Sigma=+\infty\), there are one
  fixed source \(c_*\), numbers \(M_k\uparrow\infty\), and generated
  finite pre-contact occurrence sets \(A_k\) such that
  \begin{equation}
   \label{eq:sourcewise-unbounded-amplification}
   Q^{c_*}(A_k)>M_k\mu_{\Phys}(A_k);
  \end{equation}
\item if \(\Gamma_\Sigma^\perp\ne0\), there are one fixed source
  \(c_*\), a \(\mu_{\Phys}\)-null generated occurrence atom
  \(A\) such that
  \begin{equation}
   \label{eq:sourcewise-singular-target-mass}
   \mu_{\Phys}(A)=0,
   \qquad Q^{c_*}(A)>0.
  \end{equation}
\end{enumerate}
Every selected finite occurrence set is selected before its last crossing
and carries
the complete prefix derivation, physical denominator, shell multiplicity,
and source ancestry.  Thus aggregate target amplification cannot be a
terminal mechanism: it always exposes a quantitative prospective owner in
one of the five existing channels.
\end{theorem}

\begin{proof}
Suppose first that \(f_\Sigma\) is essentially unbounded.  Since the
occurrence space is countably atomic and
\(f_\Sigma(j)=p_j^{-1}\) on every positive-price atom, for every integer
\(m\) there is an occurrence atom \(A_m=\{j_m\}\) with
\[
 \Gamma_\Sigma(A_m)>5m\mu_{\Phys}(A_m).
\]
Since \(\Gamma_\Sigma(A_m)=\sum_cQ^c(A_m)\), at least one source satisfies
\(Q^c(A_m)>m\mu_{\Phys}(A_m)\).  One of the five sources occurs for
infinitely many \(m\); take that subsequence and call it \(c_*\).  This
proves \eqref{eq:sourcewise-unbounded-amplification}.

If \(\Gamma_\Sigma^\perp\ne0\), the atomic formula
\eqref{eq:atomic-target-amplification} supplies an occurrence atom
\(A=\{j\}\) with \(p_j=0\).  Hence
\(\mu_{\Phys}(A)=0\) and \(\Gamma_\Sigma(A)=1\).  The finite-set
identity \(\Gamma_\Sigma(A)=\sum_cQ^c(A)\) gives a source \(c\) with
\(Q^c(A)>0\).  Put \(c_*=c\).  This proves
\eqref{eq:sourcewise-singular-target-mass} without extending any signed
\(Q^c\) across an infinite-variation set.  Stopping at the last index of a
finite occurrence set makes its selector pre-contact, while the cumulative
join supplies all stated coordinates.
\end{proof}

\begin{definition}[Target-normalized amplification tangent]
\label{def:target-normalized-amplification-tangent}
Let \(c_*\) and \(A_k\) be supplied by
\cref{thm:prospective-source-amplification-extraction}, and put
\(q_k:=Q^{c_*}(A_k)>0\).  On the compact target-local evaluation algebra
\(\mathscr E_{\mathfrak b}\), define
\begin{equation}
 \label{eq:target-normalized-amplification-functional}
 L_k(F):=
 \frac{Q^{c_*}(F\mathbf1_{A_k})}{q_k}.
\end{equation}
Before using positivity, apply the Hahn decomposition of
\(Q^{c_*}|_{A_k}\) and replace \(A_k\) by its positive part; the negative
part remains in the signed same-origin join, and reset
\(q_k:=Q^{c_*}(A_k)\).  Thus \(L_k\) is positive,
\(L_k(1)=1\), and it retains the source word, target shell, physical
denominator, chart, and complete prefix.

For the physical-singular alternative, use the normalized positive restriction
of \(Q^{c_*}\) to the null-carrier occurrence atom in
\eqref{eq:sourcewise-singular-target-mass} in place of
\eqref{eq:target-normalized-amplification-functional}.
Any weak-* accumulation point in the generated evaluation compactification
is called the \emph{target-normalized amplification tangent}.
\end{definition}

\begin{theorem}[Dyadic ordered-coercivity terminal rule]
\label{thm:dyadic-ordered-coercivity-terminal-rule}
Let \(L\) be a positive normalized functional on an active target-local
Archimedean cylinder algebra.  Let \(A,P,M\) be bounded nonnegative generated
coordinates, where \(A\) is target activity, \(P\) is admitted physical
production, and \(M\) is an auxiliary coefficient.  Suppose the generated
ordered relations contain, for integers \(k\ge1\), \(b\ge1\), and \(a\ge0\),
\begin{equation}
 \label{eq:dyadic-ordered-coercivity-row}
 A^{2^k}\le C M^aP^b,
 \qquad C<\infty,
 \qquad M\le M_0<\infty.
\end{equation}
Then
\begin{equation}
 \label{eq:dyadic-zero-production-nullity}
 L(P)=0\quad\Longrightarrow\quad L(A)=0.
\end{equation}
Consequently a positive target floor \(L(A)\ge\eta>0\) has exactly the
following local continuations: \(L(P)>0\), an unbounded \(M\)-coordinate, or
a nonordinary value of the coercivity, product, order, or target graph.  The
first value is charged through the physical-production row; the other values
are the corresponding coefficient or graph successors.  A zero-production,
finite-coefficient positive functional is not a terminal equality model.
\end{theorem}

\begin{proof}
Boundedness and positivity give
\[
 0\le P^b\le\|P\|_\infty^{b-1}P,
 \qquad
 0\le M^a\le M_0^a.
\]
Thus \(L(P)=0\) implies \(L(M^aP^b)=0\).  Applying \(L\) to
\eqref{eq:dyadic-ordered-coercivity-row} gives \(L(A^{2^k})=0\).
Repeated Cauchy--Schwarz yields
\[
 0\le L(A)^{2^k}
 \le L(A^2)^{2^{k-1}}
 \le\cdots\le L(A^{2^k})=0.
\]
Hence \(L(A)=0\).  Totalization of the finitely many relations in
\eqref{eq:dyadic-ordered-coercivity-row} gives the stated complementary graph
values, and the fixed target quotient makes the list exhaustive.
\end{proof}

\begin{theorem}[Amplification-to-zero-production tangent principle]
\label{thm:amplification-zero-production-tangent}
Every unbounded or physical-singular prospective source amplification
produces a positive target-local tangent \(L\) satisfying
\begin{equation}
 \label{eq:zero-production-tangent-identities}
 L(1)=1,
 \qquad L(Q^{c_*}_\Sigma)=1,
 \qquad L(P)=0
 \quad
 (P\in\mathscr P_{\mathfrak b}^{\Phys}),
\end{equation}
after the target-current normalization is included as a coordinate.  It
retains the fixed source \(c_*\) and belongs to the equality stratum
\(\mathcal E_{\mathfrak b}\).  Exactly one of the following occurs:
\begin{enumerate}[label=\textup{(Z\arabic*)}]
\item a finite equality-stratum identity makes the target projection zero;
\item an ordinary same-origin zero-production state realizes the tangent;
\item the least nonordinary equation, chart, scale, pressure, exterior,
  trace, concentration, or realization coordinate gives the next
  five-source successor.
\end{enumerate}
An amplification value is therefore never a terminal numerical infinity.
It is converted constructively into a target-retaining zero-production
state or a typed first-failure successor.
\end{theorem}

\begin{proof}
On the unbounded branch,
\[
 0\le L_k(P)
 \le \frac{\mu_{\Phys}(A_k)}{Q^{c_*}(A_k)}
 \longrightarrow0
 \qquad(P\in\mathscr P_{\mathfrak b}^{\Phys}),
\]
because every admitted production is dominated by the restriction of the
single physical measure.  On the physical-singular branch the same value is
zero identically, since the normalized positive carrier is
\(\mu_{\Phys}\)-null.  Positivity and normalization make the \(L_k\)
a weak-* compact family on the generated compact evaluation algebra.
Every accumulation point therefore satisfies
\eqref{eq:zero-production-tangent-identities}.  Expanded-derivation
saturation preserves the source word and all same-origin coordinates.

Adjoin the relations \(P=0\) to the target-local ordered algebra.  Its
finite \(-1\) identity gives \textup{(Z1)}.  Otherwise totalize the
ordinary realization graph.  An ordinary value gives \textup{(Z2)}; the
least nonordinary value gives \textup{(Z3)} with the ancestry of the failed
operation.  These are the exhaustive outputs of the total graph and introduce
no additional structural source.
\end{proof}

\begin{theorem}[Sound typing of amplification successors]
\label{thm:closed-amplification-successor-terminalization}
Start from the zero-production tangent in
\cref{thm:amplification-zero-production-tangent} and recursively append the
least nonordinary coordinate in the fixed enumeration of equation, source,
chart, scale, pressure, exterior, trace, concentration, and realization
graphs.  Retain the cumulative same-origin join and the unit target
normalization at every stage.  Then exactly one evaluation state occurs:
\begin{enumerate}[label=\textup{(T\arabic*)}]
\item a finite successor is target-null;
\item a finite successor has an ordinary zero-production realization and
  re-enters target-cyclic equality evaluation;
\item the cumulative graph-complete limit is target-null;
\item the cumulative limit is a same-origin, source-owned prospective corona
  record with unit target value and zero public physical production.
\end{enumerate}
In \textup{(T4)}, every finite generated coordinate has an ordinary value or
its declared nonordinary total-graph value, and at least one cofinal
nonordinary coordinate retains nonzero target projection.  The record remains
an internal realization and carrier--polar successor.  It is an explicit
singular mechanism in the sense of \cref{def:explicit-singular-mechanism}
only after an independent approximation-tower path supplies all ordinary
public-equation and contact rows.  A positive functional, separator, or
nonordinary graph value is not terminal.
\end{theorem}

\begin{proof}
At each finite stage, target-nullity gives \textup{(T1)} and an ordinary
realization gives \textup{(T2)}.  Otherwise append the least nonordinary
coordinate.  The cumulative joins form an increasing sequence in the compact
graph-complete evaluation product.  Choose a compatible diagonal limit
\(z_\infty\).  Every coordinate is retained permanently after it is
adjoined, so its total-graph relation passes to \(z_\infty\); the common
origin, fixed source owner, zero-production relations, and unit target
normalization pass as closed coordinates as well.

If the target projection of \(z_\infty\) vanishes, this is
\textup{(T3)}.  Otherwise no finite ordinary realization occurred, hence a
cofinal nonordinary coordinate remains target-visible.  The initial contact
record supplies the ordinary public orbit, and every restriction, chart,
normalization, and total-graph passage retained its same-origin row.
Consequently \(z_\infty\) satisfies the ancestry, hypothetical-origin,
zero-production, and nonzero-target clauses of a prospective corona record.
It fails the independent-origin clause of
\cref{def:explicit-singular-mechanism}.  This is \textup{(T4)}.  The
alternatives are disjoint by the
ordered first-terminal rule and exhaustive by compactness of the totalized
evaluation product.
\end{proof}

\begin{corollary}[Sound continuum terminal theorem]
\label{cor:binary-continuum-terminal-theorem}
Every terminal verdict produced by the continuum evaluator has exactly one of
the following two types:
\[
 \boxed{
 \text{target avoidance}
 \quad\text{or}\quad
 \text{realized five-source singular mechanism}.}
\]
The second type requires the independent-origin and ordinary-realization rows
of \cref{def:explicit-singular-mechanism}.  A contact-derived prospective
record has origin polarity \(\mathsf{Hyp}\) and is not a singular-existence
verdict.  Target-visible \(J_{\rm res}\) remains forbidden.  A nonordinary
compactness value, infinite refinement record, positive functional, or polar
separator is an internal successor rather than a third terminal type.
\end{corollary}

\begin{proof}
Finite contact is prospectively source-resolved.  Infinite contact is reduced
by \cref{thm:exact-compositional-physical-accounting,%
thm:prospective-source-amplification-extraction,%
thm:amplification-zero-production-tangent} to a zero-production successor
execution.  Apply
\cref{thm:closed-amplification-successor-terminalization}.  Its hypothetical
corona re-enters realization and carrier--polar evaluation.  A checked
target-local contradiction proves avoidance; an independently generated
ordinary contacting path proves singular existence.
\end{proof}

\begin{definition}[Evaluated covariant corona record]
\label{def:evaluated-covariant-corona-record}
An evaluated corona record consists of one active same-origin signature
\(\mathfrak b\), its fixed source owner \(Q\), record scale \(r\), scale
rate \(a=\dot r/r\), and, on every ordinary finite half-open prefix
\(I=[\tau_-,\tau_+)\), a generated Stieltjes carrier row
\begin{equation}
 \label{eq:evaluated-corona-carrier-row}
 \dd K=-\dd\mathsf D+\kappa aK\,\dd\tau+\dd\mathsf R,
 \qquad K\ge0,\qquad \dd\mathsf D\ge0.
\end{equation}
with \(a\in L^1(I)\).  Failure of this local integrability relation is a
total-graph scale coordinate and enters the same \(\Lift\)-owned corona; it
is not used in the formula below.
Put
\begin{equation}
 \label{eq:evaluated-corona-integrating-factor}
 m_I(\tau):=
 \exp\!\left(\int_{\tau_-}^{\tau}\kappa a(s)\,\dd s\right),
 \qquad
 \theta_I(\tau):=
 m_I(\tau)^{-1}\int_{\tau}^{\tau_+}m_I(s)\,\dd s .
\end{equation}
Then \(\theta_I\ge0\), \(\theta_I(\tau_+)=0\), and
\begin{equation}
 \label{eq:evaluated-corona-adjoint-weight}
 \theta_I'+\kappa a\theta_I=-1
\end{equation}
almost everywhere.  Thus the weight is generated by the carrier row and its
finite horizon; it is not an additional multiplier.  The horizon label is
retained.  When prefixes are joined, differences of their terminal weights
remain in the oriented root--leaf coboundary and are not silently identified.
It retains the expanded signed decomposition of \(\dd\mathsf R\), every compatible
interface, the physical carrier, the unnormalized shell cocycle, and every
ordinary or total-graph coordinate used below.  Its covariant production is
\begin{equation}
 \label{eq:evaluated-corona-production}
 \dd\mathcal P_{{\rm cov},I}
 :=K\,\dd\tau+\theta_I\,\dd\mathsf D,
 \qquad
 \dd(\theta_I K)
 =-\dd\mathcal P_{{\rm cov},I}+\theta_I\,\dd\mathsf R.
\end{equation}
If a carrier row, adjoint weight, product, or sign is nonordinary, its least
total-graph value is part of the record rather than an omitted premise.
\end{definition}

\begin{theorem}[Covariant corona transition evaluator]
\label{thm:complete-covariant-corona-evaluator}
For every target-visible amplification tangent, the public continuum grammar
constructs and evaluates the directed family of all records in
\cref{def:evaluated-covariant-corona-record}.  After signed interface gluing
and recursive carrier--polar saturation, exactly one of the following
evaluation states is returned:
\begin{enumerate}[label=\textup{(V\arabic*)}]
\item \emph{finite covariant financing:} for the cumulative unnormalized
  shell current \(\mathsf Q_N=\sum_{j\le N}\mathsf Q_{I_j}\), signed
  saturation generates an exact domination
  \begin{equation}
   \label{eq:covariant-cumulative-financing}
   \mathsf Q_N
   \le \delta\mathscr K_N+\mathsf P_N+\mathsf B_N+\mathsf N_N,
   \qquad \mathsf N_N\in\Null_\Sigma,
  \end{equation}
  in target-current order.  Here \(\mathsf P_N\) is the joined measure from
  \eqref{eq:evaluated-corona-production}, \(\delta\mathscr K_N\) is the
  oriented weighted root--leaf coboundary, and \(\mathsf B_N\) is genuine
  exterior input.  The same generated carrier identities give a uniform
  finite ceiling for the target pairing of the right side.  Since every
  completed crossing has \(\langle\mathsf Q_{I_j},1\rangle=1\), infinitely
  many crossings are excluded;
\item \emph{equality nullity:} every generated covariant production
  vanishes and a finite target-local ordered-coercivity row of
  \cref{thm:dyadic-ordered-coercivity-terminal-rule} makes the equality
  stratum's target projection zero;
\item \emph{target-visible corona successor:} the cumulative same-origin
  execution reaches a primal--polar fixed corona record.  The record contains
  either an unbounded weighted root/exterior financing coordinate whose
  image in the target-current quotient is nonzero, an unbounded coefficient
  or nonordinary graph of a generated ordered-coercivity row, or a
  continuous carrier-polar separator \(\lambda\) satisfying
  \begin{equation}
   \label{eq:evaluated-singular-corona-separator}
   \lambda(Q)>0,\qquad
   \lambda(\delta\mathscr K+\Null_\Sigma)=0,qquad
   \lambda(\dd\mathcal P_{{\rm cov},I}+\dd B)\le0
  \end{equation}
  for every generated covariant production and admitted exterior input.
  No positivity of \(\lambda\) is asserted: its sign is fixed by
  \(\lambda(Q)>0\), while the second and third relations say exactly that it
  belongs to the polar of the oriented finite-financing cone.
\end{enumerate}
The third state includes the full weak equation, chart cocycle, scale
weights, pressure/exterior total graphs, source ancestry, target
normalization, and physical-origin rows.  It is re-entered into the independent
approximation-realization graph.  It becomes a singular-solution terminal only
when all clauses of \cref{def:explicit-singular-mechanism} hold; otherwise its
least failed realization row is the next typed successor.
\end{theorem}

\begin{proof}
Select the nonincreasing record scale by
\cref{thm:canonical-record-scale-selector}.  Enumerate the dual-carrier
grammar of \cref{def:source-adapted-dual-carrier-grammar}; every generated
row has the form \eqref{eq:evaluated-corona-carrier-row} after adjoining its
scale exponent.  Scalar adjoint closure constructs
\eqref{eq:evaluated-corona-adjoint-weight}, and
\cref{prop:covariant-modulation-conjugation} gives
\eqref{eq:evaluated-corona-production} without a sign assumption on \(a\).
Apply \cref{thm:signed-carrier-gluing} before taking positive parts.  Internal
transport, pressure, chart, cutoff, gauge, and exact scale currents cancel;
all remaining terms keep their source owners.

For each finite prefix, scalar adjoint closure is the explicit construction
\eqref{eq:evaluated-corona-integrating-factor}; hence positivity, the terminal
condition, and \eqref{eq:evaluated-corona-production} require no solvability
hypothesis.  If the remaining signed current lies in the target-local
carrier cone, its exact cone identity telescopes on the complex containing
both selected events and their gaps.  Apply that identity to the cumulative
unit shell current, not to the positive part of any individual analytic
term.  If the generated weighted root, exterior, and physical coordinates
are uniformly finite, this is precisely
\eqref{eq:covariant-cumulative-financing} and gives \textup{(V1)}.  No
contact margin is used: the lower side equals the number of completed
normalized crossings.  If all productions vanish, restrict every equation
and target coordinate to the generated equality stratum.  Its zero target
projection is \textup{(V2)}.

Before the equality cell is retained, enumerate the generated dyadic
ordered-coercivity rows.  A finite-coefficient row closes every
zero-production target activity by
\cref{thm:dyadic-ordered-coercivity-terminal-rule}.  An unbounded coefficient
or nonordinary relation is retained in \textup{(V3)} with its source ancestry.

If the weighted root or exterior coordinate is unbounded, project its
cumulative current to the target quotient before terminalization.  A zero
projection is in \(\Null_\Sigma\) and is deleted.  A nonzero projection
retains its least cofinal total-graph value and all same-origin scale ancestry;
this is a target-financing component of \textup{(V3)}.  Thus growth
of \(\theta_I(\tau_-)\) caused only by increasing the horizon is not itself a
singular mechanism.  On the remaining
complement, geometric separation from the closed carrier cone gives
\eqref{eq:evaluated-singular-corona-separator}.  Append the separator, its
adjoint pulls, polarizations, and total-graph values and repeat.  Finite
successors return to the preceding two tests.  For a cofinal execution, form
the cumulative primal--polar join before projection.  By
\cref{thm:continuum-composition-no-creation,%
cor:no-diffuse-primal-polar-composition}, a nonzero limit has a finite
unprocessed increment or is fixed under both grammars.  The former repeats
the evaluation; the latter is terminalized by
\cref{thm:closed-amplification-successor-terminalization}.  On a
contact-seeded same-origin execution it retains hypothetical origin and is
therefore the successor in \textup{(V3)}, not an existence verdict.  An
independent ordinary realization is tested only after this transition.  These
are the exhaustive cone-membership, equality, and strict-separation values.
\end{proof}

\begin{proposition}[Finite-prefix root and exterior accounting]
\label{prop:finite-prefix-root-exterior-accounting}
Let \(\mathcal C_N\) be a finite complex of stopped corona intervals, including
the gaps between selected target crossings.  Construct on each interval the
adjoint weight of
\cref{eq:evaluated-corona-integrating-factor,eq:evaluated-corona-adjoint-weight}
and glue all oriented interfaces before taking positive parts.  Then the
sum of the generated covariant carrier rows has the exact form
\begin{equation}
 \label{eq:finite-prefix-root-exterior-accounting}
 \sum_{I\in\mathcal C_N}\dd(\theta_IK_I)
 =\delta\mathscr K_N-\mathsf P_N+\mathsf B_N+\mathsf N_N+\mathsf E_N,
 \qquad \mathsf N_N\in\Null_\Sigma .
\end{equation}
Here \(\delta\mathscr K_N\) consists of the two exterior roots and every
explicit mismatch of weights at an unglued interface,
\(\mathsf P_N\ge0\) is admitted only through the public physical measure,
\(\mathsf B_N\) is the restriction of the public forcing/boundary measure,
and \(\mathsf E_N\) is the ordered tuple of total-graph failures.  Thus every
finite-financing row contains all root, production, and exterior coordinates.

If their target pairings are bounded along \(N\), they define the
corresponding finite-prefix carrier ceiling.  If one is unbounded, take its
least rational level-crossing prefix and factor that finite readout backward
through its complete joint frontier.  Its nonzero conditional increment has
the owner of the root, scale, horizon, or exterior constructor that created
it and is evaluated by
\cref{thm:finite-prefix-source-profile-carrier-duality}.  Consequently an
unbounded adjoint weight or exterior term is expanded into a source-owned
physical-valuation cell; it is never discarded as a failed estimate.
\end{proposition}

\begin{proof}
Sum \eqref{eq:evaluated-corona-production} over the finite complex.  Opposite
oriented copies of a compatible internal face cancel.  The remaining endpoint
values are exactly \(\delta\mathscr K_N\); public exterior faces give
\(\mathsf B_N\), target-null signed terms give \(\mathsf N_N\), and the least
nonordinary operations give \(\mathsf E_N\).  This proves
\eqref{eq:finite-prefix-root-exterior-accounting}.  Bounded pairings give a
finite cylinder row.  In the unbounded case the first level-crossing prefix
is finite, so prospective backward factorization applies without a terminal
limit.  Constructor charging and finite-prefix duality give the last claim.
\end{proof}

\begin{corollary}[Sound structural transition method]
\label{cor:complete-structural-regularity-method}
For every publicly presented continuum equation and physical stratum, the
method
\[
 \begin{aligned}
 \text{prospective contact}
 &\to \text{five-source cumulative owner}
 \to \text{physical amplification}\\
 &\to \text{zero-production tangent}
 \to \text{covariant corona evaluation}
 \end{aligned}
\]
returns target avoidance by \textup{(V1)}--\textup{(V2)}, or a target-visible
successor \textup{(V3)} that is rejoined with its complete same-origin history
and evaluated in the reachable contact cell.  Every operation is generated
from the public weak equation, carrier, target shells, and their total graphs;
no source-transverse estimate or semantic compactness statement is accepted
as an input.
\end{corollary}

\begin{proof}
Prospective accountability and source exhaustion give the first two arrows.
\Cref{thm:exact-compositional-physical-accounting,%
thm:prospective-source-amplification-extraction} gives the third,
\cref{thm:amplification-zero-production-tangent} gives the fourth, and
\cref{thm:complete-covariant-corona-evaluator} gives the next evaluation
state.  Terminal typing is \cref{cor:binary-continuum-terminal-theorem}.
\end{proof}

\begin{corollary}[Continuum no-uncharged-amplification rule]
\label{cor:continuum-no-uncharged-amplification}
For every finite-budget contact, prospective source extraction followed by
target normalization gives the mandatory chain
\begin{equation*}
 \mathrm{Contact}
 \longrightarrow \mathsf{Amp}_{c_*}
 \longrightarrow \mathsf{ZP}_{c_*}
 \longrightarrow
 \left\{
 \begin{array}{l}
   \Null_\Sigma,\\
   \text{ordinary equality state},\\
   \text{typed first failure}.
 \end{array}\right.
\end{equation*}
This is applied before recurrence, averaging, compact realization, or
terminal evaluation.  Hence neither an infinite amplification ratio nor a
physical-singular numerator can be used as an unevaluated regularity
obstruction.
\end{corollary}

\begin{proof}
Combine \cref{thm:prospective-source-amplification-extraction} and
\cref{thm:amplification-zero-production-tangent}.
\end{proof}

\begin{corollary}[No unowned prospective physical amplification]
\label{cor:no-unowned-prospective-physical-amplification}
For a finite-physical-budget contact, exactly one of the following occurs:
the target is reached after finitely many shell crossings; or one fixed
source channel has the unbounded pre-contact amplification
\eqref{eq:sourcewise-unbounded-amplification}; or one fixed source channel
has the positive physical-singular pre-contact component
\eqref{eq:sourcewise-singular-target-mass}.  The latter two outputs are the
continuum analogue of the pre-hit selector carrying excess advantage.  They
are source-resolved prospective mechanisms.  Singular existence additionally
requires the independent realization in
\cref{def:explicit-singular-mechanism}.
\end{corollary}

\begin{proof}
Combine \cref{thm:exact-compositional-physical-accounting} with
\cref{thm:prospective-source-amplification-extraction}.  Cofinal contact
excludes the finite-crossing alternative.
\end{proof}

\begin{proposition}[Finite physical budget does not supply transition capacity]
\label{prop:no-free-physical-transition-capacity}
There is no equation-independent theorem deriving
\(\mathsf t^{\Phys}_\Sigma(E)<\infty\) from \(E<\infty\).  Indeed, the
scalar flow
\begin{equation}
 \label{eq:finite-budget-singular-scalar-flow}
 \dot x=x^2,
 \qquad x(0)=1,
 \qquad K(x)=x^{-1}
\end{equation}
has \(K(x(t))=1-t\), reaches the compactified singular target \(x=+\infty\)
at \(t=1\), and obeys the exact finite physical ledger
\[
 K(x(T))+\mu_{\Phys}([0,T])=K(x(0)),
 \qquad \mu_{\Phys}=\dd t,
 \qquad \mu_{\Phys}([0,1))=1.
\]
Its crossings of \(x=2^j\) have prices
\(p_j=2^{-j}-2^{-j-1}\), so their target amplification tends to infinity.
The target spine is entirely \(\Caus\)-owned.  Thus the continuum algebra
correctly returns a realized singular terminal; declaring finite transition
capacity from the physical ceiling would incorrectly erase it.
\end{proposition}

\begin{proof}
The solution is \(x(t)=(1-t)^{-1}\).  Direct differentiation gives
\(\frac{\dd}{\dd t}K(x(t))=-1\), proving the ledger.  The first hitting
time of \(2^j\) is \(1-2^{-j}\), so the physical price between consecutive
hits is the displayed \(p_j\), while every normalized shell crossing
contributes one to \(\Gamma_\Sigma\).  Hence \(p_j^{-1}\to\infty\).
The vector field is a fixed directed drift, so its ancestry is \(\Caus\).
\end{proof}

\begin{corollary}[Composition, refinement, and scale covariance]
\label{cor:composition-refinement-scale-covariance}
The decomposition \eqref{eq:target-progress-lebesgue-decomposition} commutes
with finite composition, countable refinement of the occurrence partition,
insertion of inactive gaps, and every injective chart relabelling.  Under a
noninjective quotient, the numerator, physical denominator, and fiber
multiplicity are pushed forward jointly before their Lebesgue decomposition.
Consequently neither averaging, refinement, concentration, nor scale collapse
can turn unbounded amplification or physical-singular target progress into a
zero-price record.
\end{corollary}

\begin{proof}
All four operations are pushforwards or restrictions of the pair of measures
\((\Gamma_\Sigma,\mu_{\Phys})\).  Uniqueness of Lebesgue decomposition proves
the assertion.  Keeping fiber multiplicity is exactly the unnormalized
cocycle rule of \eqref{eq:unnormalized-target-cocycle}.
\end{proof}

\begin{definition}[Universal target-aligned amplification cell]
\label{def:universal-target-amplification-cell}
Let \(Q\) be any source-owned target-current readout in a universal
cumulative record and let \(\mathscr D_{\mathfrak b}^{\Phys}\) be the
target-local physical carrier cone of
\cref{def:target-local-carrier-polar-successor}.  Evaluate all terms on the
same selected finite event or finite prefix.  Define
the target amplification relation by the total graph of
\begin{equation}
 \label{eq:universal-target-amplification}
 \operatorname{Amp}_{\Sigma}(Q;P,K,B)
 :=\frac{(Q)_+}
 {P+\lvert\delta K\rvert+B^+},
\end{equation}
with the conventions encoded as separate graph cells when the denominator
vanishes or a term is nonordinary.  Here \(P\) is admitted physical
production, \(\delta K\) is endpoint storage before telescoping, and \(B^+\)
is true positive exterior input.  The ratio is only a diagnostic coordinate:
physical charging is authorized solely by cone membership
\(Q\in\mathscr D_{\mathfrak b}^{\Phys}\).

The grammar adjoins this graph for every generated target owner, not merely
for scale or energy.  Finite amplification records a local domination row.
Unbounded or zero-denominator amplification is a target-active polarity cell
with the complete expanded ancestry of \(Q,P,K,B\); it is recursively
processed by carrier--polar exhaustion.  It is never interpreted as free
progress, target-nullity, or a residual source.
\end{definition}

\begin{theorem}[Universal saturated-composition accountability]
\label{thm:continuum-composition-no-creation}
Let \(q=h(J)\) be any finite composite generated by the public continuum
grammar, where \(J\) is the joint value of every immediate input to the
terminal constructor.  Expand the derivations of those inputs before
evaluating \(h\), and form the universal cumulative record of
\cref{def:cumulative-joint-continuum-record}.  The following statements hold.
\begin{enumerate}[label=\textup{(J\arabic*)}]
\item If \(q=h(\mathcal J_n)\) is any generated finite composite and
  \(P_q\) denotes conditional expectation onto \(\sigma(q)\), then
  \begin{equation}
   \label{eq:continuum-composition-projection-ceiling}
   P_qP_n=P_nP_q=P_q,
   \qquad
   \lVert P_qy_\Phi\rVert_2
   \le \lVert P_ny_\Phi\rVert_2
  \end{equation}
  for every bounded centered shell contrast \(y_\Phi\).
\item The spaces \(L^2(\mathcal A_n,\nu)\) increase and
  \begin{equation}
   \label{eq:continuum-joint-projection-limit}
   P_ny_\Phi\longrightarrow P_\infty y_\Phi
   \quad\text{strongly in }L^2(\nu),
   \qquad
   \lVert P_ny_\Phi\rVert_2^2\uparrow
   \lVert P_\infty y_\Phi\rVert_2^2.
  \end{equation}
  Hence a target-aligned component of the joint limit has a
  target-aligned finite joint prefix.  If every finite joint prefix is
  target-null, then the joint limit is target-null.
  More precisely, with
  \begin{equation}
   \label{eq:continuum-joint-martingale-increments}
   d_0:=P_0y_\Phi,
   \qquad d_n:=(P_n-P_{n-1})y_\Phi\quad(n\ge1),
  \end{equation}
  the \(d_n\) are pairwise orthogonal and
  \begin{equation}
   \label{eq:continuum-joint-increment-accounting}
   P_\infty y_\Phi=\sum_{n\ge0}d_n,
   \qquad
   \lVert P_\infty y_\Phi\rVert_2^2
   =\sum_{n\ge0}\lVert d_n\rVert_2^2.
  \end{equation}
  Thus no positive detection threshold is imposed on a joint contribution:
  an arbitrary countable family of small contributions is retained by its
  exact cumulative sum.
\item Finite composition and graph-complete passage to
  \(\mathcal J_\infty\) preserve the source-resolved identity
  \eqref{eq:unnormalized-target-cocycle}.  Every nonzero target component of
  a mixed product, limit, or total-graph value is therefore reprojected into
  \(\Geom,\Caus,\Abs,\Lift,\Bdry\) with the union of the ancestries of the
  joint inputs.  If \(d_n\ne0\), a generated finite cylinder readout of
  \(\mathcal J_n\) has nonzero pairing with \(d_n\); the complete derivation
  of that detecting readout supplies the responsible mixed input word.  No
  sixth source is created by composition or by the directed limit.  An
  increment is \emph{processed} once such a detecting word and its inherited
  five-source ancestry have been adjoined to the cumulative record; its
  magnitude is irrelevant.
\item On a prefix containing \(N\) completed normalized shell crossings,
  \(A_{\Phi,N}=N\).  In particular, division by \(N\), return time, scale
  depth, or the mass of an empirical recurrence law is not an admissible
  target-accounting operation.  Such a normalization may select a recurrent
  support, but the selected support is rejoined with
  \eqref{eq:unnormalized-target-cocycle} before target projection.
\item Every target-visible cross term created by nonlinear combination,
  localization, conditioning, stopping, chart change, scaling, quotient,
  weak limiting, trace, boundary passage, or recurrence is a finite
  constructor vertex in \(\operatorname{Exp}(q)\), or a total-graph boundary
  value of such a vertex.  It inherits the union of the source ancestries of
  its inputs.  Thus the fact that every input is individually target-null
  does not authorize deletion of their joint record: joint synergy is tested
  before the residual quotient.
\item Every physical quantity invoked by a composite remains attached to its
  owner and to the exact target-local balance or domination graph that uses
  it.  Finite additivity, countable additivity, a conservation law, or a
  positive sign does not by itself pay for target progress.  The composite is
  charged only after its cumulative target owner lies in
  \(\mathscr D_{\mathfrak b}^{\Phys}\).  Otherwise its universal
  amplification graph \eqref{eq:universal-target-amplification} produces the
  next finite polarity or carrier--polar successor.
\item The construction commutes with directed same-origin accumulation.  If
  a finite or countable family of regular, target-null, bounded, vanishing,
  or subthreshold pieces has target-active composite output, then a finite
  expanded joint prefix has a nonzero target increment.  That increment is
  re-entered into the five-source compiler with its complete ancestry,
  physical owners, normalization denominators, and unnormalized cocycle.
  If no such prefix exists, the composite limit is target-null.
\end{enumerate}
Thus no generated continuum construction can turn regular pieces into
unaccounted singularity-forming structure.  It may create target alignment
through their joint relation, exactly as an XOR can in Part~I, but that
alignment belongs to the saturated expanded derivation and must be bounded
as an existing five-source profile cell.  There is no separate loophole for
time accumulation, scale collapse, nonlinear interaction, weak convergence,
normalization, localization, recurrence, concentration, oscillation, tails,
or boundary passage.
\end{theorem}

\begin{proof}
Since \(q\) factors through \(\mathcal J_n\), one has
\(\sigma(q)\subseteq\mathcal A_n\).  Conditional expectations onto nested
closed subspaces give the two projection identities in
\eqref{eq:continuum-composition-projection-ceiling}; the norm bound follows
from the orthogonal Pythagorean decomposition.  The martingale convergence
theorem for the increasing sigma algebras \((\mathcal A_n)\) gives
\eqref{eq:continuum-joint-projection-limit}.  Its orthogonal-difference
decomposition gives
\eqref{eq:continuum-joint-martingale-increments}--
\eqref{eq:continuum-joint-increment-accounting}.  A positive limiting norm
has a positive finite-prefix norm, while a countable collection of small
increments is retained by the convergent orthogonal sum.  This proves
\textup{(J2)} without a detection threshold.

At finite depth, prospective accountability and the five-source normal form
give \eqref{eq:unnormalized-target-cocycle}.  The join records all mixed
coordinates before projection, so a target component selected by a composite
belongs to the joint algebra rather than to the residual of any individual
input.  Source-preserving totalization assigns a nonordinary mixed operation
to its typed graph value with the union of the input ancestries.  If
\(d_n\ne0\), density of the generated bounded cylinder readouts in
\(L^2(\mathcal A_n,\nu)\) supplies one with nonzero pairing against \(d_n\).
Because \(d_n\perp L^2(\mathcal A_{n-1},\nu)\), no readout of the old prefix
can supply this pairing; the detecting cylinder therefore identifies a finite
new joint word.  Passing to
the graph-complete limit preserves every fixed finite current identity; if a
current fails to converge in variation, its boundary owner is the retained
\(\Bdry\)-headed value.  This proves \textup{(J3)}.

Each completed normalized crossing contributes one to the prospective shell
cocycle, and the same-origin join adds these occurrences without identifying
them.  Hence \(A_{\Phi,N}=N\).  A normalized empirical law is a positive
scalar rescaling used for compact support selection.  Keeping it as a
separate coordinate and rejoining it with the cumulative cocycle proves
\textup{(J4)}.

Every operation listed in \textup{(J5)} is a constructor in the expanded
grammar.  At an ordinary value, structural induction assigns the union of
the input ancestries; at a nonordinary value, totalization assigns the same
ancestry to the graph-boundary successor.  This proves \textup{(J5)}.
The public physical quantities are owner-tagged coordinates of the joint
record.  The definition of target-aligned physical charge permits their use
only through the target-local carrier cone.  Separation from that cone is
exactly the carrier--polar successor, while totalization of
\eqref{eq:universal-target-amplification} gives all finite, infinite, and
zero-denominator values.  This proves \textup{(J6)} without inferring
domination from positivity.  Finally apply \textup{(J2)} to the increasing
sigma algebras of the expanded joint prefixes.  A nonzero limiting projection
has a nonzero orthogonal finite-prefix increment, and \textup{(J3)} assigns
its mixed source word.  This proves \textup{(J7)}.  The argument depends only
on expanded derivation, not on the analytic name of the constructor.
\end{proof}

\begin{theorem}[Infinite-chain no-emergence trichotomy]
\label{thm:infinite-chain-no-emergence-trichotomy}
For every infinite same-origin successor chain, form its cumulative joint
record, close it under every generated mixed operation, and then reproject to
the target quotient.  Exactly one of the following occurs:
\begin{enumerate}[label=\textup{(N\arabic*)}]
\item the cumulative target projection is zero;
\item a least unprocessed finite joint increment exists and is the next
  proper successor;
\item every finite increment is processed and the cumulative record is fixed
  under all generated primal and polar operations.
\end{enumerate}
In particular, a nonzero target component cannot emerge only at the infinite
limit.
\end{theorem}

\begin{proof}
Apply the orthogonal increment identity
\eqref{eq:continuum-joint-increment-accounting} of
\cref{thm:continuum-composition-no-creation} to the increasing finite joint
cylinder algebras.  If the limiting target projection vanishes, this is
\textup{(N1)}.  Otherwise some finite increment has nonzero pairing with a
generated cylinder readout; choosing the first unprocessed one gives
\textup{(N2)} with its mixed source word.  If all such increments have
already been processed, totalize every remaining primal, graph, and
carrier--polar relation.  A first failed relation is again a finite
\textup{(N2)} successor.  Only when none exists is the record fixed, giving
\textup{(N3)}.  The alternatives are disjoint by the fixed least-coordinate
order.
\end{proof}

\begin{corollary}[One rule for all continuum accumulation]
\label{cor:universal-continuum-accumulation-rule}
The following are not separate regularity principles: accumulation over
time; transfer over scales; addition over source occurrences; spatial
localization and partition gluing; nonlinear interaction; conditioning or
stopping; normalization by amplitude, mass, return number, or scale;
compactness loss, concentration, and oscillation; weak or graph limits;
recurrence; tail escape; and boundary or terminal passage.  Each is evaluated
by the same protocol:
\begin{equation}
 \label{eq:universal-continuum-accumulation-protocol}
 \boxed{
 \operatorname{Exp}\ \longrightarrow\
 \operatorname{Join}\ \longrightarrow\
 \operatorname{Sat}_{\mathsf G_\Sigma}\ \longrightarrow\
 q_\Sigma\ \longrightarrow\
 \begin{cases}
   \Null_\Sigma,\\
   \text{five-source target owner},\\
   \text{typed total-graph successor}.
 \end{cases}}
\end{equation}
For a nonnull owner, recursive carrier--polar exhaustion is mandatory before
averaging, realization, or another composition.  Hence no theorem may infer
regularity merely because every constituent is regular, small, bounded,
vanishing, summable, or target-null in isolation.
\end{corollary}

\begin{theorem}[Universal continuum saturated-profile accountability]
\label{thm:universal-continuum-saturated-profile-accountability}
Let \(q\) be any target-directed readout produced at any finite or directed
stage of a legal same-origin continuum evolution, and let \(Q_\Sigma(q)\) be
its positive target progress after universal expanded-derivation saturation.
Then
\begin{equation}
 \label{eq:universal-continuum-profile-accountability}
 Q_\Sigma(q)
 \le
 \sum_{c\in\Channels}
 \sup_{R\in\operatorname{Sat}(\operatorname{Exp}(q))}
   \bigl(\operatorname{Prog}_\Sigma^c(R)\bigr)_+
 \le \mathsf T_\Sigma^{\Phys,\Sat}(\Budget).
\end{equation}
For \(N\) completed shell crossings the same statement is applied to the
unnormalized cumulative readout and gives
\begin{equation}
 \label{eq:universal-continuum-cumulative-accountability}
 N\le\sum_{c\in\Channels}(Q^c_{\Phi,N})_+.
\end{equation}
Every target-directed effect of composition, including an effect absent from
all individual inputs, therefore appears in the same saturated profile that
is bounded in a regularity proof.  A proof that bounds only primitive or
termwise profiles is not a continuum structural proof.
\end{theorem}

\begin{proof}
Apply \cref{thm:continuum-composition-no-creation} to the expanded derivation
of \(q\).  Its target-null component vanishes, while every nonzero ordinary
or graph-boundary component has an inherited five-source owner.  The triangle
inequality for the five signed source contributions gives the first
inequality, and saturation of the profile over expanded composite records
gives the second.  For completed crossings retain
\(A_{\Phi,N}=N\) before every normalization and sum the source-resolved
identity; this is \eqref{eq:universal-continuum-cumulative-accountability}.
No constituent-wise lower bound is used.
\end{proof}

\begin{corollary}[No magical singularity from regular constituents]
\label{cor:no-magical-singularity-from-regular-constituents}
Suppose a family of continuum constituents is regular, bounded, smooth,
compact, individually target-null, or individually of zero target price.
These predicates imply nothing about a generated composite until its expanded
joint record is evaluated.  After that evaluation exactly one of the
following holds: the composite is target-null; its target owner is charged by
the public physical cone; or a source-resolved amplification successor is
generated.  Repeating the construction, at any countable depth, produces no
fourth case and no target-aligned residual.
\end{corollary}

\begin{corollary}[All previous accumulation routes are instances]
\label{cor:all-accumulation-routes-are-instances}
The unnormalized shell ledger, cumulative-before-circulation rule,
no-diffuse-composition rule, no-diffuse-primal--polar rule, covariant
microscope, entropy cocycle, signed interface gluing, recurrent return
construction, concentration/oscillation lift, tail completion, and terminal
boundary owner are restrictions of
\cref{thm:continuum-composition-no-creation,thm:universal-physical-quantity-ownership,thm:recursive-carrier-polar-exhaustion}
to particular constructor subgrammars.  None is an independent assumption or
an additional closure principle.  Any future continuum constructor is
handled by adjoining its total graph to the public grammar and applying the
same expanded-derivation theorem.
\end{corollary}

\begin{corollary}[No limit-created target alignment]
\label{cor:no-limit-created-target-alignment}
An infinite compatible compiler execution is not a terminal branch.  Form
\(\mathcal J_\infty\) and reapply the five-channel projection.  Exactly one
of the following occurs: the joint record is target-null; a finite joint
prefix has a least unprocessed target increment \(d_n\), whose complete
five-channel mixed ancestry is the next saturated compiler state; or every
increment is processed and the joint record is a fixed point of every
generated joint, source, graph, carrier, and target-cyclic operation while
retaining its nonzero unnormalized shell cocycle.  The last case is a
joint-saturated structural atom and is defined below.  A compatible inverse
limit or an ordinary path is never terminal solely because it exists.
\end{corollary}

\begin{corollary}[No diffuse composition loophole]
\label{cor:no-diffuse-composition-loophole}
Let one same-origin execution complete \(N\) normalized target-shell
crossings, and classify every mixed joint increment before forming the
residual.  If
\[
 Q^c_{\Phi,N}:=\sum_{j<N}Q^c_{\Phi,j}
 \qquad(c\in\Channels)
\]
denotes its cumulative source-resolved shell current, then
\begin{equation}
 \label{eq:cumulative-five-source-shell-accounting}
 N=A_{\Phi,N}=\sum_{c\in\Channels}Q^c_{\Phi,N},
 \qquad
 \sum_{c\in\Channels}(Q^c_{\Phi,N})_+\ge N,
 \qquad
 \max_{c\in\Channels}(Q^c_{\Phi,N})_+\ge\frac N5.
\end{equation}
Consequently there are a source \(c_*\) and a cofinal subsequence \(N_k\)
such that
\[
 (Q^{c_*}_{\Phi,N_k})_+\ge \frac{N_k}{5}.
\]
This conclusion has no per-occurrence detection threshold.  Individually
target-null inputs may have a target-aligned joint increment, and infinitely
many arbitrarily small processed increments may have nonzero total; both are
already included in the left side of
\eqref{eq:cumulative-five-source-shell-accounting}.  If every cumulative
five-source current vanished, no positive number of crossings could occur.
Thus diffuse composition cannot create a target-aligned \(J_{\rm res}\).
\end{corollary}

\begin{proof}
For each completed shell, prospective accountability gives
\(1=\sum_cQ^c_{\Phi,j}\) after that prefix has been jointly saturated.
\Cref{thm:continuum-composition-no-creation} assigns every nonzero mixed
increment to its complete inherited source word before the residual is
formed.  Sum the \(N\) identities.  The first equality in
\eqref{eq:cumulative-five-source-shell-accounting} follows, and the two
inequalities follow from \(x\le x_+\) and the five-element pigeonhole
principle.  One channel attains the last bound along a cofinal subsequence.
No lower bound on an individual increment was used.  The identity uses the
unnormalized cocycle, so normalization of a return law cannot change it.
\end{proof}

\begin{definition}[Unnormalized finite-horizon target-current problem]
\label{def:unnormalized-finite-horizon-current}
Let \(X\) be a closed target-active recurrent core and let
\(E_X\rightrightarrows X\) be its compact same-origin successor graph, with
source and range maps \(\pi_-\) and \(\pi_+\).  Retain the nonnegative
physical valuation \(p_{\Phys}\), the positive target-return coordinate
\(a_\Sigma\ge\eta\), and the signed source coordinates
\(q_\Sigma^c\), where
\(a_\Sigma=\sum_cq_\Sigma^c\).  For \(N\ge1\), let
\[
 E_X^{[N]}:=\{(e_0,\ldots,e_{N-1}):
       \pi_+(e_j)=\pi_-(e_{j+1})\text{ for }j<N-1\}
\]
with all active ideal, origin, and total-graph relations imposed at every
layer.  Define the unnormalized cumulative price
\begin{equation}
 \label{eq:unnormalized-finite-horizon-price}
 \Delta_X(N)
 :=\inf_{\boldsymbol e\in E_X^{[N]}}
       \sum_{j<N}p_{\Phys}(e_j)\in[0,+\infty],
\end{equation}
where the value is \(+\infty\) when the layered path space is empty.  The
associated prefix coordinates are
\begin{equation}
 \label{eq:unnormalized-finite-horizon-ledger}
 A_N(\boldsymbol e):=\sum_{j<N}a_\Sigma(e_j),\qquad
 Q_N^c(\boldsymbol e):=\sum_{j<N}q_\Sigma^c(e_j),\qquad
 \Pi_N(\boldsymbol e):=\sum_{j<N}p_{\Phys}(e_j).
\end{equation}
No probability normalization occurs in this definition.
\end{definition}

\begin{theorem}[Cumulative carrier duality]
\label{thm:cumulative-carrier-duality}
Let \(X\) be as in
\cref{def:unnormalized-finite-horizon-current}.  For continuous cylinder
lower envelopes \(p_m\uparrow p_{\Phys}\),
\begin{equation}
 \label{eq:cumulative-carrier-dual}
 \Delta_X(N)
 =\sup_{m,(V_j),(\delta_j)}\left\{
       \sum_{j<N}\delta_j:
 \begin{array}{l}
 V_0=V_N=0,\\
 p_m(e)+V_j(\pi_-e)-V_{j+1}(\pi_+e)-\delta_j
 \in\overline{\mathcal K_X^0}\\
 \text{on the layer-}j\text{ edge graph},\quad j<N
 \end{array}\right\}.
\end{equation}
Here \(\mathcal K_X^0\) is the local ordered equation cone before successor
coboundaries are imposed.  Every finite dual row telescopes on an actual
prefix and gives
\begin{equation}
 \label{eq:cumulative-carrier-prefix-bound}
 \Pi_N(\boldsymbol e)\ge\sum_{j<N}\delta_j.
\end{equation}
Consequently, on a physical stratum of total mass at most \(\Budget\), exactly
one of the following occurs.
\begin{enumerate}[label=\textup{(U\arabic*)}]
\item For some \(N\), a finite cylinder dual row has
\(\sum_{j<N}\delta_j>\Budget\).  No same-origin contact cascade in the
stratum has \(N\) returns.
\item The values \(\Delta_X(N)\) remain bounded by \(\Budget\).  The
cumulative joint compactification contains an infinite same-origin successor
path satisfying
\begin{equation}
 \label{eq:bounded-price-unbounded-target-current}
 \sup_N\Pi_N\le\Budget,
 \qquad A_N\ge\eta N,
 \qquad
 (Q_{N_k}^{c_*})_+\ge\frac{\eta N_k}{5}\longrightarrow\infty
\end{equation}
for a channel \(c_*\) and a cofinal subsequence.  This entire unnormalized
path ledger is the next compiler state.
\end{enumerate}
The alternatives concern cumulative expenditure.  They remain valid when
every one-step price tends to zero and when every normalized invariant law
has zero physical-expenditure density.
\end{theorem}

\begin{proof}
For fixed \(m\) and \(N\), minimize the linear functional
\(\sum_{j<N}\int p_m\,\dd\gamma_j\) over positive layer-edge measures of
mass one whose adjacent range and source marginals agree and which satisfy
the active ordered relations.  This is a linear program on a compact convex
set.  Separation gives its dual: the marginal multipliers are the cylinder
potentials \(V_j\), the layer lower bounds are \(\delta_j\), and the endpoint
normalization is \(V_0=V_N=0\).  Summing the layer inequalities cancels all
intermediate potentials and proves
\eqref{eq:cumulative-carrier-prefix-bound}.  Compactness and monotone
convergence of the lower envelopes give
\eqref{eq:cumulative-carrier-dual} for \(p_{\Phys}\).

For completeness, the layered marginal constraints admit the usual finite
superposition: disintegrate \(\gamma_0\) over its range marginal, then
successively disintegrate \(\gamma_j\) over the common source marginal.
This constructs a probability measure on \(E_X^{[N]}\) having the prescribed
layer marginals.  Conversely, every probability measure on
\(E_X^{[N]}\) has such marginals.  Hence the relaxed linear-program value is
the average of path costs and its infimum equals the deterministic path
infimum in \eqref{eq:unnormalized-finite-horizon-price}; no relaxation gap is
being hidden in \eqref{eq:cumulative-carrier-dual}.

If the first alternative holds, the displayed prefix inequality contradicts
the public physical ceiling.  Otherwise choose an \(N\)-edge prefix of cost
at most \(\Budget+N^{-1}\) in the viable recurrent core and extend it in the
closed successor graph.  Compactness of the graph-complete path space and a
diagonal subsequence give an infinite path.  Lower semicontinuity on every
fixed prefix gives \(\Pi_N\le\Budget\).  The target floor gives
\(A_N\ge\eta N\), while
\(A_N=\sum_cQ_N^c\).  Positive parts and the five-channel pigeonhole
argument produce \(c_*\) and the cofinal subsequence in
\eqref{eq:bounded-price-unbounded-target-current}.  No average has been
taken.
\end{proof}

\begin{theorem}[Finite-horizon target-current duality]
\label{thm:finite-horizon-target-current-duality}
For the \(N\)-layer graph of
\cref{def:unnormalized-finite-horizon-current}, define the sharp finite
horizon price per unit target amplitude by
\[
 \delta_X^{[N]}
 :=\inf\left\{
 \frac{\Pi_N(\boldsymbol e)}{A_N(\boldsymbol e)}:
 \boldsymbol e\in E_X^{[N]},\ A_N(\boldsymbol e)>0\right\},
\]
with value \(+\infty\) when the admissible set is empty.  Then
\[
 \delta_X^{[N]}
 =\sup\left\{\delta:
 p_{\Phys}+V_j\circ\pi_--V_{j+1}\circ\pi_+
 -\delta a_\Sigma\in\overline{\mathcal K_X^0}
 \text{ on every layer }j<N\right\},
\]
where the generated cylinder potentials have fixed endpoint values.  If
\(0<\delta<\delta_X^{[N]}\), a finite rational-cylinder row gives, for every
actual \(N\)-edge prefix,
\[
 \Pi_N\ge\delta A_N-2\max_{0\le j\le N}\lVert V_j\rVert_\infty.
\]
\end{theorem}

\begin{proof}
Minimize the physical valuation over the compact convex set of layered edge
measures satisfying the active relations and the normalization
\(\sum_{j<N}\int a_\Sigma\,\dd\gamma_j=1\).  Its Lagrange multiplier is
\(\delta\); the marginal multipliers are the layer potentials \(V_j\).
Finite-dimensional separation on each cylinder and monotone passage through
the generated lower envelopes yield the displayed dual equality, exactly as
in \cref{thm:cumulative-carrier-duality}.  Superposition identifies the
relaxed value with the path infimum.  Approximate the continuous dual
potentials by rational cylinders and absorb the error using the Archimedean
order unit.  Summing the resulting inequalities along a prefix telescopes
the intermediate potentials, leaving at most the two endpoint bounds.
\end{proof}

\begin{corollary}[Cumulative test precedes circulation]
\label{cor:cumulative-before-circulation}
Every recurrent target-active compiler branch first applies
\cref{thm:cumulative-carrier-duality} to its unnormalized layered successor
graph.  Alternative \textup{(U1)} closes the branch by a finite physical
identity.  In alternative \textup{(U2)}, a normalized circulation may select
recurrent support, but every subsequent kernel, target, and realization
operation acts on the join of that selector with the ledger
\eqref{eq:bounded-price-unbounded-target-current}.  Thus circulation is a
support reduction of an already evaluated cumulative problem; it is never
the physical-price evaluator for the original cascade.
\end{corollary}

\begin{definition}[Target-current circulation problem]
\label{def:target-current-circulation-problem}
Let \(E_{\Sigma,\mathfrak b}\) be the compact total-graph edge space of
\(\Spine_{\Sigma,\mathfrak b}\), with source and range maps
\(\pi_-,\pi_+\).  The selected signed-dyadic orientation defines a bounded
nonnegative target-return coordinate \(a_\Sigma\); normalize each active
return cell so that \(a_\Sigma\ge\eta>0\).  Write \(p_{\Phys}\ge0\) for the
lower-semicontinuous physical valuation obtained from the local
carrier-domination graph, and write \(q_\Sigma^c\) for its source-resolved
target-return coordinates, so that \(a_\Sigma=\sum_cq_\Sigma^c\).  Along an actual return sequence
retain the unnormalized occurrence measure, target cocycle, physical
expenditure, and source currents
\begin{equation}
 \label{eq:unnormalized-return-occurrence-measure}
 \Gamma_N:=\sum_{j<N}\delta_{e_j},
 \qquad
 A_{\Sigma,N}:=\int a_\Sigma\,\dd\Gamma_N,
 \qquad
 \Pi_N:=\int p_{\Phys}\,\dd\Gamma_N,
 \qquad
 Q_N^c:=\int q_\Sigma^c\,\dd\Gamma_N.
\end{equation}
When \(A_{\Sigma,N}>0\), define the empirical support selector
\begin{equation}
 \label{eq:return-density-selector}
 \gamma_N:=A_{\Sigma,N}^{-1}\Gamma_N.
\end{equation}
This is a change of scale in the observer law, not a quotient of the
structural record: \(A_{\Sigma,N}\), \(\Pi_N\), and every \(Q_N^c\) remain
separate coordinates.
The continuous cylinder lower envelopes of \(p_{\Phys}\),
\(p_N\uparrow p_{\Phys}\) are generated by the totalized carrier coordinate;
an infinite endpoint is routed directly to physical charge.  A
\emph{target-current circulation} is a finite positive Borel support-selector
measure
\(\gamma\) on \(E_{\Sigma,\mathfrak b}\) satisfying
\begin{equation}
 \label{eq:target-current-circulation-constraints}
 (\pi_-)_\#\gamma=(\pi_+)_\#\gamma,
 \qquad
 \int a_\Sigma\,\dd\gamma=1,
\end{equation}
and all ordered ideal, source, origin, and total-graph relations of the active
signature.  It is stored jointly with
\((\Gamma_N,A_{\Sigma,N},\Pi_N,(Q_N^c)_c)_N\); the normalization in
\eqref{eq:target-current-circulation-constraints} never replaces those
cumulative coordinates.  Its sharp recurrent physical price is
\begin{equation}
 \label{eq:sharp-recurrent-circulation-price}
 \delta^{\rm cyc}_{\Sigma,\mathfrak b}
 :=\inf_{\gamma}
       \int p_{\Phys}\,\dd\gamma
 \in[0,+\infty],
\end{equation}
with value \(+\infty\) when the circulation set is empty.
\end{definition}

\begin{theorem}[Multiplicity restoration for normalized circulations]
\label{thm:circulation-multiplicity-restoration}
For every actual return prefix with \(A_{\Sigma,N}>0\), normalization is
exactly reversible in the augmented structural record:
\begin{equation}
 \label{eq:circulation-denormalization-identities}
 \Gamma_N=A_{\Sigma,N}\gamma_N,
 \qquad
 \Pi_N=A_{\Sigma,N}\int p_{\Phys}\,\dd\gamma_N,
 \qquad
 Q_N^c=A_{\Sigma,N}\int q_\Sigma^c\,\dd\gamma_N,
 \qquad
 \sum_cQ_N^c=A_{\Sigma,N}.
\end{equation}
Suppose \(A_{\Sigma,N}\to\infty\) while \(\sup_N\Pi_N<\infty\).  Every
weak-* limit \(\gamma\) of the selectors satisfies
\(\int p_{\Phys}\,\dd\gamma=0\), but this conclusion has only the following
meaning: \(\gamma\) selects the zero-density recurrent support.  It implies
neither \(\Pi_N=0\) nor \(Q_N^c=0\), nor may it replace the occurrence
sequence.  Indeed, on a cofinal subsequence there is a fixed
\(c_*\in\Channels\) with
\begin{equation}
 \label{eq:circulation-restored-source-growth}
 (Q_{N_k}^{c_*})_+\ge \frac{A_{\Sigma,N_k}}5\longrightarrow\infty.
\end{equation}
Consequently a zero value of
\(\delta^{\rm cyc}_{\Sigma,\mathfrak b}\) means only that the public carrier
has no uniform positive linear lower bound per unit return on that support.
It never means zero accumulated expenditure or zero accumulated
target-aligned structure.
\end{theorem}

\begin{proof}
The first three identities are substitution of
\eqref{eq:return-density-selector}; the fourth is the source-resolved return
identity.  If the physical prefix expenditures are bounded, then
\(\int p_{\Phys}\,\dd\gamma_N=\Pi_N/A_{\Sigma,N}\to0\).  Lower
semicontinuity gives zero physical density on every weak-* limit.  No limit
has been taken in the multiplicity coordinates, however: they remain in the
augmented record and reconstruct the original prefix through the first
identity.  Applying the positive-part and five-channel pigeonhole argument
to \(\sum_cQ_N^c=A_{\Sigma,N}\) gives a channel attaining
\(A_{\Sigma,N}/5\) on a cofinal subsequence.  This proves
\eqref{eq:circulation-restored-source-growth} and forbids the inference from
zero normalized density to zero cumulative structure.
\end{proof}

\begin{theorem}[Target-current circulation duality]
\label{thm:target-current-circulation-duality}
Write
\(\mathcal K^{0}_{\mathfrak b}:=I_{\mathfrak b}+M_{\mathfrak b}\) for the
local ordered cone before shift coboundaries are adjoined.  All functions
below are functions on the target quotient; hence no current-null space is
added to this scalar cone.  On every reachable active signature,
\begin{equation}
 \label{eq:target-current-circulation-duality}
 \delta^{\rm cyc}_{\Sigma,\mathfrak b}
 =\sup\left\{\delta\in\mathbb R:
 \begin{array}{l}
 \text{there are }N\text{ and a generated continuous cylinder potential }V
 \text{ with}\\
 p_N+V\circ\pi_- -V\circ\pi_+
       -\delta a_\Sigma\in
       \overline{\mathcal K^{0}_{\mathfrak b}}
 \end{array}\right\}.
\end{equation}
Exactly one of the following values is produced.
\begin{enumerate}[label=\textup{(D\arabic*)}]
\item \(\delta^{\rm cyc}_{\Sigma,\mathfrak b}>0\).  A rational
\(0<\delta<\delta^{\rm cyc}_{\Sigma,\mathfrak b}\) and a finite cylinder
approximation of \(V\) give a finite local ordered identity.  On the first
\(N\) returns it gives
\[
 \sum_{j<N}p_{\Phys}(e_j)
 \ge \delta\sum_{j<N}a_\Sigma(e_j)-2\lVert V\rVert_\infty.
\]
Thus a recurrent sequence with unbounded accumulated unnormalized target
return has infinite public physical charge.
\item \(\delta^{\rm cyc}_{\Sigma,\mathfrak b}=0\).  Compatible minimizing
cylinder moments reconstruct a zero-density target-current circulation.  It
asserts failure of a uniform linear price and nothing about accumulated
quantities.  Its ergodic components and positive support select the next
production-density, return, or graph-boundary cell.  Before any target or
kernel operation, the selected law is augmented by the full multiplicity
record
\[
 (\Gamma_N,A_{\Sigma,N},\Pi_N,(Q_N^c)_c)_N
\]
and the identities
\eqref{eq:circulation-denormalization-identities}.  The next projection in
\cref{def:target-aligned-recurrent-closure-compiler} is applied to this
augmented record, never to the probability law alone.
\item \(\delta^{\rm cyc}_{\Sigma,\mathfrak b}=+\infty\).  A finite ordered
identity either excludes the normalized circulation constraints, or, for
each rational \(M>0\), gives recurrent physical price at least \(M\).
In the latter case choose \(M\) above the remaining public physical ceiling;
the recurrent cell is excluded by physical charge, not declared empty.
\end{enumerate}
The potential \(V\) is generated by the cylinder hierarchy.  It is not a
multiplier supplied with the PDE presentation.
\end{theorem}

\begin{proof}
Let \(\mathcal M_{\mathfrak b}\) be the convex set of positive edge measures
satisfying the first relation in
\eqref{eq:target-current-circulation-constraints} and the active ordered
relations.  On the normalized slice, \(a_\Sigma\ge\eta\) bounds total mass;
the slice is therefore weak-* compact.  First minimize each continuous lower
envelope \(p_N\).  Linear-programming separation on the compact slice gives
its continuous dual.  Compactness of minimizing measures and lower
semicontinuity show that these minima increase to the minimum of
\(p_{\Phys}\): otherwise a convergent subnet of minimizers would violate the
Portmanteau inequality.  The annihilator
of equal source and range marginals is the closed space of coboundaries
\(V\circ\pi_- -V\circ\pi_+\).  The active ordered relations contribute
\(\overline{\mathcal K^{0}_{\mathfrak b}}\).  The edge space was formed after
the target quotient, so target-null currents have already disappeared and
cannot be added to the scalar dual cone.  This proves
\eqref{eq:target-current-circulation-duality} by strong separation.

The rational cylinder functions are uniformly dense in the generated
evaluation algebra.  If the optimum is positive, choose rational \(\delta\)
strictly below it and approximate the dual potential.  Archimedean
order-unit completion absorbs the approximation error and the finite
cylinder-moment hierarchy returns the corresponding finite identity.  Its
sum along a finite path telescopes the coboundary, and its integral against
a circulation cancels it exactly.  Hence
\(\int p_{\Phys}\ge\delta\int a_\Sigma=\delta\), proving \textup{(D1)}.

If the optimum is zero, choose minimizing feasible moments at each cylinder
depth.  Compact projective consistency reconstructs a circulation with zero
physical \emph{density}.  Ergodic decomposition preserves flow conservation,
zero density, and positive target mass on at least one component.  Its
probability law selects only the support of the next recurrent compiler
state.  Apply \cref{thm:circulation-multiplicity-restoration}: the actual
occurrence, physical-expenditure, target-return, and five source-current
prefixes are reattached before joint saturation.  Then
\cref{thm:continuum-composition-no-creation} produces the cumulative joint
state.  In particular, \(\Pi_N/A_{\Sigma,N}\to0\) is never rewritten as
\(\Pi_N=0\), and the divergent cumulative source coordinate
\eqref{eq:circulation-restored-source-growth} remains live.  This proves
\textup{(D2)}.
If the feasible slice is empty, Archimedean separation supplies the finite
\(-1\)-identity in \textup{(D3)}.  If it is nonempty and the lower-semicontinuous
price is infinite on every normalized circulation, the minima of the
continuous lower envelopes increase without bound.  For any prescribed
rational \(M\), some finite envelope and its dual potential give the ordered
price identity \(\int p_{\Phys}\ge M\).  Taking \(M\) above the remaining
finite ledger excludes the recurrent branch.
\end{proof}

\begin{theorem}[Positive-price contradiction and finite certificate]
\label{thm:positive-price-contradiction-finite-certificate}
Let an active same-origin path have unbounded unnormalised target return
\(A_{\Sigma,N}\), finite public physical ceiling, and either a positive
finite-horizon price from
\cref{thm:finite-horizon-target-current-duality} or a positive recurrent
price from \cref{thm:target-current-circulation-duality}.  Then a finite
prefix has required physical price strictly larger than its remaining public
physical ceiling.  Hence the path is excluded by a finite proof certificate.

The certificate contains the target-active branch ideal, the rational
cylinder potential or layer potentials, a rational positive price
\(\delta\), the finite carrier-domination identity, the original disjoint
event indices, and the remaining physical ceiling.
\end{theorem}

\begin{proof}
Choose a rational \(0<\delta\) below the certified price.  The finite dual
row gives
\[
 \Pi_N\ge\delta A_{\Sigma,N}-2\lVert V\rVert_\infty
\]
for the relevant finite prefix, with the analogous telescoped bound for
layer potentials.  Since \(A_{\Sigma,N}\) is unbounded, choose \(N\) for
which the right side exceeds the remaining physical ceiling.  The
no-double-charging theorem identifies \(\Pi_N\) with charges on distinct
original event indices and bounds it by that ceiling, a contradiction.
All terms in the displayed finite dual row and the selected prefix are
finite presentation-generated data, so they form the stated certificate.
\end{proof}

\begin{corollary}[No averaging discharge of a recurrent contact cascade]
\label{cor:no-averaging-discharge}
Let an actual same-origin branch complete infinitely many target returns while
its public physical expenditure remains finite.  No normalized circulation
obtained from that branch can close, nullify, or quotient the branch merely
because its physical-expenditure density is zero.  More precisely, after
passing to a cofinal subsequence, there is a channel \(c_*\) carrying the
unbounded positive cumulative coordinate
\[
 (Q_{N_k}^{c_*})_+\ge A_{\Sigma,N_k}/5\longrightarrow\infty.
\]
At every refinement depth the compiler must therefore do one of the
following with the augmented pair consisting of the support selector and
this cocycle: charge a positive linear price by \textup{(D1)}, expose a least
new finite joint source word, or retain the unbounded \(c_*\)-coordinate in a
joint-saturated structural atom.  Target-nullity is impossible because it
would force \(A_{\Sigma,N}=\sum_cQ_N^c=0\) on every prefix.

Thus the implication used by the recurrent calculus is
\[
 \frac{\Pi_N}{A_{\Sigma,N}}\longrightarrow0
 \quad\Longrightarrow\quad
 \text{zero-density support plus the full cumulative cocycle},
\]
and never zero cumulative target-aligned structure.
\end{corollary}

\begin{proof}
Apply \cref{thm:circulation-multiplicity-restoration} and choose \(c_*\) by
the finite-channel pigeonhole argument.  Rejoin each prefix with the
selector before projection.  The cumulative-joint theorem
\cref{thm:continuum-composition-no-creation} says that a nonzero target
component is either detected by a finite joint word or retained in the
graph-complete fixed point, with inherited five-channel ancestry.  The
positive-price case is \textup{(D1)} of
\cref{thm:target-current-circulation-duality}.  Finally, target-nullity of
all cumulative channels contradicts the exact return identity.  These
alternatives exhaust the augmented record and use no per-return threshold.
\end{proof}

\begin{definition}[Circulation-elimination operator]
\label{def:circulation-elimination-operator}
For a target-spine cell \(X\), first compute the unnormalized horizons
\(\Delta_X(N)\) of
\cref{def:unnormalized-finite-horizon-current}.  Remove \(X\) when a finite
cumulative dual row exceeds the remaining public physical ceiling.  On the
complementary branch retain the infinite augmented ledger supplied by
\textup{(U2)} of \cref{thm:cumulative-carrier-duality}.  Apply recursive
carrier-polar exhaustion to every canonical cumulative owner; every
separator and its adjoint descendants are rejoined and reprojected.  Only on
a primal--polar fixed support compute
\(\delta^{\rm cyc}_{\Sigma,\mathfrak b}\) in the nested cylinder hierarchy
as a recurrent-support reduction.  Remove \(X\) in cases \textup{(D1)} and
\textup{(D3)}.  In case
\textup{(D2)}, restrict to the zero-density circulation support, adjoin all
production densities that vanish on that selector, retain their cumulative
prefix integrals without setting them to zero, and inspect the least unprocessed target-cyclic
covector or total graph in the fixed public enumeration.  Return its finite
zero, signed-dyadic, and graph-boundary polarity partition after deleting
the target-null members.  Rejoin every returned cell with its entire parent
transcript, saturate all mixed coordinates, and reapply the target projection
before it is embedded in the universal spine.  Thus each child retains
\((\Gamma_N,A_{\Sigma,N},\Pi_N,(Q_N^c)_c)_N\), including every accumulated
source current and physical expenditure, and is a closed joint refinement of \(X\).  Iterate
\begin{equation}
 \label{eq:circulation-elimination-iteration}
 X_0:=\Spine_{\Sigma,\mathfrak b},
 \qquad
 X_{n+1}:=\mathcal E_\Sigma(X_n),
 \qquad
 X_\infty:=\bigcap_{n\ge0}X_n,
\end{equation}
where every \(X_n\) is represented in its cumulative joint algebra and the
intersection is followed by the joint graph-complete saturation
\eqref{eq:cumulative-joint-limit-algebra}.  Thus
\(X_\infty\) is a joint fixed-point core rather than a coordinatewise
inverse-limit record.  This operator is generated solely from the target
spine, public physical valuation, and ordered equation graph.  Its priority
order is cumulative horizon, carrier-polar exhaustion, recurrent support,
polarity refinement, joint saturation, and target reprojection.
\end{definition}

\begin{definition}[Target-aligned structural branch-and-bound]
\label{def:target-aligned-structural-branch-bound}
Fix the enumeration of rational target-cylinder coordinates, public graph
relations, and source descendants.  At depth \(n\), partition the first
\(n\) coordinates into rational closed boxes of diameter at most \(2^{-n}\)
in the canonical evaluation metric.  Only coordinates in the generated
target-cyclic spine are partitioned.  For every surviving parent box \(C\),
perform the following finite operations.
\begin{enumerate}[label=\textup{(B\arabic*)}]
\item Form the saturated joint record of the complete parent prefix.  Adjoin
the first \(n\) equation, origin, carrier, carrier-polar, polarity, mixed-product,
shell-cocycle, and total-graph rows.  Form every layered horizon of length at
most \(n\) and use the continuous price envelope \(p_n\).
\item Solve the resulting finite rational moment relaxation and enumerate
its ordered dual identities.  For every retained source owner, check every
finite algebraic cone identity and every finite-cylinder separator whose
polar inequality is verified by the retained projected cone.  When neither
check succeeds, retain the total cone-boundary coordinate rather than
declaring membership or nonmembership.  Append every verified separator and
its generated descendants before solving the normalized circulation
relaxation.
Record the
largest certified cumulative lower bound through horizon \(n\), the certified recurrent dual lower price
\(\underline\delta_n(C)\) and the compact set
\(\mathsf P_n(C)\) of feasible truncated moments.
\item Delete \(C\) when a cumulative horizon lower bound exceeds the
remaining physical ceiling, when a finite \(-1\)-identity makes
\(\mathsf P_n(C)=\varnothing\), or when a positive recurrent dual identity gives
\(\underline\delta_n(C)>0\).  Attach that finite identity to the deleted
box.
\item Otherwise subdivide the least unprocessed target-visible cylinder
coordinate.  A zero value adjoins the corresponding kernel relation; a
nonzero value is divided into signed dyadic cells; and a nonordinary value
is divided by its typed total-graph boundary.  Origin and source ancestry are
copied to every child.  Each child is rejoined with the parent prefix and
reprojected; a target component produced by a combination of previously
target-null coordinates is therefore exposed at this depth.
\end{enumerate}
Let \(\mathscr U_n\) be the union of the surviving boxes, intersected with
their surviving parents, and let
\begin{equation}
 \label{eq:target-structural-computable-core}
 \mathscr U_\infty:=\bigcap_{n\ge0}\mathscr U_n.
\end{equation}
The depth \(n\) is an algebraic refinement index, not a second physical
budget.  Every operation uses rational cylinder arithmetic and the public
ordered graph; no multiplier, compactness certificate, or analytic estimate
is supplied by the user.
\end{definition}

\begin{definition}[Joint-saturated prospective source atom]
\label{def:joint-saturated-structural-atom}
A \emph{joint-saturated structural atom} on an active signature
\(\mathfrak b\) is a tuple
\begin{equation}
 \label{eq:joint-saturated-structural-atom}
 \mathfrak a
 =(X_{\mathfrak a},\mathcal J_{\mathfrak a},\Phi,
   (Q^c_{\mathfrak a})_{c\in\Channels},
   A_{\Phi,\mathfrak a},p_{\Phys},\mathbf R)
\end{equation}
with the following generated properties.
\begin{enumerate}[label=\textup{(A\arabic*)}]
\item \(\mathcal J_{\mathfrak a}\) is the cumulative join of a cofinal
  same-origin compiler execution and is fixed by joint saturation, all five
  source projections, every public mixed operation and total graph, signed
  carrier gluing, target-cyclic adjoint transport, and first-return
  formation.  It is also fixed under the recursive carrier-polar operation
  \(\mathfrak P_\Sigma\) for every canonical source or polar owner.
\item Its unnormalized cocycle is retained and nonzero on every completed
  shell prefix:
  \begin{equation}
   \label{eq:atom-unit-cocycle}
   A_{\Phi,N}=N
   =\sum_{c\in\Channels}Q^c_{\Phi,N}(1)
   \qquad(N\ge1).
  \end{equation}
  Hence its persistent source support
  \(A(\mathfrak a):=\{c:Q^c_{\mathfrak a}\ne0\}\) is nonempty.
\item Every target-visible finite joint prefix, mixed polarization,
  production, feeding, storage, carrier separator, polar descendant,
  graph-boundary, and realization value has its
  zero, signed-dyadic, or typed successor inside
  \(\mathcal J_{\mathfrak a}\).  The post-saturation residual annihilates
  every prospective target functional.
\item The physical coordinate is the restriction of the single public
  measure.  Every owner in the public physical carrier cone has been removed
  by its cumulative carrier identity; every owner outside that cone has its
  complete separating covector and generated primal--polar descendants in
  the joint record.  The retained zero-density support law is only a support
  coordinate and remains joined to \eqref{eq:atom-unit-cocycle}.
\end{enumerate}
The atom is a \emph{contact-derived joint source record} when
\(\operatorname{pol}_{\rm org}(\mathbf R)=\mathsf{Hyp}\).  It remains in the
same presentation-generated pre-contact domain and is evaluated by its
retained shell covector, source word, and physical carrier.
The atom is \emph{independently realized} when
\(\operatorname{pol}_{\rm org}(\mathbf R)=\mathsf{Ind}\) and \(\mathbf R\)
contains one ordinary same-origin public solution path whose continuation row
and cofinal shell rows realize the target contact carried by \(\Phi\).  A formal fixed-point
law, compatible moment family, or ordinary path without
\textup{(A1)}--\textup{(A4)} is not a prospective source atom.  A
contact-derived atom remains in the reachable invariant evaluator.  An
ordinary same-origin realization with the declared noncontinuation row is a
singular contact record.
\end{definition}

\begin{theorem}[Stagewise computation of the target mechanism core]
\label{thm:stagewise-target-core-computation}
The branch-and-bound construction has the following properties.
\begin{enumerate}[label=\textup{(C\arabic*)}]
\item Every depth is a finite exact calculation \emph{of the stated finite
finite prefix}.  Its output consists of a finite family of
presentation-generated target-cylinder records \(\mathscr U_n\), checked
ordered identities for every deleted incidence edge, and exact finite-prefix
invariant values for every retained record.  The families are nested.
\item The construction loses no target-active mechanism and retains no
untyped infinite branch.  Because every child contains its cumulative joint
prefix, one has
\begin{equation}
 \label{eq:branch-bound-equals-circulation-core}
 \mathscr U_\infty=X_\infty.
\end{equation}
Thus \(\mathscr U_n=\varnothing\) at some finite depth implies that the
target mechanism core is empty.  The complete core is the source-preserving
closure of these finite records, and its zero detector profile is evaluated by
\cref{thm:reachable-source-saturation-finite-occurrence}.
\item Let \(F\) be a rational cylinder observable, or a generated bounded
target-aligned structural observable whose constructor graph supplies a
verified cylinder-approximation modulus.  From the finite boxes one obtains
certified intervals
\[
 I_n(F):=
 \left[inf_{\mathscr U_n}F_n-2^{-n},
       \sup_{\mathscr U_n}F_n+2^{-n}\right],
 \]
where \(F_n\) is its generated rational cylinder approximation.  After
replacing each interval by its intersection with the preceding ones,
\begin{equation}
 \label{eq:target-observable-range-computation}
 I_n(F)\downarrow
 \left[\min_{X_\infty}F,\max_{X_\infty}F\right].
\end{equation}
Consequently every observable in this effective cylinder class is computed
by certified two-sided enclosures.  Bare continuity supplies density but not
an effective approximation rate and is not promoted to a computed result.
\item If the core is nonempty, the complete stream
\((\mathscr U_n)_n\), together with the compatible moment tables and origin
  rows, determines the compact graph-complete joint target-mechanism set.  Join
  and saturate every infinite nested box word before taking its target
  projection.  A least unprocessed limiting increment occurs on a finite
  joint prefix by \cref{thm:continuum-composition-no-creation}; that prefix is
  reclassified and refined.  If every increment is processed, the word is
  tested for joint fixedness.  Failure of fixedness exposes the next
  generated operation; a fixed word supplies the structural coordinates
  \textup{(A1)}--\textup{(A4)} of
\cref{def:joint-saturated-structural-atom}.  Its total realization value is
then ordinary, giving a prospective source atom with its unchanged origin
polarity, or nonordinary, giving the next typed realization successor.  A compatible moment family, invariant law, or
ordinary path alone is never an output.
\end{enumerate}
The algorithm is target local: refining any source-transverse or
target-null coordinate leaves every \(\mathscr U_n\) and every interval
\(I_n(F)\) unchanged.
\end{theorem}

\begin{proof}
At fixed depth there are finitely many boxes, rows, monomials, and moment
coordinates.  The retained polynomial moment relaxation and its ordered dual
are finite rational semialgebraic problems; a displayed rational dual identity
is checked by direct arithmetic.  Exponentials, resolvents, unbounded
operators, distributional products, and limiting relations enter through
their bounded total-graph coordinates and finite graph rows, not by being
misidentified as polynomial maps.  Their unprocessed graph boundaries remain
in the outer cover.  This proves \textup{(C1)}.

Every genuine circulation satisfies every finite relaxation, so it is never
deleted.  Conversely, choose a point or compatible moment table from each
box along an infinite nested word.  Compactness of the bounded total-graph
coordinates gives a diagonal limit.  Because the enumeration is exhaustive,
each equation, carrier, polarity, origin, and graph relation is imposed from
some finite depth onward.  The least-unprocessed-coordinate rule therefore
gives precisely the circulation-elimination limit, proving
\eqref{eq:branch-bound-equals-circulation-core}.  If the intersection is
empty, compactness of the decreasing covers gives a finite empty depth; the
reverse implication is immediate.

The constructor-supplied modulus gives
\(\lVert F-F_n\rVert_{\mathscr U_n}\le2^{-n}\).  Extrema of a continuous
function over decreasing nonempty compact sets converge to the extrema over
their intersection.  This proves
\eqref{eq:target-observable-range-computation}.  Finally, the compatible
moments reconstruct a measure on the total-graph inverse limit, and the
origin and cofinal-shell rows survive.  Apply cumulative joining before
interpreting that measure.  By
\cref{thm:continuum-composition-no-creation}, every nonzero limiting target
increment already occurs on a finite joint prefix.  Every unprocessed such
increment is therefore another branch-and-bound coordinate.  When no such
increment remains, apply recursive carrier--polar exhaustion to every
cumulative source and polar owner.  Target-null and physically charged
owners return to their preceding compiler rows; a nonzero separator or one
of its adjoint descendants is the next branch-and-bound coordinate.  Only
simultaneous primal--polar fixedness gives \textup{(A1)}--\textup{(A4)} of
\cref{def:joint-saturated-structural-atom}.  The totalized public-orbit row
then gives either the ordinary prospective source atom, with its origin
polarity unchanged, or its next typed
successor.  All partition coordinates lie in the target
quotient by construction, proving target locality.
\end{proof}

\begin{theorem}[Completeness of circulation elimination]
\label{thm:circulation-elimination-completeness}
The sets in \eqref{eq:circulation-elimination-iteration} are closed,
target-active, decreasing, and represented in increasing cumulative joint
algebras.  After the joint limit saturation in
\eqref{eq:cumulative-joint-limit-algebra},
\begin{equation}
 \label{eq:circulation-core-equals-source-core}
 X_\infty
 =\bigsqcup_{\mathfrak a\in
   \mathsf{Atom}^{\rm str}_\Sigma(\mathfrak b)}X_{\mathfrak a},
\end{equation}
where \(\mathsf{Atom}^{\rm str}_\Sigma(\mathfrak b)\) is the family of
joint-saturated prospective source atoms of
\cref{def:joint-saturated-structural-atom}.  A formal atom with a nonordinary
public realization coordinate is passed to its typed realization successor.
On a contact-generated same-origin execution, every atom has hypothetical
origin and supplies only a saturated source-profile lower bound.  Consequently
\begin{equation}
 \label{eq:contact-source-profile-witness}
 \text{contact on the branch}
 \quad\Longrightarrow\quad
 \mathsf{Atom}^{\mathsf{Hyp}}_\Sigma(\mathfrak b)\ne\varnothing.
\end{equation}
Independently generated singular contact is equivalent to a member of
\(\mathsf{Atom}^{\mathsf{Ind},\rm real}_\Sigma(\mathfrak b)\).
If the intersection is empty, compactness gives a finite depth at which it is
empty; concatenating the circulation-dual identities gives a finite
target-local regularity derivation.  A nonempty intersection is processed by
joint reprojection; neither the intersection nor an ordinary path is by
itself a terminal.
\end{theorem}

\begin{proof}
Closedness follows from weak-* compactness of the circulation slices and
total-graph closure.  The operator removes exactly the cells having no
normalized recurrent circulation or having positive recurrent physical
price.  On a zero-price cell it resolves the least target-cyclic descendant.
Consequently membership in every \(X_n\) is equivalent to compatibility with
the first \(n\) circulation, production, carrier--polar, polarity, graph, origin, and
cofinal-shell rows, together with every mixed coordinate generated by their
cumulative join.  Apply
\cref{thm:continuum-composition-no-creation} to an infinite compatible word.
If its joint target projection is zero, the cell is deleted in
\(\Null_\Sigma\).  If it is nonzero, a finite joint prefix already has a
nonzero target projection.  That prefix is either the next source-resolved
compiler state or has already been processed.  In the latter case, closure
under all generated joint operations, the unnormalized identity
\eqref{eq:atom-unit-cocycle}, complete source ancestry, and the retained
physical valuation are passed through
\cref{thm:recursive-carrier-polar-exhaustion}.  A target-null owner deletes
the cell, a cone identity charges it, and a separator is the next finite
coordinate.  Only a primal--polar fixed record satisfies precisely
\textup{(A1)}--\textup{(A4)}.  This proves
\eqref{eq:circulation-core-equals-source-core}.  No division by accumulated
return is used.

For decreasing closed subsets of the compact spine, empty intersection
implies that one finite intersection is empty.  The fixed enumeration records
the finitely many dual identities used before that depth.  On an execution
initialized by contact, all cumulative records retain polarity
\(\mathsf{Hyp}\), so the fixed atom gives the source-profile implication
\eqref{eq:contact-source-profile-witness} and no existence conclusion.
In the independent fiber, a realized atom contains the ordinary same-origin
path and the cofinal shell rows, hence supplies contact; conversely an
independently generated singular contact supplies those rows.  A nonordinary public realization value is a
typed successor by totalization and cannot end the compiler.
\end{proof}

\begin{definition}[Explicit singular mechanism]
\label{def:explicit-singular-mechanism}
An \emph{explicit singular mechanism} is a realized joint-saturated
structural atom \(\mathfrak a\) together with its same-origin path
\(u_\ast\) in the closed reduced public state space
\(X_{\mathfrak a}\).  Equivalently, it satisfies
\textup{(A1)}--\textup{(A4)} of
\cref{def:joint-saturated-structural-atom} and the following realization
conditions:
\begin{enumerate}[label=\textup{(E\arabic*)}]
\item \(\operatorname{pol}_{\rm org}(\mathbf R)=\mathsf{Ind}\); the public
orbit is generated without assuming target contact or noncontinuation, and
origin-polarity preservation verifies that no contact-derived record occurs
in its construction;
\item its orbit-origin and weak-equation realization rows are ordinary and
every public-equation fidelity defect vanishes; in particular the projection
of \(u_\ast\) is an actual same-origin path solving the original public weak
equation, not an equation with an independent defect forcing;
\item every auxiliary constraint, carrier, chart, covariance, boundary, and
support-transfer coordinate has either its ordinary public value or its
first typed total-graph value together with the complete successor
descendants generated from that value; no auxiliary failure is discarded or
promoted to forcing in the public equation;
\item the induced five-channel equation on \(X_\ast\) is closed under the
resaturation successor and the recurrent passivity row has been evaluated;
every nonzero feeding or storage current carries its signed-dyadic return
law, detecting target covector, source ancestry, and same-origin successor,
whereas every vanishing production is recorded in the joint kernel;
\item the path has the inherited physical ledger and origin coordinate; and
\item the shell functions satisfy
\(\sup_t\Phi_r(u_\ast(t))=1\) for every member of a cofinal target-shell
family.
\end{enumerate}
Thus an ordered positive functional, invariant law, formal profile, or
equality kernel is not an explicit singular mechanism.  An ordinary singular
path without its cumulative joint atom is likewise not an output.  A
graph-boundary coordinate becomes part of a mechanism only when its typed
descendant closure, nonzero source projection, and unnormalized shell cocycle
are retained and the ordinary same-origin weak-equation row realizes contact.
\end{definition}

\begin{theorem}[Terminal target-aligned current compiler]
\label{thm:three-output-target-aligned-current-compiler}
Run \cref{def:target-aligned-recurrent-closure-compiler} on a contact-selected
active signature.  Every maximal compiler execution has exactly one of the
following outputs:
\begin{enumerate}[label=\textup{(C\arabic*)}]
\item a cofinal collection of disjoint return events is charged to the public
physical measure with divergent total mass;
\item the retained same-origin path lies in \(\Null_\Sigma\), so it cannot
cross the selected target shells;
\item cumulative joint saturation and recursive carrier--polar exhaustion
reach a primal--polar fixed prospective source atom.  Its retained shell floor
is a lower-bound witness for one saturated five-channel source-profile cell.
\end{enumerate}
Compactness modes, fidelity defects, equality kernels, feeding classes,
storage transfers, and nonordinary graph values are internal queue
coordinates.  None is an additional output.  Consequently the full
continuum regularity track on a finite physical stratum is the target-local
decision
\begin{equation}
 \label{eq:three-output-target-aligned-decision}
 \boxed{
 \text{contact}
 \ \Longrightarrow
 \mathsf{Charge}_{\infty}
 \ \dot\cup
 \Null_\Sigma
 \ \dot\cup
 \mathsf{SrcWitness}^{\mathsf{Hyp}}_\Sigma .}
\end{equation}
The first two cells contradict contact under a finite physical ledger.  The
third is the continuum prospective-accountability output: contact has exposed
a nonzero source cell that must be bounded independently by the saturated
source profile.  It is not a singular-existence conclusion.  An infinite
finite-prefix execution is joined and reprojected.  It reaches
\textup{(C3)} only after the nonzero five-channel source, all mixed
descendants, and the unnormalized shell cocycle have been assembled in one
atom with hypothetical origin.
\end{theorem}

\begin{proof}
At a finite queue state, cumulative carrier duality, recursive
carrier--polar exhaustion, target-current circulation duality, target-local
stratification, and the ordered proof--model alternative give
\textup{(Q1)} or one of \textup{(Q2)}--\textup{(Q4)}.  Positive cycle price
or infeasibility supplies the finite dual identity in \textup{(Q1)}.  A
finite-horizon lower bound above the remaining physical ceiling supplies the
same output without a per-return price.  On its complementary branch,
\eqref{eq:bounded-price-unbounded-target-current} is retained before zero
cycle price supplies the primal circulation processed by the remaining
transitions.  The local charging rule proves that only \textup{(Q1)} spends
physical measure.  The joint-kernel descent proves that
\textup{(Q2)} either gives \textup{(C2)} or supplies another target-aligned
coordinate.  Recurrent stratification and Poincare recurrence prove that
\textup{(Q3)} retains every positive-density current, including an exact
storage transfer.  Its normalized law is rejoined with the cumulative
unnormalized current before target projection.  Fidelity-defect support transfer proves that
\textup{(Q4)} retains every target-visible failed operation.  These are the
complete polarity and total-graph values of the least coordinate, so no
finite execution has another outgoing edge.

For an infinite execution, form \(\mathcal J_\infty\).  The same-origin
witnesses and \(A_{\Phi,N}=N\) survive by joint construction.  By
\cref{thm:continuum-composition-no-creation}, every limiting target
increment occurs on a finite joint prefix.  If a least unprocessed increment
exists, its prefix is another \textup{(Q2)}--\textup{(Q4)} transition, so the
execution was not maximal.  If the complete joint target projection is zero,
the limit gives \textup{(C2)} and contradicts its unit shell cocycle.  In the
remaining maximal case every increment is processed.  Joint saturation now
either exposes another mixed increment, again contradicting maximality, or
is fixed under all generated joint operations.  Apply
\cref{thm:recursive-carrier-polar-exhaustion} to every cumulative owner.
Target-nullity gives \textup{(C2)}, a cone identity contributes to
\textup{(C1)}, and a separator contradicts maximality by creating another
queue transition.  The remaining record is primal--polar fixed and satisfies
\textup{(A1)}--\textup{(A4)}.  Its contact-selected origin has polarity
\(\mathsf{Hyp}\) by \cref{thm:no-recycled-contact-realization}; its retained
source floor gives \textup{(C3)} and cannot supply
\textup{(E1)} of \cref{def:explicit-singular-mechanism}.  A divergent cumulative carrier row is \textup{(C1)}.  The
fixed transition priority and disjoint target quotient make the outputs
mutually exclusive.
\end{proof}

\begin{corollary}[Prospective terminal typing]
\label{cor:binary-terminal-semantics}
On a finite physical stratum, a contact-seeded execution has exactly two
mathematical output types:
\begin{equation}
 \label{eq:binary-terminal-semantics}
 \boxed{
 \mathsf{Output}^{\mathsf{Hyp}}_\Sigma
 \in
 \{\mathsf{TargetUnreachable},\
   \mathsf{SrcWitness}^{\mathsf{Hyp}}_\Sigma\}.}
\end{equation}
The first value carries the local carrier or ordered identities which exclude
contact; the second carries a nonzero lower-bound witness for the saturated
prospective source profile.  It carries no singular-existence conclusion.
A successor cell, compatible cover, positive
functional, invariant law, equality kernel, normalized circulation, or
joint-saturation prefix is not an output.
\end{corollary}

\begin{proof}
The public physical ceiling turns \textup{(C1)} of
\cref{thm:three-output-target-aligned-current-compiler} into a contradiction,
and the unit shell identity does the same for \textup{(C2)}.  Both therefore
produce the target-unreachable value; \textup{(C3)} is the hypothetical
source-profile witness.
All other objects are queue states by the same theorem.
\end{proof}

\begin{corollary}[Target-spine prospective regularity criterion]
\label{cor:target-spine-regularity-criterion}
On an ordinary branch with finite public physical ledger,
\begin{equation}
 \label{eq:target-spine-reachability-equivalence}
 \boxed{
 \Reach_\Sigma^{\Phys}>0
 \quad\Longrightarrow\quad
 \mathsf T_\Sigma^{\Phys,\Sat}(\Budget)\ge\frac15 .}
\end{equation}
Consequently
\begin{equation}
 \label{eq:target-spine-regularity-equivalence}
 \boxed{
 \mathsf T_\Sigma^{\Phys,\Sat}(\Budget)<\frac15
 \quad\Longrightarrow\quad
 \Reach_\Sigma^{\Phys}=0 .}
\end{equation}
The profile is computed only on the projected target spine and decomposes
into the exhaustive source/mode cells of
\cref{thm:realization-mode-exhaustion}.  An independently realized atom still
proves singular existence, but a hypothetical atom contributes only to this
profile.
\end{corollary}

\begin{proof}
If contact occurs, initialize the compiler with its shell current.  The
finite physical ledger excludes output \textup{(C1)} of
\cref{thm:three-output-target-aligned-current-compiler}, and shell crossing
excludes \textup{(C2)}.  Output \textup{(C3)} therefore gives a
joint-saturated prospective source atom with retained floor at least
\(1/5\).  By
\cref{prop:spine-restriction-recurrent-compiler} its cumulative record lies
in the target spine, so the profile supremum is at least \(1/5\).  This proves
\eqref{eq:target-spine-reachability-equivalence}; its contrapositive is
\eqref{eq:target-spine-regularity-equivalence}.
\end{proof}

\begin{theorem}[Quantifier-faithful continuum accountability]
\label{thm:quantifier-faithful-continuum-accountability}
Five-source exhaustion and physical-budget accounting imply, for every public
continuum presentation and every finite physical stratum, exactly the
presentation-relative implication
\begin{equation}
 \label{eq:quantifier-faithful-continuum-accountability}
 \Reach_\Sigma^{\Phys}>0
 \quad\Longrightarrow\quad
 \mathsf T_\Sigma^{\Phys,\Sat}(\Budget)\ge\frac15 .
\end{equation}
They do not imply
\(\mathsf T_\Sigma^{\Phys,\Sat}(\Budget)<1/5\) uniformly over public
presentations.  Indeed, any presentation containing an independently realized
ordinary orbit which crosses the cofinal target shells has
\(\Reach_\Sigma^{\Phys}>0\), and therefore has profile at least (1/5) by
\eqref{eq:quantifier-faithful-continuum-accountability}.

Consequently a specialization has precisely two mathematically sound ways to
finish.  A generated upper profile below (1/5) proves target avoidance by
\cref{cor:target-spine-regularity-criterion}; an independently realized
ordinary contacting atom proves singular reachability.  Source
classification alone supplies neither terminal value.  This is the continuum
counterpart of a saturated-accountability bound: it is universal in the
admissible evolution and presentation-relative in the profile being
evaluated.
\end{theorem}

\begin{proof}
The displayed implication is
\eqref{eq:target-spine-reachability-equivalence}.  Apply it to an independently
realized ordinary contacting orbit to obtain the asserted obstruction to a
uniform small-profile theorem.  The two specialization conclusions are,
respectively, the contrapositive
\eqref{eq:target-spine-regularity-equivalence} and the ordinary realization
clause of \cref{def:explicit-singular-mechanism}.  The contact conclusion is
read from the ordinary same-origin realization row of the generated record.
\end{proof}

\begin{theorem}[Zero-price terminalization-to-realization]
\label{thm:no-terminal-zero-price-model}
Apply resaturation to any target-active zero-price successor.  Exactly one of
the following generated events occurs.
\begin{enumerate}[label=\textup{(Z\arabic*)}]
\item a finite dual-carrier row charges infinitely many disjoint target
crossings to positive physical production or exterior input, contradicting
the finite physical ledger;
\item a finite relation makes the branch target-null, regular, continuable,
  or incompatible with the public solution class; this includes every
  finite-coefficient dyadic ordered-coercivity row on the zero-production
  stratum;
\item a first nonordinary total-graph value, a nonzero target-local feeding
class, its positive support record, or a proper kernel restriction creates
the next target-active successor, to which resaturation is reapplied;
\item the cumulative join is fixed by every primal and carrier--polar
operation, retains the unnormalized shell cocycle and a nonzero five-channel
source, and is passed as a \emph{realization request} to the totalized
finite-prefix realization compiler.  That compiler returns a target-null or
physically charged owner, a least finite realization defect and hence another
successor, or an ordinary same-origin contacting realization.  The latter is
accepted as an explicit singular mechanism only in the independent-origin
fiber; in the contact-derived fiber it remains a hypothetical source record.
\end{enumerate}
The alternatives are read in the fixed enumeration of graph coordinates and
dual words: the first finite closing or successor event is selected, and
\textup{(Z4)} occurs only when no finite event occurs.
In particular, a positive stationary functional, target-visible equality
kernel, or hypothetical joint atom can occur only in \textup{(Z3)} or as an
intermediate realization request in \textup{(Z4)}; none is terminal by
itself.  In the contact-derived fiber, \textup{(Z4)} is not a
singular-existence statement.
\end{theorem}

\begin{proof}
At each stage first apply
\cref{thm:dyadic-ordered-coercivity-terminal-rule} to the generated target
activities, physical productions, and coefficient coordinates.  Its
finite-coefficient zero-production value is \textup{(Z2)}, while an unbounded
coefficient or nonordinary relation is the proper successor
\textup{(Z3)}.  Then run the target-cyclic joint-kernel descent of
\cref{thm:target-cyclic-joint-kernel-descent}.  Its case \textup{(K1)} is
\textup{(Z1)}, and \textup{(K2)} is \textup{(Z2)}.  A least nonzero or
nonordinary coordinate in \textup{(K3)} is exactly the proper successor
\textup{(Z3)} specified in
\cref{def:zero-price-resaturation-successor}.  Because its origin and
detecting covector are retained, the theorem applies again without adding
specialization data.  Monotonicity at finite and limit stages is
\cref{prop:monotone-zero-price-resaturation}.

If no finite closing or charged event occurs, form the cumulative joint
record and apply
\cref{thm:infinite-chain-no-emergence-trichotomy,thm:target-relative-terminal-stratification,thm:target-local-passivity-fast-track}.
Every nonzero joint feeding, fidelity, storage, support-transfer, or
nonordinary coordinate occurs on a finite prefix and returns to
\textup{(Z3)} or to its induced recurrent-current successor.  If all finite
joint prefixes are target-null, their limit is target-null.  Otherwise
maximality forces the cumulative record to be fixed under every generated
primal operation.  Recursive carrier--polar exhaustion then charges or
annihilates an owner, creates the next successor, or proves simultaneous
primal--polar fixedness.  Only the last cell continues.  There the
unnormalized shell identity, stationary carrier reduction, and source
exhaustion give \textup{(A1)}--\textup{(A4)}.  Totalize the realization row
and evaluate its least finite condition on the retained same-origin
finite-prefix records.  A null or charged owner returns to
\textup{(Z1)}--\textup{(Z2)}, and a failed condition is the successor in
\textup{(Z3)}.  If every condition is ordinary, the row is an ordinary
contacting realization request; origin-polarity preservation permits an
explicit singular mechanism only when its origin is \(\mathsf{Ind}\).
Origin polarity remains \(\mathsf{Hyp}\) on a contact-seeded execution.  The
priority convention in the joint-kernel theorem gives exclusivity.
\end{proof}

\begin{corollary}[Post-resaturation residual and terminal exhaustiveness]
\label{cor:post-resaturation-terminal-exhaustiveness}
After zero-price resaturation, the residual is target-null only after every
target-visible realization condition has been totalized.  The outputs are a
proved regularity contradiction \textup{(Z1)}--\textup{(Z2)}, a further
successor \textup{(Z3)}, or the realization request \textup{(Z4)}.  A
contact-derived request is a prospective source record; only an
independent-origin ordinary realization can become an explicit singular
mechanism.  Classification names, compactness defects, invariant laws, and
equality kernels are coordinates of this computation, not terminals.
\end{corollary}

\subsection{Diagonal first-failure closure}
\label{subsec:ranked-first-failure-closure}

The preceding protocol becomes a regularity method only after every routed
row is followed to a mathematical conclusion.  A row name records the first
failed structural test; it is not a terminal output.

\begin{definition}[Equation-generated closure graph]
\label{def:equation-generated-closure-graph}
Fix a public PDE presentation and one persistent target-aligned source support
\(A\).  Its \emph{closure graph} \(\mathfrak G_A\) has the following vertices:
the ordinary retained cells and every first-failure cell generated in
steps \textup{(I1)}--\textup{(I8)}.  There is a directed edge
\(v\to w\) precisely when the weak equation, a signed carrier identity, a
source-preserving limit operation, or a target-shell successor sends the
witness defining \(v\) to the witness defining \(w\).  The terminal vertices
of the regularity branch are exactly
\begin{enumerate}[label=\textup{(T\arabic*)}]
\item a target-null cell;
\item a regular or continuable cell;
\item a class-incompatible cell;
\item a cell whose disjoint descendants have divergent total physical price;
\end{enumerate}
Every zero-price equality vertex instead enters the resaturation successor of
\cref{def:zero-price-resaturation-successor}.  The only complementary
conclusive output is
\begin{equation*}
 \tag{S}\label{eq:explicit-singular-terminal}
 \text{a realized joint atom satisfying \textup{(A1)}--\textup{(A4)}
 and \textup{(E1)}--\textup{(E6)}.}
\end{equation*}
The singular terminal \textup{(S)} is available only in the independent
origin fiber.  In the hypothetical origin fiber the same structural rows,
without \textup{(E1)}, define the prospective source-profile output
\begin{equation*}
 \tag{P}\label{eq:hypothetical-source-profile-terminal}
 \text{a joint atom satisfying \textup{(A1)}--\textup{(A4)} with
 nonzero retained shell floor and origin \(\mathsf{Hyp}\).}
\end{equation*}
The graph and every edge relation are generated from the public equation.
They are not additional specialization data.

At depth \(k\), the canonical closure procedure evaluates the first \(k\)
generated graph coordinates and the first \(k\) dual-carrier words.  A finite
first failure creates its displayed successor.  A compatible cofinal sequence
retains all earlier coordinates, its origin tag, and its unnormalized shell
cocycle in the cumulative joint record.  It is reprojected after joint
saturation.  Thus no named row, compatible sequence, or ordinary path is
terminal merely by being classified or realized, and no closure rank is
supplied with the PDE presentation.
\end{definition}

\begin{theorem}[Diagonal structural closure]
\label{thm:ranked-structural-closure}
For every target-aligned contact branch, the canonical finite-depth closure
procedure, completed by cumulative joint saturation, has exactly one of the
following terminal outcomes:
\begin{enumerate}[label=\textup{(G\arabic*)}]
\item a finite stage reaches a terminal of type \textup{(T1)}--\textup{(T4)};
\item the primal--polar joint fixed point reconstructs the prospective
source-profile witness \eqref{eq:hypothetical-source-profile-terminal}.
\end{enumerate}
Compatible stages at every finite depth are an internal state: their
cumulative joint limit is target-null, has a least unprocessed finite joint
primal or polar increment, or is the primal--polar fixed atom in
\textup{(G2)} after every increment has been source-resolved and every
cumulative owner has passed carrier--polar exhaustion.  Hence
\textup{(T1)}--\textup{(T4)} contradict contact, while \textup{(G2)} is a
source-resolved lower-bound witness for the saturated prospective profile.
Absence of singular contact follows when the computed upper profile is below
that retained floor on every algebraically reachable active branch.  Target-local
kernel identities are generated sufficient tests, but a kernel name or
positive functional never ends the procedure.  The criterion requires no
externally specified closure rank.
\end{theorem}

\begin{proof}
Contact produces a persistent nonempty support \(A\) by
\cref{thm:unconditional-continuum-contact-alternative}.  On terminal-entry
contact, persistence is across the cofinal terminal shell records at the same
endpoint; they are one \(\Bdry\) event and are never counted repeatedly as
physical expenditure.  Run the finite cylinder and dual-word enumeration on
the active branch.  A finite terminal gives \textup{(G1)}.  Otherwise every
finite depth has a compatible retained value.  Form the cumulative joint
record rather than declaring its diagonal limit an output.  By
\cref{thm:continuum-composition-no-creation}, every target component is the
sum of its finite joint increments.  Its least unprocessed increment returns
to the closure graph; if the complete joint projection vanishes, the limit is
target-null.  After all increments are processed, the remaining fixed point
retains the original contact events and unnormalized
shell cocycle and satisfies \textup{(A1)}--\textup{(A4)}.  Totalization of
the public realization row either creates another typed successor or retains
the source floor, which is \textup{(G2)}.  Its origin remains
\(\mathsf{Hyp}\) by \cref{thm:no-recycled-contact-realization}.  The
contradictions associated with
\textup{(T1)}--\textup{(T4)} are respectively the target floor, terminal
coverage, the public branch class, and finite physical expenditure.  Finally,
on an approach branch,
\cref{thm:continuous-selector-accountability-closure} identifies the only
recurrent nonterminal support with a target-local equality kernel.  On a
terminal-entry branch, the totalized terminal-current row instead closes the
cell or retains the equation-faithful prospective \(\Bdry\)-source witness.
It is not an existence mechanism on the contact-derived origin fiber.  The
approach equality-kernel resaturation is exhaustive by
\cref{thm:no-terminal-zero-price-model}.  It either closes the branch or
reconstructs \eqref{eq:hypothetical-source-profile-terminal}.  An
independently generated construction may instead satisfy
\eqref{eq:explicit-singular-terminal}; a contact-generated construction
cannot.
\end{proof}

\section{Target-aligned geometric normal form}
\label{sec:continuum-geom-normal-form}

Fix a presentation-generated source support \(A\), a shell covector
\(\lambda\in\mathscr C^\infty_{\mathfrak b}\), and its minimal target-active
realization \(\mathsf M_{\Sigma,\mathfrak b}\).  All spaces, pairings, and
extrema in this and the next four sections are restricted to the reachable
record cell \(\mathscr S_{A,\Sigma}^{\rm reach,sat}(\Budget)\) and its finite
same-origin prefixes.
Write
\[
 \Lambda_{\Sigma,\mathfrak b}
 :=\{\lambda\in\mathscr C^\infty_{\mathfrak b}:
       \|\lambda\|_{\rm ou,*}\le1\}
\]
for the normalized retained shell-covector ball in the dual order-unit norm.
Every supremum below is taken over this ball; this normalization prevents a
nonzero pairing from being made arbitrarily large by rescaling its probe.

\begin{definition}[Geometric range--kernel decomposition]
\label{def:target-geometric-range-kernel-form}
Let \(K\ge0\) be a generated mobility or reversible form and write the
ordinary geometric current as \(J^{\Geom}=-K\zeta\), where \(\zeta\) is the
generated energy, action, or weak-equation covector.  In the form completion
\(\mathcal H_K=\overline{\operatorname{ran}K^{1/2}}\), set
\[
 \mathcal O_{K,\Sigma}
 :=\overline{\operatorname{span}}
 \{K^{1/2}\lambda:
   \lambda\in\mathscr C^\infty_{\mathfrak b}\},
 \qquad \Pi_{K,\Sigma}:\mathcal H_K\to\mathcal O_{K,\Sigma}.
\]
The ordinary target-aligned geometric projection is
\begin{equation}
 \label{eq:target-geometric-normal-form}
 K^{1/2}\zeta
 =\Pi_{K,\Sigma}K^{1/2}\zeta
  +(I-\Pi_{K,\Sigma})K^{1/2}\zeta,
 \qquad
 \operatorname{Type}_{K}(\zeta)
 \in\{\mathrm{ran},\ker/\mathrm{coker},\mathrm{gc}\}.
\end{equation}
Here \(\operatorname{Type}_{K}(\zeta)\) is a branch flag obtained by the
ordered total-graph test: first test ordinary range membership, then the
generated kernel/cokernel quotient, and finally retain the first nonordinary
graph value.  We write \(\zeta_{\rm ran}\), \(\zeta_{\ker}\), and
\(\zeta_{\rm gc}\) for the corresponding coordinates, assigning zero to a
coordinate absent on the selected graph cell.  The flag is unique; it does
not assert that the ordinary target projection and a kernel coordinate
cannot both be detected in the full invariant tuple.
The second summand in the first equality is target-orthogonal.  A graph
failure is retained with \(\Geom\) ancestry and any ancestry of the operation
whose descent failed.
\end{definition}

\begin{definition}[Geometric formation invariants]
\label{def:target-geometric-formation-invariants}
On a normalized target cell \(C\), define
\begin{align}
 G_{\Sigma}^{\Geom}(C)
 &:=\sup_{\substack{x\in C\\
              \lambda\in\Lambda_{\Sigma,\mathfrak b}}}
   \bigl|\langle K^{1/2}\lambda,
          \Pi_{K,\Sigma}K^{1/2}\zeta\rangle\bigr|,
 \label{eq:geom-alignment-invariant}\\
 \operatorname{Cap}_{\Sigma}^{\Geom}(C)
 &:=\inf\{\langle v,Kv\rangle:
       \lambda\in\Lambda_{\Sigma,\mathfrak b},\
       \langle\lambda,Kv\rangle=1\},
 \label{eq:geom-capacity-invariant}\\
 \operatorname{Ker}_{\Sigma}^{\Geom}(C)
 &:=\sup_{\substack{x\in C\\
              \lambda\in\Lambda_{\Sigma,\mathfrak b}}}
   |\langle\lambda,\zeta_{\ker}\rangle|,
 \qquad
 \operatorname{Gc}_{\Sigma}^{\Geom}(C)
 :=\sup_{\substack{x\in C\\
              \lambda\in\Lambda_{\Sigma,\mathfrak b}}}
      |\langle\lambda,\zeta_{\rm gc}\rangle|.
 \label{eq:geom-kernel-graph-invariants}
\end{align}
For a finite stopped record \(R=[s,t]\), with public shell covector
\(\lambda_{R,r}=D\Phi_R(u_r)\), define its geometric production and
shell-probe form cost on the ordinary common form-domain row by
\begin{equation}
 \label{eq:geom-production-probe-cost}
 d_g(R):=\int_s^t\|K_r^{1/2}\zeta_r\|^2\,\dd r,
 \qquad
 h_g(R):=\int_s^t\|K_r^{1/2}\lambda_{R,r}\|^2\,\dd r.
\end{equation}
The generated carrier-admission value is
\begin{equation}
 \label{eq:geom-carrier-admission-value}
 \kappa_g(R):=\inf\{\kappa\ge0:d_g(R)\le\kappa p_{\Phys}(R)\}.
\end{equation}
It is \(+\infty\) when the public carrier graph supplies no finite
domination.  In that case the first failed domination row is retained as a
\(\Geom\)-owned total-graph coordinate.
The capacity is \(+\infty\) when its normalized constraint is infeasible.
Its reciprocal is the target-conditioned geometric resistance, with the
usual conventions.  To make the record valuations into cell coordinates,
set
\[
 h_g(C):=\sup\{h_g(R):R\text{ is a finite stopped record carried by }C\},
 \qquad
 \kappa_g(C):=\sup\{\kappa_g(R):R\text{ is such a record}\},
\]
with supremum zero for the empty family.  The geometric formation tuple is
\begin{equation}
 \label{eq:geom-formation-tuple}
 \mathfrak I_{\Sigma}^{\Geom}(C)
 :=\bigl(G_{\Sigma}^{\Geom},
          \operatorname{Cap}_{\Sigma}^{\Geom},
          \operatorname{Ker}_{\Sigma}^{\Geom},
          \operatorname{Gc}_{\Sigma}^{\Geom},
          h_g,\kappa_g\bigr)(C).
\end{equation}
\(h_g\) and \(\kappa_g\) are valuation coordinates.  They are not absence
detectors: their role is to quantify the cost of a target-visible geometric
current after that current has been detected.
\end{definition}

\begin{theorem}[Exhaustion and measurement of geometric singular structure]
\label{thm:continuum-geom-normal-form-exhaustion}
Every target-aligned \(\Geom\)-headed atomic component has a unique ordered
totalized record in \eqref{eq:target-geometric-normal-form}.  Its branch flag
is the first active ordinary-range, target-visible kernel/cokernel, or
graph-completion test; all simultaneous detector coordinates remain in the
record.  On every ordinary form-domain cell, the
target-aligned/target-orthogonal form projection displayed there is unique.
Moreover:
\begin{enumerate}[label=\textup{(G\arabic*)}]
\item the finite cylinder hierarchy computes nested rational enclosures for
all coordinates of \(\mathfrak I_{\Sigma}^{\Geom}\); the capacity and
resistance are the primal and Fenchel-dual values of the generated form
program, \(h_g\) is a finite stopped form integral, and \(\kappa_g\) is the
generated carrier-domination value;
\item \(\Geom\) is absent from singular formation on \(C\) exactly when
\[
 G_{\Sigma}^{\Geom}
 =\operatorname{Ker}_{\Sigma}^{\Geom}
 =\operatorname{Gc}_{\Sigma}^{\Geom}=0;
\]
\item a nonzero ordinary range value is either admitted as physical
production by the carrier-domination graph or becomes a signed feeding
coordinate; a nonzero kernel or graph value is the corresponding proper
target-cyclic successor.
\end{enumerate}
\end{theorem}

\begin{proof}
The form completion gives the orthogonal projection in
\eqref{eq:target-geometric-normal-form}.  Totalization of \(K^{1/2}\), its
range inclusion, and its quotient gives the exhaustive ordinary, kernel, and
graph-boundary partition.  Target incidence deletes the orthogonal summand
and retains every failed deletion as a typed edge.  The capacity program is
the constrained quadratic-form problem; completing the square gives its
resistance dual by \cref{thm:resolvent-resistance}.  Cylinder approximation
and \cref{thm:stagewise-target-core-computation} give the enclosures.
Finite form evaluation computes \(h_g\), while the ordered domination
program computes \(\kappa_g\).  The
three displayed zero pairings annihilate every geometric source--shell path,
and prospective accountability proves absence.  Conversely, each nonzero
coordinate is itself such a path.  Carrier admission and total-graph polarity
give \textup{(G3)}.
\end{proof}

\section{Target-aligned directed-feeding normal form}
\label{sec:continuum-caus-normal-form}

\begin{definition}[Directed exact--cohomological--recurrent decomposition]
\label{def:target-caus-feeding-normal-form}
Glue all compatible directed transport, pressure, imposed drift, and
nonnormal response currents with their signs.  Let
\(\mathcal B_{\Caus}=\overline{\operatorname{ran}\delta}\) be the closed
shift-coboundary space on the target quotient and let \(q_{\Caus}\) be its
quotient map.  The canonical directed normal form is the ordered typed record
\begin{equation}
 \label{eq:target-caus-normal-form}
 \operatorname{NF}_{\Caus}(F_{\Caus})
 :=\bigl(q_{\Caus}F_{\Caus},K_{\Caus}^{\rm stor},
         F_{\Caus}^{\rm rec},F_{\Caus}^{\rm gc},N_\Sigma\bigr).
\end{equation}
The first coordinate is the reduced feeding-cohomology class.  When it
vanishes, the total coboundary graph returns an exact storage potential
\(K_{\Caus}^{\rm stor}\), or else its first graph value is placed in the
fourth coordinate.  The third coordinate is the positive-support
first-return law of a nonzero cohomology or storage value,
\(F_{\Caus}^{\rm gc}\) is the first nonordinary gluing, support, resolvent,
or return value, and \(N_\Sigma\in\Null_\Sigma\).  Jordan decomposition is
performed only after the signed sum and coboundary projection.
\end{definition}

\begin{definition}[Directed formation invariants]
\label{def:target-caus-formation-invariants}
Let \(D_{\Phys}\) be the admitted nonnegative physical production on the
same target cell, and let \(\mathscr R_{\Caus}(C)\) be the generated retained
\(\Caus\)-variation space of the full generator.  Define
\begin{align}
 H_{\Sigma}^{\Caus}(C)
 &:=\|[F_{\Caus}]\|_{\rm red},
 &S_{\Sigma}^{\Caus}(C)
 &:=\|K_{\Caus}^{\rm stor}\|_{\rm red},
 \label{eq:caus-cohomology-storage-invariants}\\
 \beta_{\Sigma}^{\Caus}(C)
 &:=\inf\{a\ge0:(q_{\Caus}F_{\Caus})_+\le aD_{\Phys}\},
 &\delta_{\Sigma}^{\rm cyc}(C)
 &:=\inf_\gamma\int p_{\Phys}\,\dd\gamma,
 \label{eq:caus-radius-cycle-invariants}\\
 \Psi_{\Sigma}^{\Caus}(C;\lambda)
 &:=\sup_{\substack{K\in\mathscr R_{\Caus}(C)\\
                    \|K\|_{\rm red}\le1}}
   \|P_{\mathscr O}(\lambda-L)^{-1}
      K(\lambda-L)^{-1}P_{\mathscr I}\|,
 &\operatorname{Gc}_{\Sigma}^{\Caus}(C)
 &:=\|F_{\Caus}^{\rm gc}\|_{\rm red}.
 \label{eq:caus-resolvent-graph-invariants}
\end{align}
The reduced norms are order-unit norms on the generated cylinder quotient.
The directed formation tuple is
\[
 \mathfrak I_{\Sigma}^{\Caus}
 :=(H_{\Sigma}^{\Caus},S_{\Sigma}^{\Caus},
    \beta_{\Sigma}^{\Caus},\delta_{\Sigma}^{\rm cyc},
    \Psi_{\Sigma}^{\Caus},\operatorname{Gc}_{\Sigma}^{\Caus}),
\]
together with the signed-dyadic support laws of these coordinates.
\end{definition}

\begin{theorem}[Exhaustion and measurement of directed singular structure]
\label{thm:continuum-caus-normal-form-exhaustion}
Every target-aligned \(\Caus\)-headed atomic component has a unique first
active coordinate in \eqref{eq:target-caus-normal-form} relative to the generated
coboundary space, return graph, and target kernel.  The finite evaluator
computes certified enclosures for
\(H_\Sigma^{\Caus},S_\Sigma^{\Caus},\beta_\Sigma^{\Caus},
\delta_\Sigma^{\rm cyc},\Psi_\Sigma^{\Caus}\), and
\(\operatorname{Gc}_\Sigma^{\Caus}\).  Directed structure is absent from
singular formation exactly when
\begin{equation}
 \label{eq:caus-absence-criterion}
 H_\Sigma^{\Caus}=S_\Sigma^{\Caus}
 =\operatorname{Gc}_\Sigma^{\Caus}=0
\end{equation}
and every retained resolvent input--output matrix element vanishes.
For a nonzero class, \(\beta_\Sigma^{\Caus}<1\) gives strict carrier
passivity, positive \(\delta_\Sigma^{\rm cyc}\) gives infinite recurrent
physical charge, and zero cycle price constructs the next target-aligned
return support.
\end{theorem}

\begin{proof}
The closed coboundary quotient, its total preimage graph, the recurrent
support graph, graph boundary, and target kernel form the ordered typed
partition in \eqref{eq:target-caus-normal-form}.  Signed gluing precedes all
these projections.  The sharp radius is the order domination problem of
\cref{thm:target-aligned-carrier-duality}; cycle price is
\cref{thm:target-current-circulation-duality}; and resolvent variations
factor through the same source by \cref{prop:nonnormal-source-invariance}.
Their finite primal--dual programs and
\cref{thm:stagewise-target-core-computation} compute the stated intervals.
Criterion \eqref{eq:caus-absence-criterion}, together with zero compressed
response, deletes every directed incidence path.  Conversely a nonzero
coordinate supplies its path and signed support.  The final alternatives are
the passivity and circulation dualities.
\end{proof}

\section{Target-aligned quotient and constraint normal form}
\label{sec:continuum-abs-normal-form}

\begin{definition}[Quotient horizontal--vertical normal form]
\label{def:target-abs-normal-form}
Let \(q:X\to Q\) be a generated constraint, gauge, or coarse quotient, with
pullback \(U\), generated averaging/right inverse \(P\), and totalized
projection \(\Pi_q=UP\).  For a public current \(J\), put
\[
 J_Q=PJU,
 \qquad E_q=JU-UJ_Q.
\]
On the ordinary graph the target-aligned quotient identity is
\begin{equation}
 \label{eq:target-abs-normal-form}
 J=UJ_QP+E_qP+J(I-\Pi_q).
\end{equation}
The first term is horizontal descent, the second is the intertwining defect,
and the third is vertical/fiber motion.  If an operation in this identity is
nonordinary, the ordered totalization replaces the identity by the typed
record
\[
 \operatorname{NF}_{\Abs}(J;q)
 :=(UJ_QP,E_qP,J(I-\Pi_q),J_q^{\rm coker},J_q^{\rm gc}),
\]
where \(J_q^{\rm coker}\) is the target-visible
failure of generated lifting or surjectivity, and \(J_q^{\rm gc}\) is the
first nonordinary quotient, constraint, gauge, or pressure graph value.
Each term is projected to the target incidence core after its ancestry is
fixed.
\end{definition}

\begin{definition}[Quotient formation invariants]
\label{def:target-abs-formation-invariants}
Define
\begin{align}
 A_{\Sigma}^{\rm desc}(C)&:=\|\Pi_\Sigma E_qP\|_{\rm red},
 &A_{\Sigma}^{\rm vert}(C)&:=\|\Pi_\Sigma J(I-\Pi_q)\|_{\rm red},
 \label{eq:abs-descent-vertical-invariants}\\
 A_{\Sigma}^{\rm coker}(C)&:=\|\Pi_\Sigma J_q^{\rm coker}\|_{\rm red},
 &A_{\Sigma}^{\rm gc}(C)&:=\|\Pi_\Sigma J_q^{\rm gc}\|_{\rm red},
 \label{eq:abs-cokernel-graph-invariants}\\
 A_{\Sigma}^{\rm tgt}(C)
 &:=\sup_{\lambda\in\Lambda_{\Sigma,\mathfrak b}}
       \|\lambda-\lambda\circ\Pi_q\|_{\rm ou,*},
 \label{eq:abs-target-descent-invariant}\\
 \operatorname{Cap}_{\Sigma}^{\ker q}(C)
 &:=\inf\{\mathcal E(v):v\in\ker q,
            \langle\lambda,v\rangle=1\},
 &\operatorname{Rank}_{\Sigma}(q;C)
 &:=\dim_{\rm cyl}\operatorname{span}
   \{\lambda\circ U\}_{\lambda\in\mathscr C^\infty}.
 \label{eq:abs-capacity-rank-invariants}
\end{align}
At finite depth, \(\dim_{\rm cyl}\) is the rank of the rational pairing
matrix; the limiting rank is its monotone generated value.
The quotient formation tuple is
\[
 \mathfrak I_{\Sigma}^{\Abs}
 :=(A_{\Sigma}^{\rm desc},A_{\Sigma}^{\rm vert},
 A_{\Sigma}^{\rm coker},A_{\Sigma}^{\rm gc},A_{\Sigma}^{\rm tgt},
 \operatorname{Cap}_{\Sigma}^{\ker q},\operatorname{Rank}_{\Sigma}(q)).
\]
\end{definition}

\begin{theorem}[Exhaustion and measurement of quotient singular structure]
\label{thm:continuum-abs-normal-form-exhaustion}
Every target-aligned \(\Abs\)-headed atomic component has exactly one first active value
in the ordered typed record extending \eqref{eq:target-abs-normal-form}.
The finite incidence and moment
hierarchies compute the four reduced current-defect norms, the target-descent
defect, kernel capacity, and target rank.  Quotient structure is absent from singular formation exactly
when
\begin{equation}
 \label{eq:abs-absence-criterion}
 A_{\Sigma}^{\rm desc}=A_{\Sigma}^{\rm vert}
 =A_{\Sigma}^{\rm coker}=A_{\Sigma}^{\rm gc}
 =A_{\Sigma}^{\rm tgt}=0
\end{equation}
---equivalently, the current defects vanish and every retained target
covector descends through \(q\).  On this cell the horizontal
term transfers the problem to \(Q\) without creating a source.  Every nonzero
invariant gives the corresponding intertwining, vertical, cokernel, or graph
successor with its original source ancestry.
\end{theorem}

\begin{proof}
Insert \(I=\Pi_q+(I-\Pi_q)\) on the right of \(J\), then add and subtract
\(UJ_QP\).  Totalization of \(U,P,q\), and their range graphs supplies the
cokernel and graph terms, proving exhaustiveness.  Rational pairing matrices
compute finite ranks; the remaining quantities are norms or constrained
energy minima in the Archimedean cylinder algebra and hence have the stated
enclosures.  Under \eqref{eq:abs-absence-criterion}, all vertical,
target-descent, and failed descent paths are target-null, so only exact
horizontal descent remains.
Conversely each nonzero block is detected by its defining target covector and
therefore supplies a proper incidence edge.
\end{proof}

\section{Target-aligned lift, scale, and compactness normal form}
\label{sec:continuum-lift-normal-form}

\begin{definition}[Continuum Lift normal form]
\label{def:target-lift-normal-form}
Fix the canonical evaluation compactification and a generated lift or chart
\((\Omega,\pi,\Lambda)\), with fiber averaging \(P_\kappa\) and
\(\Pi_\pi=UP_\kappa\).  Let \(J^\Omega\) denote the total current in the
lifted presentation.  On the ordinary lift graph, the first three
coordinates satisfy the exact identity
\begin{equation}
 \label{eq:target-lift-normal-form}
 J^\Omega=UJ_XP_\kappa+E_\pi P_\kappa
  +J^\Omega(I-\Pi_\pi).
\end{equation}
They are base transport, lift-intertwining defect, and within-fiber motion.
The ordinary base term retains the head of \(J_X\); the second and third
terms are \(\Lift\)-headed and retain the base ancestry.
The complete ordered typed record is
\[
 \operatorname{NF}_{\Lift}(J^\Omega;\pi)
 :=(UJ_XP_\kappa,E_\pi P_\kappa,J^\Omega(I-\Pi_\pi),
    J_{\rm conc},J_{\rm osc},J_{\rm sc}^{\rm coh},J_{\rm chart}^{\rm gc}).
\]
Its next two coordinates are the concentration and oscillation
coordinates of the fixed graph-complete compactification.  The sixth is the
nonexact scale/modulation cohomology after covariant carrier conjugation, and
the last is the first nonordinary base-descent, chart, scaling, product,
localization, or realization value.  Tail and terminal descendants are transferred to the
\(\Bdry\) normal form after retaining their \(\Lift\) ancestry.

These are sequential residual coordinates, not five independent
descriptions of the same current.  On the common ordinary graph, first
subtract the exact base/interface/fiber identity from the graph-completed
stopped incidence current.  Jordan--polar decomposition of the resulting
compactification defect defines \(J_{\rm conc}\) and then \(J_{\rm osc}\).
After subtracting both, split the scale row into its generated exact
coboundary \(\Delta_{\rm sc}K_l\) and the selected nonexact representative
\(J_{\rm sc}^{\rm coh}\).  The remaining coordinate is, by definition,
\(N_l\in\Null_\Sigma\).  Thus the
public graph-completed current has the exact reconstruction
\begin{equation}
 \label{eq:target-lift-exact-current-reconstruction}
 \operatorname{pr}_{\{\Lift\}}\operatorname{Inc}^{\Sigma}
 =\operatorname{pr}_{\{\Lift\}}\bigl(E_\pi P_\kappa
       +J^\Omega(I-\Pi_\pi)\bigr)
  +\operatorname{pr}_{\{\Lift\}}J_{\rm conc}
  +\operatorname{pr}_{\{\Lift\}}J_{\rm osc}
  +\Delta_{\rm sc}K_l
  +\operatorname{pr}_{\{\Lift\}}J_{\rm sc}^{\rm coh}+N_l,
 \qquad \Pi_\Sigma N_l=0.
\end{equation}
Every subtraction and quotient in this construction is performed in the
expanded derivation, so each nonnull summand retains its complete ancestry.
\end{definition}

\begin{definition}[Lift and compactness formation invariants]
\label{def:target-lift-formation-invariants}
Define
\begin{align}
 L_{\Sigma}^{\rm int}(C)&:=\|\Pi_\Sigma E_\pi P_\kappa\|_{\rm red},
 &L_{\Sigma}^{\rm fib}(C)&:=\|\Pi_\Sigma J^\Omega(I-\Pi_\pi)\|_{\rm red},
 \label{eq:lift-interface-fiber-invariants}\\
 L_{\Sigma}^{\rm conc}(C)&:=\|\Pi_\Sigma J_{\rm conc}\|_{\rm var},
 &L_{\Sigma}^{\rm osc}(C)&:=\|\Pi_\Sigma J_{\rm osc}\|_{\rm var},
 \label{eq:lift-concentration-oscillation-invariants}\\
 L_{\Sigma}^{\rm stor}(C)&:=\|K_l\|_{\rm red},
 &L_{\Sigma}^{\rm sc}(C)&:=\|[J_{\rm sc}]\|_{\rm red},
 \label{eq:lift-storage-scale-invariants}\\
 L_{\Sigma}^{\rm gc}(C)&:=\|\Pi_\Sigma J_{\rm chart}^{\rm gc}\|_{\rm red}.&&
 \label{eq:lift-scale-graph-invariants}
\end{align}
The variation norms are computed on the countable cylinder algebra of the
fixed compactification.  The scale valuation \(\kappa_P\) and record-scale
velocity \(a=\dot r/r\) are retained as coordinates of
\(L_\Sigma^{\rm sc}\).
The lift formation tuple is
\[
 \mathfrak I_{\Sigma}^{\Lift}
 :=(L_{\Sigma}^{\rm int},L_{\Sigma}^{\rm fib},
 L_{\Sigma}^{\rm conc},L_{\Sigma}^{\rm osc},
 L_{\Sigma}^{\rm stor},L_{\Sigma}^{\rm sc},L_{\Sigma}^{\rm gc}).
\]
\end{definition}

\begin{theorem}[Exhaustion and measurement of compactness structure]
\label{thm:continuum-lift-normal-form-exhaustion}
Relative to the fixed continuum presentation, every target-visible
\(\Lift\)-headed weak limit, concentration, oscillation, defect measure,
chart degeneration, blow-up profile, and scale escape has a unique first active coordinate in the
ordered typed record extending \eqref{eq:target-lift-normal-form}.  The seven invariants in
\cref{def:target-lift-formation-invariants} are computed by monotone
cylinder-variation and ordered moment enclosures.  Lift/compactness structure
is absent from singular formation exactly when all seven vanish.  On this cell
the ordinary base term is transferred to the normal form of its inherited
head source.  A failed base descent is part of \(L_Sigma^{\rm gc}\).
Otherwise the least nonzero invariant explicitly identifies the interface,
fiber, concentration, oscillation, scale, or graph mechanism used by the
target-active cascade.
\end{theorem}

\begin{proof}
The first three terms are the exact lift identity obtained by inserting
\(UP_\kappa\).  On the fixed canonical evaluation compactification, the
Jordan--polar decomposition of the residual graph-completed current uniquely
gives its positive concentration part and signed oscillation/polarization
part.  Subtracting these parts leaves the scale row; the generated
coboundary quotient separates the exact scale-storage current from its
selected cohomology representative.  Target projection sends the final remainder to
zero.  This proves \eqref{eq:target-lift-exact-current-reconstruction} and
exhaustiveness.  Total variation on a countable cylinder algebra
is the supremum over finite rational partitions, and the remaining norms are
ordered cylinder programs; hence the finite hierarchy computes all seven.
Their simultaneous vanishing deletes every nonbase lift incidence path, and
ordinary base transport is evaluated by its inherited source packet.  Failed
base descent is a nonordinary operation value and therefore contributes to
\(L_Sigma^{\rm gc}\).
Every nonzero value carries its detecting covector and produces the stated
successor.
\end{proof}

\section{Target-aligned boundary and terminal normal form}
\label{sec:continuum-bdry-normal-form}

\begin{definition}[Oriented boundary normal form]
\label{def:target-bdry-normal-form}
Let \(\mathscr K\) be a finite target-retaining carrier complex.  Sum its
cellwise Stieltjes identities with their orientations.  The boundary current
has the canonical decomposition
\begin{equation}
 \label{eq:target-bdry-normal-form}
 J^{\Bdry}
 =\partial_{\rm int}\mathscr K
  +J_{\rm ext}+J_{\rm tail}+J_{\rm term}+J_{\rm tr}^{\rm gc}.
\end{equation}
Paired internal faces constitute \(\partial_{\rm int}\mathscr K\) and cancel
identically.  The remaining terms are exterior physical flux, compactified
tail flux, continuation/terminal-shell flux, and the first nonordinary trace,
exhaustion, overlap, boundary-condition, or terminal-realization value.
Faces at intersections are assigned by the fixed half-open priority
\[
 \mathrm{terminal}\ \succ\ \mathrm{tail}\ \succ\ \mathrm{exterior};
\]
equivalently, use the corresponding inclusion--exclusion strata.  Hence a
corner is counted once while retaining every incident ancestry label.
Killing or support loss is the zero terminal-value specialization of
\(J_{\rm term}\).
\end{definition}

\begin{definition}[Boundary formation invariants]
\label{def:target-bdry-formation-invariants}
Define
\begin{align}
 B_{\Sigma}^{\rm ext}(C)&:=\|\Pi_\Sigma J_{\rm ext}\|_{\rm var},
 &B_{\Sigma}^{\rm tail}(C)&:=\|\Pi_\Sigma J_{\rm tail}\|_{\rm var},
 \label{eq:bdry-exterior-tail-invariants}\\
 B_{\Sigma}^{\rm term}(C)&:=\|\Pi_\Sigma J_{\rm term}\|_{\rm var},
 &B_{\Sigma}^{\rm gc}(C)&:=\|\Pi_\Sigma J_{\rm tr}^{\rm gc}\|_{\rm red},
 \label{eq:bdry-terminal-graph-invariants}\\
 \operatorname{TailCap}_\Sigma(C)
 &:=\inf\{\mathcal E(v):v|_{\partial_\infty}=1
       \text{ on a retained tail shell}\}.
 \label{eq:bdry-tail-capacity}
\end{align}
All fluxes are measured after internal-face cancellation.  A positive
exterior flux is physical expenditure only when the public carrier identity
dominates it.
The boundary formation tuple is
\[
 \mathfrak I_{\Sigma}^{\Bdry}
 :=(B_{\Sigma}^{\rm ext},B_{\Sigma}^{\rm tail},
 B_{\Sigma}^{\rm term},B_{\Sigma}^{\rm gc},
 \operatorname{TailCap}_{\Sigma}).
\]
\end{definition}

\begin{theorem}[Exhaustion and measurement of boundary singular structure]
\label{thm:continuum-bdry-normal-form-exhaustion}
Every target-aligned \(\Bdry\)-headed atomic component belongs uniquely to the exterior,
tail, terminal, or graph block of \eqref{eq:target-bdry-normal-form} after
internal cancellation.  The finite exhaustion, trace, carrier, and capacity
programs compute certified enclosures for the five invariants in
\cref{def:target-bdry-formation-invariants}.  Boundary structure is absent
from singular formation exactly when
\begin{equation}
 \label{eq:bdry-absence-criterion}
 B_{\Sigma}^{\rm ext}=B_{\Sigma}^{\rm tail}
 =B_{\Sigma}^{\rm term}=B_{\Sigma}^{\rm gc}=0.
\end{equation}
A nonzero exterior invariant is charged or returned as a \(\Bdry\)
successor; a nonzero tail invariant generates its exhaustion/compactification
successor; and a nonzero terminal invariant realizes target contact only
through the totalized terminal-flux graph.
\end{theorem}

\begin{proof}
The cellular boundary identity cancels every oppositely oriented common
face.  The fixed half-open priority partitions the remaining faces into
disjoint exterior, compactification-tail, and terminal strata; the
inclusion--exclusion identity preserves their signed sum.  Totalization
supplies the graph block.  Jordan decomposition on these disjoint supports
gives uniqueness.  Variation is computed by finite
rational partitions, capacity by its Dirichlet primal--dual program, and
graph values by the ordered hierarchy.  Criterion
\eqref{eq:bdry-absence-criterion} annihilates every boundary incidence path.
The three nonzero routes follow from carrier admission, exhaustion
totalization, and the killed-flux realization law.
\end{proof}

\section{The continuum singularity-formation profile}
\label{sec:continuum-formation-profile}

\begin{definition}[The thirty-one target-aligned formation cells]
\label{def:continuum-formation-profile}
For each nonempty support \(A\subseteq\Channels\), let
\(\mathscr C_A(E):=\mathscr S_{A,\Sigma}^{\rm reach,sat}(E)\) be the
source-preserving closure of the normalized same-origin pre-contact records
generated by the public presentation with exact support \(A\).  Ordinary path
values, first nonordinary realization values, and their origin tags remain
distinct coordinates of this reachable cell.  Its atomic invariant vector is
\begin{equation}
 \label{eq:atomic-continuum-formation-vector}
 \mathbf I_A^\Sigma
 :=\bigoplus_{c\in A}\mathfrak I_{\Sigma}^{c}\big|_{\mathscr C_A(E)},
\end{equation}
where the five summands are the invariant tuples in
\cref{def:target-geometric-formation-invariants,%
def:target-caus-formation-invariants,%
def:target-abs-formation-invariants,%
def:target-lift-formation-invariants,%
def:target-bdry-formation-invariants}.  Adjoin the aggregate signed carrier
class, recurrent cycle price, same-origin realization graph, and cofinal shell
coordinate.  The resulting invariant image is the
\emph{formation profile}
\begin{equation}
 \label{eq:continuum-formation-profile}
 \mathfrak F_A^\Sigma(E)
 :=\left\{
  (\mathbf I_A^\Sigma,[H]_A,\delta_A^{\rm cyc},
   \mathbf R,\operatorname{Shell}_\Sigma):
  C\in\mathscr C_A(E)\right\}.
\end{equation}
Scalar valuation coordinates take values in their order compactifications
\([0,\infty]\); graph and measure coordinates use the compact topologies
already fixed by the minimal realization.  Thus
\eqref{eq:continuum-formation-profile} is evaluated on the specified
source-preserving compact product, not in an ambient norm topology.
A coordinate is a \emph{detector} when its vanishing is part of the
corresponding absence criterion.  Thus capacity, resistance, rank,
domination radius, and cycle price are valuations of a detected mechanism,
not detectors by themselves.  Let \(\mathbf D_A^\Sigma\) denote the
subvector consisting of
\[
\begin{aligned}
 &(G^{\Geom},\operatorname{Ker}^{\Geom},\operatorname{Gc}^{\Geom}),\\
 &(H^{\Caus},S^{\Caus},\Psi^{\Caus},\operatorname{Gc}^{\Caus}),\\
 &(A^{\rm desc},A^{\rm vert},A^{\rm coker},A^{\rm gc},A^{\rm tgt}),\\
 &(L^{\rm int},L^{\rm fib},L^{\rm conc},L^{\rm osc},L^{\rm stor},
   L^{\rm sc},L^{\rm gc}),\\
 &(B^{\rm ext},B^{\rm tail},B^{\rm term},B^{\rm gc})
\end{aligned}
\]
for the sources belonging to \(A\).
A profile point is \emph{realized} when its origin and total-graph rows
diagonally reconstruct a public solution path.  It is \emph{formation
active} when its normalized shell coordinate is nonzero and at least one
coordinate of \(\mathbf D_A^\Sigma\) or the aggregate feeding/storage class
is nonzero.
\end{definition}

\begin{theorem}[Thirty-one-cell target-aligned reachability ledger]
\label{thm:thirty-one-cell-target-aligned-ledger}
For every nonempty \(A\subseteq\Channels\), the cell
\(\mathscr C_A(E)\) is evaluated by the detector subvector
\(\mathbf D_A^\Sigma\), its complete joint conditional increments, and its
physical valuation.  Exactly one of the following values holds:
\begin{enumerate}[label=\textup{(L\arabic*)}]
\item every target detector vanishes on \(\mathscr C_A(E)\), and the support
  \(A\) has no target-aligned reachable occurrence;
\item some target detector is positive, and there is a finite same-origin
  pre-contact record \(R_{A}\) generated by the public presentation whose
  complete joint frontier has that positive coordinate and source support
  \(A\).
\end{enumerate}
For \textup{(L2)}, the physical carrier of \(R_A\) is evaluated by its
finite occurrence atom.  Positive mass is charged, while zero mass enters
the equality record generated by the same source word.  The statement holds
separately for all \(2^5-1=31\) supports; it is not an aggregate assertion
about the five isolated channels.
\end{theorem}

\begin{proof}
Each atomic coordinate is generated by exactly one of the five normal forms;
the complete source word is retained under joining, localization, quotient,
lift, boundary passage, and source-preserving closure.  The mixed coordinates
are formed before projection, so composition cannot be discarded as five
separate null tests.  Apply the density equivalence
\eqref{eq:actual-source-zero-equivalence}--
\eqref{eq:actual-source-positive-equivalence} to every detector coordinate
of the fixed support \(A\).  The positive row returns \(R_A\) with its least
contextual occurrence.  The physical alternative is
\cref{thm:actual-source-price-equality-decision}.  Since every nonempty
subset of the five-source alphabet occurs exactly once in this construction,
the thirty-one statements are exhaustive and disjoint by support.
\end{proof}

\begin{theorem}[Cellwise public-presentation evaluation of all thirty-one supports]
\label{thm:cellwise-public-presentation-thirty-one-evaluation}
Write the five-source alphabet in its fixed order as
\(\mathsf s=(\Geom,\Caus,\Abs,\Lift,\Bdry)\), and, for each nonzero
bit vector \(\varepsilon\in\{0,1\}^{5}\setminus\{0\}\), put
\[
 A_\varepsilon:=\{\mathsf s_i:\varepsilon_i=1\},\qquad
 \mathsf{Cell}_{\varepsilon}:=
 \bigcap_{\varepsilon_i=1}\mathfrak I^{\mathsf s_i}_\Sigma
 \ \cap\!\!\bigcap_{\varepsilon_i=0}
       \ker\mathfrak I^{\mathsf s_i}_\Sigma .
\]
For each of the thirty-one values of \(\varepsilon\), the public
presentation determines the following finite-prefix calculation and no
additional analytic premise:
\begin{enumerate}[label=\textup{(E\arabic*)}]
\item form the complete joint pre-contact frontier before applying the five
  source projections, retain precisely the source word \(A_\varepsilon\),
  and evaluate every detector in
  \(\mathbf D_{A_\varepsilon}^{\Sigma}\) on its finite stopped records;
\item take the monotone source-preserving closure of those finite values;
  a zero value deletes the corresponding target-incidence edge, whereas a
  positive value returns its least contextual finite occurrence together with
  the public equation row, target covector, joint frontier, source word, and
  unnormalised shell cocycle;
\item evaluate the unique physical carrier coordinate of that occurrence.
  Positive carrier mass is charged to its source word.  Zero carrier mass is
  evaluated by the equality row of the \emph{same} joint record, whose
  target-null value deletes the edge and whose nonzero value is reinserted at
  its least unprocessed conditional increment.
\end{enumerate}
Consequently every cell has exactly one structural verdict
\[
 \operatorname{CellEval}_{\varepsilon}\in
 \left\{
   \operatorname{Absent}(A_\varepsilon),\quad
   \operatorname{Present}(R_\varepsilon,
       \lambda_\varepsilon,Q_\varepsilon,
       \operatorname{Anc}(R_\varepsilon),p_{\Phys}(R_\varepsilon))
 \right\};
\]
the second value is an actual finite same-origin target-aligned occurrence
generated by the public presentation.  Thus the displayed family is an
explicit source-by-source evaluation of
\[
 \{\Geom,\Caus,\Abs,\Lift,\Bdry\}\setminus\{\varnothing\}
 =\{A_\varepsilon:0\ne\varepsilon\in\{0,1\}^{5}\},
\]
not an aggregate test of five isolated channels.
\end{theorem}

\begin{proof}
The free five-constructor grammar assigns every finite stopped term one and
only one source word.  Joining is performed before projection, so a target
coordinate created only by composition is a coordinate of the joint frontier
and is retained in the unique \(A_\varepsilon\) containing its complete
ancestry.  The finite-record zero/positive equivalence
\eqref{eq:actual-source-zero-equivalence}--
\eqref{eq:actual-source-positive-equivalence} proves the dichotomy in
\textup{(E1)}--\textup{(E2)} for that fixed word.  Source-preserving
completion cannot add a sixth owner, while a target-visible limit coordinate
has a finite detecting prefix and hence re-enters the same word.  The
physical alternative in \textup{(E3)} is exactly
\cref{thm:actual-source-price-equality-decision}; its re-entry clause is a
conditional-joint increment rather than a new terminal.  This proves the
verdict for the chosen \(\varepsilon\).  The nonzero bit vectors partition
the nonempty source words, giving all thirty-one cases.
\end{proof}

\begin{theorem}[Explicit theorem family for the thirty-one source supports]
\label{thm:explicit-thirty-one-source-theorem-family}
For every support in the following table, let \(\mathsf T_A\) denote the
assertion that the target-aligned contribution with \emph{exact} source word
\(A\) is evaluated by its complete joint detector, finite stopped
same-origin records, and the single physical carrier.
\[
\begin{array}{c|l}
 |A| & \text{theorem instances }\mathsf T_A\\ \hline
 1 & G,\ C,\ A,\ L,\ B\\
 2 & GC,\ GA,\ GL,\ GB,\ CA,\ CL,\ CB,\ AL,\ AB,\ LB\\
 3 & GCA,\ GCL,\ GCB,\ GAL,\ GAB,\ GLB,\ CAL,\ CAB,\ CLB,\ ALB\\
 4 & GCAL,\ GCAB,\ GCLB,\ GALB,\ CALB\\
 5 & GCALB
\end{array}
\]
Here a letter abbreviates respectively
\(\Geom,\Caus,\Abs,\Lift,\Bdry\), and concatenation denotes set union,
not ordered multiplication.  Every one of these thirty-one theorem instances
has the exact conclusion
\begin{equation}
 \label{eq:thirty-one-explicit-source-verdict}
 \mathsf T_A:\qquad
 \begin{cases}
 \mathbf D_A^\Sigma=0
   &\Longleftrightarrow\ \text{no target-aligned reachable occurrence has
      exact support }A,\\
 \mathbf D_A^\Sigma\ne0
   &\Longleftrightarrow\ \text{there is a least finite same-origin
      public-presentation occurrence }R_A\\[-1mm]
   &\hspace{17mm}\text{with exact support }A\text{ and nonzero target
      pairing.}
 \end{cases}
\end{equation}
The occurrence \(R_A\) carries its complete joint frontier and source word.
Its carrier valuation is either positive and charged, or zero and processed
by the equality rule on that same joint record; no other outcome is admitted.
\end{theorem}

\begin{proof}
The five primitive detector packets are generated from the public weak and
total graphs.  For a fixed row \(A\) in the table, form their joint record,
then retain precisely the terms whose complete provenance word is \(A\).
This is a finite typed construction and is therefore a member of
\(\mathsf{Proof}_\Sigma(\mathfrak P)\).  The zero/positive equivalence for
the reachable source closure proves the two lines of
\eqref{eq:thirty-one-explicit-source-verdict}; importantly, it is applied to
the joint record and not to the five packets separately.  The physical
alternative is \cref{thm:actual-source-price-equality-decision}.  This proves
\(\mathsf T_A\) for the selected row.  The displayed table lists every
nonempty subset of a five-element alphabet exactly once, so the same finite
derivation proves all thirty-one instances and leaves no unlisted source
word.
\end{proof}

\begin{definition}[Quantitative local detector metric]
\label{def:quantitative-public-detector-metric}
For a finite public tuple \(Z=(z_1,\ldots,z_m)\) set
\[
\left\|Z\right\|_{\rm det}:=\sum_{i=1}^{m}2^{-i}\min\left\{1,d_i(z_i,0)\right\},
\]
where \(d_i\) is the bounded evaluation metric supplied by the public
presentation.  The metric is used only on finite tuples and has zero exactly
when each displayed public coordinate is zero.
\end{definition}
\begin{theorem}[Exact pure-geometric target-current accountability]
\label{thm:quantitative-source-cell-g}
Let \(R=[s,t]\in\mathfrak D_{A_g,\Sigma,n}^{\rm pre}(E)\) be a finite
same-origin stopped record with exact source word \(A_g=\{\Geom\}\), and let
\(C_R\) denote its normalized target cell.
First evaluate the totalized common-form-domain graph for
\((K_r^{1/2}\zeta_r,K_r^{1/2}D\Phi_R(u_r))\) on \([s,t]\).  If it has a
nonordinary value, retain its first such value, with its \(\Geom\) ancestry,
and stop this branch.  Thus the following formulas concern exactly the
ordinary value of that graph and do not silently assume form-domain
membership.  On this value write
\[
 J_r^{\Geom}=-K_r\zeta_r,\qquad
 \lambda_{R,r}=D\Phi_R(u_r),
\]
and define the signed target current, its oriented amplitude, and its
conditional-incidence detector by
\begin{align}
 q_g(R)&:=-\int_s^t
   \langle K_r^{1/2}\lambda_{R,r},K_r^{1/2}\zeta_r\rangle\,\dd r,
 \qquad a_g(R):=(q_g(R))_+,
 \label{eq:geom-signed-current-amplitude}\\
 \delta_g(R)&:=
 \left\|\left(G_\Sigma^{\Geom}(C_R),
 \operatorname{Ker}_\Sigma^{\Geom}(C_R),
 \operatorname{Gc}_\Sigma^{\Geom}(C_R),
 \operatorname{Inc}_{A_g}^{\Sigma}(R)\right)\right\|_{\rm det}.
 \label{eq:geom-incidence-detector}
\end{align}
The quantities \(d_g(R),h_g(R)\), and \(\kappa_g(R)\) are those of
\eqref{eq:geom-production-probe-cost}--
\eqref{eq:geom-carrier-admission-value}, and
\(p_g(R):=p_{\Phys}(R)\).

For every \(\theta>0\), the public form identity gives
\begin{equation}
 \label{eq:geom-exact-source-conversion}
 a_g(R)
 \le \frac{\theta}{2}d_g(R)
     +\frac{1}{2\theta}h_g(R).
\end{equation}
If \(\kappa_g(R)<\infty\), the generated carrier row further gives
\begin{equation}
 \label{eq:geom-admitted-source-conversion}
 a_g(R)
 \le \frac{\theta\kappa_g(R)}2p_g(R)
     +\frac{1}{2\theta}h_g(R).
\end{equation}
If \(\kappa_g(R)=\infty\), the record retains the first failed
carrier-domination coordinate as a nonzero \(\Geom\)-owned graph value.
More explicitly,
\[
 p_g(R)=0=d_g(R)\Longrightarrow\kappa_g(R)=0,
 \qquad
 p_g(R)=0<d_g(R)\Longrightarrow\kappa_g(R)=+\infty.
\]

For a pairwise disjoint ordered family
\(\mathcal F_N=(R_1,\ldots,R_N)\) with exact word \(A_g\), put
\begin{align*}
 \mathsf A_g(\mathcal F_N)&:=\sum_{j=1}^Na_g(R_j),&
 \mathsf D_g^{\rm form}(\mathcal F_N)&:=\sum_{j=1}^Nd_g(R_j),\\
 \mathsf H_g(\mathcal F_N)&:=\sum_{j=1}^Nh_g(R_j),&
 \mathsf P_g(\mathcal F_N)&:=\sum_{j=1}^Np_g(R_j),\\
 \kappa_g^*(\mathcal F_N)&:=\max_{1\le j\le N}\kappa_g(R_j).
\end{align*}
Then
\begin{equation}
 \label{eq:geom-cumulative-source-conversion}
 \mathsf A_g(\mathcal F_N)
 \le \frac{\theta}{2}\mathsf D_g^{\rm form}(\mathcal F_N)
     +\frac{1}{2\theta}\mathsf H_g(\mathcal F_N).
\end{equation}
Because the records are pairwise disjoint subrecords of one orbit in the
physical stratum of total mass at most \(E\), finite additivity of the single
physical carrier gives, without an additional hypothesis,
\begin{equation}
 \label{eq:geom-physical-subrecord-ceiling}
 \mathsf P_g(\mathcal F_N)\le E.
\end{equation}
In particular, if the public carrier graph proves
\[
 \kappa_g^*(\mathcal F_N)\le\kappa_*<\infty,\qquad
 \mathsf H_g(\mathcal F_N)\le H_g
\]
for every \(N\), then
\begin{equation}
 \label{eq:geom-soft-cascade-ceiling}
 \mathsf A_g(\mathcal F_N)
 \le \frac{\theta\kappa_*E}{2}+\frac{H_g}{2\theta}
 \qquad(N\ge1,\ \theta>0).
\end{equation}
Consequently, there is no such family satisfying in addition
\(q_g(R_j)=1\) for every \(j\).  This is precisely the pure-\(\Geom\)
normalized-crossing case; the conclusion does not identify an arbitrary
nonzero geometric detector with a unit shell increment.

The complementary cases are retained without reinterpretation:
\begin{enumerate}[label=\textup{(PG\arabic*)}]
\item \(a_g(R)>0\) or \(\delta_g(R)>0\) returns the finite
  presentation-generated \(A_g\)-occurrence \(R\) with its complete joint
  frontier;
\item if \(\sup_N\kappa_g^*(\mathcal F_N)=+\infty\), then either one finite
  record already has \(\kappa_g=+\infty\), or the record that first makes the
  running maximum cross each integer level is returned as a finite
  \(\Geom\)-owned carrier-amplification record.  If instead
  \(\kappa_g^*\le\kappa_*<\infty\),
  then unbounded \(\mathsf D_g^{\rm form}\) forces unbounded
  \(\mathsf P_g\), contradicting \eqref{eq:geom-physical-subrecord-ceiling};
\item unbounded \(\mathsf H_g\) has a least finite prefix crossing each
  integer level; that complete prefix is returned as a geometric
  shell-probe-cost record;
\item a singularity verdict is produced only when the same complete
  \(A_g\)-record lies in the ordinary same-origin realization range and has
  the noncontinuation value of the public continuation graph.
\end{enumerate}
Consequently the theorem certifies pure-\(\Geom\) unit-crossing cascade
absence from the explicit bounds in \eqref{eq:geom-soft-cascade-ceiling}.
Every failed bound remains a finite source-owned structural record.  A graph
value, profile value, or conditional increment becomes an orbit only through
the realization intersection in \textup{(PG4)}.
\end{theorem}

\begin{proof}
The ordinary form identity \(J^{\Geom}=-K\zeta\) gives
\eqref{eq:geom-signed-current-amplitude}.  Since the positive part is bounded
by the absolute value,
\[
 a_g(R)\le
 \int_s^t
 \|K_r^{1/2}\lambda_{R,r}\|\,
 \|K_r^{1/2}\zeta_r\|\,\dd r.
\]
The scalar inequality \(xy\le\theta y^2/2+x^2/(2\theta)\), integrated on
\([s,t]\), is exactly \eqref{eq:geom-exact-source-conversion}.  It uses only
the generated form pairing.  When \(\kappa_g(R)<\infty\), substitute
\(d_g(R)\le\kappa_g(R)p_g(R)\) to obtain
\eqref{eq:geom-admitted-source-conversion}.  When the infimum in
\eqref{eq:geom-carrier-admission-value} is infinite, the totalized
domination graph has its nonordinary value, with the original \(\Geom\)
ancestry.

Summing \eqref{eq:geom-exact-source-conversion} over the finite family gives
\eqref{eq:geom-cumulative-source-conversion}; no empirical normalization or
limit is taken.  The stopped form costs are additive on disjoint records,
and the physical valuation is the restriction of one
nonnegative additive carrier measure on the orbit.  Pairwise disjointness
therefore gives \eqref{eq:geom-physical-subrecord-ceiling} by monotonicity
against the total stratum mass.  Under the two displayed public bounds,
\(\mathsf D_g^{\rm form}\le\kappa_*\mathsf P_g\le\kappa_*E\).
Substitution proves \eqref{eq:geom-soft-cascade-ceiling}.  Under the additional
unit-current hypothesis \(q_g(R_j)=1\),
\(\mathsf A_g(\mathcal F_N)=N\).  The fixed right-hand side of
\eqref{eq:geom-soft-cascade-ceiling} excludes such a family for all \(N\).

If the nondecreasing sequence \(\kappa_g^*(\mathcal F_N)\) is unbounded,
then either one of its entries is infinite or its first crossing of each
integer level occurs at a finite prefix.  These are precisely the failed
domination and finite carrier-amplification records in \textup{(PG2)}.  If
it is bounded, summing
\(d_g(R_j)\le\kappa_*p_g(R_j)\) and using
\eqref{eq:geom-physical-subrecord-ceiling} bounds
\(\mathsf D_g^{\rm form}\).  Likewise, unbounded \(\mathsf H_g\) has a
first finite crossing of every integer level, proving \textup{(PG3)}.  The range--kernel--graph exhaustion theorem
\cref{thm:continuum-geom-normal-form-exhaustion} proves that a nonzero
detector or current returns the finite record asserted in \textup{(PG1)}.
Finally, ordinary realization and the noncontinuation coordinate are exactly
the conditions in
\eqref{eq:actual-contact-cell-realization-intersection}; this proves
\textup{(PG4)} without reconstructing an orbit from a closure point.
\end{proof}

\begin{theorem}[Canonical amplitude identification for exact pure geometry]
\label{thm:geom-canonical-amplitude-identification}
Let \(R=[s,t]\) be a completed normalized shell-crossing record on an
ordinary same-origin branch.  Assume that its saturated joint frontier has
exact support \(A_g=\{\Geom\}\), and let
\[
 Q_{A_g}(R):=
 \sum_{\ell:\operatorname{Prov}(v_{R,\ell})=A_g}
 \widehat\Delta_{\Gamma_R,\ell}(y_{\Phi,R})
\]
be its canonical conditional increment.  Then
\begin{equation}
 \label{eq:geom-canonical-amplitude-identity}
 Q_{A_g}(R)=q_g(R)=a_g(R)=1.
\end{equation}
If the common-form-domain graph is nonordinary, the equality is not asserted;
the first nonordinary \(\Geom\)-owned graph value is the output of the record.
\end{theorem}

\begin{proof}
On the ordinary row, apply the public weak chain rule to the stopped shell
observable:
\[
 \Phi_R(u_t)-\Phi_R(u_s)
 =\int_s^t\langle D\Phi_R(u_r),J_r\rangle\,\dd r.
\]
The normalized crossing makes the left side equal to one.  Saturation and
exact-support projection expand the right side according to
\eqref{eq:unnormalized-target-cocycle}.  The hypothesis that the complete
support is \(A_g\) deletes every other source projection, while residual
nullity deletes the target-null row.  Hence the remaining conditional sum is
both \(Q_{A_g}(R)\) and
\[
 -\int_s^t
 \langle K_r^{1/2}D\Phi_R(u_r),K_r^{1/2}\zeta_r\rangle\,\dd r
 =q_g(R).
\]
Thus \(Q_{A_g}(R)=q_g(R)=1\), and taking the positive part gives
\(a_g(R)=1\).  At a nonordinary form value the displayed chain rule pairing
has no ordinary scalar value; totalization therefore returns its first graph
coordinate instead of applying this calculation.
\end{proof}

\begin{corollary}[Pure-geometric cumulative shell identity]
\label{cor:geom-cumulative-shell-identity}
For every pairwise disjoint family \(\mathcal F_N=(R_1,\ldots,R_N)\) satisfying
the hypotheses of
\cref{thm:geom-canonical-amplitude-identification},
\begin{equation}
 \label{eq:geom-cumulative-canonical-identity}
 \sum_{j=1}^NQ_{A_g}(R_j)
 =\mathsf A_g(\mathcal F_N)=N.
\end{equation}
Consequently \eqref{eq:geom-soft-cascade-ceiling} excludes an infinite
ordinary exact-pure-\(\Geom\) normalized cascade whenever its two displayed
public bounds hold.
\end{corollary}

\begin{proof}
Sum \eqref{eq:geom-canonical-amplitude-identity} before any normalization and
apply \eqref{eq:geom-soft-cascade-ceiling}.
\end{proof}

\begin{theorem}[Exact pure-causal target-current accountability]
\label{thm:quantitative-source-cell-c}
Let \(R=[s,t]\in\mathfrak D_{A_c,\Sigma,n}^{\rm pre}(E)\) have exact
source word \(A_c=\{\Caus\}\), normalized target cell \(C_R\), and public
shell covector \(\lambda_R\).  Evaluate the totalized gluing, coboundary,
storage-preimage, and return graphs in the order of
\eqref{eq:target-caus-normal-form}.  A first nonordinary value is returned
with its \(\Caus\) ancestry.  On the ordinary row define
\begin{align}
 q_c(R)&:=\langle\!\langle\lambda_R,F_{\Caus}(R)\rangle\!\rangle_R,
 &a_c(R)&:=(q_c(R))_+,\label{eq:caus-current-amplitude}\\
 \Lambda_c(R^\pm)&:=
   \langle\lambda_{R^\pm},K_{\Caus}^{\rm stor}(R^\pm)\rangle,
 &p_c(R)&:=p_{\Phys}(R).\label{eq:caus-storage-price-readout}
\end{align}
The storage value is used only when the total coboundary-preimage graph
returns an ordinary potential.  On a nonzero cohomology branch set both
endpoint storage values to zero.  On a zero-class branch with no ordinary
preimage, return the first graph value and stop.  Define
\begin{align}
 f_c(R)&:=q_c(R)-\bigl(\Lambda_c(R^+)-\Lambda_c(R^-)\bigr),
 &c_c(R)&:=(f_c(R))_+,\label{eq:caus-feeding-readout}\\
 \kappa_c(R)&:=\inf\{\kappa\ge0:c_c(R)\le\kappa p_c(R)\}.
 &&\label{eq:caus-carrier-admission}
\end{align}
Set also
\begin{equation}
 \label{eq:caus-incidence-detector}
 \delta_c(R):=\left\|\left(
 H_\Sigma^{\Caus}(C_R),S_\Sigma^{\Caus}(C_R),
 (\Psi_\Sigma^{\Caus}(C_R;\lambda_k))_{1\le k\le m(R)},
 \operatorname{Gc}_\Sigma^{\Caus}(C_R),
 \operatorname{Inc}_{A_c}^{\Sigma}(R)
 \right)\right\|_{\rm det},
\end{equation}
where \(\lambda_1,\ldots,\lambda_{m(R)}\) are exactly the finitely many
rational resolvent probes carried by \(R\).
The equality \(\delta_c(R)=0\) says exactly that all displayed coordinates
of this finite record vanish.  Pure-\(\Caus\) structure is absent on the
whole cell precisely when these equalities hold throughout the cofinal
rational probe hierarchy.  A positive displayed value returns \(R\) as a
finite occurrence with its complete joint frontier.
Here \(\langle\!\langle\cdot,\cdot\rangle\!\rangle_R\) is the signed
stopped pairing generated by the public weak presentation.  By definition,
\begin{equation}
 \label{eq:caus-exact-residual-identity}
 q_c(R)=f_c(R)+\Lambda_c(R^+)-\Lambda_c(R^-).
\end{equation}
Since \(f_c\le(f_c)_+\),
\begin{equation}
 \label{eq:caus-local-conversion}
 a_c(R)\le c_c(R)
 +\bigl(\Lambda_c(R^+)-\Lambda_c(R^-)\bigr)_+.
\end{equation}
If \(\kappa_c(R)<+\infty\), the generated carrier row gives
\begin{equation}
 \label{eq:caus-carrier-conversion}
 a_c(R)\le\kappa_c(R)p_c(R)
 +\bigl(\Lambda_c(R^+)-\Lambda_c(R^-)\bigr)_+.
\end{equation}
If \(\kappa_c(R)=+\infty\), its failed domination row is retained as a
\(\Caus\)-owned carrier-polar coordinate.  Explicitly,
\(p_c=0=c_c\) gives \(\kappa_c=0\), while
\(p_c=0<c_c\) gives \(\kappa_c=+\infty\).

For an ordered stopped complex
\(\mathcal F_N=(R_1,\ldots,R_N)\) whose intervening gaps are retained and
whose adjacent cuts agree, assume that every retained current is
target-oriented, \(q_c(R_j)\ge0\), and set
\[
 \mathsf A_c(\mathcal F_N):=\sum_{j\le N}a_c(R_j),\quad
 \mathsf P_c(\mathcal F_N):=\sum_{j\le N}p_c(R_j),\quad
 \kappa_c^*(\mathcal F_N):=\max_{j\le N}\kappa_c(R_j).
\]
Then
\begin{equation}
 \label{eq:caus-cumulative-conversion}
 \mathsf A_c(\mathcal F_N)
 \le \kappa_c^*(\mathcal F_N)\mathsf P_c(\mathcal F_N)
 +\Lambda_c(R_N^+)-\Lambda_c(R_1^-).
\end{equation}
Moreover \(\mathsf P_c(\mathcal F_N)\le E\).

If every \(R_j\) is a completed normalized crossing with exact support
\(A_c\), then its canonical conditional amplitude satisfies
\begin{equation}
 \label{eq:caus-canonical-amplitude-identity}
 Q_{A_c}(R_j)=q_c(R_j)=a_c(R_j)=1,
 \qquad \mathsf A_c(\mathcal F_N)=N.
\end{equation}
Consequently an infinite pure-causal cascade with finite \(E\) forces
\begin{equation}
 \label{eq:caus-two-way-obstruction}
 \sup_N\kappa_c^*(\mathcal F_N)=+\infty
 \quad\text{or}\quad
 \sup_N\bigl(\Lambda_c(R_N^+)-\Lambda_c(R_1^-)\bigr)=+\infty.
\end{equation}
In the first case, either one finite record has infinite admission value or
the record which first makes the running maximum cross an integer level is a
finite carrier-amplification occurrence.  In the second case, every first
integer-crossing prefix is a finite storage occurrence.  Storage crossings
re-enter through the return-support graph and are evaluated by
\(\delta_\Sigma^{\rm cyc}\), giving positive recurrent physical charge or
the next zero-price return support with the same exact word \(A_c\).  A
singularity verdict is issued exactly when
the resulting complete record also lies in
\(\operatorname{Ran}(\mathfrak R_{\rm ord})\) and has the noncontinuation
value of the public continuation graph.
\end{theorem}

\begin{proof}
Equation \eqref{eq:caus-exact-residual-identity} is the definition of the
feeding residual after subtracting the ordinary storage endpoint difference.
Since \(f_c\le(f_c)_+=c_c\), taking positive parts proves
\eqref{eq:caus-local-conversion}.  The definition of \(\kappa_c\) gives
\(c_c\le\kappa_cp_c\) at every finite value and hence
\eqref{eq:caus-carrier-conversion}.  The zero-carrier alternatives follow
directly from the infimum.  The assertion about the complete probe hierarchy
is \eqref{eq:caus-absence-criterion} together with vanishing of every
rational compressed resolvent readout.  The conditional increment in
\eqref{eq:caus-incidence-detector} enforces exact support \(A_c\).

Sum the exact identities \eqref{eq:caus-exact-residual-identity} over the
stopped complex.  Target orientation gives
\(\mathsf A_c=\sum_jq_c(R_j)\), while agreement of adjacent cuts cancels
every internal storage value.  Bounding each \(f_c\) by
\(c_c\le\kappa_c^*p_c\) gives
\eqref{eq:caus-cumulative-conversion}.  The price bound follows from
restriction of the single nonnegative physical carrier to pairwise disjoint
subrecords of a stratum of mass at most \(E\).

For a normalized crossing, the public stopped chain rule gives unit shell
increment.  Exact-support projection leaves only the \(A_c\)-row of
\eqref{eq:unnormalized-target-cocycle}; target-null residual closure deletes
the residual row.  This identifies the canonical conditional amplitude with
\(q_c(R_j)=1\), proving \eqref{eq:caus-canonical-amplitude-identity}.
If both quantities in \eqref{eq:caus-two-way-obstruction} were bounded,
\eqref{eq:caus-cumulative-conversion} would bound \(N\), a contradiction.
First-crossing finiteness follows from monotonicity of the running maximum
and from the ordered prefix construction for storage.  The storage and
return-price successors are the second and third coordinates of
\eqref{eq:target-caus-normal-form} and
\cref{thm:continuum-caus-normal-form-exhaustion}.  Ordinary realization and
noncontinuation are the two remaining intersections in
\eqref{eq:actual-contact-cell-realization-intersection}.
\end{proof}

\begin{theorem}[Exact joint \(\Geom\Caus\) target-current accountability]
\label{thm:quantitative-source-cell-gc}
Let \(R=[s,t]\in\mathfrak D_{A_{gc},\Sigma,n}^{\rm pre}(E)\) have exact
word \(A_{gc}=\{\Geom,\Caus\}\), normalized cell \(C_R\), and shell
covector \(\lambda_{R,r}=D\Phi_R(u_r)\).  Form the complete joint record
before either source projection.  Totalize the geometric form graph and the
causal coboundary-preimage graph.  A first nonordinary value is returned with
the complete word \(A_{gc}\).  On their common ordinary row, let
\begin{align}
 q_{gc}(R)&:=
 \bigl\langle\!\bigl\langle\lambda_R,
       \operatorname{Inc}_{A_{gc}}^\Sigma(R)\bigr\rangle\!\bigr\rangle_R,
 &a_{gc}(R)&:=(q_{gc}(R))_+,\label{eq:gc-current-amplitude}\\
 g_{gc}(R)&:=-\int_s^t
 \langle K_r^{1/2}\lambda_{R,r},K_r^{1/2}\zeta_{gc,r}\rangle\,\dd r,
 &d_{gc}(R)&:=\int_s^t\|K_r^{1/2}\zeta_{gc,r}\|^2\,\dd r,\\
 h_{gc}(R)&:=\int_s^t\|K_r^{1/2}\lambda_{R,r}\|^2\,\dd r,
 &p_{gc}(R)&:=p_{\Phys}(R).\label{eq:gc-form-costs}
\end{align}
Here \(\zeta_{gc}\) is the geometric-headed coordinate after projection to
the complete word \(A_{gc}\), not the pure-\(\Geom\) coordinate.
If the exact-word causal class vanishes and the total preimage graph returns
an ordinary potential, let \(\Lambda_{gc}(R^\pm)\) be its endpoint target
readout.  If the class is nonzero, set both endpoint values to zero and leave
the class in the residual below.  If the class vanishes but the preimage
graph is nonordinary, return its first graph value and stop this branch.
Define
\begin{align}
 w_{gc}(R)&:=q_{gc}(R)-g_{gc}(R)
  -\bigl(\Lambda_{gc}(R^+)-\Lambda_{gc}(R^-)\bigr),
 &b_{gc}(R)&:=(w_{gc}(R))_+,\\
 \kappa_{gc}(R)&:=\inf\{\kappa\ge0:b_{gc}(R)\le\kappa p_{gc}(R)\},
 &\rho_{gc}(R)&:=\inf\{\rho\ge0:d_{gc}(R)\le\rho p_{gc}(R)\}.
 \label{eq:gc-admission-values}
\end{align}
Thus \(p_{gc}=0=b_{gc}\) gives \(\kappa_{gc}=0\), whereas
\(p_{gc}=0<b_{gc}\) gives \(\kappa_{gc}=+\infty\); the identical alternatives
with \(d_{gc}\) give \(\rho_{gc}=0\) or \(+\infty\).  Every infinite value
is returned as its finite failed-domination coordinate.
Set
\begin{equation}
 \label{eq:gc-incidence-detector}
 \delta_{gc}(R):=\left\|\left(
 G_\Sigma^{\Geom}(C_R),\operatorname{Ker}_\Sigma^{\Geom}(C_R),
 \operatorname{Gc}_\Sigma^{\Geom}(C_R),
 H_\Sigma^{\Caus}(C_R),S_\Sigma^{\Caus}(C_R),
 (\Psi_\Sigma^{\Caus}(C_R;\lambda_k))_{1\le k\le m(R)},
 \operatorname{Gc}_\Sigma^{\Caus}(C_R),
 \operatorname{Inc}_{A_{gc}}^\Sigma(R)\right)\right\|_{\rm det}.
\end{equation}
Its vanishing is the simultaneous zero test for the displayed finite joint
record.  Vanishing throughout the cofinal probe hierarchy is equivalent to
absence of exact-\(A_{gc}\) incidence by the two normal-form exhaustion
theorems and joint projection before source separation.
All displayed record quantities are evaluated before normalization.  For every
\(\theta>0\),
\begin{equation}
 \label{eq:gc-local-conversion}
 a_{gc}(R)\le
 \frac\theta2d_{gc}(R)+\frac1{2\theta}h_{gc}(R)+b_{gc}(R)
 +\bigl(\Lambda_{gc}(R^+)-\Lambda_{gc}(R^-)\bigr)_+.
\end{equation}

Let \(\mathcal F_N=(R_1,\ldots,R_N)\) be a target-oriented ordered stopped
complex of pairwise disjoint records with matching adjacent cuts.  Define
\begin{align*}
 \mathsf A_{gc}(\mathcal F_N)&:=\sum_{j\le N}a_{gc}(R_j),&
 \mathsf D_{gc}(\mathcal F_N)&:=\sum_{j\le N}d_{gc}(R_j),\\
 \mathsf H_{gc}(\mathcal F_N)&:=\sum_{j\le N}h_{gc}(R_j),&
 \mathsf P_{gc}(\mathcal F_N)&:=\sum_{j\le N}p_{gc}(R_j),\\
 \kappa_{gc}^*(\mathcal F_N)&:=\max_{j\le N}\kappa_{gc}(R_j),&
 \rho_{gc}^*(\mathcal F_N)&:=\max_{j\le N}\rho_{gc}(R_j).
\end{align*}
If \(\kappa_{gc}^*(\mathcal F_N)\) and
\(\rho_{gc}^*(\mathcal F_N)\) are finite, then
\begin{equation}
 \label{eq:gc-cumulative-conversion}
 \mathsf A_{gc}(\mathcal F_N)
 \le\left(\frac\theta2\rho_{gc}^*(\mathcal F_N)
              +\kappa_{gc}^*(\mathcal F_N)\right)
       \mathsf P_{gc}(\mathcal F_N)
 +\frac1{2\theta}\mathsf H_{gc}(\mathcal F_N)
 +\Lambda_{gc}(R_N^+)-\Lambda_{gc}(R_1^-),
\end{equation}
and \(\mathsf P_{gc}(\mathcal F_N)\le E\).

For completed normalized crossings with exact support \(A_{gc}\),
\begin{equation}
 \label{eq:gc-canonical-amplitude-identity}
 Q_{A_{gc}}(R_j)=q_{gc}(R_j)=a_{gc}(R_j)=1,
 \qquad \mathsf A_{gc}(\mathcal F_N)=N.
\end{equation}
Hence an infinite exact-\(A_{gc}\) cascade forces at least one of
\(\rho_{gc}^*\), \(\kappa_{gc}^*\), \(\mathsf H_{gc}\), or the terminal
storage difference in \eqref{eq:gc-cumulative-conversion} to be unbounded.
Every alternative has a finite first-crossing joint record with exact word
\(A_{gc}\).  Geometric carrier-domination failure, joint-residual
carrier-polar failure, causal storage re-entry, and zero-price return re-entry
retain that complete word.  An orbit
or singularity conclusion requires the ordinary realization and continuation
coordinates of the resulting joint record.
\end{theorem}
\begin{proof}
The definition of \(w_{gc}\) gives the exact identity
\[
 q_{gc}=g_{gc}+w_{gc}+\Lambda_{gc}(R^+)-\Lambda_{gc}(R^-).
\]
Taking positive parts, applying Cauchy--Schwarz and Young to \(g_{gc}\), and
using \(w_{gc}\le b_{gc}\) proves \eqref{eq:gc-local-conversion}.  On a
target-oriented stopped complex, sum the exact identities before taking
positive parts.  Internal storage values cancel.  At finite running maxima,
the definitions of \(\rho_{gc}\) and \(\kappa_{gc}\), followed by the
single-carrier bound, give \eqref{eq:gc-cumulative-conversion}.  An infinite
admission value is already the failed-domination branch stated in the
theorem, so no expression of the form \(+\infty\cdot0\) is used.

For a normalized crossing, the stopped shell chain rule has value one.
Projection of the saturated joint frontier to exact word \(A_{gc}\) leaves
exactly \(Q_{A_{gc}}\); the target-null residual has zero readout.  This proves
\eqref{eq:gc-canonical-amplitude-identity}.  If all four cumulative
coordinates were bounded, the right side of
\eqref{eq:gc-cumulative-conversion} would be independent of \(N\), contrary
to its left side.  Running maxima, nonnegative cumulative probe cost, and
ordered storage prefixes have finite first crossings.  Totalization and
source-preserving re-entry retain \(A_{gc}\).  Ordinary realization and
continuation are evaluated only through
\eqref{eq:actual-contact-cell-realization-intersection}.
\end{proof}

\begin{theorem}[Exact pure-quotient target-current accountability]
\label{thm:quantitative-source-cell-a}
Let \(R\in\mathfrak D_{A_a,\Sigma,n}^{\rm pre}(E)\) have exact word
\(A_a=\{\Abs\}\), normalized cell \(C_R\), and public quotient data
\((q,U,P,\Pi_q)\).  Evaluate every operation in
\eqref{eq:target-abs-normal-form} by its total graph.  A nonordinary quotient,
right-inverse, lifting, product, or projection value returns the first
\(\Abs\)-owned graph coordinate and stops that branch.  A nonzero generated
lifting/surjectivity cokernel value is likewise returned as its finite
\(\Abs\)-owned cokernel branch.  The scalar formulas below apply exactly to
the ordinary, cokernel-zero identity row.

On the common ordinary row, project the quotient identity to \(A_a\) and set
\begin{align}
 q_a(R)&:=\bigl\langle\!\bigl\langle\lambda_R,
       \operatorname{Inc}_{A_a}^{\Sigma}(R)\bigr\rangle\!\bigr\rangle_R,
 &a_a(R)&:=(q_a(R))_+,\label{eq:abs-current-amplitude}\\
 v_a(R)&:=\bigl\langle\!\bigl\langle\lambda_R,
   E_qP+J(I-\Pi_q)\bigr\rangle\!\bigr\rangle_{R,A_a},
 &p_a(R)&:=p_{\Phys}(R).\label{eq:abs-defect-current-price}
\end{align}
Complete provenance is fixed before the last pairing.  The horizontal term
\(UJ_QP\) transfers to its inherited source word and has zero exact-\(A_a\)
projection.  At the current level the exact statement is
\begin{equation}
 \label{eq:abs-exact-source-projection}
 \operatorname{pr}_{A_a}\operatorname{Inc}^{\Sigma}(R)
 =
 \operatorname{pr}_{A_a}\bigl(E_qP+J(I-\Pi_q)\bigr).
\end{equation}
Pairing this identity with the stopped shell covector gives
\begin{equation}
 \label{eq:abs-exact-current-identity}
 q_a(R)=v_a(R).
\end{equation}
Define
\begin{equation}
 \label{eq:abs-carrier-admission}
 \kappa_a(R):=\inf\{\kappa\ge0:a_a(R)\le\kappa p_a(R)\}.
\end{equation}
Thus \(p_a=0=a_a\) gives \(\kappa_a=0\), while \(p_a=0<a_a\) gives
\(\kappa_a=+\infty\).  At every finite value,
\begin{equation}
 \label{eq:abs-local-conversion}
 a_a(R)\le\kappa_a(R)p_a(R).
\end{equation}

Define
\begin{align}
 \delta_a(R)&:=\left\|\left(
 A_\Sigma^{\rm desc}(C_R),A_\Sigma^{\rm vert}(C_R),
 A_\Sigma^{\rm coker}(C_R),A_\Sigma^{\rm gc}(C_R),
 A_\Sigma^{\rm tgt}(C_R),\operatorname{Inc}_{A_a}^{\Sigma}(R)
 \right)\right\|_{\rm det},\label{eq:abs-incidence-detector}\\
 \chi_a(R)&:=\left(\operatorname{Cap}_\Sigma^{\ker q}(C_R),
 \operatorname{Rank}_\Sigma(q;C_R)\right).
\end{align}
The equality \(\delta_a(R)=0\) is the zero test for this finite record.
Vanishing throughout the cofinal incidence hierarchy is equivalent to
pure-\(\Abs\) absence by \eqref{eq:abs-absence-criterion}.  The valuation
\(\chi_a\) is retained in both rows and is not an absence detector.

For a target-oriented family
\(\mathcal F_N^{\rm ord}=(R_1,\ldots,R_N)\) of pairwise disjoint same-origin
records on the ordinary cokernel-zero row,
set
\[
 \mathsf A_a(\mathcal F_N^{\rm ord}):=\sum_{j\le N}a_a(R_j),\qquad
 \mathsf P_a(\mathcal F_N^{\rm ord}):=\sum_{j\le N}p_a(R_j),\qquad
 \kappa_a^*(\mathcal F_N^{\rm ord}):=\max_{j\le N}\kappa_a(R_j).
\]
When \(\kappa_a^*(\mathcal F_N^{\rm ord})<+\infty\),
\begin{equation}
 \label{eq:abs-cumulative-conversion}
 \mathsf A_a(\mathcal F_N^{\rm ord})
 \le\kappa_a^*(\mathcal F_N^{\rm ord})
       \mathsf P_a(\mathcal F_N^{\rm ord})
 \le\kappa_a^*(\mathcal F_N^{\rm ord})E.
\end{equation}
For completed normalized crossings on that ordinary row with exact support
\(A_a\),
\begin{equation}
 \label{eq:abs-canonical-amplitude-identity}
 Q_{A_a}(R_j)=q_a(R_j)=a_a(R_j)=1,
 \qquad \mathsf A_a(\mathcal F_N^{\rm ord})=N.
\end{equation}
For an arbitrary ordered family \(\mathcal G_N\) of completed normalized
exact-\(A_a\) crossings, let
\[
 \mathsf B_a(\mathcal G_N)
 :=\#\{j\le N:R_j\text{ exits through the cokernel or graph branch}\}.
\]
The remaining records form \(\mathcal G_N^{\rm ord}\), and the exact branch
partition gives, with \(\kappa_a^*(\varnothing):=0\),
\begin{equation}
 \label{eq:abs-ordinary-graph-partition}
 N=\mathsf B_a(\mathcal G_N)
   +\mathsf A_a(\mathcal G_N^{\rm ord}).
\end{equation}
Hence an infinite exact-\(A_a\) cascade in a finite physical stratum forces
\(\sup_N\mathsf B_a(\mathcal G_N)=+\infty\) or
\(\sup_N\kappa_a^*(\mathcal G_N^{\rm ord})=+\infty\).  In the first case,
each integer crossing returns a finite cokernel/graph prefix.  In the second,
either a finite ordinary record already has infinite admission value, or the
record which first raises the running maximum above each integer is a finite
quotient-amplification occurrence.
Zero-price equality saturation, cokernel successors, and graph successors
retain exact word \(A_a\).  A singularity verdict requires ordinary
same-origin realization of the complete record and the noncontinuation value
of the public continuation graph.
\end{theorem}
\begin{proof}
Apply exact-word projection to
\[
 J=UJ_QP+E_qP+J(I-\Pi_q).
\]
By \cref{thm:continuum-abs-normal-form-exhaustion}, the horizontal term
transfers to its inherited source and contributes no \(\Abs\)-headed defect.
Exact source projection therefore proves
\eqref{eq:abs-exact-source-projection}; pairing gives
\eqref{eq:abs-exact-current-identity}.  The admission definition proves
\eqref{eq:abs-local-conversion}, including the two zero-carrier cases.
Summing over the target-oriented family and using finite additivity of the
single physical carrier gives \eqref{eq:abs-cumulative-conversion}.

For a completed normalized crossing, the public stopped chain rule is one.
Exact-support projection of the saturated joint frontier leaves only
\(Q_{A_a}\), and target-null residual closure removes the residual row.
This proves \eqref{eq:abs-canonical-amplitude-identity}.  The ordered total
graph gives a disjoint partition into the ordinary, cokernel, and graph rows,
which proves \eqref{eq:abs-ordinary-graph-partition}.  If the nonordinary
count is bounded, the number of ordinary crossings is unbounded; a bounded
ordinary admission maximum would then contradict
\eqref{eq:abs-cumulative-conversion}.  This proves the two alternatives.
The detector equivalence is \eqref{eq:abs-absence-criterion}, and totalization
supplies the cokernel and graph successors.  Ordinary realization and
noncontinuation are evaluated only through
\eqref{eq:actual-contact-cell-realization-intersection}.
\end{proof}

\begin{theorem}[Exact joint \(\Geom\Abs\) target-current accountability]
\label{thm:quantitative-source-cell-ga}
Let \(R=[s,t]\in\mathfrak D_{A_{ga},\Sigma,n}^{\rm pre}(E)\) have exact
word \(A_{ga}=\{\Geom,\Abs\}\), normalized cell \(C_R\), and covector
\(\lambda_{R,r}=D\Phi_R(u_r)\).  Form the complete joint record before
source projection.  A nonordinary geometric or quotient operation, or a
nonzero quotient cokernel value, returns its first finite branch with exact
word \(A_{ga}\).  The formulas below apply to the ordinary cokernel-zero row.

Let \(\zeta_{ga}\) be its exact-word geometric-headed coordinate and put
\begin{align}
 q_{ga}(R)&:=\bigl\langle\!\bigl\langle\lambda_R,
 \operatorname{Inc}_{A_{ga}}^\Sigma(R)\bigr\rangle\!\bigr\rangle_R,
 &a_{ga}(R)&:=(q_{ga}(R))_+,\\
 g_{ga}(R)&:=-\int_s^t\langle K_r^{1/2}\lambda_{R,r},
                         K_r^{1/2}\zeta_{ga,r}\rangle\,\dd r,
 &d_{ga}(R)&:=\int_s^t\|K_r^{1/2}\zeta_{ga,r}\|^2\,\dd r,\\
 h_{ga}(R)&:=\int_s^t\|K_r^{1/2}\lambda_{R,r}\|^2\,\dd r,
 &p_{ga}(R)&:=p_{\Phys}(R),\\
 v_{ga}(R)&:=\bigl\langle\!\bigl\langle\lambda_R,
   E_qP+J(I-\Pi_q)\bigr\rangle\!\bigr\rangle_{R,A_{ga}},
 &c_{ga}(R)&:=(v_{ga}(R))_+.
 \label{eq:ga-record-currents}
\end{align}
Horizontal descent has zero exact-\(A_{ga}\) projection.  Define the
remaining mixed interaction by
\[
 w_{ga}(R):=q_{ga}(R)-g_{ga}(R)-v_{ga}(R),
 \qquad b_{ga}(R):=(w_{ga}(R))_+.
\]
Its three physical-admission values are
\begin{align}
 \rho_{ga}(R)&:=\inf\{\rho\ge0:d_{ga}(R)\le\rho p_{ga}(R)\},&
 \sigma_{ga}(R)&:=\inf\{\sigma\ge0:c_{ga}(R)\le\sigma p_{ga}(R)\},\\
 \kappa_{ga}(R)&:=\inf\{\kappa\ge0:b_{ga}(R)\le\kappa p_{ga}(R)\}.
 &&\label{eq:ga-admission-values}
\end{align}
An infinite value is a finite failed-domination coordinate; zero carrier and
zero numerator give admission value zero.  For every \(\theta>0\),
\begin{equation}
 \label{eq:ga-local-conversion}
 a_{ga}(R)\le\frac\theta2d_{ga}(R)+\frac1{2\theta}h_{ga}(R)
 +c_{ga}(R)+b_{ga}(R).
\end{equation}
The finite joint detector is
\begin{equation}
 \label{eq:ga-incidence-detector}
 \delta_{ga}(R):=\left\|\left(
 G_\Sigma^{\Geom}(C_R),\operatorname{Ker}_\Sigma^{\Geom}(C_R),
 \operatorname{Gc}_\Sigma^{\Geom}(C_R),
 A_\Sigma^{\rm desc}(C_R),A_\Sigma^{\rm vert}(C_R),
 A_\Sigma^{\rm coker}(C_R),A_\Sigma^{\rm gc}(C_R),
 A_\Sigma^{\rm tgt}(C_R),\operatorname{Inc}_{A_{ga}}^\Sigma(R)
 \right)\right\|_{\rm det}.
\end{equation}
Its zero value tests every displayed coordinate of \(R\).  Vanishing on the
cofinal joint hierarchy is equivalent to absence of exact-\(A_{ga}\)
incidence.  Kernel capacity and quotient rank remain attached as valuations.

Let \(\mathcal F_N^{\rm ord}\) be a target-oriented family of pairwise
disjoint ordinary records.  Define
\begin{align*}
 \mathsf A_{ga}(\mathcal F_N^{\rm ord})&:=\sum_{j\le N}a_{ga}(R_j),&
 \mathsf H_{ga}(\mathcal F_N^{\rm ord})&:=\sum_{j\le N}h_{ga}(R_j),\\
 \mathsf P_{ga}(\mathcal F_N^{\rm ord})&:=\sum_{j\le N}p_{ga}(R_j),&
 \rho_{ga}^*(\mathcal F_N^{\rm ord})&:=\max_{j\le N}\rho_{ga}(R_j),\\
 \sigma_{ga}^*(\mathcal F_N^{\rm ord})&:=\max_{j\le N}\sigma_{ga}(R_j),&
 \kappa_{ga}^*(\mathcal F_N^{\rm ord})&:=\max_{j\le N}\kappa_{ga}(R_j).
\end{align*}
When all three maxima are finite,
\begin{equation}
 \label{eq:ga-cumulative-conversion}
 \mathsf A_{ga}(\mathcal F_N^{\rm ord})
 \le\left(\frac\theta2\rho_{ga}^*(\mathcal F_N^{\rm ord})
 +\sigma_{ga}^*(\mathcal F_N^{\rm ord})
 +\kappa_{ga}^*(\mathcal F_N^{\rm ord})\right)E
 +\frac1{2\theta}\mathsf H_{ga}(\mathcal F_N^{\rm ord}).
\end{equation}
Every ordinary completed normalized exact-\(A_{ga}\) crossing satisfies
\begin{equation}
 \label{eq:ga-canonical-amplitude}
 Q_{A_{ga}}(R)=q_{ga}(R)=a_{ga}(R)=1.
\end{equation}

For an arbitrary first-\(N\) family \(\mathcal G_N\), let
\(\mathsf B_{ga}(\mathcal G_N)\) count its geometric-graph,
quotient-cokernel, and quotient-graph exits.  Its ordinary subfamily obeys
\begin{equation}
 \label{eq:ga-branch-partition}
 N=\mathsf B_{ga}(\mathcal G_N)
 +\mathsf A_{ga}(\mathcal G_N^{\rm ord}).
\end{equation}
Running maxima on an empty ordinary subfamily are defined to be zero.
Thus an infinite exact-\(A_{ga}\) cascade forces unbounded
\(\mathsf B_{ga}\), \(\rho_{ga}^*\), \(\sigma_{ga}^*\),
\(\kappa_{ga}^*\), or \(\mathsf H_{ga}\).  Each alternative has finite first
crossings retaining \(A_{ga}\).  An orbit conclusion additionally requires
ordinary same-origin realization and public noncontinuation.
\end{theorem}
\begin{proof}
Exact-word projection of the geometric form row and
\eqref{eq:target-abs-normal-form} removes horizontal descent and leaves the
geometric current, quotient defect/vertical current, and their remaining
joint interaction.  By definition,
\[
 q_{ga}=g_{ga}+v_{ga}+w_{ga}.
\]
Taking positive parts, applying Cauchy--Schwarz and Young to \(g_{ga}\), and
using \(v_{ga}\le c_{ga}\) and \(w_{ga}\le b_{ga}\) proves
\eqref{eq:ga-local-conversion}.  Sum this exact identity over the
target-oriented ordinary family.  The three finite admission inequalities,
finite additivity of the physical carrier, and its total mass bound give
\eqref{eq:ga-cumulative-conversion}.

The normalized stopped shell identity, exact-support projection, and
target-null residual closure prove \eqref{eq:ga-canonical-amplitude}.  The
ordered total graphs partition every crossing into the ordinary row or one
of the counted early exits, proving \eqref{eq:ga-branch-partition}.  If the
exit count and all four ordinary cumulative coordinates were bounded,
\eqref{eq:ga-cumulative-conversion} would bound \(N\), a contradiction.
Running maxima and cumulative probe cost have finite first crossings.
Source-preserving totalization retains \(A_{ga}\); realization and
continuation are evaluated through
\eqref{eq:actual-contact-cell-realization-intersection}.
\end{proof}

\begin{theorem}[Exact joint \(\Caus\Abs\) target-current accountability]
\label{thm:quantitative-source-cell-ca}
Let \(R=[s,t]\in\mathfrak D_{A_{ca},\Sigma,n}^{\rm pre}(E)\) have exact
word \(A_{ca}=\{\Caus,\Abs\}\), normalized cell \(C_R\), public quotient
data \((q,U,P,\Pi_q)\), and stopped shell covector \(\lambda_R\).  Form the
complete joint record before either source projection.  Evaluate every causal
gluing, coboundary-preimage, support, resolvent, and return operation, and
every quotient operation, in
\eqref{eq:target-caus-normal-form} and \eqref{eq:target-abs-normal-form} by
their total graphs.  A first nonordinary value, or a nonzero generated
quotient cokernel, returns its finite branch with exact word \(A_{ca}\).
The scalar identities below apply to the common ordinary, quotient-cokernel
zero row.

On this row define
\begin{align}
 q_{ca}(R)&:=\bigl\langle\!\bigl\langle\lambda_R,
   \operatorname{Inc}_{A_{ca}}^\Sigma(R)\bigr\rangle\!\bigr\rangle_R,
 &a_{ca}(R)&:=(q_{ca}(R))_+,\label{eq:ca-current-amplitude}\\
 v_{ca}(R)&:=\bigl\langle\!\bigl\langle\lambda_R,
   E_qP+J(I-\Pi_q)\bigr\rangle\!\bigr\rangle_{R,A_{ca}},
 &c_{ca}(R)&:=(v_{ca}(R))_+,\label{eq:ca-quotient-current}\\
 p_{ca}(R)&:=p_{\Phys}(R).&&
\end{align}
The source-pure horizontal quotient term \(UJ_QP\) transfers to its inherited
word.  Any horizontal term whose completed provenance is the joint word
\(A_{ca}\) is retained in \(w_{ca}\) below.  If the exact-word causal
class vanishes and its total preimage graph returns an ordinary potential,
let \(\Lambda_{ca}(R^\pm)\) be the two endpoint target readouts.  If the
class is nonzero, set both endpoint readouts to zero and retain that class in
the residual below.  If the class vanishes but the preimage is nonordinary,
return its first graph value and stop.  Put
\begin{equation}
 \label{eq:ca-exact-word-current-decomposition}
 \operatorname{pr}_{A_{ca}}\operatorname{Inc}^{\Sigma}(R)
 =\operatorname{pr}_{A_{ca}}\bigl(E_qP+J(I-\Pi_q)\bigr)
  +\Delta_{\rm stor}K_{ca}(R)+\mathcal E_{ca}^{=}(R).
\end{equation}
Here \(\Delta_{\rm stor}K_{ca}(R)\) is the generated stopped coboundary
current whose pairing is
\(\Lambda_{ca}(R^+)-\Lambda_{ca}(R^-)\), and
\(\mathcal E_{ca}^{=}(R)\) is the exact remainder obtained by subtraction
in the expanded joint derivation.  It retains word \(A_{ca}\), including
every joint horizontal term and every nonzero causal class.  Thus it is a
public finite-record current.  Set
\begin{align}
 w_{ca}(R)&:=\bigl\langle\!\bigl\langle\lambda_R,
                  \mathcal E_{ca}^{=}(R)\bigr\rangle\!\bigr\rangle_R,
 &b_{ca}(R)&:=(w_{ca}(R))_+,\label{eq:ca-exact-residual}\\
 \sigma_{ca}(R)&:=\inf\{\sigma\ge0:c_{ca}(R)\le\sigma p_{ca}(R)\},
 &\kappa_{ca}(R)&:=\inf\{\kappa\ge0:b_{ca}(R)\le\kappa p_{ca}(R)\}.
 \label{eq:ca-admission-values}
\end{align}
Thus the exact signed source conversion is
\begin{equation}
 \label{eq:ca-exact-current-identity}
 q_{ca}(R)=v_{ca}(R)+w_{ca}(R)
 +\Lambda_{ca}(R^+)-\Lambda_{ca}(R^-),
\end{equation}
and hence
\begin{equation}
 \label{eq:ca-local-conversion}
 a_{ca}(R)\le c_{ca}(R)+b_{ca}(R)
 +\bigl(\Lambda_{ca}(R^+)-\Lambda_{ca}(R^-)\bigr)_+.
\end{equation}
For either numerator \(x\in\{c_{ca},b_{ca}\}\), zero carrier and
\(x=0\) give admission zero, while zero carrier and \(x>0\) give admission
\(+\infty\).  Every infinite admission is retained as its finite
failed-domination coordinate; multiplication by it is never performed.

The finite joint detector is
\begin{equation}
 \label{eq:ca-incidence-detector}
 \delta_{ca}(R):=\left\|\left(
 H_\Sigma^{\Caus}(C_R),S_\Sigma^{\Caus}(C_R),
 (\Psi_\Sigma^{\Caus}(C_R;\lambda_k))_{1\le k\le m(R)},
 \operatorname{Gc}_\Sigma^{\Caus}(C_R),
 A_\Sigma^{\rm desc}(C_R),A_\Sigma^{\rm vert}(C_R),
 A_\Sigma^{\rm coker}(C_R),A_\Sigma^{\rm gc}(C_R),
 A_\Sigma^{\rm tgt}(C_R),
 \operatorname{Inc}_{A_{ca}}^\Sigma(R)
 \right)\right\|_{\rm det}.
\end{equation}
Here \(\lambda_1,\ldots,\lambda_{m(R)}\) are exactly the rational probes
stored in \(R\).  The equality \(\delta_{ca}(R)=0\) is the simultaneous
zero test for this finite record.  Vanishing throughout the cofinal joint
probe hierarchy is equivalent to absence of exact-\(A_{ca}\) incidence.
Quotient rank remains attached as a valuation and is not used as an absence
detector.

Let \(\mathcal F_N^{\rm ord}=(R_1,\ldots,R_N)\) be a target-oriented
ordered stopped complex of pairwise disjoint ordinary records whose adjacent
cuts agree and whose intervening gaps are retained.  Set
\begin{align*}
 \mathsf A_{ca}(\mathcal F_N^{\rm ord})&:=\sum_{j\le N}a_{ca}(R_j),&
 \mathsf P_{ca}(\mathcal F_N^{\rm ord})&:=\sum_{j\le N}p_{ca}(R_j),\\
 \sigma_{ca}^*(\mathcal F_N^{\rm ord})&:=\max_{j\le N}\sigma_{ca}(R_j),&
 \kappa_{ca}^*(\mathcal F_N^{\rm ord})&:=\max_{j\le N}\kappa_{ca}(R_j).
\end{align*}
When both maxima are finite, exact telescoping gives
\begin{equation}
 \label{eq:ca-cumulative-conversion}
 \mathsf A_{ca}(\mathcal F_N^{\rm ord})
 \le\bigl(\sigma_{ca}^*(\mathcal F_N^{\rm ord})
          +\kappa_{ca}^*(\mathcal F_N^{\rm ord})\bigr)
       \mathsf P_{ca}(\mathcal F_N^{\rm ord})
 +\Lambda_{ca}(R_N^+)-\Lambda_{ca}(R_1^-),
 \qquad \mathsf P_{ca}(\mathcal F_N^{\rm ord})\le E.
\end{equation}
Every completed normalized ordinary crossing with exact support \(A_{ca}\)
satisfies
\begin{equation}
 \label{eq:ca-canonical-amplitude}
 Q_{A_{ca}}(R)=q_{ca}(R)=a_{ca}(R)=1.
\end{equation}

For an arbitrary ordered family \(\mathcal G_N\) of completed normalized
exact-\(A_{ca}\) crossings, let \(\mathsf B_{ca}(\mathcal G_N)\) count the
causal gluing, preimage, support, resolvent, and return graph exits, together
with the quotient-cokernel and quotient-graph exits.  The
remaining records form \(\mathcal G_N^{\rm ord}\), and, with both empty
running maxima defined as zero,
\begin{equation}
 \label{eq:ca-branch-partition}
 N=\mathsf B_{ca}(\mathcal G_N)
  +\mathsf A_{ca}(\mathcal G_N^{\rm ord}).
\end{equation}
Consequently an infinite exact-\(A_{ca}\) cascade in a finite physical
stratum forces at least one of \(\mathsf B_{ca}\), \(\sigma_{ca}^*\),
\(\kappa_{ca}^*\), or
\(\Lambda_{ca}(R_N^+)-\Lambda_{ca}(R_1^-)\) to be unbounded.  Each
alternative has finite first-crossing prefixes retaining exact word
\(A_{ca}\).  Storage prefixes re-enter through the causal return-support
graph and are evaluated by \(\delta_\Sigma^{\rm cyc}\): positive cycle
price yields divergent recurrent physical charge, whereas zero cycle price
returns the next target-aligned support.  Zero-price quotient and residual
currents enter the equality saturation of the same complete joint word.  An
orbit or singularity conclusion additionally
requires ordinary same-origin realization and the noncontinuation value of
the public continuation graph.
\end{theorem}
\begin{proof}
Apply exact-word projection only after joining the causal and quotient
normal forms.  Quotient exhaustion transfers the source-pure part of
\(UJ_QP\) to its inherited word, leaves the projected defect/vertical
current \(v_{ca}\), and retains every joint horizontal interaction for
\(\mathcal E_{ca}^{=}\).  Causal
exhaustion either returns a graph value, supplies an ordinary endpoint
potential, or retains its nonzero class in the joint residual.  The
expanded derivation is closed under subtraction, so these exhaustive rows
give \eqref{eq:ca-exact-word-current-decomposition}.  Pairing that identity
gives \eqref{eq:ca-exact-current-identity}; no component is discarded.  Taking
positive parts proves \eqref{eq:ca-local-conversion}.

On the ordinary row with finite admissions, sum
\eqref{eq:ca-exact-current-identity} before taking positive parts.
Target orientation identifies the left side with \(\mathsf A_{ca}\), and
matching adjacent cuts cancel every internal storage value.  The two
admission inequalities and finite additivity of the single physical carrier
give \eqref{eq:ca-cumulative-conversion}.  The stated zero-carrier cases
follow directly from the two infima.

For a completed normalized crossing, the public stopped chain rule has
value one.  Projection of the saturated joint frontier to exact word
\(A_{ca}\) leaves \(Q_{A_{ca}}\), while target-null residual closure has
zero readout.  This proves \eqref{eq:ca-canonical-amplitude}.  Ordered total
graphs partition every crossing into exactly one early-exit or ordinary row,
which proves \eqref{eq:ca-branch-partition}.  If the exit count, both
admission maxima, and terminal storage difference were bounded, the
integer-valued exit count would stabilize.  After its last increase, the
remaining tail is an ordinary stopped complex; all intervening cuts and gaps
remain in that tail.  Applying \eqref{eq:ca-cumulative-conversion} there and
using the finite carrier mass would bound its crossing count, contradicting
the partition.  Running maxima and ordered storage prefixes have finite
first crossings.  The two normal-form exhaustion theorems, applied after
joint exact-word projection, also show that cofinal vanishing of
\eqref{eq:ca-incidence-detector} is equivalent to exact-\(A_{ca}\) absence.
Source-preserving totalization retains \(A_{ca}\).  The recurrent storage
alternative is the cycle-price conclusion of
\cref{thm:continuum-caus-normal-form-exhaustion}; the two zero-carrier
alternatives are processed by
\cref{thm:actual-source-price-equality-decision}.  Ordinary realization and
noncontinuation are evaluated through
\eqref{eq:actual-contact-cell-realization-intersection}.
\end{proof}

\begin{theorem}[Exact joint \(\Geom\Caus\Abs\) target-current accountability]
\label{thm:quantitative-source-cell-gca}
Let \(R=[s,t]\in\mathfrak D_{A_{gca},\Sigma,n}^{\rm pre}(E)\) have exact
word \(A_{gca}=\{\Geom,\Caus,\Abs\}\), normalized cell \(C_R\), shell
covector \(\lambda_{R,r}=D\Phi_R(u_r)\), and quotient data
\((q,U,P,\Pi_q)\).  Form the complete joint record before source
projection.  Totalize the geometric form, range, kernel/cokernel, and graph
rows, every causal gluing,
coboundary-preimage, support, resolvent, and return graph, and every quotient
operation.  A first nonordinary value, a first-active target-visible
geometric kernel/cokernel value, or a nonzero quotient cokernel returns a
finite branch with exact word \(A_{gca}\).  The formulas below apply to the
common ordinary-range, quotient-cokernel-zero row.

Let \(\zeta_{gca}\) be the geometric-headed coordinate after projection to
the complete word \(A_{gca}\), and define
\begin{align}
 q_{gca}(R)&:=\bigl\langle\!\bigl\langle\lambda_R,
   \operatorname{Inc}_{A_{gca}}^\Sigma(R)\bigr\rangle\!\bigr\rangle_R,
 &a_{gca}(R)&:=(q_{gca}(R))_+,\label{eq:gca-current-amplitude}\\
 g_{gca}(R)&:=-\int_s^t\langle K_r^{1/2}\lambda_{R,r},
                     K_r^{1/2}\zeta_{gca,r}\rangle\,\dd r,
 &d_{gca}(R)&:=\int_s^t\|K_r^{1/2}\zeta_{gca,r}\|^2\,\dd r,\\
 h_{gca}(R)&:=\int_s^t\|K_r^{1/2}\lambda_{R,r}\|^2\,\dd r,
 &p_{gca}(R)&:=p_{\Phys}(R),\label{eq:gca-geometric-costs}\\
 v_{gca}(R)&:=\bigl\langle\!\bigl\langle\lambda_R,
  E_qP+J(I-\Pi_q)\bigr\rangle\!\bigr\rangle_{R,A_{gca}},
 &c_{gca}(R)&:=(v_{gca}(R))_+.
 \label{eq:gca-quotient-current}
\end{align}
The source-pure horizontal quotient term transfers to its inherited word;
each joint horizontal interaction remains in the equality current below.
If the exact-word causal class vanishes and its total preimage graph returns
an ordinary potential, let \(\Lambda_{gca}(R^\pm)\) be its endpoint target
readouts.  For a nonzero class set both readouts to zero and retain the class
below.  A zero class with nonordinary preimage exits through its graph.

On the common ordinary row set
\[
 \mathcal G_{gca}(R):=\operatorname{pr}_{A_{gca}}
   \bigl((-K_r\zeta_{gca,r})_{r\in[s,t]}\bigr).
\]
This is a generated stopped current, and its pairing with \(\lambda_R\) is
\(g_{gca}(R)\).  The expanded public derivation has the exact current identity
\begin{equation}
 \label{eq:gca-exact-word-current-decomposition}
 \operatorname{pr}_{A_{gca}}\operatorname{Inc}^{\Sigma}(R)
 =\mathcal G_{gca}(R)
  +\operatorname{pr}_{A_{gca}}\bigl(E_qP+J(I-\Pi_q)\bigr)
  +\Delta_{\rm stor}K_{gca}(R)+\mathcal E_{gca}^{=}(R),
\end{equation}
where pairing the stopped coboundary gives
\(\Lambda_{gca}(R^+)-\Lambda_{gca}(R^-)\), and
\(\mathcal E_{gca}^{=}\) is the exact subtraction remainder in the expanded
joint derivation.  It retains every unassigned current with word
\(A_{gca}\), including joint horizontal interactions and nonzero causal
classes.  Put
\begin{align}
 w_{gca}(R)&:=\bigl\langle\!\bigl\langle\lambda_R,
                  \mathcal E_{gca}^{=}(R)\bigr\rangle\!\bigr\rangle_R,
 &b_{gca}(R)&:=(w_{gca}(R))_+,\\
 \rho_{gca}(R)&:=\inf\{\rho\ge0:d_{gca}(R)\le\rho p_{gca}(R)\},&
 \sigma_{gca}(R)&:=\inf\{\sigma\ge0:c_{gca}(R)\le\sigma p_{gca}(R)\},\\
 \kappa_{gca}(R)&:=\inf\{\kappa\ge0:b_{gca}(R)\le\kappa p_{gca}(R)\}.
 &&\label{eq:gca-admission-values}
\end{align}
For each numerator, zero numerator and zero carrier give admission zero;
positive numerator and zero carrier give \(+\infty\), retained as a finite
failed-domination coordinate.  Pairing
\eqref{eq:gca-exact-word-current-decomposition} gives
\begin{equation}
 \label{eq:gca-exact-current-identity}
 q_{gca}=g_{gca}+v_{gca}+w_{gca}
 +\Lambda_{gca}(R^+)-\Lambda_{gca}(R^-).
\end{equation}
Consequently, for every \(\theta>0\),
\begin{equation}
 \label{eq:gca-local-conversion}
 a_{gca}(R)\le \frac\theta2d_{gca}(R)
 +\frac1{2\theta}h_{gca}(R)+c_{gca}(R)+b_{gca}(R)
 +\bigl(\Lambda_{gca}(R^+)-\Lambda_{gca}(R^-)\bigr)_+.
\end{equation}

The finite exact-word detector is
\begin{equation}
 \label{eq:gca-incidence-detector}
 \delta_{gca}(R):=\left\|\left(
 G_\Sigma^{\Geom}(C_R),\operatorname{Ker}_\Sigma^{\Geom}(C_R),
 \operatorname{Gc}_\Sigma^{\Geom}(C_R),
 H_\Sigma^{\Caus}(C_R),S_\Sigma^{\Caus}(C_R),
 (\Psi_\Sigma^{\Caus}(C_R;\lambda_k))_{1\le k\le m(R)},
 \operatorname{Gc}_\Sigma^{\Caus}(C_R),
 A_\Sigma^{\rm desc}(C_R),A_\Sigma^{\rm vert}(C_R),
 A_\Sigma^{\rm coker}(C_R),A_\Sigma^{\rm gc}(C_R),
 A_\Sigma^{\rm tgt}(C_R),\operatorname{Inc}_{A_{gca}}^\Sigma(R)
 \right)\right\|_{\rm det}.
\end{equation}
The probes \(\lambda_k\) are exactly those stored in \(R\).  Vanishing on
the cofinal joint hierarchy is equivalent to absence of exact-
\(A_{gca}\) incidence.  Geometric capacity, quotient rank, domination
radii, and cycle price remain valuations rather than absence detectors.

Let \(\mathcal F_N^{\rm ord}\) be an ordered stopped complex
of pairwise disjoint ordinary records, with matching adjacent cuts and all
intervening gaps retained, and assume explicitly that
\(q_{gca}(R_j)\ge0\) for every \(j\).  Define
\begin{align*}
 \mathsf A_{gca}(\mathcal F_N^{\rm ord})&:=\sum_{j\le N}a_{gca}(R_j),&
 \mathsf H_{gca}(\mathcal F_N^{\rm ord})&:=\sum_{j\le N}h_{gca}(R_j),&
 \mathsf P_{gca}(\mathcal F_N^{\rm ord})&:=\sum_{j\le N}p_{gca}(R_j),\\
 \rho_{gca}^*(\mathcal F_N^{\rm ord})&:=\max_{j\le N}\rho_{gca}(R_j),&
 \sigma_{gca}^*(\mathcal F_N^{\rm ord})&:=\max_{j\le N}\sigma_{gca}(R_j),&
 \kappa_{gca}^*(\mathcal F_N^{\rm ord})&:=\max_{j\le N}\kappa_{gca}(R_j).
\end{align*}
If the three maxima are finite, then
\begin{equation}
 \label{eq:gca-cumulative-conversion}
 \mathsf A_{gca}(\mathcal F_N^{\rm ord})
 \le\left(\frac\theta2\rho_{gca}^*(\mathcal F_N^{\rm ord})
 +\sigma_{gca}^*(\mathcal F_N^{\rm ord})
 +\kappa_{gca}^*(\mathcal F_N^{\rm ord})\right)E
 +\frac1{2\theta}\mathsf H_{gca}(\mathcal F_N^{\rm ord})
 +\Lambda_{gca}(R_N^+)-\Lambda_{gca}(R_1^-),
 \qquad \mathsf P_{gca}(\mathcal F_N^{\rm ord})\le E.
\end{equation}
Every completed normalized ordinary exact-\(A_{gca}\) crossing satisfies
\begin{equation}
 \label{eq:gca-canonical-amplitude}
 Q_{A_{gca}}(R)=q_{gca}(R)=a_{gca}(R)=1.
\end{equation}

For an arbitrary ordered first-\(N\) family \(\mathcal G_N\), let
\(\mathsf B_{gca}(\mathcal G_N)\) count records whose unique first-active
flag is a geometric form, kernel/cokernel, or graph exit; a causal gluing,
preimage, support, resolvent, or return graph exit; or a quotient-cokernel or
quotient-graph exit.  Simultaneous detector coordinates retained behind that
flag are not counted again.  Its ordinary subfamily obeys
\begin{equation}
 \label{eq:gca-branch-partition}
 N=\mathsf B_{gca}(\mathcal G_N)
   +\mathsf A_{gca}(\mathcal G_N^{\rm ord}),
\end{equation}
with empty maxima equal to zero.  Hence an infinite exact-\(A_{gca}\)
cascade in a finite physical stratum forces unbounded
\(\mathsf B_{gca}(\mathcal G_N)\),
\(\rho_{gca}^*(\mathcal G_N^{\rm ord})\),
\(\sigma_{gca}^*(\mathcal G_N^{\rm ord})\),
\(\kappa_{gca}^*(\mathcal G_N^{\rm ord})\),
\(\mathsf H_{gca}(\mathcal G_N^{\rm ord})\), or the corresponding terminal
storage difference.  Every alternative
has finite first-crossing prefixes retaining \(A_{gca}\).  Storage re-entry
is evaluated by cycle price, giving divergent recurrent physical charge or
the next zero-price return support.  Zero-carrier form, quotient, and
equality currents enter same-word equality saturation.  An orbit or
singularity conclusion additionally requires ordinary same-origin
realization and public noncontinuation.
\end{theorem}
\begin{proof}
Join the three normal forms before exact-word projection.  Their ordinary
rows give the geometric current, quotient defect/vertical current, and
causal storage coboundary.  Source-pure horizontal descent transfers to its
inherited word.  Closure of the expanded derivation under subtraction leaves
exactly \(\mathcal E_{gca}^{=}\), proving
\eqref{eq:gca-exact-word-current-decomposition}.  Pairing gives
\eqref{eq:gca-exact-current-identity}.  Taking positive parts and applying
Cauchy--Schwarz and Young to \(g_{gca}\) proves
\eqref{eq:gca-local-conversion}.

Sum the signed identity over the ordinary stopped complex before taking
positive parts.  The sign assumption gives
\[
 \mathsf A_{gca}(\mathcal F_N^{\rm ord})
 =\sum_{j\le N}q_{gca}(R_j)
 \le \sum_{j\le N}\bigl(|g_{gca}(R_j)|+c_{gca}(R_j)+b_{gca}(R_j)\bigr)
 +\Lambda_{gca}(R_N^+)-\Lambda_{gca}(R_1^-),
\]
because matching cuts cancel every internal storage value.  Apply
Cauchy--Schwarz and Young separately on each record, then the three finite
admission inequalities.  Finite additivity of the single carrier and its
mass bound give \eqref{eq:gca-cumulative-conversion}.  Infinite admissions
are already the retained failed-domination branches, so they are not used in
that inequality.

The normalized stopped chain rule, exact-support projection, and target-null
residual closure prove \eqref{eq:gca-canonical-amplitude}.  Ordered total
graphs partition each crossing into one early exit or the ordinary row,
which proves \eqref{eq:gca-branch-partition}.  If the exit count were bounded,
it would stabilize, leaving an ordinary tail with all intervening cuts and
gaps retained.  If the remaining five coordinates were also bounded,
\eqref{eq:gca-cumulative-conversion} would bound that tail's crossing count,
contradicting the partition.  Running maxima, cumulative probe cost, and
ordered storage prefixes have finite first crossings.  For the detector
claim, cofinal vanishing deletes every atomic path by the three normal-form
exhaustion theorems and deletes every remaining mixed path through the
displayed exact conditional increment.  Conversely, any nonzero displayed
coordinate is a finite exact-\(A_{gca}\) incidence path.  This proves the
detector equivalence.  Cycle-price re-entry is
\cref{thm:continuum-caus-normal-form-exhaustion}; zero-carrier re-entry is
\cref{thm:actual-source-price-equality-decision}.  Source-preserving
totalization retains \(A_{gca}\), and realization and noncontinuation are
evaluated through \eqref{eq:actual-contact-cell-realization-intersection}.
\end{proof}

\begin{theorem}[Exact pure-\(\Lift\) target-current accountability]
\label{thm:quantitative-source-cell-l}
Let \(R\in\mathfrak D_{A_l,\Sigma,n}^{\rm pre}(E)\) have exact word
\(A_l=\{\Lift\}\), normalized cell \(C_R\), shell covector \(\lambda_R\),
and public lift data \((\Omega,\pi,U,P_\kappa,\Pi_\pi)\).  Evaluate every
base-descent, lifting, averaging, projection, compactification, scale,
localization, product, and realization operation by its ordered total graph.
A first nonordinary value returns its finite \(\Lift\)-owned graph branch.
The scalar formulas below apply precisely to the common ordinary row.

Let
\(\mathscr X_l:=\{\mathrm{int},\mathrm{fib},\mathrm{conc},
\mathrm{osc},\mathrm{sc}\}\), and define the five exact-word currents
\begin{align*}
 \mathcal L_l^{\rm int}(R)&:=\operatorname{pr}_{A_l}(E_\pi P_\kappa),&
 \mathcal L_l^{\rm fib}(R)&:=\operatorname{pr}_{A_l}
                  (J^\Omega(I-\Pi_\pi)),\\
 \mathcal L_l^{\rm conc}(R)&:=\operatorname{pr}_{A_l}J_{\rm conc},&
 \mathcal L_l^{\rm osc}(R)&:=\operatorname{pr}_{A_l}J_{\rm osc},\\
 \mathcal L_l^{\rm sc}(R)&:=\operatorname{pr}_{A_l}J_{\rm sc}^{\rm coh}.&&
\end{align*}
The source-pure base current \(UJ_XP_\kappa\) transfers to its inherited
word and has zero exact-\(A_l\) projection.  Applying
\eqref{eq:target-lift-exact-current-reconstruction} to \(R\), let
\(\Lambda_l^{\rm sc}(R^\pm)\) be the endpoint target readouts of the
ordinary scale potential \(K_l\).  If its total preimage is nonordinary,
return that graph value and stop.  On the ordinary row,
\begin{equation}
 \label{eq:lift-exact-word-current-decomposition}
 \operatorname{pr}_{A_l}\operatorname{Inc}^{\Sigma}(R)
 =\sum_{x\in\mathscr X_l}\mathcal L_l^x(R)
   +\Delta_{\rm sc}K_l(R)+N_l(R),
\end{equation}
where \(N_l(R)\in\Null_\Sigma\) is the target-null coordinate of the
exhaustive lift normal form.  Every remaining exact-\(A_l\) interaction is
therefore assigned to one of the five displayed currents and
\(\langle\!\langle\lambda_R,N_l(R)\rangle\!\rangle_R=0\).

Define, before normalization,
\begin{align}
 q_l(R)&:=\bigl\langle\!\bigl\langle\lambda_R,
       \operatorname{Inc}_{A_l}^{\Sigma}(R)\bigr\rangle\!\bigr\rangle_R,
 &a_l(R)&:=(q_l(R))_+,&p_l(R)&:=p_{\Phys}(R),
 \label{eq:lift-current-amplitude}\\
 j_l^x(R)&:=\bigl\langle\!\bigl\langle\lambda_R,
                  \mathcal L_l^x(R)\bigr\rangle\!\bigr\rangle_R,
 &c_l^x(R)&:=(j_l^x(R))_+,
 &\alpha_l^x(R)&:=\inf\{\alpha\ge0:c_l^x(R)\le\alpha p_l(R)\},
 \quad x\in\mathscr X_l.\label{eq:lift-coordinate-admissions}
\end{align}
Pairing \eqref{eq:lift-exact-word-current-decomposition} gives
\begin{equation}
 \label{eq:lift-exact-current-identity}
 q_l(R)=\sum_{x\in\mathscr X_l}j_l^x(R)
   +\Lambda_l^{\rm sc}(R^+)-\Lambda_l^{\rm sc}(R^-),
 \qquad
 a_l(R)\le\sum_{x\in\mathscr X_l}c_l^x(R)
  +\bigl(\Lambda_l^{\rm sc}(R^+)-\Lambda_l^{\rm sc}(R^-)\bigr)_+.
\end{equation}
For every one of the five numerators, zero numerator and zero carrier give
admission zero, while positive numerator and zero carrier give \(+\infty\).
An infinite value is retained as that finite coordinate's failed-domination
row and is never multiplied by zero.

The finite detector is
\begin{equation}
 \label{eq:lift-incidence-detector}
 \delta_l(R):=\left\|\left(
 L_\Sigma^{\rm int}(C_R),L_\Sigma^{\rm fib}(C_R),
 L_\Sigma^{\rm conc}(C_R),L_\Sigma^{\rm osc}(C_R),
 L_\Sigma^{\rm stor}(C_R),L_\Sigma^{\rm sc}(C_R),
 L_\Sigma^{\rm gc}(C_R),
 \operatorname{Inc}_{A_l}^{\Sigma}(R)
 \right)\right\|_{\rm det}.
\end{equation}
It is the simultaneous zero test for this finite record.  Vanishing on the
cofinal compactification, moment, and scale hierarchy is equivalent to
absence of exact-\(A_l\) incidence.  The scale valuation and scale velocity
remain stored coordinates and are not separate absence tests.

For an ordered stopped complex \(\mathcal F_N^{\rm ord}\) of pairwise
disjoint ordinary records with matching adjacent cuts, retained intervening
gaps, and \(q_l(R_j)\ge0\) for every \(j\), define
\begin{align*}
 \mathsf A_l(\mathcal F_N^{\rm ord})&:=\sum_{j\le N}a_l(R_j),&
 \mathsf P_l(\mathcal F_N^{\rm ord})&:=\sum_{j\le N}p_l(R_j),\\
 (\alpha_l^x)^*(\mathcal F_N^{\rm ord})&:=
       \max_{j\le N}\alpha_l^x(R_j)\quad(x\in\mathscr X_l).&
\end{align*}
If all five maxima are finite, then
\begin{equation}
 \label{eq:lift-cumulative-conversion}
 \mathsf A_l(\mathcal F_N^{\rm ord})
 \le\left(\sum_{x\in\mathscr X_l}(\alpha_l^x)^*
                 (\mathcal F_N^{\rm ord})\right)
       \mathsf P_l(\mathcal F_N^{\rm ord})
 +\Lambda_l^{\rm sc}(R_N^+)-\Lambda_l^{\rm sc}(R_1^-)
 \le\left(\sum_{x\in\mathscr X_l}(\alpha_l^x)^*
                 (\mathcal F_N^{\rm ord})\right)E
 +\Lambda_l^{\rm sc}(R_N^+)-\Lambda_l^{\rm sc}(R_1^-).
\end{equation}
Every completed normalized ordinary exact-\(A_l\) crossing satisfies
\begin{equation}
 \label{eq:lift-canonical-amplitude}
 Q_{A_l}(R)=q_l(R)=a_l(R)=1.
\end{equation}

For an arbitrary ordered family \(\mathcal G_N\) of completed normalized
exact-\(A_l\) crossings, let \(\mathsf B_l(\mathcal G_N)\) count records
whose unique first-active flag is a nonordinary base-descent, lift,
averaging, projection, compactification, scale, localization, product, or
realization graph value.  Its ordinary subfamily satisfies
\begin{equation}
 \label{eq:lift-branch-partition}
 N=\mathsf B_l(\mathcal G_N)+\mathsf A_l(\mathcal G_N^{\rm ord}),
\end{equation}
with every empty maximum equal to zero.  Hence an infinite exact-\(A_l\)
cascade in a finite physical stratum forces \(\mathsf B_l\) to be unbounded
or at least one of the five coordinate maxima \((\alpha_l^x)^*\) to be
unbounded, or forces the terminal scale-storage difference to be unbounded.
Every
alternative has finite first-crossing records with word \(A_l\).
Scale-storage prefixes re-enter through the total scale/modulation graph;
positive physical price is charged and zero price enters same-word equality
saturation.  Zero-carrier coordinate currents use the same equality rule.  An
orbit or singularity conclusion further
requires ordinary same-origin realization and public noncontinuation.
\end{theorem}
\begin{proof}
Project \eqref{eq:target-lift-normal-form} only after provenance is fixed.
The base term transfers to its inherited word.  The interface, fiber,
concentration, oscillation, and scale currents are precisely the five
sequential ordinary \(\Lift\)-headed rows of
\eqref{eq:target-lift-exact-current-reconstruction}.  That reconstruction
places the final row in \(\Null_\Sigma\), with zero stopped target pairing.
This proves
\eqref{eq:lift-exact-word-current-decomposition}.  Pairing and taking
positive parts prove \eqref{eq:lift-exact-current-identity}.

Sum the signed identity before any normalization.  Since \(q_l(R_j)\ge0\),
its left side is \(\sum_jq_l(R_j)=\mathsf A_l\).  Matching adjacent cuts
cancel every internal scale-storage readout.  At finite running maxima, each of the five
admission definitions bounds its numerator by the same restricted physical
carrier.  Finite additivity and total mass at most \(E\) prove
\eqref{eq:lift-cumulative-conversion}.  Infinite values are already the
failed-domination rows stated in the theorem.

The normalized stopped chain rule, exact-support projection, and target-null
residual closure prove \eqref{eq:lift-canonical-amplitude}.  The ordered
total graph partitions every crossing into one graph exit or the ordinary
row, proving \eqref{eq:lift-branch-partition}.  A bounded integer exit count
stabilizes.  Its remaining tail is an ordinary stopped complex.  Boundedness
of all five ordinary maxima and of its terminal scale-storage difference
would then make
\eqref{eq:lift-cumulative-conversion} bound the remaining unit crossings.
Thus one stated sequence is unbounded, and its running maximum has finite
first crossings; an unbounded storage difference likewise has finite first
integer-crossing prefixes.  Scale-storage re-entry and zero-carrier
coordinate re-entry are both processed by
\cref{thm:actual-source-price-equality-decision} on the retained \(A_l\)
record.  Cofinal detector equivalence is
\cref{thm:continuum-lift-normal-form-exhaustion}, with the exact increment
retaining every mixed public interaction.  Source-preserving
totalization retains \(A_l\); realization and noncontinuation are evaluated
through \eqref{eq:actual-contact-cell-realization-intersection}.
\end{proof}

\begin{theorem}[Sequential exact-word partition for \(\Geom\Lift\)]
\label{thm:geom-lift-sequential-exact-word-partition}
Let \(R\) be a finite stopped record with exact word
\(A_{gl}=\{\Geom,\Lift\}\), and suppose all geometric and lift operations
on its selected row are ordinary.  Apply the fixed source projection order
\(\Geom\) before \(\Lift\).  The geometric normal form returns a unique
exact-word current \(\mathcal G_{gl}(R)\).  Set
\[
 \mathcal R_{gl}^{(1)}(R)
 :=\operatorname{pr}_{A_{gl}}\operatorname{Inc}^{\Sigma}(R)
   -\mathcal G_{gl}(R).
\]
Apply the sequential reconstruction
\eqref{eq:target-lift-exact-current-reconstruction} to
\(\mathcal R_{gl}^{(1)}(R)\).  It returns unique currents
\(\mathcal L_{gl}^x(R)\), where
\(\mathscr X_l:=\{\mathrm{int},\mathrm{fib},\mathrm{conc},
\mathrm{osc},\mathrm{sc}\}\), an exact scale
coboundary \(\Delta_{\rm sc}K_{gl}(R)\), and
\(N_{gl}(R)\in\Null_\Sigma\), and
\begin{equation}
 \label{eq:geom-lift-sequential-partition}
 \operatorname{pr}_{A_{gl}}\operatorname{Inc}^{\Sigma}(R)
 =\mathcal G_{gl}(R)+\sum_{x\in\mathscr X_l}\mathcal L_{gl}^x(R)
  +\Delta_{\rm sc}K_{gl}(R)+N_{gl}(R).
\end{equation}
The summands are disjoint by sequential assignment: a constructor vertex is
assigned to the first active row in the fixed order and is not reconsidered
by a later row.  Every mixed vertex retains the union word \(A_{gl}\).
\end{theorem}

\begin{proof}
The geometric range--kernel--graph exhaustion gives a unique first geometric
block or a total-graph exit.  On the ordinary row subtracting that block in
the expanded derivation defines \(\mathcal R_{gl}^{(1)}\) exactly.  The lift
reconstruction then assigns, in order, its interface, fiber, concentration,
oscillation, exact scale-storage, and scale-cohomology blocks.  The
subtractions defining later blocks occur after the earlier blocks have been
removed, which proves disjointness and the displayed algebraic identity.

If the final residual had nonzero target projection, its expanded derivation
would have a first active source.  Its exact word contains only \(\Geom\) and
\(\Lift\).  A \(\Geom\)-headed residual contradicts geometric exhaustion;
a \(\Lift\)-headed residual contradicts
\eqref{eq:target-lift-exact-current-reconstruction}.  Hence its target
projection is zero and it is \(N_{gl}\in\Null_\Sigma\).  Source-preserving
projection and subtraction retain the union ancestry of every mixed vertex.
\end{proof}

\begin{theorem}[Exact joint \(\Geom\Lift\) target-current accountability]
\label{thm:quantitative-source-cell-gl}
Let \(R=[s,t]\in\mathfrak D_{A_{gl},\Sigma,n}^{\rm pre}(E)\) have exact
word \(A_{gl}=\{\Geom,\Lift\}\), normalized cell \(C_R\), shell covector
\(\lambda_{R,r}=D\Phi_R(u_r)\), and public lift data.  Form the complete
joint record before source projection.  Totalize every geometric form,
range, kernel/cokernel, and graph row and every lift operation listed in
\eqref{eq:target-lift-normal-form}.  A first-active geometric
kernel/cokernel or graph value, or a first nonordinary lift value, returns
its finite branch with word \(A_{gl}\).  The following formulas apply to the
common ordinary-range row.

Let \(\zeta_{gl}\) be the geometric-headed coordinate after exact-
\(A_{gl}\) projection and put
\begin{align}
 q_{gl}(R)&:=\bigl\langle\!\bigl\langle\lambda_R,
       \operatorname{Inc}_{A_{gl}}^\Sigma(R)\bigr\rangle\!\bigr\rangle_R,
 &a_{gl}(R)&:=(q_{gl}(R))_+,\\
 g_{gl}(R)&:=-\int_s^t\langle K_r^{1/2}\lambda_{R,r},
                       K_r^{1/2}\zeta_{gl,r}\rangle\,\dd r,
 &d_{gl}(R)&:=\int_s^t\|K_r^{1/2}\zeta_{gl,r}\|^2\,\dd r,\\
 h_{gl}(R)&:=\int_s^t\|K_r^{1/2}\lambda_{R,r}\|^2\,\dd r,
 &p_{gl}(R)&:=p_{\Phys}(R).
 \label{eq:gl-geometric-currents}
\end{align}
Let \(\mathscr X_l=\{\mathrm{int},\mathrm{fib},\mathrm{conc},
\mathrm{osc},\mathrm{sc}\}\), and set
\[
 \mathcal G_{gl}(R):=\operatorname{pr}_{A_{gl}}
   \bigl((-K_r\zeta_{gl,r})_{r\in[s,t]}\bigr).
\]
Let \(\mathcal L_{gl}^x(R)\), \(x\in\mathscr X_l\), be the five currents
returned from the residual by
\cref{thm:geom-lift-sequential-exact-word-partition}.  Thus the geometric
row is not included again in a lift coordinate.  Let
\(\Lambda_{gl}^{\rm sc}(R^\pm)\) be the endpoint readouts of its ordinary
exact scale potential.  A nonordinary scale preimage exits through its graph.
The joint ordinary row has the exact current identity
\begin{equation}
 \label{eq:gl-exact-word-current-decomposition}
 \operatorname{pr}_{A_{gl}}\operatorname{Inc}^{\Sigma}(R)
 =\mathcal G_{gl}(R)+\sum_{x\in\mathscr X_l}\mathcal L_{gl}^x(R)
 +\Delta_{\rm sc}K_{gl}(R)+N_{gl}(R),
 \qquad N_{gl}(R)\in\Null_\Sigma,
\end{equation}
This is \eqref{eq:geom-lift-sequential-partition}.  In particular, joint
geometric--lift interactions are assigned to exactly one row, and the
stopped target pairing of \(N_{gl}\) is zero.

For \(x\in\mathscr X_l\), define
\begin{align}
 j_{gl}^x(R)&:=\bigl\langle\!\bigl\langle\lambda_R,
                    \mathcal L_{gl}^x(R)\bigr\rangle\!\bigr\rangle_R,
 &c_{gl}^x(R)&:=(j_{gl}^x(R))_+,\\
 \alpha_{gl}^x(R)&:=\inf\{\alpha\ge0:c_{gl}^x(R)\le\alpha p_{gl}(R)\},&
 \rho_{gl}(R)&:=\inf\{\rho\ge0:d_{gl}(R)\le\rho p_{gl}(R)\}.
 \label{eq:gl-admission-values}
\end{align}
Pairing \eqref{eq:gl-exact-word-current-decomposition} gives
\begin{equation}
 \label{eq:gl-exact-current-identity}
 q_{gl}=g_{gl}+\sum_{x\in\mathscr X_l}j_{gl}^x
 +\Lambda_{gl}^{\rm sc}(R^+)-\Lambda_{gl}^{\rm sc}(R^-).
\end{equation}
For every \(\theta>0\),
\begin{equation}
 \label{eq:gl-local-conversion}
 a_{gl}(R)\le\frac\theta2d_{gl}(R)+\frac1{2\theta}h_{gl}(R)
 +\sum_{x\in\mathscr X_l}c_{gl}^x(R)
 +\bigl(\Lambda_{gl}^{\rm sc}(R^+)
       -\Lambda_{gl}^{\rm sc}(R^-)\bigr)_+.
\end{equation}
Each of the six admissions is zero when both numerator and carrier vanish,
and is \(+\infty\) when its numerator is positive and its carrier vanishes.
An infinite value is its finite failed-domination row.

The finite joint detector is
\begin{equation}
 \label{eq:gl-incidence-detector}
 \delta_{gl}(R):=\left\|\left(
 G_\Sigma^{\Geom}(C_R),\operatorname{Ker}_\Sigma^{\Geom}(C_R),
 \operatorname{Gc}_\Sigma^{\Geom}(C_R),
 L_\Sigma^{\rm int}(C_R),L_\Sigma^{\rm fib}(C_R),
 L_\Sigma^{\rm conc}(C_R),L_\Sigma^{\rm osc}(C_R),
 L_\Sigma^{\rm stor}(C_R),L_\Sigma^{\rm sc}(C_R),
 L_\Sigma^{\rm gc}(C_R),\operatorname{Inc}_{A_{gl}}^\Sigma(R)
 \right)\right\|_{\rm det}.
\end{equation}
Its cofinal vanishing is equivalent to absence of exact-\(A_{gl}\)
incidence.  Geometric capacity and the auxiliary scale valuations remain
attached but are not independent absence tests.

Let \(\mathcal F_N^{\rm ord}\) be an ordered stopped complex of pairwise
disjoint ordinary records with matching cuts, retained gaps, and
\(q_{gl}(R_j)\ge0\).  Define
\begin{align*}
 \mathsf A_{gl}(\mathcal F_N^{\rm ord})&:=\sum_{j\le N}a_{gl}(R_j),&
 \mathsf H_{gl}(\mathcal F_N^{\rm ord})&:=\sum_{j\le N}h_{gl}(R_j),\\
 \mathsf P_{gl}(\mathcal F_N^{\rm ord})&:=\sum_{j\le N}p_{gl}(R_j),&
 \rho_{gl}^*(\mathcal F_N^{\rm ord})&:=\max_{j\le N}\rho_{gl}(R_j),\\
 (\alpha_{gl}^x)^*(\mathcal F_N^{\rm ord})&:=
      \max_{j\le N}\alpha_{gl}^x(R_j),&&x\in\mathscr X_l.
\end{align*}
When all six maxima are finite, abbreviate their common argument only in
the following display:
\begin{equation}
 \label{eq:gl-cumulative-conversion}
 \mathsf A_{gl}
 \le\left(\frac\theta2\rho_{gl}^*
       +\sum_{x\in\mathscr X_l}(\alpha_{gl}^x)^*\right)E
 +\frac1{2\theta}\mathsf H_{gl}
 +\Lambda_{gl}^{\rm sc}(R_N^+)-\Lambda_{gl}^{\rm sc}(R_1^-),
 \qquad \mathsf P_{gl}\le E.
\end{equation}
Every completed normalized ordinary exact-\(A_{gl}\) crossing satisfies
\begin{equation}
 \label{eq:gl-canonical-amplitude}
 Q_{A_{gl}}(R)=q_{gl}(R)=a_{gl}(R)=1.
\end{equation}

For an arbitrary ordered family \(\mathcal G_N\), let
\(\mathsf B_{gl}(\mathcal G_N)\) count records whose unique first-active
flag is a geometric form, kernel/cokernel, or graph exit, or any lift graph
exit.  Simultaneous detector coordinates are not counted twice.  Then
\begin{equation}
 \label{eq:gl-branch-partition}
 N=\mathsf B_{gl}(\mathcal G_N)+\#\mathcal G_N^{\rm ord}
  =\mathsf B_{gl}(\mathcal G_N)
   +\sum_{R\in\mathcal G_N^{\rm ord}}a_{gl}(R).
\end{equation}
Thus an infinite exact-\(A_{gl}\) cascade forces unbounded
\(\mathsf B_{gl}\), \(\rho_{gl}^*\), some \((\alpha_{gl}^x)^*\),
\(\mathsf H_{gl}\), or terminal scale storage.  Every alternative has
finite first crossings retaining \(A_{gl}\).  Scale-storage and zero-carrier
records re-enter through the same-word scale graph and equality decision.
An orbit or singularity conclusion additionally requires ordinary
same-origin realization and public noncontinuation.
\end{theorem}
\begin{proof}
Apply \cref{thm:geom-lift-sequential-exact-word-partition} to the joined
record.  Its identity is \eqref{eq:gl-exact-word-current-decomposition}, and
its last row is target-null.  Pairing proves
\eqref{eq:gl-exact-current-identity}.
Cauchy--Schwarz and Young on \(g_{gl}\), followed by positive parts, prove
\eqref{eq:gl-local-conversion}.

Sum the signed identity over the ordinary complex.  Target orientation
and matching cuts give the exact telescoped inequality
\[
 \mathsf A_{gl}(\mathcal F_N^{\rm ord})
 \le\sum_{j\le N}|g_{gl}(R_j)|
   +\sum_{x\in\mathscr X_l}\sum_{j\le N}c_{gl}^x(R_j)
   +\Lambda_{gl}^{\rm sc}(R_N^+)-\Lambda_{gl}^{\rm sc}(R_1^-).
\]
Cauchy--Schwarz and Young bound the first sum.  The six finite admissions
and the single-carrier mass bound then give
\eqref{eq:gl-cumulative-conversion}.  Infinite admissions
are retained branches rather than factors in this inequality.

The normalized stopped chain rule, exact projection, and null residual prove
\eqref{eq:gl-canonical-amplitude}.  Unique first-active flags prove
\eqref{eq:gl-branch-partition}.  If the integer exit count is bounded, its
ordinary tail retains all cuts and gaps.  Boundedness of the six admissions,
probe cost, and terminal storage would make
\eqref{eq:gl-cumulative-conversion} bound its unit crossings.  Hence one
listed sequence is unbounded.  Running maxima and cumulative probe cost have
finite first crossings, and an unbounded terminal storage sequence has a
finite first prefix above every integer level.  For detector equivalence,
cofinal vanishing removes the geometric rows by geometric exhaustion, every
lift row including scale storage by lift exhaustion, and the remaining mixed
row by the exact conditional increment.  Conversely, any nonzero displayed
coordinate is a finite exact-\(A_{gl}\) incidence path.  Storage and
zero-price re-entry use
\cref{thm:actual-source-price-equality-decision}.  Totalization retains
\(A_{gl}\); realization and noncontinuation are evaluated through
\eqref{eq:actual-contact-cell-realization-intersection}.
\end{proof}

\begin{theorem}[Sequential exact-word partition for \(\Caus\Lift\)]
\label{thm:caus-lift-sequential-exact-word-partition}
Let \(R\) have exact word \(A_{cl}=\{\Caus,\Lift\}\) and lie on the common
ordinary row.  Apply \(\Caus\) before \(\Lift\).  The causal normal form
first selects, in the expanded derivation, the signed sum of vertices whose
first active source is \(\Caus\); call this current \(F_{cl}^{(0)}(R)\).
If its class is zero and the total preimage is ordinary, define
\(\Delta_{\rm caus}K_{cl}(R)\) to be that coboundary and set
\(\mathcal C_{cl}:=F_{cl}^{(0)}-\Delta_{\rm caus}K_{cl}\).  If its class is
nonzero, set \(\Delta_{\rm caus}K_{cl}:=0\) and
\(\mathcal C_{cl}:=F_{cl}^{(0)}\).  A zero class with nonordinary preimage
exits through its graph.  Subtract both ordinary currents from
\(\operatorname{pr}_{A_{cl}}\operatorname{Inc}^{\Sigma}(R)\), then apply
\eqref{eq:target-lift-exact-current-reconstruction} to the residual.  With
\(\mathscr X_l=\{\mathrm{int},\mathrm{fib},\mathrm{conc},
\mathrm{osc},\mathrm{sc}\}\), this gives
\begin{equation}
 \label{eq:caus-lift-sequential-partition}
 \operatorname{pr}_{A_{cl}}\operatorname{Inc}^{\Sigma}(R)
 =\mathcal C_{cl}(R)+\Delta_{\rm caus}K_{cl}(R)
  +\sum_{x\in\mathscr X_l}\mathcal L_{cl}^x(R)
  +\Delta_{\rm sc}K_{cl}(R)+N_{cl}(R),
 \qquad N_{cl}(R)\in\Null_\Sigma.
\end{equation}
Every mixed vertex is assigned once, to its first active row, and retains
word \(A_{cl}\).
\end{theorem}
\begin{proof}
Signed causal gluing precedes the cohomology projection, so
\(F_{cl}^{(0)}\) is a public current rather than a quotient label.  The three
class/preimage cases are disjoint and exhaustive by the causal normal-form
theorem.  Their displayed subtraction defines the remaining current exactly.
The sequential lift reconstruction partitions
that residual into its five currents, exact scale coboundary, and null row.
Any target-visible final residual would be \(\Caus\)- or \(\Lift\)-headed,
contradicting the corresponding exhaustion theorem.  Sequential subtraction
proves disjointness and preserves ancestry.
\end{proof}

\begin{theorem}[Exact joint \(\Caus\Lift\) target-current accountability]
\label{thm:quantitative-source-cell-cl}
Let \(R\in\mathfrak D_{A_{cl},\Sigma,n}^{\rm pre}(E)\) satisfy the
hypotheses of \cref{thm:caus-lift-sequential-exact-word-partition}, with
shell covector \(\lambda_R\).  Totalize every causal and lift operation; a
first nonordinary value returns its unique graph branch.  On the ordinary
row write \(\mathscr X_l:=\{\mathrm{int},\mathrm{fib},\mathrm{conc},
\mathrm{osc},\mathrm{sc}\}\) and define
\begin{align}
 q_{cl}(R)&:=\langle\!\langle\lambda_R,
      \operatorname{Inc}_{A_{cl}}^\Sigma(R)\rangle\!\rangle_R,&
 a_{cl}(R)&:=(q_{cl}(R))_+,&p_{cl}(R)&:=p_{\Phys}(R),\\
 f_{cl}(R)&:=\langle\!\langle\lambda_R,\mathcal C_{cl}(R)
                      \rangle\!\rangle_R,&
 c_{cl}(R)&:=(f_{cl}(R))_+,&
 \kappa_{cl}(R)&:=\inf\{k\ge0:c_{cl}(R)\le k p_{cl}(R)\},\\
 j_{cl}^x(R)&:=\langle\!\langle\lambda_R,\mathcal L_{cl}^x(R)
                       \rangle\!\rangle_R,&
 c_{cl}^x(R)&:=(j_{cl}^x(R))_+,&
 \alpha_{cl}^x(R)&:=\inf\{a\ge0:c_{cl}^x(R)\le a p_{cl}(R)\}.
 \label{eq:cl-amplitudes-admissions}
\end{align}
Let \(\Lambda_{cl}^{\rm caus}(R^\pm)\) and
\(\Lambda_{cl}^{\rm sc}(R^\pm)\) be the causal and scale endpoint readouts.
In the next identity \(\Delta\Lambda^y_{cl}\) means exactly
\(\Lambda^y_{cl}(R^+)-\Lambda^y_{cl}(R^-)\), for
\(y\in\{\mathrm{caus},\mathrm{sc}\}\).
Pairing \eqref{eq:caus-lift-sequential-partition} yields
\begin{equation}
 \label{eq:cl-exact-current-identity}
 q_{cl}=f_{cl}+\sum_{x\in\mathscr X_l}j_{cl}^x
 +\Delta\Lambda_{cl}^{\rm caus}+\Delta\Lambda_{cl}^{\rm sc},
 \qquad
 a_{cl}\le c_{cl}+\sum_xc_{cl}^x
 +\bigl(\Delta\Lambda_{cl}^{\rm caus}+\Delta\Lambda_{cl}^{\rm sc}\bigr)_+.
\end{equation}
All six admissions use the usual extended-real convention: zero numerator
and zero carrier give zero, while positive numerator and zero carrier give a
finite failed-domination record with value \(+\infty\).

Define
\begin{equation}
 \label{eq:cl-incidence-detector}
 \delta_{cl}(R):=\left\|\left(
 H_\Sigma^{\Caus}(C_R),S_\Sigma^{\Caus}(C_R),
 (\Psi_\Sigma^{\Caus}(C_R;\lambda_k))_{1\le k\le m(R)},
 \operatorname{Gc}_\Sigma^{\Caus}(C_R),
 L_\Sigma^{\rm int}(C_R),L_\Sigma^{\rm fib}(C_R),
 L_\Sigma^{\rm conc}(C_R),L_\Sigma^{\rm osc}(C_R),
 L_\Sigma^{\rm stor}(C_R),L_\Sigma^{\rm sc}(C_R),
 L_\Sigma^{\rm gc}(C_R),
 \operatorname{Inc}_{A_{cl}}^\Sigma(R)\right)\right\|_{\rm det}.
\end{equation}
Here \(\lambda_1,\ldots,\lambda_{m(R)}\) are exactly the probes stored in
\(R\).  Cofinal vanishing is equivalent to exact-
\(A_{cl}\) absence.

Let \(\mathcal F_N^{\rm ord}\) be an ordered complex of pairwise disjoint
ordinary records with matching cuts, retained gaps, and \(q_{cl}(R_j)\ge0\).
Define
\begin{align*}
 \mathsf A_{cl}(\mathcal F_N^{\rm ord})&:=\sum_{j\le N}a_{cl}(R_j),&
 \mathsf P_{cl}(\mathcal F_N^{\rm ord})&:=\sum_{j\le N}p_{cl}(R_j),\\
 \kappa_{cl}^*(\mathcal F_N^{\rm ord})&:=\max_{j\le N}\kappa_{cl}(R_j),&
 (\alpha_{cl}^x)^*(\mathcal F_N^{\rm ord})&:=
        \max_{j\le N}\alpha_{cl}^x(R_j),\quad x\in\mathscr X_l.
\end{align*}
When all six maxima are finite,
\begin{equation}
 \label{eq:cl-cumulative-conversion}
 \mathsf A_{cl}(\mathcal F_N^{\rm ord})
 \le\left(\kappa_{cl}^*(\mathcal F_N^{\rm ord})
     +\sum_x(\alpha_{cl}^x)^*(\mathcal F_N^{\rm ord})\right)E
 +\Lambda_{cl}^{\rm caus}(R_N^+)-\Lambda_{cl}^{\rm caus}(R_1^-)
 +\Lambda_{cl}^{\rm sc}(R_N^+)-\Lambda_{cl}^{\rm sc}(R_1^-),
 \qquad \mathsf P_{cl}(\mathcal F_N^{\rm ord})\le E.
\end{equation}
Every completed normalized ordinary crossing satisfies
\begin{equation}
 \label{eq:cl-canonical-amplitude}
 Q_{A_{cl}}(R)=q_{cl}(R)=a_{cl}(R)=1.
\end{equation}
If \(\mathsf B_{cl}(\mathcal G_N)\) counts unique first-active causal or lift
graph exits, then
\begin{equation}
 \label{eq:cl-branch-partition}
 N=\mathsf B_{cl}(\mathcal G_N)+\#\mathcal G_N^{\rm ord}
  =\mathsf B_{cl}(\mathcal G_N)+\sum_{R\in\mathcal G_N^{\rm ord}}a_{cl}(R).
\end{equation}
Thus an infinite cascade forces unbounded graph count, one of the six
admission maxima, causal storage, or scale storage.  Each has finite first
crossings retaining \(A_{cl}\).  Causal storage uses the cycle-price return;
scale storage and zero-carrier currents use the same-word equality decision.
An orbit conclusion additionally requires ordinary realization and public
noncontinuation.
\end{theorem}
\begin{proof}
Pair the sequential partition; its null row has zero readout, proving
\eqref{eq:cl-exact-current-identity}.  Sum the signed identity.  Target
orientation and matching cuts give
\[
 \mathsf A_{cl}(\mathcal F_N^{\rm ord})
 \le\sum_{j\le N}c_{cl}(R_j)
   +\sum_{x\in\mathscr X_l}\sum_{j\le N}c_{cl}^x(R_j)
   +\Lambda_{cl}^{\rm caus}(R_N^+)-\Lambda_{cl}^{\rm caus}(R_1^-)
   +\Lambda_{cl}^{\rm sc}(R_N^+)-\Lambda_{cl}^{\rm sc}(R_1^-).
\]
The six admissions plus the single-carrier bound prove
\eqref{eq:cl-cumulative-conversion}.  Infinite
admissions remain failed-domination branches.

The stopped unit chain rule and exact projection prove
\eqref{eq:cl-canonical-amplitude}; unique total-graph flags prove
\eqref{eq:cl-branch-partition}.  A bounded exit count leaves an ordinary
tail with its intervening cuts and gaps retained.  If all admissions and
both terminal storage differences were bounded above,
the cumulative inequality would bound its unit crossings.  Running maxima
and each unbounded storage sequence have finite first crossings.  For the
detector claim, cofinal vanishing removes every causal row by causal
exhaustion, every lift row including scale storage by lift exhaustion, and
the remaining mixed row by the exact conditional increment.  Conversely,
each nonzero displayed coordinate is a finite exact-\(A_{cl}\) incidence
path.  Causal storage
re-entry is \cref{thm:continuum-caus-normal-form-exhaustion}; other zero-price
re-entry is \cref{thm:actual-source-price-equality-decision}.  Realization
and noncontinuation use \eqref{eq:actual-contact-cell-realization-intersection}.
\end{proof}

\begin{theorem}[Local quantitative check: \(\Geom\Caus\Lift\)]
\label{thm:quantitative-source-cell-gcl}
For every finite same-origin stopped record \(R\in\mathfrak D_{A_{gcl},\Sigma,n}^{\rm pre}(E)\) set
\[
\delta_{gcl}(R):=\left\|\left(G^{\Geom}_{\Sigma},\operatorname{Ker}^{\Geom}_{\Sigma},\operatorname{Gc}^{\Geom}_{\Sigma},H^{\Caus},S^{\Caus},\Psi^{\Caus},\operatorname{Gc}^{\Caus},L^{\rm int},L^{\rm fib},L^{\rm conc},L^{\rm osc},L^{\rm sc},L^{\rm gc}\right)\right\|_{\rm det},
\qquad a_{gcl}(R):=\left\|\operatorname{Inc}_{A_{gcl}}^{\Sigma}(R)\right\|,
\qquad p_{gcl}(R):=p_{\Phys}(R).
\]
For a finite family \(\mathcal F=(R_1,\ldots,R_N)\) of pairwise
disjoint same-origin stopped windows with exact word \(A_{gcl}\), define the
unnormalised local aggregates
\[
\mathsf A_{gcl}(\mathcal F):=\sum_{j=1}^{N}a_{gcl}(R_j),
\qquad \mathsf D_{gcl}(\mathcal F):=\sum_{j=1}^{N}\delta_{gcl}(R_j),
\qquad \mathsf P_{gcl}(\mathcal F):=\sum_{j=1}^{N}p_{gcl}(R_j).
\]
Then the following check is exact:
\begin{enumerate}[label=\textup{(LGCL\arabic*)}]
\item \(a_{gcl}(R)=0\) exactly when this record contributes no
 target progress through exact support \(A_{gcl}\).  If
\(a_{gcl}(R)>0\), the record itself is the finite public
target-aligned occurrence, with the listed detector tuple and complete joint frontier.
\item \(\mathsf A_{gcl}(\mathcal F),\ \mathsf D_{gcl}(\mathcal F),\ \text{and }\mathsf P_{gcl}(\mathcal F)\) are the exact cumulative target amplitude,
structural detector amplitude, and physical expenditure of this support.
They are never divided by \(N\) or replaced by a limiting average.
\item If \(p_{gcl}(R)>0\), the occurrence is charged to the single
physical carrier.  If \(p_{gcl}(R)=0\), the equality row of the
same joint record is evaluated; a nonzero conditional increment is appended
to the same exact word \(A_{gcl}\), while a zero increment deletes that
record's target-incidence edge.
\end{enumerate}
\end{theorem}
\begin{proof}
The atomic target-aligned relations for \(\Geom\), \(\Caus\), \(\Lift\) are respectively
\cref{prop:ns-geometric-target-relation}, \cref{prop:ns-caus-target-relation}, \cref{prop:ns-lift-target-relation}.  Each is paired with the same public shell covector before any
source projection.  Their finite expression graphs are joined, so every
mixed product, commutator, cutoff, pressure descent, chart motion, and
trace retains the union word \(A_{gcl}\).  Restricting that joint
record to exact word \(A_{gcl}\) gives precisely the displayed
detector \(\delta_{gcl}\) and target amplitude
\(a_{gcl}\); hence its zero/nonzero test is the exact target pairing
of this source combination, not a separate test of its coordinates.
For disjoint stopped windows, linearity of the public weak pairing and
finite additivity of the physical carrier give the three unnormalised sums.
Finally, \cref{thm:actual-source-price-equality-decision} applies to
this same complete word: positive price is charged, while zero price
re-enters the equality row without losing an ancestor.  This proves the
quantitative check for exact support \(A_{gcl}\).
\end{proof}

\begin{theorem}[Local quantitative check: \(\Abs\Lift\)]
\label{thm:quantitative-source-cell-al}
For every finite same-origin stopped record \(R\in\mathfrak D_{A_{al},\Sigma,n}^{\rm pre}(E)\) set
\[
\delta_{al}(R):=\left\|\left(A^{\rm desc},A^{\rm vert},A^{\rm coker},A^{\rm gc},A^{\rm tgt},L^{\rm int},L^{\rm fib},L^{\rm conc},L^{\rm osc},L^{\rm sc},L^{\rm gc}\right)\right\|_{\rm det},
\qquad a_{al}(R):=\left\|\operatorname{Inc}_{A_{al}}^{\Sigma}(R)\right\|,
\qquad p_{al}(R):=p_{\Phys}(R).
\]
For a finite family \(\mathcal F=(R_1,\ldots,R_N)\) of pairwise
disjoint same-origin stopped windows with exact word \(A_{al}\), define the
unnormalised local aggregates
\[
\mathsf A_{al}(\mathcal F):=\sum_{j=1}^{N}a_{al}(R_j),
\qquad \mathsf D_{al}(\mathcal F):=\sum_{j=1}^{N}\delta_{al}(R_j),
\qquad \mathsf P_{al}(\mathcal F):=\sum_{j=1}^{N}p_{al}(R_j).
\]
Then the following check is exact:
\begin{enumerate}[label=\textup{(LAL\arabic*)}]
\item \(a_{al}(R)=0\) exactly when this record contributes no
 target progress through exact support \(A_{al}\).  If
\(a_{al}(R)>0\), the record itself is the finite public
target-aligned occurrence, with the listed detector tuple and complete joint frontier.
\item \(\mathsf A_{al}(\mathcal F),\ \mathsf D_{al}(\mathcal F),\ \text{and }\mathsf P_{al}(\mathcal F)\) are the exact cumulative target amplitude,
structural detector amplitude, and physical expenditure of this support.
They are never divided by \(N\) or replaced by a limiting average.
\item If \(p_{al}(R)>0\), the occurrence is charged to the single
physical carrier.  If \(p_{al}(R)=0\), the equality row of the
same joint record is evaluated; a nonzero conditional increment is appended
to the same exact word \(A_{al}\), while a zero increment deletes that
record's target-incidence edge.
\end{enumerate}
\end{theorem}
\begin{proof}
The atomic target-aligned relations for \(\Abs\), \(\Lift\) are respectively
\cref{prop:ns-abs-target-relation}, \cref{prop:ns-lift-target-relation}.  Each is paired with the same public shell covector before any
source projection.  Their finite expression graphs are joined, so every
mixed product, commutator, cutoff, pressure descent, chart motion, and
trace retains the union word \(A_{al}\).  Restricting that joint
record to exact word \(A_{al}\) gives precisely the displayed
detector \(\delta_{al}\) and target amplitude
\(a_{al}\); hence its zero/nonzero test is the exact target pairing
of this source combination, not a separate test of its coordinates.
For disjoint stopped windows, linearity of the public weak pairing and
finite additivity of the physical carrier give the three unnormalised sums.
Finally, \cref{thm:actual-source-price-equality-decision} applies to
this same complete word: positive price is charged, while zero price
re-enters the equality row without losing an ancestor.  This proves the
quantitative check for exact support \(A_{al}\).
\end{proof}

\begin{theorem}[Local quantitative check: \(\Geom\Abs\Lift\)]
\label{thm:quantitative-source-cell-gal}
For every finite same-origin stopped record \(R\in\mathfrak D_{A_{gal},\Sigma,n}^{\rm pre}(E)\) set
\[
\delta_{gal}(R):=\left\|\left(G^{\Geom}_{\Sigma},\operatorname{Ker}^{\Geom}_{\Sigma},\operatorname{Gc}^{\Geom}_{\Sigma},A^{\rm desc},A^{\rm vert},A^{\rm coker},A^{\rm gc},A^{\rm tgt},L^{\rm int},L^{\rm fib},L^{\rm conc},L^{\rm osc},L^{\rm sc},L^{\rm gc}\right)\right\|_{\rm det},
\qquad a_{gal}(R):=\left\|\operatorname{Inc}_{A_{gal}}^{\Sigma}(R)\right\|,
\qquad p_{gal}(R):=p_{\Phys}(R).
\]
For a finite family \(\mathcal F=(R_1,\ldots,R_N)\) of pairwise
disjoint same-origin stopped windows with exact word \(A_{gal}\), define the
unnormalised local aggregates
\[
\mathsf A_{gal}(\mathcal F):=\sum_{j=1}^{N}a_{gal}(R_j),
\qquad \mathsf D_{gal}(\mathcal F):=\sum_{j=1}^{N}\delta_{gal}(R_j),
\qquad \mathsf P_{gal}(\mathcal F):=\sum_{j=1}^{N}p_{gal}(R_j).
\]
Then the following check is exact:
\begin{enumerate}[label=\textup{(LGAL\arabic*)}]
\item \(a_{gal}(R)=0\) exactly when this record contributes no
 target progress through exact support \(A_{gal}\).  If
\(a_{gal}(R)>0\), the record itself is the finite public
target-aligned occurrence, with the listed detector tuple and complete joint frontier.
\item \(\mathsf A_{gal}(\mathcal F),\ \mathsf D_{gal}(\mathcal F),\ \text{and }\mathsf P_{gal}(\mathcal F)\) are the exact cumulative target amplitude,
structural detector amplitude, and physical expenditure of this support.
They are never divided by \(N\) or replaced by a limiting average.
\item If \(p_{gal}(R)>0\), the occurrence is charged to the single
physical carrier.  If \(p_{gal}(R)=0\), the equality row of the
same joint record is evaluated; a nonzero conditional increment is appended
to the same exact word \(A_{gal}\), while a zero increment deletes that
record's target-incidence edge.
\end{enumerate}
\end{theorem}
\begin{proof}
The atomic target-aligned relations for \(\Geom\), \(\Abs\), \(\Lift\) are respectively
\cref{prop:ns-geometric-target-relation}, \cref{prop:ns-abs-target-relation}, \cref{prop:ns-lift-target-relation}.  Each is paired with the same public shell covector before any
source projection.  Their finite expression graphs are joined, so every
mixed product, commutator, cutoff, pressure descent, chart motion, and
trace retains the union word \(A_{gal}\).  Restricting that joint
record to exact word \(A_{gal}\) gives precisely the displayed
detector \(\delta_{gal}\) and target amplitude
\(a_{gal}\); hence its zero/nonzero test is the exact target pairing
of this source combination, not a separate test of its coordinates.
For disjoint stopped windows, linearity of the public weak pairing and
finite additivity of the physical carrier give the three unnormalised sums.
Finally, \cref{thm:actual-source-price-equality-decision} applies to
this same complete word: positive price is charged, while zero price
re-enters the equality row without losing an ancestor.  This proves the
quantitative check for exact support \(A_{gal}\).
\end{proof}

\begin{theorem}[Local quantitative check: \(\Caus\Abs\Lift\)]
\label{thm:quantitative-source-cell-cal}
For every finite same-origin stopped record \(R\in\mathfrak D_{A_{cal},\Sigma,n}^{\rm pre}(E)\) set
\[
\delta_{cal}(R):=\left\|\left(H^{\Caus},S^{\Caus},\Psi^{\Caus},\operatorname{Gc}^{\Caus},A^{\rm desc},A^{\rm vert},A^{\rm coker},A^{\rm gc},A^{\rm tgt},L^{\rm int},L^{\rm fib},L^{\rm conc},L^{\rm osc},L^{\rm sc},L^{\rm gc}\right)\right\|_{\rm det},
\qquad a_{cal}(R):=\left\|\operatorname{Inc}_{A_{cal}}^{\Sigma}(R)\right\|,
\qquad p_{cal}(R):=p_{\Phys}(R).
\]
For a finite family \(\mathcal F=(R_1,\ldots,R_N)\) of pairwise
disjoint same-origin stopped windows with exact word \(A_{cal}\), define the
unnormalised local aggregates
\[
\mathsf A_{cal}(\mathcal F):=\sum_{j=1}^{N}a_{cal}(R_j),
\qquad \mathsf D_{cal}(\mathcal F):=\sum_{j=1}^{N}\delta_{cal}(R_j),
\qquad \mathsf P_{cal}(\mathcal F):=\sum_{j=1}^{N}p_{cal}(R_j).
\]
Then the following check is exact:
\begin{enumerate}[label=\textup{(LCAL\arabic*)}]
\item \(a_{cal}(R)=0\) exactly when this record contributes no
 target progress through exact support \(A_{cal}\).  If
\(a_{cal}(R)>0\), the record itself is the finite public
target-aligned occurrence, with the listed detector tuple and complete joint frontier.
\item \(\mathsf A_{cal}(\mathcal F),\ \mathsf D_{cal}(\mathcal F),\ \text{and }\mathsf P_{cal}(\mathcal F)\) are the exact cumulative target amplitude,
structural detector amplitude, and physical expenditure of this support.
They are never divided by \(N\) or replaced by a limiting average.
\item If \(p_{cal}(R)>0\), the occurrence is charged to the single
physical carrier.  If \(p_{cal}(R)=0\), the equality row of the
same joint record is evaluated; a nonzero conditional increment is appended
to the same exact word \(A_{cal}\), while a zero increment deletes that
record's target-incidence edge.
\end{enumerate}
\end{theorem}
\begin{proof}
The atomic target-aligned relations for \(\Caus\), \(\Abs\), \(\Lift\) are respectively
\cref{prop:ns-caus-target-relation}, \cref{prop:ns-abs-target-relation}, \cref{prop:ns-lift-target-relation}.  Each is paired with the same public shell covector before any
source projection.  Their finite expression graphs are joined, so every
mixed product, commutator, cutoff, pressure descent, chart motion, and
trace retains the union word \(A_{cal}\).  Restricting that joint
record to exact word \(A_{cal}\) gives precisely the displayed
detector \(\delta_{cal}\) and target amplitude
\(a_{cal}\); hence its zero/nonzero test is the exact target pairing
of this source combination, not a separate test of its coordinates.
For disjoint stopped windows, linearity of the public weak pairing and
finite additivity of the physical carrier give the three unnormalised sums.
Finally, \cref{thm:actual-source-price-equality-decision} applies to
this same complete word: positive price is charged, while zero price
re-enters the equality row without losing an ancestor.  This proves the
quantitative check for exact support \(A_{cal}\).
\end{proof}

\begin{theorem}[Local quantitative check: \(\Geom\Caus\Abs\Lift\)]
\label{thm:quantitative-source-cell-gcal}
For every finite same-origin stopped record \(R\in\mathfrak D_{A_{gcal},\Sigma,n}^{\rm pre}(E)\) set
\[
\delta_{gcal}(R):=\left\|\left(G^{\Geom}_{\Sigma},\operatorname{Ker}^{\Geom}_{\Sigma},\operatorname{Gc}^{\Geom}_{\Sigma},H^{\Caus},S^{\Caus},\Psi^{\Caus},\operatorname{Gc}^{\Caus},A^{\rm desc},A^{\rm vert},A^{\rm coker},A^{\rm gc},A^{\rm tgt},L^{\rm int},L^{\rm fib},L^{\rm conc},L^{\rm osc},L^{\rm sc},L^{\rm gc}\right)\right\|_{\rm det},
\qquad a_{gcal}(R):=\left\|\operatorname{Inc}_{A_{gcal}}^{\Sigma}(R)\right\|,
\qquad p_{gcal}(R):=p_{\Phys}(R).
\]
For a finite family \(\mathcal F=(R_1,\ldots,R_N)\) of pairwise
disjoint same-origin stopped windows with exact word \(A_{gcal}\), define the
unnormalised local aggregates
\[
\mathsf A_{gcal}(\mathcal F):=\sum_{j=1}^{N}a_{gcal}(R_j),
\qquad \mathsf D_{gcal}(\mathcal F):=\sum_{j=1}^{N}\delta_{gcal}(R_j),
\qquad \mathsf P_{gcal}(\mathcal F):=\sum_{j=1}^{N}p_{gcal}(R_j).
\]
Then the following check is exact:
\begin{enumerate}[label=\textup{(LGCAL\arabic*)}]
\item \(a_{gcal}(R)=0\) exactly when this record contributes no
 target progress through exact support \(A_{gcal}\).  If
\(a_{gcal}(R)>0\), the record itself is the finite public
target-aligned occurrence, with the listed detector tuple and complete joint frontier.
\item \(\mathsf A_{gcal}(\mathcal F),\ \mathsf D_{gcal}(\mathcal F),\ \text{and }\mathsf P_{gcal}(\mathcal F)\) are the exact cumulative target amplitude,
structural detector amplitude, and physical expenditure of this support.
They are never divided by \(N\) or replaced by a limiting average.
\item If \(p_{gcal}(R)>0\), the occurrence is charged to the single
physical carrier.  If \(p_{gcal}(R)=0\), the equality row of the
same joint record is evaluated; a nonzero conditional increment is appended
to the same exact word \(A_{gcal}\), while a zero increment deletes that
record's target-incidence edge.
\end{enumerate}
\end{theorem}
\begin{proof}
The atomic target-aligned relations for \(\Geom\), \(\Caus\), \(\Abs\), \(\Lift\) are respectively
\cref{prop:ns-geometric-target-relation}, \cref{prop:ns-caus-target-relation}, \cref{prop:ns-abs-target-relation}, \cref{prop:ns-lift-target-relation}.  Each is paired with the same public shell covector before any
source projection.  Their finite expression graphs are joined, so every
mixed product, commutator, cutoff, pressure descent, chart motion, and
trace retains the union word \(A_{gcal}\).  Restricting that joint
record to exact word \(A_{gcal}\) gives precisely the displayed
detector \(\delta_{gcal}\) and target amplitude
\(a_{gcal}\); hence its zero/nonzero test is the exact target pairing
of this source combination, not a separate test of its coordinates.
For disjoint stopped windows, linearity of the public weak pairing and
finite additivity of the physical carrier give the three unnormalised sums.
Finally, \cref{thm:actual-source-price-equality-decision} applies to
this same complete word: positive price is charged, while zero price
re-enters the equality row without losing an ancestor.  This proves the
quantitative check for exact support \(A_{gcal}\).
\end{proof}

\begin{theorem}[Local quantitative check: \(\Bdry\)]
\label{thm:quantitative-source-cell-b}
For every finite same-origin stopped record \(R\in\mathfrak D_{A_{b},\Sigma,n}^{\rm pre}(E)\) set
\[
\delta_{b}(R):=\left\|\left(B^{\rm ext},B^{\rm tail},B^{\rm term},B^{\rm gc}\right)\right\|_{\rm det},
\qquad a_{b}(R):=\left\|\operatorname{Inc}_{A_{b}}^{\Sigma}(R)\right\|,
\qquad p_{b}(R):=p_{\Phys}(R).
\]
For a finite family \(\mathcal F=(R_1,\ldots,R_N)\) of pairwise
disjoint same-origin stopped windows with exact word \(A_{b}\), define the
unnormalised local aggregates
\[
\mathsf A_{b}(\mathcal F):=\sum_{j=1}^{N}a_{b}(R_j),
\qquad \mathsf D_{b}(\mathcal F):=\sum_{j=1}^{N}\delta_{b}(R_j),
\qquad \mathsf P_{b}(\mathcal F):=\sum_{j=1}^{N}p_{b}(R_j).
\]
Then the following check is exact:
\begin{enumerate}[label=\textup{(LB\arabic*)}]
\item \(a_{b}(R)=0\) exactly when this record contributes no
 target progress through exact support \(A_{b}\).  If
\(a_{b}(R)>0\), the record itself is the finite public
target-aligned occurrence, with the listed detector tuple and complete joint frontier.
\item \(\mathsf A_{b}(\mathcal F),\ \mathsf D_{b}(\mathcal F),\ \text{and }\mathsf P_{b}(\mathcal F)\) are the exact cumulative target amplitude,
structural detector amplitude, and physical expenditure of this support.
They are never divided by \(N\) or replaced by a limiting average.
\item If \(p_{b}(R)>0\), the occurrence is charged to the single
physical carrier.  If \(p_{b}(R)=0\), the equality row of the
same joint record is evaluated; a nonzero conditional increment is appended
to the same exact word \(A_{b}\), while a zero increment deletes that
record's target-incidence edge.
\end{enumerate}
\end{theorem}
\begin{proof}
The atomic target-aligned relations for \(\Bdry\) are respectively
\cref{prop:ns-bdry-target-relation}.  Each is paired with the same public shell covector before any
source projection.  Their finite expression graphs are joined, so every
mixed product, commutator, cutoff, pressure descent, chart motion, and
trace retains the union word \(A_{b}\).  Restricting that joint
record to exact word \(A_{b}\) gives precisely the displayed
detector \(\delta_{b}\) and target amplitude
\(a_{b}\); hence its zero/nonzero test is the exact target pairing
of this source combination, not a separate test of its coordinates.
For disjoint stopped windows, linearity of the public weak pairing and
finite additivity of the physical carrier give the three unnormalised sums.
Finally, \cref{thm:actual-source-price-equality-decision} applies to
this same complete word: positive price is charged, while zero price
re-enters the equality row without losing an ancestor.  This proves the
quantitative check for exact support \(A_{b}\).
\end{proof}

\begin{theorem}[Local quantitative check: \(\Geom\Bdry\)]
\label{thm:quantitative-source-cell-gb}
For every finite same-origin stopped record \(R\in\mathfrak D_{A_{gb},\Sigma,n}^{\rm pre}(E)\) set
\[
\delta_{gb}(R):=\left\|\left(G^{\Geom}_{\Sigma},\operatorname{Ker}^{\Geom}_{\Sigma},\operatorname{Gc}^{\Geom}_{\Sigma},B^{\rm ext},B^{\rm tail},B^{\rm term},B^{\rm gc}\right)\right\|_{\rm det},
\qquad a_{gb}(R):=\left\|\operatorname{Inc}_{A_{gb}}^{\Sigma}(R)\right\|,
\qquad p_{gb}(R):=p_{\Phys}(R).
\]
For a finite family \(\mathcal F=(R_1,\ldots,R_N)\) of pairwise
disjoint same-origin stopped windows with exact word \(A_{gb}\), define the
unnormalised local aggregates
\[
\mathsf A_{gb}(\mathcal F):=\sum_{j=1}^{N}a_{gb}(R_j),
\qquad \mathsf D_{gb}(\mathcal F):=\sum_{j=1}^{N}\delta_{gb}(R_j),
\qquad \mathsf P_{gb}(\mathcal F):=\sum_{j=1}^{N}p_{gb}(R_j).
\]
Then the following check is exact:
\begin{enumerate}[label=\textup{(LGB\arabic*)}]
\item \(a_{gb}(R)=0\) exactly when this record contributes no
 target progress through exact support \(A_{gb}\).  If
\(a_{gb}(R)>0\), the record itself is the finite public
target-aligned occurrence, with the listed detector tuple and complete joint frontier.
\item \(\mathsf A_{gb}(\mathcal F),\ \mathsf D_{gb}(\mathcal F),\ \text{and }\mathsf P_{gb}(\mathcal F)\) are the exact cumulative target amplitude,
structural detector amplitude, and physical expenditure of this support.
They are never divided by \(N\) or replaced by a limiting average.
\item If \(p_{gb}(R)>0\), the occurrence is charged to the single
physical carrier.  If \(p_{gb}(R)=0\), the equality row of the
same joint record is evaluated; a nonzero conditional increment is appended
to the same exact word \(A_{gb}\), while a zero increment deletes that
record's target-incidence edge.
\end{enumerate}
\end{theorem}
\begin{proof}
The atomic target-aligned relations for \(\Geom\), \(\Bdry\) are respectively
\cref{prop:ns-geometric-target-relation}, \cref{prop:ns-bdry-target-relation}.  Each is paired with the same public shell covector before any
source projection.  Their finite expression graphs are joined, so every
mixed product, commutator, cutoff, pressure descent, chart motion, and
trace retains the union word \(A_{gb}\).  Restricting that joint
record to exact word \(A_{gb}\) gives precisely the displayed
detector \(\delta_{gb}\) and target amplitude
\(a_{gb}\); hence its zero/nonzero test is the exact target pairing
of this source combination, not a separate test of its coordinates.
For disjoint stopped windows, linearity of the public weak pairing and
finite additivity of the physical carrier give the three unnormalised sums.
Finally, \cref{thm:actual-source-price-equality-decision} applies to
this same complete word: positive price is charged, while zero price
re-enters the equality row without losing an ancestor.  This proves the
quantitative check for exact support \(A_{gb}\).
\end{proof}

\begin{theorem}[Local quantitative check: \(\Caus\Bdry\)]
\label{thm:quantitative-source-cell-cb}
For every finite same-origin stopped record \(R\in\mathfrak D_{A_{cb},\Sigma,n}^{\rm pre}(E)\) set
\[
\delta_{cb}(R):=\left\|\left(H^{\Caus},S^{\Caus},\Psi^{\Caus},\operatorname{Gc}^{\Caus},B^{\rm ext},B^{\rm tail},B^{\rm term},B^{\rm gc}\right)\right\|_{\rm det},
\qquad a_{cb}(R):=\left\|\operatorname{Inc}_{A_{cb}}^{\Sigma}(R)\right\|,
\qquad p_{cb}(R):=p_{\Phys}(R).
\]
For a finite family \(\mathcal F=(R_1,\ldots,R_N)\) of pairwise
disjoint same-origin stopped windows with exact word \(A_{cb}\), define the
unnormalised local aggregates
\[
\mathsf A_{cb}(\mathcal F):=\sum_{j=1}^{N}a_{cb}(R_j),
\qquad \mathsf D_{cb}(\mathcal F):=\sum_{j=1}^{N}\delta_{cb}(R_j),
\qquad \mathsf P_{cb}(\mathcal F):=\sum_{j=1}^{N}p_{cb}(R_j).
\]
Then the following check is exact:
\begin{enumerate}[label=\textup{(LCB\arabic*)}]
\item \(a_{cb}(R)=0\) exactly when this record contributes no
 target progress through exact support \(A_{cb}\).  If
\(a_{cb}(R)>0\), the record itself is the finite public
target-aligned occurrence, with the listed detector tuple and complete joint frontier.
\item \(\mathsf A_{cb}(\mathcal F),\ \mathsf D_{cb}(\mathcal F),\ \text{and }\mathsf P_{cb}(\mathcal F)\) are the exact cumulative target amplitude,
structural detector amplitude, and physical expenditure of this support.
They are never divided by \(N\) or replaced by a limiting average.
\item If \(p_{cb}(R)>0\), the occurrence is charged to the single
physical carrier.  If \(p_{cb}(R)=0\), the equality row of the
same joint record is evaluated; a nonzero conditional increment is appended
to the same exact word \(A_{cb}\), while a zero increment deletes that
record's target-incidence edge.
\end{enumerate}
\end{theorem}
\begin{proof}
The atomic target-aligned relations for \(\Caus\), \(\Bdry\) are respectively
\cref{prop:ns-caus-target-relation}, \cref{prop:ns-bdry-target-relation}.  Each is paired with the same public shell covector before any
source projection.  Their finite expression graphs are joined, so every
mixed product, commutator, cutoff, pressure descent, chart motion, and
trace retains the union word \(A_{cb}\).  Restricting that joint
record to exact word \(A_{cb}\) gives precisely the displayed
detector \(\delta_{cb}\) and target amplitude
\(a_{cb}\); hence its zero/nonzero test is the exact target pairing
of this source combination, not a separate test of its coordinates.
For disjoint stopped windows, linearity of the public weak pairing and
finite additivity of the physical carrier give the three unnormalised sums.
Finally, \cref{thm:actual-source-price-equality-decision} applies to
this same complete word: positive price is charged, while zero price
re-enters the equality row without losing an ancestor.  This proves the
quantitative check for exact support \(A_{cb}\).
\end{proof}

\begin{theorem}[Local quantitative check: \(\Geom\Caus\Bdry\)]
\label{thm:quantitative-source-cell-gcb}
For every finite same-origin stopped record \(R\in\mathfrak D_{A_{gcb},\Sigma,n}^{\rm pre}(E)\) set
\[
\delta_{gcb}(R):=\left\|\left(G^{\Geom}_{\Sigma},\operatorname{Ker}^{\Geom}_{\Sigma},\operatorname{Gc}^{\Geom}_{\Sigma},H^{\Caus},S^{\Caus},\Psi^{\Caus},\operatorname{Gc}^{\Caus},B^{\rm ext},B^{\rm tail},B^{\rm term},B^{\rm gc}\right)\right\|_{\rm det},
\qquad a_{gcb}(R):=\left\|\operatorname{Inc}_{A_{gcb}}^{\Sigma}(R)\right\|,
\qquad p_{gcb}(R):=p_{\Phys}(R).
\]
For a finite family \(\mathcal F=(R_1,\ldots,R_N)\) of pairwise
disjoint same-origin stopped windows with exact word \(A_{gcb}\), define the
unnormalised local aggregates
\[
\mathsf A_{gcb}(\mathcal F):=\sum_{j=1}^{N}a_{gcb}(R_j),
\qquad \mathsf D_{gcb}(\mathcal F):=\sum_{j=1}^{N}\delta_{gcb}(R_j),
\qquad \mathsf P_{gcb}(\mathcal F):=\sum_{j=1}^{N}p_{gcb}(R_j).
\]
Then the following check is exact:
\begin{enumerate}[label=\textup{(LGCB\arabic*)}]
\item \(a_{gcb}(R)=0\) exactly when this record contributes no
 target progress through exact support \(A_{gcb}\).  If
\(a_{gcb}(R)>0\), the record itself is the finite public
target-aligned occurrence, with the listed detector tuple and complete joint frontier.
\item \(\mathsf A_{gcb}(\mathcal F),\ \mathsf D_{gcb}(\mathcal F),\ \text{and }\mathsf P_{gcb}(\mathcal F)\) are the exact cumulative target amplitude,
structural detector amplitude, and physical expenditure of this support.
They are never divided by \(N\) or replaced by a limiting average.
\item If \(p_{gcb}(R)>0\), the occurrence is charged to the single
physical carrier.  If \(p_{gcb}(R)=0\), the equality row of the
same joint record is evaluated; a nonzero conditional increment is appended
to the same exact word \(A_{gcb}\), while a zero increment deletes that
record's target-incidence edge.
\end{enumerate}
\end{theorem}
\begin{proof}
The atomic target-aligned relations for \(\Geom\), \(\Caus\), \(\Bdry\) are respectively
\cref{prop:ns-geometric-target-relation}, \cref{prop:ns-caus-target-relation}, \cref{prop:ns-bdry-target-relation}.  Each is paired with the same public shell covector before any
source projection.  Their finite expression graphs are joined, so every
mixed product, commutator, cutoff, pressure descent, chart motion, and
trace retains the union word \(A_{gcb}\).  Restricting that joint
record to exact word \(A_{gcb}\) gives precisely the displayed
detector \(\delta_{gcb}\) and target amplitude
\(a_{gcb}\); hence its zero/nonzero test is the exact target pairing
of this source combination, not a separate test of its coordinates.
For disjoint stopped windows, linearity of the public weak pairing and
finite additivity of the physical carrier give the three unnormalised sums.
Finally, \cref{thm:actual-source-price-equality-decision} applies to
this same complete word: positive price is charged, while zero price
re-enters the equality row without losing an ancestor.  This proves the
quantitative check for exact support \(A_{gcb}\).
\end{proof}

\begin{theorem}[Local quantitative check: \(\Abs\Bdry\)]
\label{thm:quantitative-source-cell-ab}
For every finite same-origin stopped record \(R\in\mathfrak D_{A_{ab},\Sigma,n}^{\rm pre}(E)\) set
\[
\delta_{ab}(R):=\left\|\left(A^{\rm desc},A^{\rm vert},A^{\rm coker},A^{\rm gc},A^{\rm tgt},B^{\rm ext},B^{\rm tail},B^{\rm term},B^{\rm gc}\right)\right\|_{\rm det},
\qquad a_{ab}(R):=\left\|\operatorname{Inc}_{A_{ab}}^{\Sigma}(R)\right\|,
\qquad p_{ab}(R):=p_{\Phys}(R).
\]
For a finite family \(\mathcal F=(R_1,\ldots,R_N)\) of pairwise
disjoint same-origin stopped windows with exact word \(A_{ab}\), define the
unnormalised local aggregates
\[
\mathsf A_{ab}(\mathcal F):=\sum_{j=1}^{N}a_{ab}(R_j),
\qquad \mathsf D_{ab}(\mathcal F):=\sum_{j=1}^{N}\delta_{ab}(R_j),
\qquad \mathsf P_{ab}(\mathcal F):=\sum_{j=1}^{N}p_{ab}(R_j).
\]
Then the following check is exact:
\begin{enumerate}[label=\textup{(LAB\arabic*)}]
\item \(a_{ab}(R)=0\) exactly when this record contributes no
 target progress through exact support \(A_{ab}\).  If
\(a_{ab}(R)>0\), the record itself is the finite public
target-aligned occurrence, with the listed detector tuple and complete joint frontier.
\item \(\mathsf A_{ab}(\mathcal F),\ \mathsf D_{ab}(\mathcal F),\ \text{and }\mathsf P_{ab}(\mathcal F)\) are the exact cumulative target amplitude,
structural detector amplitude, and physical expenditure of this support.
They are never divided by \(N\) or replaced by a limiting average.
\item If \(p_{ab}(R)>0\), the occurrence is charged to the single
physical carrier.  If \(p_{ab}(R)=0\), the equality row of the
same joint record is evaluated; a nonzero conditional increment is appended
to the same exact word \(A_{ab}\), while a zero increment deletes that
record's target-incidence edge.
\end{enumerate}
\end{theorem}
\begin{proof}
The atomic target-aligned relations for \(\Abs\), \(\Bdry\) are respectively
\cref{prop:ns-abs-target-relation}, \cref{prop:ns-bdry-target-relation}.  Each is paired with the same public shell covector before any
source projection.  Their finite expression graphs are joined, so every
mixed product, commutator, cutoff, pressure descent, chart motion, and
trace retains the union word \(A_{ab}\).  Restricting that joint
record to exact word \(A_{ab}\) gives precisely the displayed
detector \(\delta_{ab}\) and target amplitude
\(a_{ab}\); hence its zero/nonzero test is the exact target pairing
of this source combination, not a separate test of its coordinates.
For disjoint stopped windows, linearity of the public weak pairing and
finite additivity of the physical carrier give the three unnormalised sums.
Finally, \cref{thm:actual-source-price-equality-decision} applies to
this same complete word: positive price is charged, while zero price
re-enters the equality row without losing an ancestor.  This proves the
quantitative check for exact support \(A_{ab}\).
\end{proof}

\begin{theorem}[Local quantitative check: \(\Geom\Abs\Bdry\)]
\label{thm:quantitative-source-cell-gab}
For every finite same-origin stopped record \(R\in\mathfrak D_{A_{gab},\Sigma,n}^{\rm pre}(E)\) set
\[
\delta_{gab}(R):=\left\|\left(G^{\Geom}_{\Sigma},\operatorname{Ker}^{\Geom}_{\Sigma},\operatorname{Gc}^{\Geom}_{\Sigma},A^{\rm desc},A^{\rm vert},A^{\rm coker},A^{\rm gc},A^{\rm tgt},B^{\rm ext},B^{\rm tail},B^{\rm term},B^{\rm gc}\right)\right\|_{\rm det},
\qquad a_{gab}(R):=\left\|\operatorname{Inc}_{A_{gab}}^{\Sigma}(R)\right\|,
\qquad p_{gab}(R):=p_{\Phys}(R).
\]
For a finite family \(\mathcal F=(R_1,\ldots,R_N)\) of pairwise
disjoint same-origin stopped windows with exact word \(A_{gab}\), define the
unnormalised local aggregates
\[
\mathsf A_{gab}(\mathcal F):=\sum_{j=1}^{N}a_{gab}(R_j),
\qquad \mathsf D_{gab}(\mathcal F):=\sum_{j=1}^{N}\delta_{gab}(R_j),
\qquad \mathsf P_{gab}(\mathcal F):=\sum_{j=1}^{N}p_{gab}(R_j).
\]
Then the following check is exact:
\begin{enumerate}[label=\textup{(LGAB\arabic*)}]
\item \(a_{gab}(R)=0\) exactly when this record contributes no
 target progress through exact support \(A_{gab}\).  If
\(a_{gab}(R)>0\), the record itself is the finite public
target-aligned occurrence, with the listed detector tuple and complete joint frontier.
\item \(\mathsf A_{gab}(\mathcal F),\ \mathsf D_{gab}(\mathcal F),\ \text{and }\mathsf P_{gab}(\mathcal F)\) are the exact cumulative target amplitude,
structural detector amplitude, and physical expenditure of this support.
They are never divided by \(N\) or replaced by a limiting average.
\item If \(p_{gab}(R)>0\), the occurrence is charged to the single
physical carrier.  If \(p_{gab}(R)=0\), the equality row of the
same joint record is evaluated; a nonzero conditional increment is appended
to the same exact word \(A_{gab}\), while a zero increment deletes that
record's target-incidence edge.
\end{enumerate}
\end{theorem}
\begin{proof}
The atomic target-aligned relations for \(\Geom\), \(\Abs\), \(\Bdry\) are respectively
\cref{prop:ns-geometric-target-relation}, \cref{prop:ns-abs-target-relation}, \cref{prop:ns-bdry-target-relation}.  Each is paired with the same public shell covector before any
source projection.  Their finite expression graphs are joined, so every
mixed product, commutator, cutoff, pressure descent, chart motion, and
trace retains the union word \(A_{gab}\).  Restricting that joint
record to exact word \(A_{gab}\) gives precisely the displayed
detector \(\delta_{gab}\) and target amplitude
\(a_{gab}\); hence its zero/nonzero test is the exact target pairing
of this source combination, not a separate test of its coordinates.
For disjoint stopped windows, linearity of the public weak pairing and
finite additivity of the physical carrier give the three unnormalised sums.
Finally, \cref{thm:actual-source-price-equality-decision} applies to
this same complete word: positive price is charged, while zero price
re-enters the equality row without losing an ancestor.  This proves the
quantitative check for exact support \(A_{gab}\).
\end{proof}

\begin{theorem}[Local quantitative check: \(\Caus\Abs\Bdry\)]
\label{thm:quantitative-source-cell-cab}
For every finite same-origin stopped record \(R\in\mathfrak D_{A_{cab},\Sigma,n}^{\rm pre}(E)\) set
\[
\delta_{cab}(R):=\left\|\left(H^{\Caus},S^{\Caus},\Psi^{\Caus},\operatorname{Gc}^{\Caus},A^{\rm desc},A^{\rm vert},A^{\rm coker},A^{\rm gc},A^{\rm tgt},B^{\rm ext},B^{\rm tail},B^{\rm term},B^{\rm gc}\right)\right\|_{\rm det},
\qquad a_{cab}(R):=\left\|\operatorname{Inc}_{A_{cab}}^{\Sigma}(R)\right\|,
\qquad p_{cab}(R):=p_{\Phys}(R).
\]
For a finite family \(\mathcal F=(R_1,\ldots,R_N)\) of pairwise
disjoint same-origin stopped windows with exact word \(A_{cab}\), define the
unnormalised local aggregates
\[
\mathsf A_{cab}(\mathcal F):=\sum_{j=1}^{N}a_{cab}(R_j),
\qquad \mathsf D_{cab}(\mathcal F):=\sum_{j=1}^{N}\delta_{cab}(R_j),
\qquad \mathsf P_{cab}(\mathcal F):=\sum_{j=1}^{N}p_{cab}(R_j).
\]
Then the following check is exact:
\begin{enumerate}[label=\textup{(LCAB\arabic*)}]
\item \(a_{cab}(R)=0\) exactly when this record contributes no
 target progress through exact support \(A_{cab}\).  If
\(a_{cab}(R)>0\), the record itself is the finite public
target-aligned occurrence, with the listed detector tuple and complete joint frontier.
\item \(\mathsf A_{cab}(\mathcal F),\ \mathsf D_{cab}(\mathcal F),\ \text{and }\mathsf P_{cab}(\mathcal F)\) are the exact cumulative target amplitude,
structural detector amplitude, and physical expenditure of this support.
They are never divided by \(N\) or replaced by a limiting average.
\item If \(p_{cab}(R)>0\), the occurrence is charged to the single
physical carrier.  If \(p_{cab}(R)=0\), the equality row of the
same joint record is evaluated; a nonzero conditional increment is appended
to the same exact word \(A_{cab}\), while a zero increment deletes that
record's target-incidence edge.
\end{enumerate}
\end{theorem}
\begin{proof}
The atomic target-aligned relations for \(\Caus\), \(\Abs\), \(\Bdry\) are respectively
\cref{prop:ns-caus-target-relation}, \cref{prop:ns-abs-target-relation}, \cref{prop:ns-bdry-target-relation}.  Each is paired with the same public shell covector before any
source projection.  Their finite expression graphs are joined, so every
mixed product, commutator, cutoff, pressure descent, chart motion, and
trace retains the union word \(A_{cab}\).  Restricting that joint
record to exact word \(A_{cab}\) gives precisely the displayed
detector \(\delta_{cab}\) and target amplitude
\(a_{cab}\); hence its zero/nonzero test is the exact target pairing
of this source combination, not a separate test of its coordinates.
For disjoint stopped windows, linearity of the public weak pairing and
finite additivity of the physical carrier give the three unnormalised sums.
Finally, \cref{thm:actual-source-price-equality-decision} applies to
this same complete word: positive price is charged, while zero price
re-enters the equality row without losing an ancestor.  This proves the
quantitative check for exact support \(A_{cab}\).
\end{proof}

\begin{theorem}[Local quantitative check: \(\Geom\Caus\Abs\Bdry\)]
\label{thm:quantitative-source-cell-gcab}
For every finite same-origin stopped record \(R\in\mathfrak D_{A_{gcab},\Sigma,n}^{\rm pre}(E)\) set
\[
\delta_{gcab}(R):=\left\|\left(G^{\Geom}_{\Sigma},\operatorname{Ker}^{\Geom}_{\Sigma},\operatorname{Gc}^{\Geom}_{\Sigma},H^{\Caus},S^{\Caus},\Psi^{\Caus},\operatorname{Gc}^{\Caus},A^{\rm desc},A^{\rm vert},A^{\rm coker},A^{\rm gc},A^{\rm tgt},B^{\rm ext},B^{\rm tail},B^{\rm term},B^{\rm gc}\right)\right\|_{\rm det},
\qquad a_{gcab}(R):=\left\|\operatorname{Inc}_{A_{gcab}}^{\Sigma}(R)\right\|,
\qquad p_{gcab}(R):=p_{\Phys}(R).
\]
For a finite family \(\mathcal F=(R_1,\ldots,R_N)\) of pairwise
disjoint same-origin stopped windows with exact word \(A_{gcab}\), define the
unnormalised local aggregates
\[
\mathsf A_{gcab}(\mathcal F):=\sum_{j=1}^{N}a_{gcab}(R_j),
\qquad \mathsf D_{gcab}(\mathcal F):=\sum_{j=1}^{N}\delta_{gcab}(R_j),
\qquad \mathsf P_{gcab}(\mathcal F):=\sum_{j=1}^{N}p_{gcab}(R_j).
\]
Then the following check is exact:
\begin{enumerate}[label=\textup{(LGCAB\arabic*)}]
\item \(a_{gcab}(R)=0\) exactly when this record contributes no
 target progress through exact support \(A_{gcab}\).  If
\(a_{gcab}(R)>0\), the record itself is the finite public
target-aligned occurrence, with the listed detector tuple and complete joint frontier.
\item \(\mathsf A_{gcab}(\mathcal F),\ \mathsf D_{gcab}(\mathcal F),\ \text{and }\mathsf P_{gcab}(\mathcal F)\) are the exact cumulative target amplitude,
structural detector amplitude, and physical expenditure of this support.
They are never divided by \(N\) or replaced by a limiting average.
\item If \(p_{gcab}(R)>0\), the occurrence is charged to the single
physical carrier.  If \(p_{gcab}(R)=0\), the equality row of the
same joint record is evaluated; a nonzero conditional increment is appended
to the same exact word \(A_{gcab}\), while a zero increment deletes that
record's target-incidence edge.
\end{enumerate}
\end{theorem}
\begin{proof}
The atomic target-aligned relations for \(\Geom\), \(\Caus\), \(\Abs\), \(\Bdry\) are respectively
\cref{prop:ns-geometric-target-relation}, \cref{prop:ns-caus-target-relation}, \cref{prop:ns-abs-target-relation}, \cref{prop:ns-bdry-target-relation}.  Each is paired with the same public shell covector before any
source projection.  Their finite expression graphs are joined, so every
mixed product, commutator, cutoff, pressure descent, chart motion, and
trace retains the union word \(A_{gcab}\).  Restricting that joint
record to exact word \(A_{gcab}\) gives precisely the displayed
detector \(\delta_{gcab}\) and target amplitude
\(a_{gcab}\); hence its zero/nonzero test is the exact target pairing
of this source combination, not a separate test of its coordinates.
For disjoint stopped windows, linearity of the public weak pairing and
finite additivity of the physical carrier give the three unnormalised sums.
Finally, \cref{thm:actual-source-price-equality-decision} applies to
this same complete word: positive price is charged, while zero price
re-enters the equality row without losing an ancestor.  This proves the
quantitative check for exact support \(A_{gcab}\).
\end{proof}

\begin{theorem}[Local quantitative check: \(\Lift\Bdry\)]
\label{thm:quantitative-source-cell-lb}
For every finite same-origin stopped record \(R\in\mathfrak D_{A_{lb},\Sigma,n}^{\rm pre}(E)\) set
\[
\delta_{lb}(R):=\left\|\left(L^{\rm int},L^{\rm fib},L^{\rm conc},L^{\rm osc},L^{\rm sc},L^{\rm gc},B^{\rm ext},B^{\rm tail},B^{\rm term},B^{\rm gc}\right)\right\|_{\rm det},
\qquad a_{lb}(R):=\left\|\operatorname{Inc}_{A_{lb}}^{\Sigma}(R)\right\|,
\qquad p_{lb}(R):=p_{\Phys}(R).
\]
For a finite family \(\mathcal F=(R_1,\ldots,R_N)\) of pairwise
disjoint same-origin stopped windows with exact word \(A_{lb}\), define the
unnormalised local aggregates
\[
\mathsf A_{lb}(\mathcal F):=\sum_{j=1}^{N}a_{lb}(R_j),
\qquad \mathsf D_{lb}(\mathcal F):=\sum_{j=1}^{N}\delta_{lb}(R_j),
\qquad \mathsf P_{lb}(\mathcal F):=\sum_{j=1}^{N}p_{lb}(R_j).
\]
Then the following check is exact:
\begin{enumerate}[label=\textup{(LLB\arabic*)}]
\item \(a_{lb}(R)=0\) exactly when this record contributes no
 target progress through exact support \(A_{lb}\).  If
\(a_{lb}(R)>0\), the record itself is the finite public
target-aligned occurrence, with the listed detector tuple and complete joint frontier.
\item \(\mathsf A_{lb}(\mathcal F),\ \mathsf D_{lb}(\mathcal F),\ \text{and }\mathsf P_{lb}(\mathcal F)\) are the exact cumulative target amplitude,
structural detector amplitude, and physical expenditure of this support.
They are never divided by \(N\) or replaced by a limiting average.
\item If \(p_{lb}(R)>0\), the occurrence is charged to the single
physical carrier.  If \(p_{lb}(R)=0\), the equality row of the
same joint record is evaluated; a nonzero conditional increment is appended
to the same exact word \(A_{lb}\), while a zero increment deletes that
record's target-incidence edge.
\end{enumerate}
\end{theorem}
\begin{proof}
The atomic target-aligned relations for \(\Lift\), \(\Bdry\) are respectively
\cref{prop:ns-lift-target-relation}, \cref{prop:ns-bdry-target-relation}.  Each is paired with the same public shell covector before any
source projection.  Their finite expression graphs are joined, so every
mixed product, commutator, cutoff, pressure descent, chart motion, and
trace retains the union word \(A_{lb}\).  Restricting that joint
record to exact word \(A_{lb}\) gives precisely the displayed
detector \(\delta_{lb}\) and target amplitude
\(a_{lb}\); hence its zero/nonzero test is the exact target pairing
of this source combination, not a separate test of its coordinates.
For disjoint stopped windows, linearity of the public weak pairing and
finite additivity of the physical carrier give the three unnormalised sums.
Finally, \cref{thm:actual-source-price-equality-decision} applies to
this same complete word: positive price is charged, while zero price
re-enters the equality row without losing an ancestor.  This proves the
quantitative check for exact support \(A_{lb}\).
\end{proof}

\begin{theorem}[Local quantitative check: \(\Geom\Lift\Bdry\)]
\label{thm:quantitative-source-cell-glb}
For every finite same-origin stopped record \(R\in\mathfrak D_{A_{glb},\Sigma,n}^{\rm pre}(E)\) set
\[
\delta_{glb}(R):=\left\|\left(G^{\Geom}_{\Sigma},\operatorname{Ker}^{\Geom}_{\Sigma},\operatorname{Gc}^{\Geom}_{\Sigma},L^{\rm int},L^{\rm fib},L^{\rm conc},L^{\rm osc},L^{\rm sc},L^{\rm gc},B^{\rm ext},B^{\rm tail},B^{\rm term},B^{\rm gc}\right)\right\|_{\rm det},
\qquad a_{glb}(R):=\left\|\operatorname{Inc}_{A_{glb}}^{\Sigma}(R)\right\|,
\qquad p_{glb}(R):=p_{\Phys}(R).
\]
For a finite family \(\mathcal F=(R_1,\ldots,R_N)\) of pairwise
disjoint same-origin stopped windows with exact word \(A_{glb}\), define the
unnormalised local aggregates
\[
\mathsf A_{glb}(\mathcal F):=\sum_{j=1}^{N}a_{glb}(R_j),
\qquad \mathsf D_{glb}(\mathcal F):=\sum_{j=1}^{N}\delta_{glb}(R_j),
\qquad \mathsf P_{glb}(\mathcal F):=\sum_{j=1}^{N}p_{glb}(R_j).
\]
Then the following check is exact:
\begin{enumerate}[label=\textup{(LGLB\arabic*)}]
\item \(a_{glb}(R)=0\) exactly when this record contributes no
 target progress through exact support \(A_{glb}\).  If
\(a_{glb}(R)>0\), the record itself is the finite public
target-aligned occurrence, with the listed detector tuple and complete joint frontier.
\item \(\mathsf A_{glb}(\mathcal F),\ \mathsf D_{glb}(\mathcal F),\ \text{and }\mathsf P_{glb}(\mathcal F)\) are the exact cumulative target amplitude,
structural detector amplitude, and physical expenditure of this support.
They are never divided by \(N\) or replaced by a limiting average.
\item If \(p_{glb}(R)>0\), the occurrence is charged to the single
physical carrier.  If \(p_{glb}(R)=0\), the equality row of the
same joint record is evaluated; a nonzero conditional increment is appended
to the same exact word \(A_{glb}\), while a zero increment deletes that
record's target-incidence edge.
\end{enumerate}
\end{theorem}
\begin{proof}
The atomic target-aligned relations for \(\Geom\), \(\Lift\), \(\Bdry\) are respectively
\cref{prop:ns-geometric-target-relation}, \cref{prop:ns-lift-target-relation}, \cref{prop:ns-bdry-target-relation}.  Each is paired with the same public shell covector before any
source projection.  Their finite expression graphs are joined, so every
mixed product, commutator, cutoff, pressure descent, chart motion, and
trace retains the union word \(A_{glb}\).  Restricting that joint
record to exact word \(A_{glb}\) gives precisely the displayed
detector \(\delta_{glb}\) and target amplitude
\(a_{glb}\); hence its zero/nonzero test is the exact target pairing
of this source combination, not a separate test of its coordinates.
For disjoint stopped windows, linearity of the public weak pairing and
finite additivity of the physical carrier give the three unnormalised sums.
Finally, \cref{thm:actual-source-price-equality-decision} applies to
this same complete word: positive price is charged, while zero price
re-enters the equality row without losing an ancestor.  This proves the
quantitative check for exact support \(A_{glb}\).
\end{proof}

\begin{theorem}[Local quantitative check: \(\Caus\Lift\Bdry\)]
\label{thm:quantitative-source-cell-clb}
For every finite same-origin stopped record \(R\in\mathfrak D_{A_{clb},\Sigma,n}^{\rm pre}(E)\) set
\[
\delta_{clb}(R):=\left\|\left(H^{\Caus},S^{\Caus},\Psi^{\Caus},\operatorname{Gc}^{\Caus},L^{\rm int},L^{\rm fib},L^{\rm conc},L^{\rm osc},L^{\rm sc},L^{\rm gc},B^{\rm ext},B^{\rm tail},B^{\rm term},B^{\rm gc}\right)\right\|_{\rm det},
\qquad a_{clb}(R):=\left\|\operatorname{Inc}_{A_{clb}}^{\Sigma}(R)\right\|,
\qquad p_{clb}(R):=p_{\Phys}(R).
\]
For a finite family \(\mathcal F=(R_1,\ldots,R_N)\) of pairwise
disjoint same-origin stopped windows with exact word \(A_{clb}\), define the
unnormalised local aggregates
\[
\mathsf A_{clb}(\mathcal F):=\sum_{j=1}^{N}a_{clb}(R_j),
\qquad \mathsf D_{clb}(\mathcal F):=\sum_{j=1}^{N}\delta_{clb}(R_j),
\qquad \mathsf P_{clb}(\mathcal F):=\sum_{j=1}^{N}p_{clb}(R_j).
\]
Then the following check is exact:
\begin{enumerate}[label=\textup{(LCLB\arabic*)}]
\item \(a_{clb}(R)=0\) exactly when this record contributes no
 target progress through exact support \(A_{clb}\).  If
\(a_{clb}(R)>0\), the record itself is the finite public
target-aligned occurrence, with the listed detector tuple and complete joint frontier.
\item \(\mathsf A_{clb}(\mathcal F),\ \mathsf D_{clb}(\mathcal F),\ \text{and }\mathsf P_{clb}(\mathcal F)\) are the exact cumulative target amplitude,
structural detector amplitude, and physical expenditure of this support.
They are never divided by \(N\) or replaced by a limiting average.
\item If \(p_{clb}(R)>0\), the occurrence is charged to the single
physical carrier.  If \(p_{clb}(R)=0\), the equality row of the
same joint record is evaluated; a nonzero conditional increment is appended
to the same exact word \(A_{clb}\), while a zero increment deletes that
record's target-incidence edge.
\end{enumerate}
\end{theorem}
\begin{proof}
The atomic target-aligned relations for \(\Caus\), \(\Lift\), \(\Bdry\) are respectively
\cref{prop:ns-caus-target-relation}, \cref{prop:ns-lift-target-relation}, \cref{prop:ns-bdry-target-relation}.  Each is paired with the same public shell covector before any
source projection.  Their finite expression graphs are joined, so every
mixed product, commutator, cutoff, pressure descent, chart motion, and
trace retains the union word \(A_{clb}\).  Restricting that joint
record to exact word \(A_{clb}\) gives precisely the displayed
detector \(\delta_{clb}\) and target amplitude
\(a_{clb}\); hence its zero/nonzero test is the exact target pairing
of this source combination, not a separate test of its coordinates.
For disjoint stopped windows, linearity of the public weak pairing and
finite additivity of the physical carrier give the three unnormalised sums.
Finally, \cref{thm:actual-source-price-equality-decision} applies to
this same complete word: positive price is charged, while zero price
re-enters the equality row without losing an ancestor.  This proves the
quantitative check for exact support \(A_{clb}\).
\end{proof}

\begin{theorem}[Local quantitative check: \(\Geom\Caus\Lift\Bdry\)]
\label{thm:quantitative-source-cell-gclb}
For every finite same-origin stopped record \(R\in\mathfrak D_{A_{gclb},\Sigma,n}^{\rm pre}(E)\) set
\[
\delta_{gclb}(R):=\left\|\left(G^{\Geom}_{\Sigma},\operatorname{Ker}^{\Geom}_{\Sigma},\operatorname{Gc}^{\Geom}_{\Sigma},H^{\Caus},S^{\Caus},\Psi^{\Caus},\operatorname{Gc}^{\Caus},L^{\rm int},L^{\rm fib},L^{\rm conc},L^{\rm osc},L^{\rm sc},L^{\rm gc},B^{\rm ext},B^{\rm tail},B^{\rm term},B^{\rm gc}\right)\right\|_{\rm det},
\qquad a_{gclb}(R):=\left\|\operatorname{Inc}_{A_{gclb}}^{\Sigma}(R)\right\|,
\qquad p_{gclb}(R):=p_{\Phys}(R).
\]
For a finite family \(\mathcal F=(R_1,\ldots,R_N)\) of pairwise
disjoint same-origin stopped windows with exact word \(A_{gclb}\), define the
unnormalised local aggregates
\[
\mathsf A_{gclb}(\mathcal F):=\sum_{j=1}^{N}a_{gclb}(R_j),
\qquad \mathsf D_{gclb}(\mathcal F):=\sum_{j=1}^{N}\delta_{gclb}(R_j),
\qquad \mathsf P_{gclb}(\mathcal F):=\sum_{j=1}^{N}p_{gclb}(R_j).
\]
Then the following check is exact:
\begin{enumerate}[label=\textup{(LGCLB\arabic*)}]
\item \(a_{gclb}(R)=0\) exactly when this record contributes no
 target progress through exact support \(A_{gclb}\).  If
\(a_{gclb}(R)>0\), the record itself is the finite public
target-aligned occurrence, with the listed detector tuple and complete joint frontier.
\item \(\mathsf A_{gclb}(\mathcal F),\ \mathsf D_{gclb}(\mathcal F),\ \text{and }\mathsf P_{gclb}(\mathcal F)\) are the exact cumulative target amplitude,
structural detector amplitude, and physical expenditure of this support.
They are never divided by \(N\) or replaced by a limiting average.
\item If \(p_{gclb}(R)>0\), the occurrence is charged to the single
physical carrier.  If \(p_{gclb}(R)=0\), the equality row of the
same joint record is evaluated; a nonzero conditional increment is appended
to the same exact word \(A_{gclb}\), while a zero increment deletes that
record's target-incidence edge.
\end{enumerate}
\end{theorem}
\begin{proof}
The atomic target-aligned relations for \(\Geom\), \(\Caus\), \(\Lift\), \(\Bdry\) are respectively
\cref{prop:ns-geometric-target-relation}, \cref{prop:ns-caus-target-relation}, \cref{prop:ns-lift-target-relation}, \cref{prop:ns-bdry-target-relation}.  Each is paired with the same public shell covector before any
source projection.  Their finite expression graphs are joined, so every
mixed product, commutator, cutoff, pressure descent, chart motion, and
trace retains the union word \(A_{gclb}\).  Restricting that joint
record to exact word \(A_{gclb}\) gives precisely the displayed
detector \(\delta_{gclb}\) and target amplitude
\(a_{gclb}\); hence its zero/nonzero test is the exact target pairing
of this source combination, not a separate test of its coordinates.
For disjoint stopped windows, linearity of the public weak pairing and
finite additivity of the physical carrier give the three unnormalised sums.
Finally, \cref{thm:actual-source-price-equality-decision} applies to
this same complete word: positive price is charged, while zero price
re-enters the equality row without losing an ancestor.  This proves the
quantitative check for exact support \(A_{gclb}\).
\end{proof}

\begin{theorem}[Local quantitative check: \(\Abs\Lift\Bdry\)]
\label{thm:quantitative-source-cell-alb}
For every finite same-origin stopped record \(R\in\mathfrak D_{A_{alb},\Sigma,n}^{\rm pre}(E)\) set
\[
\delta_{alb}(R):=\left\|\left(A^{\rm desc},A^{\rm vert},A^{\rm coker},A^{\rm gc},A^{\rm tgt},L^{\rm int},L^{\rm fib},L^{\rm conc},L^{\rm osc},L^{\rm sc},L^{\rm gc},B^{\rm ext},B^{\rm tail},B^{\rm term},B^{\rm gc}\right)\right\|_{\rm det},
\qquad a_{alb}(R):=\left\|\operatorname{Inc}_{A_{alb}}^{\Sigma}(R)\right\|,
\qquad p_{alb}(R):=p_{\Phys}(R).
\]
For a finite family \(\mathcal F=(R_1,\ldots,R_N)\) of pairwise
disjoint same-origin stopped windows with exact word \(A_{alb}\), define the
unnormalised local aggregates
\[
\mathsf A_{alb}(\mathcal F):=\sum_{j=1}^{N}a_{alb}(R_j),
\qquad \mathsf D_{alb}(\mathcal F):=\sum_{j=1}^{N}\delta_{alb}(R_j),
\qquad \mathsf P_{alb}(\mathcal F):=\sum_{j=1}^{N}p_{alb}(R_j).
\]
Then the following check is exact:
\begin{enumerate}[label=\textup{(LALB\arabic*)}]
\item \(a_{alb}(R)=0\) exactly when this record contributes no
 target progress through exact support \(A_{alb}\).  If
\(a_{alb}(R)>0\), the record itself is the finite public
target-aligned occurrence, with the listed detector tuple and complete joint frontier.
\item \(\mathsf A_{alb}(\mathcal F),\ \mathsf D_{alb}(\mathcal F),\ \text{and }\mathsf P_{alb}(\mathcal F)\) are the exact cumulative target amplitude,
structural detector amplitude, and physical expenditure of this support.
They are never divided by \(N\) or replaced by a limiting average.
\item If \(p_{alb}(R)>0\), the occurrence is charged to the single
physical carrier.  If \(p_{alb}(R)=0\), the equality row of the
same joint record is evaluated; a nonzero conditional increment is appended
to the same exact word \(A_{alb}\), while a zero increment deletes that
record's target-incidence edge.
\end{enumerate}
\end{theorem}
\begin{proof}
The atomic target-aligned relations for \(\Abs\), \(\Lift\), \(\Bdry\) are respectively
\cref{prop:ns-abs-target-relation}, \cref{prop:ns-lift-target-relation}, \cref{prop:ns-bdry-target-relation}.  Each is paired with the same public shell covector before any
source projection.  Their finite expression graphs are joined, so every
mixed product, commutator, cutoff, pressure descent, chart motion, and
trace retains the union word \(A_{alb}\).  Restricting that joint
record to exact word \(A_{alb}\) gives precisely the displayed
detector \(\delta_{alb}\) and target amplitude
\(a_{alb}\); hence its zero/nonzero test is the exact target pairing
of this source combination, not a separate test of its coordinates.
For disjoint stopped windows, linearity of the public weak pairing and
finite additivity of the physical carrier give the three unnormalised sums.
Finally, \cref{thm:actual-source-price-equality-decision} applies to
this same complete word: positive price is charged, while zero price
re-enters the equality row without losing an ancestor.  This proves the
quantitative check for exact support \(A_{alb}\).
\end{proof}

\begin{theorem}[Local quantitative check: \(\Geom\Abs\Lift\Bdry\)]
\label{thm:quantitative-source-cell-galb}
For every finite same-origin stopped record \(R\in\mathfrak D_{A_{galb},\Sigma,n}^{\rm pre}(E)\) set
\[
\delta_{galb}(R):=\left\|\left(G^{\Geom}_{\Sigma},\operatorname{Ker}^{\Geom}_{\Sigma},\operatorname{Gc}^{\Geom}_{\Sigma},A^{\rm desc},A^{\rm vert},A^{\rm coker},A^{\rm gc},A^{\rm tgt},L^{\rm int},L^{\rm fib},L^{\rm conc},L^{\rm osc},L^{\rm sc},L^{\rm gc},B^{\rm ext},B^{\rm tail},B^{\rm term},B^{\rm gc}\right)\right\|_{\rm det},
\qquad a_{galb}(R):=\left\|\operatorname{Inc}_{A_{galb}}^{\Sigma}(R)\right\|,
\qquad p_{galb}(R):=p_{\Phys}(R).
\]
For a finite family \(\mathcal F=(R_1,\ldots,R_N)\) of pairwise
disjoint same-origin stopped windows with exact word \(A_{galb}\), define the
unnormalised local aggregates
\[
\mathsf A_{galb}(\mathcal F):=\sum_{j=1}^{N}a_{galb}(R_j),
\qquad \mathsf D_{galb}(\mathcal F):=\sum_{j=1}^{N}\delta_{galb}(R_j),
\qquad \mathsf P_{galb}(\mathcal F):=\sum_{j=1}^{N}p_{galb}(R_j).
\]
Then the following check is exact:
\begin{enumerate}[label=\textup{(LGALB\arabic*)}]
\item \(a_{galb}(R)=0\) exactly when this record contributes no
 target progress through exact support \(A_{galb}\).  If
\(a_{galb}(R)>0\), the record itself is the finite public
target-aligned occurrence, with the listed detector tuple and complete joint frontier.
\item \(\mathsf A_{galb}(\mathcal F),\ \mathsf D_{galb}(\mathcal F),\ \text{and }\mathsf P_{galb}(\mathcal F)\) are the exact cumulative target amplitude,
structural detector amplitude, and physical expenditure of this support.
They are never divided by \(N\) or replaced by a limiting average.
\item If \(p_{galb}(R)>0\), the occurrence is charged to the single
physical carrier.  If \(p_{galb}(R)=0\), the equality row of the
same joint record is evaluated; a nonzero conditional increment is appended
to the same exact word \(A_{galb}\), while a zero increment deletes that
record's target-incidence edge.
\end{enumerate}
\end{theorem}
\begin{proof}
The atomic target-aligned relations for \(\Geom\), \(\Abs\), \(\Lift\), \(\Bdry\) are respectively
\cref{prop:ns-geometric-target-relation}, \cref{prop:ns-abs-target-relation}, \cref{prop:ns-lift-target-relation}, \cref{prop:ns-bdry-target-relation}.  Each is paired with the same public shell covector before any
source projection.  Their finite expression graphs are joined, so every
mixed product, commutator, cutoff, pressure descent, chart motion, and
trace retains the union word \(A_{galb}\).  Restricting that joint
record to exact word \(A_{galb}\) gives precisely the displayed
detector \(\delta_{galb}\) and target amplitude
\(a_{galb}\); hence its zero/nonzero test is the exact target pairing
of this source combination, not a separate test of its coordinates.
For disjoint stopped windows, linearity of the public weak pairing and
finite additivity of the physical carrier give the three unnormalised sums.
Finally, \cref{thm:actual-source-price-equality-decision} applies to
this same complete word: positive price is charged, while zero price
re-enters the equality row without losing an ancestor.  This proves the
quantitative check for exact support \(A_{galb}\).
\end{proof}

\begin{theorem}[Local quantitative check: \(\Caus\Abs\Lift\Bdry\)]
\label{thm:quantitative-source-cell-calb}
For every finite same-origin stopped record \(R\in\mathfrak D_{A_{calb},\Sigma,n}^{\rm pre}(E)\) set
\[
\delta_{calb}(R):=\left\|\left(H^{\Caus},S^{\Caus},\Psi^{\Caus},\operatorname{Gc}^{\Caus},A^{\rm desc},A^{\rm vert},A^{\rm coker},A^{\rm gc},A^{\rm tgt},L^{\rm int},L^{\rm fib},L^{\rm conc},L^{\rm osc},L^{\rm sc},L^{\rm gc},B^{\rm ext},B^{\rm tail},B^{\rm term},B^{\rm gc}\right)\right\|_{\rm det},
\qquad a_{calb}(R):=\left\|\operatorname{Inc}_{A_{calb}}^{\Sigma}(R)\right\|,
\qquad p_{calb}(R):=p_{\Phys}(R).
\]
For a finite family \(\mathcal F=(R_1,\ldots,R_N)\) of pairwise
disjoint same-origin stopped windows with exact word \(A_{calb}\), define the
unnormalised local aggregates
\[
\mathsf A_{calb}(\mathcal F):=\sum_{j=1}^{N}a_{calb}(R_j),
\qquad \mathsf D_{calb}(\mathcal F):=\sum_{j=1}^{N}\delta_{calb}(R_j),
\qquad \mathsf P_{calb}(\mathcal F):=\sum_{j=1}^{N}p_{calb}(R_j).
\]
Then the following check is exact:
\begin{enumerate}[label=\textup{(LCALB\arabic*)}]
\item \(a_{calb}(R)=0\) exactly when this record contributes no
 target progress through exact support \(A_{calb}\).  If
\(a_{calb}(R)>0\), the record itself is the finite public
target-aligned occurrence, with the listed detector tuple and complete joint frontier.
\item \(\mathsf A_{calb}(\mathcal F),\ \mathsf D_{calb}(\mathcal F),\ \text{and }\mathsf P_{calb}(\mathcal F)\) are the exact cumulative target amplitude,
structural detector amplitude, and physical expenditure of this support.
They are never divided by \(N\) or replaced by a limiting average.
\item If \(p_{calb}(R)>0\), the occurrence is charged to the single
physical carrier.  If \(p_{calb}(R)=0\), the equality row of the
same joint record is evaluated; a nonzero conditional increment is appended
to the same exact word \(A_{calb}\), while a zero increment deletes that
record's target-incidence edge.
\end{enumerate}
\end{theorem}
\begin{proof}
The atomic target-aligned relations for \(\Caus\), \(\Abs\), \(\Lift\), \(\Bdry\) are respectively
\cref{prop:ns-caus-target-relation}, \cref{prop:ns-abs-target-relation}, \cref{prop:ns-lift-target-relation}, \cref{prop:ns-bdry-target-relation}.  Each is paired with the same public shell covector before any
source projection.  Their finite expression graphs are joined, so every
mixed product, commutator, cutoff, pressure descent, chart motion, and
trace retains the union word \(A_{calb}\).  Restricting that joint
record to exact word \(A_{calb}\) gives precisely the displayed
detector \(\delta_{calb}\) and target amplitude
\(a_{calb}\); hence its zero/nonzero test is the exact target pairing
of this source combination, not a separate test of its coordinates.
For disjoint stopped windows, linearity of the public weak pairing and
finite additivity of the physical carrier give the three unnormalised sums.
Finally, \cref{thm:actual-source-price-equality-decision} applies to
this same complete word: positive price is charged, while zero price
re-enters the equality row without losing an ancestor.  This proves the
quantitative check for exact support \(A_{calb}\).
\end{proof}

\begin{theorem}[Local quantitative check: \(\Geom\Caus\Abs\Lift\Bdry\)]
\label{thm:quantitative-source-cell-gcalb}
For every finite same-origin stopped record \(R\in\mathfrak D_{A_{gcalb},\Sigma,n}^{\rm pre}(E)\) set
\[
\delta_{gcalb}(R):=\left\|\left(G^{\Geom}_{\Sigma},\operatorname{Ker}^{\Geom}_{\Sigma},\operatorname{Gc}^{\Geom}_{\Sigma},H^{\Caus},S^{\Caus},\Psi^{\Caus},\operatorname{Gc}^{\Caus},A^{\rm desc},A^{\rm vert},A^{\rm coker},A^{\rm gc},A^{\rm tgt},L^{\rm int},L^{\rm fib},L^{\rm conc},L^{\rm osc},L^{\rm sc},L^{\rm gc},B^{\rm ext},B^{\rm tail},B^{\rm term},B^{\rm gc}\right)\right\|_{\rm det},
\qquad a_{gcalb}(R):=\left\|\operatorname{Inc}_{A_{gcalb}}^{\Sigma}(R)\right\|,
\qquad p_{gcalb}(R):=p_{\Phys}(R).
\]
For a finite family \(\mathcal F=(R_1,\ldots,R_N)\) of pairwise
disjoint same-origin stopped windows with exact word \(A_{gcalb}\), define the
unnormalised local aggregates
\[
\mathsf A_{gcalb}(\mathcal F):=\sum_{j=1}^{N}a_{gcalb}(R_j),
\qquad \mathsf D_{gcalb}(\mathcal F):=\sum_{j=1}^{N}\delta_{gcalb}(R_j),
\qquad \mathsf P_{gcalb}(\mathcal F):=\sum_{j=1}^{N}p_{gcalb}(R_j).
\]
Then the following check is exact:
\begin{enumerate}[label=\textup{(LGCALB\arabic*)}]
\item \(a_{gcalb}(R)=0\) exactly when this record contributes no
 target progress through exact support \(A_{gcalb}\).  If
\(a_{gcalb}(R)>0\), the record itself is the finite public
target-aligned occurrence, with the listed detector tuple and complete joint frontier.
\item \(\mathsf A_{gcalb}(\mathcal F),\ \mathsf D_{gcalb}(\mathcal F),\ \text{and }\mathsf P_{gcalb}(\mathcal F)\) are the exact cumulative target amplitude,
structural detector amplitude, and physical expenditure of this support.
They are never divided by \(N\) or replaced by a limiting average.
\item If \(p_{gcalb}(R)>0\), the occurrence is charged to the single
physical carrier.  If \(p_{gcalb}(R)=0\), the equality row of the
same joint record is evaluated; a nonzero conditional increment is appended
to the same exact word \(A_{gcalb}\), while a zero increment deletes that
record's target-incidence edge.
\end{enumerate}
\end{theorem}
\begin{proof}
The atomic target-aligned relations for \(\Geom\), \(\Caus\), \(\Abs\), \(\Lift\), \(\Bdry\) are respectively
\cref{prop:ns-geometric-target-relation}, \cref{prop:ns-caus-target-relation}, \cref{prop:ns-abs-target-relation}, \cref{prop:ns-lift-target-relation}, \cref{prop:ns-bdry-target-relation}.  Each is paired with the same public shell covector before any
source projection.  Their finite expression graphs are joined, so every
mixed product, commutator, cutoff, pressure descent, chart motion, and
trace retains the union word \(A_{gcalb}\).  Restricting that joint
record to exact word \(A_{gcalb}\) gives precisely the displayed
detector \(\delta_{gcalb}\) and target amplitude
\(a_{gcalb}\); hence its zero/nonzero test is the exact target pairing
of this source combination, not a separate test of its coordinates.
For disjoint stopped windows, linearity of the public weak pairing and
finite additivity of the physical carrier give the three unnormalised sums.
Finally, \cref{thm:actual-source-price-equality-decision} applies to
this same complete word: positive price is charged, while zero price
re-enters the equality row without losing an ancestor.  This proves the
quantitative check for exact support \(A_{gcalb}\).
\end{proof}



\begin{definition}[Exact-support cumulative conversion ledger]
\label{def:exact-support-cumulative-conversion-ledger}
Fix a nonempty exact source word \(A\subseteq\Channels\), and let
\(R_1,\ldots,R_N\) be an ordered same-origin family of finite stopped
records in \(\mathfrak D_{A,\Sigma}^{\rm pre}(\Budget)\), with the intervening
gaps included in the stopped interval complex.  Fix the active stratum
\(X_{\mathfrak b}\) and its shell covector \(\lambda=D\Phi\).  For each
\(R_j\), first form its cumulative public joint record, restrict the
target-local factorization
\eqref{eq:target-local-fact-normal-form} to \(X_{\mathfrak b}\), apply
\(q_\Sigma\), and then select the exact provenance word \(A\) using
\cref{thm:universal-five-source-provenance}.  Let \(p_A(R_j)\) be the sum of
precisely those \(A\)-owned production or exterior-input terms admitted by
\eqref{eq:target-aligned-charge-rule} and dominated by the restriction of the
single public physical measure to the event indices of \(R_j\).  Every
carrier term without such a certificate is retained in the cone-boundary or
polar component of \(e_A^0(R_j)\), rather than being assigned a price.  Set
\[
 a_A(R_j):=\alpha_A(R_j),\qquad
 B_A(R_j):=\Lambda_A(R_j^+)-\Lambda_A(R_j^-)
              +p_A(R_j)+e_A^0(R_j).
\]
Expand the public weak identity on this complex, group every term by its
complete source word before taking positive parts, and cancel oppositely
oriented internal faces.  The signed row
\[
 q_A(R_j)=\Lambda_A(R_j^+)-\Lambda_A(R_j^-)
           +p_A(R_j)+r_A(R_j)
\]
is retained before any positive part.  Define the conversion-defect and
slack coordinates by their total positive-part graphs,
\[
 d_A(R_j):=(a_A(R_j)-B_A(R_j))_+,
 \qquad s_A(R_j):=(B_A(R_j)-a_A(R_j))_+,
 \qquad e_A(R_j):=e_A^0(R_j)+d_A(R_j).
\]
An \emph{admissible conversion row} retains these graph coordinates and the
resulting finite-cylinder identity
\begin{equation}
 \label{eq:exact-support-source-conversion}
 \Lambda_A(R_j^+)-\Lambda_A(R_j^-)
 +p_A(R_j)+e_A(R_j)
 =a_A(R_j)+s_A(R_j),
 \qquad s_A(R_j)\ge0.
\end{equation}
Here \(\Lambda_A(R)\) is the signed target pairing of the ordered tuple of
the \(A\)-owned endpoint, feeding, storage, scale, and unglued-interface
coordinates at the cut \(R\), hence the \(A\)-restriction of
\(\delta K_{\mathfrak b}+\sum_cR^c_{\mathfrak b,\lambda}\) in
\eqref{eq:target-local-fact-normal-form}.  The scalar \(e_A^0(R_j)\ge0\) is
the positive target readout of the remaining equality, cone-boundary, polar,
and total-graph coordinates in the same expanded word; \(d_A\) is retained as
the first failed conversion-graph coordinate of that remainder row.  Its
nonzero value is therefore a target-local polarity successor with the same
origin, covector, joint frontier, exact word, and carrier owner.  The full
tuple is retained alongside these scalar readouts.  The certificate is a
transcript in \(\mathsf{Proof}_{\Sigma}(\mathfrak P)\): it may use the signed
public weak row, declared positive carrier atoms, and verified
quadratic-module rows, but it may not
infer the displayed identity merely by taking the positive part of the signed
weak row.  Its final equality is instead the elementary identity
\(x+(-x)_+=x_+\), applied to \(x=B_A-a_A\).  Every omitted term is separately certified to lie in
\(\Null_\Sigma\).  A generated record for which this transcript does not
verify is a conversion-obstruction successor, not an admissible ledger row.
\end{definition}

\begin{theorem}[Uniform exact-support conversion]
\label{thm:uniform-exact-support-conversion}
For every nonempty exact support \(A\) and every finite same-origin stopped
active record \(R\in\mathfrak D_{A,\Sigma}^{\rm pre}(\Budget)\) in an accepted
presentation, the ledger of
\cref{def:exact-support-cumulative-conversion-ledger} gives
\[
 \boxed{
 \alpha_A(R)
 \le\Lambda_A(R^+)-\Lambda_A(R^-)+p_A(R)+e_A(R).}
\]
The theorem is parametric in \(A\).  Restricting it to one of the thirty-one
supports produces a specialization of the same signed current identity, not
an independent estimate.
\end{theorem}

\begin{proof}
The checker first verifies the signed weak row and all gluing cancellations
using the target-local factorization and local charging rule, then verifies
the two total positive-part graphs.  The identity
\(x+(-x)_+=x_+\), with \(x=B_A-a_A\), is exactly
\eqref{eq:exact-support-source-conversion}.  Dropping \(s_A(R)\ge0\) from
\eqref{eq:exact-support-source-conversion} gives the displayed inequality.
This step is valid even when the signed transfer difference is negative; no
invalid implication \(q\leq q_+\) in the opposite direction is used.  The
same finite checker is parametric in the symbolic support \(A\).
\end{proof}

\begin{theorem}[Sourcewise telescoping]
\label{prop:sourcewise-cascade-telescoping-reentry}
For every exact support \(A\) and every ordered finite stopped complex all of
whose rows are admissible in the sense of
\cref{def:exact-support-cumulative-conversion-ledger},
\begin{equation}
 \label{eq:sourcewise-cascade-telescoping}
 \sum_{j=1}^{N}a_A(R_j)
 \le
 \Lambda_A(R_N^+)-\Lambda_A(R_1^-)
 +\sum_{j=1}^{N}p_A(R_j)
 +\sum_{j=1}^{N}e_A(R_j).
\end{equation}
The accepted ledger has the following closure rules.
\begin{enumerate}[label=\textup{(S\arabic*)}]
\item The sums in \eqref{eq:sourcewise-cascade-telescoping} are formed before
  division by \(N\), passage to a subsequence, or construction of an
  empirical or recurrent law.
\item An unbounded target projection of \(\Lambda_A\) has a first rational
  level-crossing prefix.  Backward factorization assigns its least nonzero
  conditional increment to the constructor that created the endpoint,
  feeding, storage, scale, or interface coordinate.  That increment is
  appended to the same cumulative word and evaluated again by
  \eqref{eq:exact-support-source-conversion}.
\item On the zero-price row, a nonzero target projection of \(e_A\) has a
  least finite contextual increment and is appended to the same word.  If no
  such increment exists, the remainder lies in \(\Null_\Sigma\).
\end{enumerate}
Consequently, after saturation under \textup{(S2)}--\textup{(S3)}, the
unprocessed residual is target-null.  This conclusion is conditional only on
the displayed finite conversion certificates; it uses no boundedness,
compactness, coercivity, or terminal-realization hypothesis.
\end{theorem}

\begin{proof}
Apply \cref{thm:uniform-exact-support-conversion} to every certified row and
sum the resulting inequalities.  The stopped complex contains every gap and
records opposite orientations on its two copies of each internal face.
Consequently every internal \(\Lambda_A\)-value cancels algebraically, while
the two exterior values remain.  Dropping the sum of nonnegative certified
slacks gives \eqref{eq:sourcewise-cascade-telescoping}.  Notice that the
positive parts have already been controlled separately by the dual
transcripts; they are not moved through the signed telescoping sum.

All cumulative coordinates belong to the universal joint record of
\cref{def:cumulative-joint-continuum-record}, which proves \textup{(S1)}.
For an unbounded storage coordinate, choose its first crossing of the least
rational level above the preceding prefix.  This is a finite stopped record,
so \cref{thm:exact-prospective-continuum-backward-factorization} returns its
least nonzero contextual increment with the complete constructor ancestry.
Universal provenance places that increment in the same exact union word,
proving \textup{(S2)}.  The same argument on the equality row, followed by
\cref{thm:actual-source-price-equality-decision}, proves \textup{(S3)}.
If a target-visible residual remained after all such re-entries, then
\cref{thm:cofinal-prospective-frontier-closure} would expose a finite
unprocessed contextual increment, contrary to saturation.  Equivalently,
the residual-nullity conclusion is
\cref{cor:backward-saturated-residual-nullity} applied supportwise.
\end{proof}

\begin{theorem}[Storage and feeding re-entry]
\label{thm:storage-feeding-reentry}
Fix an exact support \(A\) and an ordered same-origin complex.  If the
target projection of its cumulative transfer coordinate \(\Lambda_A\) is
unbounded, then for every positive rational level \(m\) there is a finite
first level-crossing prefix \(R_m\), its responsible endpoint, feeding,
storage, scale, or interface constructor, and the complete cumulative
unnormalised shell cocycle.  The transfer detector on this prefix has value
at least \(m\).  Its calibrated backward-factorization successor
\(R_m^{\rm ctx}\) satisfies
\[
 \alpha_{A_m}(R_m^{\rm ctx})\ge\kappa_{A_m}m>0
\]
for its inherited support \(A_m\), which is \(A\) unless the responsible
constructor has a recorded support-transfer edge.  The successor is evaluated
again by \cref{thm:uniform-exact-support-conversion}.
\end{theorem}

\begin{proof}
Choose the least prefix whose cumulative transfer detector reaches \(m\).
It is finite because the transfer is unbounded.  The stopped complex retains
the endpoint graph and all intervening gaps, so prospective backward
factorization selects the responsible finite constructor occurrence with the
same origin and cocycle.  Apply
\cref{thm:detector-to-amplitude-calibration} to its transfer detector.  If
the selected constructor changes provenance, apply
\cref{thm:support-transfer-conservation}; otherwise the exact word remains
\(A\).  The resulting finite record is in the domain of the uniform
conversion theorem.
\end{proof}

\begin{theorem}[Equality and graph-remainder re-entry]
\label{thm:equality-graph-remainder-reentry}
Fix an exact support \(A\).  If the cumulative target readout of
\(e_A\) is nonzero, or if \(\sum_j e_A(R_j)=+\infty\), then the fixed
enumeration returns a least finite same-origin equality or total-graph
occurrence with its detecting covector, source route, and complete
predecessor transcript.  Its calibrated contextual successor has positive
canonical amplitude and re-enters
\cref{thm:uniform-exact-support-conversion}.  If no such finite occurrence
exists, the cumulative remainder lies in \(\Null_\Sigma\).
\end{theorem}

\begin{proof}
Enumerate the bounded rational truncations of all equality, cone-boundary,
polar, and total-graph coordinates in \(e_A\).  A nonzero cumulative readout
has a least positive finite truncation increment; the divergent case has a
first prefix exceeding each rational level.  Backward factorization supplies
the finite contextual occurrence, and detector calibration gives its positive
canonical amplitude.  If every finite truncation increment vanishes, then
the cumulative no-emergence theorem exposes no target-visible limit
increment; equivalently every generated shell covector annihilates the
remainder, placing it in \(\Null_\Sigma\).
\end{proof}

\begin{definition}[Canonical exact-support shell amplitude]
\label{def:canonical-exact-support-shell-amplitude}
Let \(R\) be a finite normalized same-origin public record of the canonical shell construction
\cref{def:target-shell-presentation}.  It carries the public shell covector
\(\lambda=D\Phi\), one physical-event index, origin, and the finite prefix of
its cumulative joint record.  Expand that public joint current completely
before any positive part or application of \(q_\Sigma\), and write its
target-current pairing as
\[
 1=\sum_i\operatorname{Prog}_\Sigma(T_i;R),
\]
where the sum is over the top-level typed atomic current occurrences generated
from the public weak and total graphs; the displayed equality is the
event-normalized shell identity after the target-null residual is removed by
\(q_\Sigma\).  For every nonempty exact support
\(A\subseteq\Channels\), define the signed exact-support current and its
canonical nonnegative amplitude by
\begin{equation}
 \label{eq:canonical-exact-support-amplitude}
 q_A(R):=\sum_{\operatorname{Prov}(T_i)=A}
              \operatorname{Prog}_\Sigma(T_i;R),
 \qquad
 \alpha_A(R):=\bigl(q_A(R)\bigr)_+.
\end{equation}
The sum in \(q_A\) is an ordinary finite signed sum in
\(\mathsf{Proof}_{\Sigma}(\mathfrak P)\), and its support is the value of the
public provenance morphism.  Its nonempty \(A\)-restriction is a record of
\(\mathfrak D_{A,\Sigma}^{\rm pre}(\Budget)\).  Conditional increments,
detector norms, and normalized occurrence
coordinates are retained as additional coordinates of \(R\); none is
identified with \(\alpha_A(R)\) without a stated calibration theorem.
\end{definition}

\begin{theorem}[One-currency exact-support amplitude]
\label{thm:one-currency-exact-support-amplitude}
For every finite normalized same-origin shell record \(R\),
\begin{equation}
 \label{eq:one-currency-exact-support-identity}
 1=\sum_{\varnothing\ne A\subseteq\Channels}q_A(R),
 \qquad
 \sum_{\varnothing\ne A\subseteq\Channels}\alpha_A(R)\ge1.
\end{equation}
Thus every physical-price, storage, equality, circulation, and
terminalization ledger that aggregates exact supports uses \(\alpha_A\), or
contains an explicit comparison to \(\alpha_A\).  In particular, a positive
detector or conditional increment is a location certificate rather than an
accounting amplitude until such a comparison has been proved.
\end{theorem}

\begin{proof}
Expand the canonical derived shell current into its top-level typed atomic
occurrences after public joint formation and before \(q_\Sigma\).  The
universal provenance morphism partitions this finite list by its
nonempty exact supports.  Linearity of the shell pairing before taking
positive parts gives the first identity in
\eqref{eq:one-currency-exact-support-identity}; the target-null residual has
zero pairing by \cref{def:target-null-residual}.  For real numbers \(x_A\),
\(\sum_A(x_A)_+\ge\sum_Ax_A\).  Apply this with \(x_A=q_A(R)\) and use the
first identity.  The final assertion is the definition of the common
currency: any distinct nonnegative coordinate requires a theorem comparing
it to \(\alpha_A\) before it can be summed in this ledger.
\end{proof}

\begin{definition}[Calibrated conditional-incidence detector]
\label{def:calibrated-conditional-incidence-detector}
Fix a nonempty exact support \(A\).  A conditional-incidence detector
\(d_A\) is \emph{calibrated for accounting} when its finite code contains
one of the following certificates on every normalized active shell cylinder.
The certificate must be a checked derivation from the restriction of
\(\mathsf{Proof}_{\Sigma}(\mathfrak P)\) to the same active stratum, shell
covector, cumulative joint frontier, origin, and exact word \(A\); a numerical
comparison supplied outside the public proof algebra is not a certificate.
\begin{enumerate}[label=\textup{(D\arabic*)}]
\item The exact certificate declares
\[
 d_A(R)=\alpha_A(R).
\]
\item The contextual certificate supplies a rational \(\kappa_A>0\), a
finite same-origin backward-factorization map
\(R\mapsto R_A^{\rm ctx}\), and a finite-cylinder inequality
\begin{equation}
 \label{eq:detector-contextual-amplitude-calibration}
 d_A(R)>0
 \quad\Longrightarrow\quad
 \alpha_A(R_A^{\rm ctx})\ge\kappa_A d_A(R).
\end{equation}
The certificate also supplies a modulus for this inequality on the
sourcewise closure.
\end{enumerate}
The canonical accounting detector is
\(d_A^{\rm cur}:=\alpha_A\) and carries certificate \textup{(D1)}.  A
conditional detector without either certificate may locate a finite
contextual occurrence, but it is not an accounting coordinate and is not
summed in a price, storage, circulation, or terminalization ledger.
\end{definition}

\begin{theorem}[Detector-to-amplitude calibration]
\label{thm:detector-to-amplitude-calibration}
Every calibrated conditional-incidence detector has a uniform accounting
interface along a normalized same-origin shell cascade.  The canonical
detector satisfies
\[
 d_A^{\rm cur}(R)=\alpha_A(R).
\]
For a detector carrying a contextual certificate, every positive value
returns the finite same-origin contextual occurrence \(R_A^{\rm ctx}\) and
satisfies \eqref{eq:detector-contextual-amplitude-calibration} with the same
presentation-generated constant \(\kappa_A>0\) on every shell.  Thus
\[
 \sum_j d_A(R_j)=+\infty
 \quad\Longrightarrow\quad
 \sum_j\alpha_A(R_{A,j}^{\rm ctx})=+\infty.
\]
\end{theorem}

\begin{proof}
For the canonical detector the assertion is its definition.  In the
contextual case, write the complete finite frontier coordinates as a vector
\(x\) on the normalized shell cylinder.  The detector and the exact-support
shell current are respectively the declared nonnegative readout \(d_A(x)\)
and the positive part of the linear current readout \(q_A(x)\).  The
certificate checks either their equality or the rational lower singular-value
bound which yields \eqref{eq:detector-contextual-amplitude-calibration} after
backward factorization.  The retained normalization denominator is a total
graph coordinate and is bounded away from zero on this cylinder, so the
constant does not change along the cascade.  The supplied modulus passes the
inequality to sourcewise closure.  Summing the nonnegative inequalities over
the positive detector records gives the final implication.
\end{proof}

\begin{theorem}[Head--support threshold consistency]
\label{thm:head-support-threshold-consistency}
The universal normalized shell floors are
\[
 \eta_{\rm head}=\frac15,
 \qquad
 \eta_{\rm supp}=\frac1{31}.
\]
The first applies to the five head currents of one finite shell record, and
the second applies to its thirty-one exact-support amplitudes.  A stronger
exact-support threshold may be used only on a cell carrying an additional
atomic-owner certificate.  If a selected support amplitude is normalized to a
probability or recurrent law, the compiler stores the unnormalised parent
amplitude \(\alpha_A(R)\) and its cumulative cocycle
\(\mathsf A_{A,N}\); denormalization is the exact identity
\[
 \alpha_A(R_j)=\mathsf A_{A,N}\,\nu_{A,N}(\{j\})
 \quad\text{for}\quad
 \nu_{A,N}:=\frac{\sum_{j\le N}\alpha_A(R_j)\delta_j}
                       {\mathsf A_{A,N}}
 \quad (\mathsf A_{A,N}>0).
\]
\end{theorem}

\begin{proof}
The five-head floor is
\cref{thm:deterministic-prospective-accountability}: five signed head
currents sum to one, so one positive head is at least \(1/5\).  By
\cref{thm:one-currency-exact-support-amplitude}, the thirty-one nonnegative
exact-support amplitudes have sum at least one, so one is at least \(1/31\).
The latter inference has no right to use the former larger threshold because
a head may split across several exact supports.  The displayed probability
formula is a definition when the denominator is positive; multiplying it by
the stored denominator recovers the original amplitude term by term.  Thus a
normalization changes no target-current mass.
\end{proof}

\begin{definition}[Canonical exact-support shell owner]
\label{def:canonical-exact-support-shell-owner}
Fix once and for all a total order \(\prec\) on the nonempty subsets of
\(\Channels\).  For a finite normalized shell record \(R\), set
\begin{equation}
 \label{eq:canonical-exact-support-owner}
 \operatorname{own}_{\rm supp}(R)
 :=\min\nolimits_{\prec}
 \left\{A\ne\varnothing:\alpha_A(R)\ge\frac1{31}\right\}.
\end{equation}
The owner is metadata attached to the original physical shell event; it does
not replace its complete joint record or discard the other support
amplitudes.
\end{definition}

\begin{theorem}[Canonical exact-support owner partition]
\label{thm:canonical-exact-support-owner-partition}
The owner in \eqref{eq:canonical-exact-support-owner} is defined and unique
for every finite normalized shell crossing.  On a pairwise disjoint shell
cascade, the events assigned to distinct supports are still pairwise
disjoint.  One fixed support owns infinitely many events in every infinite
cascade.  In particular, the assigned-support ledger debits each physical
shell event at most once.
\end{theorem}

\begin{proof}
Existence follows from
\cref{thm:head-support-threshold-consistency}; uniqueness follows from the
fixed total order.  Assignment changes only a finite label on the unchanged
shell interval, so pairwise disjointness is inherited from the canonical
cascade.  There are finitely many supports, and the infinite pigeonhole
principle gives one support with infinitely many assigned records.  Since an
event has one label, no two support ledgers can debit it as an assigned event.
\end{proof}

\begin{theorem}[Infinite-variation and terminal-owner calibration]
\label{thm:infinite-variation-terminal-owner-calibration}
Let \(R_\infty\) be a variation-boundary shell value or a terminal-jump shell
value.  Its joint source-complete shell graph determines one exact-support
owner \(A_\infty\ne\varnothing\), selected on a constant-owner subnet of its
finite normalized approximants.  There is an increasing
family of bounded exact-support amplitudes
\[
 \alpha_{A_\infty}^{[M]}(R_\infty)
 :=\lim_\gamma\bigl(\alpha_{A_\infty}(R_{i_\gamma})\wedge M\bigr),
 \qquad M\in\mathbb N,
\]
along a source-preserving approximating net \((R_i)\), and
\[
 \alpha_{A_\infty}^{[M]}(R_\infty)\uparrow
 \alpha_{A_\infty}^{[\infty]}(R_\infty)\in[1/31,+\infty].
\]
Every bounded member of this family is evaluated by the same
\(A_\infty\)-restricted conversion, carrier, storage, equality, and physical
price rows as a finite amplitude.  The terminal-jump case carries the one
same-origin noncontinuation event of \(R_\infty\); it is not replaced by a
sequence of distinct physical events.
\end{theorem}

\begin{proof}
The total shell-current graph compactifies all source coordinates jointly.
Before passage to the limit, retain the complete exact-support vector and
the canonical support owner of
\cref{def:canonical-exact-support-shell-owner}.  Finiteness of the support
lattice gives a cofinal subnet with constant owner \(A_\infty\).  Map this
subnet into the compact countable product \(\prod_{M\ge1}[0,M]\) of all
bounded truncation coordinates and take one further convergent subnet.  This
single subnet, denoted \((R_{i_\gamma})\), makes every displayed limit exist
simultaneously.  The pointwise inequalities between consecutive truncations
pass to the limit, so the limits increase in \(M\); the constant-owner
inequality gives the stated lower bound.  No eventual owner on the original
net is asserted.

At fixed \(M\), all current terms, carrier restrictions, endpoint rows, and
total-graph coordinates still have the one exact word \(A_\infty\).
Restricting the finite conversion and charging identities to that word gives
the asserted rows; carrier domination is established before monotone passage
in \(M\).  The completed terminal-restriction convention attaches every
terminal shell readout to the same origin and noncontinuation event.  It
therefore supplies one event coordinate, while the truncation index is only
an approximation coordinate.
\end{proof}

\begin{theorem}[Finite conditional-increment support accounting]
\label{thm:finite-exact-support-shell-accounting}
Let \((I_j,R_j)_{j\ge1}\) be pairwise disjoint normalized shell crossings
of one ordinary same-origin public branch, with physical ceiling \(E\).
For its canonical backward frontier, let
\[
 Q_{A,j}:=
 \sum_{\ell:\operatorname{Prov}(v_{j,\ell})=A}
 \widehat\Delta_{\Gamma_j,\ell}(y_{\Phi,\rho_j})\ge0
\]
be the total normalized conditional increment with exact support \(A\).
For \(N\ge1\), set
\[
 S_{A,N}:=\sum_{j=1}^{N}Q_{A,j},\qquad
 P_{A,N}:=\mathsf P_A(\mathcal F_{A,N}),\qquad
 \mathcal T_{A,N}:=(R_j,\operatorname{Inc}^{\Sigma}_A(R_j),
 p_{\Phys}(R_j))_{j\le N}.
\]
Then the public conditional-increment identity gives the exact finite relation
\begin{equation}
 \label{eq:cascade-exact-support-sum}
 N=\sum_{\varnothing\ne A\subseteq\Channels}S_{A,N}.
\end{equation}
Consequently, for every \(N\), at least one exact support satisfies
\(S_{A,N}\ge N/31\).  A nested choice of such supports has a fixed support
\(A_*\) and an unbounded subsequence \(N_k\) with
\[
 S_{A_*,N_k}\ge N_k/31.
\]
For every fixed \(A\) with \(\sup_NS_{A,N}=+\infty\), its successive integer
level-crossing prefixes are finite presentation-generated records and retain
the unnormalised \(A\)-cocycle.  This theorem concerns conditional-increment
coordinates only: it neither identifies \(Q_{A,j}\) with the canonical shell
amplitude \(\alpha_A(R_j)\) nor makes a terminal-orbit conclusion from those
prefixes.
\end{theorem}

\begin{proof}
Apply the unit shell-crossing identity to every \(I_j\), expand each current
by its exact source word before taking any limit, and sum the finite
identities.  Equivalently, group the nonnegative terms in
\eqref{eq:canonical-crossing-conditional-charge} by their exact provenance
words.  This proves
\eqref{eq:cascade-exact-support-sum}.  Since there
are thirty-one summands, the finite maximum is at least \(N/31\); an
infinite pigeonhole argument selects \(A_*\) on an unbounded subsequence.
For a prescribed support with unbounded \(S_{A,N}\), take instead its
successive first integer level-crossing prefixes.  Each is finite by
definition of an unbounded nondecreasing partial sum and retains its exact
provenance word.  No compactness or realization step is used.
\end{proof}

\begin{theorem}[Finite canonical-amplitude support accounting]
\label{thm:finite-canonical-amplitude-support-accounting}
Let \((I_j,R_j)_{j\ge1}\) be pairwise disjoint finite normalized shell
crossings of one ordinary same-origin public branch.  For each nonempty
exact support \(A\), set
\[
 \mathsf A_{A,N}:=\sum_{j=1}^N\alpha_A(R_j).
\]
Then
\begin{equation}
 \label{eq:canonical-amplitude-cascade-lower-bound}
 \sum_{\varnothing\ne A\subseteq\Channels}\mathsf A_{A,N}\ge N.
\end{equation}
Consequently, for every \(N\), some exact support satisfies
\(\mathsf A_{A,N}\ge N/31\); one fixed support satisfies this bound along an
unbounded subsequence.  All later exact-support price and terminalization
arguments use the unnormalised cocycle \(\mathsf A_{A,N}\), unless their
premises provide a displayed comparison with it.
\end{theorem}

\begin{proof}
Apply \cref{thm:one-currency-exact-support-amplitude} to each record and sum
the nonnegative inequalities over \(j\le N\).  This gives
\eqref{eq:canonical-amplitude-cascade-lower-bound}.  There are thirty-one
nonempty exact supports, so the maximum is at least \(N/31\).  The finite
pigeonhole principle selects one support on an unbounded subsequence.  The
last sentence fixes the ledger convention and does not infer a comparison
with conditional increments, detector norms, or physical price.
\end{proof}

\begin{theorem}[Sourcewise contact-route accounting]
\label{rem:unclosed-sourcewise-terminalization}
Fix a persistent exact-support route
\(A_0\subseteq A_1\subseteq\cdots\) from one same-origin contact cascade in an
engine instance generated by an effective public presentation.  If its
cumulative actionable canonical
amplitude is unbounded, then either a finite charged prefix exceeds the
generated ceiling on its affordable target-aligned carrier cell, or the
ordinary source-owned contact route remains represented in the actual-prefix
realization tree.  A target-bearing failure
of passivity, lower-carrier boundedness, exterior-input finiteness,
coherence, or carrier domination is a source-resolved successor belonging to
the second alternative; it is not discarded.  Storage, equality,
graph, support-transfer, and zero-price re-entry are likewise internal
transitions of this dichotomy.
\end{theorem}

\begin{proof}
By \cref{thm:engine-target-aligned-public-closure}, the route is the target
quotient of one cumulative same-origin public joint record: its stratum,
shell covector, origin, exact provenance word, physical-event indices, and
carrier owner remain fixed data of every transition.  Apply the uniform
exact-support conversion ledger to that record.  On the affordable cell,
accumulate the carrier-admitted prices charged only through
\eqref{eq:target-aligned-charge-rule} to the canonical public event indices.
If their cumulative value exceeds the generated ceiling, the least such
finite prefix gives the first alternative.  Otherwise the original contact
route already supplies the second alternative.  If the ceiling-generating
carrier rows fail, exhaustive
target-aligned carrier stratification closes their target-null projection or
routes their least nonzero five-source owner into the same successor queue;
the actual-prefix tree then processes it toward the second alternative.  An
unbounded transfer re-enters
through \cref{thm:storage-feeding-reentry}, and a nonzero equality or graph
remainder re-enters through
\cref{thm:equality-graph-remainder-reentry,thm:fidelity-defect-successor}.
Strict support growth is routed to its canonical superset.  If successors
continue, extend the same cumulative joint record before reapplying the
target quotient or exact-support projection.  By
\cref{thm:infinite-chain-no-emergence-trichotomy}, target-nullity contradicts
the retained amplitude and a finite unprocessed increment continues the
queue.  In the fixed case,
\cref{thm:fixed-atom-realization-or-successor} supplies a finite obstruction,
a successor, or an ordinary orbit.  The route is seeded by one actual
same-origin contact cascade.  Its restrictions belong to the
presentation-generated prospective source domain and provide an accepted
child at every finite level of the actual-prefix tree; checker soundness
therefore excludes the obstruction alternative.  A nonordinary successor
contradicts fixedness and returns to the queue.  The remaining ordinary path
supplies the represented contact route in the second alternative after the
public continuation test.  This recovers the orbit already postulated by the
contact-cascade hypothesis; it does not close the route and is not a
regularity proof.  All
re-entries retain the exact source word and the original public dependency
record, so no conditional increment is substituted for canonical amplitude
without its checked calibration and no target-relative proof cone or
realization rule is introduced.
\end{proof}

\begin{theorem}[Support-lattice stabilization]
\label{thm:support-lattice-terminalization}
Every sourcewise execution has finitely many strict support enlargements.
After the last such enlargement, all transitions remain in one exact support
and are governed by resolved-prefix, carrier--polar, and realization
refinement.
\end{theorem}

\begin{proof}
Provenance supports are nonempty subsets of the five-element channel set.
Every strict transfer appends a source letter, and cumulative joining cannot
forget one.  Thus at most four strict enlargements follow a nonempty initial
support.  Thereafter
\cref{thm:proper-successor-monotonicity} governs the fixed-support route.
\end{proof}

\begin{theorem}[Terminal-entry sourcewise disposition]
\label{thm:terminal-entry-sourcewise}
A finite noncontinuation event with a \(\Bdry\)-headed owner has exactly one
disposition: ordinary continuation, contradicting its status; a target-null
or class-incompatible terminal current; or an independently realized
singular terminal event.  It is one completed physical event, not an
infinite family of disjoint charges.
\end{theorem}

\begin{proof}
Evaluate the boundary conversion and total continuation graph directly on the
completed endpoint.  Totalization of trace, terminal flux, and germ
realization gives the first failed successor when a row is nonordinary.
Otherwise the continuation-fidelity terminal test gives the stated three
ordinary alternatives.  Repeated descriptions of the same endpoint preserve
one \(\Bdry\)-headed owner and do not create new physical events.
\end{proof}

\begin{theorem}[Uniform thirty-one-support contact accounting]
\label{thm:uniform-thirty-one-support-terminalization}
\cref{rem:unclosed-sourcewise-terminalization} applies uniformly to every
nonempty \(A\subseteq\Channels\).  Its proof is one schema over the free
provenance semilattice, instantiated with the exact word metadata; it is not
a collection of thirty-one unrelated arguments.
\end{theorem}

\begin{proof}
The conversion, transfer, carrier, no-emergence, and realization theorems are
parametric in \(A\).  Structural induction on provenance union preserves the
exact word at every constructor.  Instantiating the resulting schema over the
thirty-one nonempty subsets gives the claim.
\end{proof}

\begin{theorem}[Continuum target-aligned structural exhaustion]
\label{thm:continuum-target-structure-exhaustion}
On every nonempty ordinary finite physical stratum,
\begin{equation}
 \label{eq:contact-formation-profile-equivalence}
 \boxed{\begin{gathered}
 \Reach_\Sigma^{\Phys}(E)>0\\
 \Longleftrightarrow\
 \exists\,\varnothing\ne A\subseteq\Channels:
 \mathfrak F_A^\Sigma(E)
 \text{ contains a realized}\\[-1mm]
 \text{formation-active point}.
 \end{gathered}}
\end{equation}
There are exactly \(2^5-1=31\) possible named nonempty supports; a particular
equation may realize fewer of them.  Within each cell, the five
normal-form theorems give an exhaustive atomic word after the fixed
projection order
\begin{equation}
 \label{eq:continuum-normal-form-order}
 \Geom\longrightarrow\Caus\longrightarrow\Abs
 \longrightarrow\Lift\longrightarrow\Bdry,
\end{equation}
Within one source family the first nonzero typed coordinate is unique.
Different source families may remain simultaneously active and are retained
by \(A\), the aggregate carrier class, and the circulation graph.  No sixth
target-bearing family exists.
\end{theorem}

\begin{proof}
On a cofinal approach-contact branch, contact gives a disjoint unit-increment
shell cascade.  Finite support exhaustion selects one nonempty \(A\) on a
subsequence.  On a terminal-entry branch, the completed shell graph instead
gives a \(\Bdry\)-headed owner and the terminal-entry coordinate; take its
exact provenance support as \(A\).  Target-incidence compression retains
precisely these source--shell paths.  Apply the five
normal-form theorems in the order
\eqref{eq:continuum-normal-form-order}; every headed ordinary or nonordinary
atomic component enters exactly one block, while composite ancestry and
signed interactions remain in the exact support, aggregate carrier, and
circulation coordinates.  On a finite current record the affine shell
identity supplies a positive component; on a variation-boundary record the
closed owner supplies it.  The resulting profile is formation active.
Same-origin diagonal realization gives the required profile point.

Conversely, a realized formation-active point includes a public same-origin
path and nonzero cofinal shell coordinate.  Its shell coordinate therefore
crosses every retained neighborhood of the complete continuation target,
which is contact.  The universal provenance theorem and
\cref{thm:no-additional-continuum-singularity-source} place every ordinary,
defective, or limiting shell-current occurrence in one of these cells.
Since \(J_{\Res}\subseteq\Null_\Sigma\), no additional target-bearing family
remains.
\end{proof}


\begin{definition}[Presentation-generated prospective source domain]
\label{def:presentation-generated-prospective-source-domain}
Fix a public continuum presentation, a derived finite physical stratum
\(\Budget\), and its generated target shells.  For each nonempty
\(A\subseteq\Channels\), let
\(\mathfrak D_{A,\Sigma}^{\rm pre}(\Budget)\) contain the following records:
\begin{enumerate}[label=\textup{(D\arabic*)}]
\item every finite stopped history of an ordinary same-origin branch, cut
  before its first contact with the current shell;
\item every finite joint source occurrence obtained from such a history by
  the five source constructors, target projection, backward pullback,
  conditional joining, chart or quotient transport, signed interface gluing,
  and physical-carrier restriction;
\item every completed terminal-entry record, with its \(\Bdry\)-headed current
  and continuation-graph value; and
\item every finite prefix of a source-preserving limit, concentration,
  oscillation, tail, equality, or total-graph record generated from
  \textup{(D1)}--\textup{(D3)}.
\end{enumerate}
All records retain their origin, stopped time, target covector, complete
preceding frontier, unnormalized occurrence cocycle, source word, and physical
carrier.  Let \(\mathfrak D_{A,\Sigma,n}^{\rm pre}(\Budget)\) be the finite
subfamily generated by the first \(n\) public constructor words, rational
shell probes, and stopped-history prefixes in the fixed grammar enumeration.
Then \(\mathfrak D_{A,\Sigma,n}^{\rm pre}(\Budget)\) is increasing and its
union is \(\mathfrak D_{A,\Sigma}^{\rm pre}(\Budget)\).  With \(\iota_A\) and \(K_A\) from
\cref{def:physical-saturated-algebra}, define
\begin{equation}
 \label{eq:presentation-generated-prospective-source-domain}
 \mathscr S_{A,\Sigma}^{\rm reach,sat}(\Budget)
 :=\overline{\iota_A\bigl(\mathfrak D_{A,\Sigma}^{\rm pre}(\Budget)\bigr)}^{K_A}.
\end{equation}
The full domain is the provenance-disjoint union over nonempty \(A\).
Evaluation at a point of
\(\mathscr S_{A,\Sigma}^{\rm reach,sat}(\Budget)\) is called a
\emph{reachable source evaluation}.  The continuum structural algebra has no
second semantic domain: every source, invariant, price, and contact statement
is evaluated on \eqref{eq:presentation-generated-prospective-source-domain}.
\end{definition}

\begin{theorem}[Public-presentation closure of the structural calculus]
\label{thm:public-presentation-structural-proof-closure}
Every relation used in the continuum structural calculus is generated by one
of the following operations on the public presentation:
\begin{enumerate}[label=\textup{(S\arabic*)}]
\item a weak-equation, initial-data, boundary, continuation, target-shell, or
physical-carrier coordinate declared by the presentation;
\item a finite five-source constructor applied to earlier generated records;
\item a target projection, complete backward pullback, conditional joint
increment, signed interface gluing, or physical-carrier restriction of such
a record; or
\item a source-preserving limit of the resulting finite stopped records.
\end{enumerate}
Consequently every theorem in the structural reachability calculus is a
statement about \(\mathscr S_{A,\Sigma}^{\rm reach,sat}(\Budget)\), its
target-aligned coordinates, and the physical carrier generated by the public
presentation.  A global analytic property contributes only through the
target-projected total-graph coordinate generated by the operation that
asserts or fails it.
\end{theorem}

\begin{proof}
The primitive source compiler exhausts the public weak-expression graph by
\cref{lem:primitive-current-edge-exhaustion}.  The free provenance algebra
closes those primitive records under every finite operation in
\textup{(S2)}--\textup{(S3)}, and
\cref{thm:source-preserving-completion} proves \textup{(S4)} without adding a
source.  Prospective backward factorization and conditional-increment
accounting retain the complete joint frontier before source projection.
The physical carrier is restricted only along the same generated occurrence
record.  Structural induction on the constructor word therefore places every
premise and every conclusion of an invariant, price, equality, Hamiltonian,
entropy, or contact calculation in the stated reachable domain.
\end{proof}

\begin{theorem}[Reachable saturation and finite occurrence recovery]
\label{thm:reachable-source-saturation-finite-occurrence}
\label{thm:five-source-invariant-verification-protocol}
The set \(\mathscr S_{A,\Sigma}^{\rm reach,sat}(\Budget)\) is compact,
retains source support \(A\), and is closed under every source-preserving
constructor generated by the public presentation.  Let \(d\ge0\) be any
target-incidence detector in the complete invariant ledger.  Then
\begin{align}
 \sup_{R\in\mathscr S_{A,\Sigma}^{\rm reach,sat}(\Budget)}d(R)=0
 &\Longleftrightarrow
 d(R)=0\quad
 \text{for every }R\in\mathfrak D_{A,\Sigma}^{\rm pre}(\Budget),
 \label{eq:actual-source-zero-equivalence}\\
 \sup_{R\in\mathscr S_{A,\Sigma}^{\rm reach,sat}(\Budget)}d(R)>0
 &\Longleftrightarrow
 \exists R_0\in\mathfrak D_{A,\Sigma}^{\rm pre}(\Budget):d(R_0)>0.
 \label{eq:actual-source-positive-equivalence}
\end{align}
In the positive case, prospective backward factorization replaces \(R_0\),
when necessary, by its least nonzero finite contextual current--probe
occurrence and retains the complete preceding frontier.  Hence target-visible
source presence is always witnessed by presentation-generated finite ancestry.
\end{theorem}

\begin{proof}
Compactness follows because \(K_A=[-1,1]^{\mathbb N}\) is compact and the
right side of \eqref{eq:presentation-generated-prospective-source-domain} is
closed.  Source preservation is
\cref{thm:source-preserving-completion}; every listed constructor is already a
coordinate of the structural test system.  Each detector is a nonnegative
continuous coordinate of that system, after adjoining its bounded rational
truncations.  Density of
\(\iota_A(\mathfrak D_{A,\Sigma}^{\rm pre}(\Budget))\) in its closure proves
\eqref{eq:actual-source-zero-equivalence}.  If the supremum is positive,
choose a point where \(d>0\); continuity gives a neighborhood on which
\(d>0\), and density supplies \(R_0\).  The converse is immediate.  Exact
backward factorization gives the final assertion.
\end{proof}

\begin{theorem}[Finite contextual occurrence recovery with amplitude retention]
\label{thm:finite-contextual-occurrence-amplitude-retention}
Fix an exact support \(A\) and a calibrated conditional-incidence detector
\(d_A\).  If
\[
 \sup_{R\in\mathscr S_{A,\Sigma}^{\rm reach,sat}(\Budget)}d_A(R)>0,
\]
then the engine returns a finite same-origin pre-contact contextual occurrence
\[
 (R_0,\lambda,q_A(R_0),\alpha_A(R_0),
   \operatorname{Anc}(R_0),p_{\Phys}(R_0)).
\]
For the canonical detector, the returned record satisfies
\(d_A^{\rm cur}(R_0)=\alpha_A(R_0)>0\).  For a contextual calibration, the
returned record satisfies
\[
 \alpha_A(R_0)\ge\kappa_A d_A(R)>0
\]
for the positive finite occurrence \(R\) from which backward factorization
selects \(R_0\).  The constant \(\kappa_A\) is the uniform constant of the
detector certificate.
\end{theorem}

\begin{proof}
By \cref{thm:reachable-source-saturation-finite-occurrence}, positivity on
the reachable closure gives a finite presentation-generated occurrence
\(R\) with \(d_A(R)>0\), and backward factorization returns its least
same-origin contextual current--probe occurrence.  For the canonical
detector this occurrence is the current record itself.  In the contextual
case apply \cref{thm:detector-to-amplitude-calibration}; its certified
factorization map gives \(R_0\) and the displayed lower bound.  The
factorization transcript retains the target covector, complete preceding
frontier, ancestry, and restricted physical carrier, which are exactly the
listed fields.
\end{proof}

\begin{definition}[Exact reachable invariant profile]
\label{def:exact-reachable-invariant-profile}
For a detector coordinate \(d\) owned by \(A\), define
\begin{equation}
 \label{eq:exact-reachable-detector-profile}
 \mathcal V_{A,d}^{\Sigma}(\Budget)
 :=\max_{R\in\mathscr S_{A,\Sigma}^{\rm reach,sat}(\Budget)}d(R).
\end{equation}
For a signed scalar, graph, measure, operator, response, shell, or realization
invariant, use the squared magnitudes of its separating rational cylinder
coordinates.  For a nonnegative extended-real valuation \(v\), use all
bounded truncations \(v\wedge m\) and define its exact value by their
compatible maxima or minima on the same reachable set.  Rank is evaluated by
finite minors; concentration, oscillation, and boundary variations are
evaluated by finite rational partitions.  Mixed-source invariants are
evaluated on the complete joint record before sourcewise projection.
\end{definition}

\begin{theorem}[Complete target-aligned evaluation from the public presentation]
\label{thm:actual-source-invariant-decision}
\label{thm:invariant-complete-five-source-decision-ledger}
\label{thm:coordinatewise-source-invariant-decision}
\label{cor:total-binary-five-source-evaluation}
Every invariant generated by the five normal forms---geometric
range/kernel/graph/capacity, directed cohomology/storage/response/cycle,
quotient descent/fiber/cokernel/rank, lift interface/concentration/oscillation/
scale, boundary exterior/tail/terminal data---and every joint carrier,
conditional-increment, shell, realization, entropy, Hamiltonian, or
physical-price coordinate
has an exact value on
\(\mathscr S_{A,\Sigma}^{\rm reach,sat}(\Budget)\).  For every detector
coordinate,
\begin{equation}
 \label{eq:actual-invariant-binary-decision}
 \operatorname{Eval}_{A,d}^{\Sigma}
 =
 \begin{cases}
  \operatorname{Absent}_{A,d},&\mathcal V_{A,d}^{\Sigma}(\Budget)=0,\\
  \operatorname{Present}_{A,d}(R_0,\lambda,Q,\Anc,p_{\Phys}),
    &\mathcal V_{A,d}^{\Sigma}(\Budget)>0,
 \end{cases}
\end{equation}
where the second row contains the finite occurrence recovered by
\cref{thm:reachable-source-saturation-finite-occurrence}.  Every valuation
returns its exact extended-real value and an attaining reachable record when
the extremum is finite, or a cofinal sequence of reachable finite occurrences
when it is infinite.  Every graph-valued invariant returns its unique
cylinder-coordinate value on each reachable record.

The second row is the only positive structural value.  Its record contains
the public equation term, target covector, complete joint frontier, inherited
five-source word, and physical carrier atom that generate the positive
coordinate.  Thus \eqref{eq:actual-invariant-binary-decision} evaluates every
invariant directly in the presentation-generated algebra.
\end{theorem}

\begin{proof}
Compactness gives the maximum in
\eqref{eq:exact-reachable-detector-profile};
\cref{thm:reachable-source-saturation-finite-occurrence} gives the two rows of
\eqref{eq:actual-invariant-binary-decision}.  Rational cylinder coordinates
separate points of the fixed graph-complete product, so they determine every
graph-valued invariant.  Bounded truncations determine an extended-real
valuation uniquely.  The finite-minor and rational-partition criteria are
exact by finite-dimensional rank and total variation.  Conditional joint
increments are formed before projection, so mixed target alignment remains a
coordinate of the reachable joint record.  These arguments apply to every row
of the coordinate ledger.
\end{proof}

\begin{theorem}[Actual physical-price and equality decision]
\label{thm:actual-source-price-equality-decision}
Let \(R_0\) be a positive source occurrence returned by
\eqref{eq:actual-invariant-binary-decision}.  Restrict the same-origin physical
carrier to its finite occurrence atom.  Exactly one of the following holds:
\begin{enumerate}[label=\textup{(P\arabic*)}]
\item \(p_{\Phys}(R_0)>0\), and the occurrence is charged by that exact
  physical amount;
\item \(p_{\Phys}(R_0)=0\), and the complete occurrence belongs to the
  reachable zero-price equality cell generated by the same source word.
\end{enumerate}
In the second case, equality saturation is performed inside
\(\mathscr S_{A,\Sigma}^{\rm reach,sat}(\Budget)\).  Reprojection gives
either target absence by \eqref{eq:actual-source-zero-equivalence} or another
actual positive finite occurrence by
\eqref{eq:actual-source-positive-equivalence}.  Equality saturation therefore
does not create a formal terminal state.
\end{theorem}

\begin{proof}
The physical carrier is a nonnegative measure, so its value on the finite
occurrence atom is positive or zero.  The first value is the local charging
row.  In the zero case, every nonnegative production summand vanishes on that
atom and is adjoined to the generated equality ideal.  Source-preserving
saturation keeps the resulting descendants in the same reachable domain.
Apply
\cref{thm:reachable-source-saturation-finite-occurrence,thm:actual-source-invariant-decision}
after target reprojection.  This gives the stated exhaustive decision.
\end{proof}

\begin{theorem}[Cofinal contact retention without dilution]
\label{thm:actual-cofinal-contact-retention}
Let one ordinary same-origin branch with finite physical budget complete a
cofinal target-shell family.  Form its countable occurrence space before any
normalization.  Then one of the following actual source records exists:
\begin{enumerate}[label=\textup{(C\arabic*)}]
\item a fixed source owns finite pre-contact occurrence sets on which the
  target-current density relative to the physical carrier is unbounded;
\item a fixed source owns a positive target-current atom of zero physical
  carrier mass.
\end{enumerate}
Every record in either row belongs to
\(\mathscr S_{A,\Sigma}^{\rm reach,sat}(\Budget)\) for its inherited support
\(A\), and its invariant coordinates are evaluated by
\cref{thm:actual-source-invariant-decision,thm:actual-source-price-equality-decision}.
Thus infinitely many finite contributions cannot converge to target progress
without producing an actual positive source occurrence.
\end{theorem}

\begin{proof}
Apply exact compositional accounting
\cref{thm:exact-compositional-physical-accounting}.  Finite physical mass and
infinite unnormalized crossing mass give the amplification alternative
\eqref{eq:finite-budget-amplification-alternative}.  Prospective
source-amplification extraction assigns either value to one fixed source and
returns its finite pre-contact occurrence sets.  These are elements of
\(\mathfrak D_{A,\Sigma}^{\rm pre}(\Budget)\) by definition.  The final assertion
follows from the actual invariant and price decisions.
\end{proof}

\begin{definition}[Actual structural contact cells]
\label{def:actual-structural-contact-cells}
Let \(\mathscr C_{A,\Sigma}(\Budget)\) be the subset of
\(\mathscr S_{A,\Sigma}^{\rm reach,sat}(\Budget)\) whose total-graph
coordinates have all of the following values:
\begin{enumerate}[label=\textup{(K\arabic*)}]
\item one ordinary same-origin equation and initial-data value;
\item the cofinal unnormalized target-shell cocycle;
\item the noncontinuable value of the declared continuation graph; and
\item a nonzero target projection with its complete five-source ancestry.
\end{enumerate}
In addition, the complete record must lie in the image of the ordinary
same-origin realization map generated by the public equation and initial
data.  Thus
\begin{equation}
 \label{eq:actual-contact-cell-realization-intersection}
 \mathscr C_{A,\Sigma}(\Budget)
 =\mathscr S_{A,\Sigma}^{\rm reach,sat}(\Budget)
  \cap\operatorname{Ran}(\mathfrak R_{\rm ord})
  \cap\bigcap_{i=1}^{4}\{\textup{(K\(i\))}\}.
\end{equation}
These conditions are coordinates of the public total graphs and generated
shell family.  The cell is evaluated inside the reachable source domain, not
inside a finite-cylinder relaxation.  In particular, a compatible positive
functional, graph-boundary point, or zero-price equality model is not a
member unless one ordinary same-origin public orbit realizes the entire
joint record.
\end{definition}

\begin{theorem}[Presentation-relative structural decision of global regularity]
\label{thm:structural-global-regularity-decision}
For every derived finite physical stratum,
\begin{align}
 \Reach_\Sigma^{\Phys}(\Budget)=0
 &\Longleftrightarrow
 \mathscr C_{A,\Sigma}(\Budget)=\varnothing
 \quad\text{for every nonempty }A\subseteq\Channels,
 \label{eq:structural-regularity-empty-contact-cells}\\
 \Reach_\Sigma^{\Phys}(\Budget)>0
 &\Longleftrightarrow
 \mathscr C_{A,\Sigma}(\Budget)\ne\varnothing
 \quad\text{for some nonempty }A\subseteq\Channels.
 \label{eq:structural-singularity-nonempty-contact-cell}
\end{align}
The emptiness or nonemptiness of each contact cell is evaluated by the exact
coordinate rules of \cref{thm:actual-source-invariant-decision}: zero target
incidence deletes the cell, while a positive value returns an actual finite
pre-contact source occurrence; positive physical price is charged and zero
price enters the actual equality decision.  A cofinal family is retained by
\cref{thm:actual-cofinal-contact-retention}.  Hence the global-regularity
process uses only target-aligned, presentation-generated structural records.
Its terminal outputs are
\begin{equation}
 \label{eq:structural-global-regularity-terminal-values}
 \operatorname{GlobEval}_\Sigma(\Budget)
 \in\left\{
  \operatorname{Regular}\bigl((\mathscr C_{A,\Sigma}=\varnothing)_A\bigr),
  \operatorname{Singular}(R,A);:
    R\in\mathscr C_{A,\Sigma}(\Budget)
 \right\}.
\end{equation}
The terminal map is evaluated on ordinary same-origin contact cells with their
complete presentation-generated ancestry.
\end{theorem}

\begin{proof}
An ordinary contacting branch supplies its completed terminal or cofinal
pre-contact records.  Prospective source exhaustion gives a nonempty support
\(A\), while source-preserving saturation and the total continuation graph
give \textup{(K1)}--\textup{(K4)}.  Hence contact implies a nonempty contact
cell, and the branch itself supplies membership in
\(\operatorname{Ran}(\mathfrak R_{\rm ord})\).  Conversely,
\eqref{eq:actual-contact-cell-realization-intersection} supplies one ordinary
same-origin branch with the declared equation, initial datum, cofinal-shell
cocycle, and noncontinuation value; that branch contacts the declared target.
This proves
\eqref{eq:structural-regularity-empty-contact-cells}--
\eqref{eq:structural-singularity-nonempty-contact-cell}.

The actual invariant decision evaluates each source coordinate on the
reachable set.  The price/equality theorem processes its physical value, and
the cofinal theorem prevents dilution across infinitely many crossings.
Universal five-source provenance and target-null residual closure exhaust the
owners.  Thus the two values in
\eqref{eq:structural-global-regularity-terminal-values} are exhaustive.
\end{proof}

\begin{theorem}[Simultaneous terminal evaluation of the complete source ledger]
\label{thm:simultaneous-source-ledger-terminal-evaluation}
For every nonempty support \(A\subseteq\Channels\), saturate the complete
coordinate ledger under all mandatory successors and evaluate it on
\(\mathscr S_{A,\Sigma}^{\rm reach,sat}(\Budget)\).  The result is exactly
\begin{equation}
 \label{eq:simultaneous-support-terminal-value}
 \operatorname{Eval}^{\Sigma}_A
 =
 \begin{cases}
  \operatorname{Absent}_A,
    &\mathcal V_{A,d}^{\Sigma}(\Budget)=0
      \text{ for every target-incidence detector }d,\\
  \operatorname{Present}_A(R_0,\lambda,Q,\Anc,p_{\Phys}),
    &\mathcal V_{A,d}^{\Sigma}(\Budget)>0
      \text{ for the first positive detector }d.
 \end{cases}
\end{equation}
Every successor is already a coordinate of the saturated reachable domain.
The second value contains an actual presentation-generated finite occurrence
\(R_0\).  Boundary, cofinal, equality, and nonordinary data are coordinates
of \(R_0\) and its complete ancestry.
\end{theorem}

\begin{proof}
Apply \cref{thm:actual-source-invariant-decision} simultaneously in the fixed
ledger order.  The source alphabet is finite and the coordinate order is a
well-order, so a nonempty set of positive coordinates has a first member.
Source-preserving saturation places every successor in the same reachable
domain before evaluation.  Universal provenance and target-null residual
closure leave no target-bearing coordinate outside the five supports.
\end{proof}

\begin{theorem}[Direct target-local invariant refinement]
\label{thm:direct-continuum-invariant-protocol}
For each rational \(E>0\), support \(A\), and refinement depth \(n\), the
following finite target-local procedure constructs a finite prefix of the
presentation-generated formation profile.
\begin{enumerate}[label=\textup{(P\arabic*)}]
\item Generate the first \(n\) forward source words and backward shell
covectors, retain their conservative incidence core, and form the saturated
cumulative join with every earlier prefix.  Discard all coordinates outside
that joint source--target incidence core.
\item Apply the five canonical normal-form projections in
\eqref{eq:continuum-normal-form-order}.  Evaluate, in the displayed order of
\cref{thm:coordinatewise-source-invariant-decision}, the geometric range,
kernel, graph and capacity rows; the directed cohomology, storage, domination,
cycle, response and graph rows; the quotient descent, vertical, cokernel,
graph, target, kernel-capacity and rank rows; the lift intertwining, fiber,
concentration, oscillation, scale and graph rows; and the exterior, tail,
terminal, graph and tail-capacity rows.
\item Evaluate the mixed carrier, conditional-increment, cycle, realization,
shell, entropy, and Hamiltonian rows for every nonempty support.  Use the
exact zero/positive, finite-minor, variation, and cylinder-coordinate rule
assigned to each generated coordinate by
\cref{thm:coordinatewise-source-invariant-decision}.
\item Return the generated finite stopped records with their joint invariant,
origin, physical carrier, and total-realization coordinates.  Preserve all
finite correlations imposed by the ordered rows, and refine the least
target-visible coordinate whose zero/nonzero or ordinary/nonordinary status
can change formation activity.
\end{enumerate}
Then
\begin{equation}
 \label{eq:formation-profile-computation}
\overline{\bigcup_n\iota_A\bigl(mathfrak D_{A,\Sigma,n}^{\rm pre}(E)\bigr)}
=\mathscr S_{A,\Sigma}^{\rm reach,sat}(E).
\end{equation}
The source decision is made on
\(\mathscr S_{A,\Sigma}^{\rm reach,sat}(E)\).  More precisely, for every
detector \(d\),
\begin{equation}
 \label{eq:direct-refinement-actual-decision}
 \max_{R\in\mathscr S_{A,\Sigma}^{\rm reach,sat}(E)}d(R)=0
 \quad\text{or}\quad
 \exists R_0\in\mathfrak D_{A,\Sigma}^{\rm pre}(E):d(R_0)>0.
\end{equation}
The alternatives are exactly the absence and actual-presence rows of
\cref{thm:actual-source-invariant-decision}.  In the positive row, backward
factorization returns the least finite contextual occurrence and recursive
carrier--polar operations act on that occurrence.  In the zero row the
corresponding target-incidence edge is deleted.  Contact is decided only by
the realized cells of \cref{def:actual-structural-contact-cells}.
\end{theorem}

\begin{proof}
Finite incidence pruning is \cref{cor:finite-incidence-fast-track}.  Each
normal-form projection is a rational finite-dimensional projection or a
totalized polarity split at depth \(n\).  The invariant definitions are
quadratic-form, order-domination, variation, rank, capacity, circulation, or
graph-feasibility programs, all applied to generated finite stopped records.
Exhaustive enumeration, source-preserving closure, retention of all finite
correlations, and the five measurement theorems give
\eqref{eq:formation-profile-computation}.  Apply
\cref{thm:reachable-source-saturation-finite-occurrence}.  Continuity and
density give \eqref{eq:direct-refinement-actual-decision}; exact backward
factorization supplies the least contextual occurrence in its positive row.
The realized-contact intersection is
\eqref{eq:actual-contact-cell-realization-intersection}, which proves the last
assertion.
\end{proof}

\begin{figure}[tbp]
\centering
\resizebox{\linewidth}{!}{%
\begin{tikzpicture}[fw diagram,x=1cm,y=1cm]
  \node[fw source,text width=2.3cm] (g) at (0,2.8)
    {\(\Geom\)\\range, kernel, capacity};
  \node[fw source,text width=2.3cm] (c) at (0,1.4)
    {\(\Caus\)\\cohomology, cycles, response};
  \node[fw provenance,text width=2.3cm] (a) at (0,0)
    {\(\Abs\)\\descent, fiber, cokernel};
  \node[fw provenance,text width=2.3cm] (l) at (0,-1.4)
    {\(\Lift\)\\concentration, scale, graph};
  \node[fw provenance,text width=2.3cm] (b) at (0,-2.8)
    {\(\Bdry\)\\exterior, tail, terminal};
  \node[fw geom,text width=2.7cm] (supports) at (4.2,0)
    {\(31\) nonempty\\source-support cells \(A\)};
  \node[fw use,text width=2.8cm] (profile) at (8.1,0)
    {nested formation profiles\\
     \(\mathfrak F_{A,n}^{\Sigma}(E)\)};
  \node[fw terminal,text width=2.7cm] (regular) at (12.2,1.6)
    {reachable source profiles zero\\structural absence};
  \node[fw use,text width=2.7cm] (joint) at (12.2,-1.6)
    {cumulative joint word\\saturate and reproject};
  \node[fw terminal,text width=2.7cm] (singular) at (16.0,-1.6)
    {actual finite pre-contact\\source occurrence};
  \foreach \x in {g,c,a,l,b}
    \draw[fw arrow] (\x)--(supports);
  \draw[fw arrow] (supports)--(profile);
  \draw[fw arrow] (profile)--(regular);
  \draw[fw arrow] (profile)--(joint);
  \draw[fw dashed arrow] (joint) to[bend left=24] (profile);
  \draw[fw arrow] (joint)--(singular);
\end{tikzpicture}%
}
\caption{Direct continuum formation analysis.  Each channel has a canonical
target-local normal form and a finite invariant program.  Their simultaneous
supports produce \(31\) cells.  Outer refinement is intersected with the
presentation-generated reachable domain.  Zero on that domain proves
structural absence; positivity recovers an actual finite pre-contact source
occurrence with its complete joint frontier.  Joint reprojection returns every
newly exposed target component to its five-channel owner.  Physical contact
is read only on the ordinary same-origin realization range.}
\label{fig:continuum-formation-profile}
\end{figure}

\FloatBarrier

\begin{definition}[PDE backend]
\label{def:pde-backend}
A \emph{PDE backend} is the homomorphic compilation of a public presentation
into the source, shell, carrier, total-graph, successor, and ordered-relation
signatures used to construct \(\mathfrak G_{\mathfrak b}\).  It may compute
equation-specific coefficients, scaling weights, constraint tails, and
closure edges.  It may not add a source family, a target-relative premise, a
physical measure not generated by the equation, or a terminal declaration
without a proved relation of the ordered algebra.
Here \(\mathfrak G_{\mathfrak b}\) is the induced subgraph of
\(\mathfrak G_A\) obtained by fixing the complete active signature
\(\mathfrak b=(A,m,\Gamma,\eta,\mathbf R)\) and its earlier zero coordinates.
\end{definition}

\begin{proposition}[Backend conservativity]
\label{prop:backend-conservativity}
Every PDE backend preserves the five-channel normal form, prospective contact
accountability, source-preserving limit closure, the target-null meaning of
\(J_{\rm res}\), and the single physical-budget ledger.  Therefore a backend
can validate or refute regularity only by evaluating the vertices and closure
edges of \(\mathfrak G_{\mathfrak b}\); it cannot change the continuum theory that
generated them.
\end{proposition}

\begin{proof}
The backend is a homomorphism into the signature already constructed in
\cref{sec:pde-instantiation}.  Preservation of source normal forms follows
from \cref{thm:part-i-continuum-preservation}; preservation of contact
accountability and the residual quotient follows from
\cref{thm:unconditional-continuum-contact-alternative}; preservation of the
physical ledger follows from signed carrier gluing and
\cref{thm:carrier-physical-budget}.  The prohibited additions are not symbols
of the target signature, so they cannot occur in its image.
\end{proof}



\section{Navier--Stokes public-presentation compilation audit}
\label{sec:end-to-end-public-unit-test}

This audit verifies the specialization interface on the three-dimensional
incompressible Navier--Stokes presentation.  Every displayed object below is
compiled from the public weak equation, its local-energy carrier, the
continuation graph, and the target shells.  It does not assert a global
regularity conclusion for this presentation.

\begin{definition}[Public Navier--Stokes target presentation]
\label{def:public-ns-target-presentation}
Fix \(\nu>0\), divergence-free initial data, declared boundary data, and the
public weak relation
\[
 \partial_tu+\mathbb P((u\cdot\nabla)u)-\nu\Delta u=0,
 \qquad \nabla\cdot u=0,
\]
together with the public multiplier grammar and the continuation graph for
the declared regularity class.  The grammar derives the stopped kinetic- and
local-energy identities by testing the displayed equation; an energy bound
is not an additional input.  Let \(\Sigma_{\rm NS}\) be the set of
terminal germs at which that graph has no ordinary continuation.  The only
physical ledger is the local-energy carrier obtained from the same identity.
All cutoffs, moving charts, scale records, pressure reconstructions, and
shell covectors are generated by the continuum constructor grammar.
\end{definition}

\begin{proposition}[Generated Navier--Stokes carrier admission]
\label{prop:ns-generated-carrier-admission}
On every finite stopped ordinary Navier--Stokes branch, the public multiplier
word \(u\) generates the kinetic-energy balance.  On \(\mathbb R^3\), the
periodic domain, or a bounded domain with no-slip boundary conditions, its
internal convection and pressure contributions cancel after signed gluing,
and the admitted carrier stratum has
\[
 K_u(t)=\frac12\lVert u(t)\rVert_{L^2}^2,
 \qquad
 \mu_{\rm NS}((s,t])
 =\nu\int_s^t\lVert\nabla u(r)\rVert_{L^2}^2\,\dd r
   +\mathfrak d_{\rm en}((s,t]),
\]
where \(\mathfrak d_{\rm en}\ge0\) is the totalized weak energy-defect value
and vanishes on a smooth branch.  With legitimate positive exterior work
\(W^+\), the compiler derives
\[
 \mu_{\rm NS}([0,T))
 \le \frac12\lVert u_0\rVert_{L^2}^2+W^+([0,T))=:B_{\rm NS}.
\]
No critical norm, scale estimate, compactness statement, or regularity
certificate occurs in this admission transcript.
\end{proposition}

\begin{proof}
Pair the stopped equation with \(u\).  Incompressibility makes the convective
term a signed internal flux, the pressure pairing vanishes after constraint
descent and boundary gluing, and integration by parts turns viscosity into
\(\nu\lVert\nabla u\rVert_2^2\).  The weak total graph records failure of
energy equality as the nonnegative defect measure.  Integration of the
checked balance gives the displayed ceiling.  Every operation is one of the
finite carrier-admission rows of \cref{thm:public-carrier-admission}.
\end{proof}

\begin{theorem}[Exact Navier--Stokes scale price and mandatory lift route]
\label{thm:ns-exact-scale-price}
Let
\[
 u^{(\lambda)}(x,t)
 :=\lambda u(\lambda x,\lambda^2t),
 \qquad \lambda>0,
\]
be the public three-dimensional Navier--Stokes parabolic rescaling, and let
\(I^{(\lambda)}=\lambda^{-2}I\) be the corresponding stopped interval.
For the kinetic-energy carrier of
\cref{prop:ns-generated-carrier-admission},
\begin{align}
 \frac12\lVert u^{(\lambda)}(t)\rVert_{L^2_x}^2
 &=\lambda^{-1}\frac12\lVert u(\lambda^2t)\rVert_{L^2_x}^2,
 \\
 \nu\int_{I^{(\lambda)}}
   \lVert\nabla u^{(\lambda)}(t)\rVert_{L^2_x}^2\,\dd t
 &=\lambda^{-1}\nu\int_I
   \lVert\nabla u(t)\rVert_{L^2_x}^2\,\dd t.
\end{align}
For a scale-covariant normalized target shell, its unnormalised public shell
cocycle remains one.  Consequently the carrier price of a dyadic rescaled
copy satisfies
\[
 p_j=2^{-j}p_0,
 \qquad
 \sum_{j=0}^{\infty}p_j=2p_0<\infty,
\]
while its cumulative target amplitude equals \(N\) after \(N\) crossings.
Thus the native energy ledger does not produce a divergent-price
contradiction for this scale pattern.  The scale change and the mismatch
between unit target progress and price \(2^{-j}p_0\) are retained as a
nonzero \(\Lift\)-owned zero-price successor.  They may not be deleted,
replaced by a uniform lower price, or declared regular.
\end{theorem}

\begin{proof}
Changing variables \(y=\lambda x\) gives
\[
 \int_{\mathbb R^3}|u^{(\lambda)}(x,t)|^2\,\dd x
 =\lambda^2\lambda^{-3}
   \int_{\mathbb R^3}|u(y,\lambda^2t)|^2\,\dd y.
\]
Since \(\nabla u^{(\lambda)}=\lambda^2(\nabla u)(\lambda x,\lambda^2t)\),
the spatial dissipation density scales by \(\lambda^4\lambda^{-3}=\lambda\);
the time change \(s=\lambda^2t\) contributes \(\lambda^{-2}\), proving the
second identity.  Scale covariance of the normalized shell fixes its endpoint
increment at one.  Setting \(\lambda=2^j\) gives the displayed geometric
series.  The public provenance morphism assigns the rescaling constructor and
its price mismatch to \(\Lift\); target alignment is nonzero because the
shell cocycle remains one.  Generated carrier-price routing therefore selects
the zero-price lift successor rather than the divergent-price terminal.
\end{proof}

\begin{theorem}[Exact moving-chart Navier--Stokes lift ledger]
\label{thm:ns-moving-chart-lift-ledger}
Let \(c(t)\in\mathbb R^3\), \(r(t)>0\), and define
\[
 y=\frac{x-c(t)}{r(t)},
 \qquad \frac{\dd s}{\dd t}=r(t)^{-2},
 \qquad v(y,s)=r(t)u(x,t),
 \qquad q(y,s)=r(t)^2p(x,t).
\]
Put
\[
 a(s):=\frac{r_s}{r}=r r_t,
 \qquad b(s):=\frac{c_s}{r}=r c_t.
\]
Then the public equation in the moving chart is exactly
\begin{equation}
 \label{eq:ns-moving-chart-equation}
 \partial_sv-\nu\Delta v+(v\cdot\nabla)v+\nabla q
 -a\bigl(v+y\cdot\nabla v\bigr)-b\cdot\nabla v=0,
 \qquad \nabla\cdot v=0.
\end{equation}
On \(\mathbb R^3\), or after retaining the displayed boundary fluxes, its
kinetic-energy row is
\begin{equation}
 \label{eq:ns-moving-chart-energy}
 \frac{\dd}{\dd s}E_v(s)+\nu D_v(s)+a(s)E_v(s)
 =B_v(s),
 \quad
 E_v:=\frac12\lVert v\rVert_2^2,
 \quad D_v:=\lVert\nabla v\rVert_2^2,
\end{equation}
where \(B_v=0\) on the whole space and otherwise is the complete oriented
public boundary row.  Since \(r_s=ar\),
\begin{equation}
 \label{eq:ns-moving-chart-conjugated-energy}
 \frac{\dd}{\dd s}\bigl(rE_v\bigr)+\nu rD_v=rB_v.
\end{equation}
Thus the scale term \(aE_v\) is an exact \(\Lift\)-owned storage conjugation,
not a new positive expenditure measure.  For a shrinking chart \(a<0\), its
negative sign feeds rescaled kinetic energy, but multiplication by \(r\)
returns the original physical kinetic-energy ledger.  The framework may
charge only \(\nu rD_v\,\dd s\), together with admitted boundary work; it may
not charge \(|a|E_v\,\dd s\) as an additional budget.
\end{theorem}

\begin{proof}
At fixed \(x\), direct differentiation gives
\[
 \partial_tu
 =r^{-3}\partial_sv-r^{-2}c_t\cdot\nabla v
  -r_t r^{-2}\bigl(v+y\cdot\nabla v\bigr).
\]
Moreover
\[
 (u\cdot\nabla_x)u=r^{-3}(v\cdot\nabla_y)v,
 \qquad \Delta_xu=r^{-3}\Delta_yv,
 \qquad \nabla_xp=r^{-3}\nabla_yq.
\]
Multiplication by \(r^3\) yields
\eqref{eq:ns-moving-chart-equation}.  Pair that equation with \(v\).  The
convective, pressure, and centre-translation terms are signed internal
currents.  In three dimensions,
\[
 \int v\cdot\bigl(v+y\cdot\nabla v\bigr)\,\dd y
 =\lVert v\rVert_2^2
  +\frac12\int y\cdot\nabla|v|^2\,\dd y
 =-\frac12\lVert v\rVert_2^2=-E_v.
\]
This gives \eqref{eq:ns-moving-chart-energy}, including the oriented boundary
row when integration by parts has exterior faces.  Multiplying by \(r\) and
using \(r_s=ar\) gives
\eqref{eq:ns-moving-chart-conjugated-energy}.  Every equality is a finite
public weak-equation and signed-gluing calculation.
\end{proof}

\begin{theorem}[Exact rescaled vorticity successor]
\label{thm:ns-rescaled-vorticity-successor}
In the setting of \cref{thm:ns-moving-chart-lift-ledger}, let
\(\omega=\nabla\times v\).  Then
\begin{equation}
 \label{eq:ns-rescaled-vorticity-equation}
 \partial_s\omega-\nu\Delta\omega+(v\cdot\nabla)\omega
 -(\omega\cdot\nabla)v
 -a\bigl(2\omega+y\cdot\nabla\omega\bigr)
 -b\cdot\nabla\omega=0.
\end{equation}
Writing
\[
 E_\omega:=\frac12\lVert\omega\rVert_2^2,
 \qquad D_\omega:=\lVert\nabla\omega\rVert_2^2,
 \qquad
 \mathcal S_\omega
 :=\int (\omega\cdot\nabla)v\cdot\omega\,\dd y,
\]
the exact enstrophy row is
\begin{equation}
 \label{eq:ns-rescaled-enstrophy-row}
 \frac{\dd}{\dd s}E_\omega+\nu D_\omega-aE_\omega
 =\mathcal S_\omega+B_\omega,
\end{equation}
with the complete oriented boundary row \(B_\omega\).  Equivalently,
\begin{equation}
 \label{eq:ns-conjugated-enstrophy-row}
 \frac{\dd}{\dd s}\left(\frac{E_\omega}{r}\right)
 +\nu\frac{D_\omega}{r}
 =\frac{\mathcal S_\omega+B_\omega}{r}.
\end{equation}
The scale term is again an exact \(\Lift\)-owned storage conjugation.  After
it is removed, the only interior signed source capable of increasing
enstrophy is the explicitly displayed \(\Caus\)-owned vortex-stretching
current \(\mathcal S_\omega\).  Its positive part is not charged to the
kinetic-energy carrier unless a finite public domination identity proves that
charge on the same stopped event.
\end{theorem}

\begin{proof}
Take curl in \eqref{eq:ns-moving-chart-equation}.  Incompressibility gives
\[
 \nabla\times((v\cdot\nabla)v)
 =(v\cdot\nabla)\omega-(\omega\cdot\nabla)v,
 \qquad
 \nabla\times(v+y\cdot\nabla v)
 =2\omega+y\cdot\nabla\omega,
\]
which proves \eqref{eq:ns-rescaled-vorticity-equation}.  Pair with \(\omega\).
Transport by \(v\) and by the centre velocity \(b\) is a signed internal
current.  In three dimensions,
\[
 \int\omega\cdot(2\omega+y\cdot\nabla\omega)\,\dd y
 =2\lVert\omega\rVert_2^2-\frac{3}{2}\lVert\omega\rVert_2^2
 =E_\omega.
\]
This gives \eqref{eq:ns-rescaled-enstrophy-row}.  Since
\(r_s/r=a\), division by \(r\) gives
\eqref{eq:ns-conjugated-enstrophy-row}.  Provenance follows directly from
the public constructors: chart dilation is \(\Lift\), whereas the stretching
term descends from nonlinear convection and is \(\Caus\).
\end{proof}

\begin{theorem}[Quantitative target-aligned strain cell]
\label{thm:ns-quantitative-strain-cell}
Assume the whole-space, periodic, or no-slip zero-boundary-flux cell, and use
the public enstrophy continuation coordinate
\[
 Z(s):=\frac{E_\omega(s)}{r(s)}.
\]
Let \(I=[s_-,s_+]\) be a stopped target shell with
\(Z(s_+)-Z(s_-)=1\).  Then
\begin{equation}
 \label{eq:ns-unit-shell-stretching-floor}
 \int_I\frac{(\mathcal S_\omega(s))_+}{r(s)}\,\dd s
 \ge
 \int_I\frac{\mathcal S_\omega(s)}{r(s)}\,\dd s
 =1+\nu\int_I\frac{D_\omega(s)}{r(s)}\,\dd s
 \ge1.
\end{equation}
Let \(S=(\nabla v+\nabla v^{\mathsf T})/2\), and choose its measurable
ordered eigenvalue--projection fields
\[
 S=\sum_{k=1}^3\lambda_kP_k,
 \qquad \lambda_1+\lambda_2+\lambda_3=0.
\]
Define the three public strain-direction currents
\[
 Q_k(I):=\int_I\int
 \frac{\bigl(\lambda_k|P_k\omega|^2\bigr)_+}{r(s)}
 \,\dd y\,\dd s.
\]
Then
\begin{equation}
 \label{eq:ns-strain-cell-one-third-floor}
 \sum_{k=1}^3Q_k(I)\ge1,
 \qquad
 \max_{1\le k\le3}Q_k(I)\ge\frac13.
\end{equation}
Hence every unit enstrophy-shell crossing has an actual finite
\(\Caus\)-owned strain-direction occurrence of amplitude at least \(1/3\).
This occurrence, its eigenprojection graph, and every multiplicity or
eigenvalue-crossing defect are joined before target projection.  No
whole-solution estimate is used.
\end{theorem}

\begin{proof}
Integrate \eqref{eq:ns-conjugated-enstrophy-row} on \(I\) with
\(B_\omega=0\).  The stopped-shell increment is one and
\(D_\omega\ge0\), giving \eqref{eq:ns-unit-shell-stretching-floor}.  The
antisymmetric part of \(\nabla v\) has zero quadratic pairing with \(\omega\),
so
\[
 \mathcal S_\omega
 =\int\omega^{\mathsf T}S\omega\,\dd y
 =\sum_{k=1}^3\int\lambda_k|P_k\omega|^2\,\dd y.
\]
The positive part of a sum is bounded above by the sum of the positive parts.
After integration this gives \(\sum_kQ_k(I)\ge1\); the finite pigeonhole
bound gives \(\max_kQ_k(I)\ge1/3\).  Spectral projections are totalized public
graphs: repeated eigenvalues retain the joint eigenspace projection, while a
nonordinary spectral row is an \(\Abs/\Lift\)-owned graph successor.  Thus the
calculation retains every target-visible value.
\end{proof}

\begin{theorem}[Finite energy forces quantitative target-aligned strain amplification]
\label{thm:ns-quantitative-strain-amplification}
Let \(I_1,\ldots,I_N\) be pairwise disjoint unit enstrophy-shell intervals
from one ordinary same-origin Navier--Stokes branch, and on each interval
choose a strain-direction index \(k_j\) satisfying
\eqref{eq:ns-strain-cell-one-third-floor}.  Put
\[
 q_j:=Q_{k_j}(I_j)\ge\frac13,
 \qquad
 p_j:=\nu\int_{I_j}r(s)\lVert\omega(s)\rVert_2^2\,\dd s.
\]
The second quantity is exactly the restriction of the kinetic-energy
expenditure to the charted event.  On
\(\{\omega\ne0\}\), define
\begin{equation}
 \label{eq:ns-pointwise-strain-amplification}
 \mathcal A_{k}(y,s)
 :=\frac{(\lambda_k(y,s))_+}{\nu r(s)^2}
   \frac{|P_k(y,s)\omega(y,s)|^2}{|\omega(y,s)|^2},
\end{equation}
and put \(\mathcal A_k=0\) on \(\{\omega=0\}\).  Then
\begin{equation}
 \label{eq:ns-strain-current-carrier-density}
 q_j=\int_{I_j}\int \mathcal A_{k_j}(y,s)
          \,\nu r(s)|\omega(y,s)|^2\,\dd y\,\dd s.
\end{equation}
If the branch has generated kinetic-energy ceiling \(B_{\rm NS}<\infty\),
then
\begin{equation}
 \label{eq:ns-linear-amplification-lower-bound}
 \max_{1\le j\le N}\frac{q_j}{p_j}
 \ge\frac{N}{3B_{\rm NS}},
\end{equation}
with \(q_j/p_j=+\infty\) when \(p_j=0<q_j\).  Consequently
\begin{equation}
 \label{eq:ns-pointwise-amplification-unbounded}
 \sup_{j\ge1}\operatorname*{ess\,sup}_{(y,s)\in I_j}
 \mathcal A_{k_j}(y,s)=+\infty
\end{equation}
on every infinite unit-shell cascade.  This is the explicit
\(\Caus/\Lift\)-owned amplification coordinate passed to zero-price
resaturation.
\end{theorem}

\begin{proof}
For incompressible fields on the stated domains,
\(\lVert\nabla v\rVert_2^2=\lVert\omega\rVert_2^2\), up to the already retained
boundary row.  The change of variables in
\cref{thm:ns-moving-chart-lift-ledger} therefore identifies the physical
viscous expenditure with
\(\nu r|\omega|^2\,\dd y\,\dd s\).  Multiplying this density by
\eqref{eq:ns-pointwise-strain-amplification} gives exactly the integrand in
\(Q_k\), proving \eqref{eq:ns-strain-current-carrier-density}.

Disjointness and carrier additivity give
\(\sum_{j=1}^Np_j\le B_{\rm NS}\), whereas
\(\sum_{j=1}^Nq_j\ge N/3\).  Hence
\[
 \max_{j\le N}\frac{q_j}{p_j}
 \ge\frac{\sum_{j\le N}q_j}{\sum_{j\le N}p_j}
 \ge\frac{N}{3B_{\rm NS}},
\]
with the declared zero-denominator convention.  Each ratio is the average of
\(\mathcal A_{k_j}\) against the positive carrier measure on \(I_j\), so it
does not exceed the corresponding essential supremum.  Letting \(N\to\infty\)
proves \eqref{eq:ns-pointwise-amplification-unbounded}.
\end{proof}

\begin{theorem}[Quantitative NS strain-to-concentration successor]
\label{thm:ns-strain-concentration-successor}
In the fixed-chart whole-space or periodic setting of
\cref{thm:ns-quantitative-strain-amplification}, define
\[
 \mathcal C_j:=\int_{I_j}\int_{\mathbb R^3}|\omega|^4\,\dd x\,\dd t
\]
with the periodic integral on \(\mathbb T^3\) in that case.  Then every
selected strain-direction occurrence satisfies the public quantitative row
\begin{equation}
 \label{eq:ns-strain-concentration-lower-bound}
 \mathcal C_j
 \ge \frac{3\nu q_j^2}{p_j},
\end{equation}
with \(\mathcal C_j=+\infty\) when \(p_j=0<q_j\).  In particular, along the
zero-carrier subsequence,
\begin{equation}
 \label{eq:ns-concentration-divergence}
 \mathcal C_j\ge\nu\frac{q_j}{p_j}\longrightarrow+\infty,
\end{equation}
because \(3q_j\ge1\).  Hence the compactified concentration coordinate
\begin{equation}
 \label{eq:ns-compactified-concentration}
 c_j:=\frac1{1+\mathcal C_j}
\end{equation}
converges to zero.  The next public framework state is therefore a
\(\{\Caus,\Lift\}\)-owned vorticity-concentration successor carrying
\((I_j,k_j,q_j,p_j,\mathcal C_j,c_j)\); the \(\Lift\) head is the canonical
evaluation-compactification constructor, not an added analytic premise.
\end{theorem}

\begin{proof}
At every point, trace-freeness of \(S\) gives
\(\lambda_k^2\le\frac23|S|^2\).  On \(\mathbb R^3\) and \(\mathbb T^3\),
\(\|S\|_2^2=\frac12\|\omega\|_2^2\).  Thus
\[
 \int_{I_j}\!\!\int\lambda_{k_j}^2
 \le\frac13\int_{I_j}\!\!\int|\omega|^2
 =\frac{p_j}{3\nu}.
\]
Cauchy--Schwarz on the actual selected event gives
\[
 q_j^2
 =\left(\int_{I_j}\!\!\int
       (\lambda_{k_j}|P_{k_j}\omega|^2)_+\right)^2
 \le
 \left(\int_{I_j}\!\!\int\lambda_{k_j}^2\right)
 \left(\int_{I_j}\!\!\int|P_{k_j}\omega|^4\right)
 \le\frac{p_j}{3\nu}\mathcal C_j.
\]
This proves \eqref{eq:ns-strain-concentration-lower-bound}.  The remaining
claims follow from \(q_j\ge1/3\) and the definition of the fixed
compactification graph.  Provenance is inherited from the selected
\(\Caus\)-owned strain occurrence and the \(\Lift\)-owned concentration
constructor.
\end{proof}

\begin{theorem}[Quantitative NS concentration-to-scale successor]
\label{thm:ns-concentration-scale-successor}
For the concentration cells of
\cref{thm:ns-strain-concentration-successor}, let
\[
 M_j:=\operatorname*{ess\,sup}_{(x,t)\in\mathbb R^3\times I_j}
       |\omega(x,t)|
\]
with the analogous periodic definition.  Then
\begin{equation}
 \label{eq:ns-vorticity-maximum-lower-bound}
 M_j\ge\sqrt3\,\nu\frac{q_j}{p_j}\longrightarrow+\infty.
\end{equation}
The canonical vorticity record scale
\begin{equation}
 \label{eq:ns-canonical-vorticity-scale}
 r_j:=M_j^{-1/2}
\end{equation}
therefore satisfies \(r_j\to0\).  For every rational
\(\varepsilon\in(0,1)\), the public finite-cylinder level-set graph returns
an actual point \((x_{j,\varepsilon},t_{j,\varepsilon})\) in the selected
event with
\[
 r_j^2|\omega(x_{j,\varepsilon},t_{j,\varepsilon})|
 \ge1-\varepsilon.
\]
After centering at \(x_{j,\varepsilon}\) and applying the chart of
\cref{thm:ns-moving-chart-lift-ledger}, the rescaled vorticity has a
nonvanishing unit-scale finite occurrence.  Its charted kinetic-energy price
is \(p_j/r_j\), while conjugation by the public scale factor returns the
original physical charge \(r_j(p_j/r_j)=p_j\).  This is the explicit \(\Lift\)-owned scale
successor required by steps \textup{(I6)}--\textup{(I9)}.
\end{theorem}

\begin{proof}
Since \(|\omega|^4\le M_j^2|\omega|^2\) almost everywhere,
\[
 \mathcal C_j
 \le M_j^2\int_{I_j}\!\!\int|\omega|^2
 =M_j^2\frac{p_j}{\nu}.
\]
Combine this with
\eqref{eq:ns-strain-concentration-lower-bound} to obtain
\eqref{eq:ns-vorticity-maximum-lower-bound}.  The essential-supremum
definition gives a positive-measure \((1-\varepsilon)M_j\) level set, and
the rational finite-cylinder density returns the stated actual occurrence.
Under the NS chart, vorticity transforms as
\(\omega_v(y,s)=r_j^2\omega_u(x,t)\), while
\cref{thm:ns-exact-scale-price} with \(\lambda=r_j\) gives multiplication
of charted kinetic energy and integrated viscous price by \(r_j^{-1}\).
Equation \eqref{eq:ns-moving-chart-conjugated-energy} multiplies the charted
ledger by \(r_j\) and recovers its original physical value.  These are
exactly the scale and finite-occurrence rows claimed.
\end{proof}

\begin{theorem}[Exact NS scale-normalized successor ledger]
\label{thm:ns-scale-normalized-successor-ledger}
Choose the finite occurrences of
\cref{thm:ns-concentration-scale-successor} and set
\begin{align*}
 v_j(y,s)&:=r_j u(x_j+r_jy,t_j+r_j^2s),
 &q_j^{\rm sc}(y,s)&:=r_j^2p(x_j+r_jy,t_j+r_j^2s),\\
 \omega_j(y,s)&:=r_j^2\omega(x_j+r_jy,t_j+r_j^2s).
\end{align*}
On every rational stopped rescaled cylinder these fields satisfy the public
NS equation with the same viscosity, and
\begin{equation}
 \label{eq:ns-scale-normalized-nonvanishing}
 |\omega_j(0,0)|\ge1-\varepsilon_j,
 \qquad \varepsilon_j\downarrow0.
\end{equation}
For the rescaled event corresponding to \(I_j\), the exact carrier values are
\begin{equation}
 \label{eq:ns-scale-normalized-carrier-values}
 E_j^{\rm sc}=\frac{E_j}{r_j},
 \qquad P_j^{\rm sc}=\frac{p_j}{r_j},
 \qquad r_jE_j^{\rm sc}=E_j,
 \qquad r_jP_j^{\rm sc}=p_j.
\end{equation}
The public lift normal form therefore evaluates the scale cell in the
following fixed order:
\begin{enumerate}[label=\textup{(LS\arabic*)}]
\item a first nonzero concentration, oscillation, chart, pressure-tail, or
domain graph coordinate is the corresponding
\(\Lift/\Abs/\Bdry\)-owned finite successor;
\item if all local graph coordinates are ordinary and
\(P_j^{\rm sc}\) is locally bounded, a compatible rescaled path retains the
nonzero vorticity row \eqref{eq:ns-scale-normalized-nonvanishing} and enters
the ordinary realization graph;
\item if \(P_j^{\rm sc}\) is unbounded, its compactified value is the next
\(\Lift\)-owned scale-cohomology successor, while its physical charge remains
the already counted value \(p_j\).
\end{enumerate}
These are total-graph values of the displayed rescaling.  In particular,
\(p_j/r_j\) is never charged as a second physical budget.
\end{theorem}

\begin{proof}
Direct substitution gives the rescaled NS equation and vorticity
transformation; the selected level-set occurrence gives
\eqref{eq:ns-scale-normalized-nonvanishing}.  Spatial and time changes of
variables give the first two identities in
\eqref{eq:ns-scale-normalized-carrier-values}; the last two are their
covariant conjugations.  Apply the ordered lift record
\eqref{eq:target-lift-normal-form} and its exact reconstruction
\eqref{eq:target-lift-exact-current-reconstruction}.  A first failed local
graph is \textup{(LS1)}.  On its complement, the bounded ordinary graph gives
\textup{(LS2)}.  The remaining extended-real scale-price value is
\textup{(LS3)}.  The three cells are disjoint by the first-active-coordinate
order and exhaustive by totalization.
\end{proof}

\begin{theorem}[Exact target-local rescaled NS cylinder ledger]
\label{thm:ns-rescaled-cylinder-ledger}
Let \((v_j,\pi_j)\) denote the scale-normalized fields of
\cref{thm:ns-scale-normalized-successor-ledger}.  For every nonnegative
rational \(\chi\in C_c^\infty(B_{2R}\times(s_1,s_2])\), the public local
energy row is
\begin{align}
 \label{eq:ns-rescaled-local-energy-ledger}
 &\frac12\int |v_j(s_2)|^2\chi(s_2)\,\dd y
 -\frac12\int |v_j(s_1)|^2\chi(s_1)\,\dd y
 +\nu\int_{s_1}^{s_2}\!\!\int\chi|\nabla v_j|^2\,\dd y\,\dd s
 \\
 &\quad=
 \int_{s_1}^{s_2}\!\!\int
 \frac{|v_j|^2}{2}(\partial_s\chi+\nu\Delta\chi)
 +\left(\frac{|v_j|^2}{2}+\pi_j\right)
    v_j\cdot\nabla\chi\,\dd y\,\dd s .
\end{align}
Choose a rational spatial cutoff \(\zeta_R=1\) on \(B_R\), supported in
\(B_{2R}\).  Pressure has the exact target-local split
\begin{equation}
 \label{eq:ns-rescaled-pressure-split}
 \pi_j=\pi_{j,R}^{\rm near}+h_{j,R},
 \qquad
 \pi_{j,R}^{\rm near}:=\mathcal R_a\mathcal R_b
       (\zeta_R v_{j,a}v_{j,b}),
 \qquad \Delta h_{j,R}=0\quad\hbox{on }B_R,
\end{equation}
and the finite local operator row gives
\begin{equation}
 \label{eq:ns-rescaled-near-pressure-bound}
 \|\pi_{j,R}^{\rm near}\|_{L^{3/2}(B_R)}
 \le C_R\|v_j\|_{L^3(B_{2R})}^2.
\end{equation}
Internal cutoff faces cancel when nested cylinders are joined.  The terms in
\eqref{eq:ns-rescaled-local-energy-ledger} have exact heads
\(\Geom\) for viscous production, \(\Caus\) for advective cutoff transport,
\(\Abs\) for near pressure, and \(\{\Abs,\Bdry\}\) for the harmonic pressure
tail; time, Laplacian, and exterior cutoff currents have \(\Bdry\) ancestry.

On each rational cylinder, the ordered public graph therefore returns
exactly one of:
\begin{enumerate}[label=\textup{(LC\arabic*)}]
\item an unbounded \(L^3\) velocity coordinate, giving a \(\Lift/\Caus\)
concentration successor;
\item a least unbounded harmonic coefficient of \(h_{j,R}\), giving the
\(\Abs/\Bdry\) pressure-tail successor;
\item a nonzero local-energy, nonlinear-product, domain, or convergence graph
defect, giving its displayed five-source successor;
\item bounded local \(L^3\), near-pressure, harmonic-tail, energy, and
dissipation coordinates with zero graph defects.  Diagonal cylinder
restriction then produces an ordinary local suitable NS path.  It retains
\eqref{eq:ns-scale-normalized-nonvanishing} precisely when the vorticity
evaluation graph is ordinary; otherwise the first failed vorticity graph is
the next \(\Lift/\Geom\) concentration successor.
\end{enumerate}
The fourth value enters the continuation and realization graphs and is not a
terminal solely because a local profile exists.
\end{theorem}

\begin{proof}
Pair the rescaled NS equation with \(v_j\chi\), integrate diffusion,
transport, and pressure with their signs, and retain every derivative of
\(\chi\).  This gives \eqref{eq:ns-rescaled-local-energy-ledger}.  Applying
the pressure Poisson row to \(\zeta_Rv_{j,a}v_{j,b}\) gives the near term in
\eqref{eq:ns-rescaled-pressure-split}; the difference has zero Laplacian in
\(B_R\).  The rational-cylinder Calder\'on--Zygmund row gives
\eqref{eq:ns-rescaled-near-pressure-bound}.  These operations assign the
stated source heads before target projection.

Totalize the local \(L^3\), harmonic-coefficient, energy, product, domain,
and convergence graphs in that order.  Their first nonordinary values are
\textup{(LC1)}--\textup{(LC3)}.  On the complementary cell, the displayed
bounds give weak energy compactness, the local nonlinear and pressure rows
give distributional graph convergence, and the zero defect coordinates give
the local energy inequality.  Diagonal restriction over rational cylinders
produces \textup{(LC4)}.  Totalization of vorticity evaluation determines
whether its normalized value passes to the path or creates the stated
successor.
\end{proof}

\begin{theorem}[Target-aligned routing of unbounded rescaled dissipation]
\label{thm:ns-unbounded-rescaled-dissipation-routing}
On branch \textup{(LS3)} of
\cref{thm:ns-scale-normalized-successor-ledger}, let
\[
 \dd\mu_j^{\rm sc}:=\nu|\nabla v_j|^2\,\dd y\,\dd s,
 \qquad
 \mu_j^{\rm sc}(\mathcal I_j\times\mathbb R^3)=p_j/r_j\longrightarrow+\infty
\]
after passage to the defining subsequence.  With
\(Q_R=(-R^2,0]\times B_R\) restricted to the actual rescaled time domain,
exactly one target-aligned value occurs:
\begin{enumerate}[label=\textup{(UD\arabic*)}]
\item for some rational \(R\),
\(\sup_j\mu_j^{\rm sc}(Q_R)=+\infty\); the first rational level crossing is
an actual finite \(\Geom/\Lift\)-owned local dissipation-concentration
successor;
\item every fixed \(Q_R\) has bounded mass.  A diagonal sequence
\(R_m\uparrow\infty\) then retains the divergent complementary mass
\[
 \mu_{j_m}^{\rm sc}
   ((\mathcal I_{j_m}\times\mathbb R^3)\setminus Q_{R_m})\longrightarrow+\infty,
\]
which is the explicit \(\{\Geom,\Lift,\Bdry\}\)-owned parabolic tail
successor.
\end{enumerate}
Both values retain the conjugated physical charge
\(r_j\mu_j^{\rm sc}=p_j\); neither creates a second expenditure ledger.
\end{theorem}

\begin{proof}
The increasing rational cylinders exhaust each rescaled spacetime domain.
If their masses are not uniformly bounded, choose the least rational radius
and least rational level crossing to obtain \textup{(UD1)}.  Otherwise set
\(R_m=m\) and \(C_m:=\sup_j\mu_j^{\rm sc}(Q_m)<\infty\).  Choose \(j_m\)
with total mass larger than \(C_m+m\).  The complementary mass then exceeds
\(m\), giving \textup{(UD2)}.  Finite-cylinder occurrence recovery supplies the
actual record in the first case; exhaustion and oriented boundary
compactification supply the tail record in the second.  The last assertion
is \eqref{eq:ns-scale-normalized-carrier-values}.
\end{proof}

\begin{theorem}[Exact horizon and realization routing of ordinary NS scale profiles]
\label{thm:ns-ordinary-scale-profile-routing}
Assume branch \textup{(LC4)} of
\cref{thm:ns-rescaled-cylinder-ledger} retains the vorticity evaluation row.
For a prospective finite contact time \(T_*>0\), the rescaled time domain is
\begin{equation}
 \label{eq:ns-rescaled-time-horizons}
 \mathcal I_j=(-\tau_j^-,\tau_j^+),
 \qquad
 \tau_j^-:=\frac{t_j}{r_j^2},
 \qquad
 \tau_j^+:=\frac{T_*-t_j}{r_j^2}.
\end{equation}
Along a cofinal contact sequence, \(t_j\to T_*\) and \(r_j\to0\), hence
\(\tau_j^-\to+\infty\).  Totalization of the forward-horizon graph selects
a subsequence with
\begin{equation}
 \label{eq:ns-forward-horizon-value}
 \tau_j^+\longrightarrow\Theta\in[0,+\infty].
\end{equation}
The diagonal path of \textup{(LC4)} is therefore an ordinary nonzero suitable
NS profile \((v_\infty,\pi_\infty)\) on
\((-\infty,\Theta)\times\mathbb R^3\), with
\begin{equation}
 \label{eq:ns-ordinary-profile-nonvanishing}
 |\omega_\infty(0,0)|=1,
\end{equation}
whenever \(0<\Theta\le+\infty\); for \(\Theta=0\), the same value is retained
as the preterminal vorticity trace.  It satisfies every rational-cylinder
instance of \eqref{eq:ns-rescaled-local-energy-ledger} and
\eqref{eq:ns-rescaled-pressure-split}.

The horizon coordinate gives exactly one next framework state:
\begin{enumerate}[label=\textup{(HP\arabic*)}]
\item \(\Theta=0\): a nonzero preterminal \(\Bdry/\Lift\)-owned profile at
the collapsed terminal face;
\item \(0<\Theta<\infty\): an ordinary profile with a finite future
continuation face at \(s=\Theta\), passed to the total continuation graph;
\item \(\Theta=+\infty\): an ancient-to-eternal recurrent profile, passed to
target-cyclic return and carrier--polar evaluation on every finite cylinder.
\end{enumerate}
The profile retains origin \(\mathsf{Hyp}\).  It enters a singular terminal
only after an independent public path realizes the same equation, horizon,
nonvanishing, and continuation-failure rows.
\end{theorem}

\begin{proof}
Equation \eqref{eq:ns-rescaled-time-horizons} is the affine time change in
the public scaling chart.  Cofinality gives \(t_j\to T_*\); eventually
\(t_j\ge T_*/2\), so \(\tau_j^-\ge T_*/(2r_j^2)\to\infty\).  Compactification
of \([0,+\infty]\) gives \eqref{eq:ns-forward-horizon-value}.

Branch \textup{(LC4)} supplies compatible ordinary local graphs on every
rational cylinder.  Diagonal restriction constructs the profile and passes
the exact local energy and pressure rows.  Ordinary vorticity evaluation and
\eqref{eq:ns-scale-normalized-nonvanishing} give
\eqref{eq:ns-ordinary-profile-nonvanishing}; when the forward domain shrinks
to its terminal face, the total trace graph retains the preterminal value.
The three values of \(\Theta\) give \textup{(HP1)}--\textup{(HP3)}.  Source
and origin coordinates are closed under the same diagonal restriction, and
the independent-origin condition is exactly the singular realization rule.
\end{proof}

\begin{theorem}[Positive local geometric occurrence in every ordinary NS scale profile]
\label{thm:ns-profile-positive-geometric-occurrence}
Let \((v_\infty,\pi_\infty)\) be an ordinary profile from
\cref{thm:ns-ordinary-scale-profile-routing}.  The ordinary vorticity graph
and \(|\omega_\infty(0,0)|=1\) generate a rational \(\delta>0\) such that
\(|\omega_\infty|\ge1/2\) on the retained one-sided cylinder
\[
 Q_\delta^-:=(-\delta^2,0]\times B_\delta.
\]
Consequently its local viscous occurrence has the explicit lower bound
\begin{equation}
 \label{eq:ns-profile-local-dissipation-floor}
 d_\delta^{\Geom}
 :=\nu\int_{Q_\delta^-}|\nabla v_\infty|^2\,\dd y\,\dd s
 \ge\frac{\nu}{8}|Q_\delta^-|>0.
\end{equation}
On the approximating scale records the corresponding original physical
charge is
\begin{equation}
 \label{eq:ns-profile-pullback-price}
 p_{j,\delta}^{\rm pull}=r_jd_{j,\delta}^{\Geom}\longrightarrow0
\end{equation}
whenever \(d_{j,\delta}^{\Geom}\to d_\delta^{\Geom}<\infty\).  If this
convergence is nonordinary or the local dissipation is unbounded, its first
graph value is respectively the \(\Geom/\Lift\) compactness successor or the
local branch \textup{(UD1)} of
\cref{thm:ns-unbounded-rescaled-dissipation-routing}.

Thus every ordinary nonvanishing scale profile produces an actual positive
\(\Geom\)-headed finite occurrence.  Projection to the retained target
spine either deletes its incidence in \(\Null_{\Sigma_{\rm NS}}\) or appends
the \(\{\Geom,\Lift\}\)-owned zero-physical-price successor with the exact
floor \eqref{eq:ns-profile-local-dissipation-floor}; the occurrence is never
discarded because its pullback price tends to zero.
\end{theorem}

\begin{proof}
Ordinary vorticity evaluation gives continuity at the retained point, hence
the rational cylinder with \(|\omega_\infty|\ge1/2\).  Pointwise,
\(|\operatorname{curl}v|^2\le2|\nabla v|^2\), which proves
\eqref{eq:ns-profile-local-dissipation-floor}.  Ordinary local dissipation
convergence gives \(d_{j,\delta}^{\Geom}\to d_\delta^{\Geom}\), while the
scale conjugation in \eqref{eq:ns-scale-normalized-carrier-values} gives
\eqref{eq:ns-profile-pullback-price}.  The remaining graph values are the
totalized convergence and extended-real alternatives.  Finally apply the
target quotient before carrier--polar evaluation, yielding the null or
zero-price successor cells stated.
\end{proof}

\begin{theorem}[NS profile continuity-radius and scale-summability transition]
\label{thm:ns-profile-scale-summability-transition}
For each ordinary nonvanishing rescaled record, let \(\delta_j\) be the
largest dyadic radius \(2^{-m}\le1\) for which its ordinary vorticity graph
verifies that the cylinder lies inside the rescaled shell event and
\begin{equation}
 \label{eq:ns-profile-continuity-radius}
 |\omega_j(y,s)|\ge\frac12
 \quad\text{on }Q_{\delta_j}^-=(-\delta_j^2,0]\times B_{\delta_j}.
\end{equation}
Exactly one target-aligned transition occurs after passage to a subsequence:
\begin{enumerate}[label=\textup{(CR\arabic*)}]
\item \(\delta_j\to0\).  The compactified radius
\(\delta_j/(1+\delta_j)\to0\) is the next
\(\{\Geom,\Lift,\Bdry\}\)-owned vorticity-modulus/domain successor, carrying
the first failed dyadic continuity or event-domain row.
\item There is a dyadic \(\delta>0\) with \(\delta_j\ge\delta\) for all
records in the subsequence.  Their disjoint original pullbacks satisfy
\begin{equation}
 \label{eq:ns-profile-scale-price-floor}
 p_j\ge r_j\frac{\nu}{8}|Q_\delta^-|,
\end{equation}
and hence
\begin{equation}
 \label{eq:ns-profile-scale-summability}
 \sum_j r_j
 \le\frac{8B_{\rm NS}}{\nu|Q_\delta^-|}<\infty.
\end{equation}
The scale current therefore enters the \(\Lift\) recurrent closure with its
exact summable physical prices and unbounded logarithmic depth
\(-\log r_j\to+\infty\).
\end{enumerate}
Neither branch is a regularity premise: \textup{(CR1)} is a proper graph
successor, while \textup{(CR2)} is the quantitative scale-price record used
by the next recurrent carrier--polar evaluation.
\end{theorem}

\begin{proof}
Every ordinary record has a positive dyadic radius by continuity, its
normalized vorticity value, and the interior finite occurrence selected from
the positive-measure level set.  Either these radii have a positive dyadic lower
bound on a subsequence or they converge to zero.  The latter is
\textup{(CR1)} with its least failed continuity cylinder.

In the former case, pointwise
\(|\operatorname{curl}v_j|^2\le2|\nabla v_j|^2\) and
\eqref{eq:ns-profile-continuity-radius} give rescaled dissipation at least
\(\nu|Q_\delta^-|/8\).  Covariant pullback multiplies this value by \(r_j\),
proving \eqref{eq:ns-profile-scale-price-floor}.  The pulled-back cylinders
lie in the disjoint shell events \(I_j\), so carrier additivity and
\(\sum_jp_j\le B_{\rm NS}\) give
\eqref{eq:ns-profile-scale-summability}.  Finally \(r_j\to0\) gives the
unbounded logarithmic scale coordinate.
\end{proof}

\begin{theorem}[Quantitative recurrent NS scale-cohomology successor]
\label{thm:ns-recurrent-scale-cohomology-successor}
On branch \textup{(CR2)} of
\cref{thm:ns-profile-scale-summability-transition}, pass to the canonical
strictly decreasing record-scale subsequence and put
\begin{equation}
 \label{eq:ns-logarithmic-scale-increments}
 g_j:=\log\frac{r_j}{r_{j+1}}>0.
\end{equation}
Then
\begin{equation}
 \label{eq:ns-logarithmic-scale-divergence}
 \sum_jg_j=+\infty,
 \qquad
 \sum_jp_j\le B_{\rm NS},
 \qquad
 \sup_j\frac{g_j}{p_j}=+\infty,
\end{equation}
with the usual infinite value when \(p_j=0<g_j\).  If the public moving chart
interpolates monotonically between \(r_j\) and \(r_{j+1}\), its scale rate
\(a=r_s/r\le0\) satisfies the exact cohomology row
\begin{equation}
 \label{eq:ns-scale-rate-cocycle}
 -\int_{s_j}^{s_{j+1}}a(s)\,\dd s=g_j.
\end{equation}
Along a subsequence with \(g_j/p_j\to\infty\), the probability laws
\begin{equation}
 \label{eq:ns-normalized-scale-rate-law}
 \dd\zeta_j:=\frac{-a(s)}{g_j}\,\dd s
\end{equation}
carry unit \(\Lift\) scale current and normalized physical price
\(p_j/g_j\to0\).  Every graph-complete limit is therefore the explicit
zero-physical-price \(\Lift\)-owned scale-cohomology tangent.  It is rejoined
with the NS chart equation \eqref{eq:ns-moving-chart-equation}, conjugated
energy row \eqref{eq:ns-moving-chart-conjugated-energy}, vorticity row
\eqref{eq:ns-conjugated-enstrophy-row}, and all five-normal-form descendants
before carrier--polar evaluation.
\end{theorem}

\begin{proof}
Since \(r_j\downarrow0\), the logarithmic increments telescope to
\(+\infty\).  If \(g_j/p_j\) were bounded by \(C\), then
\(\sum_jg_j\le C\sum_jp_j\le CB_{\rm NS}\), a contradiction.  This proves
\eqref{eq:ns-logarithmic-scale-divergence}.  The fundamental theorem of
calculus applied to \(a=\partial_s\log r\) gives
\eqref{eq:ns-scale-rate-cocycle}; monotonicity makes
\eqref{eq:ns-normalized-scale-rate-law} positive with total mass one.
The normalized carrier price tends to zero on the selected subsequence.
Source-preserving completion retains the chart word, the unnormalized
logarithmic cocycle, and the cited NS identities, giving the asserted
zero-price successor.
\end{proof}

\begin{theorem}[Covariant elimination of the whole-space NS pressure tail]
\label{thm:ns-covariant-pressure-tail-elimination}
On the whole-space cell, normalize pressure by the public rational mean-zero
gauge on \(B_1\), and let \(h_{j,R}\) be the harmonic tail in
\eqref{eq:ns-rescaled-pressure-split}.  For \(R\ge2\), its target-visible
part is its gradient, and the pressure kernel gives
\begin{equation}
 \label{eq:ns-covariant-pressure-tail-bound}
 r_j\|\nabla h_{j,R}(s)\|_{L^\infty(B_{R/2})}
 \le \frac{C}{R^4}\,r_j\|v_j(s)\|_2^2
 =\frac{2C}{R^4}E_u(t_j+r_j^2s)
 \le\frac{2CB_{\rm NS}}{R^4}.
\end{equation}
The same estimate holds for every rational time localization contained in
the rescaled domain.  Consequently the iterated public tail coordinate
obeys
\begin{equation}
 \label{eq:ns-covariant-pressure-tail-vanishing}
 \lim_{R\to\infty}\sup_j
 r_j\|\nabla h_{j,R}\|_{L^\infty(B_{R/2})}=0.
\end{equation}
Thus a persistent value of \textup{(LC2)} has exactly one of the following
dispositions: its constant coefficient is pressure-gauge null; its
covariantly valued whole-space tail is target-null by
\eqref{eq:ns-covariant-pressure-tail-vanishing}; or the pressure
reconstruction, energy, domain, or boundary graph is nonordinary and its
least failed row is the corresponding \(\Abs/\Lift/\Bdry\)-owned successor.
On the torus there is no harmonic tail after the zero-mean pressure gauge.
On a bounded domain the uncancelled term is retained in the public
\(\Bdry\) row rather than being assigned the whole-space estimate.
\end{theorem}

\begin{proof}
Let \(K_{ab}\) be the kernel of \(\mathcal R_a\mathcal R_b\).  After the
near field has been removed, differentiation of the exterior convolution
gives, for \(|y|\le R/2\),
\[
 |\nabla h_{j,R}(y,s)|
 \le C\int_{|z|\ge R}\frac{|v_j(z,s)|^2}{|z|^4}\,\dd z
 \le \frac{C}{R^4}\|v_j(s)\|_2^2.
\]
The cutoff-transition annulus satisfies the same estimate and is part of the
near/tail joint record.  The scale identity
\(r_j\|v_j(s)\|_2^2=2E_u(t_j+r_j^2s)\), followed by the generated kinetic
energy ceiling, proves \eqref{eq:ns-covariant-pressure-tail-bound}.  Taking
the supremum and then \(R\to\infty\) proves
\eqref{eq:ns-covariant-pressure-tail-vanishing}.  Pressure constants vanish
in the constraint quotient.  Every operation used above is totalized, so a
failed kernel, cutoff, energy, domain, or trace row returns the stated typed
successor instead of being covered by the estimate.
\end{proof}

\begin{theorem}[NS dilation-kernel terminal rule for the scale tangent]
\label{thm:ns-dilation-kernel-terminal-rule}
Define the public spatial dilation generators
\begin{equation}
 \label{eq:ns-dilation-generators}
 \mathcal D_1v:=v+y\cdot\nabla v,
 \qquad
 \mathcal D_2\omega:=2\omega+y\cdot\nabla\omega,
 \qquad
 \nabla\times\mathcal D_1v=\mathcal D_2\omega.
\end{equation}
On the zero-price scale-cohomology tangent of
\cref{thm:ns-recurrent-scale-cohomology-successor}, target-cyclic return
evaluation has exactly the following ordered values:
\begin{enumerate}[label=\textup{(DK\arabic*)}]
\item a nonzero ordinary target projection of \(\mathcal D_1v\) or
\(\mathcal D_2\omega\), giving the next \(\Lift\)-owned dilation-defect
successor;
\item a first nonordinary dilation, origin-evaluation, product, domain, or
trace graph, giving its \(\Lift/\Geom/\Bdry\)-owned successor;
\item the ordinary fixed dilation kernel
\begin{equation}
 \label{eq:ns-fixed-dilation-kernel}
 \mathcal D_1v=0,
 \qquad \mathcal D_2\omega=0
\end{equation}
on a neighborhood of the retained origin.  This cell is incompatible with
\(|\omega(0,0)|=1\) and is therefore target-null.
\end{enumerate}
Thus an ordinary nonzero NS scale profile cannot survive as a fixed
zero-price dilation atom.  A surviving scale route has a displayed proper
successor from \textup{(DK1)} or \textup{(DK2)}.
\end{theorem}

\begin{proof}
The curl identity in \eqref{eq:ns-dilation-generators} is a direct product
calculation.  The scale term in
\eqref{eq:ns-moving-chart-equation} is \(-a\mathcal D_1v\), and the scale
term in \eqref{eq:ns-rescaled-vorticity-equation} is
\(-a\mathcal D_2\omega\).  Hence normalized return evaluation first tests
their target projections and total graphs, giving \textup{(DK1)} and
\textup{(DK2)}.

On \textup{(DK3)}, integration of the radial equations gives
\[
 v(\lambda y)=\lambda^{-1}v(y),
 \qquad
 \omega(\lambda y)=\lambda^{-2}\omega(y)
 \quad(\lambda>0)
\]
whenever the ray remains in the ordinary neighborhood.  Ordinary evaluation
at \(y=0\) makes both fields locally bounded.  Letting \(\lambda\downarrow0\)
forces \(v(y)=0\) and \(\omega(y)=0\) there, contradicting
\eqref{eq:ns-scale-normalized-nonvanishing}.  The fixed kernel therefore has
zero incidence with the retained nonvanishing target.
\end{proof}

\begin{theorem}[Forced NS dilation balance and public five-source handoff]
\label{thm:ns-forced-dilation-balance}
For every ordinary scale-normalized record,
\begin{equation}
 \label{eq:ns-forced-origin-dilation}
 \mathcal D_2\omega_j(0,0)=2\omega_j(0,0),
 \qquad
 |\mathcal D_2\omega_j(0,0)|\ge2(1-\varepsilon_j).
\end{equation}
Hence the ordinary value of the scale tangent never enters
\textup{(DK3)}: it enters \textup{(DK1)} with the displayed quantitative
amplitude.  This value is not declared terminal.  Let \(e_j\) be the unit
direction of \(\omega_j(0,0)\), and choose the first rational mollifier
\(\phi_{\rho_j}\) for which
\begin{equation}
 \label{eq:ns-mollified-dilation-floor}
 \left\langle \mathcal D_2\omega_j,
 e_j\phi_{\rho_j}\right\rangle\ge1.
\end{equation}
Such a mollifier is generated by the ordinary origin and derivative graphs.
On every moving-scale interval \(J_j\), weak testing of
\eqref{eq:ns-rescaled-vorticity-equation} gives the exact balance
\begin{align}
 \label{eq:ns-dilation-five-source-balance}
 &\int_{J_j}(-a)
   \left\langle\mathcal D_2\omega,e_j\phi_{\rho_j}\right\rangle\,\dd s
 \notag\\
 &\quad=-\int_{J_j}\bigl\langle
   \partial_s\omega-\nu\Delta\omega
   +(v\cdot\nabla)\omega-(\omega\cdot\nabla)v
   -b\cdot\nabla\omega,
   e_j\phi_{\rho_j}\bigr\rangle\,\dd s.
\end{align}
If \(-a\ge0\), \(-\int_{J_j}a\,\dd s=g_j\), and
\eqref{eq:ns-mollified-dilation-floor} persists on \(J_j\), then at least
one of the five terms on the right has normalized absolute current at least
\(1/5\):
\begin{equation}
 \label{eq:ns-dilation-handoff-floor}
 \max_{1\le\ell\le5}
 \frac{|\mathfrak J_{j,\ell}|}{g_j}\ge\frac15.
\end{equation}
Their public heads, in the displayed order, are storage/trace
\(\{\Lift,\Bdry\}\), diffusion \(\Geom\), advective transport \(\Caus\),
vortex stretching \(\Caus\), and chart translation
\(\{\Lift,\Bdry\}\).  Failure of persistence of
\eqref{eq:ns-mollified-dilation-floor} is itself the least
origin, direction, derivative, trace, or domain graph successor.  Therefore
the nonzero dilation value is converted by an exact public identity into a
quantitative five-source successor; it is neither an assumed estimate nor a
fixed-point conclusion.
\end{theorem}

\begin{proof}
At the retained center, the term \(y\cdot\nabla\omega_j\) vanishes, so
\eqref{eq:ns-forced-origin-dilation} follows from
\eqref{eq:ns-dilation-generators} and
\eqref{eq:ns-scale-normalized-nonvanishing}.  Ordinary local derivative and
origin evaluation make mollification converge to that value, which supplies
the first rational radius satisfying
\eqref{eq:ns-mollified-dilation-floor}.  Move the scale term in
\eqref{eq:ns-rescaled-vorticity-equation} to the opposite side, pair with
the displayed rational test, and integrate in \(s\).  This is
\eqref{eq:ns-dilation-five-source-balance}.  Under the persistence branch,
the left side is at least \(g_j\).  The triangle inequality applied to its
five right-hand summands proves \eqref{eq:ns-dilation-handoff-floor}.  The
source heads are read directly from the five public constructors.  If any
pairing, direction, mollification, interval, or trace operation is
nonordinary, totalization selects its first failed graph row before the
inequality is used.
\end{proof}

\begin{theorem}[Cumulative NS dilation handoff after chart gluing]
\label{thm:ns-cumulative-dilation-handoff}
Concatenate consecutive moving-scale intervals from
\cref{thm:ns-forced-dilation-balance}, and put
\(G_N:=\sum_{j=1}^Ng_j\).  Transport each rational vector test through the
public center chart so that
\begin{equation}
 \label{eq:ns-center-adjoint-test}
 \partial_s\psi_j=b\cdot\nabla\psi_j
\end{equation}
and compatible endpoint tests agree.  Then the time and center-translation
terms glue exactly:
\begin{equation}
 \label{eq:ns-time-translation-gluing}
 \int_{J_j}\left\langle
   \partial_s\omega-b\cdot\nabla\omega,\psi_j
 \right\rangle\dd s
 =\langle\omega(s_j^+),\psi_j(s_j^+)\rangle
  -\langle\omega(s_j^-),\psi_j(s_j^-)\rangle.
\end{equation}
Assume the dilation floor persists and the cumulative endpoint storage is
\(o(G_N)\).  Then
\begin{equation}
 \label{eq:ns-three-interior-handoff}
 \limsup_{N\to\infty}
 \max\left\{
  \frac{|\mathfrak J_{\Geom,N}|}{G_N},
  \frac{|\mathfrak J_{\rm adv,N}|}{G_N},
  \frac{|\mathfrak J_{\rm str,N}|}{G_N}
 \right\}\ge\frac13,
\end{equation}
where the three cumulative currents are given without suppressed terms by
\begin{align}
 \label{eq:ns-three-dilation-currents}
 \mathfrak J_{\Geom,N}
 &:=-\nu\sum_{j\le N}\int_{J_j}
       \langle\Delta\omega,\psi_j\rangle\,\dd s,
 &
 \mathfrak J_{\rm adv,N}
 &:=\sum_{j\le N}\int_{J_j}
       \langle(v\cdot\nabla)\omega,\psi_j\rangle\,\dd s,\notag\\
 \mathfrak J_{\rm str,N}
 &:=-\sum_{j\le N}\int_{J_j}
       \langle(\omega\cdot\nabla)v,\psi_j\rangle\,\dd s.
\end{align}
Integration by parts identifies their exact public localizations:
\begin{equation}
 \label{eq:ns-three-dilation-localizations}
 \langle\Delta\omega,\psi_j\rangle
   =\langle\omega,\Delta\psi_j\rangle,
 \qquad
 \langle(v\cdot\nabla)\omega,\psi_j\rangle
   =-\langle\omega,(v\cdot\nabla)\psi_j\rangle,
\end{equation}
with every exterior face retained.  Thus the three surviving owners are,
respectively, \(\{\Geom,\Bdry\}\),
\(\{\Caus,\Bdry\}\), and \(\Caus\).  If endpoint storage is not
\(o(G_N)\), its least unbounded normalized endpoint value is the explicit
\(\{\Lift,\Bdry\}\)-owned storage/trace successor.  If tests do not glue,
their least transport, domain, or trace defect is the corresponding proper
successor.  Hence chart motion cannot absorb the logarithmic scale current:
after signed gluing it is a quantitative interior source current or a
displayed storage/graph successor.
\end{theorem}

\begin{proof}
Differentiate \(\langle\omega,\psi_j\rangle\), use
\eqref{eq:ns-center-adjoint-test}, and integrate the translation term by
parts.  This proves \eqref{eq:ns-time-translation-gluing}.  Summing over
\(j\le N\) cancels every compatible internal endpoint.  The persistent
dilation floor makes the summed left side of
\eqref{eq:ns-dilation-five-source-balance} at least \(G_N\).  Divide the
summed identity by \(G_N\).  The two glued endpoint values tend to zero by
the storage hypothesis, so the triangle inequality leaves the three currents
in \eqref{eq:ns-three-dilation-currents}; finite pigeonholing proves
\eqref{eq:ns-three-interior-handoff}.  The identities in
\eqref{eq:ns-three-dilation-localizations} follow from incompressibility and
two integrations by parts.  Totalization before those operations gives the
stated storage, transport, domain, and trace successors on every
nonordinary value.
\end{proof}

\begin{theorem}[Exact two-owner reduction of the NS dilation current]
\label{thm:ns-two-owner-dilation-reduction}
On the ordinary glued cell of
\cref{thm:ns-cumulative-dilation-handoff}, define
\begin{equation}
 \label{eq:ns-curl-nonlinear-current}
 \mathfrak J_{{\rm nl},N}
 :=\sum_{j\le N}\int_{J_j}
 \left\langle
  (v\cdot\nabla)\omega-(\omega\cdot\nabla)v,\psi_j
 \right\rangle\dd s.
\end{equation}
The public curl identity gives the exact representations
\begin{align}
 \label{eq:ns-curl-nonlinear-representations}
 \mathfrak J_{{\rm nl},N}
 &=\sum_{j\le N}\int_{J_j}
   \left\langle\nabla\times((v\cdot\nabla)v),\psi_j\right\rangle\dd s
 \notag\\
 &=\sum_{j\le N}\int_{J_j}
   \left\langle(v\cdot\nabla)v,\nabla\times\psi_j\right\rangle\dd s
 \notag\\
 &=-\sum_{j\le N}\int_{J_j}\int
   (v\otimes v):\nabla(\nabla\times\psi_j)\,\dd y\,\dd s,
\end{align}
including the oriented exterior faces on nonclosed supports.  Consequently
\begin{equation}
 \label{eq:ns-two-owner-dilation-floor}
 \limsup_{N\to\infty}
 \max\left\{
  \frac{|\mathfrak J_{\Geom,N}|}{G_N},
  \frac{|\mathfrak J_{{\rm nl},N}|}{G_N}
 \right\}\ge\frac12.
\end{equation}
The first current has exact owner \(\{\Geom,\Bdry\}\); the second has exact
owner \(\{\Caus,\Bdry\}\).  No pressure, scale, or abstract source term is
left in this interior handoff.  If the curl, product, integration-by-parts,
or exterior-face graph is nonordinary, its least failed finite row is the
corresponding typed successor.
\end{theorem}

\begin{proof}
For divergence-free \(v\),
\[
 \nabla\times((v\cdot\nabla)v)
 =(v\cdot\nabla)\omega-(\omega\cdot\nabla)v.
\]
Pairing with \(\psi_j\), integrating curl once and convection once, proves
\eqref{eq:ns-curl-nonlinear-representations} with its boundary faces.  After
the time--translation endpoint row has been removed, the summed dilation
identity contains exactly
\(\mathfrak J_{\Geom,N}+\mathfrak J_{{\rm nl},N}\), up to a term
\(o(G_N)\).  Its absolute value is at least \(G_N-o(G_N)\).  Divide by
\(G_N\), apply the triangle inequality, and use two-term pigeonholing to
obtain \eqref{eq:ns-two-owner-dilation-floor}.  The source labels follow
from diffusion and nonlinear convection before target projection.
\end{proof}

\begin{theorem}[Maximum-centered NS scale current is forced strain]
\label{thm:ns-maximum-centered-strain-handoff}
On an ordinary maximum-selector cell, let
\[
 M(t):=\max_x|\omega_u(x,t)|,
 \qquad r(t):=M(t)^{-1/2},
 \qquad |\omega_u(c(t),t)|=M(t),
\]
and use the public moving chart of
\cref{thm:ns-moving-chart-lift-ledger}.  With
\(\rho=|\omega_v|\) and \(\xi=\omega_v/\rho\), its normalization gives
\begin{equation}
 \label{eq:ns-maximum-chart-normalization}
 \rho(0,s)=1,
 \qquad \rho(y,s)\le1,
 \qquad \nabla\rho(0,s)=0,
 \qquad \Delta\rho(0,s)\le0
\end{equation}
at every ordinary differentiability time.  The moving-chart magnitude
equation is
\begin{align}
 \label{eq:ns-moving-vorticity-magnitude}
 (\partial_s+(v-b)\cdot\nabla)\rho
 ={}&\nu(\Delta\rho-\rho|\nabla\xi|^2)
   +\rho\,\xi^{\mathsf T}S\xi
   +a(2\rho+y\cdot\nabla\rho).
\end{align}
Evaluation at the generated maximum therefore yields the pointwise public
inequality
\begin{equation}
 \label{eq:ns-maximum-strain-scale-rate}
 \xi(0,s)^{\mathsf T}S(0,s)\xi(0,s)
 =-\nu\Delta\rho(0,s)+\nu|\nabla\xi(0,s)|^2-2a(s)
 \ge-2a(s).
\end{equation}
For every shrinking maximum-chart interval \(J_j\),
\begin{equation}
 \label{eq:ns-maximum-strain-logarithmic-floor}
 \int_{J_j}
 \bigl(\xi^{\mathsf T}S\xi\bigr)_+(0,s)\,\dd s
 \ge2\int_{J_j}(-a(s))\,\dd s=2g_j.
\end{equation}
Thus the sign-aligned ordinary value of the dilation handoff is the actual
\(\Caus\)-owned strain-direction current with threshold \(2g_j\); diffusion
and chart transport cannot finance it.  A failure of maximum attainment,
measurable center selection, differentiability, direction evaluation,
Laplacian evaluation, or persistence of the maximum chart is not omitted:
the first failed row is respectively the generated tail, selector, temporal,
direction, \(\Geom\), or \(\Lift/\Bdry\) successor.  Rational continuity
localization turns every strict instance of
\eqref{eq:ns-maximum-strain-scale-rate} into an actual finite-cylinder
occurrence; collapse of those localization radii is the modulus successor
\textup{(CR1)}.
\end{theorem}

\begin{proof}
Vorticity scales as \(\omega_v=r^2\omega_u\), so the definitions of \(r\)
and \(c\) give \eqref{eq:ns-maximum-chart-normalization}.  Take the scalar
product of \eqref{eq:ns-rescaled-vorticity-equation} with \(\xi\), use
\(\xi\cdot\Delta\omega=\Delta\rho-\rho|\nabla\xi|^2\), and move the chart
terms to the right.  This proves
\eqref{eq:ns-moving-vorticity-magnitude}.  At \(y=0\), the normalized maximum
has zero time derivative and zero spatial transport; substitution of
\eqref{eq:ns-maximum-chart-normalization} gives the equality and sign in
\eqref{eq:ns-maximum-strain-scale-rate}.  Integrating its positive part and
using \eqref{eq:ns-scale-rate-cocycle} gives
\eqref{eq:ns-maximum-strain-logarithmic-floor}.  The remaining alternatives
are precisely the total graphs of the operations used in this calculation;
rational finite-occurrence recovery supplies the last assertion.
\end{proof}

\begin{theorem}[NS strain-modulus or inverse-scale-horizon transition]
\label{thm:ns-strain-modulus-horizon-transition}
On a shrinking maximum-centered interval \(J_j\), let \(\sigma_j\) be the
largest dyadic radius such that the ordinary strain graph verifies
\begin{equation}
 \label{eq:ns-strain-radius-floor}
 |S(y,s)|\ge -a(s)
 \quad\text{for }y\in B_{\sigma_j}\text{ and almost every }s\in J_j
 \text{ with }a(s)<0.
\end{equation}
The radius is generated from
\eqref{eq:ns-maximum-strain-scale-rate}; it is not input data.  After
passage to a subsequence exactly one of the following occurs:
\begin{enumerate}[label=\textup{(SR\arabic*)}]
\item \(\sigma_j\to0\).  The first failed dyadic strain-continuity,
center, time, or domain row is the
\(\{\Caus,\Lift,\Bdry\}\)-owned strain-modulus successor.
\item \(\sigma_j\ge\sigma>0\).  With the actual physical carrier price and
inverse-scale chart horizon
\begin{equation}
 \label{eq:ns-strain-horizon-definitions}
 \mathfrak p_j:=\nu\int_{J_j}r(s)\|\omega(s)\|_2^2\,\dd s,
 \qquad
 H_j:=\int_{J_j}\frac{\dd s}{r(s)},
\end{equation}
the public strain--vorticity identity and weighted Cauchy--Schwarz give
\begin{equation}
 \label{eq:ns-strain-horizon-price}
 \mathfrak p_j
 \ge2\nu|B_\sigma|\int_{J_j}r(s)a(s)^2\,\dd s
 \ge2\nu|B_\sigma|\frac{g_j^2}{H_j}.
\end{equation}
Consequently every zero-price scale subsequence
\(\mathfrak p_j/g_j\to0\) satisfies the forced quantitative transition
\begin{equation}
 \label{eq:ns-inverse-scale-horizon-divergence}
 \frac{H_j}{g_j}
 \ge2\nu|B_\sigma|\frac{g_j}{\mathfrak p_j}
 \longrightarrow+\infty.
\end{equation}
Its next coordinate is therefore the compactified
\(\{\Lift,\Bdry\}\)-owned inverse-scale-horizon successor, joined with the
same strain occurrence and physical event.  It is not an assumed long-time
estimate and it is not counted as a second budget.
\end{enumerate}
\end{theorem}

\begin{proof}
The strict center inequality in
\eqref{eq:ns-maximum-strain-scale-rate}, ordinary strain evaluation, and
dyadic finite-cylinder recovery generate a positive radius at each finite
record.  Its radii either collapse or have a positive dyadic lower bound,
giving \textup{(SR1)}--\textup{(SR2)}.  On \textup{(SR2)},
\(\|S\|_2^2=\frac12\|\omega\|_2^2\) and
\eqref{eq:ns-strain-radius-floor} imply the first inequality in
\eqref{eq:ns-strain-horizon-price}.  Since
\(g_j=\int_{J_j}(-a)\,\dd s\), weighted Cauchy--Schwarz gives
\[
 g_j^2
 \le\left(\int_{J_j}r a^2\,\dd s\right)
      \left(\int_{J_j}r^{-1}\,\dd s\right),
\]
which proves the second inequality.  Rearrangement gives
\eqref{eq:ns-inverse-scale-horizon-divergence}.  Totalization of the radius,
integral, and quotient graphs supplies the stated successor on every
nonordinary value.
\end{proof}

\begin{theorem}[Inverse-scale horizon law and slow-collapse tangent]
\label{thm:ns-inverse-horizon-slow-collapse}
On branch \textup{(SR2)} of
\cref{thm:ns-strain-modulus-horizon-transition}, use physical time
\(\dd s=r^{-2}\dd t\) and define
\begin{equation}
 \label{eq:ns-inverse-horizon-physical-form}
 H_j=\int_{J_j}\frac{\dd s}{r(s)}
 =\int_{\widehat J_j}\frac{\dd t}{r(t)^3}
 =\int_{\widehat J_j}M(t)^{3/2}\,\dd t,
 \qquad M=r^{-2}.
\end{equation}
The public probability law and covariant collapse-speed coordinate are
\begin{equation}
 \label{eq:ns-inverse-horizon-law}
 \dd\vartheta_j
 :=\frac{r(s)^{-1}\,\dd s}{H_j}
 =\frac{r(t)^{-3}\,\dd t}{H_j},
 \qquad
 \kappa:=-a r=-r^2r_t\ge0.
\end{equation}
They satisfy the exact rows
\begin{equation}
 \label{eq:ns-slow-collapse-values}
 \int\kappa\,\dd\vartheta_j=\frac{g_j}{H_j},
 \qquad
 \frac{\mathfrak p_j}{H_j}
 =\frac{\mathfrak p_j}{g_j}\frac{g_j}{H_j}.
\end{equation}
Hence the zero-price value forced by
\eqref{eq:ns-inverse-scale-horizon-divergence} obeys
\begin{equation}
 \label{eq:ns-slow-collapse-limit}
 \int\kappa\,\dd\vartheta_j\longrightarrow0,
 \qquad
 \frac{\mathfrak p_j}{H_j}\longrightarrow0.
\end{equation}
Every graph-complete limit is the explicit
\(\{\Lift,\Bdry\}\)-owned slow-collapse horizon tangent, while the
unnormalized logarithmic cocycle \(g_j\) and the \(\Caus\)-owned strain floor
\eqref{eq:ns-maximum-strain-logarithmic-floor} remain in the cumulative
record.  A nonordinary time change, maximum graph, Radon--Nikodym quotient,
or endpoint is its least typed successor; none is replaced by a lifespan or
integrability assumption.
\end{theorem}

\begin{proof}
The first two equalities in
\eqref{eq:ns-inverse-horizon-physical-form} are the public time chart, and
the last uses \(M=r^{-2}\).  The law in
\eqref{eq:ns-inverse-horizon-law} has mass one.  Since
\(a=r r_t\) and \(\dd s=r^{-2}\dd t\), direct multiplication gives
\[
 \int\kappa\,\dd\vartheta_j
 =\frac1{H_j}\int_{J_j}(-ar)r^{-1}\,\dd s
 =\frac1{H_j}\int_{J_j}(-a)\,\dd s
 =\frac{g_j}{H_j}.
\]
The price identity in \eqref{eq:ns-slow-collapse-values} is algebraic.
On the zero-price subsequence of branch \textup{(SR2)}, both factors tend to
zero, proving
\eqref{eq:ns-slow-collapse-limit}.  Source-preserving completion retains the
unnormalized scale and strain coordinates and totalizes every operation just
listed.
\end{proof}

\begin{theorem}[Forced physical horizon divergence and dyadic maximum ledger]
\label{thm:ns-horizon-dyadic-maximum-ledger}
Suppose the maximum-centered intervals \(\widehat J_j\subset[0,T_*)\) are
disjoint and the slow-collapse value of
\cref{thm:ns-inverse-horizon-slow-collapse} persists.  Put
\(\mathcal T:=\bigcup_j\widehat J_j\).  This is the generated target-bearing
physical event union, and
\begin{equation}
 \label{eq:ns-three-halves-maximum-divergence}
 \int_{\mathcal T}M(t)^{3/2}\,\dd t
 =\sum_jH_j=+\infty.
\end{equation}
For its public dyadic maximum cells
\[
 E_k^{\mathcal T}:=\{t\in\mathcal T:2^k\le M(t)<2^{k+1}\},
\]
this is equivalent, up to the fixed factors \(1\) and \(2^{3/2}\), to
\begin{equation}
 \label{eq:ns-dyadic-maximum-duration-divergence}
 \sum_{k\in\mathbb Z}2^{3k/2}|E_k^{\mathcal T}|=+\infty.
\end{equation}
Every finite level crossing and duration in this series is an actual
same-origin record.  Thus the inverse-horizon successor is not a symbolic
``long-time'' branch: it is the explicit
\(\{\Caus,\Lift,\Bdry\}\)-owned divergence of a physical-time dyadic
maximum ledger.  A failed maximum measurability, dyadic partition, duration,
or disjoint-event row is its least typed graph successor.
\end{theorem}

\begin{proof}
By \eqref{eq:ns-inverse-scale-horizon-divergence}, eventually
\(H_j/g_j\ge1\).  Since \(\sum_jg_j=+\infty\), one has
\(\sum_jH_j=+\infty\).  Disjointness and
\eqref{eq:ns-inverse-horizon-physical-form} prove
\eqref{eq:ns-three-halves-maximum-divergence}.  On \(E_k^{\mathcal T}\),
\(2^{3k/2}\le M^{3/2}<2^{3(k+1)/2}\); summing over the disjoint dyadic cells
proves the equivalence with
\eqref{eq:ns-dyadic-maximum-duration-divergence}.  The cells and their
durations belong to the rational stopped-event grammar, so totalization
gives the last assertion.
\end{proof}

\begin{theorem}[Maximum-core energy budget and dyadic duration corridor]
\label{thm:ns-maximum-core-duration-corridor}
On the maximum-centered chart, let \(\delta(t)\) be the largest dyadic radius
for which the public vorticity graph verifies
\begin{equation}
 \label{eq:ns-maximum-vorticity-core}
 |\omega_v(y,s(t))|\ge\frac12
 \qquad(y\in B_{\delta(t)}).
\end{equation}
On every physical tail interval exactly one target-aligned value occurs:
\begin{enumerate}[label=\textup{(MC\arabic*)}]
\item the dyadic radii have no positive lower bound.  Their first collapsing
level is the \(\{\Geom,\Lift,\Bdry\}\)-owned vorticity-core modulus successor;
\item there is a dyadic \(\delta>0\) with \(\delta(t)\ge\delta\).  The native
kinetic-energy expenditure then gives the computed physical-time bound
\begin{equation}
 \label{eq:ns-square-root-maximum-budget}
 \frac{\nu|B_\delta|}{4}
 \int_{\mathcal T}M(t)^{1/2}\,\dd t
 \le
 \nu\int_{\mathcal J}r(s)\|\omega_v(s)\|_2^2\,\dd s
 \le B_{\rm NS},
\end{equation}
where \(\mathcal J\) is the joined rescaled interval family and
\(\mathcal T\) its physical pullback.
\end{enumerate}
If \textup{(MC2)} and the slow-collapse value of
\cref{thm:ns-horizon-dyadic-maximum-ledger} both persist on the same cofinal
tail, their public dyadic durations satisfy the explicit corridor
\begin{equation}
 \label{eq:ns-dyadic-duration-corridor}
 \sum_k2^{k/2}|E_k^{\mathcal T}|<+\infty,
 \qquad
 \sum_k2^{3k/2}|E_k^{\mathcal T}|=+\infty.
\end{equation}
Thus the surviving route is localized to a quantitatively specified
amplitude--duration distribution.  Neither side of
\eqref{eq:ns-dyadic-duration-corridor} is supplied as a regularity
criterion; both are derived from the public maximum chart, the generated
vorticity core, and the single native energy ledger.
\end{theorem}

\begin{proof}
The radii are positive at every ordinary finite record by
\eqref{eq:ns-maximum-chart-normalization} and ordinary vorticity
continuity.  Dyadic totalization gives \textup{(MC1)}--\textup{(MC2)}.
On \textup{(MC2)},
\[
 \nu r(s)\|\omega_v(s)\|_2^2
 \ge\frac{\nu|B_\delta|}{4}r(s).
\]
Since \(\dd s=r^{-2}\dd t\) and \(r=M^{-1/2}\),
\(r\,\dd s=M^{1/2}\dd t\).  Integration and carrier additivity prove
\eqref{eq:ns-square-root-maximum-budget}.  On \(E_k^{\mathcal T}\),
\(M^{1/2}\) is comparable to \(2^{k/2}\), so the first statement in
\eqref{eq:ns-dyadic-duration-corridor} follows.  The second is
\eqref{eq:ns-dyadic-maximum-duration-divergence}.
\end{proof}

\begin{theorem}[Carrier--polar amplification on the NS duration corridor]
\label{thm:ns-duration-corridor-polar-amplification}
On the common \textup{(MC2)} slow-collapse cell, use the actual dyadic
target events \(E_k^{\mathcal T}\) and define
\begin{align}
 \label{eq:ns-duration-cell-current-price}
 q_k^{\rm hor}
 &:=\int_{E_k^{\mathcal T}}M(t)^{3/2}\,\dd t,
 &
 p_k^{\rm core}
 &:=\nu\int_{s(E_k^{\mathcal T})}
      r(s)\int_{B_\delta}|\omega_v(y,s)|^2\,\dd y\,\dd s .
\end{align}
These are respectively the inverse-horizon target current and the restriction
of the single native physical carrier to the same event.  They satisfy
\begin{equation}
 \label{eq:ns-duration-current-price-sums}
 \sum_kq_k^{\rm hor}=+\infty,
 \qquad
 \sum_kp_k^{\rm core}\le B_{\rm NS}.
\end{equation}
Let \(K_N\) be the least finite set of dyadic levels for which
\(\sum_{k\in K_N}q_k^{\rm hor}\ge N\).  Then
\begin{equation}
 \label{eq:ns-duration-polar-amplification}
 \max_{k\in K_N}\frac{q_k^{\rm hor}}{p_k^{\rm core}}
 \ge\frac{N}{B_{\rm NS}},
\end{equation}
with the usual infinite value for a positive current with zero price.
Consequently a subsequence of actual dyadic maximum-duration events has
\(q_k^{\rm hor}/p_k^{\rm core}\to+\infty\).  This is the explicit
\(\{\Caus,\Geom,\Lift,\Bdry\}\)-owned carrier--polar successor of the
slow-collapse route.  A failed same-event restriction, duration, carrier, or
quotient row is its least typed graph successor; no global norm or external
regularity criterion is used.
\end{theorem}

\begin{proof}
The first identity in \eqref{eq:ns-duration-current-price-sums} is
\eqref{eq:ns-three-halves-maximum-divergence}.  The dyadic events are
disjoint subsets of the generated physical event union, so carrier
additivity gives the second.  Divergence of the nonnegative current sum makes
every \(K_N\) finite.  If all ratios on \(K_N\) were below \(N/B_{\rm NS}\),
then
\[
 \sum_{k\in K_N}q_k^{\rm hor}
 <\frac{N}{B_{\rm NS}}\sum_{k\in K_N}p_k^{\rm core}
 \le N,
\]
contradicting the definition of \(K_N\).  This proves
\eqref{eq:ns-duration-polar-amplification}; finite-prefix occurrence recovery
and provenance preservation give the stated successor.
\end{proof}

\begin{theorem}[Explicit zero-carrier horizon-density tangent]
\label{thm:ns-horizon-density-tangent}
On an amplified duration cell from
\cref{thm:ns-duration-corridor-polar-amplification}, set
\begin{equation}
 \label{eq:ns-horizon-price-measure-density}
 \dd m_k^{\rm core}
 :=\nu r(s)|\omega_v(y,s)|^2
    \mathbf 1_{B_\delta\times s(E_k^{\mathcal T})}\,\dd y\,\dd s,
 \qquad
 \mathcal A_k^{\rm hor}
 :=\frac{1}{\nu|B_\delta|r(s)^2|\omega_v(y,s)|^2}.
\end{equation}
The core floor makes the denominator nonzero.  The time and scale charts give
the exact same-event identity
\begin{equation}
 \label{eq:ns-horizon-current-carrier-density}
 q_k^{\rm hor}
 =\int\mathcal A_k^{\rm hor}\,\dd m_k^{\rm core}.
\end{equation}
For \(q_k^{\rm hor}>0\), define the public probability law
\[
 \dd\eta_k^{\rm hor}
 :=\frac{\mathcal A_k^{\rm hor}}{q_k^{\rm hor}}\,
   \dd m_k^{\rm core}.
\]
Along the carrier--polar subsequence,
\begin{equation}
 \label{eq:ns-horizon-reciprocal-tangent}
 \int\frac{1}{\mathcal A_k^{\rm hor}}\,\dd\eta_k^{\rm hor}
 =\frac{p_k^{\rm core}}{q_k^{\rm hor}}\longrightarrow0,
 \qquad
 \eta_k^{\rm hor}\{\mathcal A_k^{\rm hor}\le L\}
 \le L\frac{p_k^{\rm core}}{q_k^{\rm hor}}\longrightarrow0
\end{equation}
for every finite \(L\).  Moreover the maximum-chart bounds
\(1/2\le|\omega_v|\le1\) on the core give
\begin{equation}
 \label{eq:ns-horizon-density-maximum-comparison}
 \frac{M(t)}{\nu|B_\delta|}
 \le\mathcal A_k^{\rm hor}
 \le\frac{4M(t)}{\nu|B_\delta|}.
\end{equation}
Thus the duration corridor produces a computed zero-carrier tangent
supported at infinite horizon density, quantitatively equivalent on its
actual core to the dyadic vorticity maximum.  Every coordinate is generated
from the same event restriction and native carrier; a nonordinary density,
division, or probability graph is its least typed successor.
\end{theorem}

\begin{proof}
Multiplying the two expressions in
\eqref{eq:ns-horizon-price-measure-density} gives
\[
 \mathcal A_k^{\rm hor}\,\dd m_k^{\rm core}
 =\frac{1}{|B_\delta|r(s)}
   \mathbf 1_{B_\delta\times s(E_k^{\mathcal T})}\,\dd y\,\dd s.
\]
Spatial integration gives \(\int r^{-1}\dd s\).  By
\(\dd t=r^2\dd s\) and \(M=r^{-2}\), this equals
\(\int_{E_k^{\mathcal T}}M^{3/2}\dd t=q_k^{\rm hor}\), proving
\eqref{eq:ns-horizon-current-carrier-density}.  The reciprocal identity and
finite-level estimate in \eqref{eq:ns-horizon-reciprocal-tangent} are the
same direct probability calculation as in
\cref{thm:ns-zero-carrier-strain-tangent}.  Finally \(r^{-2}=M\) and the
core bounds prove \eqref{eq:ns-horizon-density-maximum-comparison}.
\end{proof}

\begin{theorem}[Horizon-density re-entry into the generated scale corona]
\label{thm:ns-horizon-density-scale-resaturation}
Retain the same amplified duration cell, maximum chart, and core radius as in
\cref{thm:ns-horizon-density-tangent}.  The centre trace of the public density
and its bounded corona coordinate are
\begin{equation}
 \label{eq:ns-horizon-centre-corona}
 \mathcal A_k^{\rm ctr}(s)
 :=\frac{1}{\nu|B_\delta|r(s)^2}
 =\frac{M(t(s))}{\nu|B_\delta|},
 \qquad
 \zeta_k(s):=\frac{1}{1+\nu|B_\delta|\mathcal A_k^{\rm ctr}(s)}
 =\frac{r(s)^2}{1+r(s)^2}.
\end{equation}
On the actual vorticity core,
\begin{equation}
 \label{eq:ns-horizon-core-centre-comparison}
 \mathcal A_k^{\rm ctr}
 \le \mathcal A_k^{\rm hor}
 \le4\mathcal A_k^{\rm ctr}.
\end{equation}
Consequently the horizon tangent laws satisfy
\begin{equation}
 \label{eq:ns-horizon-centre-corona-vanishing}
 0\le\int\zeta_k\,\dd\eta_k^{\rm hor}
 \le\frac{4}{\nu|B_\delta|}
      \frac{p_k^{\rm core}}{q_k^{\rm hor}}
 \longrightarrow0.
\end{equation}

The coordinate is not a new uncomputed asymptotic parameter.  With the
already generated scale rate \(a=\partial_s\log r\), it obeys the exact
public graph identity
\begin{equation}
 \label{eq:ns-horizon-centre-corona-transport}
 \partial_s\zeta_k=2a\,\zeta_k(1-\zeta_k).
\end{equation}
Thus its graph-complete boundary has exactly two dispositions.  If \(a\) is
nonordinary, its least finite failure is the existing \(\Lift\)-owned scale
successor.  If \(a\) is ordinary, \(\zeta=0\) is invariant and the branch
re-enters the normalized-origin dilation record of
\cref{thm:ns-forced-dilation-balance}; the persistent logarithmic scale
current is therefore handed, with the explicit \(1/5\) floor, to diffusion,
advection, stretching, time storage, or centre translation in
\eqref{eq:ns-dilation-five-source-balance}.  In particular the
horizon-density construction neither imports a whole-interval estimate nor
creates an unevaluated terminal atom: it resaturates a previously generated
public coordinate and resumes the prescribed source evaluation.
\end{theorem}

\begin{proof}
At the maximum-chart origin \(|\omega_v(0,s)|=1\), so the first identity in
\eqref{eq:ns-horizon-centre-corona} is the centre trace of
\eqref{eq:ns-horizon-price-measure-density}.  On the persistent core,
\(1/2\le|\omega_v|\le1\); hence
\(\mathcal A_k^{\rm hor}=\mathcal A_k^{\rm ctr}/|\omega_v|^2\), proving
\eqref{eq:ns-horizon-core-centre-comparison}.  It follows that
\[
 \zeta_k\le
 \frac{1}{\nu|B_\delta|\mathcal A_k^{\rm ctr}}
 \le\frac{4}{\nu|B_\delta|\mathcal A_k^{\rm hor}}.
\]
Integrating and using
\eqref{eq:ns-horizon-reciprocal-tangent} proves
\eqref{eq:ns-horizon-centre-corona-vanishing}.

Since \(\partial_s(r^2)=2ar^2\), direct differentiation of
\(r^2/(1+r^2)\) gives
\eqref{eq:ns-horizon-centre-corona-transport}.  Totalizing the multiplication
graph before restricting to \(\zeta=0\) yields the stated nonordinary
\(a\)-row.  On the ordinary cell the right side vanishes at \(\zeta=0\).
The same chart retains \(|\omega_v(0,s)|=1\), so
\cref{thm:ns-forced-dilation-balance} applies without a change of origin,
carrier, or event restriction and gives the five-term handoff.
\end{proof}

\begin{theorem}[Public Biot--Savart direction decomposition at the maximum]
\label{thm:ns-maximum-biot-savart-direction-decomposition}
On an ordinary maximum chart in \(\mathbb R^3\), put
\(e(s):=\xi(0,s)\).  The constraint-reconstruction graph gives the exact
principal-value identity
\begin{equation}
 \label{eq:ns-maximum-biot-savart-direction-identity}
 e^{\mathsf T}S(0,s)e
 =-\frac{3}{4\pi}\operatorname{p.v.}
   \int_{\mathbb R^3}
   \frac{(e\cdot\widehat y)
   \bigl((e\mathbin\times\omega_v(-y,s))\cdot\widehat y\bigr)}{|y|^3}
   \,\dd y .
\end{equation}
Equivalently, on \(\{\omega_v(-y,s)\ne0\}\), the integrand is
\begin{equation}
 \label{eq:ns-maximum-biot-savart-direction-only}
 |\omega_v(-y,s)|
 (e\cdot\widehat y)
 \bigl((e\mathbin\times\xi(-y,s))\cdot\widehat y\bigr)|y|^{-3}.
\end{equation}
It is zero on the zero-vorticity cell.  Hence magnitude variation parallel to
the centre direction contributes no maximum stretching: every nonzero
integrand records a generated vorticity-direction mismatch.

For \(R\ge1\), let \(F_R(s)\) be the part of
\eqref{eq:ns-maximum-biot-savart-direction-identity} over \(|y|>R\).
Then
\begin{equation}
 \label{eq:ns-biot-savart-far-l2-bound}
 |F_R(s)|\le C_{\rm BS}R^{-3/2}\|\omega_v(s)\|_2.
\end{equation}
On \(\mathbb T^3\), the same statements hold with the periodized public
Biot--Savart kernel: its Euclidean singular part gives
\eqref{eq:ns-maximum-biot-savart-direction-only}, while its smooth
periodic remainder is included in \(F_R\); the constant
\(C_{\rm BS,\mathbb T^3}\) is the explicit \(L^2\)-norm of that declared
kernel outside the generated ball.  A failed principal value, periodization,
direction, or kernel-integrability row is its least \(\Abs/\Geom\)-owned
total-graph successor.
\end{theorem}

\begin{proof}
For \(z=x-y\), the whole-space public reconstruction is
\[
 u(x)=\frac1{4\pi}\int
       \frac{\omega(y)\mathbin\times z}{|z|^3}\,\dd y.
\]
Differentiate, symmetrize, and contract twice with a unit vector \(e\).
The Kronecker term is antisymmetric and drops out, leaving
\[
 e^{\mathsf T}S(x)e
 =-\frac3{4\pi}\operatorname{p.v.}\int
   \frac{(e\cdot\widehat z)
   ((e\mathbin\times\omega(y))\cdot\widehat z)}{|z|^3}\,\dd y.
\]
The maximum-chart scaling and \(x=0\) give
\eqref{eq:ns-maximum-biot-savart-direction-identity}.  Writing
\(\omega_v=|\omega_v|\xi\) proves
\eqref{eq:ns-maximum-biot-savart-direction-only}; its value is defined as
zero when \(\omega_v=0\).  Finally Cauchy--Schwarz and
\(\||y|^{-3}\mathbf1_{|y|>R}\|_2=C R^{-3/2}\) give
\eqref{eq:ns-biot-savart-far-l2-bound}.  Periodization separates the same
local singular kernel from a smooth declared remainder, whose \(L^2\) norm
is a public kernel coordinate.
\end{proof}

\begin{theorem}[Generated near--far handoff to direction misalignment]
\label{thm:ns-generated-direction-misalignment-handoff}
Retain a maximum-centred interval \(J_j\), logarithmic scale current \(g_j\),
native price \(\mathfrak p_j\), and inverse-scale horizon \(H_j\) from
\cref{thm:ns-maximum-centered-strain-handoff,thm:ns-strain-modulus-horizon-transition}.
Let \(R_j\) be the least public dyadic radius \(R\ge1\) for which
\begin{equation}
 \label{eq:ns-generated-biot-savart-radius}
 \frac{C_{\rm BS}}{\sqrt\nu\,R^{3/2}}
 \sqrt{\mathfrak p_jH_j}\le\frac{g_j}{2}.
\end{equation}
This radius is computed from the same event record.  Define the nonnegative
near direction current
\begin{align}
 \label{eq:ns-near-direction-current}
 q_j^{\rm dir}:={}&\frac3{4\pi}\int_{J_j}\int_{0<|y|\le R_j}
 \frac{|\omega_v(-y,s)|}{|y|^3}
 |e(s)\cdot\widehat y|
 |(e(s)\mathbin\times\xi(-y,s))\cdot\widehat y|
 \,\dd y\,\dd s .
\end{align}
Then the exact maximum strain floor and the native carrier imply
\begin{equation}
 \label{eq:ns-near-direction-current-floor}
 q_j^{\rm dir}\ge\frac32g_j.
\end{equation}

More precisely, on the same near event put
\begin{align}
 \label{eq:ns-direction-carrier-density}
 \dd m_j^{\rm dir}
 &:=
 \nu r(s)|\omega_v(-y,s)|^2
 \mathbf1_{0<|y|\le R_j}\,\dd y\,\dd s,\\
 \mathcal A_j^{\rm dir}
 &:=
 \frac{3|e\cdot\widehat y|
 |(e\mathbin\times\xi(-y,s))\cdot\widehat y|}
 {4\pi\nu r(s)|\omega_v(-y,s)||y|^3},
 \end{align}
with density zero on \(\{\omega_v=0\}\).  These are public total-graph
coordinates and satisfy
\begin{equation}
 \label{eq:ns-direction-current-carrier-identity}
 q_j^{\rm dir}=\int\mathcal A_j^{\rm dir}\,\dd m_j^{\rm dir},
 \qquad
 m_j^{\rm dir}(1)\le\mathfrak p_j.
\end{equation}
Consequently, along the recurrent scale route
\(\sum_jg_j=\infty\), \(\sum_j\mathfrak p_j\le B_{\rm NS}\), the least
finite prefixes force
\begin{equation}
 \label{eq:ns-direction-polar-amplification}
 \sup_j\frac{q_j^{\rm dir}}{m_j^{\rm dir}(1)}=+\infty.
\end{equation}
The associated normalized current laws therefore satisfy
\begin{equation}
 \label{eq:ns-direction-reciprocal-tangent}
 \int\frac1{\mathcal A_j^{\rm dir}}\,\dd\eta_j^{\rm dir}
 =\frac{m_j^{\rm dir}(1)}{q_j^{\rm dir}}\longrightarrow0,
 \qquad
 \eta_j^{\rm dir}\{\mathcal A_j^{\rm dir}\le L\}\longrightarrow0
\end{equation}
for every finite \(L\), where
\(\dd\eta_j^{\rm dir}=(\mathcal A_j^{\rm dir}/q_j^{\rm dir})
\dd m_j^{\rm dir}\).
Thus the scale-resaturated route produces an explicit same-event
direction-misalignment polar successor.  Its next framework step is the
generated polarity split of the four displayed factors
\(r,|\omega_v|,|y|\), and
\(|(e\times\xi)\cdot\widehat y|\); no direction-coherence estimate is an
input hypothesis.
\end{theorem}

\begin{proof}
Integrating \eqref{eq:ns-biot-savart-far-l2-bound} and applying weighted
Cauchy--Schwarz gives
\[
 \int_{J_j}|F_R(s)|\,\dd s
 \le\frac{C_{\rm BS}}{\sqrt\nu R^{3/2}}
 \left(\nu\int_{J_j}r\|\omega_v\|_2^2\,\dd s\right)^{1/2}
 \left(\int_{J_j}r^{-1}\,\dd s\right)^{1/2}
 =\frac{C_{\rm BS}}{\sqrt\nu R^{3/2}}
  \sqrt{\mathfrak p_jH_j}.
\]
The dyadic radius in \eqref{eq:ns-generated-biot-savart-radius} exists and
makes this at most \(g_j/2\).  Pointwise, the positive part of the full
signed integral is bounded by the absolute near integrand plus \(|F_R|\).
Integrate and use
\eqref{eq:ns-maximum-strain-logarithmic-floor} to obtain
\(2g_j\le q_j^{\rm dir}+g_j/2\), proving
\eqref{eq:ns-near-direction-current-floor}.

Multiplication in \eqref{eq:ns-direction-carrier-density} gives exactly the
integrand of \eqref{eq:ns-near-direction-current}; restriction of the native
carrier gives the mass bound in
\eqref{eq:ns-direction-current-carrier-identity}.  Since the direction
currents have divergent sum and their carrier masses have sum at most
\(B_{\rm NS}\), finite-prefix carrier--polar pigeonholing proves
\eqref{eq:ns-direction-polar-amplification}.  The reciprocal identity and
finite-level estimate follow by the same direct probability calculation as
\eqref{eq:ns-horizon-reciprocal-tangent}.
\end{proof}

\begin{theorem}[Exhaustive polarity of the direction-density tangent]
\label{thm:ns-direction-density-polarity}
Pass to the direction-polar subsequence of
\cref{thm:ns-generated-direction-misalignment-handoff}.  For every
\(\varepsilon>0\), let
\[
 G_{j,\varepsilon}
 :=\{r(s)\ge\varepsilon,
       |\omega_v(-y,s)|\ge\varepsilon,\ |y|\ge\varepsilon\}.
\]
Since both angular factors in
\eqref{eq:ns-direction-carrier-density} are at most one,
\begin{equation}
 \label{eq:ns-direction-density-good-cell-bound}
 \mathcal A_j^{\rm dir}
 \le\frac{3}{4\pi\nu\varepsilon^5}
 \quad\hbox{on }G_{j,\varepsilon},
 \qquad
 \eta_j^{\rm dir}(G_{j,\varepsilon})\longrightarrow0.
\end{equation}
Consequently
\begin{equation}
 \label{eq:ns-direction-density-three-collapse}
 \eta_j^{\rm dir}
 \bigl(
  \{r<\varepsilon\}\cup
  \{|\omega_v(-y,s)|<\varepsilon\}\cup
  \{|y|<\varepsilon\}
 \bigr)\longrightarrow1.
\end{equation}
Diagonal extraction and finite sourcewise pigeonholing give one persistent
public polarity cell:
\begin{enumerate}[label=\textup{(DP\arabic*)}]
\item \(r\to0\) in direction-current law.  This is the already generated
  \(\Lift\)-owned scale corona and is cumulatively rejoined with
  \eqref{eq:ns-dilation-five-source-balance}.
\item \(|\omega_v(-y,s)|\to0\) in direction-current law.  The current
  approaches the zero-vorticity boundary through the displayed quotient
  density; totalization of \(\omega/|\omega|\) gives the least
  \(\Geom/\Caus\)-owned magnitude--direction graph successor.
\item \(|y|\to0\) in direction-current law.  The two-point direction graph
  meets its diagonal boundary.  Its generated dyadic modulus either has a
  positive lower radius, in which case the diagonal contribution is evaluated
  by the ordinary modulus and re-enters the carrier--polar calculation, or
  the radii collapse and give the \(\Geom/\Caus/\Lift\)-owned
  direction-modulus successor.
\end{enumerate}
These cells are exhaustive for the amplified density.  In particular,
direction coherence is neither assumed nor used as an unproved terminal:
its modulus is generated only on branch \textup{(DP3)}, after the scale and
zero-vorticity graph boundaries have been retained.
\end{theorem}

\begin{proof}
On \(G_{j,\varepsilon}\), the denominator of
\(\mathcal A_j^{\rm dir}\) is at least
\(4\pi\nu\varepsilon^5\), proving the pointwise bound in
\eqref{eq:ns-direction-density-good-cell-bound}.  The finite-level estimate
in \eqref{eq:ns-direction-reciprocal-tangent} then makes the
\(\eta_j^{\rm dir}\)-mass of this cell tend to zero.  Taking complements
gives \eqref{eq:ns-direction-density-three-collapse}.

Apply the generated zero/positive and dyadic polarity split to the three
bounded compactifications of \(r\), \(|\omega_v|\), and \(|y|\).
At least one of the three cells retains positive limiting current mass.
The first is the scale coordinate already present in the cumulative record.
The second is the total-graph boundary of the magnitude--direction
factorization used in
\eqref{eq:ns-maximum-biot-savart-direction-only}.  The third is the diagonal
boundary of the public two-point direction graph; its dyadic continuity
radius is generated by rational finite-cylinder recovery.  These are exactly
\textup{(DP1)}--\textup{(DP3)}.
\end{proof}

\begin{theorem}[Generated Dini direction current on the diagonal cell]
\label{thm:ns-generated-dini-direction-current}
For the same maximum chart and generated radius \(R_j\), define directly
from the public two-point direction graph
\begin{equation}
 \label{eq:ns-generated-direction-modulus}
 \mu_j(s,\rho)
 :=\operatorname*{ess\,sup}_{\substack{0<|y|\le\rho\\
                         \omega_v(-y,s)\ne0}}
   |\xi(-y,s)-e(s)|,
 \qquad 0<\rho\le R_j,
\end{equation}
and its event-local Dini current
\begin{equation}
 \label{eq:ns-generated-dini-current}
 \mathfrak D_j^{\rm dir}
 :=\int_{J_j}\int_0^{R_j}
     \mu_j(s,\rho)\,\frac{\dd\rho}{\rho}\,\dd s.
\end{equation}
The value is taken in the total graph \([0,+\infty]\); it is not an
admitted regularity modulus.  Maximum normalization
\(|\omega_v|\le1\), spherical disintegration of
\eqref{eq:ns-near-direction-current}, and
\(|e\times\xi|\le|e-\xi|\) give
\begin{equation}
 \label{eq:ns-direction-current-dini-bound}
 q_j^{\rm dir}\le3\mathfrak D_j^{\rm dir}.
\end{equation}
Therefore every persistent scale event satisfies the explicit target-aligned
floor
\begin{equation}
 \label{eq:ns-dini-direction-floor}
 \mathfrak D_j^{\rm dir}\ge\frac12g_j,
 \qquad
 \sum_j\mathfrak D_j^{\rm dir}=+\infty.
\end{equation}

If \(\mathfrak D_j^{\rm dir}=+\infty\) at a finite event, the least
failed Dini-integrability or direction graph is the
\(\Geom/\Caus\)-owned nonordinary successor.  If all values are finite,
the normalized currents
\(g_j^{-1}\mu_j(s,\rho)\rho^{-1}\dd\rho\dd s\) retain mass at least
\(1/2\); cumulative joining gives the ordinary recurrent
direction-oscillation successor.  Thus branch \textup{(DP3)} is evaluated
quantitatively: a singular scale route must finance a divergent sum of
generated Dini direction currents, even though no global direction modulus
is requested.
\end{theorem}

\begin{proof}
Use polar coordinates \(y=\rho\theta\) in
\eqref{eq:ns-near-direction-current}.  On the maximum chart,
\(|\omega_v(-\rho\theta,s)|\le1\), and
\[
 |e\cdot\theta|\,
 |(e\times\xi(-\rho\theta,s))\cdot\theta|
 \le|e-\xi(-\rho\theta,s)|
 \le\mu_j(s,\rho).
\]
The surface measure of \(\mathbb S^2\) is \(4\pi\); multiplying by
the coefficient \(3/(4\pi)\) proves
\eqref{eq:ns-direction-current-dini-bound}.  Combine it with
\eqref{eq:ns-near-direction-current-floor} to obtain the first inequality
in \eqref{eq:ns-dini-direction-floor}.  The recurrent scale cocycle has
\(\sum_jg_j=+\infty\), proving the divergent sum.  Totalization of the
essential supremum and Dini integral gives the nonordinary branch; otherwise
the displayed normalized measures and cumulative same-origin join give the
ordinary recurrent successor.
\end{proof}

\begin{theorem}[Dini-to-gradient handoff on the ordinary direction graph]
\label{thm:ns-dini-direction-gradient-handoff}
Restrict branch \textup{(DP3)} to the ordinary cell on which every radial
segment from \(0\) to \(-y\), \(0<|y|\le R_j\), remains in
\(\{\omega_v\ne0\}\).  A segment meeting \(\{\omega_v=0\}\) is
already the magnitude--direction graph successor \textup{(DP2)}.  On the
ordinary cell define the generated radial maximal gradient
\begin{equation}
 \label{eq:ns-generated-direction-gradient-maximal}
 \mathcal M_j^{\nabla\xi}(s,\rho)
 :=\operatorname*{ess\,sup}_{0<|z|\le\rho}|\nabla\xi(-z,s)|
\end{equation}
and its event current
\begin{equation}
 \label{eq:ns-generated-direction-gradient-current}
 \mathfrak G_j^{\rm dir}
 :=\int_{J_j}\int_0^{R_j}
 \mathcal M_j^{\nabla\xi}(s,\rho)\,\dd\rho\,\dd s.
\end{equation}
The fundamental theorem of calculus on each public radial segment gives
\begin{equation}
 \label{eq:ns-direction-modulus-gradient-bound}
 \mu_j(s,\rho)
 \le\rho\mathcal M_j^{\nabla\xi}(s,\rho),
 \qquad
 \mathfrak D_j^{\rm dir}\le\mathfrak G_j^{\rm dir}.
\end{equation}
Consequently every persistent ordinary diagonal event satisfies
\begin{equation}
 \label{eq:ns-direction-gradient-current-floor}
 \mathfrak G_j^{\rm dir}\ge\frac12g_j,
 \qquad
 \sum_j\mathfrak G_j^{\rm dir}=+\infty.
\end{equation}
An infinite value at a finite event is the least nonordinary
\(\Geom\)-owned gradient/maximal-graph successor.  Finite values define the
normalized ordinary gradient-current law
\[
 \dd\gamma_j^{\rm dir}
 :=\frac{
 \mathcal M_j^{\nabla\xi}(s,\rho)
 \,\dd\rho\,\dd s}{\mathfrak G_j^{\rm dir}},
\]
joined with the same scale, carrier, two-point direction, and source records.
Thus the direction-modulus branch is converted into an explicit derivative
current rather than retained as a qualitative coherence condition.
\end{theorem}

\begin{proof}
For \(|y|\le\rho\), ordinary evaluation along the segment gives
\[
 \xi(-y,s)-\xi(0,s)
 =-\int_0^1 y\cdot\nabla\xi(-\tau y,s)\,\dd\tau.
\]
Taking norms and the essential supremum proves the first inequality in
\eqref{eq:ns-direction-modulus-gradient-bound}.  Divide by \(\rho\),
integrate in \(\rho\) and \(s\), and use
\eqref{eq:ns-dini-direction-floor} to obtain
\eqref{eq:ns-direction-gradient-current-floor}.  The remaining statements
are the total-graph and normalized-law values of the displayed public
coordinates.
\end{proof}

\begin{theorem}[Exact two-point evolution of the generated direction graph]
\label{thm:ns-two-point-direction-evolution}
On the ordinary moving maximum chart put
\[
 U(y,s):=v(y,s)-b(s)-a(s)y,
 \qquad
 P_\xi:=I-\xi\otimes\xi,
 \qquad
 \mathcal N_\xi
 :=|\nabla\xi|^2\xi
   +2\nabla(\log|\omega_v|)\cdot\nabla\xi .
\]
The rescaled vorticity equation gives the exact direction row
\begin{equation}
 \label{eq:ns-moving-direction-equation}
 (\partial_s+U\cdot\nabla)\xi
 =\nu\Delta\xi+\nu\mathcal N_\xi+P_\xi S\xi.
\end{equation}
For \(\delta_yf:=f(-y,s)-f(0,s)\), direct subtraction and polarization
give the complete two-point identity
\begin{align}
 \label{eq:ns-two-point-direction-ledger}
 \frac12\partial_s|\delta_y\xi|^2
 ={}&
 \nu\,\delta_y\xi\cdot\delta_y(\Delta\xi)
 +\nu\,\delta_y\xi\cdot\delta_y\mathcal N_\xi
 +\delta_y\xi\cdot\delta_y(P_\xi S\xi)\\
 &-\delta_y\xi\cdot\delta_y(U\cdot\nabla\xi).
\end{align}
The four displayed currents have, in order, source supports
\[
 \{\Geom\},\qquad
 \{\Geom\},\qquad
 \{\Caus\},\qquad
 \{\Caus,\Lift\}.
\]
The logarithmic magnitude factor is evaluated before division; its zero or
nonordinary value is branch \textup{(DP2)}.  A failed Laplacian, gradient,
strain, transport, endpoint, or subtraction row is the corresponding least
typed total-graph successor.

On the finite ordinary cell, rational selectors realizing one half of the
essential suprema in
\eqref{eq:ns-generated-direction-gradient-maximal} turn
\eqref{eq:ns-two-point-direction-ledger} into actual finite-cylinder
identities.  Signed time gluing cancels compatible direction-storage
endpoints.  Every recurrent direction-gradient current is therefore
re-evaluated through the four written source currents and the uncancelled
endpoint row before the next carrier--polar projection.  No additional
evolution equation or direction estimate is supplied to this evaluator.
\end{theorem}

\begin{proof}
Rewrite \eqref{eq:ns-rescaled-vorticity-equation} as
\[
 (\partial_s+U\cdot\nabla)\omega_v
 =\nu\Delta\omega_v+S\omega_v+2a\omega_v.
\]
Its polar decomposition cancels the scalar dilation term and gives
\eqref{eq:ns-moving-direction-equation}.  Evaluate that identity at
\(-y\) and \(0\), subtract, and pair with \(\delta_y\xi\).  This proves
\eqref{eq:ns-two-point-direction-ledger}.  Diffusion and the polar
Laplacian terms descend from viscosity, stretching descends from convection,
and \(U=v-b-ay\) records convection together with centre and scale lift.
Rational near-maximizers exist by the definition of essential supremum;
localization, polarization, and signed gluing are finite public operations.
\end{proof}

\begin{theorem}[Exact localized evolution of the direction-gradient current]
\label{thm:ns-direction-gradient-localized-ledger}
On the ordinary nonzero-vorticity cell of
\cref{thm:ns-dini-direction-gradient-handoff}, set
\[
 B:=\nabla\xi,\qquad h:=\frac12|B|^2.
\]
Differentiating the public moving direction equation gives the exact matrix
row
\begin{equation}
 \label{eq:ns-direction-gradient-matrix-equation}
 (\partial_s+U\cdot\nabla)B
 =\nu\Delta B+\nu\nabla\mathcal N_\xi
   +\nabla(P_\xi S\xi)-B\nabla U .
\end{equation}
Its scalar square completion is
\begin{equation}
 \label{eq:ns-direction-gradient-pointwise-ledger}
 (\partial_s+U\cdot\nabla)h-\nu\Delta h
 +\nu|\nabla B|^2
 =\nu B:\nabla\mathcal N_\xi
  +B:\nabla(P_\xi S\xi)
  -B:(B\nabla U).
\end{equation}
For every rational compactly supported space--time cutoff \(\phi\), signed
integration gives the finite public carrier row
\begin{align}
 \label{eq:ns-localized-direction-gradient-ledger}
 \frac{\dd}{\dd s}\int\phi^2h\,\dd y
 +\nu\int\phi^2|\nabla B|^2\,\dd y
 ={}&
 \nu\int\phi^2 B:\nabla\mathcal N_\xi\,\dd y
 +\int\phi^2B:\nabla(P_\xi S\xi)\,\dd y\\
 &-\int\phi^2B:(B\nabla U)\,\dd y
 +\int h\bigl(
   2\phi\phi_s+\nu\Delta(\phi^2)
   +\nabla\cdot(\phi^2U)\bigr)\,\dd y .
\end{align}
The positive square has \(\Geom\) head.  The three interior currents on the
right have supports
\[
 \{\Geom\},\qquad
 \{\Geom,\Caus\},\qquad
 \{\Caus,\Lift\},
\]
and the final displayed row retains the complete
\(\{\Bdry,\Caus,\Lift\}\) cutoff and chart support.  In particular
\(\nabla\cdot U=-3a\), so the scale contribution remains an explicit
\(\Lift\)-owned term inside the last row.

The rational near-maximizers used for
\(\mathcal M_j^{\nabla\xi}\) generate cutoffs \(\phi\) with nonzero
direction-gradient incidence.  Equation
\eqref{eq:ns-localized-direction-gradient-ledger} is therefore the complete
next source ledger on every finite ordinary diagonal record.  A first
nonordinary Hessian, polar-diffusion, strain-gradient, deformation, cutoff,
or trace value is its typed successor.  When all values are ordinary, signed
gluing first cancels matching storage and cutoff faces; the remaining
interior currents and positive square enter the same native-carrier and
carrier--polar evaluation.  No unlisted global derivative estimate is part
of this transition.
\end{theorem}

\begin{proof}
For \(D_U:=\partial_s+U\cdot\nabla\),
\[
 D_U(\partial_k\xi_i)
 =\partial_k(D_U\xi_i)
  -(\partial_kU_\ell)(\partial_\ell\xi_i).
\]
Substitution of \eqref{eq:ns-moving-direction-equation} proves
\eqref{eq:ns-direction-gradient-matrix-equation}.  Pair with \(B\) and use
\(B:\Delta B=\Delta(|B|^2/2)-|\nabla B|^2\) to obtain
\eqref{eq:ns-direction-gradient-pointwise-ledger}.  Multiply by \(\phi^2\)
and integrate diffusion and transport by parts.  The time derivative of the
cutoff contributes \(2\phi\phi_s h\), diffusion contributes
\(\nu h\Delta(\phi^2)\), and transport contributes
\(h\nabla\cdot(\phi^2U)\), proving
\eqref{eq:ns-localized-direction-gradient-ledger}.  The source supports
follow from the displayed public constructors.
\end{proof}

\begin{theorem}[Explicit zero-carrier strain-amplification tangent]
\label{thm:ns-zero-carrier-strain-tangent}
Along every infinite cascade in
\cref{thm:ns-quantitative-strain-amplification}, there is a subsequence,
still indexed by \(j\), for which \(q_j/p_j\to\infty\).  Let
\[
 \dd m_j:=\nu r(s)|\omega(y,s)|^2\,\dd y\,\dd s
 \quad\text{on }I_j,
 \qquad
 \dd\eta_j:=\frac{\mathcal A_{k_j}}{q_j}\,\dd m_j.
\]
Then \(\eta_j\) is a probability measure and
\begin{equation}
 \label{eq:ns-reciprocal-amplification-tangent}
 \int_{\{\mathcal A_{k_j}>0\}}
 \frac{1}{\mathcal A_{k_j}}\,\dd\eta_j
 \le\frac{p_j}{q_j}\longrightarrow0.
\end{equation}
Consequently, for every finite \(M>0\),
\begin{equation}
 \label{eq:ns-finite-amplification-mass-vanishes}
 \eta_j\{\mathcal A_{k_j}\le M\}
 \le M\frac{p_j}{q_j}\longrightarrow0.
\end{equation}
Every weak graph-complete limit of the joined records is therefore supported
on the infinite-amplification boundary, has unit normalized strain current,
and has zero normalized kinetic-energy expenditure.  Its immutable source
support contains \(\Caus\) from strain and \(\Lift\) whenever the cofinal
scales vary.  This is the explicit NS zero-carrier tangent passed to the
primal--polar equality algebra.
\end{theorem}

\begin{proof}
The lower bound \eqref{eq:ns-linear-amplification-lower-bound} selects a
subsequence with \(q_j/p_j\to\infty\).  Identity
\eqref{eq:ns-strain-current-carrier-density} gives
\(\eta_j(I_j)=q_j^{-1}\int\mathcal A_{k_j}\,\dd m_j=1\).  On the positive
amplification set,
\[
 \int\mathcal A_{k_j}^{-1}\,\dd\eta_j
 =\frac{m_j\{\mathcal A_{k_j}>0\}}{q_j}
 \le\frac{p_j}{q_j},
\]
which proves \eqref{eq:ns-reciprocal-amplification-tangent}.  On
\(\{0<\mathcal A_{k_j}\le M\}\), one has
\(1\le M/\mathcal A_{k_j}\); integration gives
\eqref{eq:ns-finite-amplification-mass-vanishes}.  The zero-amplification set
has zero \(\eta_j\)-mass by definition.  Source-preserving graph completion
retains the strain eigenprojection and scale-chart words, proving the final
claim.
\end{proof}

\begin{theorem}[Exact Navier--Stokes target-spine source collapse]
\label{thm:ns-target-spine-source-collapse}
On \(\mathbb R^3\) or \(\mathbb T^3\), use the fixed physical chart and the
ordinary \(H^1\)-continuation coordinate
\[
 E_\omega(t)=\frac12\lVert\omega(t)\rVert_2^2,
 \qquad \omega=\nabla\times u.
\]
For a stopped shell interval \(I=[t_-,t_+]\) with
\(E_\omega(t_+)-E_\omega(t_-)=1\), expanded derivation of the public equation
gives the exact five-source row
\begin{align}
 q_{\{\Geom\}}(I)
 &:=-\nu\int_I\lVert\nabla\omega\rVert_2^2\,\dd t\le0,
 \\
 q_{\{\Caus\}}(I)
 &:=\int_I\int\omega^{\mathsf T}S\omega\,\dd x\,\dd t
   =1+\nu\int_I\lVert\nabla\omega\rVert_2^2\,\dd t\ge1,
 \\
 q_{\{\Abs\}}(I)&=q_{\{\Lift\}}(I)=q_{\{\Bdry\}}(I)=0.
\end{align}
Every exact-support current with \(|A|\ge2\), and every singleton not shown
above, is zero on this target spine.  Consequently the canonical actionable
amplitudes satisfy
\begin{equation}
 \label{eq:ns-exact-thirty-one-amplitudes}
 \alpha_{\{\Caus\}}(I)=q_{\{\Caus\}}(I)\ge1,
 \qquad
 \alpha_A(I)=0
 \quad(A\ne\{\Caus\}).
\end{equation}
Thus the full thirty-one-support evaluation is explicit: the only positive
interior owner of unit enstrophy-shell progress is vortex stretching.
\end{theorem}

\begin{proof}
Curl the fixed-chart equation.  The pressure term vanishes identically and
the solenoidal constraint has zero descent defect, giving the \(\Abs\) row.
There is no chart motion and no boundary on the stated domains, giving the
\(\Lift\) and \(\Bdry\) rows.  Pairing the vorticity equation with \(\omega\)
gives
\[
 E_\omega(t_+)-E_\omega(t_-)
 +\nu\int_I\lVert\nabla\omega\rVert_2^2\,\dd t
 =\int_I\int\omega^{\mathsf T}S\omega\,\dd x\,\dd t.
\]
The left increment is one.  Viscosity has primitive head \(\Geom\), and the
stretching term has primitive head \(\Caus\).  No product of distinct public
constructors remains after curl, integration by parts, and signed gluing, so
all mixed exact-support rows vanish.  Taking positive parts after this signed
grouping gives \eqref{eq:ns-exact-thirty-one-amplitudes}.
\end{proof}

\begin{theorem}[Explicit Navier--Stokes target-cyclic consequence hierarchy]
\label{thm:ns-target-cyclic-consequence-hierarchy}
On the fixed-chart whole-space or periodic target spine, every finite
differential, adjoint-transport, polarization, and commutator consequence
generated from the enstrophy shell is obtained from the following public
multiindex recurrence.  For every \(\alpha\in\mathbb N^3\),
\begin{equation}
 \label{eq:ns-multiindex-consequence}
 (\partial_t-\nu\Delta+u\cdot\nabla)\partial^\alpha u
 =-\partial^\alpha\nabla p
  -\sum_{0<\beta\le\alpha}
    \binom{\alpha}{\beta}
    (\partial^\beta u\cdot\nabla)
    \partial^{\alpha-\beta}u,
\end{equation}
with constraint rows
\begin{equation}
 \label{eq:ns-multiindex-pressure}
 \nabla\cdot\partial^\alpha u=0,
 \qquad
 -\Delta\partial^\alpha p
 =\partial^\alpha\partial_i\partial_j(u_i u_j).
\end{equation}
Pairing \eqref{eq:ns-multiindex-consequence} with
\(\partial^\alpha u\), or with any rational localization generated from the
shell covector, gives the exact signed row
\begin{align}
 \label{eq:ns-multiindex-energy-row}
 \frac12\frac{\dd}{\dd t}\lVert\partial^\alpha u\rVert_2^2
 +\nu\lVert\nabla\partial^\alpha u\rVert_2^2
 ={}&-\sum_{0<\beta\le\alpha}\binom{\alpha}{\beta}
 \left\langle
  (\partial^\beta u\cdot\nabla)
  \partial^{\alpha-\beta}u,
  \partial^\alpha u
 \right\rangle,
\end{align}
plus the explicitly oriented cutoff row for a localized pairing.  The
top-order transport and pressure pairings vanish after signed gluing and
constraint descent.  Thus every finite target-cyclic NS consequence has an
explicit finite sum, coefficient, source head, and pressure row; no
``adjoint compiler'' is invoked to supply its value.
\end{theorem}

\begin{proof}
Apply \(\partial^\alpha\) to the public equation and use the finite Leibniz
identity
\[
 \partial^\alpha((u\cdot\nabla)u)
 =(u\cdot\nabla)\partial^\alpha u
 +\sum_{0<\beta\le\alpha}\binom{\alpha}{\beta}
  (\partial^\beta u\cdot\nabla)
  \partial^{\alpha-\beta}u.
\]
Taking divergence of the original equation and using
\(\nabla\cdot u=0\) gives the pressure Poisson row; differentiation gives
\eqref{eq:ns-multiindex-pressure}.  Pair with
\(\partial^\alpha u\).  Divergence-free transport is skew, pressure is
orthogonal to the solenoidal field, and viscosity gives the displayed square.
Rational localization adds only the derivatives of its explicitly stored
cutoff.  These calculations generate precisely
\eqref{eq:ns-multiindex-energy-row}.
\end{proof}

\begin{theorem}[Exact strain--pressure successor equations]
\label{thm:ns-strain-pressure-successor}
Let \(A=\nabla u=S+\Omega\), with \(S=S^{\mathsf T}\) and
\(\Omega=-\Omega^{\mathsf T}\).  On the fixed-chart whole-space or periodic
NS target spine,
\begin{align}
 \label{eq:ns-gradient-matrix-equation}
 (\partial_t+u\cdot\nabla-\nu\Delta)A+A^2+\nabla^2p&=0,
 \\
 \label{eq:ns-strain-matrix-equation}
 (\partial_t+u\cdot\nabla-\nu\Delta)S+S^2+\Omega^2+\nabla^2p&=0,
 \\
 \label{eq:ns-pressure-strain-poisson}
 -\Delta p=\operatorname{tr}(A^2)
 =|S|^2-|\Omega|^2&.
\end{align}
On a simple eigenvalue cell \(SP_k=\lambda_kP_k\), choose a unit eigenvector
\(e_k\) in \(P_k\).  Its exact material derivative is
\begin{equation}
 \label{eq:ns-strain-eigenvalue-equation}
 (\partial_t+u\cdot\nabla)\lambda_k
 =\nu\,e_k^{\mathsf T}(\Delta S)e_k
  -\lambda_k^2
  -e_k^{\mathsf T}\Omega^2e_k
  -e_k^{\mathsf T}(\nabla^2p)e_k.
\end{equation}
For \(\ell\ne k\), the off-diagonal eigenprojection transport is
\begin{equation}
 \label{eq:ns-strain-eigenprojection-equation}
 P_\ell(\partial_t+u\cdot\nabla)P_k
 =\frac{P_\ell
  \bigl(\nu\Delta S-S^2-\Omega^2-\nabla^2p\bigr)P_k}
 {\lambda_k-\lambda_\ell}.
\end{equation}
The numerator terms have, respectively, \(\Geom\), \(\Caus\), \(\Caus\),
and \(\Abs\) heads.  A vanishing spectral gap is not divided out: the joint
eigenspace projection and its multiplicity flag form the finite
\(\Abs/\Lift\) total-graph row.  Equations
\eqref{eq:ns-strain-eigenvalue-equation}--
\eqref{eq:ns-strain-eigenprojection-equation} therefore explicitly evaluate
the next target-cyclic successor of the strain-amplification cell.
\end{theorem}

\begin{proof}
Differentiate the velocity equation.  The product rule gives
\(\nabla((u\cdot\nabla)u)=u\cdot\nabla A+A^2\), proving
\eqref{eq:ns-gradient-matrix-equation}.  Taking symmetric parts and using
\(\operatorname{sym}(A^2)=S^2+\Omega^2\) proves
\eqref{eq:ns-strain-matrix-equation}.  Taking divergence of NS gives the
pressure Poisson equation, while
\(\operatorname{tr}(A^2)=|S|^2-|\Omega|^2\).

On a simple eigenvalue cell, differentiating
\(Se_k=\lambda_ke_k\) and pairing with \(e_k\) gives
\((\partial_t+u\cdot\nabla)\lambda_k
=e_k^{\mathsf T}(\partial_t+u\cdot\nabla)Se_k\), which yields
\eqref{eq:ns-strain-eigenvalue-equation}.  Projecting the differentiated
eigenvalue equation from \(P_k\) to \(P_\ell\) gives
\eqref{eq:ns-strain-eigenprojection-equation}.  The stated source heads are
the primitive public constructors of the four displayed terms.
\end{proof}

\begin{theorem}[Exact NS vorticity-direction, strain, pressure, and localization descendants]
\label{thm:ns-complete-target-cyclic-descendants}
On the fixed-chart whole-space or periodic cell, put
\(\rho=|\omega|\) and \(\xi=\omega/\rho\) on \(\{\rho>0\}\).  The public
vorticity graph has the exact polar decomposition
\begin{align}
 \label{eq:ns-vorticity-magnitude-equation}
 D_t\rho
 &=\nu\bigl(\Delta\rho-\rho|\nabla\xi|^2\bigr)
   +\rho\,\xi^{\mathsf T}S\xi,
 \\
 \label{eq:ns-vorticity-direction-equation}
 D_t\xi
 &=\nu\bigl(\Delta\xi+|\nabla\xi|^2\xi
       +2\nabla(\log\rho)\cdot\nabla\xi\bigr)
   +(I-\xi\otimes\xi)S\xi .
\end{align}
The zero-vorticity cell is evaluated before division and is target-null for
every strain-direction current.  Diffusion and direction-gradient terms in
these equations have \(\Geom\) head; the two strain terms have \(\Caus\)
head.

The strain and pressure descendants satisfy the exact integral identities
\begin{align}
 \label{eq:ns-betchov-identity}
 \int\operatorname{tr}(S^3)\,\dd x
 &=-\frac34\int\omega^{\mathsf T}S\omega\,\dd x,
 \\
 \label{eq:ns-strain-energy-identity}
 \frac12\frac{\dd}{\dd t}\lVert S\rVert_2^2
 +\nu\lVert\nabla S\rVert_2^2
 &=\frac12\int\omega^{\mathsf T}S\omega\,\dd x,
 \\
 \label{eq:ns-pressure-hessian-identity}
 \lVert\nabla^2p\rVert_2^2
 &=\lVert\Delta p\rVert_2^2
 =\left\lVert |S|^2-\frac12|\omega|^2\right\rVert_2^2 .
\end{align}
Thus the cubic strain row is the same \(\Caus\)-owned target current as
vortex stretching, with coefficient \(-3/4\), while pressure reconstruction
is an explicit \(\Abs\)-owned quadratic graph coordinate rather than a
discarded solenoidal term.

For every rational smooth compactly supported cutoff \(\phi(x,t)\), the
localized vorticity carrier row is
\begin{align}
 \label{eq:ns-localized-vorticity-identity}
 \frac12\frac{\dd}{\dd t}\int\phi^2|\omega|^2\,\dd x
 +\nu\int\phi^2|\nabla\omega|^2\,\dd x
 ={}&\int\phi^2\omega^{\mathsf T}S\omega\,\dd x
 \\
 &+\int
 \bigl(\phi\phi_t+\nu\nabla\!\cdot(\phi\nabla\phi)
       +\phi u\cdot\nabla\phi\bigr)|\omega|^2\,\dd x .
\end{align}
Here the stretching term is \(\Caus\), the localized positive square is
\(\Geom\), the time and viscous cutoff terms are \(\Bdry\), and the
advected cutoff term has support \(\{\Caus,\Bdry\}\).  Moving-chart pullback
adds exactly the \(\Lift\) rows already displayed in
\eqref{eq:ns-moving-chart-equation} and
\eqref{eq:ns-rescaled-vorticity-equation}.  Consequently
\eqref{eq:ns-vorticity-magnitude-equation}--
\eqref{eq:ns-localized-vorticity-identity}, their rational derivatives,
polarizations, chart pullbacks, and pressure reconstructions instantiate the
full target-cyclic and five-normal-form successor family required by
steps \textup{(I7)}--\textup{(I9)}.  A fixed-point or rank conclusion may be
drawn only after these coordinates and their total-graph values have been
adjoined.
\end{theorem}

\begin{proof}
Write \(\omega=\rho\xi\) in
\(D_t\omega=\nu\Delta\omega+S\omega\).  Dotting with \(\xi\), using
\(\xi\cdot\Delta\xi=-|\nabla\xi|^2\), gives
\eqref{eq:ns-vorticity-magnitude-equation}.  Subtracting the resulting
magnitude component from the vector equation and dividing on \(\{\rho>0\}\)
gives \eqref{eq:ns-vorticity-direction-equation}.

For trace-free \(S\), \(\operatorname{tr}(S^3)=3\det S\).  Signed periodic
or whole-space integration by parts gives the Betchov relation
\(-4\int\det S=\int\omega^{\mathsf T}S\omega\), proving
\eqref{eq:ns-betchov-identity}.  Pair
\eqref{eq:ns-strain-matrix-equation} with \(S\).  The pressure-Hessian
pairing vanishes after signed integration by parts, and
\(S:\Omega^2=\frac14\omega^{\mathsf T}S\omega\); substitution of
\eqref{eq:ns-betchov-identity} gives
\eqref{eq:ns-strain-energy-identity}.  Fourier evaluation on
\(\mathbb R^3\) or \(\mathbb T^3\) gives
\(\|\nabla^2p\|_2=\|\Delta p\|_2\); use
\eqref{eq:ns-pressure-strain-poisson} and
\(|\Omega|^2=|\omega|^2/2\) to obtain
\eqref{eq:ns-pressure-hessian-identity}.

Finally pair the vorticity equation with \(\phi^2\omega\).  Integrate
diffusion and transport by parts without discarding cutoff derivatives.  The
result is \eqref{eq:ns-localized-vorticity-identity}.  Its source assignments
are exactly the public constructors that create its four terms.  Closure
under differentiation, polarization, chart pullback, pressure reconstruction,
and total graphs is the defining recursion of
\cref{def:target-cyclic-covector-module}, proving the final assertion.
\end{proof}

\begin{theorem}[Exact transport of the NS amplification coordinate]
\label{thm:ns-amplification-transport}
On a fixed-chart cell with \(\lambda_k>0\), simple spectral projection
\(P_k\), \(\omega\ne0\), and \(P_k\omega\ne0\), put
\[
 w_k:=P_k\omega,
 \qquad \theta_k:=\frac{|w_k|^2}{|\omega|^2},
 \qquad \mathcal A_k:=\nu^{-1}\lambda_k\theta_k,
 \qquad D_t:=\partial_t+u\cdot\nabla.
\]
Then
\begin{align}
 \label{eq:ns-amplification-log-transport}
 D_t\log\mathcal A_k
={}&\frac{\nu e_k^{\mathsf T}(\Delta S)e_k
 -\lambda_k^2-e_k^{\mathsf T}\Omega^2e_k
 -e_k^{\mathsf T}(\nabla^2p)e_k}{\lambda_k}
 \\
 &+2\frac{w_k\cdot
 \bigl((D_tP_k)\omega+P_k(\nu\Delta\omega+S\omega)\bigr)}{|w_k|^2}
 -2\frac{\omega\cdot(\nu\Delta\omega+S\omega)}{|\omega|^2},
\end{align}
where \(D_tP_k\) is the finite sum of the off-diagonal rows
\eqref{eq:ns-strain-eigenprojection-equation}.  In a moving chart the right
side acquires the explicit scale contribution \(-2a\) and the transported
versions of the displayed rows, because
\(\mathcal A_k=(\nu r^2)^{-1}\lambda_k\theta_k\).

Equation \eqref{eq:ns-amplification-log-transport} is the complete ordinary
amplification row.  The complementary cells
\[
 \lambda_k=0,
 \quad \lambda_k=\lambda_\ell,
 \quad \omega=0,
 \quad P_k\omega=0,
\]
are evaluated before division: the first and last are target-null for this
strain cell, the repeated-eigenvalue cell uses the joint projection, and
\(\omega=0\) has zero enstrophy-shell incidence.  Every nonordinary
Laplacian, pressure-Hessian, or projection value retains its displayed
\(\Geom/\Abs/\Lift\) graph row.
\end{theorem}

\begin{proof}
The vorticity equation is
\(D_t\omega=\nu\Delta\omega+S\omega\).  Therefore
\[
 D_tw_k=(D_tP_k)\omega+P_k(\nu\Delta\omega+S\omega).
\]
Logarithmic differentiation of
\(\theta_k=|w_k|^2/|\omega|^2\) gives the last two terms of
\eqref{eq:ns-amplification-log-transport}.  Logarithmic differentiation of
\(\lambda_k\), followed by
\eqref{eq:ns-strain-eigenvalue-equation}, gives the first line.  The moving
chart formula follows from \(D_s\log r=a\).  The total-graph dispositions are
the zero-denominator branches of these same finite identities.
\end{proof}

\begin{theorem}[Explicit NS equality-boundary and corona transition]
\label{thm:ns-equality-corona-transition}
On every ordinary amplification cell of
\cref{thm:ns-amplification-transport}, define the bounded public coordinate
\[
 z_k:=\frac{1}{1+\mathcal A_k}\in(0,1],
 \qquad
 F_k:=D_t\log\mathcal A_k,
\]
where \(F_k\) is the complete right side of
\eqref{eq:ns-amplification-log-transport}.  Then
\begin{equation}
 \label{eq:ns-compactified-amplification-transport}
 D_tz_k=-z_k(1-z_k)F_k.
\end{equation}
For the tangent laws of
\cref{thm:ns-zero-carrier-strain-tangent},
\begin{equation}
 \label{eq:ns-compactified-amplification-vanishing}
 0\le\int z_{k_j}\,\dd\eta_j
 \le\int_{\{\mathcal A_{k_j}>0\}}
       \mathcal A_{k_j}^{-1}\,\dd\eta_j
 \le\frac{p_j}{q_j}\longrightarrow0.
\end{equation}
Hence every graph-complete tangent has \(z=0\), unit normalized strain
current, and zero normalized kinetic-energy carrier.  Its next framework
transition is exactly one of:
\begin{enumerate}[label=\textup{(EC\arabic*)}]
\item a Laplacian, pressure-Hessian, eigenprojection, vorticity-direction, or
  scale term displayed in \(F_k\) is nonordinary or unbounded after target
  normalization; its least such finite row is the next
  \(\Geom/\Abs/\Caus/\Lift\)-owned successor;
\item every term of \(F_k\) has an ordinary finite graph value; then
  \eqref{eq:ns-compactified-amplification-transport} restricts to the exact
  identity \(D_tz=0\) on \(z=0\), so the infinite-amplification boundary is
  invariant and retains its unit target value.
\end{enumerate}
There is no third equality-cell value, and no target-nullity or carrier
contradiction is inferred in \textup{(EC2)} without a further displayed
identity.
\end{theorem}

\begin{proof}
Since \(\mathcal A_k=(1-z_k)/z_k\), logarithmic differentiation gives
\(D_t\log\mathcal A_k=-D_tz_k/[z_k(1-z_k)]\), proving
\eqref{eq:ns-compactified-amplification-transport} before compactification.
Moreover \((1+\mathcal A)^{-1}\le\mathcal A^{-1}\) for
\(\mathcal A>0\); the tangent estimate
\eqref{eq:ns-reciprocal-amplification-tangent} gives
\eqref{eq:ns-compactified-amplification-vanishing}.  Totalize each finite
term of the displayed formula for \(F_k\) in its written order.  A first
nonordinary value gives \textup{(EC1)}.  If every value is finite, the factor
\(z(1-z)\) vanishes at \(z=0\), yielding \textup{(EC2)}.  These are the two
values of the total graph of
\eqref{eq:ns-compactified-amplification-transport}.
\end{proof}

\begin{theorem}[Explicit NS primal--polar tangent functional]
\label{thm:ns-primal-polar-tangent}
Let \(\eta_j\) be the zero-carrier tangent laws, evaluated on the bounded
public cylinder algebra containing the compactified amplification \(z_j\),
the normalized selected strain-current coordinate \(Q_j\), the normalized
kinetic-energy price \(P_j\), and every bounded total-graph coordinate in
\cref{thm:ns-target-cyclic-consequence-hierarchy,thm:ns-strain-pressure-successor,thm:ns-amplification-transport}.
Every weak-* cluster functional \(L_{\rm NS}\) satisfies
\begin{equation}
 \label{eq:ns-primal-polar-tangent-values}
 L_{\rm NS}(1)=1,
 \qquad L_{\rm NS}(z)=0,
 \qquad L_{\rm NS}(Q)=1,
 \qquad L_{\rm NS}(P)=0,
\end{equation}
is nonnegative on every generated square and admitted production, and
annihilates every bounded exact NS graph identity and its negative.
It retains \(\Caus\) ownership and the \(\Lift\) word of every nonconstant
cofinal chart.

Let \(\mathcal K_{\rm NS}^{\rm tan}\) be the finite ordered cone generated by
those identities, squares, carrier productions, and the equality rows
\(z=0\), \(Q=1\), and \(P=0\).  Then
\begin{equation}
 \label{eq:ns-no-finite-tangent-contradiction}
 -1\notin\mathcal K_{\rm NS}^{\rm tan}.
\end{equation}
If a later normalized NS coordinate has no finite ordinary graph value, its
least such row is the explicit successor \textup{(EC1)}.  If all finite rows
are ordinary, \(L_{\rm NS}\) is the contact-derived primal--polar fixed
zero-carrier strain atom.  It is a prospective source-profile value with
origin \(\mathsf{Hyp}\), not a regularity contradiction and not an
independently realized singular orbit.
\end{theorem}

\begin{proof}
The bounded cylinder state space is compact.  Pass to a subnet on which all
of its countably many coordinates converge.  Positivity and every exact
finite graph identity pass to the limit.  Normalization of \(\eta_j\) gives
\(L_{\rm NS}(1)=L_{\rm NS}(Q)=1\).  Equation
\eqref{eq:ns-compactified-amplification-vanishing} gives
\(L_{\rm NS}(z)=0\), while \(p_j/q_j\to0\) gives
\(L_{\rm NS}(P)=0\).  Source words and origin are discrete closed
coordinates, so they also pass to the limit.

Every generator of \(\mathcal K_{\rm NS}^{\rm tan}\) has nonnegative
\(L_{\rm NS}\)-value, while exact identities have value zero.  If
\(-1\in\mathcal K_{\rm NS}^{\rm tan}\), applying \(L_{\rm NS}\) would give
\(-1\ge0\), a contradiction.  This proves
\eqref{eq:ns-no-finite-tangent-contradiction}.  Total-graph ordering gives
the stated nonordinary successor; absence of such a row gives the fixed
functional with its inherited hypothetical origin.
\end{proof}

\begin{theorem}[Explicit NS covariant-corona resaturation]
\label{thm:ns-covariant-corona-resaturation}
Let \(I_j=[t_j^-,t_j^+]\) be the disjoint fixed-chart shell intervals
selected in \cref{thm:ns-quantitative-strain-amplification}, contained in a
finite prospective contact interval \([0,T_*)\), and put
\(\ell_j=t_j^+-t_j^-\).  On \(\mathbb R^3\) or \(\mathbb T^3\), with no
exterior work, write
\[
 E(t):=\frac12\lVert u(t)\rVert_2^2,
 \qquad D(t):=\lVert\nabla u(t)\rVert_2^2,
 \qquad \theta_j(t):=t_j^+-t.
\]
The public kinetic-energy row and its generated terminal adjoint are
\begin{equation}
 \label{eq:ns-fixed-corona-carrier-adjoint}
 \frac{\dd E}{\dd t}+\nu D=0,
 \qquad \theta_j'=-1,
 \qquad \theta_j(t_j^+)=0.
\end{equation}
Consequently the covariant product row prescribed by
\cref{def:evaluated-covariant-corona-record} is the exact identity
\begin{equation}
 \label{eq:ns-fixed-corona-product}
 \frac{\dd}{\dd t}(\theta_jE)
 =-E-\theta_j\nu D,
 \qquad
 \ell_jE(t_j^-)
 =\int_{I_j}E\,\dd t
  +\nu\int_{I_j}\theta_jD\,\dd t.
\end{equation}
Normalize the three nonnegative coordinates in this row by the selected
strain current \(q_j\):
\begin{equation}
 \label{eq:ns-normalized-corona-coordinates}
 H_j:=\frac1{q_j}\int_{I_j}E\,\dd t,
 \quad
 V_j:=\frac{\nu}{q_j}\int_{I_j}\theta_jD\,\dd t,
 \quad
 R_j:=\frac{\ell_jE(t_j^-)}{q_j}=H_j+V_j.
\end{equation}
Along the zero-carrier subsequence of
\cref{thm:ns-zero-carrier-strain-tangent},
\begin{equation}
 \label{eq:ns-corona-coordinate-vanishing}
 0\le H_j\le3B_{\rm NS}\ell_j\longrightarrow0,
 \qquad
 0\le V_j\le\ell_j\frac{p_j}{q_j}\longrightarrow0,
 \qquad R_j\longrightarrow0.
\end{equation}
Thus the graph-complete tangent functional satisfies
\begin{equation}
 \label{eq:ns-corona-tangent-values}
 L_{\rm NS}(Q)=1,
 \qquad
 L_{\rm NS}(H)=L_{\rm NS}(V)=L_{\rm NS}(R)=0.
\end{equation}
The finite-financing and equality-null alternatives \textup{(V1)} and
\textup{(V2)} of
\cref{thm:complete-covariant-corona-evaluator} are therefore impossible on
this tangent: applying \(L_{\rm NS}\) to either claimed target domination
would give \(1\le0\).  The evaluator necessarily returns \textup{(V3)}.  A
least nonordinary finite NS coordinate is its typed successor.  When the
scalar corona coordinates are ordinary, the evaluator next adjoins, in the
fixed public order, every magnitude--direction, strain, pressure,
localization, chart, adjoint, polarization, and total-graph coordinate in
\cref{thm:ns-complete-target-cyclic-descendants}.  Each nonzero or
nonordinary target projection is re-entered with its exact source owner.
Equations \eqref{eq:ns-corona-tangent-values} alone therefore establish the
\textup{(V3)} transition; they establish neither a fixed point nor failure of
the generated Noetherian condition.
\end{theorem}

\begin{proof}
Equation \eqref{eq:ns-fixed-corona-carrier-adjoint} is the fixed-chart
specialization of the admitted kinetic-energy identity, and multiplication
by the explicitly solved adjoint \(\theta_j\) gives
\eqref{eq:ns-fixed-corona-product}.  Since \(q_j\ge1/3\),
\(E(t)\le B_{\rm NS}\), and \(0\le\theta_j\le\ell_j\),
\[
 H_j\le3B_{\rm NS}\ell_j,
 \qquad
 V_j\le\ell_j p_j/q_j.
\]
The intervals are disjoint subsets of \([0,T_*)\), so
\(\sum_j\ell_j\le T_*\) and \(\ell_j\to0\).  The selected subsequence has
\(p_j/q_j\to0\), proving
\eqref{eq:ns-corona-coordinate-vanishing}; the product identity gives the
same limit for \(R_j\).

Pass these bounded normalized coordinates to the cluster functional of
\cref{thm:ns-primal-polar-tangent}.  This proves
\eqref{eq:ns-corona-tangent-values}.  Every finite financing or
ordered-coercivity identity belongs to the ordered tangent algebra and is
annihilated by exact equality rows.  Positivity of \(L_{\rm NS}\) would turn
a domination of the unit coordinate \(Q\) by the zero coordinates
\(H,V,R\) into \(1\le0\); hence neither closing alternative occurs.
The remaining alternative of the covariant-corona evaluator is \textup{(V3)}.
The mandatory successor family is
\cref{thm:ns-complete-target-cyclic-descendants}; recursive carrier--polar
exhaustion is applied after those coordinates are joined and before any
fixed-point or rank test.  This is exactly the transition priority in
\cref{thm:complete-covariant-corona-evaluator}.
\end{proof}

\begin{theorem}[Explicit NS realization and terminal disposition]
\label{thm:ns-tangent-realization-disposition}
Every ordinary smooth finite stopped NS record has finite
\(\lambda_k\), \(r>0\), and finite graph values, hence
\[
 z_k=(1+\mathcal A_k)^{-1}>0
\]
on each positive strain-direction cell.  The fixed tangent of
\cref{thm:ns-primal-polar-tangent} has \(z=0\).  Consequently it is not an
ordinary evaluation state of the original NS continuation graph.

Its realization tree has the following exact disposition.  Every finite
cylinder neighborhood used to construct the tangent contains an actual
same-origin stopped NS strain record from the selected cascade.  A compatible
cofinal path therefore retains the weak equation, datum, shell order, source
word, and carrier ledger, but its amplification coordinate converges to the
nonordinary boundary \(z=0\).  The total continuation graph then returns
exactly one of:
\begin{enumerate}[label=\textup{(NR\arabic*)}]
\item an ordinary continuation of the same origin, in which case the
  continuation-failure target is not reached;
\item a same-origin noncontinuable terminal germ whose preterminal
  restrictions cross the cofinal enstrophy shells; this is semantic singular
  contact by the original branch;
\item a least failed domain, convergence, pressure, trace, or nonlinear-graph
  row, which remains the displayed \(\Geom/\Caus/\Abs/\Lift/\Bdry\)
  realization successor.
\end{enumerate}
The contact-derived tangent retains origin \(\mathsf{Hyp}\) in all three
cases.  It cannot itself satisfy the independent-origin clause of an accepted
singular certificate.
\end{theorem}

\begin{proof}
Smoothness on a finite stopped interval makes every pointwise strain,
vorticity, chart, and pressure coordinate finite, proving \(z_k>0\).
Equation \eqref{eq:ns-primal-polar-tangent-values} gives \(z=0\) on the
tangent, so ordinary evaluation is impossible.  By construction of the
cluster functional, every finite cylinder neighborhood with positive tangent
mass contains one of the actual stopped records used to define \(\eta_j\).
Diagonal restriction preserves the displayed public rows.  Totalization of
the continuation and realization graphs gives ordinary continuation,
ordinary noncontinuation with cofinal shell contact, or the least failed
finite row.  These are \textup{(NR1)}--\textup{(NR3)}.  Origin preservation
is \cref{thm:no-recycled-contact-realization}.
\end{proof}

\begin{theorem}[Navier--Stokes realization of the complete I1--I9 protocol]
\label{thm:ns-i1-i9-public-instantiation}
The public presentation of \cref{def:public-ns-target-presentation} realizes
every operation in the protocol of \cref{sec:pde-instantiation} as follows.
\begin{enumerate}[label=\textup{(I\arabic*)}]
\item The equation, divergence constraint, datum, domain, weak tests,
continuation graph, target, and rational approximation words are the declared
public NS graphs.
\item Products, pressure reconstruction, spectral projections, charts,
cutoffs, and weak limits use their total graphs; their ordinary equations and
first-failure values are displayed in
\cref{thm:ns-target-cyclic-consequence-hierarchy,thm:ns-strain-pressure-successor}.
\item Testing by \(u\) gives the unique admitted kinetic-energy carrier and
ceiling in \cref{prop:ns-generated-carrier-admission}; chart conjugation is
\eqref{eq:ns-moving-chart-conjugated-energy}.
\item Primitive and composite currents have the exact five-source ownership
of \cref{thm:ns-thirty-one-cell-public-compilation}; the fixed enstrophy spine
has the explicit thirty-one-cell values
\eqref{eq:ns-exact-thirty-one-amplitudes}.
\item The sourcewise closure is generated parametrically by the multiindex
equation \eqref{eq:ns-multiindex-consequence}, pressure row
\eqref{eq:ns-multiindex-pressure}, rational localization
\eqref{eq:ns-localized-vorticity-identity}, and their total graphs.
\item The shell current is the conjugated enstrophy identity
\eqref{eq:ns-conjugated-enstrophy-row}; every unit shell produces the actual
finite strain occurrence \eqref{eq:ns-strain-cell-one-third-floor}, the
concentration bound \eqref{eq:ns-strain-concentration-lower-bound}, and the
canonical collapsing scale \eqref{eq:ns-canonical-vorticity-scale}.
\item The minimal target spine deletes the zero \(\Abs,\Lift,\Bdry\) fixed-chart
incidences and retains the \(\Caus\)-owned stretching current, its viscous
\(\Geom\) square, and every descendant in
\eqref{eq:ns-vorticity-magnitude-equation}--
\eqref{eq:ns-localized-vorticity-identity}.
\item The five normal forms are instantiated respectively by vorticity
diffusion and direction gradients; stretching, strain determinant, and
eigenframe transport; pressure Poisson and Hessian reconstruction; the exact
moving-chart conjugations and concentration-scale graph; and the oriented cutoff row.  Their equations are
\eqref{eq:ns-vorticity-magnitude-equation}--
\eqref{eq:ns-localized-vorticity-identity} and
\eqref{eq:ns-moving-chart-equation}--
\eqref{eq:ns-conjugated-enstrophy-row}, with the scale realization values
\eqref{eq:ns-scale-normalized-carrier-values}, local cylinder row
\eqref{eq:ns-rescaled-local-energy-ledger}, and pressure split
\eqref{eq:ns-rescaled-pressure-split}.
\item Recurrent closure uses the quantitative amplification
\eqref{eq:ns-linear-amplification-lower-bound}, zero-carrier tangent
\eqref{eq:ns-reciprocal-amplification-tangent}, compactified equality row
\eqref{eq:ns-compactified-amplification-transport}, primal--polar values
\eqref{eq:ns-primal-polar-tangent-values}, and covariant-corona values
\eqref{eq:ns-corona-tangent-values}.  It then processes the least unprocessed
coordinate from \textup{(I7)}--\textup{(I8)} before realization or a rank
test.  Unbounded scale price is routed by
\cref{thm:ns-unbounded-rescaled-dissipation-routing}; an ordinary scale path
is routed by
\cref{thm:ns-ordinary-scale-profile-routing,thm:ns-profile-positive-geometric-occurrence,thm:ns-profile-scale-summability-transition,thm:ns-recurrent-scale-cohomology-successor,thm:ns-covariant-pressure-tail-elimination,thm:ns-dilation-kernel-terminal-rule,thm:ns-forced-dilation-balance,thm:ns-cumulative-dilation-handoff,thm:ns-two-owner-dilation-reduction,thm:ns-maximum-centered-strain-handoff,thm:ns-strain-modulus-horizon-transition,thm:ns-inverse-horizon-slow-collapse,thm:ns-horizon-dyadic-maximum-ledger,thm:ns-maximum-core-duration-corridor,thm:ns-duration-corridor-polar-amplification,thm:ns-horizon-density-tangent,thm:ns-horizon-density-scale-resaturation,thm:ns-maximum-biot-savart-direction-decomposition,thm:ns-generated-direction-misalignment-handoff}.
\end{enumerate}
Every listed quantity is derived from the public NS equation, datum, domain,
and target shell.  The protocol accepts no additional global estimate,
regularity criterion, carrier, rank, or realization certificate.
\end{theorem}

\begin{proof}
Items \textup{(I1)}--\textup{(I3)} are the defining public graphs and the
finite multiplier calculation of
\cref{prop:ns-generated-carrier-admission,thm:ns-moving-chart-lift-ledger}.
Items \textup{(I4)}--\textup{(I6)} follow from exact source compilation,
source collapse, and the unit-shell calculation.  The differentiated,
polarized, localized, pressure, direction, and chart identities in
\cref{thm:ns-target-cyclic-consequence-hierarchy,thm:ns-strain-pressure-successor,thm:ns-complete-target-cyclic-descendants}
give \textup{(I7)}--\textup{(I8)} with the displayed source heads.
Finally
\cref{thm:ns-quantitative-strain-amplification,thm:ns-strain-concentration-successor,thm:ns-concentration-scale-successor,thm:ns-scale-normalized-successor-ledger,thm:ns-rescaled-cylinder-ledger,thm:ns-unbounded-rescaled-dissipation-routing,thm:ns-ordinary-scale-profile-routing,thm:ns-profile-positive-geometric-occurrence,thm:ns-profile-scale-summability-transition,thm:ns-recurrent-scale-cohomology-successor,thm:ns-covariant-pressure-tail-elimination,thm:ns-dilation-kernel-terminal-rule,thm:ns-forced-dilation-balance,thm:ns-cumulative-dilation-handoff,thm:ns-two-owner-dilation-reduction,thm:ns-maximum-centered-strain-handoff,thm:ns-strain-modulus-horizon-transition,thm:ns-inverse-horizon-slow-collapse,thm:ns-horizon-dyadic-maximum-ledger,thm:ns-maximum-core-duration-corridor,thm:ns-duration-corridor-polar-amplification,thm:ns-horizon-density-tangent,thm:ns-horizon-density-scale-resaturation,thm:ns-maximum-biot-savart-direction-decomposition,thm:ns-generated-direction-misalignment-handoff,thm:ns-zero-carrier-strain-tangent,thm:ns-equality-corona-transition,thm:ns-primal-polar-tangent,thm:ns-covariant-corona-resaturation}
gives \textup{(I9)} in the required transition order.
\end{proof}

\begin{theorem}[Navier--Stokes instantiation without side certificates]
\label{thm:ns-no-side-certificate-instantiation}
Starting only from \cref{def:public-ns-target-presentation}, the engine
generates the carrier of \cref{prop:ns-generated-carrier-admission}, the
complete continuation target, its shell covectors, and all thirty-one exact
source cells.  On every cofinal enstrophy-shell contact route with finite
kinetic-energy ceiling \(B_{\rm NS}\), the framework computes actual disjoint
stopped events \(I_j\), strain-direction indices \(k_j\), carrier prices
\(p_j\), and target currents \(q_j\) satisfying
\begin{equation}
 \label{eq:ns-instantiated-quantitative-output}
 q_j\ge\frac13,
 \qquad
 \sum_{j=1}^{\infty}p_j\le B_{\rm NS},
 \qquad
 \max_{j\le N}\frac{q_j}{p_j}
 \ge\frac{N}{3B_{\rm NS}}.
\end{equation}
The corresponding pointwise amplification is exactly
\eqref{eq:ns-pointwise-strain-amplification}, and its essential supremum is
unbounded along the cascade.  On the subsequence selected by
\cref{thm:ns-zero-carrier-strain-tangent}, the normalized strain-current laws
also satisfy, for every finite \(M\),
\begin{equation}
 \label{eq:ns-instantiated-tangent-transition}
\eta_j\{\mathcal A_{k_j}\le M\}
 \le M\frac{p_j}{q_j}\longrightarrow0.
\end{equation}
The same actual events satisfy
\begin{equation}
 \label{eq:ns-instantiated-concentration-scale-output}
 \mathcal C_j\ge\frac{3\nu q_j^2}{p_j},
 \qquad
 M_j\ge\sqrt3\,\nu\frac{q_j}{p_j},
 \qquad
 r_j=M_j^{-1/2}\longrightarrow0.
\end{equation}
Thus the branch transition is quantitatively forced through the
infinite-amplification boundary into the explicit vorticity-concentration and
collapsing-scale cells.  At every finite
stage exactly one of the following values is returned:
\begin{enumerate}[label=\textup{(NS\arabic*)}]
\item a finite prefix whose generated carrier price exceeds
  \(B_{\rm NS}\);
\item a target-null cumulative joint record;
\item the explicit \(\Caus/\Lift\)-owned strain-amplification record
  \((I_j,k_j,q_j,p_j,\mathcal A_{k_j},\eta_j,
    \mathcal C_j,M_j,r_j)\), satisfying
  \eqref{eq:ns-instantiated-tangent-transition}, together with the
  concentration-scale row
  \eqref{eq:ns-instantiated-concentration-scale-output}, the
  scale-normalized carrier row
  \eqref{eq:ns-scale-normalized-carrier-values}, and the
  moving-chart identities \eqref{eq:ns-moving-chart-conjugated-energy} and
  \eqref{eq:ns-conjugated-enstrophy-row}; or
\item the least nonordinary pressure, spectral-projection, chart, boundary,
  or approximation row, with its displayed finite graph defect and exact
  five-source owner.
\end{enumerate}
The first two values close the route.  In the third, every finite
target-cyclic consequence is the explicit multiindex row
\eqref{eq:ns-multiindex-consequence}; strain and pressure evolve by
\eqref{eq:ns-strain-matrix-equation}--
\eqref{eq:ns-strain-eigenprojection-equation}; amplification evolves by
\eqref{eq:ns-amplification-log-transport}; and the equality boundary is
evaluated by \eqref{eq:ns-compactified-amplification-transport}.  Its
primal--polar value is \eqref{eq:ns-primal-polar-tangent-values}, and its
covariant-corona resaturation is the explicit calculation
\eqref{eq:ns-fixed-corona-carrier-adjoint}--
\eqref{eq:ns-corona-tangent-values}.  Before a fixed-point test, the
magnitude--direction, strain, pressure, localization, and chart descendants
are the explicit identities
\eqref{eq:ns-vorticity-magnitude-equation}--
\eqref{eq:ns-localized-vorticity-identity}.  Its ordinary/nonordinary
realization is exactly
\textup{(NR1)}--\textup{(NR3)} of
\cref{thm:ns-tangent-realization-disposition}.  The fourth value is processed
through its actual failed graph row.  No stage may be replaced by a global
estimate, user-supplied certificate, or unnamed compiler result.

The collapsing-scale component of \textup{(NS3)} is processed by the exact
local energy and pressure rows
\eqref{eq:ns-rescaled-local-energy-ledger}--
\eqref{eq:ns-rescaled-near-pressure-bound}.  The whole-space harmonic
pressure tail is eliminated in the covariant target valuation by
\eqref{eq:ns-covariant-pressure-tail-bound}--
\eqref{eq:ns-covariant-pressure-tail-vanishing}; a surviving pressure value
is therefore a displayed reconstruction, domain, or boundary successor.
Unbounded rescaled dissipation
returns \textup{(UD1)} or \textup{(UD2)}.  A bounded ordinary local path
returns one of the horizon values \textup{(HP1)}--\textup{(HP3)} and carries
the actual geometric occurrence
\begin{equation}
 \label{eq:ns-instantiated-profile-geometric-floor}
 d_\delta^{\Geom}\ge\frac{\nu}{8}|Q_\delta^-|>0,
 \qquad
 p_{j,\delta}^{\rm pull}=r_jd_{j,\delta}^{\Geom}\longrightarrow0.
\end{equation}
Its target projection is then null or the next
\(\{\Geom,\Lift\}\)-owned zero-price successor.
The ordinary-profile branch next evaluates the canonical continuity radii.
Their collapse gives \textup{(CR1)}.  A uniform radius gives
\begin{equation}
 \label{eq:ns-instantiated-scale-recurrence-output}
 p_j\ge r_j\frac{\nu}{8}|Q_\delta^-|,
 \qquad
 \sum_jr_j\le\frac{8B_{\rm NS}}{\nu|Q_\delta^-|},
 \qquad
 \sup_j\frac{\log(r_j/r_{j+1})}{p_j}=+\infty,
\end{equation}
and hence the normalized zero-price \(\Lift\) scale-rate law
\eqref{eq:ns-normalized-scale-rate-law}.  This law is rejoined with all NS
chart, vorticity, pressure, localization, and realization coordinates before
the next carrier--polar test.  Its ordinary fixed dilation kernel is
target-null by \cref{thm:ns-dilation-kernel-terminal-rule}.  In fact the
ordinary normalized origin forces \textup{(DK1)} with amplitude
\eqref{eq:ns-forced-origin-dilation}; the exact weak vorticity balance then
hands at least \(1/5\) of each persistent logarithmic scale current to one of
the five public terms in \eqref{eq:ns-dilation-five-source-balance}.  Signed
chart gluing then cancels compatible time and center-translation storage; on
the bounded-storage cell at least one of the diffusion, advection, or
stretching currents in \eqref{eq:ns-three-dilation-currents} retains
normalized cumulative amplitude at least \(1/3\).  The curl identity then
combines advection and stretching into the single public convection current
\eqref{eq:ns-curl-nonlinear-current}; diffusion or convection carries at
least \(1/2\) of the normalized logarithmic current.  A
failed origin or mollification row is the explicit graph successor
\textup{(DK2)}.
On an ordinary maximum-centered chart the magnitude equation preserves sign:
diffusion opposes the normalized maximum, and the strain direction satisfies
\eqref{eq:ns-maximum-strain-scale-rate}--
\eqref{eq:ns-maximum-strain-logarithmic-floor}.  Collapse of its strain
localization radius is \textup{(SR1)}.  A uniform radius gives the computed
price--horizon inequality \eqref{eq:ns-strain-horizon-price}; hence a
zero-price strain route forces the inverse-scale-horizon divergence
\eqref{eq:ns-inverse-scale-horizon-divergence} in \textup{(SR2)}.
The corresponding public horizon law has density
\(r^{-1}\dd s/H_j=r^{-3}\dd t/H_j\), and its covariant collapse speed
\(\kappa=-ar\) and normalized physical price both tend to zero by
\eqref{eq:ns-slow-collapse-values}--\eqref{eq:ns-slow-collapse-limit}.
On disjoint physical intervals this forces the concrete maximum-duration
ledger \eqref{eq:ns-three-halves-maximum-divergence}--
\eqref{eq:ns-dyadic-maximum-duration-divergence}.
The generated maximum-vorticity core either collapses in \textup{(MC1)} or
the native energy carrier proves
\eqref{eq:ns-square-root-maximum-budget}.  On the common persistent cell the
two public moment rows give the exact duration corridor
\eqref{eq:ns-dyadic-duration-corridor}.
Restricting the horizon current and native carrier to each same dyadic event
then gives \eqref{eq:ns-duration-current-price-sums}; its finite prefixes
force the carrier--polar amplification
\eqref{eq:ns-duration-polar-amplification}.
Its same-event Radon--Nikodym density is the explicit
\(\mathcal A_k^{\rm hor}\) in
\eqref{eq:ns-horizon-price-measure-density}; the normalized laws satisfy
\eqref{eq:ns-horizon-reciprocal-tangent} and are quantitatively comparable
to the dyadic maximum by
\eqref{eq:ns-horizon-density-maximum-comparison}.  Its centre trace
compactifies to \(\zeta=r^2/(1+r^2)\), satisfies
\eqref{eq:ns-horizon-centre-corona-transport}, and therefore re-enters the
already generated dilation balance rather than opening an unrelated global
estimate branch.  Constraint reconstruction on that same maximum chart gives
the direction-only Biot--Savart identity
\eqref{eq:ns-maximum-biot-savart-direction-identity}.  The generated dyadic
radius \eqref{eq:ns-generated-biot-savart-radius} charges the far field to
the native carrier and leaves the actual near direction current
\eqref{eq:ns-near-direction-current} with amplitude at least \(3g_j/2\).
Its exact carrier density
\eqref{eq:ns-direction-carrier-density} produces the polar tangent
\eqref{eq:ns-direction-reciprocal-tangent}.  Thus direction coherence is
evaluated as a generated polarity coordinate; it is not supplied as an NS
hypothesis.  The exhaustive scale, zero-vorticity, and diagonal values of
that coordinate are \textup{(DP1)}--\textup{(DP3)} in
\cref{thm:ns-direction-density-polarity}.  On the diagonal value, the
generated two-point modulus has the quantitative Dini current
\eqref{eq:ns-generated-dini-current}, forced by
\eqref{eq:ns-dini-direction-floor}; no global coherence modulus is appended.
On its ordinary nonzero-vorticity cell, radial differentiation produces the
\(\Geom\)-owned current
\eqref{eq:ns-generated-direction-gradient-current} with floor
\eqref{eq:ns-direction-gradient-current-floor}.  Its public evolution is
the four-source two-point ledger
\eqref{eq:ns-two-point-direction-ledger}.  Differentiation and rational
localization then give the complete gradient-energy row
\eqref{eq:ns-localized-direction-gradient-ledger}, including its positive
Hessian square and all cutoff and lift terms.

If every currently evaluated row in \textup{(NS3)} remains ordinary,
recursive target-aligned continuation adjoins the next coordinate of
\cref{thm:ns-complete-target-cyclic-descendants}, cumulatively joins its
complete history, and repeats carrier--polar evaluation.  No fixed atom,
self-edge, or rank value is inferred before that generated family has been
processed.
\end{theorem}

\begin{proof}
Carrier admission is
\cref{prop:ns-generated-carrier-admission}.  Exact scale accounting and lift
conjugation are
\cref{thm:ns-exact-scale-price,thm:ns-moving-chart-lift-ledger}.  On every
unit enstrophy shell,
\cref{thm:ns-rescaled-vorticity-successor,thm:ns-quantitative-strain-cell}
give an actual \(\Caus\)-owned strain-direction event with current at least
\(1/3\).  Carrier additivity and
\cref{thm:ns-quantitative-strain-amplification} give
\eqref{eq:ns-instantiated-quantitative-output} and its pointwise density.
The trace-free strain calculation and essential-supremum scale selector in
\cref{thm:ns-strain-concentration-successor,thm:ns-concentration-scale-successor}
give \eqref{eq:ns-instantiated-concentration-scale-output}.
The rescaled equation and its exact covariant physical ownership are
\cref{thm:ns-scale-normalized-successor-ledger}.
Local scale closure, unbounded-dissipation routing, horizon routing, and the
profile occurrence floor are respectively
\cref{thm:ns-rescaled-cylinder-ledger,thm:ns-unbounded-rescaled-dissipation-routing,thm:ns-ordinary-scale-profile-routing,thm:ns-profile-positive-geometric-occurrence}.
The continuity-radius, summable-scale, recurrent scale-cohomology,
covariant pressure-tail, and forced dilation-handoff values
are
\cref{thm:ns-profile-scale-summability-transition,thm:ns-recurrent-scale-cohomology-successor,thm:ns-covariant-pressure-tail-elimination,thm:ns-dilation-kernel-terminal-rule,thm:ns-forced-dilation-balance,thm:ns-cumulative-dilation-handoff,thm:ns-two-owner-dilation-reduction,thm:ns-maximum-centered-strain-handoff,thm:ns-strain-modulus-horizon-transition,thm:ns-inverse-horizon-slow-collapse,thm:ns-horizon-dyadic-maximum-ledger,thm:ns-maximum-core-duration-corridor,thm:ns-duration-corridor-polar-amplification,thm:ns-horizon-density-tangent,thm:ns-horizon-density-scale-resaturation,thm:ns-maximum-biot-savart-direction-decomposition,thm:ns-generated-direction-misalignment-handoff}.
The reciprocal-density computation in
\cref{thm:ns-zero-carrier-strain-tangent} gives
\eqref{eq:ns-instantiated-tangent-transition} directly.
Exact support collapse is \cref{thm:ns-target-spine-source-collapse}.  The
remaining finite consequence, strain--pressure, amplification, equality,
primal--polar, and realization transitions are respectively
\cref{thm:ns-target-cyclic-consequence-hierarchy,%
thm:ns-strain-pressure-successor,%
thm:ns-complete-target-cyclic-descendants,%
thm:ns-amplification-transport,%
thm:ns-equality-corona-transition,%
thm:ns-primal-polar-tangent,%
thm:ns-covariant-corona-resaturation,%
thm:ns-tangent-realization-disposition}.
If a required public row is nonordinary, its total graph returns the least
failed finite row, giving \textup{(NS4)}.  Otherwise the quantitative record
is \textup{(NS3)}.  A target-null joint projection is \textup{(NS2)}, and a
least over-budget prefix is \textup{(NS1)}.  These cases are exhaustive by
expanded-derivation no-creation and the public graph order.
On the all-ordinary branch, the next recursive transition is the least
unprocessed coordinate in
\eqref{eq:ns-vorticity-magnitude-equation}--
\eqref{eq:ns-localized-vorticity-identity}; it is evaluated before any
Noetherian rank comparison.
\end{proof}

\begin{theorem}[Navier--Stokes 31-cell source compilation]
\label{thm:ns-thirty-one-cell-public-compilation}
For the public presentation of
\cref{def:public-ns-target-presentation}, every finite same-origin
pre-contact shell record has an exact source word in
\[
 \mathcal P(\{\Geom,\Caus,\Abs,\Lift,\Bdry\})\setminus\{\varnothing\}.
\]
The primitive assignments are: viscosity and its positive local production
are \(\Geom\); convection is \(\Caus\); the solenoidal constraint and
pressure reconstruction are \(\Abs\); a moving centre, scale, or normalized
chart is \(\Lift\); and cutoff, exterior, initial/terminal, and boundary
fluxes are \(\Bdry\).  A record formed from several of these operations is
assigned their union before target projection.  Thus each target-visible
Navier--Stokes shell contribution belongs to one of the same thirty-one cells
and is evaluated by the public-presentation cell rule
\cref{thm:cellwise-public-presentation-thirty-one-evaluation}.  Equivalently,
the Navier--Stokes specialization realizes the explicit theorem family
\(\mathsf T_A\) of
\cref{thm:explicit-thirty-one-source-theorem-family} for each of its
thirty-one possible exact supports.
\end{theorem}

\begin{proof}
Pair the public weak equation with the finite shell covector and expand the
local-energy identity before taking a positive part.  The viscous term is the
positive symmetric form and hence has head \(\Geom\); the advective term is
the fixed nonsymmetric transport current and hence has head \(\Caus\).
Applying the Leray constraint or reconstructing pressure is precisely the
quotient/constraint constructor, while every cutoff trace is a boundary
constructor and every moving normalization is a lift constructor.  These are
the five primitive clauses of the source compiler.  Structural induction on
the finite expression graph assigns the union word to every product,
commutator, localization, and signed gluing before target projection.
Source-preserving completion retains that word for weak defects,
concentrations, oscillations, tails, and failed graph values.  Therefore
\cref{thm:universal-five-source-provenance} and
\cref{thm:no-additional-continuum-singularity-source} give the stated
disjoint thirty-one-cell routing.  Finally,
\cref{thm:cellwise-public-presentation-thirty-one-evaluation} evaluates each
cell on finite same-origin records and applies its physical carrier rule.
The tabled theorem family applies word-by-word, so no combined
Navier--Stokes contribution is discharged by a theorem for an isolated
channel.
\end{proof}

\begin{theorem}[Navier--Stokes local shell-current computation]
\label{thm:ns-local-shell-current-computation}
Let \(I=[s,t]\) be a finite stopped interval on an ordinary same-origin
Navier--Stokes branch, and let \(\Phi\) be its public normalized target-shell
observable, with \(\Phi(u_s)=0\) and \(\Phi(u_t)=1\).  In the unprojected
public weak equation
\[
\partial_t u-\nu\Delta u+(u\cdot\nabla)u+\nabla p=0,
\qquad \nabla\cdot u=0,
\]
the exact target-current computation is
\begin{equation}
\label{eq:ns-exact-local-shell-ledger}
1=Q^{\Geom}_{I}+Q^{\Caus}_{I}+Q^{\Abs}_{I}+Q^{\Lift}_{I}+Q^{\Bdry}_{I}.
\end{equation}
where
\begin{align}
Q^{\Geom}_{I}&:=\nu\int_I\langle D\Phi(u_r),\Delta u_r\rangle\,\dd r,
\label{eq:ns-geom-current}\\
Q^{\Caus}_{I}&:=-\int_I\langle D\Phi(u_r),(u_r\cdot\nabla)u_r\rangle\,\dd r,
\label{eq:ns-caus-current}\\
Q^{\Abs}_{I}&:=-\int_I\langle D\Phi(u_r),\nabla p_r\rangle\,\dd r,
\label{eq:ns-abs-current}\\
Q^{\Lift}_{I}&:=\int_I\partial_r^{\rm chart}\Phi(u_r)\,\dd r,
\label{eq:ns-lift-current}\\
Q^{\Bdry}_{I}&:=\int_I\operatorname{Flux}_{\partial\Omega,\Phi}(u_r,p_r)\,\dd r.
\label{eq:ns-bdry-current}
\end{align}
Here lift and boundary rows are zero when the selected shell has no moving
chart and no boundary/exterior interface.  Its physical price is
\(p_I:=\mu_{\rm NS}^{\Phys}(I)\), the restriction of the public local-energy
carrier.  For pairwise disjoint stopped shell intervals \(I_1,\ldots,I_N\),
all five currents and all thirty-one exact-support currents are accumulated
by unnormalised sums.
\end{theorem}
\begin{proof}
Pair the weak equation with \(D\Phi(u_r)\), apply the chain rule on
the finite stopped interval, and retain pressure before constraint descent.
The moving-chart derivative is the lift correction; integration by parts on
the declared domain gives the oriented boundary flux.  This proves the
identity.  The local-energy identity defines the sole physical carrier.
Additivity on disjoint intervals proves the final assertion.  Every mixed
term is joined before source projection and is evaluated by its exact-support
local theorem.
\end{proof}

\begin{proposition}[Navier--Stokes geometric target relation]
\label{prop:ns-geometric-target-relation}
For every finite stopped shell interval \(I\), the geometric target current is
exactly \(Q_I^{\Geom}\) in \eqref{eq:ns-geom-current}.  Its target detector is
\[
 \delta_{\Geom}(I)=\left\|
 (G^{\Geom}_\Sigma,\operatorname{Ker}^{\Geom}_\Sigma,
   \operatorname{Gc}^{\Geom}_\Sigma,
   \operatorname{Inc}^{\Sigma}_{\Geom}(I))\right\|_{\rm det}.
\]
The positive physical production entering its record is the viscosity term in
the public local-energy carrier; no target alignment is inferred from its
unsigned size before the displayed incidence coordinate is nonzero.
\end{proposition}

\begin{proof}
The weak Laplacian is paired with \(D\Phi\) in
\eqref{eq:ns-geom-current}.  The source compiler assigns this symmetric
mobility term to \(\Geom\), and the local-energy identity assigns its
nonnegative viscous production to the physical carrier.  The detector is the
geometric packet together with the complete-frontier target incidence, which
is precisely the finite local \(\Geom\)-check.
\end{proof}

\begin{proposition}[Navier--Stokes directed target relation]
\label{prop:ns-caus-target-relation}
For every finite stopped shell interval \(I\), the directed target current is
exactly \(Q_I^{\Caus}\) in \eqref{eq:ns-caus-current}, with detector
\[
 \delta_{\Caus}(I)=\left\|
 (H^{\Caus},S^{\Caus},\Psi^{\Caus},\operatorname{Gc}^{\Caus},
   \operatorname{Inc}^{\Sigma}_{\Caus}(I))\right\|_{\rm det}.
\]
Its physical effect is evaluated only after signed gluing with the other
currents in the same complete local record.
\end{proposition}

\begin{proof}
Substitution of the convective term from the public weak equation gives
\eqref{eq:ns-caus-current}.  It is the fixed nonsymmetric transport clause of
the source compiler, hence has \(\Caus\) head.  Since transport is signed,
the target incidence is retained before the physical positive part is formed.
This is exactly the stated detector and gluing rule.
\end{proof}

\begin{proposition}[Navier--Stokes constraint target relation]
\label{prop:ns-abs-target-relation}
For every finite stopped shell interval \(I\), the constraint/pressure target
current is exactly \(Q_I^{\Abs}\) in \eqref{eq:ns-abs-current}, with detector
\[
 \delta_{\Abs}(I)=\left\|
 (A^{\rm desc},A^{\rm vert},A^{\rm coker},A^{\rm gc},A^{\rm tgt},
   \operatorname{Inc}^{\Sigma}_{\Abs}(I))\right\|_{\rm det}.
\]
It is target-null exactly when this complete quotient/constraint incidence
vanishes; pressure may not be discarded merely because a solenoidal test
would annihilate it.
\end{proposition}

\begin{proof}
The pressure term is the weak constraint descent of the public equation and
therefore has \(\Abs\) head.  Pairing it with the actual shell covector gives
\eqref{eq:ns-abs-current}.  The five displayed quotient coordinates and the
joint incidence are exactly its public finite record, proving the claim.
\end{proof}

\begin{proposition}[Navier--Stokes lift target relation]
\label{prop:ns-lift-target-relation}
For every finite stopped shell interval \(I\), chart, centre, scale, and
normalization motion contributes exactly \(Q_I^{\Lift}\) in
\eqref{eq:ns-lift-current}, with detector
\[
 \delta_{\Lift}(I)=\left\|
 (L^{\rm int},L^{\rm fib},L^{\rm conc},L^{\rm osc},L^{\rm sc},L^{\rm gc},
   \operatorname{Inc}^{\Sigma}_{\Lift}(I))\right\|_{\rm det}.
\]
When the public shell has a fixed chart this row is identically zero; a
moving normalization is charged only through its actual target incidence.
\end{proposition}

\begin{proof}
Differentiate the public shell observable in its declared moving chart.  The
additional chain-rule term is \eqref{eq:ns-lift-current}; it is exactly the
lift constructor.  The stated detector is its complete finite lift packet,
and the fixed-chart specialization gives zero directly.
\end{proof}

\begin{proposition}[Navier--Stokes boundary target relation]
\label{prop:ns-bdry-target-relation}
For every finite stopped shell interval \(I\), boundary, exterior, cutoff,
and terminal-flux contributions are exactly \(Q_I^{\Bdry}\) in
\eqref{eq:ns-bdry-current}, with detector
\[
 \delta_{\Bdry}(I)=\left\|
 (B^{\rm ext},B^{\rm tail},B^{\rm term},B^{\rm gc},
   \operatorname{Inc}^{\Sigma}_{\Bdry}(I))\right\|_{\rm det}.
\]
The physical boundary work is the boundary/exterior component of the same
public local-energy carrier, after cancellation of internal faces.
\end{proposition}

\begin{proof}
Integration by parts in the public local-energy identity gives the oriented
flux in \eqref{eq:ns-bdry-current}.  The trace/exterior constructor assigns
its \(\Bdry\) head, and the displayed boundary packet records every possible
target-visible terminal or tail component.  Internal-face cancellation occurs
before the carrier is restricted, proving the asserted accounting relation.
\end{proof}

\begin{corollary}[Navier--Stokes finite-cascade quantitative ledger]
\label{cor:ns-finite-cascade-quantitative-ledger}
For every finite same-origin shell cascade \(\mathcal F_N=(I_1,\ldots,I_N)\),
the public presentation computes
\begin{equation}
\label{eq:ns-finite-cascade-exact-ledger}
N=\sum_{\varnothing\ne A\subseteq\Channels}\sum_{j=1}^{N}Q_{A,I_j},
\qquad \mu_{\rm NS}^{\Phys}(\bigcup_{j=1}^{N}I_j)
=\sum_{\varnothing\ne A\subseteq\Channels}\mathsf P_A(\mathcal F_{A,N}).
\end{equation}
Each \(\mathsf A_A(\mathcal F_{A,N})\), \(\mathsf D_A(\mathcal F_{A,N})\),
and \(\mathsf P_A(\mathcal F_{A,N})\) is the explicit local quantity
in the theorem labelled by \(A\).  No combination of subthreshold currents
is averaged away or left outside the physical ledger.
\end{corollary}
\begin{proof}
Sum the local shell identity over the finite disjoint family, then partition
each expanded current by exact provenance support.  Finite additivity of the
single positive physical carrier gives the second identity.  The thirty-one
local theorems retain the complete joint frontier of every summand.
\end{proof}

\begin{corollary}[End-to-end Navier--Stokes contact audit]
\label{cor:ns-end-to-end-contact-audit}
An ordinary Navier--Stokes branch reaches \(\Sigma_{\rm NS}\) only if its
public presentation generates a finite same-origin target-aligned occurrence
in at least one of the thirty-one cells of
\cref{thm:ns-thirty-one-cell-public-compilation}.  Conversely, an ordinary
same-origin record satisfying the continuation-graph contact rows is a
target-reaching branch.  In either direction, the contribution is evaluated
solely by its complete joint frontier and the single local-energy physical
carrier; no whole-solution estimate, compactness premise, or external
regularity certificate enters the audit.
\end{corollary}

\begin{proof}
For contact, the normalized shell identity supplies unit pre-contact
progress.  The no-additional-source theorem and the preceding compilation
place a positive contribution in a unique exact-support cell; the cellwise
theorem returns its finite source-owned occurrence and carrier valuation.
Conversely, the ordinary realization and noncontinuation rows are precisely
the public continuation-graph conditions defining \(\Sigma_{\rm NS}\).
This is the general target-contact equivalence restricted to the displayed
public presentation.
\end{proof}

\section{Engine acceptance suite}
\label{sec:engine-acceptance-suite}

\begin{theorem}[Acceptance suite for the binary engine]
\label{thm:binary-engine-acceptance-suite}
The following eleven test families are discharged by the engine contract.
\begin{enumerate}[label=\textup{(AT\arabic*)}]
\item The dissipative scalar flow \(\dot x=-x\), with target escape to
  infinity, has the finite regularity certificate furnished by its decreasing
  Lyapunov carrier.
\item The scalar flow \(\dot x=x^2\), \(x(0)=1\), has finite physical budget
  for the declared clock carrier \(\mu^{\Phys}=\dd t\) on \([0,1)\), and the
  independently checkable contacting orbit
  \(x(t)=(1-t)^{-1}\); it is assigned the singular output.
\item A zero-price recurrent record whose generated target-cyclic pairings
  all vanish is removed in \(\Null_\Sigma\) after multiplicity restoration
  and kernel restriction.
\item A zero-price equality record with a nonzero feeding or storage
  coordinate re-enters the resaturation queue and is not an output.
\item A nonordinary graph value retains its least finite detector, source
  ancestry, and actual contextual successor.
\item A mixed operation that is target-active although its separate inputs
  are target-null is retained on a finite joint prefix with its complete
  support and physical owner.
\item Each subtoken of a current acquiring \(\Lift\) or \(\Bdry\) ancestry is
  debited from its parent support and credited once to the superset determined
  by that constructor.
\item Every public approximation, extension, path-selection, or realization
  algorithm is expanded and joined before target projection.  Its complete
  contribution is target-null, finitely over budget, a typed five-source
  successor, or part of a fixed actual-prefix tree; algorithm failure
  is never an external admission condition.
\item Local germs, carrier ceilings, affordability, graph and shell moduli,
  approximation coverage, realization, continuation, and contact enter only
  through their native public graphs.  Every unfavorable target-bearing
  value is null, finitely contradictory, or a typed successor; none is a
  certificate attached to the datum.
\item The public multiplier grammar admits an energy, entropy, action, clock,
  or other Lyapunov carrier only through a finite target-blind balance,
  positivity, lower-bound, localizability, and exterior-work transcript.  No
  user-supplied or target-dependent estimate can authorize charge.
\item Effective terminalization is asserted exactly on cells whose public
  successor graph carries the checked well-founded rank of
  \cref{def:generated-noetherian-target-cell}.  Other affordable cells retain
  semantic target accountability without a halting claim.
\end{enumerate}
\end{theorem}

\begin{proof}
For \textup{(AT1)}, direct integration gives
\(x(t)=x(0)e^{-t}\); the identity
\(\frac{\dd}{\dd t}x^2=-2x^2\) is the finite Lyapunov transcript excluding
escape.  For \textup{(AT2)}, direct differentiation verifies
\(x(t)=(1-t)^{-1}\), its initial value, and its rational approximation program
on \([0,1-2^{-n}]\).  The clock carrier has mass one, while
\(x(t)\to+\infty\) as \(t\uparrow1\), so the same program verifies every
cofinal shell crossing and noncontinuation at \(t=1\).  Its construction uses
only the equation and datum and therefore has independent origin.  Statements
\textup{(AT3)}--\textup{(AT5)} follow respectively from
\cref{thm:zero-production-kernel-soundness,thm:target-aligned-passivity-support,thm:fidelity-defect-successor}.
Statement \textup{(AT6)} is the finite-prefix conclusion of
\cref{thm:continuum-composition-no-creation}; \textup{(AT7)} is
\cref{thm:support-transfer-conservation,thm:no-double-charging}; and
\textup{(AT8)} is
\cref{thm:target-aligned-public-algorithm-saturation}; and \textup{(AT9)} is
\cref{thm:public-target-aligned-premise-closure}.  Statement \textup{(AT10)}
is \cref{thm:public-carrier-admission}, and \textup{(AT11)} is
\cref{thm:generated-noetherian-target-terminalization}.  The final
output type in the two scalar cases is determined by
\cref{thm:binary-proof-carrying-terminalization}.
\end{proof}

\section{Framework conclusion}

The continuous extension is the target-aligned structural calculus generated
by a public continuum presentation.  The presentation supplies the equation,
initial and boundary data, continuation relation, target shells, and
permitted transformations.  Its target-blind multiplier grammar generates
and checks the admissible carrier strata.  Its finite expression graph is
compiled into the five source families
\(\Geom,\Caus,\Abs,\Lift,\Bdry\).  The compiler is closed under signed
composition, quotient, chart, boundary, and source-preserving limit
constructors, so every dynamically generated concentration, oscillation,
defect, tail, and graph failure retains inherited ancestry.

For a physical stratum, the framework evaluates only the closure of finite
same-origin pre-contact records generated by that presentation.  Each record
carries the complete backward frontier, unnormalized shell cocycle, target
covector, source word, and physical carrier.  The joint frontier is formed
before source projection.  Orthogonal conditional increments therefore retain
target alignment created by composition, repeated subthreshold events, or
scale and time accumulation.  A positive invariant has a least finite
pre-contact occurrence with complete five-source ancestry; zero deletes its
target-incidence edge.  No analytic object outside this reachable algebra is
an additional source.

Within an admitted carrier-tagged stratum, its public carrier is the single
budget and no other identity may be mixed into its charges.  Shell contact gives a unit
pre-contact crossing, hence a target-aligned source occurrence.  The exact
compositional accounting law charges its unnormalized cumulative current to
the carrier or routes a zero-price equality record through the same
source-preserving closure.  Cofinal contact cannot be erased by normalization:
it produces an unbounded carrier density or a target-current atom of zero
carrier mass, each with a fixed source owner and a finite contextual
occurrence.

Accordingly, target avoidance on a physical stratum is the emptiness of the
ordinary same-origin target-contact cells in the presentation-generated
five-source closure; target reachability is their nonemptiness.  This is the
continuous counterpart of Part~I saturation: target-aligned structure is
classified, jointly retained, physically charged, and evaluated entirely
inside the closed algebra generated by the public presentation.

\bibliographystyle{plainnat}
\bibliography{continuous_hamiltonian_structural_complexity}

\end{document}
